\documentclass[11pt]{article}
\usepackage[margin=1in]{geometry}
\usepackage{amsmath,amssymb,amsthm}
\usepackage{enumitem}
\usepackage[hidelinks]{hyperref}

\newtheorem{theorem}{Theorem}[section]
\newtheorem{lemma}[theorem]{Lemma}
\newtheorem{proposition}[theorem]{Proposition}
\newtheorem{corollary}[theorem]{Corollary}
\newtheorem{definition}[theorem]{Definition}
\newtheorem{assessment}[theorem]{Assessment}
\newtheorem{firewall}[theorem]{Firewall}

\title{Renewal Standard Model--Einstein Continuum Research Notes\\
Sixth Series, V1}
\author{}
\date{\today}

\begin{document}
\maketitle
\tableofcontents

\section{Context, method, and claim boundary}
\label{sec:v6v1-context}

These notes continue a constructive programme aimed at a full
Standard Model--Einstein continuum.  The preceding active manuscript had
again reached one hundred versions and had become difficult to navigate.
It has therefore been frozen, without alteration, as
\begin{center}
\texttt{Renewal\_SM\_Einstein\_Continuum\_Research\_Notes\_V100\_f5.tex},
\end{center}
with SHA--256
\begin{center}
\texttt{7072C98999EAC60F2378C35B8F137F2C40347EB1744894A14E48E4AE4B81F386}.
\end{center}
This compact V1 starts the sixth series.  It records the main mathematical
interfaces already proved, states the live obstructions, and then continues
the attack.  Later versions in this series will be appended to this file and
will replace their active predecessors.  At V100 the same freeze-and-restart
rule will be applied again.

The intended endpoint is a cutoff-independent, anomaly-free,
gauge- and diffeomorphism-covariant Schwinger functional for the coupled
Standard Model and gravity, with controlled fermion zero modes, clustering,
reflection positivity, and Osterwalder--Schrader reconstruction.  That
endpoint has not been reached.  In this series, a theorem called a
``card'' is a proved implication from explicitly listed hypotheses.  It is
not a claim that those hypotheses already hold for the full interacting
continuum measure.

\section{Major mathematical results inherited from the frozen volumes}
\label{sec:v6v1-inherited}

\subsection{Fermion topology, anomalies, and exact zero-mode retention}

The archived programme organizes the anomaly-cancelled Standard Model
fermion determinant as a determinant-line section rather than as a globally
chosen complex number.  It derives the chiral index and finite Grassmann
saturation rules and uses exact Schur/Feshbach identities to retain every
light determinant zero.  A uniformly invertible high block may be integrated
and its modulus placed in a positive reference weight; the low Schur factor
remains a terminating Grassmann polynomial.  The most recent frozen volume
adds smooth gauge-covariant Riesz charts for singular subspaces and an exact
crossing-transfer identity.  Crossing modes are moved to the retained block
before their gap closes, without division by their determinant.

\subsection{Metric normalization and local ultraviolet owners}

A metric change alters the Dirac principal symbol.  The programme inserts an
elliptic order-zero microlocal normalizer before taking a relative
determinant; the remaining relative perturbation is then of order minus one
and lies in the critical weak Schatten class in four dimensions.  Heat and
symbol expansions assign the cosmological, Einstein--Hilbert, gauge, Higgs,
nonminimal, curvature-squared, and boundary divergences to local owners.  The
marginal fourth-order logarithm is the local Wodzicki residue.  Finite
nonlocal Born amplitudes are retained rather than subtracted with the local
counterterms.

\subsection{Finite-depth Wilsonian consistency and compactness}

Exact conditional-density and transfer-product identities telescope along a
finite cutoff chain.  Under compact admissible shell spaces, compatible
finite-depth ancestors have an inverse-limit trajectory.  One positive
derivative of smoothing, a strict exponential localization margin, and an
arity gap give compactness from a strong rooted kernel norm to a weaker one.
Covariant Taylor forests and a Banach-valued
Koteck\'y--Preiss estimate provide conditional summation cards.  These are
functional-analytic implications; the full coupled action has not yet been
shown to satisfy every uniform hypothesis.

\subsection{Einstein constraints and the massless polynomial grade}

Near flat constant-mean-curvature data, a nonlinear York--Lichnerowicz chart
solves the mean-zero local constraint by the analytic implicit-function
theorem.  The global scalar constraint remains.  Elliptic constraint
substitution produces Coulomb tails, so it cannot be placed in an
exponentially local sector.  A divergence source is paid in
\(\dot H^{-1}\).  Vanishing monopole and dipole moments produce a double
difference and an \(r^{-4}\) response in three spatial dimensions, the first
such multipole response with an absolutely summable row.

The electromagnetic photon belongs to the same polynomial grade.  A
compactly supported differentiated gauge-invariant observable is an exact
lattice curl.  Gauge invariance forces componentwise neutrality and makes
its first moment antisymmetric, but does not force that moment to vanish.
The surviving coefficient is a local field-strength insertion.  Cubic
reflection symmetry or an explicit coherent matching condition removes it.
For a three-dimensional Coulomb kernel, cancellation orders \(r,s\) at the
two endpoints give decay \(R^{-1-r-s}\); the absolute row sum is finite
exactly when \(r+s>2\).

\subsection{Massive one-form covariance and rooted phase control}

After exact gauge quotienting, a normalized block-mixing bound \(q<1\)
gives a positive one-form Schur gap on the radial Higgs, \(W^\pm\), \(Z\),
massive auxiliary, and Pauli--Villars directions.  A Combes--Thomas argument
turns that gap plus finite-range or exponentially weighted locality into an
exponentially decaying inverse with a volume-uniform \(\ell^1\) row.  The
photon and gravitational constraints are excluded from this massive card.

The high Dirac trace logarithm has a covariant finite-support approximation:
range truncation of an endpoint-covariant kernel and truncation by closed
loop order preserve the Ward identity exactly.  The differentiated omitted
loop tail is geometric.  Rooted phase covariance is then controlled by the
massive inverse and by the massless moment cards.  Independent vacuum factors
cancel from normalized rooted expectations, so distant vacuum denominators
are not bounded separately.

\subsection{Cutoff exponents, owned shells, and positive block comparison}

The invariant object is the gradient--inverse-Hessian--gradient sandwich.
In four dimensions, changing from a dimensionless link angle to an
\(L^2\)-cell field coordinate costs one power of the lattice spacing in the
gradient and supplies the inverse power in the covariance.  If the inverse
is fixed in the latter coordinate, a differentiated defect of order
\(a^\beta\) has exponent \(\beta-1\).  A smooth-band estimate
\(a^b\Lambda^r\) and a high-shell estimate
\(a^{-h}\Lambda^{-s}\) are summable along geometric cutoffs precisely when
\(sb>rh\).  Assigning every high annulus to one cutoff step prevents a
nested-tail incidence loss.

A high background momentum in a translation-covariant local Dirac vertex
forces at least one adjacent fermion line high and gives one inverse-momentum
resolvent gain.  Rough backgrounds create an explicit low--low leakage.  A
good-curvature block admits a smooth local gauge, while an \(L^m\) roughness
moment and an \(L^p\) bad-field envelope interpolate to a positive high-shell
power.  For the positive normalized block measure, a centered exponential
moment of the real log weight controls its density relative to a positive
gauge reference without inserting a raw partition-function lower bound.
Uniform conditional bad-block estimates imply product domination by
sequential conditioning.

\subsection{Reflection positivity and the Gaussian all-order frontier}

Reflection positivity is formulated on finite reflected cylinder matrices
and is closed under entrywise limits.  A positive transfer factorization, or
an exact contour deformation preserving every reflected matrix element,
would supply it.  No such construction has been proved for the complete
SM--Einstein measure.

The latest companion Yang--Mills audit proves a connected Gaussian BKAR
forest formula and an all-order Gaussian tangent tree card with exponential
arity constants.  Applied to the massive-safe covariance above, the same
argument gives the all-order rooted tree card for a Gaussian SM reference.
The remaining problem is to transport that card to the positive interacting
gauge--Higgs--PV determinant-modulus measure while keeping constants
exponential in arity.

\section{Live obstruction ledger}
\label{sec:v6v1-obstructions}

\begin{theorem}[Current claim boundary]
\label{thm:v6v1-boundary}
The archived results do not yet construct the full SM--Einstein continuum.
The following gates remain open:
\begin{enumerate}
\item a regulator-level proof of local and global anomaly cancellation and
      determinant-line trivialization on every spectral-chart overlap;
\item an interacting, cutoff-uniform differentiated improvement estimate
      and rough-block probability bound satisfying the two-regime criterion;
\item an all-order transport or cluster card from the massive Gaussian
      reference to the coupled positive modulus measure;
\item complete interacting photon Ward moments and gravitational stress
      monopole--dipole matching, including all contact terms;
\item volume- and cutoff-uniform control of retained light-fermion ranks and
      spectral crossing charts;
\item a nonperturbative ultraviolet completion or replacement for the
      hypercharge Gaussian-boundary trajectory;
\item a gravitational construction accommodating the nonclosed higher-
      curvature tower and the Euclidean conformal instability; and
\item reflection positivity, clustering of the full mixed massless theory,
      cutoff removal, and Osterwalder--Schrader reconstruction.
\end{enumerate}
\end{theorem}

\begin{proof}
Each item is an explicit hypothesis or firewall at the endpoint of the
frozen fifth volume.  None is implied by the finite-cutoff algebraic
factorizations, the massive Gaussian covariance theorem, or fixed-order
perturbative cancellation.  Therefore the continuum conclusion cannot yet
be drawn.
\end{proof}

\begin{firewall}[Known route obstructions]
\label{fw:v6v1-known-route}
Under the perturbative Gaussian-boundary assumptions recorded in the
archives, the hypercharge inverse coupling reaches zero after finitely many
ultraviolet shells.  A finite curvature list is not closed by perturbative
Einstein gravity, and the real Euclidean conformal direction has an
unbounded-below kinetic action.  These results delimit the present route;
they are not a theorem excluding every ultraviolet completion compatible
with low-energy SM--Einstein physics.
\end{firewall}

\section{New attack: exact cumulant transport under a positive tilt}
\label{sec:v6v1-transport}

Let \(\gamma\) be the massive Gaussian reference from the frozen V100 card.
Let the real interacting correction be a local sum
\begin{equation}
 W=\sum_{A\in\mathcal P}W_A,
 \label{eq:v6v1-interaction-sum}
\end{equation}
and define the positive interpolation
\begin{equation}
 d\mu_t=\frac{e^{-tW}}{\mathbb E_\gamma e^{-tW}}\,d\gamma,
 \qquad 0\le t\le1.
 \label{eq:v6v1-positive-interpolation}
\end{equation}
The low Grassmann factor and the residual determinant phase are retained as
Banach-algebra-valued observables; they are not included in the positive
normalization.

\begin{lemma}[Cumulant derivative identity]
\label{lem:v6v1-cumulant-derivative}
Suppose the exponential moments needed below are finite.  For observables
\(F_1,\ldots,F_r\),
\begin{equation}
 \frac d{dt}\kappa_r^{\mu_t}(F_1,\ldots,F_r)
 =-\kappa_{r+1}^{\mu_t}(F_1,\ldots,F_r,W).
 \label{eq:v6v1-cumulant-flow}
\end{equation}
More generally,
\begin{equation}
 \frac {d^m}{dt^m}\kappa_r^{\mu_t}(F_1,\ldots,F_r)
 =(-1)^m\kappa_{r+m}^{\mu_t}
 (F_1,\ldots,F_r,\underbrace{W,\ldots,W}_{m}).
 \label{eq:v6v1-iterated-flow}
\end{equation}
\end{lemma}

\begin{proof}
Introduce sources \(z=(z_1,\ldots,z_r)\) and write
\begin{equation}
 \Phi(t,z)=\log\mathbb E_\gamma
 \exp\left(-tW+\sum_i z_iF_i\right)
 -\log\mathbb E_\gamma e^{-tW}.
 \label{eq:v6v1-generating-function}
\end{equation}
The joint cumulant is
\(\partial_{z_1}\cdots\partial_{z_r}\Phi(t,0)\).  Differentiating once in
\(t\) inserts \(-W\) into the connected derivative.  Iteration proves
\eqref{eq:v6v1-iterated-flow}.  The second logarithm is independent of
\(z\), so the normalization denominator cancels before any estimate.
\end{proof}

\begin{theorem}[Exact positive-tilt cumulant expansion]
\label{thm:v6v1-tilt-expansion}
Assume \(t\mapsto\Phi(t,z)\) is analytic for \(|t|<1+\epsilon\) and the
following series is absolutely convergent after the local sum
\eqref{eq:v6v1-interaction-sum}.  Then
\begin{align}
 \kappa_r^{\mu_1}(F_1,\ldots,F_r)
 ={}&\sum_{m=0}^{\infty}\frac{(-1)^m}{m!}
 \sum_{A_1,\ldots,A_m}
 \kappa_{r+m}^{\gamma}
 (F_1,\ldots,F_r,W_{A_1},\ldots,W_{A_m}).
 \label{eq:v6v1-exact-tilt-series}
\end{align}
For every finite \(M\), without assuming convergence of the infinite
series, the exact remainder is
\begin{align}
 \kappa_r^{\mu_1}(F_1,\ldots,F_r)
 ={}&\sum_{m=0}^{M-1}\frac{(-1)^m}{m!}
 \kappa_{r+m}^{\gamma}(F_1,\ldots,F_r,W^{\otimes m})
 \notag\\
 &+\frac{(-1)^M}{(M-1)!}\int_0^1(1-t)^{M-1}
 \kappa_{r+M}^{\mu_t}(F_1,\ldots,F_r,W^{\otimes M})\,dt.
 \label{eq:v6v1-exact-remainder}
\end{align}
\end{theorem}

\begin{proof}
Taylor-expand the left side as a function of \(t\) at zero and use
Lemma~\ref{lem:v6v1-cumulant-derivative}.  Absolute convergence permits
expanding each copy of \(W\) by \eqref{eq:v6v1-interaction-sum} and
interchanging sums.  Taylor's integral remainder gives
\eqref{eq:v6v1-exact-remainder}.
\end{proof}

\subsection{A sufficient all-order transport card}

Let \(\mathfrak T_r(x_1,\ldots,x_r)\) be the rooted spatial tree weight
formed from the V100 massive covariance kernel, with the polynomial
photon--constraint edges inserted only when their endpoint moments satisfy
the inherited summability threshold.

\begin{definition}[Summed interaction-insertion card]
\label{def:v6v1-SIIC}
The positive interaction satisfies the SIIC with constants \(K\) and
\(0\le\eta<1\) if for every \(r\ge1\) and \(m\ge0\),
\begin{align}
 &\sum_{A_1,\ldots,A_m}
 \left\|
 \kappa_{r+m}^{\gamma}
 (F_{x_1},\ldots,F_{x_r},W_{A_1},\ldots,W_{A_m})
 \right\|
 \notag\\
 &\hspace{4em}\le
 r!\,m!\,K^r\eta^m
 \mathfrak T_r(x_1,\ldots,x_r).
 \label{eq:v6v1-SIIC}
\end{align}
The norm may be a finite Grassmann/determinant-line coefficient norm.
\end{definition}

\begin{theorem}[All-order transport from the Gaussian reference]
\label{thm:v6v1-all-order-transport}
If the SIIC holds and the analytic hypothesis of
Theorem~\ref{thm:v6v1-tilt-expansion} is valid, then
\begin{equation}
 \left\|\kappa_r^{\mu_1}
 (F_{x_1},\ldots,F_{x_r})\right\|
 \le \frac{r!K^r}{1-\eta}
 \mathfrak T_r(x_1,\ldots,x_r)
 \label{eq:v6v1-transported-tree-card}
\end{equation}
for every arity.  If the covariance edge kernels have uniform
\(\ell^1\) rows, the rooted spatial sums are volume-uniform.
\end{theorem}

\begin{proof}
Take norms in \eqref{eq:v6v1-exact-tilt-series}, use
\eqref{eq:v6v1-SIIC}, and cancel the Taylor factor \(1/m!\).  The remaining
interaction series is \(\sum_{m\ge0}\eta^m=(1-\eta)^{-1}\).  Root each
tree at one observable and remove its leaves successively.  Each spatial
edge sum costs at most its uniform \(\ell^1\) row norm, proving the last
statement.
\end{proof}

\begin{proposition}[What the Gaussian BKAR theorem reduces]
\label{prop:v6v1-BKAR-reduction}
Suppose the complete derivative rows of every \(F_x\) and \(W_A\) obey
factorial-exponential bounds in the massive Gaussian coordinate.  Then the
Gaussian BKAR formula bounds each joint cumulant in
\eqref{eq:v6v1-SIIC} by a sum over labelled trees with one covariance edge
per contraction.  Consequently, proving the SIIC reduces to a rooted
polymer-tree sum over the interaction supports \(A_1,\ldots,A_m\), with an
exponential activity small enough to supply \(\eta^m\); there is no further
Gaussian pairing factorial.
\end{proposition}

\begin{proof}
The connected BKAR formula differentiates once along every tree edge before
taking a single Gaussian expectation.  Complete row--column contraction
bounds all derivative indices without a tensor-dimension factor.  The sum
of vertex degrees is \(2(r+m-1)\), so exponential derivative constants
remain exponential in \(r+m\).  The tree-degree factorial product is at
most \((r+m-1)!\); the SIIC allocates the interaction-label part to the
explicit \(m!\) and the observable part to \(r!\).  What remains is exactly
the spatial/support tree sum and its activities.
\end{proof}

\begin{firewall}[The SIIC is the live theorem, not an assumption to hide]
\label{fw:v6v1-SIIC-open}
Theorem~\ref{thm:v6v1-all-order-transport} proves that the SIIC is
sufficient.  It does not prove the SIIC for the coupled positive measure.
The next notes must derive its derivative rows and polymer activities from
the gauge, Higgs, PV, high-determinant-modulus, and counterterm factors.
\end{firewall}

\begin{firewall}[Massless edges require their moment cards]
\label{fw:v6v1-massless-tree}
A photon or constraint covariance is not an \(\ell^1\) edge before its
rooted endpoints supply enough Ward or multipole differences.  Such an edge
may enter \(\mathfrak T_r\) only with the inherited endpoint moment
certificates.  The massive Gaussian tree proof cannot assign a false mass
to it.
\end{firewall}

\begin{firewall}[Positive transport excludes the residual phase]
\label{fw:v6v1-phase-excluded}
The interpolation \eqref{eq:v6v1-positive-interpolation} applies to a real
positive modulus measure.  The determinant-line phase and retained low
Grassmann factor are observables governed by the rooted connected series.
Putting the phase into \(W\) would invalidate the order and moment estimates
used for positive comparison.
\end{firewall}

\begin{theorem}[Sixth-series V1 result]
\label{thm:v6v1-status}
The normalized positive interpolation obeys an exact cumulant flow and an
exact expansion in Gaussian joint cumulants with local interaction
insertions.  No separate normalization-denominator bound occurs.  A summed
interaction-insertion estimate with activity \(\eta<1\) transports the
all-order Gaussian rooted tree card to the interacting positive measure with
only the factor \((1-\eta)^{-1}\).  The companion Gaussian BKAR theorem
removes pairing-factorial growth and reduces the live problem to a local
polymer-tree activity bound for the complete positive interaction.  That
SIIC, the massless endpoint moments, and the residual phase estimates remain
to be proved for the full coupled measure.
\end{theorem}

\begin{proof}
Use Lemma~\ref{lem:v6v1-cumulant-derivative},
Theorems \ref{thm:v6v1-tilt-expansion} and
\ref{thm:v6v1-all-order-transport}, and
Proposition~\ref{prop:v6v1-BKAR-reduction}, subject to the three firewalls.
\end{proof}

\section{Next acquisition order}
\label{sec:v6v1-next}

V2 will decompose the complete positive interaction into local
gauge--Higgs, PV, high-determinant-modulus, gauge-fixing, and counterterm
polymers.  It will prove a sharp SIIC for the bounded massive-safe chart if
possible, and will isolate any factorial derivative or spatial activity
that prevents \(\eta<1\).  The determinant phase, photon, gravitational
constraint, and spectral crossing factors will remain explicit rather than
being absorbed into that positive massive card.

\clearpage
\section{V2: adjacent-cutoff all-order transport}
\label{sec:v6v2}

V1 formulated a sufficient all-order transport card for one positive tilt.
Applying it once to the complete interacting action would demand that the
entire non-Gaussian interaction be small.  The Wilsonian identities in the
frozen volumes suggest the correct upstream organization: transport only
the adjacent-cutoff increment at step \(n\), prove a small local activity
\(\eta_n\), and use summability of \((\eta_n)\).

\subsection{Local decomposition of one positive increment}

Let \(\mu_n\) be the normalized positive modulus measure after cutoff step
\(n\).  Write the next conditional density as
\begin{equation}
 d\mu_{n+1}=
 \frac{e^{-W_n}}{\mathbb E_{\mu_n}e^{-W_n}},d\mu_n,
 \qquad
 W_n=\sum_{A\in\mathcal P_n}W_{n,A}.
 \label{eq:v6v2-adjacent-tilt}
\end{equation}
On the massive-safe good-field chart, the local terms are assigned to
\begin{equation}
 W_{n,A}=W_{n,A}^{\rm gauge/Higgs}
 +W_{n,A}^{\rm PV}
 +W_{n,A}^{\rm high\,mod}
 +W_{n,A}^{\rm gf/ghost}
 +W_{n,A}^{\rm ct}.
 \label{eq:v6v2-local-owners}
\end{equation}
The determinant-line phase, retained low Grassmann factor, photon grade,
and gravitational constraint grade remain observables and are not included
in this positive massive increment.

For a local functional \(U_A\), define its complete derivative-row activity
in the Gaussian normal coordinate by
\begin{equation}
 z(U_A):=sup_{d\ge0}\frac{1}{d!L^d}
 \sup_\phi\|\nabla^dU_A(\phi)\|_{\rm row},
 \label{eq:v6v2-derivative-activity}
\end{equation}
where the supremum is restricted to the stated bounded chart and \(L\) is
a fixed derivative scale.

\begin{lemma}[Local derivative rows of the positive owners]
\label{lem:v6v2-local-rows}
On a bounded massive-safe chart:
\begin{enumerate}
\item a local polynomial owner of degree \(k\) has finite activity
      \eqref{eq:v6v2-derivative-activity};
\item a plaquette or link exponential with analytic chart radius \(R>L\)
      has activity bounded by its supremum on the radius-\(R\) complex
      chart times a constant depending on \(L/R\);
\item if a local high fermion block \(E(\phi)\) has
      \(\|E^{-1}\|\le m_H^{-1}\) throughout that complex chart, then
      \(\log|\det E|\) has derivative rows bounded by
\begin{equation}
 \|\nabla^d\log|\det E|\|_{\rm row}
 \le d!\,N_E\,C_E^d,
 \label{eq:v6v2-logdet-rows}
\end{equation}
where \(N_E\) is the local block rank and \(C_E\) depends on the analytic
norm of \(E\) and on \(m_H^{-1}\).
\end{enumerate}
\end{lemma}

\begin{proof}
Polynomial derivatives vanish above degree \(k\), proving the first item.
Cauchy's integral estimate in each direction gives
\(\|\nabla^dU\|\le d!R^{-d}\sup_{\|z\|<R}|U(z)|\); complete row norms are
obtained by polarization with an exponential change of the constant.
For the last item, choose a continuous logarithm of the high determinant on
the gap chart.  Jacobi's formula begins with
\(d\log\det E=\operatorname{Tr}(E^{-1}dE)\).  Repeated differentiation
produces ordered products of derivatives of \(E\) separated by factors of
\(E^{-1}\).  Cauchy's estimate applied directly to the analytic scalar
\(\log\det E\), together with
\(|\operatorname{Tr}X|\le N_E\|X\|\), yields
\eqref{eq:v6v2-logdet-rows}.  Taking the real part gives the modulus.
\end{proof}

The rank factor is harmless only after it is assigned to a local block
activity.  A global determinant rank would reintroduce a volume factor.

\subsection{A rooted interaction-tree summation lemma}

Let \(C(x,y)\ge0\) be a covariance majorant with
\begin{equation}
 C(x,y)\le C_0e^{-\mu d(x,y)},
 \qquad
 \sup_x\sum_yC(x,y)\le C_1.
 \label{eq:v6v2-covariance-majorant}
\end{equation}
Every interaction support \(A\) has diameter at most \(R_0\), an anchor
\(a(A)\), and activity \(z_A\ge0\).  Put
\begin{equation}
 \zeta:=\sup_x\sum_{A:a(A)=x}z_Ae^{\tau|A|}.
 \label{eq:v6v2-local-zeta}
\end{equation}

\begin{lemma}[Summation of inserted interaction vertices]
\label{lem:v6v2-insertion-tree-sum}
Fix \(0<\mu'<\mu\).  There is
\(K_*=K_*(C_0,C_1,\mu-\mu',R_0,\tau)\) such that, for every \(r\ge1\) and
\(m\ge0\), the sum of all Gaussian BKAR tree majorants with fixed observable
roots \(x_1,\ldots,x_r\) and \(m\) labelled interaction vertices satisfies
\begin{align}
 &\sum_{A_1,\ldots,A_m}
 \sum_{T\in\mathfrak T_{r+m}}
 \prod_{\{i,j\}\in E(T)}C(v_i,v_j)
 \prod_{j=1}^m z_{A_j}
 \notag\\
 &\hspace{3em}\le
 r!\,m!\,K_*^{r}(K_*\zeta)^m
 \mathfrak T_r^{\mu'}(x_1,\ldots,x_r),
 \label{eq:v6v2-insertion-tree-bound}
\end{align}
where \(v_i=x_i\) for observable vertices, \(v_{r+j}=a(A_j)\), and
\(\mathfrak T_r^{\mu'}\) is the sum of spanning-tree weights built from
\(e^{-\mu'd(x,y)}\).  For \(r=1\), set
\(\mathfrak T_1^{\mu'}=1\).
\end{lemma}

\begin{proof}
Replace every covariance edge by \(C_0e^{-\mu d}\).  Contract each maximal
connected subtree consisting only of interaction vertices toward the first
observable vertex it meets.  Summing an interaction anchor along the parent
edge costs at most \(C_1\zeta\); every additional incident edge is retained
with exponent \(\mu'\), while the margin \(\mu-\mu'\) pays the convolution.
The bounded support radius changes this cost by at most
\(e^{\mu R_0}\).  Repeating removes all interaction vertices and leaves a
connected graph on the observable roots.  Choose one of its spanning trees
and bound surplus edges by \(C_0\).

It remains to count labelled contraction orders and tree choices.  Cayley's
bound and Stirling's inequality give
\begin{equation}
 (r+m)^{r+m-2}\le r!\,m!\,K_2^{r+m}
 \max\{1,r^{r-2}\}
 \label{eq:v6v2-Cayley-factorial}
\end{equation}
for a numerical \(K_2\); the finitely many small \(r,m\) are absorbed by
enlarging it.  The remaining \(r^{r-2}\) choices are precisely dominated
by the observable spanning-tree sum.  Collecting the convolution and
counting constants gives \eqref{eq:v6v2-insertion-tree-bound}.
\end{proof}

\subsection{One-step and all-scale transport}

\begin{theorem}[One adjacent-cutoff SIIC]
\label{thm:v6v2-one-step-SIIC}
Assume at step \(n\):
\begin{enumerate}
\item the massive reference covariance obeys
      \eqref{eq:v6v2-covariance-majorant};
\item all rooted observables and local owners obey complete derivative-row
      bounds with common exponential scale \(L\); and
\item the owner activities satisfy \eqref{eq:v6v2-local-zeta} with
      \(\zeta=\zeta_n\).
\end{enumerate}
Then the summed interaction-insertion card of V1 holds with
\begin{equation}
 \eta_n=K_*\zeta_n.
 \label{eq:v6v2-eta-n}
\end{equation}
In particular, if \(\eta_n<1\), the all-order rooted tree card transports
from \(\mu_n\) to \(\mu_{n+1}\) with loss at most
\((1-\eta_n)^{-1}\), after a fixed enlargement of the exponential arity
constant.
\end{theorem}

\begin{proof}
Apply the connected Gaussian/BKAR derivative formula to the observable and
interaction vertices.  The complete derivative-row estimates contract
along tree edges without a tensor-dimension factor.  The product of vertex
degree factorials is bounded by \((r+m-1)!\).  Since
\begin{equation}
 (r+m-1)!\le(r+m)!\le2^{r+m}r!m!,
 \label{eq:v6v2-factorial-split}
\end{equation}
the extra combinatorics is absorbed into exponential constants.  Lemma
\ref{lem:v6v2-insertion-tree-sum} supplies the spatial and support sums,
giving the V1 SIIC with \(\eta_n=K_*\zeta_n\).  The V1 all-order transport
theorem gives the last statement.
\end{proof}

The same proof applies when \(\mu_n\) is already non-Gaussian if it carries
the identical mixed insertion tree card uniformly along the interpolation.
Thus the induction datum is the tree card itself, not Gaussianity at every
step.

\begin{theorem}[Summable product of cutoff transports]
\label{thm:v6v2-all-scale-product}
Suppose every adjacent step has \(0\le\eta_n\le1/2\) and
\begin{equation}
 \sum_{n=0}^{\infty}\eta_n<\infty.
 \label{eq:v6v2-summable-eta}
\end{equation}
Then the cumulative all-order transport loss is finite:
\begin{equation}
 \prod_{n=0}^{\infty}(1-\eta_n)^{-1}
 \le\exp\left(2\sum_{n=0}^{\infty}\eta_n\right)<\infty.
 \label{eq:v6v2-finite-product}
\end{equation}
Consequently, if the finite-depth cylinder expectations are compatible and
the inherited compactness hypotheses hold, the limiting cylinder
functional retains the all-order rooted tree bound.
\end{theorem}

\begin{proof}
For \(0\le x\le1/2\), \(-\log(1-x)\le2x\).  Sum this inequality and
exponentiate to obtain \eqref{eq:v6v2-finite-product}.  Apply the one-step
transport successively.  The inherited finite-depth compatibility and
compactness then permit passage to the cylinder limit; the uniform tree
bound is closed under coefficientwise convergence.
\end{proof}

\begin{corollary}[Connection with the two-regime cutoff exponent]
\label{cor:v6v2-two-regime-activity}
If the complete adjacent positive interaction satisfies
\begin{equation}
 \zeta_n\le Ca_n^\varepsilon,
 \qquad \varepsilon>0,
 \label{eq:v6v2-zeta-power}
\end{equation}
then \(\sum_n\eta_n<\infty\) along a geometric cutoff sequence.  The V95
smooth/high-shell certificate supplies such an \(\varepsilon\) whenever
its interacting hypotheses hold and \(sb>rh\).
\end{corollary}

\begin{proof}
Use \eqref{eq:v6v2-eta-n} and sum the geometric sequence.  The last
statement is the balanced exponent theorem in the frozen fifth volume.
\end{proof}

\subsection{What has and has not been reduced}

\begin{firewall}[The full action is not a small adjacent increment]
\label{fw:v6v2-full-versus-increment}
The activity \(\zeta_n\) belongs to the difference between two consecutive
normalized effective actions, with all relevant local owners matched.  The
bare Higgs quartic, gauge action, or complete high determinant need not be
small.  Inserting the full action into
\eqref{eq:v6v2-zeta-power} would change the theorem and is not justified.
\end{firewall}

\begin{firewall}[Bad fields require their own polymer activity]
\label{fw:v6v2-bad-field-activity}
Lemma~\ref{lem:v6v2-local-rows} holds on the bounded massive-safe chart.
The complement must be represented by bad-block polymers whose activities
are controlled by the V97--V98 conditional probability comparison.  A
one-block probability estimate is insufficient unless it is uniform under
conditioning and yields the product domination used in the polymer tree
sum.
\end{firewall}

\begin{firewall}[Relevant local flow is matched, not summed as irrelevant]
\label{fw:v6v2-relevant-flow}
Relevant and marginal coefficients generally do not satisfy
\eqref{eq:v6v2-zeta-power} before their running couplings and counterterms
are matched.  Their finite-dimensional renormalization flow must be solved
first.  Only the matched residual belongs to the summable interaction
activity.
\end{firewall}

\begin{firewall}[Gravity and the massless grades remain separate]
\label{fw:v6v2-gravity-separate}
The Euclidean conformal direction is not a positive massive Gaussian
coordinate, and photon or constraint propagators do not satisfy
\eqref{eq:v6v2-covariance-majorant}.  They may enter the tree only through
the inherited Ward/multipole endpoint cards or through a separately
justified contour and reflection-positive construction.
\end{firewall}

\begin{theorem}[Sixth-series V2 result]
\label{thm:v6v2-status}
On a bounded massive-safe chart, the local gauge--Higgs, PV, high-determinant
modulus, gauge-fixing/ghost, and counterterm owners have factorial-
exponential derivative rows under an analytic high-block gap.  A rooted
interaction-tree summation converts their local activity \(\zeta_n\) into
the V1 all-order insertion parameter \(\eta_n=K_*\zeta_n\).  Transporting
one adjacent-cutoff increment at a time has a finite cumulative loss when
\(\sum_n\eta_n<\infty\); a positive two-regime cutoff exponent is sufficient.
The remaining acquisition is a cutoff-uniform proof of the matched activity
bound, including conditionally dominated bad-field polymers and the finite
relevant flow.
\end{theorem}

\begin{proof}
Use Lemmas \ref{lem:v6v2-local-rows} and
\ref{lem:v6v2-insertion-tree-sum}, Theorems
\ref{thm:v6v2-one-step-SIIC} and
\ref{thm:v6v2-all-scale-product}, and
Corollary~\ref{cor:v6v2-two-regime-activity}, subject to the four firewalls.
\end{proof}

V3 will write the matched local activity \(\zeta_n\) as an explicit sum of
smooth-band, owned high-shell, bad-block, and relevant-flow contributions.
It will derive the exact exponent inequalities needed for each contribution
to be summable and identify the first one not supplied by the present
SM--Einstein estimates.

\clearpage
\section{V3: the complete adjacent-activity exponent ledger}
\label{sec:v6v3}

V2 reduces all-order transport to the summability of the local adjacent
activity \(\zeta_n\).  This note expands that activity into its four
analytic sources and solves the common crossover inequalities.

\subsection{Four positive contributions}

Let \(a_n=a_0\rho^n\), \(0<\rho<1\), and choose
\begin{equation}
 \Lambda_n=a_n^{-\eta},\qquad\eta>0.
 \label{eq:v6v3-crossover}
\end{equation}
After all relevant local coefficients at step \(n\) have been matched,
decompose the absolute local activity as
\begin{equation}
 \zeta_n\le
 \zeta_n^{\rm sm}+\zeta_n^{\rm hi}
 +\zeta_n^{\rm bad}+\zeta_n^{\rm rel}.
 \label{eq:v6v3-four-activities}
\end{equation}
Assume the separate estimates
\begin{align}
 \zeta_n^{\rm sm}&\le C a_n^b\Lambda_n^r,
 & b,r&>0,
 \label{eq:v6v3-smooth-activity}\\
 \zeta_n^{\rm hi}&\le C a_n^{-h}\Lambda_n^{-s},
 & h&\ge0,\ s>0,
 \label{eq:v6v3-high-activity}\\
 \mathbb P_n({\rm bad};\Lambda_n)&\le
 C a_n^\pi\Lambda_n^{-\tau},
 & \tau&>0,
 \label{eq:v6v3-bad-probability}\\
 \|X_n^{\rm bad}\|_{L^p}&\le Ca_n^{-\chi},
 & p&>1,
 \label{eq:v6v3-bad-envelope}\\
 \zeta_n^{\rm rel}&\le Ca_n^\delta,
 & \delta&>0.
 \label{eq:v6v3-relevant-residual}
\end{align}
The exponent \(\pi\) may be zero; a positive value records extra cutoff
rarity beyond the high-momentum tail.  Put \(p'=p/(p-1)\).

\begin{lemma}[Bad-block activity exponent]
\label{lem:v6v3-bad-exponent}
Under \eqref{eq:v6v3-bad-probability}--
\eqref{eq:v6v3-bad-envelope},
\begin{equation}
 \zeta_n^{\rm bad}
 \le C a_n^{\varepsilon_{\rm bad}(\eta)},
 \qquad
 \varepsilon_{\rm bad}(\eta)
 :=\frac\pi{p'}-\chi+\frac{\eta\tau}{p'}.
 \label{eq:v6v3-bad-activity}
\end{equation}
\end{lemma}

\begin{proof}
H\"older's inequality gives
\[
 \mathbb E({\bf1}_{\rm bad}|X_n^{\rm bad}|)
 \le \mathbb P_n({\rm bad};\Lambda_n)^{1/p'}
      \|X_n^{\rm bad}\|_{L^p}.
\]
Insert \eqref{eq:v6v3-bad-probability},
\eqref{eq:v6v3-bad-envelope}, and
\(\Lambda_n=a_n^{-\eta}\).
\end{proof}

\subsection{The exact crossover window}

Define the four affine exponents
\begin{align}
 \varepsilon_{\rm sm}(\eta)&=b-r\eta,
 &\varepsilon_{\rm hi}(\eta)&=s\eta-h,
 \notag\\
 \varepsilon_{\rm bad}(\eta)&=pi/p'-\chi+\tau\eta/p',
 &\varepsilon_{\rm rel}&=\delta.
 \label{eq:v6v3-four-exponents}
\end{align}

\begin{theorem}[Complete exponent feasibility criterion]
\label{thm:v6v3-feasibility}
There exists a crossover \(\eta>0\) for which every contribution in
\eqref{eq:v6v3-four-activities} is bounded by a positive power of \(a_n\)
if and only if \(\delta>0\) and
\begin{equation}
 \max\left\{
 0,\frac hs,\frac{p'\chi-\pi}{\tau}
 \right\}<\frac br.
 \label{eq:v6v3-window-condition}
\end{equation}
For every \(\eta\) in this open interval,
\begin{equation}
 \zeta_n\le4C a_n^{\varepsilon(\eta)},
 \qquad
 \varepsilon(\eta):=min\{
 b-r\eta,s\eta-h,
 \pi/p'-\chi+\tau\eta/p',\delta\}>0.
 \label{eq:v6v3-total-epsilon}
\end{equation}
Consequently the V2 all-order transport product is finite.
\end{theorem}

\begin{proof}
After substitution of \eqref{eq:v6v3-crossover}, the smooth exponent is
positive exactly when \(\eta<b/r\), the high exponent exactly when
\(\eta>h/s\), and the bad exponent exactly when
\(\eta>(p'\chi-\pi)/\tau\).  Include \(\eta>0\).  Their intersection is
nonempty exactly under \eqref{eq:v6v3-window-condition}; the relevant
residual separately requires \(\delta>0\).  Taking the minimum proves
\eqref{eq:v6v3-total-epsilon}.  A geometric cutoff makes
\(\sum_na_n^{\varepsilon(\eta)}\) finite, so V2 applies.
\end{proof}

\begin{proposition}[Optimal crossover is a finite linear programme]
\label{prop:v6v3-optimal-crossover}
On the feasible interval, the best certified power is
\begin{equation}
 \varepsilon_*=max_{\eta}
 \min\{b-r\eta,s\eta-h,
 \pi/p'-\chi+\tau\eta/p',\delta\}.
 \label{eq:v6v3-linear-programme}
\end{equation}
It is attained at an intersection of two of the three nonconstant affine
functions, or on the plateau \(\varepsilon=\delta\).  In the absence of an
active bad-block constraint, it reduces to
\begin{equation}
 \eta_*=(b+h)/(r+s),
 \qquad
 \varepsilon_*=min\{(sb-rh)/(r+s),\delta\}.
 \label{eq:v6v3-old-optimum}
\end{equation}
\end{proposition}

\begin{proof}
The minimum of finitely many affine functions is continuous and piecewise
affine.  An interior strict maximum can occur only where at least two active
affine functions meet, unless the constant function \(\delta\) is active.
Equating the smooth and high functions gives
\eqref{eq:v6v3-old-optimum} when the bad function lies above their
intersection.
\end{proof}

\subsection{Owner-by-owner interpretation}

The four activities in \eqref{eq:v6v3-four-activities} have distinct
mathematical owners.

\begin{enumerate}
\item The smooth term contains the complete differentiated Symanzik
      remainder after the gauge, Higgs, fermion, measure, ghost, and stress
      contact packet is matched.  In the four-dimensional \(L^2\)-cell
      convention, a second-order improved defect gives the candidate
      value \(b=1\).
\item The high term contains only annuli owned by this cutoff step.  The
      one-high-fermion-line identity gives the candidate \(s=1\) on a
      translation-covariant good background; global inverse and coordinate
      factors determine \(h\).
\item The bad term is paid by a conditional curvature-block probability
      and an unrestricted \(L^p\) envelope.  The normalized positive
      comparison transfers a reference plaquette tail only after its
      centered log-weight moment is bounded.
\item The relevant term is the error left after solving the finite
      running-coupling and counterterm matching equations.  A marginal
      coefficient of size \(1/n\) is not represented by any \(a_n^\delta\)
      and must be incorporated into the running trajectory before this
      ledger is used.
\end{enumerate}

\begin{lemma}[Borderline logarithmic summation]
\label{lem:v6v3-log-borderline}
If one contribution has no positive lattice power but obeys
\begin{equation}
 \zeta_n^{(0)}\le C(1+n)^{-q},
 \label{eq:v6v3-log-decay}
\end{equation}
then it is summable exactly when \(q>1\).  A marginal \(q=1\) flow produces
a logarithmically divergent all-order transport loss unless it is absorbed
into the reference trajectory or cancelled.
\end{lemma}

\begin{proof}
This is the convergence criterion for the \(p\)-series.  The V2 transport
loss contains \(\sum_n\eta_n\), so no additional cancellation occurs after
absolute activities have been taken.
\end{proof}

\subsection{Current evidence placed in the ledger}

\begin{assessment}[Which inequalities are presently supplied]
\label{ass:v6v3-current-evidence}
The frozen programme supplies the following implications, but not their
complete interacting hypotheses:
\begin{enumerate}
\item the free fully improved smooth symbol has the favorable order behind
      \(b=1\), while the complete interacting differentiated packet is
      unproved;
\item the exact high-transfer resolvent identity supplies one power
      \(s=1\) on translation-covariant good fields, while rough leakage and
      uniform global inverse factors are unproved;
\item centered-log comparison turns a uniform conditional reference tail
      into the bad-block exponent, while the coupled reference tail and
      light-singular-value moments are unproved; and
\item finite-depth local matching identities exist, while the ultraviolet
      hypercharge trajectory and the unbounded gravitational coupling tower
      prevent a closed finite relevant-flow card for the stated route.
\end{enumerate}
Thus no common numerical \(\eta\) has yet been certified for the full
SM--Einstein measure.
\end{assessment}

\begin{firewall}[Absolute owner sums cannot borrow a cancellation]
\label{fw:v6v3-no-cross-cancellation}
The activity \(\zeta_n\) is an absolute derivative-row norm.  A negative
smooth contribution cannot cancel a positive bad-block activity after this
norm is taken.  Any Ward or counterterm cancellation must be proved before
the owner is entered into \eqref{eq:v6v3-four-activities}.
\end{firewall}

\begin{firewall}[Expectation bounds do not automatically give SIIC]
\label{fw:v6v3-expectation-versus-SIIC}
Lemma~\ref{lem:v6v3-bad-exponent} is a one-insertion estimate.  The V2
interaction-tree sum requires the corresponding conditional product or
polymer activity at all insertion numbers.  The sequential domination card
from the frozen V98 supplies that implication only under a uniform
conditional bad-block bound.
\end{firewall}

\begin{firewall}[A feasible exponent does not solve the relevant flow]
\label{fw:v6v3-feasible-not-flow}
Theorem~\ref{thm:v6v3-feasibility} is a sharp bookkeeping theorem.  Choosing
formal values satisfying \eqref{eq:v6v3-window-condition} does not prove
that the Standard Model and Einstein beta functions admit the required
matched all-scale trajectory.
\end{firewall}

\begin{theorem}[Sixth-series V3 result]
\label{thm:v6v3-status}
The complete adjacent positive activity is summable if the smooth, owned
high-shell, bad-block, and matched relevant residual exponents admit the
single crossover window
\[
 \max\{0,h/s,(p'\chi-\pi)/\tau\}<b/r
\]
and \(\delta>0\).  The optimal gain is the solution of the finite linear
programme \eqref{eq:v6v3-linear-programme}; any zero-power contribution
must instead decay faster than \(1/n\).  Present SM--Einstein estimates do
not yet provide a common numerical point in this window because the
interacting improved packet, rough-field conditional tail, and closed
relevant trajectory remain open.
\end{theorem}

\begin{proof}
Use Lemma~\ref{lem:v6v3-bad-exponent},
Theorem~\ref{thm:v6v3-feasibility},
Proposition~\ref{prop:v6v3-optimal-crossover}, and
Lemma~\ref{lem:v6v3-log-borderline}, subject to the assessment and three
firewalls.
\end{proof}

V4 will attack the first analytic entry, the complete differentiated
smooth-band packet.  It will formulate the lattice background Ward and
translation identities for the sum of vertex, self-energy, measure, ghost,
and counterterm insertions, and determine the exact dimension-five
coefficients whose vanishing is equivalent to \(b=1\).

\clearpage
\section{V4: a finite improvement test for the complete smooth packet}
\label{sec:v6v4}

The frozen fifth volume computed the free Wilson smooth-band order and
identified the need for the complete dimension-five improvement packet.
This note does not repeat that symbol expansion.  It supplies the missing
finite-dimensional equivalence: after all Ward contacts are assembled, a
second-order smooth defect is equivalent to the vanishing, to one higher
order, of finitely many coefficients in the dimension-five local quotient.

\subsection{The Ward-complete differentiated packet}

Let \(\Gamma_a\) be the finite-cutoff effective functional before the
residual determinant phase is estimated.  Its complete background gauge
vertex and stress vertex include derivatives of the action, regulator,
measure, gauge fixing, ghosts, PV factors, and local counterterms.  In
momentum notation the background gauge identity has the form
\begin{equation}
 \widehat k_\mu\Gamma_{a}^{\mu}(p+k,p)
 =S_a^{-1}(p+k)Q-QS_a^{-1}(p)+\mathcal A_a(p,k),
 \label{eq:v6v4-gauge-Ward}
\end{equation}
where \(\mathcal A_a\) is the complete regulated anomaly defect.  The
translation identity has the schematic exact form
\begin{equation}
 \widehat k_\mu\Gamma_{T,a}^{\mu\nu}(p+k,p)
 =\mathcal C_a^\nu[S_a^{-1}](p,k)+\mathcal M_a^\nu(p,k),
 \label{eq:v6v4-translation-Ward}
\end{equation}
where \(\mathcal C_a^\nu\) denotes the two external-leg contact terms and
\(\mathcal M_a^\nu\) contains measure, gauge-fixing, and regulator contacts.
The exact formula depends on the chosen lattice difference and field
representations; all its terms are retained in the object below.

\begin{definition}[Complete smooth packet]
\label{def:v6v4-complete-packet}
Fix a physical momentum band \(|p_i|,|k_j|\le P\) and a bounded smooth
background chart.  The complete defect \(\mathscr D_a\) is the direct sum
of the lattice-minus-target differences of:
\begin{enumerate}
\item the inverse propagators and all background gauge, Higgs/Yukawa, and
      stress vertices needed by the differentiated observable;
\item every contact on the right sides of
      \eqref{eq:v6v4-gauge-Ward} and
      \eqref{eq:v6v4-translation-Ward}; and
\item the local measure, ghost, PV, and counterterm insertions required to
      make those identities exact.
\end{enumerate}
Its norm \(\|\cdot\|_{P,k}\) controls derivatives through a fixed order
\(k\) on that band.
\end{definition}

\begin{lemma}[Ward defects are part of the packet]
\label{lem:v6v4-Ward-packet}
If the target continuum identities have vanishing gauge anomaly, then
\begin{equation}
 \|\mathcal A_a\|_{P,k}
 +\|\mathcal M_a-\mathcal M_{\rm target}\|_{P,k}
 \le C\|\mathscr D_a\|_{P,k+1}.
 \label{eq:v6v4-Ward-defect-bound}
\end{equation}
In particular, an \(O(a^2)\) claim for a vertex with an omitted contact
does not imply an \(O(a^2)\) bound for \(\mathscr D_a\).
\end{lemma}

\begin{proof}
Subtract the target Ward identities from
\eqref{eq:v6v4-gauge-Ward}--\eqref{eq:v6v4-translation-Ward}.  Every term
on the resulting right sides is a component of
Definition~\ref{def:v6v4-complete-packet}.  Multiplication by the bounded
lattice momentum and taking one extra derivative are continuous on the
fixed band, giving \eqref{eq:v6v4-Ward-defect-bound}.
\end{proof}

\subsection{The dimension-five local quotient}

Let \(\mathcal L_5\) be the finite-dimensional vector space of local
canonical-dimension-five operators allowed by the exact lattice gauge,
flavour, discrete, and background symmetries on the chosen chart.  Let
\(\mathcal N_5\) be the subspace generated by total differences, target
equations of motion when on-shell improvement is intended, BRST-exact
terms in a proved BRST formulation, and local field redefinitions already
fixed by normalization.  Define
\begin{equation}
 \mathcal Q_5:=\mathcal L_5/\mathcal N_5,
 \qquad N_5:=\dim\mathcal Q_5<\infty.
 \label{eq:v6v4-Q5}
\end{equation}

For orientation, the fermion sector before quotienting contains the Wilson
kinetic insertion \(\bar\psi D^2\psi\), Pauli/clover insertions
\(\bar\psi\sigma^{\mu\nu}F_{\mu\nu}\psi\), mass-dressed kinetic and scalar
terms, Yukawa/Higgs-dressed terms allowed by the exact chiral/flavour
symmetries, and curvature-dressed fermion bilinears on a metric background.
Operators such as the lepton-number-violating Weinberg term are included
only if the regulator and target symmetry permit them.  The quotient,
rather than this illustrative list, fixes the actual independent basis.

Choose representatives \(O_{5,1},\ldots,O_{5,N_5}\) of a basis of
\(\mathcal Q_5\), and let \(J_{5,\ell}\) be their contributions to the
complete packet.

\begin{definition}[Uniform Symanzik packet expansion]
\label{def:v6v4-Symanzik-expansion}
The complete packet has a uniform dimension-five expansion if
\begin{equation}
 \mathscr D_a
 =a\sum_{\ell=1}^{N_5}c_\ell(a)J_{5,\ell}+a^2R_{6,a},
 \qquad
 \sup_{a<a_0}\|R_{6,a}\|_{P,k}<\infty.
 \label{eq:v6v4-packet-expansion}
\end{equation}
The finite nonlocal continuum amplitude is subtracted in defining
\(\mathscr D_a\); it is not included in \(R_{6,a}\).
\end{definition}

\subsection{Dual extraction of every improvement coefficient}

\begin{theorem}[Finite improvement test]
\label{thm:v6v4-finite-improvement-test}
There exist \(N_5\) bounded linear test functionals
\(\ell_1,\ldots,\ell_{N_5}\) on the smooth packet space such that the
matrix
\begin{equation}
 M_{j\ell}:=\ell_j(J_{5,\ell})
 \label{eq:v6v4-matching-matrix}
\end{equation}
is invertible.  Under
\eqref{eq:v6v4-packet-expansion}, the following are equivalent:
\begin{align}
 \|\mathscr D_a\|_{P,k}&\le Ca^2,
 \label{eq:v6v4-second-order-packet}\\
 |c_\ell(a)|&\le C_\ell a
 \quad(1\le\ell\le N_5).
 \label{eq:v6v4-coefficient-order}
\end{align}
If the coefficients have limits as \(a\to0\), second-order improvement is
therefore equivalent to the finite matching conditions
\(c_\ell(0)=0\), together with the quantitative remainder
\eqref{eq:v6v4-coefficient-order}.
\end{theorem}

\begin{proof}
Linear independence of the classes \(J_{5,\ell}\) in the finite-dimensional
quotient permits a dual basis on their span.  Extend it to bounded
functionals on the ambient normed packet space by Hahn--Banach; in this
basis \(M=I\), while any other independent test family gives an invertible
matrix.

If \eqref{eq:v6v4-coefficient-order} holds, substitute it in
\eqref{eq:v6v4-packet-expansion} and use finite dimensionality to obtain
\eqref{eq:v6v4-second-order-packet}.  Conversely, apply every \(\ell_j\) to
\eqref{eq:v6v4-packet-expansion} and divide by \(a\):
\begin{equation}
 \sum_\ell M_{j\ell}c_\ell(a)
 =a^{-1}\ell_j(\mathscr D_a)-a\ell_j(R_{6,a})=O(a).
 \label{eq:v6v4-extracted-coefficients}
\end{equation}
Multiplication by \(M^{-1}\) proves
\eqref{eq:v6v4-coefficient-order}.
\end{proof}

\begin{corollary}[The smooth activity exponent]
\label{cor:v6v4-smooth-b}
In four-dimensional \(L^2\)-cell coordinates, if the differentiated
version of \eqref{eq:v6v4-second-order-packet} holds uniformly on the fixed
physical band, then the smooth adjacent activity has the V3 exponent
\begin{equation}
 b=1.
 \label{eq:v6v4-b-equals-one}
\end{equation}
Conversely, a nonzero limiting coefficient in \(\mathcal Q_5\) prevents
this route from proving \(b=1\).
\end{corollary}

\begin{proof}
The packet is \(O(a^2)\) in the dimensionless link convention.  The
four-dimensional \(L^2\)-cell derivative costs one power of \(a\), as
proved in the frozen normalization theorem, leaving \(O(a)\).  If some
\(c_\ell(0)\ne0\), Theorem~\ref{thm:v6v4-finite-improvement-test} supplies
a bounded test on which the packet is order \(a\), so the canonical
derivative is only order one.
\end{proof}

\subsection{Tree-level orientation and interacting matching}

For the conventional dimensionful Wilson operator, the free expansion
contains
\begin{equation}
 D_W=D_{\rm cont}+\frac{ra}{2}D^2+O(a^2P^3)
 \label{eq:v6v4-free-Wilson-orientation}
\end{equation}
on a fixed band.  Using the squared Dirac identity, the on-shell class of
this term contains the Pauli/clover operator.  The familiar tree-level
clover choice cancels that class in the corresponding convention.  This
orients one coordinate of \(\mathcal Q_5\); it does not determine the
interacting matching matrix for the entire Standard Model packet.

\begin{proposition}[Ward-compatible matching system]
\label{prop:v6v4-Ward-compatible-matching}
Assume the regulated anomaly defect vanishes and the counterterm family
contains one adjustable coefficient for every nonredundant class in
\(\mathcal Q_5\).  Then the improvement conditions can be written as the
finite system
\begin{equation}
 \ell_j(\mathscr D_a)=O(a^2),
 \qquad 1\le j\le N_5,
 \label{eq:v6v4-finite-matching-system}
\end{equation}
with the same coefficients used simultaneously in the propagator, vertex,
stress, measure, ghost, and contact components.  Solving separate systems
for those components can violate the Ward identities and does not prove
complete improvement.
\end{proposition}

\begin{proof}
Theorem~\ref{thm:v6v4-finite-improvement-test} gives the finite system.
Lemma~\ref{lem:v6v4-Ward-packet} shows that Ward defects are linear
components of the same packet.  A single local counterterm functional
differentiates into all components with fixed relative coefficients; using
different coefficients would no longer be the derivative of that
functional and need not satisfy the identities.
\end{proof}

\begin{firewall}[Existence of the Symanzik expansion is an input]
\label{fw:v6v4-Symanzik-input}
The finite-dimensional argument begins with the uniform expansion
\eqref{eq:v6v4-packet-expansion}.  Establishing it for the interacting
chiral gauge theory on the bounded chart requires locality, differentiable
renormalization, and control of logarithmic coefficient dependence.  The
Hahn--Banach extraction theorem does not prove that expansion.
\end{firewall}

\begin{firewall}[Anomaly cancellation precedes improvement]
\label{fw:v6v4-anomaly-first}
If \(\mathcal A_a\ne0\), the gauge Ward defect can represent a nontrivial
cohomology class rather than an adjustable element of \(\mathcal Q_5\).
Regulator-level anomaly cancellation and determinant-line compatibility
must be proved before the finite counterterm matching system is invoked.
\end{firewall}

\begin{firewall}[On-shell and off-shell quotients differ]
\label{fw:v6v4-onshell-offshell}
Using equations of motion in \(\mathcal N_5\) proves on-shell improvement.
The rooted covariance and Wilsonian density may require off-shell
background derivatives, for which those operators cannot be discarded.
The quotient must be chosen to match the actual observable norm.
\end{firewall}

\begin{firewall}[Improvement must preserve the positivity route]
\label{fw:v6v4-improvement-positivity}
A counterterm that solves \eqref{eq:v6v4-finite-matching-system} may alter
the transfer matrix or reflection operation.  The final improved lattice
action must still satisfy the finite-cutoff positivity certificate; this is
not implied by smooth-band matching.
\end{firewall}

\begin{theorem}[Sixth-series V4 result]
\label{thm:v6v4-status}
Once the complete differentiated lattice-minus-continuum packet has a
uniform Symanzik expansion, its order-\(a\) obstruction is a
finite-dimensional class in \(\mathcal Q_5\).  A finite dual family of
smooth test backgrounds extracts every coefficient, and the canonical
smooth activity exponent \(b=1\) is equivalent to those coefficients being
\(O(a)\).  The matching must use one Ward-compatible counterterm functional
across propagator, vertex, stress, measure, ghost, and contact insertions.
The interacting expansion, regulator anomaly cancellation, correct
off-shell quotient, and positivity of the improved action remain open.
\end{theorem}

\begin{proof}
Use Lemma~\ref{lem:v6v4-Ward-packet},
Theorem~\ref{thm:v6v4-finite-improvement-test},
Corollary~\ref{cor:v6v4-smooth-b}, and
Proposition~\ref{prop:v6v4-Ward-compatible-matching}, subject to the four
firewalls.
\end{proof}

V5 is the next compulsory companion audit.  It will then attack the
existence of the interacting packet expansion by a finite-volume analytic
interpolation, separating logarithmic running coefficients from the local
dimension-five remainder and checking whether the companion all-order
transport estimates control its derivative rows.

\clearpage
\section{V5: companion audit and integrated Gevrey transport}
\label{sec:v6v5}

\subsection{Compulsory companion audit}

The active Yang--Mills file was
\texttt{Renewal\_Yang\_Mills\_Mass\_Gap\_Research\_Notes\_V5.tex}, with
\(81\,365\) bytes and SHA--256
\begin{center}
\texttt{7799D4A189B42E94CBDCD4742274B5A654578DAB0966D12E41F408084AD0F973}.
\end{center}
Its new result is an integrated Gevrey-one expansion of the exact normalized
scaled Wilson density.  If \(N\epsilon\) is sufficiently small, its
remainder is bounded by
\((KN\epsilon)^N+Ke^{-c/\epsilon}\).  The coefficient of order \(j\) is a
centered radial polynomial of degree at most \(4j\) with integrated
coefficient norm at most \(K^jj!\).  It also proves that connected source
expansion assigns the total density-defect order exactly once.  It does not
yet prove the derivative-row forest compiler needed for the Wilson
all-order tree card.

The Navier--Stokes companion remains
\texttt{Renewal\_NS\_Stability\_Research\_Notes\_V31.tex}, with
\(887\,537\) bytes and SHA--256
\begin{center}
\texttt{F1121B62238AD5491BB7CF03635EA9934942EDF573ED8FAAA9EE79E1D38A6696}.
\end{center}
Its one-use shell lesson agrees with both the owned high annuli in the
frozen SM volume and the owned total defect order in the new YM density
expansion.

The audit changes the V4 route.  A uniform pointwise complex chart for the
complete interacting density is sufficient but may be unnecessarily
strong near compact-group chart boundaries.  An integrated Gevrey
expansion can replace it at the expectation level, provided its coefficient
bound is upgraded to the complete derivative-row activity used in V2.

\subsection{Optimal truncation of a Gevrey activity series}

Let a normalized local density relative to a positive reference have the
finite expansion
\begin{equation}
 h_\epsilon=1+\sum_{j=1}^{N-1}\epsilon^jV_j+\mathcal R_N,
 \qquad \mathbb E_\gamma V_j=0.
 \label{eq:v6v5-density-expansion}
\end{equation}
Let \(z(V_j)\) denote the complete local derivative-row activity of V2.

\begin{theorem}[Gevrey coefficients give a small adjacent activity]
\label{thm:v6v5-Gevrey-activity}
Assume
\begin{equation}
 z(V_j)\le K^jj!,
 \qquad
 z(\mathcal R_N)\le (KN\epsilon)^N+Ke^{-c/\epsilon}
 \label{eq:v6v5-activity-Gevrey}
\end{equation}
whenever \(N\epsilon\le\delta\).  Choose
\begin{equation}
 N_\epsilon=\left\lfloor\frac{\delta_0}{K\epsilon}\right\rfloor,
 \qquad 0<\delta_0<\min\{\delta K,1/4\}.
 \label{eq:v6v5-optimal-N}
\end{equation}
Then for sufficiently small \(\epsilon\),
\begin{equation}
 \sum_{j=1}^{N_\epsilon-1}\epsilon^jz(V_j)
 +z(\mathcal R_{N_\epsilon})
 \le C_K\epsilon+C e^{-c'/\epsilon}.
 \label{eq:v6v5-small-zeta}
\end{equation}
Thus a local density parameter \(\epsilon_n\) contributes
\(\zeta_n=O(\epsilon_n)\) to the V2 adjacent activity.
\end{theorem}

\begin{proof}
Let \(a_j=(K\epsilon)^jj!\).  For
\(1\le j<N_\epsilon\),
\begin{equation}
 \frac{a_{j+1}}{a_j}=K\epsilon(j+1)le2\delta_0<1/2
 \label{eq:v6v5-ratio}
\end{equation}
after decreasing the small-\(\epsilon\) range to absorb the floor.
Therefore \(\sum_{j=1}^{N_\epsilon-1}a_j\le2K\epsilon\).
Moreover \(KN_\epsilon\epsilon\le\delta_0\), so
\begin{equation}
 (KN_\epsilon\epsilon)^{N_\epsilon}
 \le\delta_0^{\delta_0/(2K\epsilon)}
 \le e^{-c_1/\epsilon}.
 \label{eq:v6v5-optimal-remainder}
\end{equation}
Combine the two estimates and the second remainder term.
\end{proof}

\begin{corollary}[All-scale consequence]
\label{cor:v6v5-Gevrey-all-scale}
If \(\epsilon_n\le Ca_n^\omega(1+|\log a_n|)^M\) for some
\(\omega>0\) and fixed \(M\), then the activities in
\eqref{eq:v6v5-small-zeta} are summable on a geometric cutoff sequence.
If instead \(\epsilon_n\asymp n^{-1}\), this argument gives a divergent
harmonic activity and the coefficient must be included in the running
reference trajectory.
\end{corollary}

\begin{proof}
Every positive geometric power dominates a fixed logarithm, so
\(\sum_n\epsilon_n<\infty\).  The exponential remainders are smaller.  The
last statement is the divergence of the harmonic series.
\end{proof}

\subsection{Stability of cumulants under the integrated remainder}

The audited YM estimate controls its remainder in \(L^1\), not yet in the
V2 derivative activity.  An \(L^1\) remainder still has a useful bounded-
observable consequence.

\begin{lemma}[Cumulant stability in total variation]
\label{lem:v6v5-cumulant-TV}
Let \(\mu,\nu\) be probability measures with
\(\|\mu-\nu\|_{\rm TV}\le\Delta\), and let
\(|F_i|\le1\).  Then
\begin{equation}
 |\kappa_r^\mu(F_1,\ldots,F_r)
  -\kappa_r^\nu(F_1,\ldots,F_r)|
 \le \Delta\,r!\,C_0^r
 \label{eq:v6v5-TV-cumulant}
\end{equation}
for a numerical constant \(C_0\).
\end{lemma}

\begin{proof}
The moment--cumulant formula is
\begin{equation}
 \kappa_r(F_1,\ldots,F_r)=
 \sum_{\pi\in\Pi_r}(|\pi|-1)!(-1)^{|\pi|-1}
 \prod_{B\in\pi}\mathbb E\prod_{i\in B}F_i.
 \label{eq:v6v5-moment-cumulant}
\end{equation}
Every moment has modulus at most one and changes by at most a fixed
multiple of \(\Delta\).  A product indexed by \(\pi\) therefore changes by
at most \(C|\pi|\Delta\).  The exponential generating function for set
partitions, or the elementary bound obtained by assigning every element to
an ordered block, gives
\begin{equation}
 \sum_{\pi\in\Pi_r}|\pi|!le r!C_0^r.
 \label{eq:v6v5-partition-bound}
\end{equation}
Insert these estimates in \eqref{eq:v6v5-moment-cumulant}.
\end{proof}

\begin{theorem}[Arity window for an exponentially small density remainder]
\label{thm:v6v5-remainder-window}
Suppose a normalized density approximation has total-variation error
\begin{equation}
 \Delta_\epsilon\le Ce^{-c/\epsilon}.
 \label{eq:v6v5-exponential-TV}
\end{equation}
For bounded rooted observables, its cumulants differ from the exact ones by
at most
\begin{equation}
 r!\exp\left[-\frac c\epsilon+r\log C_1\right].
 \label{eq:v6v5-arity-remainder}
\end{equation}
Hence for every \(0<c_1<c/\log C_1\), the error remains exponentially small
uniformly for \(r\le c_1/\epsilon\).
\end{theorem}

\begin{proof}
Apply Lemma~\ref{lem:v6v5-cumulant-TV} and absorb the fixed prefactor into
\(C_1^r\).  In the stated arity window, the exponent in
\eqref{eq:v6v5-arity-remainder} is at most
\(-[c-c_1\log C_1]/\epsilon\).
\end{proof}

This theorem does not control arbitrary arity at fixed \(\epsilon\).  A
separate bounded-variable or exact compact-law estimate must handle the
complementary arity range, as the YM companion programme emphasizes.

\subsection{Owned defect order in the SM transport series}

Introduce formal sources \(t_i\) for rooted observables and auxiliary
sources \(u_j\) for the centered density defects \(V_j\).  The logarithm
\begin{equation}
 \log\mathbb E_\gamma
 \exp\left(\sum_it_iF_i+\sum_ju_jV_j\right)
 \label{eq:v6v5-auxiliary-source}
\end{equation}
has joint Gaussian cumulants as coefficients.

\begin{proposition}[One-use total Gevrey order]
\label{prop:v6v5-owned-order}
After substituting the polynomial density
\(1+\sum_{j<N}\epsilon^jV_j\) into the normalized connected source
functional, the coefficient of
\(\epsilon^qt_1\cdots t_r\) is a finite sum of connected Gaussian
cumulants whose defect labels have positive orders summing exactly to
\(q\).  No order is charged at more than one cutoff owner or more than once
inside that owner.
\end{proposition}

\begin{proof}
Expand \eqref{eq:v6v5-auxiliary-source} as a multivariate formal power
series.  Its coefficients are connected cumulants by definition of the
logarithm.  Substitute each auxiliary source by the corresponding
coefficient in the finite polynomial density and collect the ordinary
power of \(\epsilon\).  Multiplication of monomials adds their positive
orders once.  The cutoff owner is unique by the adjacent-density
factorization in V2.
\end{proof}

\subsection{The precise missing bridge}

The companion integrated coefficient norm controls
\(\|V_j\|_{L^1(\gamma)}\).  The V2 activity instead controls every complete
derivative row that can occur at a vertex of a BKAR tree.

\begin{definition}[Integrated derivative Gevrey card]
\label{def:v6v5-IDGC}
The IDGC is the estimate
\begin{equation}
 \sup_{d\ge0}\frac1{d!L^d}
 \mathbb E_\gamma\|\nabla^dV_j\|_{\rm row}
 \le K^jj!
 \label{eq:v6v5-IDGC}
\end{equation}
uniformly in local block dimension, cutoff, and every positive Gaussian
forest interpolation used by the connected formula.
\end{definition}

\begin{proposition}[IDGC closes the local Gevrey contribution]
\label{prop:v6v5-IDGC-closes}
If the IDGC and the corresponding remainder estimate hold for every local
positive owner, then Theorem~\ref{thm:v6v5-Gevrey-activity} supplies
\(\zeta_n^{\rm Gevrey}=O(\epsilon_n)\).  The V2 interaction-tree lemma and
Proposition~\ref{prop:v6v5-owned-order} then give the summed
interaction-insertion card without an additional defect-order incidence
factor.
\end{proposition}

\begin{proof}
The IDGC is precisely the integrated form of the activity hypothesis in
\eqref{eq:v6v5-activity-Gevrey}.  Sum the optimally truncated coefficient
activities with Theorem~\ref{thm:v6v5-Gevrey-activity}; the connected
Gaussian contractions are bounded by the V2 tree lemma.  The total
\(\epsilon\)-order is assigned once by
Proposition~\ref{prop:v6v5-owned-order}.
\end{proof}

\begin{firewall}[An \(L^1\) coefficient bound is not the IDGC]
\label{fw:v6v5-L1-not-IDGC}
Differentiation of a degree-\(4j\) polynomial can introduce powers of \(j\)
depending on its incident degree.  The audited Yang--Mills Gevrey norm does
not, by itself, prove \eqref{eq:v6v5-IDGC}.  Its announced defect-vertex
forest compiler is directly relevant but remains unfinished at V5.
\end{firewall}

\begin{firewall}[Integrated improvement differs from uniform packet
improvement]
\label{fw:v6v5-integrated-versus-uniform}
An integrated density expansion can prove expectation and cumulant bounds
without giving the uniform pointwise packet expansion assumed in V4.
These are two valid sufficient routes with different conclusions.  The
integrated route cannot be quoted as a pointwise Symanzik theorem.
\end{firewall}

\begin{firewall}[Chart tails and bad blocks must have the same owner]
\label{fw:v6v5-chart-tail-owner}
The exponentially small compact-group chart tail belongs to the bad-block
or density-remainder owner once, not to both.  Charging it in both V3 terms
would duplicate the gain and invalidate the activity ledger.
\end{firewall}

\begin{firewall}[The conformal direction has no positive density expansion]
\label{fw:v6v5-conformal-no-density}
The Gevrey density argument uses a normalized positive reference.  It does
not apply to the unrotated Euclidean conformal factor, whose action is
unbounded below.  No integrated series repairs the missing normalization.
\end{firewall}

\begin{theorem}[Sixth-series V5 result]
\label{thm:v6v5-status}
The companion audit replaces a needlessly strong pointwise-analytic target
by an integrated Gevrey alternative.  If its coefficient and remainder
bounds hold in the complete derivative activity, optimal truncation at
\(N\asymp\epsilon^{-1}\) gives a local adjacent activity \(O(\epsilon)\)
plus an exponentially small remainder.  Total defect order is owned once,
and total-variation stability controls bounded-observable cumulants through
arity \(O(\epsilon^{-1})\).  The exact remaining bridge is the integrated
derivative Gevrey card, uniformly over Gaussian forest interpolations; the
current companion result proves only the scalar integrated coefficient
norm.
\end{theorem}

\begin{proof}
Use Theorem~\ref{thm:v6v5-Gevrey-activity},
Corollary~\ref{cor:v6v5-Gevrey-all-scale},
Lemma~\ref{lem:v6v5-cumulant-TV},
Theorem~\ref{thm:v6v5-remainder-window}, and Propositions
\ref{prop:v6v5-owned-order} and
\ref{prop:v6v5-IDGC-closes}, subject to the four firewalls.
\end{proof}

V6 will attack the IDGC for a general radial polynomial coefficient.  It
will use Hermite/Ornstein--Uhlenbeck derivative identities to track the
incident degree at a defect vertex and test whether the order weight
\(\epsilon^j\) absorbs that degree with only factorial-exponential growth.

\clearpage
\section{V6: Hermite derivative rows and the correlated-product gate}
\label{sec:v6v6}

V5 proposed an integrated derivative Gevrey card.  Before using it in a
forest with many defect vertices, one must distinguish a one-vertex
derivative estimate from a correlated product estimate.  This note proves
the former and shows by example that it does not imply the latter.  It then
formulates the assembled total-order card that is actually sufficient for
the V2 transport compiler.

\subsection{A Gaussian chaos derivative identity}

Let \(\gamma_q\) be standard Gaussian measure on \(\mathbb R^q\), and let
\(P\) be a polynomial of degree at most \(D\).  Write its orthogonal Wiener
chaos decomposition as
\begin{equation}
 P=\sum_{n=0}^{D}P_n.
 \label{eq:v6v6-chaos-decomposition}
\end{equation}

\begin{lemma}[Complete derivative rows from chaos degree]
\label{lem:v6v6-chaos-derivative}
For every \(d\ge0\),
\begin{equation}
 \mathbb E_{\gamma_q}|\nabla^dP\|_{\rm HS}^2
 \le D^d\|P\|_{L^2(\gamma_q)}^2,
 \label{eq:v6v6-chaos-derivative}
\end{equation}
where the left side is zero for \(d>D\).  Consequently, for every \(L>0\),
\begin{equation}
 \sup_{d\ge0}\frac{1}{d!L^d}
 \mathbb E_{\gamma_q}\|\nabla^dP\|_{\rm row}
 \le e^{\sqrt D/L}\|P\|_{L^2(\gamma_q)}.
 \label{eq:v6v6-one-vertex-IDGC}
\end{equation}
The constants are independent of \(q\).
\end{lemma}

\begin{proof}
For a homogeneous chaos \(P_n\), the Gaussian Malliavin derivative identity
gives
\begin{equation}
 \mathbb E\|\nabla^dP_n\|_{\rm HS}^2
 =(n)_d\|P_n\|_2^2,
 \qquad (n)_d=n(n-1)\cdots(n-d+1).
 \label{eq:v6v6-falling-factorial}
\end{equation}
Different chaos orders remain orthogonal after tensor differentiation, so
summing \eqref{eq:v6v6-falling-factorial} and using
\((n)_d\le D^d\) proves \eqref{eq:v6v6-chaos-derivative}.  The complete row
norm is bounded by the Hilbert--Schmidt norm, and Cauchy--Schwarz changes
the \(L^2\) estimate into an \(L^1\) estimate.  Finally,
\[
 \sup_{d\ge0}\frac{D^{d/2}}{d!L^d}
 \le\sum_{d\ge0}\frac{(\sqrt D/L)^d}{d!}
 =e^{\sqrt D/L}.
\]
\end{proof}

\subsection{From an absolute monomial norm to \(L^2\)}

For a polynomial \(P(y)=\sum_\alpha c_\alpha y^\alpha\), define
\begin{equation}
 \|P\|_{\mathfrak G,1}:=
 \sum_\alpha|c_\alpha|\mathbb E_{\gamma_q}|Y^\alpha|.
 \label{eq:v6v6-absolute-monomial-norm}
\end{equation}
The radial norm in the audited YM V5 is a special case.

\begin{lemma}[Gaussian moment doubling]
\label{lem:v6v6-moment-doubling}
There is a numerical \(C\) such that for every multi-index \(\alpha\),
\begin{equation}
 \left(\mathbb E|Y^\alpha|^2\right)^{1/2}
 \le C^{|\alpha|}\mathbb E|Y^\alpha|.
 \label{eq:v6v6-moment-doubling}
\end{equation}
Consequently, if \(\deg P\le D\), then
\begin{equation}
 \|P\|_2\le C^D\|P\|_{\mathfrak G,1}.
 \label{eq:v6v6-G-to-L2}
\end{equation}
\end{lemma}

\begin{proof}
Gaussian coordinates are independent, so it is enough to prove the
one-dimensional estimate.  The exact moments are
\begin{equation}
 \mathbb E|Z|^m=2^{m/2}
 \frac{\Gamma((m+1)/2)}{\sqrt\pi}.
 \label{eq:v6v6-Gaussian-moment}
\end{equation}
The duplication formula and Stirling bounds give
\((\mathbb E|Z|^{2m})^{1/2}\le C^m\mathbb E|Z|^m\) uniformly in \(m\).
Multiply over coordinates.  The triangle inequality in \(L^2\), followed
by \(|\alpha|\le D\), gives \eqref{eq:v6v6-G-to-L2}.
\end{proof}

\begin{theorem}[One-vertex Gevrey derivative upgrade]
\label{thm:v6v6-one-vertex-upgrade}
Suppose \(V_j\) has degree at most \(4j\) and
\begin{equation}
 \|V_j\|_{\mathfrak G,1}\le K^jj!.
 \label{eq:v6v6-Gevrey-coefficient}
\end{equation}
Then for every fixed \(L>0\),
\begin{equation}
 \sup_{d\ge0}\frac{1}{d!L^d}
 \mathbb E\|\nabla^dV_j\|_{\rm row}
 \le K_1^jj!,
 \label{eq:v6v6-one-vertex-Gevrey}
\end{equation}
with \(K_1\) independent of \(j\) and of the Gaussian dimension.
The same estimate holds at one replica of every Gaussian forest
interpolation whose marginal covariance at that replica is the identity.
\end{theorem}

\begin{proof}
Lemmas \ref{lem:v6v6-chaos-derivative} and
\ref{lem:v6v6-moment-doubling} give the left side bounded by
\begin{equation}
 e^{2\sqrt j/L}C^{4j}K^jj!.
 \label{eq:v6v6-upgrade-bound}
\end{equation}
Since \(e^{2\sqrt j/L}\le C_L^j\) for \(j\ge1\), all factors are absorbed
into \(K_1^j\).  The law of one replica depends only on its diagonal
covariance block, which remains the identity in the forest interpolation.
\end{proof}

This proves the one-vertex form of the V5 IDGC for the audited radial
coefficients.  It is not yet the product bound used when several defect
vertices occur in one connected Gaussian expectation.

\subsection{Why individual integrated rows do not tensorize}

\begin{proposition}[Perfect-correlation counterexample]
\label{prop:v6v6-perfect-correlation}
Let \(Y_1=\cdots=Y_m=Z\) be perfectly correlated standard Gaussians, and
put
\begin{equation}
 P_j(z):=\frac{|z|^{2j}}{\mathbb E|Z|^{2j}}.
 \label{eq:v6v6-normalized-monomial}
\end{equation}
Then \(\mathbb E|P_j(Y_i)|=1\) at every vertex, but
\begin{equation}
 \mathbb E\prod_{i=1}^mP_j(Y_i)
 =\frac{\mathbb E|Z|^{2jm}}
 {\bigl(\mathbb E|Z|^{2j}\bigr)^m},
 \label{eq:v6v6-correlated-product}
\end{equation}
whose logarithm grows at least as \(cjm\log m\) for fixed \(j\) and large
\(m\).  Hence no product estimate with a constant \(K^{jm}\) follows from
the individual \(L^1\) bounds alone.
\end{proposition}

\begin{proof}
The identity is immediate from perfect correlation.  Using
\eqref{eq:v6v6-Gaussian-moment} and Stirling's formula,
\begin{equation}
 \log\mathbb E|Z|^{2n}=n\log n+O(n).
 \label{eq:v6v6-moment-asymptotic}
\end{equation}
Apply it with \(n=jm\) and with \(n=j\) multiplied by \(m\).  Their
difference is \(jm\log m+O(jm)\), proving the claim.
\end{proof}

Forest interpolation covariances can approach perfect correlation when
edge parameters approach one.  Therefore a proof uniform in those
parameters must use the exact coefficient structure and connected
assembly; generalized H\"older applied to separately normed vertices loses
too much.

\begin{assessment}[Correction to the V5 closing statement]
\label{ass:v6v6-V5-correction}
Proposition V5.6 is valid only when the phrase ``IDGC'' is read as including
the assembled multi-vertex product bound below.  The one-vertex estimate
\eqref{eq:v6v5-IDGC}, by itself, does not justify its appeal to the V2 tree
lemma, whose original hypothesis was a pointwise derivative activity.  The
distinction is now explicit and no all-order conclusion is drawn from the
one-vertex upgrade.
\end{assessment}

\subsection{The sufficient assembled total-order card}

Let defect vertices be labelled by positive orders
\(j_1,\ldots,j_m\), let \(J=\sum_{ell=1}^mj_\ell\), and let
\(d_\ell\) be their incident derivative degrees in a Gaussian forest.

\begin{definition}[Assembled product Gevrey card]
\label{def:v6v6-APGC}
The APGC is the uniform estimate
\begin{align}
 &\mathbb E_{C^F}
 \prod_{\ell=1}^m
 \frac{\|\nabla^{d_\ell}V_{j_\ell}(Y_\ell)\|_{\rm row}}
 {d_\ell!L^{d_\ell}}
 \le K^{J+m}J!
 \label{eq:v6v6-APGC}
\end{align}
after summing all coefficient compositions that produce the same total
order \(J\).  The bound must hold uniformly in the forest, its interpolation
parameters, local block dimension, and cutoff.
\end{definition}

The order is combined before taking the norm.  This is stronger than the
product of one-vertex Gevrey bounds and is compatible with the one-use
defect-order principle audited at V5.

\begin{theorem}[APGC gives a small all-order density activity]
\label{thm:v6v6-APGC-closes}
Assume the APGC and an activity remainder
\((KN\epsilon)^N+Ke^{-c/\epsilon}\).  Then, for
\(N\asymp\epsilon^{-1}\) chosen with a sufficiently small proportionality
constant, the sum of every connected density-defect contribution of total
positive order \(J<N\) is bounded by
\begin{equation}
 \sum_{J=1}^{N-1}(K_2\epsilon)^JJ!
 \le C\epsilon,
 \label{eq:v6v6-total-order-sum}
\end{equation}
and the remainder is exponentially small.  The resulting adjacent
interaction parameter satisfies \(\eta_n\le C\epsilon_n\).
\end{theorem}

\begin{proof}
The APGC has already summed the compositions and correlated defect
vertices of total order \(J\).  The number of remaining connected
attachment patterns and derivative contractions is exponential in \(J\)
after the Gaussian BKAR tree is assembled; absorb it into \(K_2^J\).  The
ratio of successive terms in \eqref{eq:v6v6-total-order-sum} is
\(K_2\epsilon(J+1)\).  Choose \(N=\lfloor\delta/(K_2\epsilon)\rfloor\)
with \(2\delta<1\); the same geometric-ratio proof as in V5 gives the
\(C\epsilon\) bound.  Optimal truncation makes the stated remainder
exponential.  The V1--V2 transport theorem identifies this summed activity
with \(\eta_n\).
\end{proof}

\begin{firewall}[The APGC is not proved by hypercontractivity]
\label{fw:v6v6-hypercontractivity}
Applying Gaussian hypercontractivity or H\"older separately to \(m\)
vertices introduces a power of \(m\) proportional to their total degree,
as the perfect-correlation example predicts.  The APGC must exploit exact
coefficient denominators, connected half-edge incidence, and summation of
all compositions before norms.
\end{firewall}

\begin{firewall}[The one-vertex upgrade covers coefficients, not the
remainder]
\label{fw:v6v6-remainder-derivatives}
Theorem~\ref{thm:v6v6-one-vertex-upgrade} applies to polynomial
coefficients \(V_j\).  An \(L^1\) estimate for the nonpolynomial truncation
remainder does not control all of its derivatives.  For bounded observables
one may use the V5 total-variation arity window; a derivative remainder
requires a separate analytic or integration-by-parts estimate.
\end{firewall}

\begin{firewall}[The SM local density must have an audited Gevrey structure]
\label{fw:v6v6-SM-Gevrey}
The radial Wilson density has explicit action and Haar series with helpful
factorial denominators.  The coupled SM positive density also contains
Higgs, PV, determinant-modulus, ghost, and counterterm factors.  Each must
be placed in a common total-order expansion before the APGC can be tested.
\end{firewall}

\begin{theorem}[Sixth-series V6 result]
\label{thm:v6v6-status}
A Gaussian polynomial of degree \(D\) has complete integrated derivative
rows bounded by \(e^{\sqrt D/L}\) times its \(L^2\) norm.  The audited
absolute monomial Gevrey bound for degree at most \(4j\) therefore implies
a one-vertex derivative bound \(K^jj!\), uniformly in Gaussian dimension
and in one-replica forest marginals.  This estimate does not tensorize over
correlated defect vertices; a perfect-correlation example has growth
\(e^{cjm\log m}\).  The sufficient replacement is the assembled product
Gevrey card at total owned order \(J\), which would give an
\(O(\epsilon)\) adjacent activity after optimal truncation.  Proving that
assembled card for the complete SM local density is the live gate.
\end{theorem}

\begin{proof}
Use Lemmas \ref{lem:v6v6-chaos-derivative} and
\ref{lem:v6v6-moment-doubling},
Theorem~\ref{thm:v6v6-one-vertex-upgrade},
Proposition~\ref{prop:v6v6-perfect-correlation}, and
Theorem~\ref{thm:v6v6-APGC-closes}, subject to the assessment and three
firewalls.
\end{proof}

V7 will attack the APGC by retaining the exact exponential generating
function of each local positive owner.  Instead of norming separate
\(V_j\)'s, it will differentiate the assembled coefficient of total order
\(J\) and test whether the action/Haar/high-determinant factorial
denominators absorb the perfect-correlation moment growth.

\clearpage
\section{V7: bounded occupancy closes the polynomial APGC}
\label{sec:v6v7}

\paragraph{Source correction.}
In the first paragraph of the V6 subsection defining the assembled
total-order card, the displayed definition of \(J\) contains the literal
text \(\sum_{ell=1}^m j_\ell\).  It is to be read as
\begin{equation}
 J=\sum_{\ell=1}^m j_\ell.
 \label{eq:v6v7-V6-correction}
\end{equation}
No proof or estimate is changed.

V6 showed that independent one-vertex norms cannot control arbitrarily many
copies of the same perfectly correlated Gaussian variable.  The local
density expansion has more structure than that counterexample: after it is
written as a product of labelled owner factors, each factor is selected at
most once.  If the number of labels anchored at one block is uniformly
bounded, spatial \(\ell^1\) covariance controls the full replicated
covariance operator independently of the number of selected factors.

\subsection{Owner-factor expansion and occupancy}

At one cutoff step, group the positive local increment into a family of
normalized owner factors
\begin{equation}
 e^{-W_n}=\prod_{A\in\mathcal P_n}h_{n,A},
 \qquad
 h_{n,A}=1+K_{n,A}.
 \label{eq:v6v7-owner-product}
\end{equation}
The algebraic expansion is
\begin{equation}
 \prod_{A\in\mathcal P_n}(1+K_{n,A})
 =\sum_{\mathcal S\subset\mathcal P_n}
   \prod_{A\in\mathcal S}K_{n,A}.
 \label{eq:v6v7-subset-expansion}
\end{equation}
Every labelled factor occurs zero or one time.  Overlapping supports are
allowed.

\begin{definition}[Bounded anchor occupancy]
\label{def:v6v7-occupancy}
Choose an anchor \(a(A)\) for every owner.  The owner presentation has
occupancy \(M_0\) if
\begin{equation}
 \#\{A\in\mathcal P_n:a(A)=x\}\le M_0
 \label{eq:v6v7-occupancy}
\end{equation}
for every block \(x\), cutoff step, and volume.
\end{definition}

Finite-range interaction shapes have bounded occupancy.  An infinite
trace-loop family must first be summed into finitely many exponentially
local anchored factors; treating every loop as a separate owner does not
satisfy Definition~\ref{def:v6v7-occupancy}.

\subsection{A replicated covariance bound}

Let \(C(x,y)\) be the block covariance of the massive Gaussian reference,
normalized so that its diagonal fibre covariance is the identity.  Assume
\begin{equation}
 C_1:=\sup_x\sum_y\|C(x,y)\|<\infty.
 \label{eq:v6v7-C-row}
\end{equation}
For a selected family \(A_1,\ldots,A_m\), let \(\Sigma^F\) be any BKAR
forest interpolation covariance.  Its \((i,j)\) block is
\begin{equation}
 \Sigma^F_{ij}=s^F_{ij}C(a(A_i),a(A_j)),
 \qquad 0\le s^F_{ij}\le1,
 \label{eq:v6v7-forest-covariance}
\end{equation}
with identity diagonal blocks.

\begin{lemma}[Occupancy bounds the replicated covariance]
\label{lem:v6v7-replicated-covariance}
If the selected owners obey \eqref{eq:v6v7-occupancy}, then
\begin{equation}
 \|\Sigma^F\|_{\ell^2(\{1,\ldots,m\};\mathbb R^q)}
 \le p_0:=M_0C_1,
 \label{eq:v6v7-Sigma-bound}
\end{equation}
uniformly in \(m\), the forest, and its interpolation parameters.
\end{lemma}

\begin{proof}
For a fixed selected owner \(i\), at most \(M_0\) selected labels have any
given anchor \(y\).  Hence
\begin{equation}
 \sum_j\|\Sigma^F_{ij}\|
 \le M_0\sum_y\|C(a(A_i),y)\|
 \le M_0C_1.
 \label{eq:v6v7-row-bound}
\end{equation}
The symmetric column bound is the same.  The block Schur test proves
\eqref{eq:v6v7-Sigma-bound}.
\end{proof}

The perfect-correlation example of V6 repeats one anchor \(m\) times, so
its occupancy and covariance norm both grow like \(m\).  It is excluded by
the exact subset structure \eqref{eq:v6v7-subset-expansion}.

\subsection{Gaussian Brascamp--Lieb and polynomial rows}

\begin{lemma}[Gaussian product inequality]
\label{lem:v6v7-Gaussian-BL}
Let \((Y_1,\ldots,Y_m)\) be centered jointly Gaussian with identity
marginals.  If its block covariance satisfies \(\Sigma\le p_0I\) for some
\(p_0\ge1\), then for nonnegative measurable \(f_i\),
\begin{equation}
 \mathbb E_\Sigma\prod_{i=1}^mf_i(Y_i)
 \le\prod_{i=1}^m\|f_i\|_{L^{p_0}(\gamma_q)}.
 \label{eq:v6v7-Gaussian-BL}
\end{equation}
\end{lemma}

\begin{proof}
This is the Gaussian Brascamp--Lieb inequality for the coordinate maps
\(B_i(Y_1,\ldots,Y_m)=Y_i\).  Its matrix condition is
\begin{equation}
 \Sigma\le\operatorname{diag}(p_0I_q,\ldots,p_0I_q),
 \label{eq:v6v7-BL-condition}
\end{equation}
which is the hypothesis.  For a direct proof, regularize the functions and
apply the Gaussian heat semigroup to
\(\log\mathbb E\prod_i(P_tf_i)^{1/p_0}\).  Gaussian integration by parts
makes its derivative a quadratic form with coefficient
\(p_0I-\Sigma\), hence nonnegative.  At infinite time the variables
decouple into their means; reversing the interpolation yields
\eqref{eq:v6v7-Gaussian-BL}.  Approximation removes regularity and strict
positivity assumptions.
\end{proof}

For \(p_0<2\), replace it by \(2\); the covariance inequality remains
valid.  Gaussian hypercontractivity for a polynomial \(P\) of degree
\(D\) gives
\begin{equation}
 \|P\|_{L^{p_0}(\gamma_q)}
 \le(p_0-1)^{D/2}\|P\|_{L^2(\gamma_q)}.
 \label{eq:v6v7-hypercontractivity}
\end{equation}

\begin{lemma}[One defect vertex in the Brascamp--Lieb norm]
\label{lem:v6v7-one-vertex-p}
If \(V_j\) has degree at most \(4j\) and satisfies the V6 absolute-monomial
bound \(\|V_j\|_{\mathfrak G,1}\le K^jj!\), then for every derivative
degree \(d\),
\begin{equation}
 \frac{\|\nabla^dV_j\|_{L^{p_0}(\gamma_q);\,\rm row}}
 {d!L^d}
 \le K_1^jj!,
 \label{eq:v6v7-Lp-row}
\end{equation}
where \(K_1\) depends on \(p_0,L\) but not on \(j,d,q\).
\end{lemma}

\begin{proof}
The derivative is a polynomial of degree at most \(4j\).  Apply
\eqref{eq:v6v7-hypercontractivity} to its complete row norm by first using
the Hilbert-valued version, whose proof is the same chaos estimate.
The factor \((p_0-1)^{2j}\) is exponential in \(j\).  The V6 chaos
derivative and moment-doubling lemmas bound the remaining normalized
\(L^2\) row by another exponential in \(j\) times \(j!\).  Absorb both
into \(K_1^j\).
\end{proof}

\subsection{Closure of the polynomial assembled card}

\begin{theorem}[Bounded-occupancy polynomial APGC]
\label{thm:v6v7-polynomial-APGC}
Assume:
\begin{enumerate}
\item the owner-factor representation has occupancy \(M_0\);
\item the massive covariance has the row bound
      \eqref{eq:v6v7-C-row}; and
\item every order-\(j\) polynomial coefficient of every owner satisfies
      the degree and absolute-monomial hypotheses of
      Lemma~\ref{lem:v6v7-one-vertex-p}, uniformly in its owner type.
\end{enumerate}
Then, for any selected defect vertices with orders
\(j_1,\ldots,j_m\), total order \(J=\sum_i j_i\), incident degrees \(d_i\),
and any forest interpolation,
\begin{equation}
 \mathbb E_{\Sigma^F}
 \prod_{i=1}^m
 \frac{\|\nabla^{d_i}V_{j_i}(Y_i)\|_{\rm row}}
 {d_i!L^{d_i}}
 \le K_2^J\prod_{i=1}^mj_i!
 \le K_2^JJ!.
 \label{eq:v6v7-APGC-result}
\end{equation}
After summing the at most \(2^{J-1}\) ordered compositions of \(J\), the
assembled product Gevrey card of V6 holds with an enlarged exponential
constant.
\end{theorem}

\begin{proof}
Lemma~\ref{lem:v6v7-replicated-covariance} gives
\(\Sigma^F\le p_0I\).  Apply
Lemma~\ref{lem:v6v7-Gaussian-BL} to the nonnegative row norms and then
Lemma~\ref{lem:v6v7-one-vertex-p} at every vertex.  This proves the first
bound.  The multinomial coefficient
\(J!/(j_1!\cdots j_m!)\) is an integer at least one, proving the second.
The number of compositions of \(J\) is \(2^{J-1}\), which is absorbed into
the exponential base.
\end{proof}

\begin{corollary}[Polynomial adjacent density activity]
\label{cor:v6v7-polynomial-activity}
Under Theorem~\ref{thm:v6v7-polynomial-APGC}, the optimally truncated
polynomial part of an integrated Gevrey owner density has adjacent activity
\begin{equation}
 \eta_n^{\rm poly}\le C\epsilon_n
 \label{eq:v6v7-eta-poly}
\end{equation}
uniformly in arity, provided the connected owner-support tree sum satisfies
the V2 spatial activity bound.
\end{corollary}

\begin{proof}
Theorem~\ref{thm:v6v7-polynomial-APGC} supplies the V6 APGC.  Apply the V6
optimal total-order sum and then the V2 owner-support tree lemma.
\end{proof}

\subsection{Why the density-factor representation matters}

The exponential tilt series in V1 contains repeated copies of the same
\(W_A\).  It is exact and useful for formal identities, but a direct
absolute estimate permits arbitrarily high occupancy at one anchor.  The
factor expansion \eqref{eq:v6v7-subset-expansion} is also exact and is
better suited to the integrated Gevrey compiler.

\begin{proposition}[No repeated owner in the normalized source ratio]
\label{prop:v6v7-no-repeated-owner}
For finite \(\mathcal P_n\), both the numerator with rooted sources and the
vacuum denominator are multilinear polynomials in the labelled variables
\(K_{n,A}\).  Expanding their ratio by connected set partitions or a
polymer logarithm uses each label at most once in every monomial.  Vacuum
components cancel as in the frozen rooted cancellation theorem.
\end{proposition}

\begin{proof}
Equation \eqref{eq:v6v7-subset-expansion} is square-free in every labelled
variable.  The logarithm of a square-free polymer generating function is a
sum over connected clusters of distinct labels; repeated labels arise only
as an intermediate formal power-series representation and cancel under
multilinear coefficient extraction.  In the normalized ratio, any
component disjoint from all rooted sources appears identically in numerator
and denominator and cancels.
\end{proof}

\begin{firewall}[Grouping infinitely many loops must be proved]
\label{fw:v6v7-loop-grouping}
The high determinant logarithm has infinitely many exponentially local
loops before truncation.  Bounded occupancy requires summing them into a
finite family of anchored owner functions with complete derivative Gevrey
bounds.  Exponential kernel locality suggests this grouping but does not by
itself prove the required coefficient norm.
\end{firewall}

\begin{firewall}[Massless covariance has no Brascamp--Lieb row constant]
\label{fw:v6v7-massless-BL}
The covariance row bound \eqref{eq:v6v7-C-row} is the massive V100 input.
A photon or gravitational constraint covariance has a divergent bare row
sum.  Endpoint Ward and multipole differences must be inserted before any
analogue of Lemma~\ref{lem:v6v7-replicated-covariance} is used.
\end{firewall}

\begin{firewall}[The nonpolynomial remainder is still open]
\label{fw:v6v7-remainder-open}
Theorem~\ref{thm:v6v7-polynomial-APGC} controls the optimally truncated
polynomial density.  The integrated remainder has only the V5
total-variation arity window unless a complete derivative or bounded-owner
cluster estimate is proved.  It cannot yet be included in
\eqref{eq:v6v7-eta-poly} at arbitrary arity.
\end{firewall}

\begin{firewall}[Owner occupancy is a representation property]
\label{fw:v6v7-occupancy-property}
Artificially duplicating one owner into many identical labels increases
\(M_0\) and destroys the uniform covariance bound.  The bounded-occupancy
presentation must follow from a canonical local grouping of the actual
adjacent density, not from a convenient relabelling.
\end{firewall}

\begin{theorem}[Sixth-series V7 result]
\label{thm:v6v7-status}
For a massive Gaussian reference with a uniform spatial covariance row,
an exact square-free owner-factor expansion with bounded anchor occupancy
has a replicated covariance operator norm independent of the number of
selected factors.  Gaussian Brascamp--Lieb and polynomial
hypercontractivity then upgrade the one-vertex Gevrey estimate to the
assembled product card \(K^JJ!\).  The optimally truncated polynomial
density therefore contributes \(O(\epsilon_n)\) to the all-order adjacent
activity.  Canonical grouping of the infinite determinant loops and
all-arity control of the nonpolynomial remainder remain open.
\end{theorem}

\begin{proof}
Use Lemmas \ref{lem:v6v7-replicated-covariance},
\ref{lem:v6v7-Gaussian-BL}, and
\ref{lem:v6v7-one-vertex-p},
Theorem~\ref{thm:v6v7-polynomial-APGC}, and
Proposition~\ref{prop:v6v7-no-repeated-owner}, subject to the four
firewalls.
\end{proof}

V8 will attack the nonpolynomial remainder by keeping each exact bounded
compact-gauge owner as one square-free factor.  It will derive an
all-arity cumulant/cluster estimate directly from its sup norm and spatial
support, and will separate the unbounded Higgs and determinant-modulus
remainders that require a different tail estimate.

\clearpage
\section{V8: sparse-color control of the exact compact-gauge remainder}
\label{sec:v6v8}

V7 closes the polynomial part of the assembled Gevrey card but leaves the
nonpolynomial compact-chart remainder.  Its \(L^1\) smallness cannot be
multiplied over correlated replicas.  The remedy is to retain an
\(L^p\) estimate with \(p>1\), and to make the Gaussian product inequality
use that same \(p\) by sparsifying the owner anchors.

\subsection{The compact-density \(L^p\) threshold}

Consider first the normalized scaled one-link density on a compact
exponential chart.  Let \(\theta=\sqrt\epsilon,y\) range in a bounded
star-shaped set \(\Omega\subset\mathbb R^q\), and write
\begin{equation}
 h_\epsilon(y)=Z_\epsilon^{-1}
 J(\sqrt\epsilon,y)
 \exp\left[-\epsilon^{-1}s(\sqrt\epsilon,y)
            +\frac12|y|^2\right]
 {\bf1}_{\{\sqrt\epsilon y\in\Omega\}}
 \label{eq:v6v8-scaled-density}
\end{equation}
relative to standard Gaussian measure.  Assume
\begin{equation}
 0\le J(\theta)\le C_J,
 \qquad
 s(\theta)\ge c_s|\theta|^2,
 \qquad 0<c_s<\frac12,
 \label{eq:v6v8-action-coercivity}
\end{equation}
and \(Z_\epsilon\) is bounded above and below for small \(\epsilon\).

\begin{theorem}[Uniform compact-density \(L^p\) interval]
\label{thm:v6v8-Lp-density}
For every
\begin{equation}
 1<p<p_*:=\frac{1}{1-2c_s},
 \label{eq:v6v8-p-star}
\end{equation}
there is \(C_p\), independent of sufficiently small \(\epsilon\), such that
\begin{equation}
 \|h_\epsilon\|_{L^p(\gamma_q)}\le C_p.
 \label{eq:v6v8-uniform-Lp}
\end{equation}
For the elementary \(SU(2)\) radial Wilson inequality
\(1-\cos\theta\ge2\theta^2/\pi^2\) on \(|\theta|\le\pi\), this gives
\begin{equation}
 p_* =\frac{\pi^2}{\pi^2-4}>1.
 \label{eq:v6v8-SU2-pstar}
\end{equation}
\end{theorem}

\begin{proof}
Raise \eqref{eq:v6v8-scaled-density} to the power \(p\) and multiply by the
Gaussian density.  Apart from bounded normalization and Jacobian factors,
the exponent is
\begin{equation}
 -\frac p\epsilon s(\sqrt\epsilon y)
 +\frac{p-1}{2}|y|^2
 \le-\left(pc_s-\frac{p-1}{2}\right)|y|^2.
 \label{eq:v6v8-Lp-exponent}
\end{equation}
The coefficient is positive exactly when
\(p(1-2c_s)<1\), which is \eqref{eq:v6v8-p-star}.  The resulting Gaussian
integral is uniform in \(\epsilon\).  Substitute \(c_s=2/\pi^2\) for the
last formula.
\end{proof}

The interval can be narrow.  Replacing its exponent by \(2\) in the V7
Brascamp--Lieb bound would discard the useful information; the actual
exponent \(p_0\) must be retained.

\subsection{Sparse coloring makes \(p_0\) close to one}

Let the normalized massive covariance obey
\begin{equation}
 \|C(x,y)\|\le C_0e^{-\mu d(x,y)},
 \qquad C(x,x)=I.
 \label{eq:v6v8-massive-C}
\end{equation}
Color the block lattice with finitely many colors so that distinct anchors
of one color are at distance at least \(L\).  There are at most
\(N_L\le C(1+L)^d\) colors on a bounded-degree \(d\)-dimensional block
graph.

\begin{lemma}[Sparse-color covariance exponent]
\label{lem:v6v8-sparse-covariance}
Assume one grouped owner per anchor in a fixed color substep.  Then every
BKAR replica covariance for selected anchors of that color satisfies
\begin{equation}
 \Sigma^F\le p_LI,
 \qquad
 p_L:=1+\rho_L,
 \qquad
 \rho_L\le C_Le^{-\mu L/2},
 \label{eq:v6v8-pL}
\end{equation}
where \(C_L\) grows at most polynomially in \(L\).  In particular,
\(p_L\downarrow1\) as \(L\to\infty\).
\end{lemma}

\begin{proof}
The diagonal block contributes one.  For a selected anchor \(x\), at most
polynomially many same-color anchors lie in a shell of radius \(r\), and
none lies closer than \(L\).  Therefore
\begin{equation}
 \sum_{y\ne x}\|C(x,y)\|
 \le C\sum_{r\ge L}(1+r)^{d-1}e^{-\mu r}
 \le C_Le^{-\mu L/2}.
 \label{eq:v6v8-offdiagonal-row}
\end{equation}
BKAR interpolation multipliers have modulus at most one.  The symmetric
block Schur test gives \eqref{eq:v6v8-pL}.
\end{proof}

\begin{corollary}[Compatibility of density and covariance exponents]
\label{cor:v6v8-compatible-p}
For every compact owner satisfying
Theorem~\ref{thm:v6v8-Lp-density}, one can choose a fixed \(L\), independent
of cutoff and volume, such that
\begin{equation}
 1<p_L<p_*.
 \label{eq:v6v8-p-compatible}
\end{equation}
The Gaussian product inequality of V7 then uses the uniform density
\(L^{p_L}\) norm.
\end{corollary}

\begin{proof}
The interval \((1,p_*)\) is open and
\(p_L\to1\).  Choose \(L\) once so that
\eqref{eq:v6v8-p-compatible} holds.
\end{proof}

Splitting one cutoff step into \(N_L\) color substeps multiplies its
activity by a fixed constant only.  It does not alter cutoff summability.

\subsection{An \(L^p\) Gevrey remainder}

The proof of the integrated compact-density expansion uses a central
Taylor series and the coercive chart tail.  The same proof works in the
strict interval \(p<p_*\), with a smaller Gaussian decay constant.

\begin{theorem}[Compact-density Gevrey remainder in \(L^p\)]
\label{thm:v6v8-Lp-Gevrey}
Assume in addition to \eqref{eq:v6v8-action-coercivity} that \(s\) and
\(J\) have convergent even chart expansions whose order-\(j\) coefficients
have the factorial denominators and degree bound \(D_j\le c_Dj\) used in
the compact Wilson density.  For every fixed \(1<p<p_*\), there are
centered polynomials \(V_j\) and constants \(K_p,c_p,\delta_p>0\) such that
\begin{align}
 \deg V_j&\le c_D'j,
 \qquad
 \|V_j\|_{\mathfrak G,p}\le K_p^jj!,
 \label{eq:v6v8-Lp-coefficients}\\
 \left\|h_\epsilon-left(1+sum_{j=1}^{N-1}epsilon^jV_j\right)
 \right\|_{L^p(\gamma_q)}
 &\le(K_pN\epsilon)^N+K_pe^{-c_p/\epsilon}
 \label{eq:v6v8-Lp-remainder}
\end{align}
whenever \(N\epsilon\le\delta_p\).  Here
\(\|\cdot\|_{\mathfrak G,p}\) is the absolute coefficient norm weighted by
the corresponding \(L^p\) Gaussian monomial norms.
\end{theorem}

\begin{proof}
In the \(p\)-th power estimate, equation
\eqref{eq:v6v8-Lp-exponent} leaves a Gaussian weight
\(e^{-c(p)|y|^2}\) with \(c(p)>0\).  Gaussian moment convolution with this
fixed variance changes every coefficient estimate by only an exponential
base.  Expanding the analytic action and Jacobian through total order
\(N-1\), the factorial denominators absorb the moment factorials exactly as
in the \(L^1\) argument.  Taylor's integral remainder gives
\((K_pN\epsilon)^N\) on the central region.  On the chart tail,
\eqref{eq:v6v8-Lp-exponent} yields
\(e^{-c_p/\epsilon}\), including all polynomials of degree \(O(N)\) when
\(N\epsilon\le\delta_p\).  Normalizing the finite series uses formal
inversion of a Gevrey-one scalar series, which preserves the coefficient
class and centers every \(V_j\).
\end{proof}

The theorem is an abstract compact-density statement.  The audited
one-link Wilson density satisfies its explicit trigonometric hypotheses.
Every SM compact gauge owner must be checked in its actual representation
and chart.

\subsection{All-arity control of the exact compact remainder}

Choose \(p=p_L<p_*\) and the optimal order
\(N_\epsilon\asymp\epsilon^{-1}\).  Put
\begin{equation}
 R_{\epsilon,A}:=
 h_{\epsilon,A}-1-sum_{j=1}^{N_\epsilon-1}epsilon^jV_{j,A}.
 \label{eq:v6v8-exact-remainder}
\end{equation}

\begin{theorem}[Square-free compact remainder card]
\label{thm:v6v8-remainder-card}
In one sparse-color substep, assume each anchor carries at most one compact
owner satisfying Theorem~\ref{thm:v6v8-Lp-Gevrey}.  Then for every selected
set \(\mathcal S\) of exact remainder factors and every BKAR covariance,
\begin{equation}
 \mathbb E_{\Sigma^F}
 \prod_{A\in\mathcal S}|R_{\epsilon,A}(Y_A)|
 \le\prod_{A\in\mathcal S}\|R_{\epsilon,A}\|_{L^{p_L}}
 \le e^{-c'|\mathcal S|/\epsilon}
 \label{eq:v6v8-remainder-product}
\end{equation}
for sufficiently small \(\epsilon\).  Hence the exact nonpolynomial compact
remainder has an all-arity polymer activity exponentially small in
\(1/\epsilon\).
\end{theorem}

\begin{proof}
Lemma~\ref{lem:v6v8-sparse-covariance} gives
\(\Sigma^F\le p_LI\).  Apply the Gaussian Brascamp--Lieb product inequality
to the absolute remainders.  At optimal truncation,
\eqref{eq:v6v8-Lp-remainder} is at most \(e^{-c'/\epsilon}\), after reducing
\(c'\).  Multiply the factors.  Square-free owner selection ensures that
no anchor is repeated.
\end{proof}

\begin{corollary}[Complete compact-gauge adjacent activity]
\label{cor:v6v8-complete-compact-activity}
Under the polynomial hypotheses of V7 and the exact remainder hypotheses
above, the complete compact-gauge owner contributes
\begin{equation}
 \eta_n^{\rm compact}le C\epsilon_n+Ce^{-c/\epsilon_n}
 \label{eq:v6v8-compact-eta}
\end{equation}
to the V2 all-order adjacent transport parameter.  This contribution is
summable whenever \(\sum_n\epsilon_n<\infty\).
\end{corollary}

\begin{proof}
V7 gives \(C\epsilon_n\) for the polynomial part.  Theorem
\ref{thm:v6v8-remainder-card} supplies an exponentially small polymer
activity for the exact remainder.  Sum over the fixed number of colors.
\end{proof}

\begin{firewall}[The color split must preserve the exact density]
\label{fw:v6v8-color-exactness}
Color substeps may reorder scalar positive owner factors because they
commute.  Gauge constraints, retained Grassmann factors, or matrix-valued
transfer operators cannot be reordered without an exact factorization.
Those objects are outside the present positive scalar density card.
\end{firewall}

\begin{firewall}[Unbounded Higgs owners need a different \(L^p\) proof]
\label{fw:v6v8-Higgs-unbounded}
Compactness of the gauge chart and inequality
\eqref{eq:v6v8-action-coercivity} control the exact gauge density.  The
Higgs field is unbounded.  Its quartic confinement can provide stronger
tails, but the coupled gauge--Higgs density and its normalized
\(L^{p_L}\) remainder require a separate proof.
\end{firewall}

\begin{firewall}[Determinant zeros are not compact-density tails]
\label{fw:v6v8-determinant-zero}
The high determinant modulus is analytic only on the V100 high-gap chart.
Modes approaching a zero must first be transferred to the retained
Grassmann block.  Sparse coloring cannot replace that spectral crossing
operation.
\end{firewall}

\begin{firewall}[Massive correlation length must match the block lattice]
\label{fw:v6v8-correlation-length}
The fixed color spacing \(L\) is uniform only after the Gaussian reference
is expressed on blocks comparable to its massive correlation scale.  If
the correlation length diverges in block units, the number of colors can
grow with the cutoff and destroy summability.
\end{firewall}

\begin{theorem}[Sixth-series V8 result]
\label{thm:v6v8-status}
A coercive compact link density has a uniform \(L^p\) interval
\(1<p<p_*\).  Sparse coloring of a massive block lattice makes the
replicated covariance Brascamp--Lieb exponent \(p_L\) lie inside that
interval.  The integrated Gevrey expansion then has an exponentially small
\(L^{p_L}\) remainder, and square-free owner selection multiplies those
remainders at arbitrary arity.  Together with V7, the complete compact-gauge
adjacent activity is \(O(\epsilon_n)+O(e^{-c/\epsilon_n})\).  Unbounded
Higgs owners, high-determinant gap charts, and uniformity of the block
correlation scale remain open.
\end{theorem}

\begin{proof}
Use Theorem~\ref{thm:v6v8-Lp-density},
Lemma~\ref{lem:v6v8-sparse-covariance},
Corollary~\ref{cor:v6v8-compatible-p},
Theorems \ref{thm:v6v8-Lp-Gevrey} and
\ref{thm:v6v8-remainder-card}, and
Corollary~\ref{cor:v6v8-complete-compact-activity}, subject to the four
firewalls.
\end{proof}

V9 will treat the unbounded Higgs owner.  It will use the positive quartic
radial potential to derive a normalized \(L^p\) density interval and a
Gevrey remainder uniform under sparse-color Gaussian interpolation, while
keeping the critical or tunnelling charts outside the massive-safe sector.

\clearpage
\section{V9: exact transport of a stable Higgs quartic}
\label{sec:v6v9}

The compact-gauge owner in V8 is controlled by a bounded chart and an
\(L^p\) coercivity interval.  The Higgs field is unbounded, but a positive
quartic tail supplies stronger exponential moments.  Rather than expand a
fixed physical quartic coupling as one small perturbation, this note
transports along the quartic coupling in finitely many exact normalized
increments.

\subsection{Normalized ratios along a quartic path}

Let \(\nu_0\) be a positive local reference measure on
\(\mathbb R^q\) with a nondegenerate Gaussian tail, and let
\begin{equation}
 P(\phi)=(|\phi|^2-v^2)^2\ge0.
 \label{eq:v6v9-Higgs-polynomial}
\end{equation}
For \(\lambda\ge0\), define
\begin{equation}
 d\nu_\lambda=Z(\lambda)^{-1}e^{-\lambda P(\phi)}d\nu_0,
 \qquad
 Z(\lambda)=\mathbb E_{\nu_0}e^{-\lambda P}.
 \label{eq:v6v9-quartic-family}
\end{equation}
For an increment \(\delta\) with \(\lambda+\delta\ge0\), its normalized
density relative to \(\nu_\lambda\) is
\begin{equation}
 h_{\lambda,\delta}(\phi)
 =\frac{Z(\lambda)}{Z(\lambda+\delta)}e^{-\delta P(\phi)}.
 \label{eq:v6v9-Higgs-ratio}
\end{equation}

\begin{lemma}[Exact \(L^p\) norm]
\label{lem:v6v9-exact-Lp}
For \(p>1\), if \(\lambda+p\delta\ge0\), then
\begin{equation}
 \|h_{\lambda,\delta}\|_{L^p(\nu_\lambda)}^p
 =\frac{Z(\lambda+p\delta)Z(\lambda)^{p-1}}
 {Z(\lambda+\delta)^p}.
 \label{eq:v6v9-exact-Lp}
\end{equation}
In particular, on every compact path
\(0\le\lambda\le\lambda_+\) and for
\(|\delta|\le\delta_0\) with
\(\lambda+p\delta\ge0\), this norm is finite and locally uniform.
\end{lemma}

\begin{proof}
Raise \eqref{eq:v6v9-Higgs-ratio} to the \(p\)-th power and integrate
against \eqref{eq:v6v9-quartic-family}.  The exponent becomes
\(-(lambda+p\delta)P\), giving the ratio of partition functions in
\eqref{eq:v6v9-exact-Lp}.  Positivity and continuity of \(Z\) on a compact
nonnegative interval give the uniform statement.
\end{proof}

The condition for a negative increment is real: an \(L^p\) comparison
spends \(p\delta\) units of quartic confinement, not merely \(\delta\).

\subsection{Smallness of the exact normalized increment}

Let \(h_t=h_{\lambda,t}\).  Differentiating normalization gives
\begin{equation}
 \partial_t h_t
 =-\bigl(P-\mathbb E_{\nu_{\lambda+t}}P\bigr)h_t.
 \label{eq:v6v9-density-derivative}
\end{equation}

\begin{theorem}[Exact Higgs increment in \(L^p\)]
\label{thm:v6v9-Higgs-Lp-increment}
Fix \(p>1\), \(\lambda_+<\infty\), and a stability margin
\(\lambda+t\ge0\), \(\lambda+pt\ge0\) along the increment.  Then
\begin{equation}
 \|h_{\lambda,\delta}-1\|_{L^p(\nu_\lambda)}
 \le C_{p,\lambda_+,\delta_0}|\delta|
 \label{eq:v6v9-Higgs-Lp-small}
\end{equation}
uniformly for \(0\le\lambda\le\lambda_+\) and
\(|\delta|\le\delta_0\) satisfying the margins.
\end{theorem}

\begin{proof}
By the fundamental theorem of calculus,
\(h_{\lambda,\delta}-1=\int_0^\delta\partial_th_tdt\).  Equation
\eqref{eq:v6v9-density-derivative} and H\"older give
\begin{align}
 \|\partial_th_t\|_{L^p(\nu_\lambda)}^p
 &\le C_p\left[
 \mathbb E_{\nu_\lambda}P^p h_t^p
 +\bigl(\mathbb E_{\nu_{\lambda+t}}P\bigr)^p
 \mathbb E_{\nu_\lambda}h_t^p\right].
 \label{eq:v6v9-dh-Lp}
\end{align}
By Lemma~\ref{lem:v6v9-exact-Lp}, these terms are ratios of partition
functions and moments of \(P\) under
\(\nu_{\lambda+pt}\).  The stability margins and the Gaussian tail of
\(\nu_0\) make them uniformly finite on the compact parameter interval.
Integrate the resulting uniform bound over \(t\).
\end{proof}

\subsection{Complete field-derivative activity}

For the V2 compiler, one needs field derivatives as well as the density
itself.  Since \(P\) has degree four, repeated differentiation of
\(e^{-tP}\) has controlled polynomial coefficients.

\begin{lemma}[Quartic exponential derivative]
\label{lem:v6v9-quartic-derivatives}
For every \(d\ge1\),
\begin{equation}
 \nabla^de^{-tP}=e^{-tP}\mathcal P_{d,t}(\phi),
 \label{eq:v6v9-Faa-form}
\end{equation}
where \(\mathcal P_{d,t}\) is a tensor polynomial of degree at most \(3d\)
and every monomial contains at least one factor of \(t\).  For
\(|t|\le\delta_0\),
\begin{equation}
 \|\mathcal P_{d,t}(\phi)\|_{\rm row}
 \le C_0^dd!\sum_{b=\lceil d/4\rceil}^{d}
 \frac{|t|^b}{b!}(1+|\phi|)^{4b-d}.
 \label{eq:v6v9-Faa-bound}
\end{equation}
after enlarging \(C_0\).
\end{lemma}

\begin{proof}
Fa\`a di Bruno's formula expresses the derivative as a sum over set
partitions of the \(d\) derivative labels.  Each block differentiates
\(-tP\), and derivatives of \(P\) vanish above order four.  Thus every term
with \(b=|\pi|\) blocks contains \(t^b\), has degree at most \(4b-d\),
and requires \(b\ge\lceil d/4\rceil\).  The exponential formula for set
partitions, including the block factorials, bounds the sum of coefficients
at fixed \(b\) by \(C_0^dd!/b!\).  Summing over \(b\) gives
\eqref{eq:v6v9-Faa-bound}.  Since \(d\ge1\), every admissible \(b\) is at
least one, so every term contains a factor of \(t\).
\end{proof}

\begin{theorem}[Exact Higgs derivative activity]
\label{thm:v6v9-Higgs-activity}
Under the stability margins of
Theorem~\ref{thm:v6v9-Higgs-Lp-increment}, and assuming
\(\lambda+p\delta\ge\lambda_*>0\), there are \(L,C<\infty\) such
that
\begin{equation}
 \sup_{d\ge0}\frac1{d!L^d}
 \|\nabla^d(h_{\lambda,\delta}-1)\|_{L^p(\nu_\lambda);\,\rm row}
 \le C|\delta|.
 \label{eq:v6v9-Higgs-activity}
\end{equation}
The constants are uniform on a compact stable coupling path.
\end{theorem}

\begin{proof}
The case \(d=0\) is
Theorem~\ref{thm:v6v9-Higgs-Lp-increment}.  For \(d\ge1\), normalization is
independent of \(\phi\), so differentiate
\eqref{eq:v6v9-Higgs-ratio}.  Lemma
\ref{lem:v6v9-quartic-derivatives} gives a factor \(|\delta|d!C_0^d\)
times a polynomial moment under a density with quartic coefficient
\(\lambda+p\delta\).  At the endpoint where this coefficient can be zero,
the Gaussian tail of \(\nu_0\) still has every polynomial moment; however
the order-\(3d\) moment grows faster than a fixed \(L^d\).  To obtain the
displayed uniform-in-\(d\) bound, retain a strict quartic margin
\begin{equation}
 \lambda+p\delta\ge\lambda_*>0.
 \label{eq:v6v9-strict-margin}
\end{equation}
Then the \(L^p\) norm of \((1+|\phi|)^{4b-d}\) is bounded by
\(C^d b!\): this follows from the gamma integral after the substitution
\(u=\lambda_*r^4\) and Stirling's bound.  In
\eqref{eq:v6v9-Faa-bound} that moment cancels the denominator \(b!\).
Because \(b\ge1\) and \(|\delta|\le\delta_0\), the remaining finite sum is
bounded by \(|\delta|C_1^dd!\).  Divide by \(d!L^d\) and choose
\(L>C_1\).
\end{proof}

The strict margin is needed for the complete all-derivative activity.  The
weaker \(L^p\) density estimate remains true at \(\lambda=0\), but a single
fixed derivative scale does not follow from Gaussian moments of degrees
growing with \(d\) by this proof.

\subsection{Finite-coupling transport without small physical \(\lambda\)}

Let \([\lambda_-,\lambda_+]\) be a compact interval with
\(\lambda_->0\).  Partition a path of total variation
\(V_\lambda\) into increments \(\delta_k\) small enough that the strict
margin and the one-step activity threshold hold.

\begin{theorem}[Finite quartic path loss]
\label{thm:v6v9-finite-path}
Assume the current positive measure carries the V2 mixed insertion tree
card uniformly for \(\lambda\in[\lambda_-,\lambda_+]\), and the owner
supports obey the bounded-occupancy spatial sum.  Then the exact Higgs
increments transport the all-order rooted tree card along the entire path,
with cumulative loss bounded by
\begin{equation}
 \prod_k(1-C|\delta_k|)^{-1}
 \le e^{2CV_\lambda}
 \label{eq:v6v9-path-loss}
\end{equation}
when \(C|\delta_k|\le1/2\).  The endpoint quartic coupling need not be
small.
\end{theorem}

\begin{proof}
Theorem~\ref{thm:v6v9-Higgs-activity} gives one-step interaction activity
\(\eta_k\le C|\delta_k|\).  Apply the V1--V2 transport theorem at every
substep.  Since \(-\log(1-x)\le2x\) for \(x\le1/2\),
\[
 \log\prod_k(1-C|\delta_k|)^{-1}
 \le2C\sum_k|\delta_k|=2CV_\lambda.
\]
\end{proof}

For an ultraviolet cutoff chain, an additional running increment
\(\delta\lambda_n\) contributes a summable activity precisely when its
absolute variation is summable, or when its nonsummable part is absorbed
into the chosen reference trajectory.

\subsection{Broken-phase and coupled-field firewalls}

\begin{firewall}[Positive quartic tail is indispensable]
\label{fw:v6v9-positive-quartic}
If the running effective quartic coefficient crosses zero, the strict
margin \eqref{eq:v6v9-strict-margin} fails and the local all-derivative
activity theorem no longer applies.  A higher stabilizing operator, a new
phase, or a different ultraviolet trajectory must then be included
explicitly.  Analytic continuation of \(\lambda\) does not define a
positive normalized measure.
\end{firewall}

\begin{firewall}[The broken potential is not globally convex]
\label{fw:v6v9-broken-nonconvex}
Although \((|\phi|^2-v^2)^2\) has a positive tail, its Hessian is not
uniformly positive near every field.  The V90 massive one-form inverse
holds only on a selected broken-phase chart with gauge/Goldstone directions
treated correctly.  Tunnelling and chart changes require a phase/polymer
decomposition.
\end{firewall}

\begin{firewall}[Gauge and determinant couplings alter the conditional
tail]
\label{fw:v6v9-coupled-tail}
The exact formulas above concern a positive quartic family relative to
\(\nu_0\).  In the coupled measure, covariant gradients are positive but
the determinant modulus and counterterms change the conditional density.
The centered-log comparison must prove that a uniform quartic margin and
moment bound survive conditioning.
\end{firewall}

\begin{firewall}[The initial zero-to-positive path needs a separate first
step]
\label{fw:v6v9-initial-step}
The complete derivative activity uses \(\lambda_*>0\).  Transport from a
pure Gaussian reference at \(\lambda=0\) to the first positive quartic
reference is not covered uniformly in all derivative orders by
Theorem~\ref{thm:v6v9-Higgs-activity}.  It may be handled by a finite
low-arity construction plus a positive-quartic restart, or by a stronger
Gaussian-polynomial compiler; neither is proved here.
\end{firewall}

\begin{theorem}[Sixth-series V9 result]
\label{thm:v6v9-status}
The normalized density ratio between two stable Higgs quartic couplings has
an exact \(L^p\) norm and differs from one by \(O(|\delta\lambda|)\).  Under
a strict positive quartic margin, every field-derivative row has the same
linear activity bound.  Subdividing a finite positive coupling path
therefore transports an existing all-order tree card with loss at most
\(e^{CV_\lambda}\), so the physical quartic itself need not be small.  A
zero or negative ultraviolet quartic, global broken-phase nonconvexity,
coupled conditional tails, and the initial Gaussian-to-quartic step remain
open.
\end{theorem}

\begin{proof}
Use Lemma~\ref{lem:v6v9-exact-Lp},
Theorems \ref{thm:v6v9-Higgs-Lp-increment},
\ref{thm:v6v9-Higgs-activity}, and
\ref{thm:v6v9-finite-path}, subject to the four firewalls.
\end{proof}

V10 is the next compulsory companion audit.  It will then attack the
remaining high-determinant-modulus owner, using the V100 crossing atlas and
singular gap to derive an exact normalized \(L^p\) increment and complete
derivative activity without multiplying by the global fermion rank.

\clearpage
\section{V10: companion audit and rooted high-determinant owners}
\label{sec:v6v10}

\subsection{Compulsory companion audit}

The active Yang--Mills file was
\texttt{Renewal\_Yang\_Mills\_Mass\_Gap\_Research\_Notes\_V9.tex}, with
\(134\,582\) bytes and SHA--256
\begin{center}
\texttt{386CF89D2B4F9F60EA2CF8D480D6D1D105547B27880663EEFC337953B50BF8B7}.
\end{center}
It proves an all-order all-thermal selected \(SU(2)\) gate uniform in
thermal anisotropy.  The proof uses a continuous weighted AM--GM load
allocation and an exact two-regime Wilson tail, avoiding scale-bin counts.
Its remaining selected local sector contains a thermally hot row coupled to
a hard private bank.  It does not prove the full Yang--Mills mass gap.

The Navier--Stokes file remained
\texttt{Renewal\_NS\_Stability\_Research\_Notes\_V31.tex}, with
\(887\,537\) bytes and SHA--256
\begin{center}
\texttt{F1121B62238AD5491BB7CF03635EA9934942EDF573ED8FAAA9EE79E1D38A6696}.
\end{center}
Its one-use ancestry and signed-correlation bottleneck are unchanged.

For the determinant-modulus owner, the companion lesson is to preserve
rowwise spectral information.  Bounding a high determinant by its global
rank loses volume, just as binning every private Yang--Mills scale would
lose anisotropy.  The construction below assigns one trace-log density to
each spatial root and uses a weighted kernel algebra.

\subsection{Locality of the high spectral bundle}

Let \(H=D^*D\) be a finite-range or exponentially local positive operator.
On a V100 chart, its low/high threshold is separated from the spectrum by
\(\delta>0\), and the high projector is
\begin{equation}
 Q=\frac1{2\pi i}\oint_{\Gamma_H}(z-H)^{-1}\,dz.
 \label{eq:v6v10-high-projector}
\end{equation}

\begin{lemma}[Exponential locality of a gapped spectral projector]
\label{lem:v6v10-projector-locality}
Assume the local row and column norms of \(H\) are volume-uniform and
\(\operatorname{dist}(\Gamma_H,\sigma(H))\ge\delta\).  Then for sufficiently
small \(\mu>0\),
\begin{equation}
 \sup_x\sum_y e^{\mu d(x,y)}\|Q_{xy}\|+
 \sup_y\sum_x e^{\mu d(x,y)}\|Q_{xy}\|
 \le C_{\delta,\Gamma_H}.
 \label{eq:v6v10-Q-local}
\end{equation}
The same holds for a finite number of field derivatives of \(Q\), with
the corresponding powers of the gap distance.
\end{lemma}

\begin{proof}
For every \(z\in\Gamma_H\), conjugate \(z-H\) by the exponential distance
weight used in the frozen Combes--Thomas theorem.  The conjugation error is
smaller than \(\delta/2\) when \(\mu\) is small, so
\((z-H)^{-1}\) has a uniform exponentially weighted Schur norm.  Integrate
that estimate around \(\Gamma_H\) to obtain
\eqref{eq:v6v10-Q-local}.  Differentiating the resolvent inserts factors
\((z-H)^{-1}(dH)(z-H)^{-1}\); the weighted kernel algebra and the same
contour bound prove the derivative statement.
\end{proof}

Thus the high map and its inverse can be embedded into the original local
site fibres without replacing their spatial trace by the dimension of the
entire high space.

\subsection{The exact adjacent high determinant ratio}

Use the crossing atlas to identify the high bundles at two adjacent cutoff
steps after every near-crossing cluster has been transferred to the
retained sector.  Let \(E_n\) and \(E_{n+1}\) be the identified high maps
and define
\begin{equation}
 R_n:=E_n^{-1}(E_{n+1}-E_n),
 \qquad
 E_{n+1}=E_n(I+R_n).
 \label{eq:v6v10-relative-high-defect}
\end{equation}
Assume
\begin{equation}
 \|R_n\|_{\mu,1}\le\kappa_n<1.
 \label{eq:v6v10-R-small}
\end{equation}
Define the rooted log density
\begin{equation}
 \ell_{n,x}:=operatorname{Re}
 \sum_{m\ge1}\frac{(-1)^{m+1}}m
 \operatorname{tr}_{\rm int}(R_n^m)_{xx}.
 \label{eq:v6v10-rooted-log}
\end{equation}

\begin{theorem}[Exact rooted determinant factorization]
\label{thm:v6v10-rooted-determinant}
The series \eqref{eq:v6v10-rooted-log} converges absolutely and
\begin{align}
 \log\frac{|\det E_{n+1}|}{|\det E_n|}
 &=\sum_x\ell_{n,x},
 \label{eq:v6v10-log-sum}\\
 \frac{|\det E_{n+1}|}{|\det E_n|}
 &=\prod_x e^{\ell_{n,x}}.
 \label{eq:v6v10-owner-product}
\end{align}
Moreover, if \(d_*\) is the fixed internal fibre dimension,
\begin{equation}
 \sup_x|\ell_{n,x}|
 \le d_*[-\log(1-\kappa_n)].
 \label{eq:v6v10-root-size}
\end{equation}
No global high-space rank occurs.
\end{theorem}

\begin{proof}
The weighted Schur norm is a Banach algebra, so
\(\|R_n^m\|_{\mu,1}\le\kappa_n^m\).  The diagonal internal trace at one
root is bounded by \(d_*\kappa_n^m\), proving absolute convergence and
\eqref{eq:v6v10-root-size}.  Since \(\|R_n\|<1\), the finite-dimensional
identity
\begin{equation}
 \log\det(I+R_n)=operatorname{Tr}log(I+R_n)
 =\sum_{m\ge1}\frac{(-1)^{m+1}}m\operatorname{Tr}R_n^m
 \label{eq:v6v10-trace-log}
\end{equation}
holds on the principal analytic branch.  Split the total trace into its
site-diagonal internal traces and take real parts to obtain
\eqref{eq:v6v10-log-sum}.  Exponentiation proves
\eqref{eq:v6v10-owner-product}.
\end{proof}

Each \(e^{\ell_{n,x}}\) is a positive scalar owner.  Gauge covariance of
\(R_{n,xy}\) by endpoint conjugation makes every diagonal trace, and hence
every owner, gauge invariant.

\subsection{Complete derivative rows of a rooted trace logarithm}

Let \(\mathfrak A_\mu\) denote the weighted kernel Banach algebra.  Suppose
the adjacent defect is analytic in the complex field ball
\(\|z\|<r_0\) and obeys
\begin{equation}
 \sup_{\|z\|<r_0}\|R_n(z)\|_{\mu,1}le C_R\epsilon_n,
 \qquad C_R\epsilon_n\le\kappa<1.
 \label{eq:v6v10-analytic-R}
\end{equation}

\begin{lemma}[Banach analytic trace-log activity]
\label{lem:v6v10-log-activity}
There are \(L,C<\infty\), independent of volume and \(n\), such that
\begin{equation}
 \sup_{d\ge0}\frac1{d!L^d}
 \sup_x\|\nabla^d\ell_{n,x}\|_{\rm row}
 \le C\epsilon_n.
 \label{eq:v6v10-ell-activity}
\end{equation}
The same bound holds for the positive owner increment
\begin{equation}
 K_{n,x}:=e^{\ell_{n,x}}-1.
 \label{eq:v6v10-K-owner}
\end{equation}
\end{lemma}

\begin{proof}
The analytic map
\begin{equation}
 \mathcal F(R)=\sum_{m\ge1}\frac{(-1)^{m+1}}mR^m
 \label{eq:v6v10-Banach-log}
\end{equation}
maps the ball \(\|R\|_{\mu,1}\le\kappa\) into
\(\mathfrak A_\mu\), with norm at most \(-\log(1-\kappa)\).  More sharply,
\(\|\mathcal F(R_n(z))\|_{\mu,1}\le C\epsilon_n\) under
\eqref{eq:v6v10-analytic-R}, because every term contains at least one
factor of \(R_n\).  Cauchy's estimate on a smaller field ball gives
\begin{equation}
 \|\nabla^d\mathcal F(R_n)\|_{\mu,1;\,\rm row}
 \le C\epsilon_n d!L^d.
 \label{eq:v6v10-Cauchy-log}
\end{equation}
Taking one diagonal internal trace is a bounded functional with norm
\(d_*\), proving \eqref{eq:v6v10-ell-activity}.  Composition with the
scalar exponential preserves the same factorial-exponential derivative
class, and \(e^u-1\) has no constant term, so the factor \(\epsilon_n\)
remains.
\end{proof}

\begin{theorem}[High-modulus owner card]
\label{thm:v6v10-high-owner-card}
Assume the high spectral projector and inverse are exponentially local, the
crossing atlas has removed all modes below the fixed gap, and
\eqref{eq:v6v10-analytic-R} holds for the matched adjacent high defect.
Then the entire high determinant-modulus ratio has the exact bounded-
occupancy owner representation \eqref{eq:v6v10-owner-product}, with one
owner per spatial root and local derivative activity
\begin{equation}
 z(K_{n,x})\le C\epsilon_n.
 \label{eq:v6v10-high-zeta}
\end{equation}
If its exponentially local tails satisfy the V2 support-tree sum, it
contributes \(\eta_n^{\rm high\,mod}\le C\epsilon_n\) to the all-order
adjacent transport parameter.
\end{theorem}

\begin{proof}
Theorem~\ref{thm:v6v10-rooted-determinant} gives the exact product and one
root per site.  Lemma~\ref{lem:v6v10-log-activity} gives the derivative
activity.  The exponential weighted norm assigns a summable spatial tail
to each root.  Apply the V2 rooted interaction-tree lemma.
\end{proof}

The infinite trace-loop family is now grouped canonically into one rooted
analytic owner.  This supplies the grouping required by the V7
bounded-occupancy card.

\subsection{Rowwise gaps and the companion load lesson}

The preceding theorem uses one uniform high threshold.  A refinement can
retain nonuniform local singular margins.  Let \(w_x>0\) and define the
similarity-weighted norm
\begin{equation}
 \|R\|_{\mu,1;w}:=
 \max\left\{
 \sup_x\sum_y e^{\mu d(x,y)}\frac{w_y}{w_x}\|R_{xy}\|,
 \sup_y\sum_x e^{\mu d(x,y)}\frac{w_x}{w_y}\|R_{xy}\|
 \right\}.
 \label{eq:v6v10-row-weighted-norm}
\end{equation}

\begin{proposition}[Continuous row-load invariance]
\label{prop:v6v10-row-load}
All trace-log and derivative-activity conclusions above remain valid with
\(\|\cdot\|_{\mu,1}\) replaced by
\(\|\cdot\|_{\mu,1;w}\), provided the weighted defect norm is below one.
No count of the number of distinct weights appears.
\end{proposition}

\begin{proof}
The weighted norm is the ordinary Schur norm after conjugation by the
diagonal multiplication operator \(W:f_x\mapsto w_xf_x\).  It is
submultiplicative, and diagonal traces are invariant under this similarity.
Repeat the Banach-algebra proofs.  Since weights enter continuously, no
dyadic partition or bin count is introduced.
\end{proof}

This is the determinant analogue of the audited continuous row-load
allocation.  Constructing useful weights from the actual local high
singular margins remains a separate estimate; the proposition only shows
that a scale-bin loss is unnecessary.

\subsection{Normalization and crossing firewalls}

\begin{firewall}[The global normalized ratio is not an \(L^p\) owner]
\label{fw:v6v10-global-Lp}
Even when every \(K_{n,x}=O(\epsilon_n)\), the global product can have an
\(L^p\) norm exponential in volume.  Normalization must be handled by the
rooted connected/polymer expansion, where vacuum components cancel.  A
global determinant \(L^p\) bound would reintroduce the rank loss removed by
Theorem~\ref{thm:v6v10-rooted-determinant}.
\end{firewall}

\begin{firewall}[The analytic high gap is required on a complex ball]
\label{fw:v6v10-complex-gap}
A real singular gap does not automatically imply
\eqref{eq:v6v10-analytic-R} on the complex field ball used by Cauchy's
estimate.  One needs a pseudospectral resolvent bound for the complexified
high operator, or a real-variable derivative proof with the same
factorial-exponential constants.
\end{firewall}

\begin{firewall}[Crossing modes must move before the logarithm]
\label{fw:v6v10-crossing-first}
If a high singular value approaches zero, \(E_n^{-1}\) and \(R_n\) cease to
obey the hypotheses.  The V100 crossing transfer must move that cluster to
the retained Grassmann sector before defining
\eqref{eq:v6v10-relative-high-defect}.  The trace logarithm cannot be used
through the crossing.
\end{firewall}

\begin{firewall}[The phase is a determinant-line section]
\label{fw:v6v10-phase-line}
The rooted owners use the real part of the trace logarithm and reconstruct
the high modulus.  The imaginary part has chart transitions governed by
the determinant line and must be combined with the low phase, measure, and
anomaly counterterms.  It is not a positive owner.
\end{firewall}

\begin{firewall}[Matched smallness is not supplied by the high gap]
\label{fw:v6v10-matched-smallness}
The gap makes \(E_n^{-1}\) bounded; it does not prove that the adjacent
relative defect is \(O(\epsilon_n)\).  Relevant heat coefficients,
marginal logarithms, and the first local Born terms must be assigned to
their running couplings and counterterms before
\eqref{eq:v6v10-analytic-R} can hold.
\end{firewall}

\begin{theorem}[Sixth-series V10 result]
\label{thm:v6v10-status}
On a gapped covariant spectral chart, the high projector and inverse are
exponentially local.  A norm-small matched adjacent high defect has an exact
determinant-modulus factorization into one gauge-invariant rooted trace-log
owner per spatial block.  Every owner and all its derivative rows have
activity \(O(\epsilon_n)\), with the fixed internal fibre dimension rather
than the global high rank.  Similarity row weights preserve the proof
without scale-bin counts.  Complex-gap control, crossing transfer,
determinant-line phase compatibility, and matched adjacent smallness remain
required.
\end{theorem}

\begin{proof}
Use Lemma~\ref{lem:v6v10-projector-locality},
Theorem~\ref{thm:v6v10-rooted-determinant},
Lemma~\ref{lem:v6v10-log-activity},
Theorem~\ref{thm:v6v10-high-owner-card}, and
Proposition~\ref{prop:v6v10-row-load}, subject to the five firewalls.
\end{proof}

V11 will assemble the compact-gauge, stable-Higgs, PV, and rooted
high-modulus owner cards into one positive massive adjacent-density theorem.
It will state the exact common smallness and sparse-color conditions and
identify which remaining term---bad blocks, relevant flow, or complex high
gap---first prevents an unconditional all-scale SIIC.

\clearpage
\section{V11: the assembled positive massive adjacent density}
\label{sec:v6v11}

V8--V10 provide separate owner estimates.  This note assembles them into a
single all-order transport theorem and identifies the remaining inputs by
their actual logical order.  All factors in this section are real and
positive.  Determinant phases, possible ghost signs, retained light
Grassmann factors, photons, and gravitational constraints remain outside
the positive normalization.

\subsection{The exact positive owner product}

At cutoff step \(n\), after matching relevant local coefficients and
transferring every near-crossing fermion mode to the retained block, write
the positive massive density ratio as
\begin{equation}
 \mathcal H_n=
 \prod_{A\in\mathcal P_n^{\rm cg}}h_{n,A}^{\rm cg}
 \prod_{A\in\mathcal P_n^{\rm H}}h_{n,A}^{\rm H}
 \prod_{x}h_{n,x}^{\rm PV}
 \prod_{x}h_{n,x}^{\rm F,hi}
 \prod_{A\in\mathcal P_n^{\rm loc}}h_{n,A}^{\rm loc}.
 \label{eq:v6v11-positive-product}
\end{equation}
The superscripts denote compact gauge, Higgs, Pauli--Villars, high fermion
modulus, and remaining real local matched owners.  Set \(K=h-1\) for every
labelled factor.  Group each infinite trace-loop family into the single
rooted analytic owner constructed in V10.

\begin{definition}[Positive massive step data]
\label{def:v6v11-step-data}
A cutoff step has positive massive data
\((M_0,N_L,p_L,\mathbf z_n)\) if:
\begin{enumerate}
\item the owner presentation has at most \(M_0\) grouped labels per anchor;
\item a fixed \(N_L\)-color decomposition gives same-color anchor spacing
      \(L\) and Gaussian product exponent \(p_L\);
\item the one-form covariance has the massive exponential row bound of the
      frozen V90 theorem in the block units used by the coloring; and
\item the complete derivative activities are bounded by
\begin{equation}
 \mathbf z_n=(z_n^{\rm cg},z_n^{\rm H},z_n^{\rm PV},
 z_n^{\rm F,hi},z_n^{\rm loc},z_n^{\rm bad}).
 \label{eq:v6v11-z-vector}
\end{equation}
\end{enumerate}
The bad component is a conditionally dominated polymer activity, not merely
a one-block probability.
\end{definition}

\subsection{Pauli--Villars owner}

The PV factor can be placed beside the V10 high modulus when its regulated
operator is positive.  Let \(P_n>0\), with
\(P_{n+1}=P_n(I+R_n^{\rm PV})\) and
\(\|R_n^{\rm PV}\|_{\mu,1}<1\).

\begin{lemma}[Rooted PV activity]
\label{lem:v6v11-PV-activity}
Suppose the matched relative PV defect is analytic on a fixed complex field
ball and
\begin{equation}
 \sup_{\|z\|<r_0}\|R_n^{\rm PV}(z)\|_{\mu,1}
 \le C\epsilon_n^{\rm PV}.
 \label{eq:v6v11-PV-defect}
\end{equation}
For any fixed real determinant power \(\sigma_{\rm PV}\), the ratio
\((\det P_{n+1}/\det P_n)^{\sigma_{\rm PV}}\) factors exactly into one
positive rooted owner per site, and
\begin{equation}
 z_n^{\rm PV}\le C_{\sigma_{\rm PV}}\epsilon_n^{\rm PV}.
 \label{eq:v6v11-PV-z}
\end{equation}
\end{lemma}

\begin{proof}
Apply the rooted trace-log construction of V10 to
\(R_n^{\rm PV}\) and multiply every rooted logarithm by
\(\sigma_{\rm PV}\).  Positivity of \(P_n\) makes the real logarithm
single-valued.  The Banach analytic derivative proof is unchanged.
\end{proof}

\subsection{One-step assembly}

Let
\begin{equation}
 Z_n:=z_n^{\rm cg}+z_n^{\rm H}+z_n^{\rm PV}
      +z_n^{\rm F,hi}+z_n^{\rm loc}+z_n^{\rm bad}.
 \label{eq:v6v11-total-z}
\end{equation}

\begin{theorem}[Positive massive adjacent SIIC]
\label{thm:v6v11-positive-SIIC}
Assume a step has the data of
Definition~\ref{def:v6v11-step-data}.  Assume further that:
\begin{enumerate}
\item \(p_L\) lies in every compact-gauge \(L^p\) interval used in V8;
\item the Higgs increments remain inside a strict positive quartic margin;
\item the PV and high-fermion operators obey their complex high-gap cards;
\item all local relevant and marginal coefficients have been assigned to
      the reference trajectory, leaving only the residual activity
      \(z_n^{\rm loc}\); and
\item each color substep obeys \(C_*Z_n<1\).
\end{enumerate}
Then the complete positive product
\eqref{eq:v6v11-positive-product} satisfies the summed
interaction-insertion card with
\begin{equation}
 \eta_n\le C_*N_L Z_n.
 \label{eq:v6v11-eta-total}
\end{equation}
It transports the all-order rooted tree card across the cutoff step with
loss at most
\begin{equation}
 (1-C_*N_LZ_n)^{-1}
 \label{eq:v6v11-one-step-loss}
\end{equation}
when the right side is defined, after subdividing a color step further if
necessary.
\end{theorem}

\begin{proof}
Expand every positive scalar owner as \(1+K\).  The compact-gauge
polynomial and exact remainder activities are V7--V8.  The exact Higgs
increment activity is V9.  The high fermion modulus is V10, and
Lemma~\ref{lem:v6v11-PV-activity} gives the PV term.  Real local polynomial
owners use the V7 derivative-row card.  Conditionally dominated bad
polymers enter with their declared activity.

Within one color, bounded occupancy and the massive covariance row give the
Gaussian Brascamp--Lieb product estimate and the rooted owner-support tree
sum.  Activities add by the triangle inequality, so the one-color
interaction parameter is at most \(C_*Z_n\).  Positive scalar factors
commute, hence the \(N_L\) color substeps reproduce the exact product.  Sum
their parameters to get \eqref{eq:v6v11-eta-total}, and apply the V1--V2
transport theorem.  If one substep is too large but lies on a finite stable
coupling path, subdivide it as in V9.
\end{proof}

\begin{corollary}[All-scale positive massive card]
\label{cor:v6v11-all-scale}
Suppose \(N_L,C_*\) are cutoff-uniform,
\(C_*N_LZ_n\le1/2\), and
\begin{equation}
 \sum_{n\ge0} Z_n<\infty.
 \label{eq:v6v11-summable-Z}
\end{equation}
Then the positive massive all-order transport loss is finite:
\begin{equation}
 \prod_n(1-C_*N_LZ_n)^{-1}
 \le\exp\left(2C_*N_L\sum_nZ_n\right).
 \label{eq:v6v11-all-scale-loss}
\end{equation}
The limiting compatible cylinder functional inherits the massive rooted
tree bound.
\end{corollary}

\begin{proof}
Use \(-\log(1-x)\le2x\) and the compactness/compatibility passage already
proved in V2.
\end{proof}

\subsection{Finite coupling paths and infinite cutoff tails}

Some owners need not have summable coupling values.  What enters the
adjacent transport is their matched increment.

\begin{proposition}[Variation plus irrelevant tail]
\label{prop:v6v11-variation-tail}
Suppose
\begin{equation}
 Z_n\le C\left(
 |g_{n+1}-g_n|+|\lambda_{n+1}-\lambda_n|
 +|c_{n+1}-c_n|+a_n^\varepsilon+b_n
 \right),
 \label{eq:v6v11-Z-decomposition}
\end{equation}
where \(g,\lambda,c\) range in compact stable coupling charts,
\(\varepsilon>0\), and \(b_n\) is the bad-polymer activity.  If the three
coupling trajectories have finite total variation and
\(\sum_nb_n<\infty\), then \eqref{eq:v6v11-summable-Z} holds.
\end{proposition}

\begin{proof}
Sum \eqref{eq:v6v11-Z-decomposition}.  Total variations and
\(\sum_nb_n\) are finite by hypothesis, while a positive power of a
geometric cutoff is summable.
\end{proof}

This proposition permits a nonzero limiting Higgs or gauge coupling.  It
does not permit a trajectory that leaves every compact stable chart or has
infinite variation.

\subsection{The first unproved inputs in logical order}

\begin{assessment}[Positive-sector dependency order]
\label{ass:v6v11-dependency-order}
At a fixed finite cutoff, the first unproved measure input is the uniform
conditional bad-block/product-domination estimate of the frozen V97--V98
cards.  On a good block, the first unproved analytic input is a complex
high-gap or equivalent real all-derivative estimate for the rooted high
fermion and PV owners.  After those local inputs, the first all-scale
structural obstruction is the relevant trajectory: the stated hypercharge
Gaussian-boundary flow reaches its Landau boundary, the Higgs quartic may
lose its positive margin, and a finite gravitational coupling list is not
closed.  Therefore the hypotheses of
Corollary~\ref{cor:v6v11-all-scale} have not been verified for the full
SM--Einstein programme.
\end{assessment}

\begin{firewall}[Faddeev--Popov sign]
\label{fw:v6v11-FP-sign}
A Faddeev--Popov determinant belongs to
\eqref{eq:v6v11-positive-product} only on a chart where its sign and zero
set are controlled, for example inside a proved positive Gribov region.
Otherwise its sign or phase must remain in the rooted observable sector.
Writing ``ghost owner'' does not make a signed determinant positive.
\end{firewall}

\begin{firewall}[Anomaly phase and low Grassmann block]
\label{fw:v6v11-phase-low}
The determinant modulus theorem does not control the determinant-line phase
or the retained light Schur factor.  Their Ward identities, spectral
crossing transitions, and connected rooted sums remain separate gates.
\end{firewall}

\begin{firewall}[Massless and conformal sectors]
\label{fw:v6v11-massless-conformal}
The photon and gravitational constraints enter only through their
Ward/multipole polynomial grades.  The unrotated Euclidean conformal mode
has no positive normalized owner.  None is included in
Theorem~\ref{thm:v6v11-positive-SIIC}.
\end{firewall}

\begin{firewall}[Finite variation is not established by perturbative beta
functions]
\label{fw:v6v11-finite-variation}
Proposition~\ref{prop:v6v11-variation-tail} is a sufficient criterion.
Perturbative running over a finite energy window does not prove an
all-scale nonperturbative trajectory with finite variation, positive
quartic margin, and uniform spectral charts.
\end{firewall}

\begin{firewall}[Reflection positivity is not stable under arbitrary
factor subdivision]
\label{fw:v6v11-RP-subdivision}
Positive scalar density factors commute in Euclidean expectation, but an
arbitrary subdivision need not correspond to a positive transfer operator
across a reflection plane.  The final assembled regulator must satisfy the
separate finite-cutoff reflection-positivity certificate.
\end{firewall}

\begin{theorem}[Sixth-series V11 result]
\label{thm:v6v11-status}
The compact-gauge, stable-Higgs, positive PV, rooted high-fermion-modulus,
real local, and conditionally dominated bad-block factors assemble into one
positive massive adjacent-density SIIC with parameter
\(\eta_n\le C N_LZ_n\).  A summable activity or finite-variation coupling
path gives a finite all-scale transport loss and preserves the massive
all-order rooted tree card.  The first missing local inputs are conditional
bad-block domination and complex high-gap derivative control; the first
all-scale structural obstruction is the unresolved hypercharge--Higgs--
gravity trajectory.  Phases, light Grassmann modes, massless fields, the
conformal direction, and reflection positivity remain outside this card.
\end{theorem}

\begin{proof}
Use Lemma~\ref{lem:v6v11-PV-activity},
Theorem~\ref{thm:v6v11-positive-SIIC},
Corollary~\ref{cor:v6v11-all-scale}, and
Proposition~\ref{prop:v6v11-variation-tail}, subject to the assessment and
five firewalls.
\end{proof}

V12 will attack the first missing local input: a conditional bad-block
estimate for the coupled positive modulus measure.  It will combine the
V98 centered-log comparison with the new rooted owner decomposition and
derive a finite-volume criterion whose constants are local rather than
proportional to the exterior volume.

\clearpage
\section{V12: local conditional comparison and bad-block domination}
\label{sec:v6v12}

V11 identifies uniform conditional bad-block domination as the first
missing measure input.  A global comparison can depend on the exterior
volume.  The rooted owner representation removes that loss: changing
fields in one block affects a distant rooted owner only exponentially.

\subsection{Exterior-volume cancellation}

Let \(B\) be a block and let \(\phi_B,\phi'_B\) have the same exterior.
Suppose a rooted log owner satisfies
\begin{equation}
 \sup_{e\subset B}\sup_\phi
 |\partial_{\phi_e}\ell_x(\phi)|
 \le z e^{-\mu d(x,e)}.
 \label{eq:v6v12-root-sensitivity}
\end{equation}
Assume the joining coordinate path has length at most \(L_0\) per field
coordinate.

\begin{lemma}[Local oscillation of an infinite rooted sum]
\label{lem:v6v12-local-oscillation}
On a block graph of polynomial volume growth,
\begin{equation}
 \left|\sum_x\ell_x(\phi_B,\phi_{B^c})
 -\sum_x\ell_x(\phi'_B,\phi_{B^c})\right|
 \le C_\mu L_0z|B|.
 \label{eq:v6v12-local-oscillation}
\end{equation}
The constant is independent of total volume and exterior data.
\end{lemma}

\begin{proof}
Integrate along the coordinate path and sum absolutely:
\[
 \sum_{e\subset B}\sum_xL_0ze^{-\mu d(x,e)}
 \le C_\mu L_0z|B|.
\]
The inner exponential row sum is finite uniformly by polynomial volume
growth.  Exterior-only owners have zero difference and cancel from the
normalized conditional density.
\end{proof}

\begin{proposition}[Conditional bounded-owner log]
\label{prop:v6v12-bounded-log}
Let \(L_B^{\rm bd}\) be the sum of compact-gauge, positive PV,
high-modulus, and bounded local log owners depending on \(B\).  If every
rooted family obeys \eqref{eq:v6v12-root-sensitivity} and there are
finitely many owner types, then
\begin{equation}
 \operatorname*{ess\,osc}_{\phi_B}L_B^{\rm bd}
 \le C_{\rm bd}|B|
 \label{eq:v6v12-bounded-oscillation}
\end{equation}
uniformly in the exterior volume.
\end{proposition}

\begin{proof}
Apply Lemma~\ref{lem:v6v12-local-oscillation} to each rooted family and add
the finitely many constants.  Strictly local owners are the special case
with zero sensitivity beyond a fixed neighbourhood.
\end{proof}

\subsection{The unbounded Higgs centered moment}

Let the conditional positive reference have a stable tail
\begin{equation}
 V_{H,B}(\phi_B\mid\phi_{B^c})
 \ge\lambda_*\sum_{x\in B}|\phi_x|^4-C_0|B|,
 \qquad\lambda_*>0,
 \label{eq:v6v12-Higgs-tail}
\end{equation}
and suppose
\begin{equation}
 |L_B^{\rm H}|
 \le a_H\sum_{x\in B}(1+|\phi_x|^4)+C_1|B|.
 \label{eq:v6v12-Higgs-log-growth}
\end{equation}

\begin{lemma}[Conditional centered Higgs moment]
\label{lem:v6v12-Higgs-mgf}
If \(qa_H<\lambda_*/2\), then
\begin{equation}
 \log\mathbb E_{\nu_B}
 e^{q(L_B^{\rm H}-\mathbb E_{\nu_B}L_B^{\rm H})}
 \le qC_H|B|.
 \label{eq:v6v12-Higgs-centered}
\end{equation}
The constant is exterior-uniform whenever
\eqref{eq:v6v12-Higgs-tail} and a local conditional partition lower bound
are exterior-uniform.
\end{lemma}

\begin{proof}
The \(q\)-th exponential spends at most \(qa_H\) of the quartic tail.
The remaining \(\lambda_*/2\) tail makes every one-site integral finite,
and dropping positive covariant-gradient terms gives an upper bound
\(e^{C|B|}\).  A fixed product neighbourhood of one admissible bounded
interior configuration gives a partition lower bound \(e^{-C'|B|}\).
The same quartic moment estimate bounds
\(|\mathbb E_{\nu_B}L_B^{\rm H}|\) by \(C''|B|\).  Combine the three
extensive estimates.
\end{proof}

\subsection{The full centered-log comparison}

Put \(L_B=L_B^{\rm bd}+L_B^{\rm H}\) and
\begin{equation}
 d\mu_B=\frac{e^{L_B}}{\mathbb E_{\nu_B}e^{L_B}}\,d\nu_B.
 \label{eq:v6v12-muB}
\end{equation}

\begin{theorem}[Exterior-uniform centered-log comparison]
\label{thm:v6v12-local-comparison}
Under Proposition~\ref{prop:v6v12-bounded-log} and
Lemma~\ref{lem:v6v12-Higgs-mgf},
\begin{equation}
 \log\mathbb E_{\nu_B}
 e^{q(L_B-\mathbb E_{\nu_B}L_B)}
 \le qC_L|B|.
 \label{eq:v6v12-full-centered-mgf}
\end{equation}
Consequently, with \(p=q/(q-1)\),
\begin{equation}
 \mu_B(E)\le e^{C_L|B|}\nu_B(E)^{1/p}
 \label{eq:v6v12-local-event-comparison}
\end{equation}
for every interior event \(E\).
\end{theorem}

\begin{proof}
After centering, the bounded log contributes at most
\(e^{qC_{\rm bd}|B|}\).  Multiply by the Higgs centered exponential and
apply Lemma~\ref{lem:v6v12-Higgs-mgf}.  The denominator-free
centered-log theorem from the frozen V98 volume then gives
\eqref{eq:v6v12-local-event-comparison}.
\end{proof}

\subsection{A conditional reference action tail}

Let \(S_B\) contain the plaquette action terms meeting an interior link and
set
\begin{equation}
 S_{B,\min}(\phi_{B^c})=\inf_{\phi_B}S_B(\phi_B,\phi_{B^c}).
 \label{eq:v6v12-conditional-minimum}
\end{equation}

\begin{lemma}[Compact reference excess-action tail]
\label{lem:v6v12-reference-tail}
Assume uniformly in admissible exterior data that a minimizer exists and,
for every fixed \(u>0\), a product neighbourhood \(N_u\) of a minimizer
has Haar volume at least \(e^{-c_u|B|}\) and satisfies
\begin{equation}
 S_B\le S_{B,\min}+\frac u2|B|\quad\hbox{on }N_u.
 \label{eq:v6v12-minimizer-neighbourhood}
\end{equation}
Then the conditional Wilson reference obeys
\begin{equation}
 \nu_B^{\rm g}\{S_B-S_{B,\min}\ge u|B|\}
 \le\exp\left[-\left(\frac{\beta u}{2}-c_u\right)|B|\right].
 \label{eq:v6v12-reference-excess-tail}
\end{equation}
\end{lemma}

\begin{proof}
The numerator is at most
\(e^{-\beta(S_{B,\min}+u|B|)}\) times normalized Haar volume.
The partition function is at least its integral over \(N_u\), namely
\(e^{-c_u|B|}e^{-\beta(S_{B,\min}+u|B|/2)}\).  Divide.
\end{proof}

Use of excess action makes the estimate insensitive to the numerical value
of the boundary-frustrated minimum.  To control curvature, badness must
still force a positive excess above that minimum.

\subsection{Bad-block domination}

Assume the curvature event satisfies
\begin{equation}
 E_B^{\rm bad}\subset
 \{S_B-S_{B,\min}\ge u|B|\}.
 \label{eq:v6v12-bad-implies-excess}
\end{equation}

\begin{theorem}[Uniform conditional bad-block probability]
\label{thm:v6v12-bad-probability}
Under Theorem~\ref{thm:v6v12-local-comparison},
Lemma~\ref{lem:v6v12-reference-tail}, and
\eqref{eq:v6v12-bad-implies-excess},
\begin{equation}
 \mu_B(E_B^{\rm bad}\mid\mathcal F_{B^c})
 \le e^{-I_{\rm bad}|B|},
 \qquad
 I_{\rm bad}:=\frac1p\left(\frac{\beta u}{2}-c_u\right)-C_L.
 \label{eq:v6v12-interacting-bad-tail}
\end{equation}
If \(I_{\rm bad}>0\), the bound is uniform in total volume and exterior
configuration.
\end{theorem}

\begin{proof}
Apply \eqref{eq:v6v12-local-event-comparison} to the reference probability
in \eqref{eq:v6v12-reference-excess-tail} and collect exponents.
\end{proof}

\begin{corollary}[Product domination and polymer convergence]
\label{cor:v6v12-product-domination}
If \(I_{\rm bad}>0\) and the conditional estimate remains valid while
other blocks are exposed, then every finite family \(\mathcal S\) obeys
\begin{equation}
 \mathbb E_\mu\prod_{B\in\mathcal S}{\bf1}_{E_B^{\rm bad}}
 \le p_0^{|\mathcal S|},
 \qquad p_0=e^{-I_{\rm bad}|B|}.
 \label{eq:v6v12-bad-product}
\end{equation}
The rooted bad-polymer sum converges whenever
\begin{equation}
 C_{\rm an}w e^{-I_{\rm bad}|B|}<1.
 \label{eq:v6v12-polymer-condition}
\end{equation}
\end{corollary}

\begin{proof}
Expose blocks one at a time and use the tower property; each conditional
indicator costs at most \(p_0\).  The rooted animal count then gives the
geometric condition \eqref{eq:v6v12-polymer-condition}.
\end{proof}

\subsection{Scope of the criterion}

\begin{firewall}[Good minimizers and flux sectors]
\label{fw:v6v12-good-minimizer}
Excess action controls curvature only if every admissible exterior/flux
sector has a good conditional minimizer, or if unavoidable topological
action is subtracted and the remaining excess controls local roughness.
A frustrated boundary can make an absolute small-curvature event
impossible.
\end{firewall}

\begin{firewall}[The rate balance is quantitative]
\label{fw:v6v12-rate-balance}
One needs \(I_{\rm bad}>0\): the reference action rate must beat the
centered-log cost \(C_L\).  Extensivity alone is insufficient.  Light
determinant singularities or a weakening Higgs quartic can make \(C_L\)
diverge.
\end{firewall}

\begin{firewall}[Gauge fixing and ghost boundaries]
\label{fw:v6v12-gauge-fixing}
Uniform conditioning may cross a Gribov horizon or gauge-chart boundary.
A Faddeev--Popov factor enters the positive comparison only where its sign
and centered log moment are controlled; otherwise the boundary is a
separate chart-transition polymer.
\end{firewall}

\begin{firewall}[The conformal factor is excluded]
\label{fw:v6v12-conformal-tail}
The negative Euclidean conformal kinetic direction cannot improve the
reference action tail and has no positive centered-log comparison.  This
theorem concerns only the positive massive gauge--Higgs--PV--modulus
sector.
\end{firewall}

\begin{theorem}[Sixth-series V12 result]
\label{thm:v6v12-status}
Exponential sensitivity of rooted positive owners makes the conditional
centered-log cost proportional to the tested block volume and independent
of exterior volume.  A stable Higgs quartic supplies the unbounded-field
centered moment.  If curvature badness forces a uniform reference
excess-action \(u|B|\), the interacting conditional bad rate is
\(I_{\rm bad}=p^{-1}(\beta u/2-c_u)-C_L\).  A positive rate implies
all-family product domination and the V11 bad-polymer sum.  Good-minimizer
flux control, a positive rate balance, gauge-chart boundaries, and the
conformal sector remain open.
\end{theorem}

\begin{proof}
Use Lemma~\ref{lem:v6v12-local-oscillation},
Proposition~\ref{prop:v6v12-bounded-log},
Lemma~\ref{lem:v6v12-Higgs-mgf},
Theorem~\ref{thm:v6v12-local-comparison},
Lemma~\ref{lem:v6v12-reference-tail},
Theorem~\ref{thm:v6v12-bad-probability}, and
Corollary~\ref{cor:v6v12-product-domination}, subject to the four
firewalls.
\end{proof}

V13 will attack the rate-balance inputs.  It will isolate the local
high-modulus and Higgs contributions to \(C_L\), compare them with the
weak-coupling plaquette rate, and determine whether a block size and
curvature threshold can make \(I_{\rm bad}>0\) uniformly along a proposed
cutoff trajectory.

\clearpage
\section{V13: the bad-block rate threshold and gauge trajectories}
\label{sec:v6v13}

V12 reduces conditional bad-block domination to positivity of one explicit
rate.  This note expresses the excess \(u\) in terms of a curvature
threshold, proves the critical inverse-coupling condition, and separates
what block size can accomplish from what only the running coupling can
accomplish.

\subsection{Curvature badness forces an action excess}

Let \(s(U_p)\) be the one-plaquette action and suppose that on the relevant
flux chart
\begin{equation}
 s(U_p)-s(U_p^*)\ge c_g
 \operatorname{dist}(U_p,U_p^*)^2
 \label{eq:v6v13-plaquette-coercivity}
\end{equation}
relative to a good conditional minimizer \(U^*\).  Declare a block bad if
at least a fraction \(\rho\in(0,1]\) of its plaquettes satisfy
\begin{equation}
 \operatorname{dist}(U_p,U_p^*)\ge\delta.
 \label{eq:v6v13-bad-threshold}
\end{equation}
Let the number of tested plaquettes be at least \(c_P|B|\).

\begin{lemma}[Curvature threshold to excess action]
\label{lem:v6v13-curvature-excess}
Under the preceding assumptions,
\begin{equation}
 E_B^{\rm bad}\subset
 \{S_B-S_{B,\min}\ge u_{\rho,\delta}|B|\},
 \qquad
 u_{\rho,\delta}:=c_gc_P\rho\delta^2.
 \label{eq:v6v13-u-rho-delta}
\end{equation}
\end{lemma}

\begin{proof}
Every bad plaquette contributes at least \(c_g\delta^2\) above the
minimizer by \eqref{eq:v6v13-plaquette-coercivity}.  There are at least
\(\rho c_P|B|\) such plaquettes.  Sum their nonnegative excesses.
\end{proof}

This proof requires a plaquettewise good minimizer and nonnegative
decomposition of the excess.  For a frustrated non-Abelian boundary, the
hypothesis is a substantive flux-chart condition.

\subsection{The critical inverse coupling}

Insert \(u_{\rho,\delta}\) into the V12 rate:
\begin{equation}
 I_{\rm bad}(\beta;\rho,\delta)
 =\frac1p\left(
 \frac{\beta c_gc_P\rho\delta^2}{2}
 -c_{u_{\rho,\delta}}\right)-C_L.
 \label{eq:v6v13-I-explicit}
\end{equation}

\begin{theorem}[Critical inverse-coupling criterion]
\label{thm:v6v13-beta-critical}
For fixed \(p,\rho,\delta\) and local comparison cost \(C_L\),
\(I_{\rm bad}>0\) exactly when
\begin{equation}
 \beta>\beta_{\rm crit}(\rho,\delta)
 :=\frac{2\,[c_{u_{\rho,\delta}}+pC_L]}
 {c_gc_P\rho\delta^2}.
 \label{eq:v6v13-beta-critical}
\end{equation}
If this inequality holds with margin \(I_{\rm bad}\ge I_*>0\), then the
bad-polymer condition is satisfied for every block volume obeying
\begin{equation}
 |B|>\frac{\log(C_{\rm an}w)}{I_*}
 \label{eq:v6v13-block-size}
\end{equation}
when \(C_{\rm an}w>1\); if \(C_{\rm an}w\le1\), every positive block volume
suffices.
\end{theorem}

\begin{proof}
Solve the affine inequality
\eqref{eq:v6v13-I-explicit}>0 for \(\beta\).  Under a positive margin, the
V12 polymer condition is
\(C_{\rm an}w e^{-I_*|B|}<1\).  Taking logarithms gives
\eqref{eq:v6v13-block-size}.
\end{proof}

\begin{corollary}[Block size cannot change the sign]
\label{cor:v6v13-block-sign}
If \(I_{\rm bad}\le0\), increasing \(|B|\) cannot prove bad-polymer
convergence by the V12 estimate.  If \(I_{\rm bad}>0\), increasing
\(|B|\) suppresses the polymer activity exponentially.
\end{corollary}

\begin{proof}
The estimated bad probability is \(e^{-I_{\rm bad}|B|}\).  Its exponent
decays with volume only for positive \(I_{\rm bad}\).
\end{proof}

\subsection{Local comparison cost from rooted owners}

For bounded rooted owner families \(r\), assume
\begin{equation}
 \sup_{e}\sum_x
 |\partial_{\phi_e}\ell_x^{(r)}|
 \le c_r.
 \label{eq:v6v13-owner-sensitivity-sum}
\end{equation}
Let \(C_H\) be the centered Higgs constant and let \(C_{\rm chart}\) contain
the finite gauge-chart and local partition costs.

\begin{proposition}[Local cost ledger]
\label{prop:v6v13-CL-ledger}
The V12 comparison constant may be chosen as
\begin{equation}
 C_L\le C_{\rm chart}+C_H+
 L_0\sum_r c_r.
 \label{eq:v6v13-CL-ledger}
\end{equation}
For the high determinant and PV factors, \(c_r=O(\epsilon_r)\) for a
matched adjacent increment by V10--V11.  For a finite path of bounded
rooted owners, the corresponding contributions add in total variation and
remain finite without a global-rank factor.
\end{proposition}

\begin{proof}
Integrating \eqref{eq:v6v13-owner-sensitivity-sum} along an interior path
gives oscillation at most \(L_0c_r|B|\).  Sum owner types and add the Higgs
centered moment and chart constants.  The rooted trace-log activity theorem
gives \(c_r=O(\epsilon_r)\) at an adjacent high-modulus or PV step.
Oscillations of real scalar logs add along a finite path, so their local
constants add as well.
\end{proof}

The determinant rank never appears because
\eqref{eq:v6v13-owner-sensitivity-sum} sums spatially decaying rooted
derivatives before the block volume is counted.

\subsection{Consequences for a cutoff trajectory}

Write \(\beta_{G,n}=1/g_{G,n}^2\) in a fixed lattice normalization for a
gauge factor \(G\).  Let
\begin{equation}
 \overline C_L:=\sup_n C_{L,n}<\infty
 \label{eq:v6v13-uniform-CL}
\end{equation}
on the proposed positive massive trajectory, and fix one admissible
curvature threshold.

\begin{theorem}[Trajectory alternatives for the bad-block mechanism]
\label{thm:v6v13-trajectory-alternatives}
Under the uniform local assumptions above:
\begin{enumerate}
\item If \(\beta_{G,n}\to+\infty\), then the rate criterion
      \eqref{eq:v6v13-beta-critical} holds with a positive uniform margin
      for every sufficiently large \(n\).  The finitely many earlier steps
      may be handled separately.
\item If \(\inf_n\beta_{G,n}>\sup_n\beta_{{\rm crit},n}\), the rate is
      positive uniformly at every step.
\item If \(\beta_{G,n}\to0\), while
      \(c_{u_{\rho,\delta}}+pC_{L,n}\) has a positive lower bound, then
      \(I_{{\rm bad},n}<0\) eventually.  Fixed-curvature bad blocks are not
      rare by this comparison mechanism.
\end{enumerate}
\end{theorem}

\begin{proof}
All three statements follow directly from the affine formula
\eqref{eq:v6v13-I-explicit} and boundedness of its other constants.
\end{proof}

For asymptotically free \(SU(3)\) and \(SU(2)\) candidate trajectories,
item 1 is the desired direction, conditional on a nonperturbative
trajectory and the uniform owner estimates.  For the Gaussian-boundary
hypercharge trajectory recorded in the archives, the perturbative flow
drives \(g_Y\) upward and \(\beta_Y\) downward; item 3 expresses the same
obstruction at the level of good-field probability.  This is not an
independent proof of a Landau pole, but it shows that the present
small-curvature polymer mechanism cannot survive a vanishing inverse
coupling.

\subsection{Optimizing the curvature threshold}

The entropy constant \(c_u\) generally worsens as the minimizer
neighbourhood shrinks.  Define
\begin{equation}
 \mathcal R(\rho,\delta)
 :=\frac{c_{u_{\rho,\delta}}+pC_L}
 {\rho\delta^2}.
 \label{eq:v6v13-threshold-objective}
\end{equation}

\begin{proposition}[Threshold optimization]
\label{prop:v6v13-threshold-optimization}
Within any compact admissible set of \((\rho,\delta)\) on which
\(c_{u_{\rho,\delta}}\) is continuous, the best rate criterion is attained
at a minimizer of \(\mathcal R\).  No choice of threshold can remove the
need for a positive lower bound on \(\beta\) if
\(\inf(c_u+pC_L)>0\).
\end{proposition}

\begin{proof}
Continuity and compactness give a minimizer.  Equation
\eqref{eq:v6v13-beta-critical} differs from \(\mathcal R\) only by the
fixed factor \(2/(c_gc_P)\).  Since \(\rho\delta^2\) is bounded above on
the compact group chart, a positive numerator gives a positive critical
inverse coupling.
\end{proof}

\subsection{Metric and multi-factor effects}

On a bounded-geometry metric chart, plaquette counts, distances, and
coercivity constants are uniformly comparable.  Let
\(c_g(g),c_P(g),C_L(g)\) denote their metric-dependent values.

\begin{corollary}[Bounded-geometry stability]
\label{cor:v6v13-metric-stability}
If
\begin{equation}
 \inf_g c_g(g)c_P(g)>0,
 \qquad
 \sup_g C_L(g)<\infty
 \label{eq:v6v13-metric-uniformity}
\end{equation}
on a positive bounded-geometry chart, then one common
\(\beta_{\rm crit}\) works throughout that chart.
\end{corollary}

\begin{proof}
Take the worst numerator and denominator in
\eqref{eq:v6v13-beta-critical}.
\end{proof}

For a product gauge group, a block is good only if every required factor is
good.  A union bound preserves exponential rarity with the minimum positive
rate and a fixed group-factor multiplicity.  Thus one factor whose rate
becomes nonpositive obstructs the joint good-block card.

\begin{firewall}[Perturbative asymptotic freedom is not the trajectory]
\label{fw:v6v13-perturbative-AF}
The signs of the perturbative \(SU(3)\) and \(SU(2)\) beta functions suggest
\(\beta_{G,n}\to\infty\), but do not construct the nonperturbative
finite-cutoff measures with the uniform local constants assumed above.
Item 1 of Theorem~\ref{thm:v6v13-trajectory-alternatives} is conditional.
\end{firewall}

\begin{firewall}[Hypercharge requires a changed ultraviolet mechanism]
\label{fw:v6v13-hypercharge-mechanism}
If the hypercharge inverse coupling tends to zero, neither a larger block
nor a different fixed curvature threshold makes the present reference
tail positive.  Progress then requires an ultraviolet completion, a
non-Gaussian fixed trajectory, embedding into another gauge sector, or a
different control mechanism.  None is selected or proved here.
\end{firewall}

\begin{firewall}[Metric bounded geometry excludes the conformal runaway]
\label{fw:v6v13-metric-chart}
Corollary~\ref{cor:v6v13-metric-stability} assumes a positive
bounded-geometry chart.  It does not control the unbounded real Euclidean
conformal direction or sum over all metric topologies.
\end{firewall}

\begin{firewall}[A reference tail does not prove reflection positivity]
\label{fw:v6v13-tail-not-RP}
Bad-block rarity is a Euclidean probability estimate.  It does not imply a
positive reflected transfer matrix for the improved chiral gauge and
gravity regulator.
\end{firewall}

\begin{theorem}[Sixth-series V13 result]
\label{thm:v6v13-status}
If a fraction \(\rho\) of plaquettes exceeds curvature threshold
\(\delta\), plaquette coercivity gives the explicit action excess
\(u=c_gc_P\rho\delta^2\).  The conditional bad-block rate is positive
exactly above the critical inverse coupling
\[
 \beta_{\rm crit}
 =2(c_u+pC_L)/(c_gc_P\rho\delta^2).
\]
Block size can enforce polymer convergence after this rate is positive,
but cannot change its sign.  A bounded local rooted-owner cost is therefore
compatible with a candidate asymptotically free trajectory and
incompatible with a trajectory whose inverse coupling tends to zero.
Nonperturbative gauge trajectories, the hypercharge ultraviolet
replacement, unbounded metric directions, and reflection positivity
remain open.
\end{theorem}

\begin{proof}
Use Lemma~\ref{lem:v6v13-curvature-excess},
Theorem~\ref{thm:v6v13-beta-critical},
Proposition~\ref{prop:v6v13-CL-ledger},
Theorem~\ref{thm:v6v13-trajectory-alternatives}, and
Proposition~\ref{prop:v6v13-threshold-optimization}, subject to the four
firewalls.
\end{proof}

V14 will go upstream at the hypercharge obstruction.  It will formulate
precise continuum-construction alternatives as changes of ultraviolet
data, derive their matching requirements at a finite transition scale, and
separate mathematical compatibility conditions from unproved model
selection.


\clearpage
\section{V14: finite-scale ultraviolet matching at the hypercharge obstruction}
\label{sec:v6v14}

\subsection{Purpose and logical position}

V13 showed that the present small-curvature comparison has a sharp sign
condition.  In particular, it cannot control a hypercharge trajectory whose
inverse coupling tends to zero.  The correct upstream move is therefore to
identify what ultraviolet data could replace that trajectory and what must be
proved at the scale where the replacement is joined to the Standard Model.
This note gives an exact finite-cutoff sewing theorem, derives gauge-coupling,
representation, anomaly, and reflection-positivity matching conditions, and
separates three logically different ultraviolet alternatives.  It does not
select or construct an ultraviolet completion.

Fix a transition slice \(\Sigma_*\) at lattice spacing \(a_*\).  All spaces in
this note are finite-dimensional because the statement is first made at fixed
volume and fixed cutoff.  Let \((X_*,\nu_*)\) be the interface configuration
space with its reference measure.  Gauge fixing, when used, is restricted to a
sign chart on which the density is real and nonnegative, as in V11.

\subsection{Exact sewing of two cutoff measures}

Let \(X_-\), \(X_+\), and \(X_*\) be standard Borel spaces.  Write the infrared
and ultraviolet half-space weights in disintegrated form
\begin{equation}
 d\mu_-(x_-,q)=K_-(x_-,q)\,d\lambda_-(x_-)\,d\nu_*(q),
 \qquad
 d\mu_+(q,x_+)=K_+(q,x_+)\,d\nu_*(q)\,d\lambda_+(x_+),
 \label{eq:v6v14-half-densities}
\end{equation}
with nonnegative measurable kernels.  Define their interface amplitudes
\begin{equation}
 z_-(q):=\int K_-(x_-,q)\,d\lambda_-(x_-),
 \qquad
 z_+(q):=\int K_+(q,x_+)\,d\lambda_+(x_+).
 \label{eq:v6v14-interface-amplitudes}
\end{equation}

\begin{theorem}[Finite-cutoff sewing]
\label{thm:v6v14-sewing}
Assume \(0<Z_*:=\int z_-(q)z_+(q)\,d\nu_*(q)<\infty\).  Then
\begin{equation}
 d\mu_*(x_-,q,x_+)
 :=Z_*^{-1}K_-(x_-,q)K_+(q,x_+)
 \,d\lambda_-\,d\nu_*\,d\lambda_+
 \label{eq:v6v14-sewn-measure}
\end{equation}
is a probability measure.  Conditional on \(q\), the two open halves are
independent.  Moreover, \(\mu_*\) is the unique probability measure with this
conditional-independence property, the two conditional kernels obtained by
normalizing \(K_-\) and \(K_+\), and interface marginal
\begin{equation}
 d\pi_*(q)=Z_*^{-1}z_-(q)z_+(q)\,d\nu_*(q).
 \label{eq:v6v14-interface-marginal}
\end{equation}
\end{theorem}

\begin{proof}
Tonelli's theorem gives total mass
\(Z_*^{-1}\int z_-z_+\,d\nu_*=1\).  Wherever \(z_-z_+>0\), division by the
two amplitudes gives the conditional laws
\[
 K_-(x_-,q)d\lambda_-(x_-)/z_-(q),\qquad
 K_+(q,x_+)d\lambda_+(x_+)/z_+(q).
\]
Their product and \eqref{eq:v6v14-interface-marginal} recover
\eqref{eq:v6v14-sewn-measure}.  On the null complement the conditional laws
are irrelevant.  Conversely, disintegration and conditional independence
force precisely this product, proving uniqueness up to null sets.
\end{proof}

The theorem says that matching does not mean equality of two bulk actions.
It means that both sides use the same interface variables and reference
measure, and that their boundary amplitudes have an integrable product.
Local counterterms at \(\Sigma_*\) may be split between the two kernels; their
sum is what enters the sewn measure.

\begin{proposition}[Quantitative stability under interface mismatch]
\label{prop:v6v14-sewing-stability}
Let \(w=z_-z_+\) and \(\widetilde w=\widetilde z_-\widetilde z_+\) be two
nonnegative interface weights with positive finite masses \(Z\) and
\(\widetilde Z\).  Their normalized interface laws satisfy
\begin{equation}
 \|\pi-\widetilde\pi\|_{\rm TV}
 \le \frac{\|w-\widetilde w\|_{L^1(\nu_*)}}
 {\min\{Z,\widetilde Z\}}.
 \label{eq:v6v14-TV-stability}
\end{equation}
Consequently, the error in every interface observable \(F\) with
\(\|F\|_\infty\le1\) is at most twice the right-hand side under the convention
\(\|\pi-\widetilde\pi\|_{\rm TV}=\frac12\int|d\pi-d\widetilde\pi|\).
\end{proposition}

\begin{proof}
Writing \(d\pi=Z^{-1}w\,d\nu_*\) and adding and subtracting
\(Z^{-1}\widetilde w\),
\[
 \int\left|\frac wZ-\frac{\widetilde w}{\widetilde Z}\right|d\nu_*
 \le \frac{\|w-\widetilde w\|_1}{Z}
 +\frac{|Z-\widetilde Z|}{Z}.
\]
Since \(|Z-\widetilde Z|\le\|w-\widetilde w\|_1\), this gives the claim
with denominator \(Z\); interchange the two laws and take the better bound.
The observable estimate is the defining dual bound for total variation.
\end{proof}

Thus approximate matching is meaningful only with a norm and an error that
vanishes along the continuum sequence.  Agreement of finitely many fitted
couplings alone does not imply \eqref{eq:v6v14-TV-stability}.

\subsection{Kinetic matching for an embedded hypercharge generator}

Let \(\mathfrak g_1,\ldots,\mathfrak g_s\) be compact Lie algebras with
positive invariant inner products \(\langle\cdot,\cdot\rangle_i\), and let
\(Y_i\in\mathfrak g_i\) be commuting elements.  Embed the infrared
hypercharge connection by
\begin{equation}
 A^{(i)}=a\,Y_i,
 \qquad
 F^{(i)}=f\,Y_i,
 \label{eq:v6v14-Y-embedding}
\end{equation}
where \(a\) is an Abelian one-form and \(f=da\).

\begin{proposition}[Tree-level kinetic matching]
\label{prop:v6v14-kinetic-matching}
If the ultraviolet gauge action restricted to these directions is
\begin{equation}
 S_{\rm UV}^{(2)}=\frac14\sum_{i=1}^s g_i^{-2}
 \int\langle F^{(i)},F^{(i)}\rangle_i,
 \label{eq:v6v14-UV-kinetic}
\end{equation}
then the infrared normalization \(S_Y^{(2)}=(4g_Y^2)^{-1}\int f^2\)
requires
\begin{equation}
 g_Y^{-2}=\sum_{i=1}^s
 \langle Y_i,Y_i\rangle_i\,g_i^{-2}.
 \label{eq:v6v14-coupling-match}
\end{equation}
The formula is invariant under the simultaneous rescaling
\(Y_i\mapsto cY_i\), \(a\mapsto a/c\), once the physical charges are kept
fixed.
\end{proposition}

\begin{proof}
Substitute \eqref{eq:v6v14-Y-embedding} into
\eqref{eq:v6v14-UV-kinetic}.  Invariance and bilinearity give
\(\langle F^{(i)},F^{(i)}\rangle_i=f^2\langle Y_i,Y_i\rangle_i\).
Comparison with the normalized Abelian quadratic form yields
\eqref{eq:v6v14-coupling-match}.  Under the stated rescaling, the components
of the connection and the charge eigenvalues transform inversely, leaving
the covariant derivative and the quadratic action unchanged.
\end{proof}

Loop-level matching replaces \eqref{eq:v6v14-coupling-match} by the same
identity plus finite threshold counterterms obtained after integrating out
fields with masses near the transition scale.  Those counterterms belong to
the common interface action in Theorem~\ref{thm:v6v14-sewing}; they cannot be
discarded when comparing the two boundary amplitudes.

\subsection{Representation and anomaly matching}

Let \(R\) be a finite-dimensional unitary ultraviolet representation.  Its
restriction to the one-parameter subgroup generated by \(Y\) decomposes as
\begin{equation}
 R|_{U(1)_Y}=\bigoplus_j R_j,
 \qquad Y|_{R_j}=y_j\mathbf 1.
 \label{eq:v6v14-charge-restriction}
\end{equation}
The multiset of light eigenvalues \(y_j\), including multiplicities and
their \(SU(3)\times SU(2)\) representations, must equal the Standard Model
charge data below the threshold after vectorlike massive pairs are removed.

\begin{lemma}[Necessary infinitesimal anomaly conditions]
\label{lem:v6v14-anomaly-necessary}
For left-handed Weyl fields with hypercharges \(y_j\) and non-Abelian
representations \(r_{3,j}\), \(r_{2,j}\), gauge invariance of the chiral
determinant requires
\begin{align}
 &\sum_j d(r_{3,j})d(r_{2,j})y_j^3=0,
 &&\sum_j d(r_{3,j})d(r_{2,j})y_j=0,
 \label{eq:v6v14-abelian-anomalies}\\
 &\sum_j d(r_{2,j})T_3(r_{3,j})y_j=0,
 &&\sum_j d(r_{3,j})T_2(r_{2,j})y_j=0,
 \label{eq:v6v14-mixed-anomalies}
\end{align}
as well as vanishing of the pure non-Abelian local anomalies and of every
applicable global anomaly.  Here \(T_k\) is the quadratic index in the chosen
normalization.
\end{lemma}

\begin{proof}
The infinitesimal phase variation of a four-dimensional chiral determinant
is a local degree-six invariant polynomial.  On the product algebra, its
coefficients are traces over the Weyl representation.  The \(U(1)^3\) and
mixed gravitational--\(U(1)\) coefficients are respectively
\(\operatorname{tr}Y^3\) and \(\operatorname{tr}Y\), which give
\eqref{eq:v6v14-abelian-anomalies}.  The
\(SU(k)^2U(1)\) coefficient is
\(\operatorname{tr}_{r_k}(T^aT^b)\,y_j
=T_k(r_{k,j})y_j\delta^{ab}\), with the spectator dimension multiplying
it; summing gives \eqref{eq:v6v14-mixed-anomalies}.  A gauge-invariant
determinant has zero infinitesimal variation, so every independent
coefficient vanishes.  The local polynomial does not detect global
anomalies, hence their separate condition.
\end{proof}

\begin{firewall}[Vanishing local coefficients is not determinant construction]
\label{fw:v6v14-anomaly-firewall}
Lemma~\ref{lem:v6v14-anomaly-necessary} supplies necessary local
conditions.  It neither constructs a globally smooth gauge-invariant
determinant section nor proves cancellation of torsion or global anomalies.
The phase problem isolated before V10 therefore remains a separate local
input to the continuum program.
\end{firewall}

\subsection{Relevant-coordinate matching}

At the transition scale, let \(u\in\mathbb R^m\) denote ultraviolet relevant
parameters and let \(c\in\mathbb R^r\) denote the infrared couplings and
counterterms to be matched.  Let \(M(u)\) be the finite-cutoff matching map
defined by integrating one side to the common interface.

\begin{proposition}[Local solvability of the matching equations]
\label{prop:v6v14-IFT-matching}
Suppose \(M\) is \(C^1\) near \(u_0\), \(M(u_0)=c_0\), and
\(D M(u_0):\mathbb R^m\to\mathbb R^r\) is surjective.  Then every \(c\) in a
neighbourhood of \(c_0\) is attained by parameters \(u\) near \(u_0\).
After choosing a complement to \(\ker D M(u_0)\), the solutions form locally
a \(C^1\) manifold of dimension \(m-r\).
\end{proposition}

\begin{proof}
Choose an \(r\)-dimensional subspace on which \(D M(u_0)\) is an
isomorphism onto \(\mathbb R^r\).  The implicit-function theorem solves the
\(r\) matching equations for those coordinates as \(C^1\) functions of the
remaining \(m-r\) coordinates and of \(c\).
\end{proof}

This proposition identifies the required transversality calculation.  A
parameter count without surjectivity does not establish matching, and a
formal beta function does not establish differentiability of the exact map.

\subsection{Reflection-positive sewing}

Let \(\mathcal H_*\) be the interface Hilbert space and suppose the two
regions are represented by bounded operators \(T_-\) and \(T_+\) on
\(\mathcal H_*\).

\begin{lemma}[Positive transfer through the interface]
\label{lem:v6v14-positive-transfer}
If \(T_-=A^*A\) and \(T_+=B^*B\) are positive self-adjoint operators, then
for every \(f\in\mathcal H_*\),
\begin{equation}
 \langle f,T_-^{1/2}T_+T_-^{1/2}f\rangle
 =\|BT_-^{1/2}f\|^2\ge0.
 \label{eq:v6v14-positive-composition}
\end{equation}
Hence sewing preserves positivity when the common interface operation is
the symmetric product \(T_-^{1/2}T_+T_-^{1/2}\).
\end{lemma}

\begin{proof}
Functional calculus gives the positive square root of \(T_-\).  Substitute
\(T_+=B^*B\) and use the definition of the adjoint.
\end{proof}

The unsymmetrized product \(T_-T_+\) need not be self-adjoint.  Therefore
positivity of each half cannot be transferred by a purely formal
multiplication of Euclidean densities; the common reflection slice and its
Hilbert-space identification are part of the matching data.

\subsection{The three ultraviolet alternatives}

The rate obstruction of V13 leaves the following exhaustive alternatives
for the present strategy.

\begin{enumerate}
\item \emph{Embedded ultraviolet gauge sector.}  Above \(a_*\), hypercharge
is the restriction of one or more compact gauge factors whose inverse
couplings stay above the critical rate.  One must prove
\eqref{eq:v6v14-coupling-match}, the charge restriction
\eqref{eq:v6v14-charge-restriction}, anomaly-line compatibility, threshold
matching, a positive transfer construction, and the uniform local estimates
of V8--V13.
\item \emph{Non-Gaussian Abelian fixed trajectory.}  Hypercharge remains an
Abelian factor, but its exact ultraviolet trajectory stays in a region
\(\inf_n\beta_{Y,n}>\sup_n\beta_{{\rm crit},n}\).  One must construct the
fixed measure and prove its relevant manifold, reflection positivity, and
uniform owner bounds.  Naming an asymptotically safe fixed point is not such
a construction.
\item \emph{Finite effective cutoff.}  Stop at a smallest spacing
\(a_{\min}>0\).  This may define a controlled effective theory on a finite
range of scales, but it is not a continuum limit because no sequence with
\(a\downarrow0\) is supplied.
\end{enumerate}

\begin{proposition}[A finite cutoff cannot satisfy the continuum goal]
\label{prop:v6v14-finite-cutoff}
Suppose a family of Schwinger functionals is defined only for lattice
spacings \(a\ge a_{\min}>0\).  Then it cannot, by restriction alone, define
a continuum Schwinger functional as a limit along \(a_n\downarrow0\).
\end{proposition}

\begin{proof}
No sequence in the stated domain has \(a_n\downarrow0\), since every term is
bounded below by \(a_{\min}\).  A continuum extension would be additional
data, not a consequence of the finite-cutoff family.
\end{proof}

\subsection{Gravity at the same transition slice}

For a metric interface, the common variable includes the induced metric
\(h\).  If the two actions contain boundary terms giving a well-posed
Dirichlet variation, their first variations have the form
\begin{equation}
 \delta S_\pm=(\hbox{bulk equations})
 +\int_{\Sigma_*}\langle\pi_\pm,\delta h\rangle.
 \label{eq:v6v14-gravity-boundary-variation}
\end{equation}

\begin{lemma}[Stationary gravitational sewing]
\label{lem:v6v14-gravity-sewing}
With the outward-normal convention on each half, a sewn stationary point
requires a common induced metric and
\begin{equation}
 \pi_-+\pi_+=0
 \label{eq:v6v14-momentum-match}
\end{equation}
in the absence of an interface stress tensor.  An interface action
\(S_*(h)\) changes the right side to
\(-\delta S_*/\delta h\).
\end{lemma}

\begin{proof}
The induced metric is the variable over which the two half amplitudes are
sewn.  Add their first variations and vary this common \(h\) arbitrarily.
The boundary coefficient must vanish, giving
\eqref{eq:v6v14-momentum-match}; adding \(S_*\) adds its derivative.
\end{proof}

\begin{firewall}[Stationary matching does not cure the conformal problem]
\label{fw:v6v14-gravity-firewall}
Lemma~\ref{lem:v6v14-gravity-sewing} is a boundary compatibility condition.
It does not make the Euclidean Einstein action bounded below, construct a
positive metric measure, or control changes of topology.
\end{firewall}

\begin{theorem}[Conditional ultraviolet-interface criterion]
\label{thm:v6v14-interface-criterion}
An ultraviolet replacement at a finite transition scale is compatible with
the V8--V13 positive massive program if the following data are constructed
uniformly in volume and along the cutoff sequence:
\begin{enumerate}
\item common interface variables and integrable half amplitudes satisfying
Theorem~\ref{thm:v6v14-sewing};
\item kinetic and threshold matching, including
\eqref{eq:v6v14-coupling-match} in an embedded realization;
\item the correct restricted light representations and a globally
gauge-invariant determinant line, whose necessary local conditions include
\eqref{eq:v6v14-abelian-anomalies}--\eqref{eq:v6v14-mixed-anomalies};
\item a surjective exact relevant-coordinate matching map as in
Proposition~\ref{prop:v6v14-IFT-matching};
\item a common reflection-positive interface transfer construction as in
Lemma~\ref{lem:v6v14-positive-transfer};
\item a positive bad-block rate and summable local owner activities on both
sides of the transition;
\item the gravitational boundary matching
\eqref{eq:v6v14-momentum-match} together with an independently constructed
positive, tight metric sector.
\end{enumerate}
Under these hypotheses, the two finite-cutoff measures sew exactly at
\(a_*\), and every already-proved local activity and bad-block estimate may
be applied on its own side with the interface constants added to the finite
local cost ledger.  If their all-scale activities are summable, the V11
transport argument crosses the transition with finite total cost.
\end{theorem}

\begin{proof}
Items 1 and 4 give a well-defined matched cutoff measure and locally solve
the finite set of relevant matching equations.  Items 2 and 3 identify the
same gauge normalization and light chiral content on the common interface.
Item 5 preserves positivity across that interface.  Item 6 supplies the
uniform probabilistic domain on which the local cluster estimates apply.
Item 7 supplies the corresponding metric boundary data and explicitly
assumes the missing positive metric measure.  The transition occupies only
a fixed finite number of scale steps, so its finite local constants add to
the V13 ledger.  Summability away from it then gives the finite total
transport bound of V11.
\end{proof}

\begin{firewall}[Compatibility is not existence]
\label{fw:v6v14-compatibility-existence}
Theorem~\ref{thm:v6v14-interface-criterion} is an exact implication from
listed hypotheses.  No embedded model or non-Gaussian Abelian fixed measure
has been constructed here, and no claim is made that all seven hypotheses
hold for a single theory.
\end{firewall}

\begin{theorem}[Sixth-series V14 result]
\label{thm:v6v14-status}
The hypercharge obstruction can be crossed at a finite transition scale
without changing the infrared Standard Model only if the ultraviolet and
infrared measures share compatible interface variables, kinetic
normalization, light representations, determinant line, relevant
coordinates, and reflection-positive transfer data.  These conditions are
sufficient for exact finite-cutoff sewing when supplemented by positive
integrable half amplitudes, a positive bad-block rate, summable local
activities, and a positive tight metric sector.  An embedded compact sector
or a constructed non-Gaussian Abelian trajectory could meet the rate
requirement; a theory stopped at finite cutoff cannot meet the continuum
goal.
\end{theorem}

\begin{proof}
Combine Theorem~\ref{thm:v6v14-sewing},
Propositions~\ref{prop:v6v14-kinetic-matching} and
\ref{prop:v6v14-IFT-matching}, Lemma~\ref{lem:v6v14-anomaly-necessary},
Lemma~\ref{lem:v6v14-positive-transfer},
Proposition~\ref{prop:v6v14-finite-cutoff}, and
Theorem~\ref{thm:v6v14-interface-criterion}, subject to the two firewalls.
\end{proof}

V15 will perform the required companion-program audit before choosing the
next bottleneck.  It will compare the newest Yang--Mills and Navier--Stokes
notes with the interface criterion above, import only genuinely new local
mechanisms, and determine whether the next attack should target the complex
chiral determinant, the nonperturbative gauge trajectory, or the positive
metric sector.


\clearpage
\section{V15: compulsory companion audit and a weighted-tree repair}
\label{sec:v6v15}

\subsection{Read-only audit record}

The newest active companion files inspected at this checkpoint were
\begin{center}
\begin{tabular}{p{0.62\textwidth}r}
\hline
file & bytes\\
\hline
\texttt{Renewal\_Yang\_Mills\_Mass\_Gap\_Research\_Notes\_V12.tex}
& \(176\,208\)\\
\texttt{Renewal\_NS\_Stability\_Research\_Notes\_V31.tex}
& \(887\,537\)\\
\hline
\end{tabular}
\end{center}
Their SHA--256 digests, in the same order, were
\begin{align}
 &\texttt{e49d5fd3495e82f6cc6813e1d1e494bfab5eab0a385ed1d76c5683ce6deca422},
 \label{eq:v6v15-YM-hash}\\
 &\texttt{f1121b62238ad5491bb7cf03635ea9934942edf573ed8faaa9ee79e1d38a6696}.
 \label{eq:v6v15-NS-hash}
\end{align}
Neither file was modified.

The Yang--Mills V12 note makes two relevant advances.  First, it corrects
the grade of its V7--V11 marked-tree atlas: an analytic \(m!\) at a fixed
tree cannot be followed by another factorial from rooted Cayley
enumeration.  Second, it restores the exact \(SU(2)\) Haar series before
taking derivative norms and proves an integrated bound exponential in the
derivative order.  Its exact action factor and strong Wilson-density card
remain open.

The Navier--Stokes V31 file is unchanged from the V10 audit.  Its useful
methodological boundary remains that a restricted recent-entry formation
quantity must not be replaced prematurely by the global continuation
stock, and already-paid persistent intervals must not be charged again.
For the present program this means retaining conditional bad-block events
and the exact transition interface, rather than replacing them by an
unconditional global oscillation norm.

\begin{firewall}[No theorem is transferred merely by analogy]
\label{fw:v6v15-no-analogy-transfer}
The companion notes concern different measures and target statements.  The
audit below re-proves the only imported combinatorial identity in the SM
notation.  The Yang--Mills exact-Haar theorem can be used for an SM
\(SU(2)\) link only after the chart and normalization are identified; it
does not supply the missing full action card, hypercharge trajectory,
chiral phase, or metric measure.
\end{firewall}

\subsection{Source corrections to V8}

Three V8 displays lost backslashes during source insertion.  Their intended
and henceforth authoritative readings are
\begin{align}
 \left\|h_\epsilon-
 \left(1+\sum_{j=1}^{N-1}\epsilon^jV_j\right)\right\|_{L^p(\gamma_q)}
 &\le (K_pN\epsilon)^N+K_pe^{-c_p/\epsilon},
 \label{eq:v6v15-correct-V8-remainder}\\
 R_{\epsilon,A}
 &:=h_{\epsilon,A}-1-
 \sum_{j=1}^{N_\epsilon-1}\epsilon^jV_{j,A},
 \label{eq:v6v15-correct-V8-R}\\
 \eta_n^{\rm compact}
 &\le C\epsilon_n+Ce^{-c/\epsilon_n}.
 \label{eq:v6v15-correct-V8-eta}
\end{align}
These corrections change no intended estimate.

\subsection{The factorial defect in the V2 proof}

The V2 proof first bounded an unweighted labelled-tree sum by a factor
\(r!m!\), and then bounded the product of derivative-degree factorials by
another factorial of order \(r+m\).  Those two estimates cannot be
multiplied and absorbed into a fixed exponential base.  Consequently, the
displayed proof of Theorem~\ref{thm:v6v2-one-step-SIIC} is incomplete as
written.  The theorem can nevertheless be repaired because the tree count
and all vertex-degree factorials admit one joint factorial bound.

For a labelled tree \(T\) on \([N]\), let \(d_i(T)\) be the degree of
vertex \(i\).

\begin{lemma}[Exact degree-weighted Cayley census]
\label{lem:v6v15-weighted-Cayley}
For every \(N\ge2\),
\begin{equation}
 \sum_{T\in\mathfrak T_N}\prod_{i=1}^N d_i(T)!
 =(N-2)!\binom{3N-3}{N-2}.
 \label{eq:v6v15-weighted-Cayley}
\end{equation}
In particular,
\begin{equation}
 \sum_{T\in\mathfrak T_N}\prod_i d_i(T)!
 \le 8^N N!.
 \label{eq:v6v15-weighted-Cayley-bound}
\end{equation}
\end{lemma}

\begin{proof}
A Pruefer word of length \(N-2\) contains label \(i\) exactly
\(n_i=d_i(T)-1\) times.  For fixed nonnegative counts with
\(\sum_i n_i=N-2\), the number of words is
\((N-2)!/\prod_i n_i!\).  Therefore the contribution of these counts is
\[
 (N-2)!\prod_i\frac{(n_i+1)!}{n_i!}
 =(N-2)!\prod_i(n_i+1).
\]
Summing over the counts extracts the coefficient of \(x^{N-2}\) from
\[
 \left(\sum_{n\ge0}(n+1)x^n\right)^N
 =(1-x)^{-2N}.
\]
That coefficient is \(\binom{2N+(N-2)-1}{N-2}\), proving
\eqref{eq:v6v15-weighted-Cayley}.  Finally
\(\binom{3N-3}{N-2}\le2^{3N-3}\) and \((N-2)!\le N!\).
\end{proof}

The exact formula is the point: the largest vertex factorial and the number
of trees are correlated through the same degree sequence.  Bounding them
separately manufactures a second factorial.

\subsection{Spatially weighted joint compiler}

Retain the V2 covariance majorant
\(C(x,y)\le C_0e^{-\mu d(x,y)}\), bounded support radius \(R_0\), and
local activity \(\zeta\).  Let \(v_1,\ldots,v_r\) be observable anchors and
let \(A_1,\ldots,A_m\) be interaction owners.  Repetitions are permitted in
the ordered upper bound; a square-free sum is smaller after division by
\(m!\).

\begin{lemma}[Degree-weighted insertion-tree sum]
\label{lem:v6v15-weighted-insertion}
For \(0<\mu'<\mu\), there is \(K<\infty\), depending only on the V2 local
geometric constants, such that
\begin{align}
 &\sum_{A_1,\ldots,A_m}
 \sum_{T\in\mathfrak T_{r+m}}
 \left(\prod_{u=1}^{r+m}d_u(T)!\right)
 \left(\prod_{ij\in E(T)}C(v_i,v_j)\right)
 \prod_{j=1}^m z_{A_j}
 \notag\\
 &\hspace{4em}\le
 r!m!K^{r+m}\zeta^m
 \mathfrak T_r^{\mu'}(v_1,\ldots,v_r).
 \label{eq:v6v15-weighted-insertion}
\end{align}
For \(r=1\), the last factor is one.
\end{lemma}

\begin{proof}
Fix a labelled tree and orient every component of its interaction-vertex
forest toward the first observable vertex reached.  Sum each interaction
anchor along its unique parent edge.  The covariance row sum, the activity
definition, bounded support displacement, and the exponent margin
\(\mu-\mu'\) give at most \(K^m\zeta^m\), independently of the degrees.
After contracting those oriented branches, retain a spanning tree on the
observable anchors and sum all surplus choices into \(K^{r+m}\).  This is
the same leaf-contraction estimate as V2, performed before counting the
labelled tree.

It remains to sum the degree weights over the labelled tree choices.
Lemma~\ref{lem:v6v15-weighted-Cayley} bounds their total by
\(8^{r+m}(r+m)!\).  The elementary binomial inequality
\[
 (r+m)!\le2^{r+m}r!m!
\]
gives the displayed factorials.  The sum over the surviving observable
trees is precisely dominated by
\(\mathfrak T_r^{\mu'}\).  Enlarging \(K\) absorbs all fixed geometric and
exponential factors.
\end{proof}

\begin{theorem}[Repaired adjacent-cutoff SIIC]
\label{thm:v6v15-repaired-SIIC}
Assume the three hypotheses of
Theorem~\ref{thm:v6v2-one-step-SIIC}.  Expand an exponential interaction
with its coefficient \(1/m!\), or equivalently write a square-free
unordered owner selection as at most \(1/m!\) times the ordered distinct
sum.  Then the complete Gaussian BKAR insertion series satisfies
\begin{equation}
 \sum_{m\ge0}(K\zeta)^m
 \label{eq:v6v15-geometric-insertion}
\end{equation}
times the required \(r!K^r\) observable tree bound.  Thus for
\(K\zeta<1\) the one-step all-order transport theorem remains valid with
\(\eta=K\zeta\).
\end{theorem}

\begin{proof}
At a tree vertex of degree \(d\), the V2 normalized derivative-row
hypothesis gives an unnormalized factor at most \(d!L^dz_A\).  Covariance
edges absorb the fixed powers of \(L\).  Instead of multiplying a
degree-factorial estimate by an unweighted Cayley estimate, apply
Lemma~\ref{lem:v6v15-weighted-insertion} directly.  Its \(m!\) is canceled
by the exponential coefficient or by unordered-to-ordered conversion.
The remaining series is \eqref{eq:v6v15-geometric-insertion}, and the
remaining \(r!\) is exactly the allowed observable arity factorial.
\end{proof}

\begin{assessment}[Append-only correction to V2]
\label{ass:v6v15-V2-correction}
The conclusion of Theorem~\ref{thm:v6v2-one-step-SIIC} is retained, but its
proof is replaced by Theorem~\ref{thm:v6v15-repaired-SIIC}.  Any use of V2
must include the unordered-owner or \(1/m!\) hypothesis explicitly.  No
argument may multiply the old unweighted tree factorial by a separately
estimated derivative-degree factorial.
\end{assessment}

\subsection{Effect on the V7--V8 compact activity}

The polynomial APGC of V7 estimates the \emph{normalized} derivative rows
at total owned order \(J\) by \(K^JJ!\).  The repaired tree compiler pays
all degree factorials through
Lemma~\ref{lem:v6v15-weighted-insertion}; it does not introduce a second
\(J!\).  Hence the V7 optimally truncated series still has the single
Gevrey-one sum
\[
 \sum_{J<N}(K\epsilon)^JJ!.
\]

\begin{proposition}[Corrected factorial grade of the polynomial route]
\label{prop:v6v15-polynomial-grade}
Assume the V7 APGC, bounded owner occupancy, the V2 spatial hypotheses, and
unordered square-free selection.  Then the polynomial adjacent activity is
still \(O(\epsilon)\) at truncation order \(N\asymp\epsilon^{-1}\), with
one allowed observable arity factorial.  The proof must use the joint
compiler \eqref{eq:v6v15-weighted-insertion}.
\end{proposition}

\begin{proof}
For fixed total order \(J\), apply the APGC to the normalized vertex rows
and Lemma~\ref{lem:v6v15-weighted-insertion} to the spatial trees and their
degree factorials.  Cancel its \(m!\) by unordered selection.  The remaining
coefficient is \((K\epsilon)^JJ!\) times the allowed rooted-observable
factor.  The ratio of successive terms is \(K\epsilon(J+1)\), so the V6
optimal-truncation argument gives \(O(\epsilon)\).
\end{proof}

The Yang--Mills audit nevertheless changes the claim boundary of V8.  Its
abstract compact-density theorem is conditional on an assembled
coefficient estimate for the exact action and Jacobian.  The companion V12
proves the exact \(SU(2)\) Haar part at the stronger exponential derivative
grade, but explicitly leaves the assembled exact action factor open.

\begin{assessment}[Correct scope of the V8 compact conclusion]
\label{ass:v6v15-V8-scope}
Equations \eqref{eq:v6v15-correct-V8-remainder}--
\eqref{eq:v6v15-correct-V8-eta} and the V8 \(L^p\) product argument remain
valid under their displayed coefficient hypotheses.  The statement that
the complete physical compact-gauge owner has been checked must be read
conditionally: for the actual \(SU(2)\) Wilson owner, the exact Haar factor
is now supported by the companion calculation, while the exact assembled
action-density derivative card remains to be proved.  The \(SU(3)\) and
hypercharge factors require their own group-specific cards.
\end{assessment}

\begin{firewall}[FGMTA and adjacent Gevrey transport have different grades]
\label{fw:v6v15-grade-distinction}
The Yang--Mills V12 correction concerns a fixed-tree estimate intended for
a strong MTA, where no analytic factorial is allowed before Cayley rooting.
The present adjacent-cutoff route permits one \(J!\) in a small
Gevrey-order series and one \(r!\) in observable arity.  The joint compiler
proves that tree degrees do not add another factorial.  These are distinct
norm statements and must not be interchanged.
\end{firewall}

\subsection{Direction selected after the audit}

The three major open candidates named at the end of V14 are the complex
chiral determinant, a nonperturbative ultraviolet gauge trajectory, and a
positive metric sector.  The audit exposes a logically earlier local gate:
the actual compact Wilson action factor must satisfy the assembled
derivative card used by the adjacent transport.  Without it, neither a
candidate non-Abelian hypercharge embedding nor the existing \(SU(2)\)
sector has a fully audited local activity.

The next attack will therefore remain upstream and local.  For one
\(SU(2)\) link it will keep
\[
 \exp\{-\epsilon^{-1}s(\sqrt\epsilon y)+|y|^2/2\}
\]
assembled, derive its exact radial series and integrated complete
derivative rows, and combine it with the exact Haar card without expanding
their logarithms separately.  Only after this card is established should
the program return to the global gauge trajectory or determinant phase.

\begin{theorem}[Sixth-series V15 result]
\label{thm:v6v15-status}
The compulsory audit yields four rigorous conclusions.
\begin{enumerate}
\item The newest Yang--Mills note corrects a factorial-grade overclaim and
      proves an exact-Haar derivative card; the Navier--Stokes endpoint is
      unchanged.
\item The V2 SM transport proof also separated tree and derivative
      factorials incorrectly.  The exact degree-weighted Cayley identity
      \eqref{eq:v6v15-weighted-Cayley} repairs it: after the unordered-owner
      coefficient is used, only the allowed observable factorial remains.
\item The V7 polynomial APGC therefore retains its
      \(O(\epsilon)\) adjacent activity with the repaired joint compiler.
\item The V8 compact activity is conditional on an actual assembled
      action-density card.  The exact \(SU(2)\) Haar factor is available,
      while the exact action factor, other gauge groups, hypercharge
      trajectory, chiral phase, and positive metric measure remain open.
\end{enumerate}
\end{theorem}

\begin{proof}
Item 1 is the recorded read-only audit.  Item 2 is
Lemmas~\ref{lem:v6v15-weighted-Cayley} and
\ref{lem:v6v15-weighted-insertion} with
Theorem~\ref{thm:v6v15-repaired-SIIC}.  Item 3 is
Proposition~\ref{prop:v6v15-polynomial-grade}.  Item 4 is
Assessment~\ref{ass:v6v15-V8-scope} and the audit claim boundary.
\end{proof}

V16 will derive the exact assembled \(SU(2)\) Wilson action-factor card in
the scaled radial variable, test its constants at every derivative order,
and combine it with the exact Haar factor.  If an exponential-grade card is
false, the first obstructing coefficient will be isolated and the program
will move further upstream rather than reusing the invalid factorial split.


\clearpage
\section{V16: the exact scaled \(SU(2)\) Wilson action factor}
\label{sec:v6v16}

\subsection{Normalization and exact radial series}

V15 selected the missing local calculation exposed by the Yang--Mills
audit.  We now perform it in the elementary one-link radial normalization
\begin{equation}
 s(\theta)=1-\cos\theta,
 \qquad 0\le\theta\le\pi.
 \label{eq:v6v16-Wilson-action}
\end{equation}
Put \(r^2=y\cdot y\) and define the action correction relative to the
standard Gaussian by
\begin{equation}
 A_\epsilon(y)
 :=\exp Q_\epsilon(y),
 \qquad
 Q_\epsilon(y)
 :=\frac{r^2}{2}-\frac{1-\cos(\sqrt\epsilon r)}{\epsilon}.
 \label{eq:v6v16-action-factor}
\end{equation}
The expression is entire in the Cartesian components because it is an
entire power series in \(y\cdot y\).

\begin{lemma}[Exact cancellation of the Gaussian quadratic part]
\label{lem:v6v16-Q-series}
For every complex \(y\),
\begin{equation}
 Q_\epsilon(y)
 =\sum_{k=2}^{\infty}
 \frac{(-1)^k\epsilon^{k-1}(y\cdot y)^k}{(2k)!}.
 \label{eq:v6v16-Q-series}
\end{equation}
Consequently,
\begin{equation}
 A_\epsilon(y)=1+\sum_{J=1}^{\infty}\epsilon^JP_J(y),
 \label{eq:v6v16-A-series}
\end{equation}
where
\begin{equation}
 P_J(y)=
 \sum_{m=1}^{J}\frac1{m!}
 \sum_{\substack{k_1,\ldots,k_m\ge2\\
                  \sum_{i=1}^m(k_i-1)=J}}
 \prod_{i=1}^m
 \frac{(-1)^{k_i}(y\cdot y)^{k_i}}{(2k_i)!}.
 \label{eq:v6v16-PJ-exact}
\end{equation}
In particular, \(\deg P_J\le4J\).
\end{lemma}

\begin{proof}
Expand \(1-\cos z\) and substitute \(z=\sqrt\epsilon r\).  Its quadratic
term is \(\epsilon r^2/2\), which cancels the first term of
\eqref{eq:v6v16-action-factor}; this gives
\eqref{eq:v6v16-Q-series}.  Expanding the exponential and collecting total
power \(J=\sum_i(k_i-1)\) gives \eqref{eq:v6v16-PJ-exact}.  Since
\(m\le J\) and \(\sum_i k_i=J+m\), the polynomial degree is
\(2(J+m)\le4J\).
\end{proof}

The order \(J\) in \eqref{eq:v6v16-PJ-exact} is owned exactly once.  No
logarithmic action vertices need to be introduced.

\subsection{A Gevrey-one integrated coefficient bound}

Let \(\gamma_q\) be standard Gaussian measure in fixed dimension \(q\),
and fix \(p<\infty\).  The complete row norm is the same one used in
V6--V8.

\begin{lemma}[Uniform factorial comparison]
\label{lem:v6v16-factorial-comparison}
There is a numerical \(C\) such that, whenever
\(m\le J\), \(k_i\ge2\),
\(\sum_i(k_i-1)=J\), and \(K=\sum_i k_i=J+m\),
\begin{equation}
 \frac{K^K}{m!\prod_{i=1}^m(2k_i)!}
 \le C^J J!.
 \label{eq:v6v16-factorial-comparison}
\end{equation}
\end{lemma}

\begin{proof}
Stirling's two-sided bounds give
\[
 (2k_i)!\ge c^{k_i}k_i^{2k_i},
 \qquad
 m!\ge c^m m^m,
 \qquad
 J!\ge c^J J^J.
\]
Convexity of \(x\log x\), with \(\sum_i k_i=K\), gives
\[
 \prod_i k_i^{k_i}\ge(K/m)^K.
\]
It is therefore enough, up to \(C^J\), to bound
\[
 \frac{K^K m^{2K}}{m^mK^{2K}J^J}
 =\frac{m^{2K-m}}{K^KJ^J}.
\]
Write \(t=m/J\in(0,1]\), so \(K=J(1+t)\).  All powers of \(J\) cancel:
\[
 \frac{m^{2K-m}}{K^KJ^J}
 =\left\{
 \frac{t^{,2+t}}{(1+t)^{1+t}}
 \right\}^{J}.
\]
The bracket is continuous and bounded on \([0,1]\), with limiting value
zero at \(t=0\).  Enlarging the fixed exponential constant proves the
claim.
\end{proof}

\begin{theorem}[Integrated action-coefficient card]
\label{thm:v6v16-action-coefficient-card}
For every fixed \(p<\infty\) and derivative scale \(L>0\), there is
\(K_{p,L}<\infty\) such that, for every \(J\ge1\) and \(d\ge0\),
\begin{equation}
 \frac{\|\nabla^dP_J\|_{L^p(\gamma_q);\,\rm row}}
 {d!L^d}
 \le K_{p,L}^J J!.
 \label{eq:v6v16-action-coefficient-card}
\end{equation}
The constant is uniform in \(J,d\) in fixed fibre dimension.
\end{theorem}

\begin{proof}
For \(K\ge0\), fixed-dimensional Gaussian moments obey
\begin{equation}
 \||Y|^{2K}\|_{L^p(\gamma_q)}
 \le C_p^K K^K.
 \label{eq:v6v16-Gaussian-radial-moment}
\end{equation}
Apply this to a summand of \eqref{eq:v6v16-PJ-exact} and use
Lemma~\ref{lem:v6v16-factorial-comparison}.  For a fixed \(m\), the inner
indices are compositions of \(J\) into \(m\) positive parts
\(k_i-1\), so their number is \(\binom{J-1}{m-1}\).  Summing over \(m\)
costs \(2^{J-1}\), which changes only the exponential base.  Thus
\begin{equation}
 \|P_J\|_{L^p(\gamma_q)}\le K_p^JJ!.
 \label{eq:v6v16-PJ-Lp}
\end{equation}

Each \(P_J\) has degree at most \(4J\).  Apply the V6 complete Gaussian
chaos derivative-row estimate, or its \(L^p\) hypercontractive version used
in V7.  After division by \(d!L^d\), the derivative operation costs at most
an exponential in the degree, hence at most \(C_{p,L}^J\).  Combining this
with \eqref{eq:v6v16-PJ-Lp} proves
\eqref{eq:v6v16-action-coefficient-card}.
\end{proof}

\subsection{Why the stronger exponential derivative grade is false}

The card just proved is Gevrey one in the owned order \(J\).  It is not the
factorial-free fixed-tree card sought by the companion Yang--Mills strong
MTA program.

\begin{proposition}[Quartic obstruction to a pointwise strong card]
\label{prop:v6v16-strong-card-no-go}
There is no constant \(C\), independent of \(0<\epsilon\le\epsilon_0\),
such that
\begin{equation}
 |D^vA_\epsilon(0)|\le C^v
 \label{eq:v6v16-false-strong-card}
\end{equation}
for every derivative order \(v\).  More precisely, along any Cartesian
coordinate and for every \(n\ge1\),
\begin{equation}
 \left|\frac{d^{4n}}{dt^{4n}}A_\epsilon(te_1)
 \bigg|_{t=0}\right|
 \ge\frac{(4n)!}{n!(4!)^n}\epsilon^n.
 \label{eq:v6v16-quartic-lower-bound}
\end{equation}
\end{proposition}

\begin{proof}
Evaluate the entire function on the imaginary line.  By
\eqref{eq:v6v16-Q-series},
\begin{equation}
 A_\epsilon(ite_1)
 =\exp\left\{
 \sum_{k=2}^{\infty}
 \frac{\epsilon^{k-1}t^{2k}}{(2k)!}
 \right\}.
 \label{eq:v6v16-imaginary-positive-series}
\end{equation}
Every Taylor coefficient on the right is nonnegative.  Selecting \(n\)
copies of the quartic term \(\epsilon t^4/4!\) gives coefficient at least
\(\epsilon^n/(n!(4!)^n)\) for \(t^{4n}\).  Since \(i^{4n}=1\), coefficient
comparison proves \eqref{eq:v6v16-quartic-lower-bound}.

Stirling's lower bound shows that the \((4n)\)-th root of its right side is
at least \(c\epsilon^{1/4}n^{3/4}\).  Given a proposed \(C\), first fix
\(\epsilon>0\) and then take \(n\) sufficiently large; the root exceeds
\(C\), contradicting \eqref{eq:v6v16-false-strong-card}.
\end{proof}

\begin{assessment}[The obstruction identifies the correct compiler]
\label{ass:v6v16-correct-compiler}
The exact action factor must not be installed as a single pointwise
analytic owner with derivative activity \(\sup_d\|\nabla^dA\|/(d!L^d)\)
uniform in all orders.  Its quartic exponential produces the unavoidable
growth in \eqref{eq:v6v16-quartic-lower-bound}.  The valid SM route is the
owned total-order expansion \eqref{eq:v6v16-A-series}, followed by the V7
APGC and the V15 joint weighted-tree compiler.  The small factor
\(\epsilon^J\) then pays the single \(J!\).
\end{assessment}

\subsection{Combination with the exact Haar factor}

In the same radial convention, the normalized \(SU(2)\) Haar factor is
\begin{equation}
 H_\epsilon(y)
 =\left\{
 \frac{\sin(\sqrt\epsilon r)}{\sqrt\epsilon r}
 \right\}^{2}
 =\sum_{k=0}^{\infty}
 \frac{(-1)^k2^{2k+1}}{(2k+2)!}
 \epsilon^kr^{2k}.
 \label{eq:v6v16-Haar-series}
\end{equation}
The companion Yang--Mills V12 calculation proves that its integrated
complete derivative rows are exponential-grade after the exact series is
assembled.  The following weaker coefficient statement also follows
directly from \eqref{eq:v6v16-Haar-series} and is sufficient here.

\begin{lemma}[Action--Haar coefficient convolution]
\label{lem:v6v16-action-Haar-convolution}
Write
\begin{equation}
 A_\epsilon H_\epsilon
 =1+\sum_{J\ge1}\epsilon^JW_J.
 \label{eq:v6v16-action-Haar-product}
\end{equation}
Then, for fixed \(p,L\),
\begin{equation}
 \frac{\|\nabla^dW_J\|_{L^p(\gamma_q);\,\rm row}}
 {d!L^d}
 \le K_{p,L}^JJ!.
 \label{eq:v6v16-action-Haar-card}
\end{equation}
\end{lemma}

\begin{proof}
The order-\(k\) Haar coefficient is a radial polynomial of degree \(2k\)
with denominator \((2k+2)!\).  The Gaussian moment and derivative estimate
used in Theorem~\ref{thm:v6v16-action-coefficient-card} bounds its
normalized complete derivative row by \(C^k\).  Since
\(W_J=\sum_{j=0}^JP_jH_{J-j}\), with \(P_0=H_0=1\), the complete Leibniz
rule and the normalization by \(d!\) turn the derivative sum into a
binomial convolution with total coefficient one.  Hence
\[
 \frac{\|\nabla^dW_J\|}{d!L^d}
 \le\sum_{j=0}^JK^jj!C^{J-j}
 \le K_1^JJ!,
\]
after enlarging the base.
\end{proof}

\subsection{The actual \(SU(2)\) compact owner}

On \(0\le\theta\le\pi\),
\begin{equation}
 1-\cos\theta\ge\frac{2}{\pi^2}\theta^2.
 \label{eq:v6v16-coercivity}
\end{equation}
Therefore the V8 uniform \(L^p\) interval applies for
\begin{equation}
 1<p<\frac{\pi^2}{\pi^2-4}.
 \label{eq:v6v16-p-interval}
\end{equation}
The coefficient hypotheses that were abstract in V8 are now supplied for
the action by Theorem~\ref{thm:v6v16-action-coefficient-card} and for the
action--Haar product by
Lemma~\ref{lem:v6v16-action-Haar-convolution}.

\begin{theorem}[Conditional closure of the one-link \(SU(2)\) compact card]
\label{thm:v6v16-SU2-compact-card}
Assume the massive block covariance, sparse-color occupancy, and exact
owner factorization hypotheses of V7--V8.  For the one-link \(SU(2)\)
Wilson density with action \eqref{eq:v6v16-Wilson-action} and Haar factor
\eqref{eq:v6v16-Haar-series}, the polynomial coefficients obey the APGC
input at Gevrey-one grade, while the optimally truncated remainder obeys
the V8 \(L^p\) product card.  With the V15 weighted-tree repair, its
adjacent compact activity satisfies
\begin{equation}
 \eta_n^{SU(2)}
 \le C\epsilon_n+Ce^{-c/\epsilon_n}.
 \label{eq:v6v16-SU2-activity}
\end{equation}
\end{theorem}

\begin{proof}
The action--Haar coefficient card is
\eqref{eq:v6v16-action-Haar-card}; the scalar normalization is a
Gevrey-one series with nonzero constant term and its local reciprocal has
the same grade by the standard coefficient recursion.  Coercivity
\eqref{eq:v6v16-coercivity} gives the strict \(L^p\) interval
\eqref{eq:v6v16-p-interval}.  Apply the V8 central Taylor and coercive-tail
remainder theorem, the V7 APGC, and
Proposition~\ref{prop:v6v15-polynomial-grade}.  Sparse coloring supplies a
Brascamp--Lieb exponent inside the interval.  Adding the polynomial and
exact-remainder activities gives \eqref{eq:v6v16-SU2-activity}.
\end{proof}

\begin{firewall}[One-link normalization is not the complete gauge theory]
\label{fw:v6v16-one-link-scope}
Theorem~\ref{thm:v6v16-SU2-compact-card} checks the local radial action and
Haar density under the stated factorization and massive covariance
hypotheses.  It does not construct the nonperturbative \(SU(2)\) cutoff
trajectory, prove gauge-invariant gluing of link charts, establish the
analogous \(SU(3)\) or hypercharge card, or prove reflection positivity
after chiral and gravitational regulators are included.
\end{firewall}

\begin{theorem}[Sixth-series V16 result]
\label{thm:v6v16-status}
The exact scaled \(SU(2)\) Wilson action correction has an owned
total-order expansion whose order-\(J\) coefficient has degree at most
\(4J\) and integrated derivative rows bounded by \(K^JJ!\).  Combined with
the exact Haar series, it supplies the concrete Gevrey-one coefficient
input left conditional in V8 and yields the adjacent activity
\eqref{eq:v6v16-SU2-activity} under the inherited massive factorization
hypotheses.  A factorial-free pointwise strong derivative card is false:
the exact quartic lower bound \eqref{eq:v6v16-quartic-lower-bound} grows
superexponentially in derivative order.  Thus the valid SM mechanism is
total owned order plus the V15 joint tree compiler, not a strong-MTA claim.
\end{theorem}

\begin{proof}
Use Lemma~\ref{lem:v6v16-Q-series},
Theorem~\ref{thm:v6v16-action-coefficient-card},
Proposition~\ref{prop:v6v16-strong-card-no-go},
Lemma~\ref{lem:v6v16-action-Haar-convolution}, and
Theorem~\ref{thm:v6v16-SU2-compact-card}, subject to the scope firewall.
\end{proof}

V17 will test the same total-order mechanism for a compact semisimple group
in exponential coordinates.  It will identify the exact invariant-tensor
and Weyl-Jacobian hypotheses needed for \(SU(3)\), because the
finite-transition alternative of V14 cannot be used until the ultraviolet
non-Abelian sector has a group-specific local card.


\clearpage
\section{V17: compact semisimple action and Weyl-Jacobian cards}
\label{sec:v6v17}

\subsection{Fixed-group setup}

V16 proves the concrete \(SU(2)\) action card.  The ultraviolet-interface
criterion of V14 also needs a local card for any non-Abelian sector used
above the transition, and the physical color factor requires \(SU(3)\).
This note isolates hypotheses that depend only on a fixed compact group and
verifies them for a faithful finite-dimensional Wilson representation.

Let \(G\) be a fixed compact connected semisimple group, \(\mathfrak g\)
its Lie algebra with an invariant Euclidean norm, and
\(R:G\to U(d_R)\) a faithful finite-dimensional representation.  Normalize
the one-link action so that in exponential coordinates
\begin{equation}
 s_R(X)=\frac12|X|^2+\sum_{k=2}^{\infty}s_{2k}(X),
 \label{eq:v6v17-action-expansion}
\end{equation}
where \(s_{2k}\) is an invariant homogeneous polynomial of degree \(2k\).
This normalization is obtained from a positive multiple of
\(d_R-\operatorname{Re}\operatorname{tr}R(e^X)\).

Choose a bounded fundamental exponential region \(\Omega\subset\mathfrak g\)
whose boundary has Haar measure zero and which contains exactly one
representative of almost every group element.  If the representation has a
finite kernel, use one translated region around each action minimum; below
we state the single-minimum case.

\subsection{Uniform invariant-tensor coefficients}

\begin{lemma}[Wilson character coefficient bound]
\label{lem:v6v17-character-bound}
There is \(C_R<\infty\) such that
\begin{equation}
 |s_{2k}(X)|
 \le\frac{C_R^{2k}|X|^{2k}}{(2k)!}
 \label{eq:v6v17-character-bound}
\end{equation}
for every \(k\ge2\) and \(X\in\mathfrak g\).  The same bound holds for the
complete coefficient-tensor row norm after changing \(C_R\).
\end{lemma}

\begin{proof}
Write \(R(e^X)=e^{R_*(X)}\).  Since the real part removes odd powers,
the degree-\(2k\) term is a fixed normalization times
\(\operatorname{Re}\operatorname{tr}R_*(X)^{2k}/(2k)!\).  In finite
dimension,
\[
 |\operatorname{tr}R_*(X)^{2k}|
 \le d_R\|R_*(X)\|^{2k}
 \le d_R C^{2k}|X|^{2k}.
\]
Absorb \(d_R\) into the exponential base.  Polarization in the fixed
dimension bounds the complete tensor row by another exponential factor.
\end{proof}

Define the scaled action correction
\begin{equation}
 A^G_\epsilon(y)
 :=\exp Q^G_\epsilon(y),
 \qquad
 Q^G_\epsilon(y)
 :=-\sum_{k=2}^{\infty}\epsilon^{k-1}s_{2k}(y).
 \label{eq:v6v17-scaled-action}
\end{equation}

\begin{theorem}[Fixed-group action coefficient card]
\label{thm:v6v17-action-card}
Write
\begin{equation}
 A^G_\epsilon(y)=1+\sum_{J\ge1}\epsilon^JP^G_J(y).
 \label{eq:v6v17-action-coefficients}
\end{equation}
Then \(\deg P^G_J\le4J\), and for every fixed \(p<\infty\) and \(L>0\),
\begin{equation}
 \frac{\|\nabla^dP^G_J\|_{L^p(\gamma_{\dim G});\,\rm row}}
 {d!L^d}
 \le K_{G,p,L}^JJ!.
 \label{eq:v6v17-action-card}
\end{equation}
\end{theorem}

\begin{proof}
Expanding \eqref{eq:v6v17-scaled-action} gives the same composition formula
as \eqref{eq:v6v16-PJ-exact}, with
\((y\cdot y)^{k_i}/(2k_i)!\) replaced by \(s_{2k_i}(y)\).
Lemma~\ref{lem:v6v17-character-bound} bounds the latter by the former up to
\(C_G^{\sum k_i}\).  The constraint
\(\sum_i(k_i-1)=J\) again gives \(m\le J\), total degree
\(2(J+m)\le4J\), and the factorial comparison
Lemma~\ref{lem:v6v16-factorial-comparison}.  Fixed-dimensional Gaussian
moments, the \(2^{J-1}\) composition count, and the V6 complete derivative
row estimate then prove \eqref{eq:v6v17-action-card}.
\end{proof}

\subsection{Exact Weyl-Jacobian factor}

Choose positive roots \(\Delta_+\).  On the exponential region the
normalized Haar Jacobian is
\begin{equation}
 J_G(X)=
 \prod_{\alpha\in\Delta_+}
 \left\{
 \frac{\sin(\alpha(X)/2)}{\alpha(X)/2}
 \right\}^{2},
 \label{eq:v6v17-Weyl-Jacobian}
\end{equation}
with every factor defined continuously at zero.  Formula
\eqref{eq:v6v17-Weyl-Jacobian} is the determinant of the exponential map;
it is used as a Cartesian analytic product, not as a logarithmic root
sum.

\begin{lemma}[Fixed-root Haar coefficient card]
\label{lem:v6v17-Haar-card}
For fixed \(G\), write
\begin{equation}
 J_G(\sqrt\epsilon y)
 =1+\sum_{J\ge1}\epsilon^JH^G_J(y).
 \label{eq:v6v17-Haar-coefficients}
\end{equation}
Then \(\deg H^G_J=2J\), and
\begin{equation}
 \frac{\|\nabla^dH^G_J\|_{L^p(\gamma_{\dim G});\,\rm row}}
 {d!L^d}
 \le C_{G,p,L}^J.
 \label{eq:v6v17-Haar-card}
\end{equation}
\end{lemma}

\begin{proof}
Each positive-root factor has the exact scalar series
\[
 \left(\frac{\sin z}{z}\right)^2
 =\sum_{n\ge0}\frac{(-1)^n2^{2n+1}}{(2n+2)!}z^{2n}.
\]
After \(z=\sqrt\epsilon\,\alpha(y)/2\), the order-\(n\) coefficient has
denominator \((2n+2)!\) and numerator bounded by
\(C_G^n|y|^{2n}\).  Gaussian moments and the normalized derivative row
cost at most \(C^n n^n\), which is absorbed by the factorial denominator.
There are finitely many positive roots.  For total order \(J\), their
weak-composition count is polynomial in \(J\), hence exponential, and the
product rule normalized by \(d!\) is a binomial convolution.  This proves
\eqref{eq:v6v17-Haar-card}.
\end{proof}

For \(SU(3)\), \(|\Delta_+|=3\); all constants in the lemma are therefore
fixed numerical constants once the root normalization is chosen.

\begin{corollary}[Exact action--Jacobian coefficient card]
\label{cor:v6v17-action-Jacobian}
If
\begin{equation}
 A^G_\epsilon(y)J_G(\sqrt\epsilon y)
 =1+\sum_{J\ge1}\epsilon^JW^G_J(y),
 \label{eq:v6v17-full-compact-series}
\end{equation}
then \(\deg W^G_J\le4J\) and
\begin{equation}
 \frac{\|\nabla^dW^G_J\|_{L^p;\,\rm row}}
 {d!L^d}
 \le K_{G,p,L}^JJ!.
 \label{eq:v6v17-full-compact-card}
\end{equation}
\end{corollary}

\begin{proof}
Use \(W^G_J=\sum_{j=0}^JP^G_jH^G_{J-j}\),
Theorem~\ref{thm:v6v17-action-card},
Lemma~\ref{lem:v6v17-Haar-card}, and the normalized complete Leibniz rule.
The convolution is bounded by
\(\sum_{j=0}^JK^jj!C^{J-j}\le K_1^JJ!\).
\end{proof}

\subsection{Coercivity on a fundamental region}

\begin{lemma}[Compact faithful coercivity]
\label{lem:v6v17-faithful-coercivity}
Assume \(R\) is faithful and the Wilson action is normalized as above.
There is \(c_G>0\) such that
\begin{equation}
 s_R(X)\ge c_G|X|^2,
 \qquad X\in\Omega.
 \label{eq:v6v17-faithful-coercivity}
\end{equation}
After reducing \(c_G\) if necessary, one may take \(0<c_G<1/2\).
\end{lemma}

\begin{proof}
The action is nonnegative.  Equality means every eigenvalue of
\(R(e^X)\) is one, hence \(e^X\) lies in the kernel of the faithful
representation and equals the identity.  The chosen fundamental region
contains only \(X=0\) above the identity.  The ratio
\(s_R(X)/|X|^2\) extends continuously to \(X=0\) with value \(1/2\).
On the compact complement of a small ball it has a positive minimum.
Taking the smaller of the two local minima proves
\eqref{eq:v6v17-faithful-coercivity}.
\end{proof}

\begin{corollary}[Fixed-group \(L^p\) interval]
\label{cor:v6v17-Lp-interval}
The scaled compact density has a cutoff-uniform \(L^p\) bound relative to
the normalized Gaussian for every
\begin{equation}
 1<p<p_{*,G}:=(1-2c_G)^{-1}.
 \label{eq:v6v17-Lp-interval}
\end{equation}
\end{corollary}

\begin{proof}
The exponent after taking the \(p\)-th power is bounded by
\(-[pc_G-(p-1)/2]|y|^2\), exactly as in V8.  It is integrable precisely in
the stated open interval.  The finite Weyl-Jacobian product is bounded on
the compact fundamental region.
\end{proof}

\subsection{The \(SU(3)\) local activity}

\begin{theorem}[Conditional \(SU(3)\) compact activity]
\label{thm:v6v17-SU3-activity}
For the faithful fundamental Wilson action of \(SU(3)\), assume the V7--V8
massive block covariance, bounded owner occupancy, sparse-color exact
factorization, and cutoff-uniform chart hypotheses.  Then the exact local
action--Haar density has polynomial APGC input
\eqref{eq:v6v17-full-compact-card}, an exponentially small optimal
\(L^p\) remainder, and adjacent activity
\begin{equation}
 \eta_n^{SU(3)}
 \le C\epsilon_n+Ce^{-c/\epsilon_n}.
 \label{eq:v6v17-SU3-activity}
\end{equation}
The constants are independent of volume and cutoff.
\end{theorem}

\begin{proof}
The fundamental representation is faithful, so
Lemma~\ref{lem:v6v17-faithful-coercivity} and
Corollary~\ref{cor:v6v17-Lp-interval} apply.  The action and Weyl-Jacobian
coefficient hypotheses are
Corollary~\ref{cor:v6v17-action-Jacobian}.  Apply the V8 integrated
Gevrey remainder argument at a fixed \(p<p_{*,G}\), choose the sparse color
spacing so its Brascamp--Lieb exponent lies below \(p_{*,G}\), and apply
the V7 APGC through the repaired V15 weighted-tree compiler.  The
polynomial and exact-remainder contributions give
\eqref{eq:v6v17-SU3-activity}.
\end{proof}

\begin{corollary}[Local compact color and weak cards]
\label{cor:v6v17-color-weak-cards}
Under a common massive-safe factorization, the product
\(SU(3)\times SU(2)\) compact owner has adjacent activity bounded by the
sum of \eqref{eq:v6v17-SU3-activity} and
\eqref{eq:v6v16-SU2-activity}, hence remains \(O(\epsilon_n)\) when both
small parameters are comparable and summable.
\end{corollary}

\begin{proof}
Expand the two commuting positive scalar owner factors.  Assign each
factor its own label and use the finite product version of the square-free
owner expansion.  Activities add after enlarging the fixed local constant.
\end{proof}

\begin{firewall}[A fixed-group local card is not a running trajectory]
\label{fw:v6v17-local-not-trajectory}
Theorem~\ref{thm:v6v17-SU3-activity} proves a conditional local weak-coupling
card.  It does not construct a nonperturbative asymptotically free
trajectory, show that the correlation length is uniform in block units, or
prove the all-scale continuum limit.  The Abelian action has an analogous
local card while its inverse coupling is large; this does not address the
V13 hypercharge obstruction when that inverse coupling decreases.
\end{firewall}

\begin{firewall}[Faithfulness and vacuum charts are part of the data]
\label{fw:v6v17-faithfulness}
If \(R\) has a nontrivial kernel, the action has several minima and the
single-chart proof of Lemma~\ref{lem:v6v17-faithful-coercivity} is false.
One must use a finite vacuum atlas and prove compatible overlap and
topological-sector estimates.  Those estimates are not supplied by the
coefficient card.
\end{firewall}

\begin{theorem}[Sixth-series V17 result]
\label{thm:v6v17-status}
For a fixed compact semisimple group and faithful finite-dimensional Wilson
representation, the exact character expansion supplies factorial
denominators for every invariant action tensor, while the exact
Weyl-Jacobian is a finite product of squared sinc factors.  Their assembled
order-\(J\) coefficient has degree at most \(4J\) and integrated derivative
rows \(K^JJ!\).  Compact faithful coercivity supplies a nontrivial uniform
\(L^p\) interval.  Consequently, under the inherited massive sparse-color
factorization, the physical \(SU(3)\) local owner has the same
\(O(\epsilon)+O(e^{-c/\epsilon})\) adjacent activity as the \(SU(2)\)
owner.  The nonperturbative running trajectory, multiple vacuum charts,
hypercharge ultraviolet replacement, chiral phase, and gravity remain
open.
\end{theorem}

\begin{proof}
Use Lemma~\ref{lem:v6v17-character-bound},
Theorem~\ref{thm:v6v17-action-card},
Lemma~\ref{lem:v6v17-Haar-card},
Corollary~\ref{cor:v6v17-action-Jacobian},
Lemma~\ref{lem:v6v17-faithful-coercivity}, and
Theorem~\ref{thm:v6v17-SU3-activity}, subject to both firewalls.
\end{proof}

V18 will return to the complex chiral determinant.  It will prove the
complete local trace-log derivative card on a simply connected high-gap
chart, then isolate the exact additional hypotheses needed to transport a
phase through a positive modulus measure without confusing local
analyticity with global anomaly cancellation.


\clearpage
\section{V18: local chiral phase owners and complex polymer transport}
\label{sec:v6v18}

\subsection{The high-gap ratio as a complex trace log}

V10 treated the modulus of a high fermion determinant by rooted trace-log
owners.  The same local series controls the phase, provided one works on a
simply connected chart on which the determinant does not vanish.  This
closes the local analytic part of the complex high-gap problem; it does not
trivialize the determinant line globally.

Let \(E\) and \(E'=E+\Delta E\) be finite-cutoff high fermion blocks and put
\begin{equation}
 R:=E^{-1}\Delta E,
 \qquad
 E'=E(I+R).
 \label{eq:v6v18-relative-block}
\end{equation}
Use the exponentially weighted kernel norm
\begin{equation}
 \|K\|_\mu:=
 \max\left\{
 \sup_x\sum_y e^{\mu d(x,y)}\|K_{xy}\|,
 \sup_y\sum_x e^{\mu d(x,y)}\|K_{xy}\|
 \right\}.
 \label{eq:v6v18-kernel-norm}
\end{equation}

\begin{lemma}[Weighted kernel algebra]
\label{lem:v6v18-kernel-algebra}
The norm \eqref{eq:v6v18-kernel-norm} is submultiplicative:
\begin{equation}
 \|K_1K_2\|_\mu\le\|K_1\|_\mu\|K_2\|_\mu.
 \label{eq:v6v18-kernel-algebra}
\end{equation}
\end{lemma}

\begin{proof}
The triangle inequality for the distance gives
\(e^{\mu d(x,z)}\le e^{\mu d(x,y)}e^{\mu d(y,z)}\).  Insert the kernel
product, sum first over \(z\), and then over \(y\).  This proves the row
bound; reversing the order proves the column bound.
\end{proof}

Assume henceforth that
\begin{equation}
 \|R\|_\mu\le q<1.
 \label{eq:v6v18-small-R}
\end{equation}
Then the principal local logarithm is
\begin{equation}
 \ell:=\operatorname{Tr}\log(I+R)
 =\sum_{k\ge1}\frac{(-1)^{k+1}}{k}\operatorname{Tr}R^k.
 \label{eq:v6v18-trace-log}
\end{equation}

\begin{theorem}[Exact rooted complex determinant factorization]
\label{thm:v6v18-rooted-factorization}
Let the internal fibre dimension at one site be \(d_f<\infty\), independent
of volume.  Define
\begin{equation}
 \ell_x:=\sum_{k\ge1}\frac{(-1)^{k+1}}{k}
 \operatorname{tr}_{\rm fib}(R^k)_{xx}.
 \label{eq:v6v18-rooted-log}
\end{equation}
The series converges absolutely and
\begin{align}
 \ell&=\sum_x\ell_x,
 \label{eq:v6v18-log-root-sum}\\
 \frac{\det E'}{\det E}
 &=\prod_x e^{\ell_x},
 \label{eq:v6v18-complex-factorization}\\
 |\ell_x|&\le d_f[-\log(1-q)].
 \label{eq:v6v18-root-bound}
\end{align}
If \(\|R\|_\mu\le C\epsilon\le1/2\), the last bound is
\(C'd_f\epsilon\).
\end{theorem}

\begin{proof}
Lemma~\ref{lem:v6v18-kernel-algebra} gives
\(\|R^k\|_\mu\le q^k\).  In particular,
\(|\operatorname{tr}_{\rm fib}(R^k)_{xx}|\le d_fq^k\), so
\eqref{eq:v6v18-rooted-log} converges absolutely and gives
\eqref{eq:v6v18-root-bound}.  At finite volume, summing diagonal blocks is
the full trace, proving \eqref{eq:v6v18-log-root-sum}.  The convergent
matrix logarithm satisfies
\(\exp\operatorname{Tr}\log(I+R)=\det(I+R)\), which proves
\eqref{eq:v6v18-complex-factorization}.  Finally
\(-\log(1-C\epsilon)\le2C\epsilon\) when \(C\epsilon\le1/2\).
\end{proof}

Write
\begin{equation}
 \ell_x=u_x+iv_x.
 \label{eq:v6v18-real-imaginary}
\end{equation}
Then \(e^{u_x}\) is the modulus owner and \(e^{iv_x}\) is a unit-modulus
phase owner.  Both are parts of one analytic function and must use the same
branch on their common chart.

\subsection{Complete derivative rows}

Assume \(R(\phi)\) is analytic on a complex field chart of fixed radius
\(\rho>0\) and, in the complete kernel derivative-row norm,
\begin{equation}
 \frac{\|D^dR(\phi)\|_{\mu,\rm row}}
 {d!L^d}\le C_R\epsilon
 \qquad(d\ge0).
 \label{eq:v6v18-R-derivatives}
\end{equation}
The row norm includes the exponentially weighted sum over the locations of
the \(d\) field insertions.

\begin{theorem}[Local complex trace-log derivative card]
\label{thm:v6v18-complex-derivative-card}
Under \eqref{eq:v6v18-small-R} and
\eqref{eq:v6v18-R-derivatives}, after enlarging \(L\) by a constant
depending on \((1-q)^{-1}\),
\begin{equation}
 \sup_x
 \frac{\|D^d\ell_x\|_{\mu,\rm row}}
 {d!L_1^d}
 \le C\epsilon
 \qquad(d\ge0).
 \label{eq:v6v18-log-derivatives}
\end{equation}
The same estimate holds for \(u_x\), \(v_x\),
\(e^{\ell_x}-1\), and \(e^{iv_x}-1\), with a changed constant.
No global determinant rank appears.
\end{theorem}

\begin{proof}
The map \(R\mapsto\log(I+R)\) is analytic in the Banach algebra ball
\(\|R\|_\mu<1\).  Its \(a\)-th Frechet derivative is a sum of cyclically
ordered products containing \(a\) perturbations and \(a\) resolvent
factors; equivalently it follows by differentiating the absolutely
convergent series \eqref{eq:v6v18-trace-log}.  Submultiplicativity gives
the bound
\begin{equation}
 \|D_R^a\log(I+R)\|\le(a-1)!(1-q)^{-a}.
 \label{eq:v6v18-Frechet-log-bound}
\end{equation}
The complete Faa di Bruno formula composes this estimate with
\eqref{eq:v6v18-R-derivatives}.  For a partition into \(a\) derivative
blocks, its coefficient \(d!/\prod b_j!\) cancels the derivative
denominators from the \(R\) rows; summing set partitions with weight
\((a-1)!(1-q)^{-a}\) changes only the exponential derivative scale.
At least one differentiated or undifferentiated \(R\) factor remains, so
the adjacent factor \(C\epsilon\) is retained.

Root the fibre trace at \(x\).  Every intermediate spatial sum is paid by
the weighted kernel algebra, and the final trace costs only \(d_f\).
This proves \eqref{eq:v6v18-log-derivatives}.  Real and imaginary parts are
norm contractions.  Finally compose with the scalar exponential; since
\(|\ell_x|=O(\epsilon)\), the same normalized Faa di Bruno estimate applies
to \(e^{\ell_x}-1\) and \(e^{iv_x}-1\).
\end{proof}

\begin{corollary}[The local high-gap phase gate is closed]
\label{cor:v6v18-local-phase-gate}
On a simply connected high-gap chart satisfying the stated analytic kernel
hypotheses, the determinant phase increment has the same exponentially
local factorial-exponential owner activity \(O(\epsilon)\) as the modulus
increment.
\end{corollary}

\begin{proof}
Use \(v_x=\operatorname{Im}\ell_x\) in
Theorem~\ref{thm:v6v18-complex-derivative-card}.
\end{proof}

\subsection{Complex normalization and noncancellation}

Let \(\mu^+\) be the positive normalized measure containing the determinant
modulus and all other positive owners.  If the total retained determinant
phase is \(U=e^{i\Theta}\), the normalized complex functional is
\begin{equation}
 \omega(F)=\frac{\mathbb E_{\mu^+}[FU]}
 {\mathbb E_{\mu^+}[U]},
 \label{eq:v6v18-complex-functional}
\end{equation}
provided its denominator is nonzero.

\begin{lemma}[One-step denominator stability]
\label{lem:v6v18-denominator-stability}
Suppose \(|\mathbb E_{mu^+}U|\ge\delta>0\), \(|U|=|V|=1\), and
\begin{equation}
 a:=\mathbb E_{mu^+}|V-1|<\delta.
 \label{eq:v6v18-phase-distance}
\end{equation}
Then
\begin{equation}
 |\mathbb E_{mu^+}[UV]|\ge\delta-a.
 \label{eq:v6v18-new-denominator}
\end{equation}
For \(\|F\|_\infty\le1\), the two normalized phase functionals differ by
at most
\begin{equation}
 \frac{a}{\delta-a}
 +\frac{a}{\delta(\delta-a)}.
 \label{eq:v6v18-ratio-stability}
\end{equation}
\end{lemma}

\begin{proof}
The reverse triangle inequality and \(|U|=1\) give
\[
 |\mathbb E(UV)|
 \ge|\mathbb EU|-|\mathbb E[U(V-1)]|
 \ge\delta-a.
\]
For the ratio, add and subtract
\(\mathbb E(FU)/\mathbb E(UV)\).  The numerator change is at most \(a\),
the denominator change is at most \(a\), and
\(|\mathbb E(FU)|\le1\).  Insert the two denominator lower bounds.
\end{proof}

This elementary lemma is useful at fixed volume, but \(a\) for a product
of phase owners may grow with volume.  Uniform local correlations require
a complex polymer expansion rather than a global \(L^1\) estimate of the
total phase.

\subsection{A conditional complex polymer gate}

For local phase owners set
\begin{equation}
 K_A^{\rm ph}:=e^{iv_A}-1,
 \qquad
 \prod_Ae^{iv_A}=\prod_A(1+K_A^{\rm ph}).
 \label{eq:v6v18-phase-owner-product}
\end{equation}
This expansion is square-free in the owner labels.  Let
\(z(\Gamma)\in\mathbb C\) denote the resulting connected polymer
activities after expectation in \(\mu^+\).

\begin{theorem}[Complex Kotecky--Preiss noncancellation gate]
\label{thm:v6v18-complex-KP}
Suppose there is \(a_0>0\) such that, for every polymer \(\Gamma\),
\begin{equation}
 \sum_{\Gamma'\not\sim\Gamma}
 |z(\Gamma')|e^{a_0|\Gamma'|}
 \le a_0|\Gamma|.
 \label{eq:v6v18-KP-condition}
\end{equation}
Then the finite-volume phase partition function is nonzero, its logarithm
has an absolutely convergent connected cluster expansion, and connected
correlations with fixed local insertions are bounded uniformly in the
ambient volume by the corresponding rooted cluster sums.
\end{theorem}

\begin{proof}
Expand the square-free product into compatible polymer families.  The
formal logarithm is the sum over connected incompatibility graphs with
Ursell coefficient.  The tree-graph inequality bounds the absolute Ursell
coefficient by a sum over spanning trees.  Root such a tree at one polymer
and remove its leaves.  Condition \eqref{eq:v6v18-KP-condition} pays each
removed incompatible child together with the factor
\(e^{a_0|\Gamma'|}\); induction leaves a convergent geometric majorant at
the root.  Thus the logarithmic series converges absolutely.  Its
exponential equals the polymer partition function near zero activity and,
by convergence throughout the stated domain, never vanishes there.
Rooting the same expansion at a local insertion proves the uniform
correlation bound.
\end{proof}

\begin{theorem}[Conditional all-scale high-phase transport]
\label{thm:v6v18-all-scale-phase}
Assume at every adjacent cutoff step:
\begin{enumerate}
\item the positive modulus measure satisfies the V11--V13 conditional
      massive and bad-block estimates;
\item the high-gap phase owners satisfy
      Theorem~\ref{thm:v6v18-complex-derivative-card};
\item their complex polymer activities satisfy
      \eqref{eq:v6v18-KP-condition} with a uniform margin and rooted norm
      \(\eta_n^{\rm ph}\le C\epsilon_n\); and
\item \(\sum_n\epsilon_n<\infty\).
\end{enumerate}
Then normalized local complex correlations transport across all high-gap
adjacent steps with finite cumulative rooted-cluster loss.
\end{theorem}

\begin{proof}
Theorem~\ref{thm:v6v18-complex-KP} gives a nonzero normalized partition and
a rooted connected expansion at each step.  The local phase derivative
card and the V15 weighted-tree compiler bound its insertion norm by
\(C\epsilon_n\).  The same logarithmic product estimate as V2 gives
\[
 \prod_n(1-K\eta_n^{\rm ph})^{-1}<\infty
\]
after discarding finitely many initial steps and choosing the tail so that
\(K\eta_n^{\rm ph}\le1/2\).  The finitely many initial normalized
functionals are part of the base data.
\end{proof}

\subsection{The global determinant-line boundary}

On overlapping high-gap charts, two logarithms differ by
\(2\pi i\) times an integer winding.  A globally defined chiral determinant
requires these transition functions to define a trivializable determinant
line compatible with gauge transformations.

\begin{firewall}[Local logarithms do not cancel anomalies]
\label{fw:v6v18-local-not-global}
Theorem~\ref{thm:v6v18-complex-derivative-card} is local on a simply
connected zero-free chart.  It does not prove that the phase patches
globally, that large gauge transformations act trivially, or that global
anomalies vanish.  The necessary local trace conditions in V14 and a
global determinant-line trivialization remain required.
\end{firewall}

\begin{firewall}[A complex polymer expansion is not reflection positivity]
\label{fw:v6v18-complex-not-RP}
Absolute convergence and noncancellation of a complex polymer partition do
not imply Osterwalder--Schrader positivity.  Reflection positivity must be
proved in a formulation retaining the fermionic degrees of freedom or an
equivalent positive transfer structure; it cannot be inferred from the
phase-weighted bosonic marginal.
\end{firewall}

\begin{firewall}[Gap crossings belong to the retained low sector]
\label{fw:v6v18-gap-crossings}
When an eigenvalue approaches zero, \(q<1\) and the local logarithm may
fail.  Such modes must be moved into the retained Grassmann block before
the crossing.  No trace-log estimate can replace that exact spectral
transfer.
\end{firewall}

\begin{theorem}[Sixth-series V18 result]
\label{thm:v6v18-status}
On a simply connected high-gap chart, the full complex adjacent determinant
ratio has an exact exponentially local rooted trace-log factorization.  Its
modulus and phase owners both have complete factorial-exponential derivative
activity \(O(\epsilon)\), with only the fixed internal fibre dimension and
no global-rank factor.  A uniform complex Kotecky--Preiss bound is sufficient
to prevent finite-volume normalization zeros and transport local complex
correlations through summable adjacent phase increments.  Global
determinant-line trivialization, spectral crossings, and reflection
positivity remain independent open gates.
\end{theorem}

\begin{proof}
Use Theorems~\ref{thm:v6v18-rooted-factorization},
\ref{thm:v6v18-complex-derivative-card},
\ref{thm:v6v18-complex-KP}, and
\ref{thm:v6v18-all-scale-phase}, subject to the three firewalls.
\end{proof}

V19 will go upstream at determinant-line trivialization.  It will formulate
the finite-lattice cocycle problem, prove an exact patching criterion, and
separate perturbative anomaly cancellation from the remaining global and
topological obstruction before the program turns to the Euclidean metric
sector.


\clearpage
\section{V19: determinant-line patching and the locality of the phase}
\label{sec:v6v19}

\subsection{Finite-lattice configuration quotient}

Fix a finite lattice and remove the determinant-zero set and all chart
boundaries not handled by the retained low-mode sector.  Let
\(\mathcal C^{\rm good}\) be the resulting good configuration space,
\(\mathcal G\) the lattice gauge group, and
\begin{equation}
 \mathcal B^{\rm good}:=\mathcal C^{\rm good}/\mathcal G
 \label{eq:v6v19-orbit-space}
\end{equation}
the good orbit space, with stabilizer strata treated by invariant local
charts.  A chiral determinant is naturally a section of a complex line
bundle over \(\mathcal B^{\rm good}\), rather than an a priori scalar
function.

Choose an open cover \(\{U_\alpha\}\) and nonvanishing local determinant
sections \(D_\alpha\).  On a nonempty overlap set
\begin{equation}
 D_\alpha=g_{\alpha\beta}D_\beta,
 \qquad
 g_{\alpha\beta}:U_\alpha\cap U_\beta\to U(1).
 \label{eq:v6v19-transition-functions}
\end{equation}
They obey
\begin{equation}
 g_{\alpha\alpha}=1,
 \qquad
 g_{\alpha\beta}=g_{\beta\alpha}^{-1},
 \qquad
 g_{\alpha\beta}g_{\beta\gamma}g_{\gamma\alpha}=1.
 \label{eq:v6v19-Cech-cocycle}
\end{equation}

\subsection{Exact patching criterion}

\begin{theorem}[Scalar determinant patching]
\label{thm:v6v19-patching}
There is a global nowhere-zero determinant scalar \(D\) obtained from the
local sections by phase changes if and only if there are continuous
functions \(h_\alpha:U_\alpha\to U(1)\) such that
\begin{equation}
 g_{\alpha\beta}=h_\alpha h_\beta^{-1}
 \label{eq:v6v19-coboundary}
\end{equation}
on every overlap.  In that case
\begin{equation}
 D|_{U_\alpha}=h_\alpha^{-1}D_\alpha
 \label{eq:v6v19-global-determinant}
\end{equation}
is well defined.  Two such global choices differ by a global
\(U(1)\)-valued function on \(\mathcal B^{\rm good}\).
\end{theorem}

\begin{proof}
If \eqref{eq:v6v19-coboundary} holds, then on an overlap
\[
 h_\alpha^{-1}D_\alpha
 =h_\alpha^{-1}g_{\alpha\beta}D_\beta
 =h_\beta^{-1}D_\beta,
\]
so \eqref{eq:v6v19-global-determinant} patches.  Conversely, if a global
nonzero \(D\) exists, define \(h_\alpha:=D_\alpha/D\).  These functions
have unit modulus after the common determinant modulus is fixed, and
\(D_\alpha/D_\beta=g_{\alpha\beta}\) proves
\eqref{eq:v6v19-coboundary}.  The quotient of two global nonzero sections is
a global nonvanishing scalar; after matching their modulus, it is
\(U(1)\)-valued.
\end{proof}

Thus the exact global question is whether the transition cocycle is a
coboundary.  Choosing independent branches of \(\log D_\alpha\) without
checking \eqref{eq:v6v19-coboundary} does not solve it.

\subsection{Connection, curvature, and global holonomy}

On a local chart define the real phase connection
\begin{equation}
 \mathcal A_\alpha:=\operatorname{Im}d\log D_\alpha.
 \label{eq:v6v19-local-connection}
\end{equation}
On overlaps,
\begin{equation}
 \mathcal A_\alpha-\mathcal A_\beta
 =-i g_{\alpha\beta}^{-1}dg_{\alpha\beta}.
 \label{eq:v6v19-connection-overlap}
\end{equation}
Therefore
\begin{equation}
 \mathcal F:=d\mathcal A_\alpha
 \label{eq:v6v19-curvature}
\end{equation}
is a globally defined closed two-form.

\begin{lemma}[Curvature and holonomy obstructions]
\label{lem:v6v19-curvature-holonomy}
If a global phase convention makes \(D\) a constant-phase nowhere-zero
section, then \(\mathcal F=0\) and the phase connection has trivial holonomy
around every closed loop.  Conversely, on a connected component, a flat
phase connection with trivial holonomy defines a global parallel
nowhere-zero section after one fibre value is chosen.
\end{lemma}

\begin{proof}
A constant-phase global section gauges the phase connection to zero, so
its curvature and holonomy vanish.  Conversely, parallel transport a
chosen nonzero vector from a base point.  Flatness makes the result
invariant under contractible path deformations, and trivial holonomy makes
it invariant under changing the path by an arbitrary loop.  It therefore
defines a single-valued global parallel section.
\end{proof}

The second condition cannot be dropped: a flat line bundle can have
nontrivial holonomy on a nonsimply-connected orbit space.  Perturbative
anomaly cancellation controls infinitesimal curvature in gauge directions;
it does not alone prove trivial holonomy for large gauge loops.

\subsection{Tensor products and multiplet cancellation}

For several Weyl multiplets, the full determinant line is the tensor
product of their individual lines.

\begin{proposition}[Exact addition laws for determinant lines]
\label{prop:v6v19-tensor-cancellation}
Let species \(s=1,\ldots,N_f\) have transitions
\(g^{(s)}_{\alpha\beta}\), local phase connections
\(\mathcal A^{(s)}_\alpha\), curvatures \(\mathcal F^{(s)}\), and loop
holonomies \(H^{(s)}(\gamma)\).  The product determinant has
\begin{align}
 g^{\rm tot}_{\alpha\beta}
 &=\prod_sg^{(s)}_{\alpha\beta},
 \label{eq:v6v19-product-transition}\\
 \mathcal A^{\rm tot}_\alpha
 &=\sum_s\mathcal A^{(s)}_\alpha,
 &
 \mathcal F^{\rm tot}&=\sum_s\mathcal F^{(s)},
 \label{eq:v6v19-product-curvature}\\
 H^{\rm tot}(\gamma)&=\prod_sH^{(s)}(\gamma).
 \label{eq:v6v19-product-holonomy}
\end{align}
Thus cancellation is a statement about the complete chiral multiplet; no
individual anomalous determinant need be trivial.
\end{proposition}

\begin{proof}
Tensor-product transition functions multiply.  Taking the logarithmic
differential gives the sum of connections, exterior differentiation gives
the sum of curvatures, and parallel transport in a tensor product multiplies
the scalar holonomies.
\end{proof}

The trace conditions in V14 are the local polynomial consequences of
\(\mathcal F^{\rm tot}\) vanishing in gauge directions.  The remaining
global statement is \eqref{eq:v6v19-product-holonomy}=1 for all gauge loops,
together with an actual local construction of the coboundary
\eqref{eq:v6v19-coboundary}.

\subsection{Gauge-equivariant form of the cocycle}

Before quotienting, let a local determinant transform as
\begin{equation}
 D_\alpha(g\phi)=a_\alpha(g,\phi)D_\alpha(\phi),
 \qquad g\in\mathcal G,
 \label{eq:v6v19-gauge-cocycle}
\end{equation}
whenever both points lie in the chart.  Associativity gives
\begin{equation}
 a_\alpha(g_1g_2,\phi)
 =a_\alpha(g_1,g_2\phi)a_\alpha(g_2,\phi),
 \label{eq:v6v19-group-cocycle}
\end{equation}
up to the chart transition needed when \(g_2\phi\) leaves \(U_\alpha\).

\begin{proposition}[Gauge-invariant trivialization criterion]
\label{prop:v6v19-gauge-trivialization}
Suppose \(h_\alpha\) satisfies the overlap equation
\eqref{eq:v6v19-coboundary}.  The patched determinant
\eqref{eq:v6v19-global-determinant} is gauge invariant if and only if
\begin{equation}
 a_\alpha(g,\phi)
 =\frac{h_\alpha(g\phi)}{h_\alpha(\phi)}
 \label{eq:v6v19-gauge-coboundary}
\end{equation}
whenever one chart is used, with the corresponding transition factors on
mixed-chart pairs.
\end{proposition}

\begin{proof}
Substitute \eqref{eq:v6v19-gauge-cocycle} into
\(D(g\phi)=h_\alpha(g\phi)^{-1}D_\alpha(g\phi)\) and compare it with
\(D(\phi)=h_\alpha(\phi)^{-1}D_\alpha(\phi)\).  Equality is equivalent to
\eqref{eq:v6v19-gauge-coboundary}.  The overlap equation makes the result
independent of the local chart.
\end{proof}

\subsection{Topological triviality is not a locality estimate}

The V18 cluster expansion needs rooted derivative and support bounds for
the patched phase.  The existence of continuous \(h_\alpha\) in
Theorem~\ref{thm:v6v19-patching} gives no such quantitative control.

\begin{definition}[Locality-compatible determinant trivialization]
\label{def:v6v19-local-trivialization}
A trivialization \(\{h_\alpha\}\) is locality compatible at cutoff step
\(n\) if on each chart
\begin{equation}
 h_\alpha=e^{i\chi_\alpha},
 \qquad
 \chi_\alpha=\sum_x\chi_{\alpha,x},
 \label{eq:v6v19-local-phase-sum}
\end{equation}
and the rooted pieces obey a uniform complete derivative row
\begin{equation}
 \sup_x
 \frac{\|D^d\chi_{\alpha,x}\|_{\mu,\rm row}}
 {d!L^d}
 \le C\epsilon_n,
 \qquad d\ge0,
 \label{eq:v6v19-local-trivialization-bound}
\end{equation}
with compatible overlap representatives and a complex polymer norm meeting
the V18 Kotecky--Preiss condition.
\end{definition}

\begin{lemma}[Arbitrary trivializations have no uniform derivative bound]
\label{lem:v6v19-topology-not-locality}
Even for the trivial line bundle over an interval, topological triviality
does not bound derivatives of a chosen phase convention.
\end{lemma}

\begin{proof}
The functions \(h_N(t)=e^{iNt}\), \(0\le t\le1\), are global
trivializations of the same trivial line.  Yet
\(\sup_t|h_N'(t)|=N\).  Thus a topological existence statement cannot
imply a cutoff-uniform locality norm.
\end{proof}

\begin{theorem}[Conditional global chiral-phase criterion]
\label{thm:v6v19-global-phase-criterion}
The V18 high-gap complex transport extends across all gauge charts if the
following are constructed uniformly along the cutoff sequence:
\begin{enumerate}
\item the full Standard Model determinant transition cocycle obeys the
      overlap coboundary equation \eqref{eq:v6v19-coboundary};
\item its gauge cocycle obeys the equivariant coboundary equation
      \eqref{eq:v6v19-gauge-coboundary};
\item the chosen trivialization is locality compatible in the sense of
      Definition~\ref{def:v6v19-local-trivialization}; and
\item determinant-zero crossings are transferred to the retained low
      Grassmann sector before leaving a high-gap chart.
\end{enumerate}
Under these hypotheses, the local rooted phase owners of V18 patch to a
global gauge-invariant complex polymer family, and the conditional
all-scale phase theorem \ref{thm:v6v18-all-scale-phase} applies.
\end{theorem}

\begin{proof}
Items 1 and 2 with
Theorem~\ref{thm:v6v19-patching} and
Proposition~\ref{prop:v6v19-gauge-trivialization} give a global
gauge-invariant determinant on the good orbit space.  Item 3 makes the
chart corrections additional rooted phase owners with the same local
activity grade as the trace-log owners.  Compatibility on overlaps makes
their polymer expansions identical after rephasing.  Item 4 preserves the
high-gap hypothesis needed by V18.  The uniform Kotecky--Preiss margin in
Definition~\ref{def:v6v19-local-trivialization} then invokes
Theorem~\ref{thm:v6v18-all-scale-phase}.
\end{proof}

\begin{firewall}[Local anomaly cancellation is necessary, not sufficient]
\label{fw:v6v19-local-anomaly-boundary}
Vanishing of the polynomial coefficients in V14 is necessary for the
curvature part of the gauge anomaly to cancel.  It does not construct the
coboundaries \eqref{eq:v6v19-coboundary} and
\eqref{eq:v6v19-gauge-coboundary}, rule out nontrivial flat holonomy, or
prove the locality bound \eqref{eq:v6v19-local-trivialization-bound}.
\end{firewall}

\begin{firewall}[The low chiral sector remains dynamical]
\label{fw:v6v19-low-sector}
The determinant-line discussion applies to the high integrated block.
The retained light fermions must still have a gauge-covariant finite-cutoff
Grassmann action, compatible spectral transfer, correct continuum
chirality, and reflection-positive reconstruction.  They cannot be
replaced by a high-gap determinant.
\end{firewall}

\begin{theorem}[Sixth-series V19 result]
\label{thm:v6v19-status}
The global chiral obstruction is exactly a combined overlap and gauge
cocycle problem for the tensor-product Standard Model determinant line.
A scalar gauge-invariant determinant exists precisely when the transition
and gauge phases are coboundaries.  Vanishing curvature does not exclude a
global flat holonomy, and topological triviality does not supply the rooted
locality needed by the complex polymer norm.  A locality-compatible
equivariant trivialization, together with exact low-mode transfer, is
sufficient to globalize the V18 high-gap phase transport.  Construction of
that trivialization and of the retained reflection-positive chiral sector
remains open.
\end{theorem}

\begin{proof}
Use Theorem~\ref{thm:v6v19-patching},
Lemma~\ref{lem:v6v19-curvature-holonomy},
Propositions~\ref{prop:v6v19-tensor-cancellation} and
\ref{prop:v6v19-gauge-trivialization},
Lemma~\ref{lem:v6v19-topology-not-locality}, and
Theorem~\ref{thm:v6v19-global-phase-criterion}, subject to both firewalls.
\end{proof}

V20 will perform the next compulsory companion audit.  It will then turn
to the positive metric sector, beginning with a finite-dimensional
conformal-mode no-go and an exact statement of what a contour or transfer
construction would have to prove before gravity can enter the positive
massive transport architecture.


\clearpage
\section{V20: companion audit and the positive metric boundary}
\label{sec:v6v20}

\subsection{Compulsory read-only audit}

The active companion files inspected at this checkpoint were
\begin{center}
\begin{tabular}{p{0.62\textwidth}r}
\hline
file & bytes\\
\hline
\texttt{Renewal\_Yang\_Mills\_Mass\_Gap\_Research\_Notes\_V15.tex}
& \(221\,304\)\\
\texttt{Renewal\_NS\_Stability\_Research\_Notes\_V31.tex}
& \(887\,537\)\\
\hline
\end{tabular}
\end{center}
Their SHA--256 digests were
\begin{align}
 &\texttt{800a5a73a79c6b36d8a1ad164b43acfa240f0032869600cd840f41d2459e9c63},
 \label{eq:v6v20-YM-hash}\\
 &\texttt{f1121b62238ad5491bb7cf03635ea9934942edf573ed8faaa9ee79e1d38a6696}.
 \label{eq:v6v20-NS-hash}
\end{align}
Neither file was modified.

Yang--Mills V13--V15 proves three boundaries relevant here.  The exact
\(SU(2)\) density has a sharp finite uniform \(L^p\) window.  Repeated
copies of that density at one fully correlated link therefore cannot be
controlled at arbitrary arity by one fixed \(L^p\) exponent.  A
density-free random-cluster forest identity avoids that multiplication,
but its individual ternary tree terms contain lower-cumulant leakage and
fail the desired strong tree weight; only the sum over compensated trees
recovers the complete ternary cancellation.

The NS V31 endpoint is unchanged.  It again warns that a restricted
correlated quantity cannot be replaced by a larger absolute stock before
its cancellation is used.  In the metric problem below, this forbids
replacing lapse, conformal, and constraint correlations by independent
positive Gaussian owners unless an exact factorization has first been
proved.

\begin{assessment}[Scope correction for V8 and V16--V17]
\label{ass:v6v20-density-scope}
The SM sparse-color argument concerns a square-free family of adjacent
\emph{density owners}: each owner label occurs once, and within one color
their distinct anchors are separated.  Under that hypothesis the
replicated covariance exponent can be made close to one.  It does not
control several copies of the same density at one anchor, and it does not
prove an intrinsic strong Wilson MTA for arbitrary same-link observables.
Accordingly, the V16--V17 \(SU(2)\) and \(SU(3)\) results are local
adjacent-owner activity cards under the exact no-duplication factorization;
they are not Yang--Mills mass-gap or strong-MTA results.
\end{assessment}

\begin{firewall}[The repaired Gaussian compiler has a narrower input]
\label{fw:v6v20-Gaussian-compiler-scope}
The V15 degree-weighted Cayley compiler applies after an exact Gaussian
BKAR insertion formula with unordered owners has been established.  The
new Yang--Mills counterexample shows that it cannot be applied to raw
intrinsic Wilson random-cluster trees whose proper-cumulant pieces cancel
only after summing different trees.
\end{firewall}

\subsection{The conformal direction of the Euclidean Einstein action}

Use the Euclidean sign convention
\begin{equation}
 S_E(g)=-\frac1{16\pi G_N}
 \int_M\sqrt g\,(R(g)-2\Lambda).
 \label{eq:v6v20-Einstein-action}
\end{equation}
Let \(g=e^{2\phi}\bar g\) on a closed four-manifold.  The conformal
transformation formula is
\begin{equation}
 R(e^{2\phi}\bar g)
 =e^{-2\phi}
 \left(\bar R-6\bar\Delta\phi-6|\bar\nabla\phi|^2\right).
 \label{eq:v6v20-conformal-curvature}
\end{equation}

\begin{lemma}[Exact conformal action formula]
\label{lem:v6v20-conformal-action}
On a closed four-manifold,
\begin{equation}
 S_E(e^{2\phi}\bar g)
 =-\frac1{16\pi G_N}
 \int_M\sqrt{\bar g}\,
 \left(e^{2\phi}\bar R
 +6e^{2\phi}|\bar\nabla\phi|^2
 -2\Lambda e^{4\phi}\right).
 \label{eq:v6v20-conformal-action}
\end{equation}
\end{lemma}

\begin{proof}
Insert \eqref{eq:v6v20-conformal-curvature} and
\(\sqrt g=e^{4\phi}\sqrt{\bar g}\) into
\eqref{eq:v6v20-Einstein-action}.  On a closed manifold,
\[
 \int e^{2\phi}\bar\Delta\phi
 =-2\int e^{2\phi}|\bar\nabla\phi|^2.
\]
Combining this with the original gradient term gives the coefficient
\(+6\) inside the bracket, hence a negative kinetic term in \(S_E\).
\end{proof}

\begin{theorem}[Cutoff-uniform conformal instability]
\label{thm:v6v20-conformal-no-go}
Let \(M\) contain a flat coordinate patch and let \(\bar g\) be smooth.
There are smooth conformal factors \(\phi_k\), uniformly bounded in
\(L^\infty\), such that
\begin{equation}
 S_E(e^{2\phi_k}\bar g)\longrightarrow-\infty
 \qquad(k\to\infty).
 \label{eq:v6v20-conformal-divergence}
\end{equation}
The scalar-curvature and cosmological terms without derivatives remain
bounded along this sequence.  Hence the real Euclidean Einstein action has
no cutoff-uniform lower bound on a bounded conformal-amplitude chart.
\end{theorem}

\begin{proof}
Choose a nonzero smooth cutoff \(\chi\) supported in the coordinate patch
and fix a sufficiently small amplitude \(a>0\).  In one coordinate set
\begin{equation}
 \phi_k(x)=a\chi(x)\sin(kx_1).
 \label{eq:v6v20-conformal-packet}
\end{equation}
Then \(\|\phi_k\|_\infty\le a\|\chi\|_\infty\), so
\(e^{2\phi_k}\) and \(e^{4\phi_k}\) are bounded above and below uniformly.
The terms
\(\int e^{2\phi_k}\bar R\) and
\(\int e^{4\phi_k}\) are therefore bounded uniformly in \(k\).

On a smaller patch where \(|\chi|\) is bounded below, the leading part of
\(\partial_1\phi_k\) is \(ak\chi\cos(kx_1)\).  The cutoff derivative is
order one.  Averaging \(\cos^2(kx_1)\), or integrating explicitly and using
the Riemann--Lebesgue lemma, gives
\begin{equation}
 \int e^{2\phi_k}|\bar\nabla\phi_k|^2
 \ge c a^2k^2-Cak-C.
 \label{eq:v6v20-gradient-growth}
\end{equation}
The negative kinetic term in
\eqref{eq:v6v20-conformal-action} therefore tends to minus infinity and
dominates the bounded nondifferential terms.
\end{proof}

\begin{corollary}[The positive massive metric ansatz is unavailable]
\label{cor:v6v20-no-positive-metric-Gaussian}
The real Euclidean Einstein weight \(e^{-S_E(g)}Dg\) cannot be placed in
the V2 positive massive Gaussian architecture with cutoff-uniform
coercivity merely by restricting the conformal amplitude to a fixed
bounded interval.
\end{corollary}

\begin{proof}
The configurations \eqref{eq:v6v20-conformal-packet} lie in every fixed
amplitude bound but make \(S_E\) arbitrarily negative as the available
frequency increases with cutoff.  A uniform positive Gaussian comparison
would imply a lower quadratic coercivity bound, contradicting
\eqref{eq:v6v20-conformal-divergence}.
\end{proof}

This is an upstream obstruction to V12's conditional bad-block comparison:
without a positive reference probability, conditional tail probabilities
are not defined.

\subsection{What a conformal contour would have to prove}

At finite cutoff let \(z\in\mathbb C^N\) denote the complexified conformal
coordinates and let \(S_N(z)\) and an observable \(F_N(z)\) be holomorphic
away from a specified singular set.

\begin{theorem}[Finite-dimensional admissible contour criterion]
\label{thm:v6v20-contour-criterion}
Let \(\Gamma_N\) be an oriented real \(N\)-cycle avoiding the singular set.
Assume
\begin{equation}
 \operatorname{Re}S_N(z)\ge c\|z\|^2-C
 \quad(z\in\Gamma_N)
 \label{eq:v6v20-contour-coercivity}
\end{equation}
with \(c>0\), and assume \(F_N\) has at most polynomial growth on the
cycle.  Then
\begin{equation}
 \int_{\Gamma_N}F_N(z)e^{-S_N(z)}\,dz_1\cdots dz_N
 \label{eq:v6v20-contour-integral}
\end{equation}
is absolutely convergent.  If two such cycles are connected by a homology
through the holomorphic region whose boundary at infinity obeys the same
decay, their integrals agree.
\end{theorem}

\begin{proof}
The coercive real part bounds the absolute integrand by a polynomial times
\(e^{-c\|z\|^2+C}\), which is integrable.  For the deformation statement,
truncate the homology at radius \(R\) and apply Stokes' theorem to the
closed holomorphic top form \(F_Ne^{-S_N}dz_1\wedge\cdots\wedge dz_N\).
The boundary at infinity tends to zero by the assumed Gaussian decay.
Letting \(R\to\infty\) proves equality.
\end{proof}

\begin{firewall}[Contour convergence does not imply positivity]
\label{fw:v6v20-contour-not-positive}
Theorem~\ref{thm:v6v20-contour-criterion} constructs a complex integral.
It does not make its weight a positive probability, prove independence of
all allowed contour choices, recover real Lorentzian observables, or imply
reflection positivity.  A formal rotation of the conformal coordinate is
therefore not yet a metric sector for the V14 interface theorem.
\end{firewall}

\subsection{A positive transfer alternative}

Let \((X_h,\nu_h)\) be a cutoff spatial-geometry slice space after gauge
constraints, and let \(K(h,h')\) be a one-step kernel.  Define
\begin{equation}
 (Tf)(h'):=\int K(h,h')f(h)\,d\nu_h(h).
 \label{eq:v6v20-transfer-operator}
\end{equation}

\begin{theorem}[Metric transfer criterion for reflection positivity]
\label{thm:v6v20-transfer-positivity}
Assume \(T\) is a bounded positive self-adjoint operator on
\(L^2(X_h,\nu_h)\), and that finite-time cylinder amplitudes are formed by
composing copies of \(T\) with nonnegative boundary vectors.  Reflection
about a time slice then obeys the Osterwalder--Schrader quadratic inequality
\begin{equation}
 \langle\Theta F,F\rangle_{\rm cyl}\ge0
 \label{eq:v6v20-OS-quadratic}
\end{equation}
for every cylinder functional \(F\) supported on one side of the slice.
If \(T=e^{-aH}\) is a strongly continuous contraction family, its
generator \(H\) is nonnegative self-adjoint.
\end{theorem}

\begin{proof}
Reflecting the negative-time half changes its ordered kernel product into
the Hilbert-space adjoint of the positive-time product.  The glued
quadratic form is therefore
\begin{equation}
 \langle\Theta F,F\rangle_{\rm cyl}
 =\|T^{m/2}\Psi_F\|_{L^2}^2
 \label{eq:v6v20-OS-square}
\end{equation}
for an even number of steps, with the central positive factor treated by
the square root for an odd number.  It is nonnegative.  For a strongly
continuous self-adjoint contraction semigroup, the spectral theorem writes
\(T(a)=e^{-aH}\) with \(H\ge0\).
\end{proof}

The theorem is deliberately a criterion, not a construction.  A metric
kernel must still implement the Hamiltonian and momentum constraints,
handle lapse and shift, avoid negative-norm gauge modes, and be compatible
with the gauge--matter interface data of V14.

\subsection{Exhaustive metric alternatives for the present route}

The conformal no-go leaves three mathematically distinct possibilities.
\begin{enumerate}
\item Construct a complex contour satisfying
      Theorem~\ref{thm:v6v20-contour-criterion}, then separately prove an
      Osterwalder--Schrader reconstruction theorem for the resulting
      complex representation.
\item Construct a positive constrained transfer operator satisfying
      Theorem~\ref{thm:v6v20-transfer-positivity}; derive Euclidean
      correlations from that operator rather than from the unbounded real
      Einstein weight.
\item Modify the ultraviolet gravitational action or fixed point so its
      cutoff measures are coercive and positive, then prove finite-scale
      matching to Einstein gravity by the V14 sewing criterion.  This is a
      changed ultraviolet theory and requires its own relevant-coordinate
      and positivity proof.
\end{enumerate}

\begin{theorem}[Conditional metric-entry criterion]
\label{thm:v6v20-metric-entry}
The gravitational sector can enter the positive SM transport architecture
only after one constructs either:
\begin{enumerate}
\item a cutoff-uniform positive constrained transfer family with local
      rooted kernel bounds, tight slice measures, and the reflection
      property of Theorem~\ref{thm:v6v20-transfer-positivity}; or
\item an ultraviolet positive metric measure satisfying the same local
      bounds and the finite-scale matching conditions of V14.
\end{enumerate}
A convergent complex contour can replace these inputs only if an additional
theorem reconstructs a positive Hilbert space and compatible local
transport from that contour.
\end{theorem}

\begin{proof}
V11--V13 use conditional probabilities and positive comparison measures,
so positivity is required for direct entry.  The first alternative supplies
it through the transfer kernel and
Theorem~\ref{thm:v6v20-transfer-positivity}; the second supplies it by
hypothesis and sews it through V14.  Theorem~\ref{thm:v6v20-contour-criterion}
alone supplies only a complex functional, so the stated reconstruction is
necessary before the same positive arguments apply.
\end{proof}

\begin{firewall}[Higher derivatives and fixed points are new hypotheses]
\label{fw:v6v20-gravity-UV-hypotheses}
Adding curvature-squared terms or naming a non-Gaussian gravitational fixed
point does not by itself meet
Theorem~\ref{thm:v6v20-metric-entry}.  One must prove cutoff-uniform
coercivity or transfer positivity, constraint compatibility, tightness,
local owner bounds, and matching to the Einstein regime.
\end{firewall}

\begin{theorem}[Sixth-series V20 result]
\label{thm:v6v20-status}
The companion audit restricts the SM compact-density result to exact
square-free adjacent owners at separated anchors and forbids its use as a
same-link intrinsic Wilson MTA.  In the gravitational sector, the exact
four-dimensional conformal transformation makes the real Euclidean
Einstein action tend to minus infinity along bounded-amplitude,
high-frequency conformal packets.  The positive massive comparison route
therefore cannot include the real Einstein conformal mode.  A viable
replacement must construct either a positive constrained transfer operator
or a positive ultraviolet metric measure matched at finite scale; a
coercive complex contour is sufficient for convergence but not for
reflection positivity.
\end{theorem}

\begin{proof}
The audit scope is Assessment~\ref{ass:v6v20-density-scope} and its
firewall.  The conformal statement is
Theorem~\ref{thm:v6v20-conformal-no-go} and
Corollary~\ref{cor:v6v20-no-positive-metric-Gaussian}.  The two replacement
criteria are Theorems~\ref{thm:v6v20-contour-criterion},
\ref{thm:v6v20-transfer-positivity}, and
\ref{thm:v6v20-metric-entry}, subject to their firewalls.
\end{proof}

V21 will attack the positive-transfer alternative.  It will separate the
conformal coordinate from lapse and shift at finite cutoff, formulate a
Schur-complement test for positivity of the physical slice kernel, and
identify whether the first obstruction is the Hamiltonian constraint, the
indefinite conformal block, or failure of locality after constraint
reduction.


\clearpage
\section{V21: constrained Gaussian transfer and the locality Schur complement}
\label{sec:v6v21}

\subsection{Finite-dimensional quadratic reduction}

V20 shows that the unconstrained real conformal integral is not a positive
cutoff-uniform measure.  Canonical gravity suggests a different order of
operations: impose lapse and shift constraints, reduce to physical slice
variables, and only then ask for a positive transfer kernel.  At quadratic
order the exact test is a Schur complement.

Let \(x\in\mathbb R^n\) denote variables retained on a slice and
\(y\in\mathbb R^m\) denote variables to be integrated out.  Consider
\begin{equation}
 Q(x,y)=\frac12x^TAx+x^TBy+\frac12y^TCy,
 \label{eq:v6v21-block-quadratic}
\end{equation}
with \(A=A^T\) and \(C=C^T\).

\begin{theorem}[Gaussian Schur-complement criterion]
\label{thm:v6v21-Schur}
Assume \(C>0\).  Then
\begin{equation}
 \int_{\mathbb R^m}e^{-Q(x,y)}\,dy
 =(2\pi)^{m/2}(\det C)^{-1/2}
 \exp\left[-\frac12x^TSx\right],
 \label{eq:v6v21-Gaussian-elimination}
\end{equation}
where
\begin{equation}
 S:=A-BC^{-1}B^T.
 \label{eq:v6v21-Schur-complement}
\end{equation}
Moreover, the full quadratic form is positive definite if and only if
\(C>0\) and \(S>0\).
\end{theorem}

\begin{proof}
Complete the square:
\[
 Q(x,y)=\frac12(y+C^{-1}B^Tx)^TC(y+C^{-1}B^Tx)
 +\frac12x^T(A-BC^{-1}B^T)x.
\]
Translation and the standard Gaussian integral prove
\eqref{eq:v6v21-Gaussian-elimination}.  The same identity shows positivity
of the full form when \(C,S>0\).  Conversely, restriction to \(x=0\) gives
\(C>0\), and minimization over \(y\) gives \(x^TSx>0\) for nonzero \(x\).
\end{proof}

\begin{corollary}[A negative auxiliary block cannot be integrated on the
real contour]
\label{cor:v6v21-negative-block}
If \(C\) has a negative eigenvalue, the integral over \(y\) in
\eqref{eq:v6v21-Gaussian-elimination} diverges for some fixed \(x\).  A real
positive Gaussian reduction is then impossible before a constraint removes
that direction or the integration cycle is changed.
\end{corollary}

\begin{proof}
Along a negative eigenvector \(v\), set \(y=tv\).  The exponent contains a
positive quadratic multiple of \(t^2\), which is not integrable.
\end{proof}

\subsection{Exact linear-constraint test}

Let \(L:\mathbb R^N\to\mathbb R^r\) have constant rank and impose the
linear constraint \(Lq=0\).  Let \(P\) be an isometric column matrix whose
range is \(\ker L\), so \(P^TP=I\).

\begin{lemma}[Positivity after exact constraint reduction]
\label{lem:v6v21-constraint-positivity}
For a symmetric quadratic matrix \(M\), the constrained Gaussian on
\(\ker L\) is integrable and positive if and only if
\begin{equation}
 P^TMP>0.
 \label{eq:v6v21-restricted-positivity}
\end{equation}
Negative directions of \(M\) outside \(\ker L\) are irrelevant to the
constrained integral; any negative or zero direction inside \(\ker L\)
prevents a nondegenerate positive Gaussian.
\end{lemma}

\begin{proof}
Every constrained vector has the unique form \(q=Pz\).  The restricted
quadratic form is \(z^TP^TMPz/2\).  The ordinary finite-dimensional
Gaussian criterion is exactly positive definiteness of this matrix.
\end{proof}

For a nonlinear constraint surface \(\mathcal C(q)=0\), the same test
applies to the tangent kernel of \(D\mathcal C(q)\) at a stationary point.
Uniform nonlinear control additionally requires a global parametrization,
Jacobian bounds, and exclusion of rank-changing strata.

\begin{firewall}[Constraint counting is not constraint positivity]
\label{fw:v6v21-counting-not-positivity}
The statement that lapse and shift impose four constraints does not imply
\eqref{eq:v6v21-restricted-positivity}.  One must identify the actual
finite-cutoff constraint surface, its gauge quotient, and the restricted
quadratic form.  Faddeev--Popov determinants and stabilizer strata are part
of this calculation.
\end{firewall}

\subsection{A sufficient positive Gaussian transfer kernel}

Let the reduced one-step kernel on \(\mathbb R^n\) be
\begin{equation}
 K(x,x')=c\exp\left[-\frac12x^TAx-\frac12{x'}^TAx'
 +x^TBx'\right],
 \label{eq:v6v21-Gaussian-kernel}
\end{equation}
with \(A=A^T>0\), \(B=B^T\).

\begin{theorem}[Quadratic transfer positivity]
\label{thm:v6v21-quadratic-transfer}
If
\begin{equation}
 0\le B<A
 \label{eq:v6v21-transfer-matrix-order}
\end{equation}
as quadratic forms, then \(K\) defines a positive self-adjoint
Hilbert--Schmidt operator on \(L^2(\mathbb R^n,dx)\).  It therefore meets
the finite-step reflection-positivity input of
Theorem~\ref{thm:v6v20-transfer-positivity}.
\end{theorem}

\begin{proof}
Symmetry gives self-adjointness.  Since \(B\ge0\), let
\(B^{1/2}\) be its positive square root.  The power series
\begin{equation}
 e^{x^TBx'}
 =\sum_{k\ge0}\frac1{k!}
 \left\langle(B^{1/2}x)^{\otimes k},
 (B^{1/2}x')^{\otimes k}\right\rangle
 \label{eq:v6v21-feature-expansion}
\end{equation}
is a sum of positive rank-one kernels.  Multiplication on both sides by
\(c^{1/2}e^{-x^TAx/2}\) preserves kernel positivity, so
\(\langle f,Kf\rangle\ge0\).

The square of the kernel is integrable precisely when the block quadratic
matrix
\(\left(\begin{smallmatrix}A&-B\\-B&A\end{smallmatrix}\right)\) is positive.
Its symmetric and antisymmetric sectors are \(A-B\) and \(A+B\), both
positive under \eqref{eq:v6v21-transfer-matrix-order}.  Hence the operator
is Hilbert--Schmidt.
\end{proof}

The condition is sufficient and directly checkable after constraint and
gauge reduction.  A convergent Gaussian kernel with a negative eigenvalue
of \(B\) can have alternating Hermite-sector eigenvalues and need not be a
positive operator.

\subsection{Locality of the Schur complement}

For lattice kernels replace matrices by exponentially weighted operators.
If \(C\) has a spectral gap and exponentially local entries, a
Combes--Thomas estimate gives exponential locality of \(C^{-1}\), and
\eqref{eq:v6v21-Schur-complement} stays local.  A massless constraint block
has a different boundary.

\begin{proposition}[Massless inverse obstruction]
\label{prop:v6v21-massless-inverse}
On a translation-invariant lattice sequence, suppose the Fourier symbol of
\(C\) satisfies
\begin{equation}
 \widehat C(k)=c|k|^2+O(|k|^4)
 \qquad(k\to0),
 \label{eq:v6v21-massless-symbol}
\end{equation}
with \(c>0\).  Then \(C^{-1}\), after removing a possible zero mode, has
no volume-uniform exponentially weighted kernel norm.  Consequently the
Schur term \(BC^{-1}B^T\) is exponentially local only if its numerator
contains a cancellation that removes the pole at \(k=0\).
\end{proposition}

\begin{proof}
A volume-uniform exponentially decaying kernel has a Fourier transform
that extends analytically to a fixed complex neighbourhood of the real
torus.  The inverse symbol
\(\widehat C(k)^{-1}\sim(c|k|^2)^{-1}\) is singular at \(k=0\), so no such
analytic extension and hence no uniform exponential kernel bound exists.
For the Schur term the symbol is
\begin{equation}
 \widehat B(k)\widehat C(k)^{-1}\widehat B(k)^*.
 \label{eq:v6v21-Schur-symbol}
\end{equation}
It can be analytic at the origin only if the factors \(\widehat B(k)\)
vanish or project away the singular directions with sufficient order.
\end{proof}

\begin{corollary}[Ward cancellation required for local reduction]
\label{cor:v6v21-Ward-cancellation}
If \(C\) is massless of Laplace type, a local physical transfer kernel can
result from eliminating its variable only after an exact constraint, Ward,
or multipole identity makes \eqref{eq:v6v21-Schur-symbol} regular at zero
momentum.  Bounding \(B\), \(C^{-1}\), and \(B^T\) separately cannot prove
the required locality.
\end{corollary}

\begin{proof}
Separate absolute bounds retain the \(|k|^{-2}\) pole identified in
Proposition~\ref{prop:v6v21-massless-inverse}.  Regularity can therefore
come only from a cancellation in the assembled product.
\end{proof}

This is the metric analogue of the massless Ward/multipole firewall in V2
and of the companion programs' insistence on retaining correlated sums
before taking absolute norms.

\subsection{Conformal rotation and transfer positivity}

At quadratic order, rotating a negative conformal variable \(y=i\widetilde
y\) changes the sign of its pure quadratic block.  The mixed term becomes
\(ix^TB\widetilde y\).

\begin{lemma}[Rotation alone does not produce a positive kernel]
\label{lem:v6v21-rotation-complex}
If \(B\ne0\), the rotated quadratic weight contains a nonconstant complex
phase.  It is not a real nonnegative density and cannot satisfy
Theorem~\ref{thm:v6v21-quadratic-transfer} without an additional exact
cancellation or reconstruction.
\end{lemma}

\begin{proof}
For real \(x,\widetilde y\), the factor
\(e^{-ix^TB\widetilde y}\) has nonconstant argument when \(B\ne0\).
Therefore the pointwise weight is not real nonnegative.  The feature proof
of Theorem~\ref{thm:v6v21-quadratic-transfer} requires a real positive
cross matrix and does not apply.
\end{proof}

The lemma does not exclude a more subtle positive reconstruction after
integrating the rotated variable.  It identifies the extra theorem that
would be needed.

\subsection{The finite-cutoff gravity gate}

\begin{theorem}[Quadratic constrained-transfer gate]
\label{thm:v6v21-gravity-gate}
At a finite cutoff, a gravitational slice kernel can enter the V20 positive
transfer route if all of the following hold uniformly on a bounded-geometry
chart:
\begin{enumerate}
\item the lapse, shift, gauge, and Hamiltonian constraints have constant
      rank and define a controlled quotient chart;
\item the quadratic form restricted to the physical tangent space obeys
      \eqref{eq:v6v21-restricted-positivity};
\item after all auxiliary variables are eliminated in assembled Schur
      complements, the reduced kernel has the form
      \eqref{eq:v6v21-Gaussian-kernel} with \(0\le B<A\);
\item every massless inverse in the reduction is canceled inside an exact
      Ward or constraint product so that the reduced kernels have uniform
      exponential or admitted endpoint locality; and
\item nonlinear remainders have summable rooted activities and preserve the
      same transfer positivity.
\end{enumerate}
Under these hypotheses the quadratic transfer is positive and local, and
the nonlinear theory has precisely the local inputs required for the V14
interface and V20 reflection criteria.
\end{theorem}

\begin{proof}
Items 1 and 2 use
Lemma~\ref{lem:v6v21-constraint-positivity} to define a positive reduced
Gaussian.  Theorem~\ref{thm:v6v21-Schur} performs each gapped auxiliary
elimination.  Item 4 and
Corollary~\ref{cor:v6v21-Ward-cancellation} preserve the required locality
for massless blocks.  Item 3 with
Theorem~\ref{thm:v6v21-quadratic-transfer} gives a positive self-adjoint
one-step operator, so V20 gives reflection positivity.  Item 5 is exactly
the nonlinear stability hypothesis needed to transport the construction
and sew it to the matter sector.
\end{proof}

\begin{firewall}[The gate is not verified for nonlinear Einstein gravity]
\label{fw:v6v21-gravity-gate-open}
No claim is made that the five hypotheses of
Theorem~\ref{thm:v6v21-gravity-gate} hold for a known finite-cutoff chiral
Standard Model--Einstein regulator.  In particular, constant-rank global
constraint charts, the sign of the reduced cross-time block, massless
locality after solving constraints, and nonlinear transfer positivity are
open.
\end{firewall}

\begin{theorem}[Sixth-series V21 result]
\label{thm:v6v21-status}
At quadratic finite cutoff, imposing constraints removes the Euclidean
conformal instability only if the action is positive on the exact
constraint tangent space.  Integrating auxiliary variables produces the
Schur complement \(A-BC^{-1}B^T\).  A reduced symmetric Gaussian transfer
kernel is positive when its cross-time block satisfies \(0\le B<A\).
If the eliminated constraint block is massless, its inverse is not
uniformly exponentially local; locality survives only through an assembled
Ward or constraint cancellation of the zero-momentum pole.  These sign,
constraint-rank, and locality tests form the first concrete positive-metric
gate, which remains unverified for the nonlinear coupled theory.
\end{theorem}

\begin{proof}
Use Theorem~\ref{thm:v6v21-Schur},
Lemma~\ref{lem:v6v21-constraint-positivity},
Theorem~\ref{thm:v6v21-quadratic-transfer},
Proposition~\ref{prop:v6v21-massless-inverse}, and
Theorem~\ref{thm:v6v21-gravity-gate}, subject to its firewall.
\end{proof}

V22 will test the locality condition on the linearized scalar constraint
block.  It will compute the small-momentum orders of the constraint
couplings before inversion, determine whether the \(|k|^{-2}\) pole
cancels on gauge-invariant slice observables, and stop at the first
uncanceled projector rather than assuming a local reduced metric measure.


\clearpage
\section{V22: linearized gravitational constraints and the residual projector}
\label{sec:v6v22}

\subsection{Vacuum constraint symbols on a flat slice}

V21 reduced the locality question to the small-momentum structure of the
constraint Schur complement.  We now compute that structure for linearized
Einstein gravity on a three-dimensional flat spatial slice.  This is a
quadratic diagnostic; nonlinear constraint charts and the coupled matter
source remain separate.

Let \(h_{ij}\) be a symmetric spatial metric perturbation and
\(\pi_{ij}\) its conjugate momentum.  At nonzero Fourier momentum
\(k\in\mathbb R^3\), the linearized vacuum Hamiltonian and momentum
constraints have symbols
\begin{align}
 \mathcal H(k)h
 &=Q_{ij}(k)h_{ij},
 &
 Q_{ij}(k)&:=k_ik_j-|k|^2\delta_{ij},
 \label{eq:v6v22-Hamiltonian-symbol}\\
 \mathcal M_i(k)\pi
 &=-2ik_j\pi_{ij}.
 \label{eq:v6v22-momentum-symbol}
\end{align}
Overall convention-dependent nonzero constants do not affect the kernels
or projectors below.

Put
\begin{equation}
 P_{ij}(k):=\delta_{ij}-\frac{k_ik_j}{|k|^2}.
 \label{eq:v6v22-transverse-projector}
\end{equation}
Then \(Q=-|k|^2P\), \(P^2=P\), and \(\operatorname{tr}P=2\).

\begin{lemma}[Hamiltonian normal operator]
\label{lem:v6v22-H-normal}
For \(k\ne0\),
\begin{equation}
 Q_{ij}(k)Q_{ij}(k)=2|k|^4.
 \label{eq:v6v22-H-normal}
\end{equation}
The orthogonal projection onto the kernel of the scalar constraint is
\begin{equation}
 (\Pi_H)_{ij,kl}(k)
 =I^{\rm sym}_{ij,kl}-\frac12P_{ij}(k)P_{kl}(k),
 \label{eq:v6v22-H-projector}
\end{equation}
where
\(I^{\rm sym}_{ij,kl}=(\delta_{ik}\delta_{jl}+\delta_{il}\delta_{jk})/2\).
\end{lemma}

\begin{proof}
Since \(Q=-|k|^2P\),
\(Q:Q=|k|^4\operatorname{tr}(P^2)=2|k|^4\).  For a nonzero linear
functional \(h\mapsto Q:h\), the orthogonal projection onto its kernel is
\(I-Q\otimes Q/(Q:Q)\).  Substitution gives
\eqref{eq:v6v22-H-projector}.
\end{proof}

The factors \(|k|^4\) cancel between the two constraint couplings and the
inverse normal operator.  Nevertheless, the resulting projector contains
\(k_ik_j/|k|^2\); magnitude cancellation is not analytic cancellation.

\subsection{The physical transverse--traceless projector}

After quotienting the linearized spatial diffeomorphisms and imposing the
scalar constraint, the vacuum configuration variable is the
transverse--traceless part.  Its symbol is
\begin{equation}
 (\Pi_{TT})_{ij,kl}(k)
 =\frac12\left(P_{ik}P_{jl}+P_{il}P_{jk}\right)
 -\frac12P_{ij}P_{kl}.
 \label{eq:v6v22-TT-projector}
\end{equation}

\begin{lemma}[Algebraic properties of \(\Pi_{TT}\)]
\label{lem:v6v22-TT-properties}
For every \(k\ne0\), \(\Pi_{TT}\) is a self-adjoint orthogonal projector,
and
\begin{align}
 k_i(\Pi_{TT})_{ij,kl}&=0,
 &
 \delta_{ij}(\Pi_{TT})_{ij,kl}&=0.
 \label{eq:v6v22-TT-transverse-trace}
\end{align}
Its range has dimension two.
\end{lemma}

\begin{proof}
Use \(P^2=P\), \(P^T=P\), \(Pk=0\), and \(\operatorname{tr}P=2\).
The first term symmetrically projects both tensor indices to the
two-dimensional transverse plane; subtracting
\(P\otimes P/2\) removes its one-dimensional trace.  The symmetric tensors
on a two-dimensional plane have dimension three, so the traceless range has
dimension two.  The same identities give idempotence and
\eqref{eq:v6v22-TT-transverse-trace}.
\end{proof}

\begin{theorem}[The constraint projection is bounded but nonlocal]
\label{thm:v6v22-projector-nonlocal}
The symbols \(\Pi_H(k)\) and \(\Pi_{TT}(k)\) are uniformly bounded for real
\(k\ne0\), but neither has a direction-independent limit at \(k=0\).
Consequently, their position-space kernels have no volume-uniform
exponential decay bound.
\end{theorem}

\begin{proof}
Orthogonal projectors have operator norm one, proving boundedness.  Along
\(k=(t,0,0)\), the transverse projector tends to
\(\operatorname{diag}(0,1,1)\); along \(k=(0,t,0)\), it tends to
\(\operatorname{diag}(1,0,1)\).  Thus \(P\), and hence the displayed
constraint projectors, have different directional limits at the origin.
A uniformly exponentially decaying lattice kernel has a Fourier symbol
analytic in a fixed complex strip and therefore continuous at the origin.
The directional discontinuity rules out such a bound.
\end{proof}

This identifies the first residual projector promised in V21.  The
constraint pole is canceled enough to produce a bounded projection, but
not enough to produce a local metric coordinate.

\subsection{Endpoint cancellation on physical observables}

Let \(O(k)\) be the Fourier symbol of a linear local observable acting on
symmetric tensors.

\begin{proposition}[Exact endpoint criterion]
\label{prop:v6v22-endpoint-criterion}
If \(O(k)\) is polynomial in \(k\) and obeys
\begin{equation}
 O(k)\Pi_{TT}(k)=O(k)
 \label{eq:v6v22-endpoint-identity}
\end{equation}
for \(k\ne0\), then every occurrence of the constraint projector adjacent
to that endpoint cancels exactly:
\begin{equation}
 O\Pi_{TT}K\Pi_{TT}O^*=OKO^*
 \label{eq:v6v22-endpoint-cancellation}
\end{equation}
for any operator \(K\) acting within the physical range.  No inverse power
of \(|k|\) is introduced by the endpoint projection.
\end{proposition}

\begin{proof}
Equation \eqref{eq:v6v22-endpoint-identity} and self-adjointness give
\(\Pi_{TT}O^*=O^*\).  Substitute these two identities in the left side of
\eqref{eq:v6v22-endpoint-cancellation}.
\end{proof}

A sufficient form of \eqref{eq:v6v22-endpoint-identity} is that the
observable already acts only on transverse traceless tensors.  In raw
metric coordinates that property is a nontrivial Ward identity; it is not
implied by placing derivatives somewhere in the observable.

\begin{corollary}[Metric fields are not admissible local endpoints]
\label{cor:v6v22-metric-not-endpoint}
The identity observable on \(h_{ij}\) fails
\eqref{eq:v6v22-endpoint-identity}.  Therefore unconstrained metric
components cannot be assigned the local endpoint norm of
Proposition~\ref{prop:v6v22-endpoint-criterion}.
\end{corollary}

\begin{proof}
For the identity operator, the condition would read \(\Pi_{TT}=I\), which
is false because the physical range has dimension two inside the
six-dimensional symmetric tensor space.
\end{proof}

\subsection{Matter sourcing and the Newton pole}

With a linearized energy-density source \(\rho(k)\), the scalar constraint
takes the form
\begin{equation}
 Q_{ij}(k)h_{ij}(k)=\kappa\rho(k).
 \label{eq:v6v22-sourced-constraint}
\end{equation}
The minimum Euclidean-norm solution normal to the vacuum constraint
surface is
\begin{equation}
 h^{\perp}_{ij}(k)
 =\frac{\kappa Q_{ij}(k)}{2|k|^4}\rho(k)
 =-\frac{\kappa}{2|k|^2}P_{ij}(k)\rho(k).
 \label{eq:v6v22-sourced-solution}
\end{equation}

\begin{lemma}[Exact source solution]
\label{lem:v6v22-source-solution}
Equation \eqref{eq:v6v22-sourced-solution} solves
\eqref{eq:v6v22-sourced-constraint} and is orthogonal to
\(\ker\mathcal H(k)\).  Every other solution differs from it by an element
of that kernel.
\end{lemma}

\begin{proof}
Contract with \(Q\) and use
\eqref{eq:v6v22-H-normal}.  The result is \(\kappa\rho\).  The normal space
to the kernel of a nonzero linear functional is spanned by its coefficient
tensor \(Q\), proving orthogonality and uniqueness modulo the kernel.
\end{proof}

\begin{theorem}[The coupled scalar constraint is genuinely long range]
\label{thm:v6v22-Newton-pole}
For a general local matter energy density with
\(\rho(0)\ne0\), the solution \eqref{eq:v6v22-sourced-solution} has a
\(|k|^{-2}\) pole.  Eliminating the scalar metric constraint therefore
produces a nonlocal matter interaction and cannot yield a uniformly
exponentially local metric-only owner family.
\end{theorem}

\begin{proof}
The transverse projector is bounded away from zero in some tensor
direction, while the scalar coefficient in
\eqref{eq:v6v22-sourced-solution} is \(\rho(k)/|k|^2\).  If
\(\rho(k)\to\rho(0)\ne0\), this diverges quadratically.  The Fourier symbol
therefore has no analytic extension through zero, whereas uniform
exponential locality would provide one.
\end{proof}

The pole represents the physical long-range gravitational constraint; it
should not be canceled by an invalid absolute bound.

\subsection{Exact divisibility versus finite moment cancellation}

On a translation-invariant lattice let \(\widehat\Delta(k)\) be the
positive Laplacian symbol, with
\(\widehat\Delta(k)=|k|^2+O(|k|^4)\).

\begin{proposition}[Local inverse criterion]
\label{prop:v6v22-local-inverse}
Suppose a source has the exact factorization
\begin{equation}
 \widehat\rho(k)=\widehat\Delta(k)\widehat\sigma(k),
 \label{eq:v6v22-Laplacian-divisibility}
\end{equation}
where \(\widehat\sigma\) extends analytically and boundedly to a fixed
complex strip.  Then
\(\widehat\Delta^{-1}\widehat\rho=\widehat\sigma\) has an exponentially
local kernel with the same strip width.  Conversely, vanishing of only
finitely many moments of \(\rho\) does not by itself imply the exact
analytic divisibility \eqref{eq:v6v22-Laplacian-divisibility}.
\end{proposition}

\begin{proof}
The first statement is algebraic cancellation followed by the standard
Fourier characterization of exponential decay.  For the converse, finite
moment conditions prescribe only finitely many derivatives of
\(\widehat\rho\) at zero.  Exact division by the analytic function
\(\widehat\Delta\) requires the quotient to extend through the complete
zero set with compatible analytic jets; finitely many scalar conditions do
not enforce this in several momentum variables.
\end{proof}

Thus the massless endpoint card must be formulated either with exact Ward
divisibility or in a deliberately power-law topology.  Calling a source
``neutral'' is not enough to infer exponential locality.

\subsection{The correct free metric carrier}

For vacuum linearized gravity, retain the TT field as a massless carrier
rather than integrate its projector into local owners.  Let its covariance
symbol have the schematic form
\begin{equation}
 \widehat C_{TT}(k_0,k)
 =\frac{\Pi_{TT}(k)}{k_0^2+|k|^2}.
 \label{eq:v6v22-TT-covariance}
\end{equation}

\begin{theorem}[Linearized sector decomposition]
\label{thm:v6v22-linearized-decomposition}
At nonzero momentum, the quadratic coupled metric program must separate:
\begin{enumerate}
\item the two TT graviton polarizations, carried by the massless covariance
      \eqref{eq:v6v22-TT-covariance};
\item pure spatial-gauge directions, removed only after the gauge quotient
      and ghost determinant are controlled;
\item the vacuum scalar constraint normal, projected by
      \eqref{eq:v6v22-H-projector}; and
\item the matter-sourced scalar solution
      \eqref{eq:v6v22-sourced-solution}, which carries a physical Newton
      pole unless an exact source identity removes it.
\end{enumerate}
None of these four pieces may be inserted as a positive massive local
owner.  The TT and sourced scalar pieces require massless endpoint norms;
the other two require exact quotient and constraint operations.
\end{theorem}

\begin{proof}
The TT decomposition and its dimension are
Lemma~\ref{lem:v6v22-TT-properties}.  Gauge directions are the range of the
linearized spatial-diffeomorphism symbol.  The vacuum scalar normal and its
projection are Lemma~\ref{lem:v6v22-H-normal}.  The sourced solution and
its pole are Lemma~\ref{lem:v6v22-source-solution} and
Theorem~\ref{thm:v6v22-Newton-pole}.  Their Fourier nonanalyticity excludes
the massive exponential owner norm by
Theorem~\ref{thm:v6v22-projector-nonlocal}.
\end{proof}

\begin{firewall}[Linearized constraint closure is not nonlinear gravity]
\label{fw:v6v22-linear-not-nonlinear}
The calculation is around a flat fixed-topology slice and excludes the
zero-momentum moduli.  Nonlinear constraints couple all Fourier modes, can
change rank at symmetric metrics, and receive matter and cosmological
terms.  The linearized decomposition supplies the correct norm classes but
does not construct the nonlinear constrained transfer operator.
\end{firewall}

\begin{firewall}[Power-law carriers still need reflection positivity]
\label{fw:v6v22-massless-RP}
Identifying the TT and Newton poles does not prove that a chosen lattice
discretization has a positive transfer matrix.  The cross-time kernel must
still satisfy the V21 sign test on the physical quotient, and its endpoint
Ward identities must be compatible with reflection.
\end{firewall}

\begin{theorem}[Sixth-series V22 result]
\label{thm:v6v22-status}
The linearized Hamiltonian constraint couples at order \(|k|^2\), while its
normal inverse is order \(|k|^{-4}\).  Their assembled product is a bounded
but direction-dependent projector, not an exponentially local kernel.
After the spatial gauge quotient, the physical projector is
\(\Pi_{TT}\), whose nonanalyticity disappears at an endpoint only under an
exact transverse--traceless Ward identity.  Coupling matter produces the
uncanceled Newton factor \(\rho(k)/|k|^2\) unless the source is exactly
divisible by the lattice Laplacian.  Therefore the metric sector must use
massless TT and constraint carriers with endpoint norms; it cannot be
folded into the positive massive owner family of V11--V13.
\end{theorem}

\begin{proof}
Use Lemma~\ref{lem:v6v22-H-normal},
Theorem~\ref{thm:v6v22-projector-nonlocal},
Proposition~\ref{prop:v6v22-endpoint-criterion},
Theorem~\ref{thm:v6v22-Newton-pole},
Proposition~\ref{prop:v6v22-local-inverse}, and
Theorem~\ref{thm:v6v22-linearized-decomposition}, subject to both
firewalls.
\end{proof}

V23 will formulate the mixed massive--massless transport norm suggested by
this decomposition.  It will keep TT and Newton kernels as distinguished
carriers, attach local positive owners only through Ward-annihilating
endpoints, and derive the exact power-counting condition required for a
summable adjacent cutoff step.


\clearpage
\section{V23: a mixed massive--massless transport norm}
\label{sec:v6v23}

\subsection{Why the target norm must change}

V22 proves that physical graviton and Newton carriers cannot satisfy the
massive exponential row bound used in V2.  They must remain visible in the
target correlation tree.  Local owner subtrees may still be summed if each
massless attachment is paid either by enough exact endpoint differences or
by decay of its scale activity.  This note makes that alternative precise.

Work on a bounded-degree lattice of Euclidean dimension \(D\ge3\).  Let
\(C^{\rm m}\) be a massive covariance with
\begin{equation}
 |C^{\rm m}(x,y)|\le C e^{-\mu d(x,y)}.
 \label{eq:v6v23-massive-covariance}
\end{equation}
Let a massless covariance have a shell decomposition
\begin{equation}
 C^0=\sum_{j\ge0}C_j^0,
 \qquad L_j:=2^j,
 \label{eq:v6v23-shell-decomposition}
\end{equation}
where a shell is concentrated at distances and wavelengths comparable to
\(L_j\).

\begin{definition}[Massless shell derivative card]
\label{def:v6v23-shell-card}
The decomposition has the order-\(q\) row card if, for every split
\(q=q_x+q_y\),
\begin{equation}
 \sup_x\sum_y
 |D_x^{q_x}D_y^{q_y}C_j^0(x,y)|
 \le K_q L_j^{2-q}.
 \label{eq:v6v23-shell-row}
\end{equation}
Here \(D^q\) is a product of exact lattice differences specified by the
endpoint Ward identity, not an absolute derivative inserted after the
fact.
\end{definition}

The power \(L_j^2\) is universal for a second-order massless Green kernel:
its pointwise size is \(L_j^{2-D}\) and a shell contains \(O(L_j^D)\)
sites.  Every exact difference gains one inverse scale.

\begin{lemma}[Continuum-power shell estimate]
\label{lem:v6v23-shell-estimate}
Suppose, away from the diagonal,
\begin{equation}
 |D_x^{q_x}D_y^{q_y}C^0(x,y)|
 \le K_q(1+d(x,y))^{2-D-q}.
 \label{eq:v6v23-massless-derivative-decay}
\end{equation}
If \(C_j^0\) is obtained with a uniformly smooth dyadic spatial or Fourier
partition, then \eqref{eq:v6v23-shell-row} holds after changing \(K_q\).
\end{lemma}

\begin{proof}
On the \(j\)-th shell the pointwise derivative is
\(O(L_j^{2-D-q})\), while the number of lattice points in the shell is
\(O(L_j^D)\).  Their product is \(O(L_j^{2-q})\).  Smooth overlap changes
the bound by a fixed finite multiplicity.
\end{proof}

\begin{corollary}[Absolute row threshold]
\label{cor:v6v23-row-threshold}
The full differentiated massless kernel has a finite absolute row sum if
\begin{equation}
 q>2.
 \label{eq:v6v23-q-threshold}
\end{equation}
At \(q=2\) the shell majorant is constant and its sum is logarithmically
divergent; for \(q<2\) it grows by a power.
\end{corollary}

\begin{proof}
Sum \(K_q2^{j(2-q)}\) over \(j\ge0\).  It is geometric precisely when
\(2-q<0\).
\end{proof}

The threshold is independent of \(D\) because the volume factor and the
Green-kernel dimension cancel to two powers.

\subsection{Hybrid rooted tree weights}

For local observables \(F_1,\ldots,F_r\) at anchors \(x_i\), define a
hybrid tree monomial by assigning to every edge either a massive weight
\(w^{\rm m}(x,y)=Ke^{-\mu' d(x,y)}\) or a retained massless endpoint
weight \(w^0(x,y)\).  The latter may be the undifferentiated Green kernel
or a Ward-dressed derivative and need not have a finite row sum.

\begin{definition}[Hybrid all-order tree card]
\label{def:v6v23-hybrid-card}
A functional has a hybrid rooted tree card if
\begin{equation}
 |\kappa_r(F_1,\ldots,F_r)|
 \le r!K^r
 \sum_{T\in\mathfrak T_r}
 \prod_{e\in E(T)}w_e(x_{e_-},x_{e_+}),
 \label{eq:v6v23-hybrid-card}
\end{equation}
where each \(w_e\) records its massive or physical massless carrier and all
tensor endpoint contractions are retained.  Massless root--root edges are
not summed over space in this norm.
\end{definition}

This target permits the power-law correlations required by a graviton or
photon.  It asks only that additional internal owner vertices be summable
without degrading those explicit root carriers.

\subsection{Scale-paid owner attachments}

Let an owner \(A\) at scale \(j\) have local activity \(z_{A,j}\) and
anchor \(a(A)\).  If its parent edge is a massless shell with total exact
endpoint order \(q(A,j)\), define
\begin{equation}
 \zeta_j^{(q)}
 :=\sup_x\sum_{A:a(A)=x}z_{A,j}e^{\tau|A|},
 \qquad
 \eta^0:=K\sum_{j\ge0}\zeta_j^{(q_j)}L_j^{2-q_j}.
 \label{eq:v6v23-massless-activity}
\end{equation}
Massive parent edges contribute the V2 parameter \(K\zeta^{\rm m}\).

\begin{theorem}[Mixed insertion-tree compiler]
\label{thm:v6v23-mixed-compiler}
Assume:
\begin{enumerate}
\item the massive edges obey \eqref{eq:v6v23-massive-covariance};
\item every internal massless edge to be removed in leaf contraction obeys
      the shell card \eqref{eq:v6v23-shell-row} with its actual endpoint
      order;
\item the owner expansion is unordered or carries its exponential
      coefficient \(1/m!\); and
\item the normalized derivative rows obey the V15 degree-weighted tree
      hypotheses.
\end{enumerate}
Then the sum over all internal owner vertices preserves the hybrid root
tree card with one-step insertion parameter
\begin{equation}
 \eta=K\zeta^{\rm m}+\eta^0.
 \label{eq:v6v23-mixed-eta}
\end{equation}
If \(\eta<1\), the loss is at most \((1-\eta)^{-1}\).
\end{theorem}

\begin{proof}
Apply the Gaussian or already-transported mixed tree formula and orient
every internal owner subtree toward the first observable root it meets.
For a massive parent edge, summing the child anchor costs its exponential
row norm and the local activity.  For a massless shell edge, Definition
\ref{def:v6v23-shell-card} makes the corresponding cost
\(Kz_{A,j}L_j^{2-q_j}\).  Sum the scale label before proceeding to the next
leaf; this gives \(\eta^0\).

After all internal owners are removed, the unsummed root--root edges are
exactly the hybrid weights in \eqref{eq:v6v23-hybrid-card}.  Count the
tree choices and all derivative-degree factorials jointly by the V15
degree-weighted Cayley identity.  Its owner factorial is canceled by
\(1/m!\) or by unordered-to-ordered conversion, leaving the allowed root
factorial.  Each additional owner therefore costs at most \(\eta\), and
summing their number gives the geometric factor \((1-\eta)^{-1}\).
\end{proof}

\begin{firewall}[Physical massless edges are not leaf-summed]
\label{fw:v6v23-no-root-row-sum}
Theorem~\ref{thm:v6v23-mixed-compiler} does not take an absolute row sum of
an undifferentiated graviton, photon, or Newton edge joining observable
roots.  Such an operation would diverge and erase the required power-law
correlation.  Those edges remain explicit in the target tree.
\end{firewall}

\subsection{The exact scale inequality}

Suppose for large \(j\)
\begin{equation}
 \zeta_j^{(q)}\le C L_j^{-\sigma},
 \qquad \sigma\in\mathbb R.
 \label{eq:v6v23-scale-activity}
\end{equation}

\begin{proposition}[Massless attachment power count]
\label{prop:v6v23-power-count}
The shell contribution to \(\eta^0\) is summable if and only if the
majorant has strictly negative scale exponent,
\begin{equation}
 \boxed{\qquad \sigma+q>2.\qquad}
 \label{eq:v6v23-power-condition}
\end{equation}
At equality the absolute estimate is logarithmic and needs an additional
cancellation or running improvement.
\end{proposition}

\begin{proof}
Equations \eqref{eq:v6v23-massless-activity} and
\eqref{eq:v6v23-scale-activity} give the geometric series
\[
 \sum_jL_j^{2-q-\sigma}.
\]
It converges exactly when \(2-q-\sigma<0\).  At zero exponent every shell
has a nondecaying upper bound.
\end{proof}

Three cases are now sharply separated:
\begin{enumerate}
\item \(q>2\): endpoint Ward differences alone give an absolute row sum;
\item \(q\le2\) and \(\sigma>2-q\): scale decay of the adjacent activity
      pays the remaining long-range row growth;
\item \(q\le2\) and \(\sigma\le2-q\): the present absolute compiler fails,
      and the carrier must remain explicit or acquire a new signed
      cancellation.
\end{enumerate}

\subsection{Exact Ward differences cannot be inferred from conservation}

Let a local matter vertex coupled to a graviton have tensor symbol
\(V_{ij}(k;\mathbf p)\), where \(k\) is the graviton momentum and
\(\mathbf p\) denotes the other legs.  Linearized diffeomorphism invariance
can give
\begin{equation}
 k_iV_{ij}(k;\mathbf p)=0
 \label{eq:v6v23-transverse-Ward}
\end{equation}
after the complete gauge-invariant family of diagrams or owners is
assembled.

\begin{lemma}[Transversality removes longitudinal projectors, not the
physical pole]
\label{lem:v6v23-Ward-pole}
If \eqref{eq:v6v23-transverse-Ward} holds, contraction with
\(P_{ij}=\delta_{ij}-k_ik_j/|k|^2\) removes its longitudinal fraction at
that endpoint.  It does not imply \(V_{ij}(k)=O(|k|)\), and therefore does
not by itself supply a positive endpoint order \(q\) in
\eqref{eq:v6v23-shell-row}.
\end{lemma}

\begin{proof}
Transversality gives \(V_{ij}P_{ik}=V_{kj}\), so the projector denominator
cancels algebraically.  A nonzero tensor independent of \(|k|\) on the
transverse subspace satisfies the same identity.  Hence no radial vanishing
order follows.
\end{proof}

For the physical TT carrier, a conserved stress tensor may therefore have
\(q=0\).  Its internal attachment then costs \(L_j^2\), and the adjacent
activity must decay with \(\sigma>2\) or the TT edge must remain part of the
external hybrid skeleton.

\subsection{All-scale mixed transport}

Let \(\eta_n^{\rm mix}\) denote the parameter
\eqref{eq:v6v23-mixed-eta} at cutoff step \(n\), including all carrier
shells used at that step.

\begin{theorem}[Summable mixed massive--massless transport]
\label{thm:v6v23-all-scale}
Assume every finite-cutoff functional has a compatible hybrid root tree
card, every adjacent step satisfies
Theorem~\ref{thm:v6v23-mixed-compiler}, and
\begin{equation}
 0\le\eta_n^{\rm mix}\le\frac12,
 \qquad
 \sum_n\eta_n^{\rm mix}<\infty.
 \label{eq:v6v23-mixed-summability}
\end{equation}
Then the cutoff-limit cylinder functional retains the hybrid tree card
with finite cumulative loss.  Its retained massless edges may have
power-law decay.
\end{theorem}

\begin{proof}
Iterate the one-step bound and use
\(-\log(1-x)\le2x\) for \(x\le1/2\):
\[
 \prod_n(1-\eta_n^{\rm mix})^{-1}
 \le\exp\left(2\sum_n\eta_n^{\rm mix}\right)<\infty.
\]
Compatibility and the inherited tightness assumption pass each fixed
cylinder coefficient to the limit.  Since the massless root weights were
never row-summed, their stated power-law form is preserved.
\end{proof}

\begin{firewall}[A hybrid norm is not a constructed gravity measure]
\label{fw:v6v23-hybrid-not-construction}
Theorems~\ref{thm:v6v23-mixed-compiler} and
\ref{thm:v6v23-all-scale} are conditional transport results.  They do not
construct the TT covariance on the nonlinear constraint quotient, prove
the actual SM--gravity vertices have \(\sigma+q>2\), establish tightness,
or prove reflection positivity.
\end{firewall}

\begin{firewall}[Absolute borderline summation cannot use a formal Ward
identity]
\label{fw:v6v23-borderline}
At \(\sigma+q=2\), the scale sum is logarithmic.  Cancellation requires an
exact identity for the complete owner or diagram family before norms are
taken.  A perturbative conservation equation imposed on one vertex after
absolute values does not remove the logarithm.
\end{firewall}

\begin{theorem}[Sixth-series V23 result]
\label{thm:v6v23-status}
A second-order massless shell has absolute row size \(L_j^2\); a total of
\(q\) exact endpoint differences changes this to \(L_j^{2-q}\).  The
correct SM--Einstein correlation target therefore leaves physical massless
root edges explicit and sums internal owners in a hybrid tree norm.  If an
owner shell has activity \(L_j^{-\sigma}\), its attachment is summable
exactly under the strict inequality \(\sigma+q>2\).  Stress-tensor
transversality cancels longitudinal projectors but supplies no radial
vanishing by itself.  Verification of the strict inequality for the actual
adjacent gravitational vertices is the next live gate.
\end{theorem}

\begin{proof}
Use Corollary~\ref{cor:v6v23-row-threshold},
Theorem~\ref{thm:v6v23-mixed-compiler},
Proposition~\ref{prop:v6v23-power-count},
Lemma~\ref{lem:v6v23-Ward-pole}, and
Theorem~\ref{thm:v6v23-all-scale}, subject to both firewalls.
\end{proof}

V24 will compute \((\sigma,q)\) for the leading Einstein--matter and pure
graviton adjacent vertices under canonical engineering scaling.  It will
separate a genuine irrelevant activity gain from mere tensor
transversality and determine whether the perturbative Gaussian fixed point
can meet \(\sigma+q>2\).


\clearpage
\section{V24: canonical scaling of the leading massless vertices}
\label{sec:v6v24}

\subsection{Engineering scale convention}

V23 reduces an internal massless attachment at length scale \(L\) to the
strict condition \(\sigma+q>2\).  We now compute the canonical value of
\(\sigma\) for the leading Einstein--matter, pure-gravity, and photon
vertices.  The conclusion concerns the absolute leaf-summation compiler;
it is not a beta-function calculation.

In four Euclidean dimensions, a canonically normalized free boson has
engineering dimension one, a Dirac field has dimension \(3/2\), and a
local operator \(\mathcal O\) of dimension \(\Delta\) with coefficient
\(\lambda_{\mathcal O}\) has dimensionless strength at length \(L\)
\begin{equation}
 g_{\mathcal O}(L)
 :=|\lambda_{\mathcal O}|L^{4-\Delta}.
 \label{eq:v6v24-dimensionless-strength}
\end{equation}

\begin{lemma}[Canonical shell exponent]
\label{lem:v6v24-canonical-exponent}
If \(\Delta>4\) and the dimensionful coefficient is held fixed, then
\begin{equation}
 g_{\mathcal O}(L)=O(L^{-\sigma}),
 \qquad
 \sigma=\Delta-4.
 \label{eq:v6v24-sigma}
\end{equation}
For \(\Delta=4\), \(\sigma=0\); a relevant operator has negative
\(\sigma\).
\end{lemma}

\begin{proof}
This is the definition \eqref{eq:v6v24-dimensionless-strength} rewritten
as a power of \(L\).
\end{proof}

Combining the lemma with V23 gives the canonical absolute attachment test
\begin{equation}
 \Delta-4+q>2,
 \qquad\hbox{equivalently}\qquad
 \Delta+q>6.
 \label{eq:v6v24-dimension-test}
\end{equation}

\subsection{Expansion of the Einstein action}

Write the metric perturbatively as
\begin{equation}
 g_{\mu\nu}=\delta_{\mu\nu}+\kappa h_{\mu\nu},
 \qquad [\kappa]=-1,
 \label{eq:v6v24-metric-expansion}
\end{equation}
where \(h\) has canonical dimension one.  After gauge fixing and removal
of total derivatives, the two-derivative Einstein action has the schematic
homogeneous expansion
\begin{equation}
 S_{EH}
 =S_2(h)+\sum_{m\ge3}\kappa^{m-2}
 \int \mathcal V_m(h,\partial h),
 \qquad
 \mathcal V_m\sim h^{m-2}(\partial h)^2.
 \label{eq:v6v24-Einstein-expansion}
\end{equation}

\begin{lemma}[Dimension of the \(m\)-graviton vertex]
\label{lem:v6v24-graviton-dimension}
The operator \(\mathcal V_m\) has engineering dimension
\begin{equation}
 \Delta_m=m+2,
 \label{eq:v6v24-graviton-dimension}
\end{equation}
and dimensionless strength
\begin{equation}
 g_m(L)\asymp\left(\frac{|\kappa|}{L}\right)^{m-2}.
 \label{eq:v6v24-graviton-strength}
\end{equation}
Thus its canonical shell exponent is \(\sigma_m=m-2\).
\end{lemma}

\begin{proof}
There are \(m\) fields of dimension one and two derivatives, giving
\(m+2\).  The coefficient \(\kappa^{m-2}\) has dimension \(-(m-2)\), so
\eqref{eq:v6v24-graviton-strength} is the corresponding dimensionless
combination at length \(L\).
\end{proof}

\begin{theorem}[Cubic and quartic gravity fail the absolute leaf test]
\label{thm:v6v24-gravity-leaf-failure}
For a physical TT graviton attachment, tensor transversality supplies no
radial endpoint order, so \(q=0\) unless an additional soft identity is
proved.  Under canonical fixed-\(\kappa\) scaling:
\begin{align}
 m=3:&\qquad \sigma+q=1<2,
 \label{eq:v6v24-cubic-failure}\\
 m=4:&\qquad \sigma+q=2.
 \label{eq:v6v24-quartic-borderline}
\end{align}
The cubic attachment has a power divergence in the V23 absolute shell
sum, and the quartic attachment is logarithmically borderline.  Hence the
Einstein interaction cannot be treated entirely as leaf-summed local
owners around a free massless TT carrier.
\end{theorem}

\begin{proof}
Lemma~\ref{lem:v6v24-graviton-dimension} gives
\(\sigma_3=1\) and \(\sigma_4=2\).  V23 proves that transversality cancels
longitudinal projectors but does not imply vanishing with \(|k|\), so the
physical TT value is \(q=0\) without further input.  Apply the strict
criterion \(\sigma+q>2\).
\end{proof}

Higher vertices \(m\ge5\) have \(\sigma_m>2\) canonically, but the cubic
vertex already prevents closure of the absolute owner family.

\subsection{The graviton--matter vertex}

For a matter action with stress tensor \(T_{\mu\nu}\), the leading metric
coupling is
\begin{equation}
 S_{hT}=\frac\kappa2\int h_{\mu\nu}T_{\mu\nu}.
 \label{eq:v6v24-hT-coupling}
\end{equation}
The stress tensor has engineering dimension four, so the operator \(hT\)
has dimension five.

\begin{proposition}[Stress-tensor attachment exponent]
\label{prop:v6v24-stress-exponent}
At fixed \(\kappa\), the graviton--matter vertex has
\begin{equation}
 \sigma_{hT}=1.
 \label{eq:v6v24-stress-sigma}
\end{equation}
The Ward identity \(k_\mu T_{\mu\nu}=0\) gives transverse tensor
contraction but does not imply \(T_{\mu\nu}=O(|k|)\).  Hence the physical
TT attachment has \(q=0\) in general and fails
\(\sigma+q>2\).
\end{proposition}

\begin{proof}
The dimension-five operator and
Lemma~\ref{lem:v6v24-canonical-exponent} give \(\sigma=1\).  A constant
nonzero tensor on the transverse subspace obeys the conservation identity
at a fixed nonzero momentum direction, so conservation supplies no radial
power.  Apply Proposition~\ref{prop:v6v23-power-count}.
\end{proof}

This is consistent with the Newton pole in V22: a generic nonzero total
energy density sources a long-range scalar field rather than a locally
summable perturbation.

\subsection{The photon--current vertex}

The unbroken electromagnetic field is also massless.  Its leading matter
coupling is
\begin{equation}
 S_{AJ}=e\int A_\mu J_\mu,
 \label{eq:v6v24-photon-current}
\end{equation}
where \([A]=1\), \([J]=3\), and \(e\) is dimensionless.

\begin{proposition}[Marginal photon attachment]
\label{prop:v6v24-photon-exponent}
At canonical scaling the photon--current vertex has \(\sigma=0\).
Current conservation removes the longitudinal photon projector but does
not supply radial vanishing, so generically \(q=0\).  Its absolute internal
shell attachment has row cost \(L_j^2\) and is not summable by the V23
owner compiler.
\end{proposition}

\begin{proof}
The operator \(A_\mu J_\mu\) has dimension four, so
Lemma~\ref{lem:v6v24-canonical-exponent} gives \(\sigma=0\).  The same
argument as V23's transversality lemma gives \(q=0\).  The shell cost is
therefore \(L_j^{2-0-0}\).
\end{proof}

Thus a full SM--Einstein continuum cannot be obtained by adding photon and
graviton propagators to a massive cluster expansion as if all their
interaction vertices were small local owners.  Both require an interacting
massless skeleton or a genuinely multiscale renormalization argument.

\subsection{Running couplings and the required improvement}

Let the dimensionless Newton coupling at length \(L\) be
\begin{equation}
 \widehat\kappa(L):=\frac{|\kappa(L)|}{L}.
 \label{eq:v6v24-running-kappa}
\end{equation}

\begin{corollary}[Running required by absolute cubic attachment]
\label{cor:v6v24-running-required}
With \(q=0\), the cubic graviton or \(hT\) attachment can satisfy the V23
absolute shell criterion only if
\begin{equation}
 \widehat\kappa(L)=O(L^{-2-\delta})
 \label{eq:v6v24-running-condition}
\end{equation}
for some \(\delta>0\) along the shell direction being summed.  Fixed
dimensionful \(\kappa\) gives only \(\widehat\kappa(L)=O(L^{-1})\).
\end{corollary}

\begin{proof}
Equation \eqref{eq:v6v24-running-condition} is
\(\sigma>2\), the V23 condition at \(q=0\).  The fixed-coupling comparison
is immediate from \eqref{eq:v6v24-running-kappa}.
\end{proof}

This condition is sufficient for this particular absolute infrared shell
sum; it is not asserted to be the physical running of Newton's constant.
The natural alternative is to stop trying to leaf-sum the leading
interactions and retain them inside the massless skeleton.

\subsection{Ultraviolet Gaussian scaling}

At momentum scale \(\mu=L^{-1}\), the dimensionless Newton coupling is
\begin{equation}
 g_N(\mu)\asymp G_N\mu^2
 \asymp(\kappa\mu)^2.
 \label{eq:v6v24-dimensionless-Newton}
\end{equation}

\begin{theorem}[Gaussian ultraviolet power-counting obstruction]
\label{thm:v6v24-Gaussian-UV}
If \(G_N\) tends to a fixed nonzero dimensionful constant as
\(\mu\to\infty\), then \(g_N(\mu)\to\infty\).  Hence the ultraviolet
cutoff cannot be removed by a perturbation series whose small parameter is
the Gaussian-fixed-point Newton coupling.  A continuum construction must
instead provide a non-Gaussian ultraviolet trajectory, a changed
ultraviolet theory matched through V14, or a nonperturbative mechanism not
organized by small \(g_N\).
\end{theorem}

\begin{proof}
Equation \eqref{eq:v6v24-dimensionless-Newton} grows quadratically in
\(\mu\) when \(G_N\) stays nonzero.  A perturbative expansion requiring
uniformly small \(g_N\) therefore loses its control parameter.
\end{proof}

\begin{firewall}[Power counting is not an RG nonexistence theorem]
\label{fw:v6v24-power-not-RG}
Theorem~\ref{thm:v6v24-Gaussian-UV} excludes only the stated Gaussian
small-coupling construction.  It does not rule out an interacting fixed
point, prove or disprove asymptotic safety, or construct a UV completion.
Those alternatives require exact cutoff measures, relevant manifolds, and
the V14 matching data.
\end{firewall}

\subsection{Promotion to the massless skeleton}

\begin{definition}[Non-summable skeleton vertex]
\label{def:v6v24-skeleton-vertex}
A vertex family is promoted to the massless skeleton if its incident
physical massless edges and scale labels remain explicit in the target
connected object, rather than being removed by the absolute leaf sum of
Theorem~\ref{thm:v6v23-mixed-compiler}.
\end{definition}

\begin{proposition}[Necessary promotion at canonical scaling]
\label{prop:v6v24-promotion}
At canonical fixed-coupling scaling, the cubic graviton,
graviton--stress-tensor, and photon--current vertices must be promoted to
the massless skeleton unless an additional exact endpoint identity or
running estimate changes their \((\sigma,q)\) values so that
\(\sigma+q>2\).
\end{proposition}

\begin{proof}
Theorem~\ref{thm:v6v24-gravity-leaf-failure} and
Propositions~\ref{prop:v6v24-stress-exponent} and
\ref{prop:v6v24-photon-exponent} show that all three fail the strict V23
criterion.  The absolute leaf compiler is therefore unavailable by its own
hypotheses.  Retaining the vertices explicitly is the definition of
promotion.
\end{proof}

The new target is an intrinsic connected expansion for the interacting
massless skeleton with its complete Ward identities imposed before norms.
This is analogous in proof architecture, though not in analytic content,
to the Yang--Mills companion's need to compensate raw forest trees before
estimating them.

\begin{firewall}[Skeleton promotion is a reformulation, not closure]
\label{fw:v6v24-promotion-not-proof}
Definition~\ref{def:v6v24-skeleton-vertex} prevents an invalid divergent
summation.  It does not bound the resulting interacting massless skeleton,
prove its renormalization flow, establish positivity, or construct its
continuum limit.
\end{firewall}

\begin{theorem}[Sixth-series V24 result]
\label{thm:v6v24-status}
For a four-dimensional operator of dimension \(\Delta\), canonical
large-length activity has \(\sigma=\Delta-4\), so the V23 absolute
massless attachment condition is \(\Delta+q>6\).  Cubic gravity and the
graviton--stress-tensor vertex have \((\sigma,q)=(1,0)\); quartic gravity
is borderline \((2,0)\); and the photon--current vertex has \((0,0)\).
The leading interactions therefore cannot be leaf-summed as local owners.
They must remain in an interacting massless skeleton unless new exact Ward
orders or running estimates improve the inequality.  At short distances,
fixed Newton constant also makes the dimensionless Gaussian coupling grow,
so Gaussian perturbative cutoff removal is unavailable.
\end{theorem}

\begin{proof}
Use Lemmas~\ref{lem:v6v24-canonical-exponent} and
\ref{lem:v6v24-graviton-dimension},
Theorem~\ref{thm:v6v24-gravity-leaf-failure},
Propositions~\ref{prop:v6v24-stress-exponent} and
\ref{prop:v6v24-photon-exponent},
Theorem~\ref{thm:v6v24-Gaussian-UV}, and
Proposition~\ref{prop:v6v24-promotion}, subject to both firewalls.
\end{proof}

V25 will perform the next compulsory companion audit and then formulate the
first intrinsic massless-skeleton norm.  The norm must retain scale labels,
impose complete Ward projections before absolute values, and distinguish
renormalizable photon vertices from canonically irrelevant but
UV-growing Einstein vertices.


\clearpage
\section{V25: companion audit and an intrinsic scale-score skeleton}
\label{sec:v6v25}

\subsection{Compulsory read-only audit}

The active companion files inspected at this checkpoint were
\begin{center}
\begin{tabular}{p{0.62\textwidth}r}
\hline
file & bytes\\
\hline
\texttt{Renewal\_Yang\_Mills\_Mass\_Gap\_Research\_Notes\_V19.tex}
& \(282\,336\)\\
\texttt{Renewal\_NS\_Stability\_Research\_Notes\_V31.tex}
& \(887\,537\)\\
\hline
\end{tabular}
\end{center}
Their SHA--256 digests were
\begin{align}
 &\texttt{936d939fdc70f9c609cc19cce4b906b1a484b2b46e3821c020d41b5b8ad3cfb2},
 \label{eq:v6v25-YM-hash}\\
 &\texttt{f1121b62238ad5491bb7cf03635ea9934942edf573ed8faaa9ee79e1d38a6696}.
 \label{eq:v6v25-NS-hash}
\end{align}
Neither file was modified.

Yang--Mills V16--V19 proves that raw forest terms carry proper-cumulant
leakage, constructs an exact balancing projection, and then finds a more
intrinsic comparison: differentiating a complete cumulant along a path of
probability laws inserts one centered score.  This removes ordered-history
factorials and keeps the complete connected cancellation intact.  Its
remaining analytic input is a uniform interpolation gap and a
representation-uniform connected score card.

The NS endpoint remains unchanged.  Its relevant discipline is the same:
retain the restricted correlated quantity until the signed cancellation is
used.  For the present massless skeleton, the complete scale score must be
inserted before photon, graviton, constraint, and counterterm pieces are
normed separately.

\begin{firewall}[The companion score estimate is not imported]
\label{fw:v6v25-no-score-transfer}
Only the elementary differentiation identity is re-proved below.  No
Wilson heat-kernel gap, Yang--Mills representation estimate, or
Navier--Stokes cancellation is imported.  The SM--Einstein score contains
different massless fields, constraints, chiral phases, and running
couplings and needs its own analytic card.
\end{firewall}

\subsection{The complete connected scale-score identity}

Let \((\Omega,\lambda)\) be a finite-cutoff configuration space and let
\begin{equation}
 d\nu_t=w_t\,d\lambda,
 \qquad 0\le t\le1,
 \label{eq:v6v25-law-path}
\end{equation}
be a differentiable path of strictly positive normalized densities.  Put
\begin{equation}
 \mathcal S_t:=\partial_t\log w_t.
 \label{eq:v6v25-score}
\end{equation}
Normalization gives \(\mathbb E_t\mathcal S_t=0\).

\begin{theorem}[Connected scale-score insertion]
\label{thm:v6v25-score-cumulant}
For bounded, \(t\)-independent observables \(F_1,\ldots,F_r\),
\begin{equation}
 \boxed{\qquad
 \frac d{dt}\kappa_r^{\nu_t}(F_1,\ldots,F_r)
 =\kappa_{r+1}^{\nu_t}
 (F_1,\ldots,F_r,\mathcal S_t).
 \qquad}
 \label{eq:v6v25-score-cumulant}
\end{equation}
If the observables also depend on \(t\), add the \(r\) multilinear terms
obtained by replacing one \(F_i\) by \(\dot F_i\).
\end{theorem}

\begin{proof}
For one moment,
\[
 \frac d{dt}\mathbb E_tF
 =\mathbb E_t(F\mathcal S_t).
\]
Differentiate the partition formula for the cumulant.  The product rule
chooses one block and inserts \(\mathcal S_t\) in its moment.  In the
partition expansion of the right side of
\eqref{eq:v6v25-score-cumulant}, a singleton score block vanishes because
the score is centered.  Every remaining partition is obtained uniquely by
adjoining the score label to one block of a partition of \([r]\), with the
same Möbius coefficient.  This proves the identity.  Moving inputs follow
from multilinearity.
\end{proof}

\begin{corollary}[Exact adjacent functional comparison]
\label{cor:v6v25-adjacent-comparison}
If \(\nu_0=\nu_n\) and \(\nu_1=\nu_{n+1}\), then
\begin{equation}
 \kappa_r^{\nu_{n+1}}(F_1,\ldots,F_r)
 -\kappa_r^{\nu_n}(F_1,\ldots,F_r)
 =\int_0^1
 \kappa_{r+1}^{\nu_t}
 (F_1,\ldots,F_r,\mathcal S_t)\,dt.
 \label{eq:v6v25-adjacent-comparison}
\end{equation}
\end{corollary}

\begin{proof}
Integrate Theorem~\ref{thm:v6v25-score-cumulant} over \(t\).
\end{proof}

There is one score insertion at every arity.  No expansion into an ordered
family of cubic gravitational or photon--current vertices is required.

\subsection{Normalized action paths and matched scores}

For a positive action path
\begin{equation}
 w_t=Z_t^{-1}e^{-S_t},
 \label{eq:v6v25-action-path}
\end{equation}
the score is
\begin{equation}
 \mathcal S_t
 =-\dot S_t+\mathbb E_t\dot S_t.
 \label{eq:v6v25-centered-action-score}
\end{equation}

\begin{lemma}[Vacuum cancellation is automatic]
\label{lem:v6v25-vacuum-cancellation}
Adding an arbitrary scalar function \(c(t)\) to \(S_t\) does not change the
score \eqref{eq:v6v25-centered-action-score}.  More generally, every part
of \(\dot S_t\) independent of the fields cancels before the connected
cumulant is estimated.
\end{lemma}

\begin{proof}
Both \(-\dot S_t\) and \(\mathbb E_t\dot S_t\) change by \(-\dot c(t)\)
and \(+\dot c(t)\), respectively.  Their sum is unchanged.
\end{proof}

Relevant and marginal field-dependent components do not cancel this way.
They must be matched by the finite-dimensional running couplings before the
residual score can be expected to be summable.

\begin{definition}[Matched scale score]
\label{def:v6v25-matched-score}
Let \(\mathcal R_t\) be the chosen finite-dimensional span of relevant and
marginal local operators and \(P_{\mathcal R_t}\) its renormalization
projection.  The path is matched if its running parameters make
\begin{equation}
 P_{\mathcal R_t}\dot S_t=0
 \label{eq:v6v25-relevant-matching}
\end{equation}
after the prescribed beta-function and counterterm coordinates are included.
The centered residual of \((I-P_{\mathcal R_t})\dot S_t\) is the matched
score \(\mathcal S_t^{\rm mat}\).
\end{definition}

Equation \eqref{eq:v6v25-relevant-matching} is an exact renormalization
condition, not a claim that the required trajectory exists.

\subsection{Ward-first massless skeleton amplitudes}

Let \(\mathfrak C_r\) denote the vector space of scale-resolved connected
amplitudes with \(r\) external local observables.  An amplitude retains all
photon, TT-graviton, and Newton shell labels and all tensor indices.  Let
\(\mathcal W_r\) be the complete gauge, diffeomorphism, and constraint Ward
projection acting on the sum of amplitudes at fixed external data.

\begin{definition}[Intrinsic massless-skeleton seminorm]
\label{def:v6v25-IMS-norm}
For \(A\in\mathfrak C_r\), define
\begin{equation}
 \|A\|_{\rm IMS}
 :=\sup_{\mathbf x,\mathbf j}
 \frac{\|\mathcal W_rA(\mathbf x,\mathbf j)\|}
 {r!K^r\mathcal T_r^{\rm hyb}(\mathbf x,\mathbf j)},
 \label{eq:v6v25-IMS-norm}
\end{equation}
where \(\mathcal T_r^{\rm hyb}\) is the sum of retained scale-labelled
hybrid root-tree weights of V23.  Proper-cumulant and gauge-related
microscopic terms are summed inside \(\mathcal W_rA\) before the numerator
norm is taken.
\end{definition}

The expression is a seminorm before quotienting Ward-null amplitudes.  It
becomes a norm on the corresponding quotient.  It deliberately does not
require an individual Feynman graph, forest tree, or constraint component
to satisfy the final bound.

\begin{lemma}[Ward-null changes do not affect the seminorm]
\label{lem:v6v25-Ward-quotient}
If \(A-A'\in\ker\mathcal W_r\), then
\(\|A\|_{\rm IMS}=\|A'\|_{\rm IMS}\).
\end{lemma}

\begin{proof}
The numerators agree because \(\mathcal W_rA=\mathcal W_rA'\).
\end{proof}

\subsection{The connected scale-score card}

\begin{definition}[Intrinsic massless scale-score card]
\label{def:v6v25-IMSSC}
The IMSSC holds at adjacent step \(n\) if there is
\(\eta_n^{\rm sk}\ge0\) such that, uniformly in \(0\le t\le1\),
\begin{equation}
 \left\|
 \kappa_{r+1}^{\nu_{n,t}}
 (F_1,\ldots,F_r,\mathcal S_{n,t}^{\rm mat})
 \right\|_{\rm IMS}
 \le\eta_n^{\rm sk}
 \prod_{i=1}^r\|F_i\|_{\rm end}
 \label{eq:v6v25-IMSSC}
\end{equation}
for all \(r\ge1\).  The complete score and the Ward projection are formed
before sectorwise absolute values.
\end{definition}

\begin{theorem}[IMSSC transports the interacting massless skeleton]
\label{thm:v6v25-IMSSC-transport}
Assume the finite-cutoff laws and observables are compatible, the endpoint
norms are uniform, and the IMSSC holds with
\begin{equation}
 \sum_n\eta_n^{\rm sk}<\infty.
 \label{eq:v6v25-score-summability}
\end{equation}
Then all connected cylinder cumulants are Cauchy in the IMS seminorm and
converge to a cutoff-limit functional with the same hybrid massless
skeleton bound.
\end{theorem}

\begin{proof}
Corollary~\ref{cor:v6v25-adjacent-comparison}, the triangle inequality
after the complete Ward projection, and
\eqref{eq:v6v25-IMSSC} give
\begin{equation}
 \|\kappa_r^{n+1}-\kappa_r^n\|_{\rm IMS}
 \le\eta_n^{\rm sk}\prod_i\|F_i\|_{\rm end}.
 \label{eq:v6v25-Cauchy-increment}
\end{equation}
The summability hypothesis makes the telescoping sequence Cauchy.  The
uniform hybrid tree denominator passes to the coefficientwise cylinder
limit.
\end{proof}

Unlike the V23 owner theorem, this transport bound is additive rather than
multiplicative because it compares the complete cumulants directly.  A
separate uniform base bound supplies their absolute size.

\begin{firewall}[Using the desired card at arity \(r+1\) is circular]
\label{fw:v6v25-score-circularity}
The IMSSC cannot be proved by treating \(\mathcal S_{n,t}^{\rm mat}\) as an
arbitrary extra observable and invoking the desired all-order IMS bound at
arity \(r+1\).  That assumes the conclusion.  A proof must exploit the
explicit centered scale score, exact Ward identities, interpolation gaps,
and scale matching.
\end{firewall}

\subsection{Sector decomposition without premature norms}

Formally decompose the matched score as
\begin{equation}
 \mathcal S^{\rm mat}
 =\mathcal S^\gamma+\mathcal S^{TT}
 +\mathcal S^{\rm N}+\mathcal S^{\rm mix}
 +\mathcal S^{\rm massive}+\mathcal S^{\rm ph},
 \label{eq:v6v25-score-sectors}
\end{equation}
for photon, TT graviton, Newton constraint, mixed massless, already-treated
massive, and determinant-phase sectors.  Equation
\eqref{eq:v6v25-score-sectors} is inserted into the complete cumulant before
the Ward projection.  Only after cancellations may one use a triangle
inequality between genuinely disjoint residual sectors.

\begin{proposition}[Separate analytic requirements]
\label{prop:v6v25-sector-requirements}
The leading sectors require different inputs:
\begin{enumerate}
\item the photon score is marginal by V24 and needs an exact running charge,
      gauge Ward identities, and an intrinsic massless connected estimate;
\item the TT and Newton scores need the physical constraint quotient,
      gravitational Ward identities, and either a non-Gaussian UV
      trajectory or another nonperturbative scale estimate;
\item the massive score may use the V15 repaired owner compiler;
\item the high determinant phase may use V18--V19 only after a
      locality-compatible determinant trivialization is constructed.
\end{enumerate}
No single massive covariance gap supplies all four requirements.
\end{proposition}

\begin{proof}
Items 1 and 2 are the canonical failures proved in V24 and the constraint
decomposition of V22.  Item 3 is the domain of the V15 theorem.  Item 4 is
the claim boundary of Theorems~\ref{thm:v6v18-all-scale-phase} and
\ref{thm:v6v19-global-phase-criterion}.
\end{proof}

\subsection{The first upstream analytic gate}

The score identity makes the next acquisition order explicit.  A useful
resolvent or covariance representation for the score insertion requires a
uniform spectral or coercive inequality along the interpolation path.

\begin{definition}[Skeleton interpolation gap]
\label{def:v6v25-skeleton-gap}
A positive skeleton path has a scale-uniform gap if its Markov or transfer
generator \(L_{n,t}\), on the gauge-constraint quotient and orthogonal to
constants, obeys
\begin{equation}
 \langle f,-L_{n,t}f\rangle_{\nu_{n,t}}
 \ge\lambda_n\|f\|_2^2
 \label{eq:v6v25-skeleton-gap}
\end{equation}
with a quantitatively controlled \(\lambda_n>0\), or an explicitly graded
massless substitute in which the zero modes are retained in the IMS
skeleton.
\end{definition}

A literal volume-uniform positive gap is impossible for physical massless
photons and gravitons.  Therefore the second alternative in the definition
is essential: invert only on a shell complement, retain the zero and low
modes, and prove shellwise bounds compatible with the IMS denominator.

\begin{theorem}[Conditional score-resolvent route]
\label{thm:v6v25-score-resolvent-route}
Suppose for every shell and interpolation parameter:
\begin{enumerate}
\item the generator inverse on the centered shell complement has norm at
      most the shell power assigned in \(\mathcal T_r^{\rm hyb}\);
\item the matched score has a complete Ward-projected derivative row with
      one additional summable scale clock; and
\item inserting that row into a connected covariance identity does not
      invoke the IMSSC recursively.
\end{enumerate}
Then the IMSSC follows with \(\eta_n^{\rm sk}\) equal to the summed score
clock.
\end{theorem}

\begin{proof}
Represent the centered score insertion by applying the shell inverse to
the score and contracting its gradient with the complete connected
observable gradient.  Item 1 pays the retained massless carrier, item 2
pays the adjacent step, and the Ward projection is applied before the
absolute norm.  The connected tree or covariance expansion then has one
marked score row and the same hybrid root denominator.  Item 3 ensures that
the argument supplies, rather than assumes, the bound
\eqref{eq:v6v25-IMSSC}.
\end{proof}

\begin{firewall}[The scale-score card remains unproved]
\label{fw:v6v25-IMSSC-open}
No interpolation path satisfying the shellwise gap, exact SM--Einstein Ward
identities, chiral phase patching, nonlinear constraint reduction, and
summable matched score has been constructed here.  The IMSSC is a precise
sufficient target, not a completed continuum theorem.
\end{firewall}

\begin{theorem}[Sixth-series V25 result]
\label{thm:v6v25-status}
The compulsory audit imports no analytic estimate but identifies an exact
algebraic improvement: an adjacent path derivative of a complete connected
cumulant is one cumulant with a centered scale score.  This avoids
leaf-summing the canonically non-summable photon and Einstein vertices.
The Ward-first IMS seminorm retains all physical massless scale labels and
takes norms only after proper-cumulant, gauge, and constraint cancellations.
A summable intrinsic massless scale-score card makes the cutoff cumulants
Cauchy in that seminorm.  The first upstream analytic gate is a shellwise
interpolation inverse and a matched score row that proves the card without
using it recursively.
\end{theorem}

\begin{proof}
Use Theorem~\ref{thm:v6v25-score-cumulant},
Corollary~\ref{cor:v6v25-adjacent-comparison},
Lemma~\ref{lem:v6v25-vacuum-cancellation},
Theorem~\ref{thm:v6v25-IMSSC-transport},
Proposition~\ref{prop:v6v25-sector-requirements}, and
Theorem~\ref{thm:v6v25-score-resolvent-route}, subject to the two
firewalls.
\end{proof}

V26 will construct the shellwise score identity first for the Abelian
photon sector with conserved current and a fixed finite volume.  It will
separate gauge zero modes, prove the exact shell Poincare estimate, and test
whether charge renormalization supplies a summable matched score before
coupling the result to the TT-gravity score.


\clearpage
\section{V26: the transverse Abelian shell score}
\label{sec:v6v26}

\subsection{Gauge-fixed finite-volume shell}

V25 asks first for a shellwise inverse and a noncircular score row.  We
construct both for the free transverse Abelian sector.  Let \(K_j\) be the
positive Maxwell quadratic operator restricted to nonzero transverse
Fourier modes in one shell of wavelength \(L_j\).  Assume
\begin{equation}
 cL_j^{-2}I\le K_j\le CL_j^{-2}I.
 \label{eq:v6v26-shell-spectrum}
\end{equation}
For inverse coupling \(\beta>0\), define
\begin{equation}
 d\gamma_{\beta,j}(A)
 =Z_{\beta,j}^{-1}
 \exp\left[-\frac\beta2\langle A,K_jA\rangle\right]dA,
 \qquad
 C_{\beta,j}=\beta^{-1}K_j^{-1}.
 \label{eq:v6v26-shell-Gaussian}
\end{equation}
The zero mode and harmonic torons are not included in this Gaussian; they
must be retained as separate compact or global variables.

\begin{theorem}[Exact shell Poincare inequality]
\label{thm:v6v26-shell-Poincare}
For every smooth \(F\),
\begin{equation}
 \operatorname{Var}_{\gamma_{\beta,j}}F
 \le \lambda_{\max}(C_{\beta,j})
 \mathbb E_{\gamma_{\beta,j}}\|\nabla F\|^2
 \le C\beta^{-1}L_j^2
 \mathbb E_{\gamma_{\beta,j}}\|\nabla F\|^2.
 \label{eq:v6v26-shell-Poincare}
\end{equation}
Equivalently, the Ornstein--Uhlenbeck generator with Dirichlet form
\(\mathbb E\|\nabla F\|^2\) has gap at least
\begin{equation}
 \lambda_{\beta,j}\ge c\beta L_j^{-2}.
 \label{eq:v6v26-shell-gap}
\end{equation}
\end{theorem}

\begin{proof}
Diagonalize the positive covariance.  The tensorized one-dimensional
Gaussian Poincare inequality gives
\(\operatorname{Var}F\le\mathbb E\langle\nabla F,C\nabla F\rangle\).
Bound this by \(\lambda_{\max}(C)\mathbb E\|\nabla F\|^2\).  From
\eqref{eq:v6v26-shell-spectrum},
\(\lambda_{\max}(C_{\beta,j})\le C\beta^{-1}L_j^2\).  The reciprocal is
the stated spectral-gap bound.
\end{proof}

The gap degenerates at large \(L_j\), exactly as required for a massless
field.  The factor \(L_j^2\) is assigned to the retained photon carrier in
the IMS norm; it is not bounded by a fictitious uniform gap.

\subsection{The centered quadratic score}

Let \(\beta_t>0\) be a differentiable path.  If \(N_j\) is the dimension
of the transverse shell, the normalized score is
\begin{equation}
 \mathcal S_{t,j}^{\gamma}
 =-\frac{\dot\beta_t}{2}
 \left(\langle A,K_jA\rangle-\frac{N_j}{\beta_t}\right).
 \label{eq:v6v26-photon-score}
\end{equation}
Its direct \(L^2\) norm grows like \(\sqrt{N_j}\) and is unsuitable for a
volume-uniform rooted estimate.  Gaussian integration by parts removes
that false multiplicity.

\begin{theorem}[Score insertion is one covariance-variation edge]
\label{thm:v6v26-score-edge}
For every smooth \(F\) with integrable derivatives,
\begin{align}
 \mathbb E_{\gamma_{\beta_t,j}}
 [F\mathcal S_{t,j}^{\gamma}]
 &=\frac12\operatorname{Tr}
 \left(\dot C_{t,j}\,
 \mathbb E_{\gamma_{\beta_t,j}}D^2F\right),
 \label{eq:v6v26-score-edge}\\
 \dot C_{t,j}
 &=-\frac{\dot\beta_t}{\beta_t}C_{t,j}.
 \label{eq:v6v26-covariance-variation}
\end{align}
Thus the complete centered score inserts two endpoint derivatives joined
by one retained photon shell, with relative clock
\begin{equation}
 \vartheta_t^\gamma
 :=|\partial_t\log\beta_t|.
 \label{eq:v6v26-photon-clock}
\end{equation}
No factor \(N_j\) remains.
\end{theorem}

\begin{proof}
For a centered Gaussian with covariance \(C\), the second-order Stein
identity gives
\begin{equation}
 \mathbb E\left[
 (\langle A,KA\rangle-\operatorname{Tr}(KC))F
 \right]
 =\operatorname{Tr}(CKC\,\mathbb ED^2F).
 \label{eq:v6v26-second-Stein}
\end{equation}
Here \(C=\beta_t^{-1}K_j^{-1}\), so
\(\operatorname{Tr}(K_jC)=N_j/\beta_t\) and
\(CK_jC=\beta_t^{-2}K_j^{-1}\).  Multiply
\eqref{eq:v6v26-second-Stein} by \(-\dot\beta_t/2\).  Since
\(\dot C=-\dot\beta_t\beta_t^{-2}K_j^{-1}\), this proves
\eqref{eq:v6v26-score-edge}.  Equation
\eqref{eq:v6v26-covariance-variation} is immediate.
\end{proof}

\begin{corollary}[Connected photon score identity]
\label{cor:v6v26-connected-score}
In the complete cumulant of V25, inserting
\(\mathcal S_{t,j}^{\gamma}\) is exactly equivalent to summing one marked
photon covariance-variation edge over all pairs of differentiated legs of
the complete connected observable product.  The vacuum trace term has
already canceled through score centering.
\end{corollary}

\begin{proof}
Apply Theorem~\ref{thm:v6v26-score-edge} to every moment in the partition
formula for the cumulant and use the proof of
Theorem~\ref{thm:v6v25-score-cumulant}.  Terms in which both derivatives
act inside a vacuum component cancel with the centering term; the remaining
terms are precisely the marked connected insertions.
\end{proof}

\subsection{Ward projection at the score edge}

The transverse shell covariance has symbol
\begin{equation}
 \widehat C_{\beta,j}^{\mu\nu}(k)
 =\beta^{-1}\frac{\chi_j(k)}{\widehat\Delta(k)}
 P_T^{\mu\nu}(k),
 \label{eq:v6v26-transverse-covariance}
\end{equation}
where the shell cutoff excludes \(k=0\).

\begin{lemma}[Conserved endpoints remove only the longitudinal part]
\label{lem:v6v26-Ward-projection}
If the complete endpoint rows obey
\(k_\mu J_\mu(k)=0\), then
\begin{equation}
 J_\mu\widehat C^{\mu\nu}_{\beta,j}J'_\nu
 =\beta^{-1}\frac{\chi_j}{\widehat\Delta}
 J_\mu J'_\mu.
 \label{eq:v6v26-Ward-projection}
\end{equation}
The Ward identity removes the projector denominator but leaves the physical
massless Green factor.
\end{lemma}

\begin{proof}
Contract
\(P_T^{\mu\nu}=\delta^{\mu\nu}-k^\mu k^\nu/|k|^2\) with either conserved
endpoint.  The longitudinal term vanishes.  The scalar Green denominator
is unaffected.
\end{proof}

This edge is therefore retained in the IMS skeleton, consistently with
V23--V24.

\subsection{Finite adjacent variation}

Choose the geometric path
\begin{equation}
 \log\beta_t=(1-t)\log\beta_n+t\log\beta_{n+1}.
 \label{eq:v6v26-geometric-beta-path}
\end{equation}
Then \(\vartheta_t^\gamma\) is constant and
\begin{equation}
 \int_0^1\vartheta_t^\gamma\,dt
 =|\log\beta_{n+1}-\log\beta_n|.
 \label{eq:v6v26-adjacent-clock}
\end{equation}

\begin{theorem}[Free-photon scale-score transport]
\label{thm:v6v26-free-photon-transport}
For the gauge-fixed nonzero transverse Gaussian shells, assume the endpoint
derivative rows satisfy the IMS tree bound uniformly in \(t\).  Then the
photon contribution to the V25 scale-score card is bounded by
\begin{equation}
 \eta_n^\gamma
 \le K|\log\beta_{n+1}-\log\beta_n|.
 \label{eq:v6v26-photon-eta}
\end{equation}
Consequently, the free transverse sector transports through every coupling
path of finite logarithmic total variation,
\begin{equation}
 \sum_n|\log\beta_{n+1}-\log\beta_n|<\infty.
 \label{eq:v6v26-log-variation}
\end{equation}
\end{theorem}

\begin{proof}
Theorems~\ref{thm:v6v25-score-cumulant} and
\ref{thm:v6v26-score-edge} replace the complete score by one marked
covariance edge.  Normalize this edge by the same retained photon weight in
the IMS denominator.  Equation \eqref{eq:v6v26-covariance-variation}
leaves only the relative clock \(\vartheta_t^\gamma\).  Integrating the
geometric path gives \eqref{eq:v6v26-photon-eta}.  Sum over adjacent steps
and apply Theorem~\ref{thm:v6v25-IMSSC-transport}.
\end{proof}

\begin{corollary}[A vanishing inverse coupling has infinite logarithmic
distance]
\label{cor:v6v26-zero-beta-distance}
If \(\beta_n\to0\), then every positive path satisfies
\begin{equation}
 \sum_n|\log\beta_{n+1}-\log\beta_n|=\infty.
 \label{eq:v6v26-zero-beta-infinite}
\end{equation}
The free Gaussian scale-score estimate therefore cannot cross such a
trajectory with finite total loss.
\end{corollary}

\begin{proof}
Finite total variation would make \(\log\beta_n\) a Cauchy sequence with a
finite limit.  But \(\beta_n\to0\) implies
\(\log\beta_n\to-\infty\).
\end{proof}

This is the score-metric version of the hypercharge rate obstruction in
V13.  It does not prove that the physical coupling follows such a
trajectory.

\subsection{What charged matter changes}

With charged matter, integrating one photon shell changes the matter
effective action, the wave-function factors, and the running charge.  The
complete matched score is no longer the quadratic expression
\eqref{eq:v6v26-photon-score} alone.

\begin{firewall}[The free shell identity is not interacting QED]
\label{fw:v6v26-not-QED}
Theorem~\ref{thm:v6v26-free-photon-transport} proves the Gaussian transverse
carrier calculation.  It does not construct the interacting Abelian
continuum, prove a charge beta function, control fermion vacuum
polarization, patch the chiral determinant, or establish reflection
positivity of the coupled regulator.
\end{firewall}

\begin{firewall}[Zero modes are separate data]
\label{fw:v6v26-zero-modes}
Gauge zero modes, harmonic torons, and global electric-flux sectors were
excluded before \eqref{eq:v6v26-shell-spectrum}.  They have no shell gap and
must be integrated as compact/global variables with their own Gauss-law
and transfer analysis.
\end{firewall}

\begin{theorem}[Sixth-series V26 result]
\label{thm:v6v26-status}
On a nonzero transverse photon shell of wavelength \(L_j\), the exact
Gaussian Poincare gap is order \(\beta L_j^{-2}\).  The centered quadratic
coupling score, whose direct norm is volume dependent, becomes by Gaussian
integration by parts one marked covariance-variation edge with no shell
multiplicity.  After Ward projection, the longitudinal projector cancels
but the physical massless Green edge remains in the IMS skeleton.  Along a
geometric coupling path the adjacent clock is
\(|\Delta\log\beta|\); finite logarithmic total variation suffices for the
free carrier, while \(\beta\to0\) lies at infinite score distance.
Interacting charge matching, zero modes, chiral patching, and positivity
remain open.
\end{theorem}

\begin{proof}
Use Theorem~\ref{thm:v6v26-shell-Poincare},
Theorem~\ref{thm:v6v26-score-edge},
Lemma~\ref{lem:v6v26-Ward-projection},
Theorem~\ref{thm:v6v26-free-photon-transport}, and
Corollary~\ref{cor:v6v26-zero-beta-distance}, subject to both firewalls.
\end{proof}

V27 will add the one-loop-form vacuum-polarization structure only as an
exact transverse two-point score decomposition, without assuming its sign
or asymptotic trajectory.  The aim is to identify the precise scalar
running kernel whose total variation must be controlled and to keep the
chiral determinant phase outside the positive Gaussian calculation.


\clearpage
\section{V27: transverse vacuum polarization and the running photon kernel}
\label{sec:v6v27}

\subsection{Gauge-invariant two-point structure}

V26 treats a free transverse shell with scalar inverse coupling \(\beta\).
After charged fields are integrated over a scale step, the quadratic photon
kernel acquires vacuum polarization.  We first derive its exact transverse
decomposition from gauge invariance, without assigning a perturbative sign
or asymptotic trajectory.

Let \(\Gamma_n[A]\) be a finite-cutoff gauge-invariant effective action on
a periodic lattice, differentiable at \(A=0\), and let
\begin{equation}
 \Gamma^{(2)}_{n,\mu\nu}(k)
 :=\frac{\delta^2\Gamma_n}
 {\delta\widehat A_\mu(-k)\delta\widehat A_\nu(k)}\bigg|_{A=0}.
 \label{eq:v6v27-two-point-kernel}
\end{equation}
Write \(\widehat k_\mu\) for the exact lattice-difference symbol.

\begin{lemma}[Exact lattice Ward transversality]
\label{lem:v6v27-Ward-transversality}
If \(\Gamma_n[A+\nabla\lambda]=\Gamma_n[A]\), then for every nonzero
momentum
\begin{equation}
 \widehat k_\mu\Gamma^{(2)}_{n,\mu\nu}(k)=0,
 \qquad
 \Gamma^{(2)}_{n,\mu\nu}(k)\widehat k_\nu=0.
 \label{eq:v6v27-Ward-transversality}
\end{equation}
\end{lemma}

\begin{proof}
Differentiate gauge invariance once with respect to \(\lambda\) and once
with respect to \(A\), then evaluate at \(A=0\).  Fourier transformation
turns the lattice gradient into multiplication by \(i\widehat k\).  The
second identity follows from symmetry of the Hessian.
\end{proof}

Let
\begin{equation}
 P^T_{\mu\nu}(k)
 :=\delta_{\mu\nu}-
 \frac{\widehat k_\mu\widehat k_\nu}{\widehat k^2}.
 \label{eq:v6v27-photon-projector}
\end{equation}
On its \((D-1)\)-dimensional transverse range define
\begin{align}
 b_n(k)&:=\frac1{(D-1)\widehat k^2}
 \operatorname{Tr}_T\Gamma_n^{(2)}(k),
 \label{eq:v6v27-scalar-kernel}\\
 R_n(k)&:=\Gamma_n^{(2)}(k)
 -\widehat k^2b_n(k)P^T(k).
 \label{eq:v6v27-anisotropic-remainder}
\end{align}

\begin{proposition}[Exact scalar--anisotropic decomposition]
\label{prop:v6v27-decomposition}
For every \(k\ne0\),
\begin{equation}
 \Gamma_n^{(2)}(k)
 =\widehat k^2b_n(k)P^T(k)+R_n(k),
 \label{eq:v6v27-decomposition}
\end{equation}
where
\begin{equation}
 P^TR_nP^T=R_n,
 \qquad
 \operatorname{Tr}_TR_n=0.
 \label{eq:v6v27-remainder-properties}
\end{equation}
The decomposition is unique.
\end{proposition}

\begin{proof}
Ward transversality makes the Hessian an operator on the transverse range.
Split that operator into its normalized trace times the identity on the
range and its transverse traceless part.  This gives the displayed formula
and uniqueness.
\end{proof}

Rotational symmetry in a continuum limit forces \(R_n=0\).  At finite
lattice spacing hypercubic artifacts may contribute \(R_n\), so it must be
bounded rather than discarded.

\subsection{Exact Gaussian score for a varying quadratic kernel}

Let \(G_t\) be any positive symmetric matrix on one nonzero transverse
shell and let \(C_t=G_t^{-1}\).  Define
\begin{equation}
 d\gamma_t(A)=Z_t^{-1}
 e^{-\langle A,G_tA\rangle/2}\,dA.
 \label{eq:v6v27-general-Gaussian}
\end{equation}

\begin{theorem}[Quadratic-kernel score identity]
\label{thm:v6v27-kernel-score}
The centered score of \eqref{eq:v6v27-general-Gaussian} is
\begin{equation}
 \mathcal S_t^G
 =-\frac12\left(
 \langle A,\dot G_tA\rangle
 -\operatorname{Tr}(C_t\dot G_t)\right).
 \label{eq:v6v27-kernel-score}
\end{equation}
For every smooth integrable \(F\),
\begin{align}
 \mathbb E_t(F\mathcal S_t^G)
 &=\frac12\operatorname{Tr}
 \left(\dot C_t\,\mathbb E_tD^2F\right),
 \label{eq:v6v27-kernel-score-edge}\\
 \dot C_t&=-C_t\dot G_tC_t.
 \label{eq:v6v27-inverse-derivative}
\end{align}
\end{theorem}

\begin{proof}
Differentiate the normalized Gaussian density to obtain
\eqref{eq:v6v27-kernel-score}.  The second-order Gaussian integration by
parts identity is
\[
 \mathbb E[(\langle A,HA\rangle-\operatorname{Tr}(HC))F]
 =\operatorname{Tr}(CHC\,\mathbb ED^2F).
\]
Set \(H=\dot G_t\), multiply by \(-1/2\), and differentiate
\(G_tC_t=I\) to get \eqref{eq:v6v27-inverse-derivative}.
\end{proof}

Thus the correct clock is the relative covariance variation, not the direct
norm of the extensive quadratic score.

\begin{definition}[Relative photon-kernel clock]
\label{def:v6v27-relative-clock}
On a shell define
\begin{equation}
 \vartheta_t^{\rm pol}
 :=\left\|C_t^{1/2}\dot G_tC_t^{1/2}\right\|_{
m op}
 =\left\|C_t^{-1/2}\dot C_tC_t^{-1/2}\right\|_{
m op}.
 \label{eq:v6v27-relative-clock}
\end{equation}
\end{definition}

\begin{lemma}[Scalar kernel clock]
\label{lem:v6v27-scalar-clock}
If \(R_t=0\) and
\(G_t(k)=\widehat k^2b_t(k)P^T(k)\) with \(b_t(k)>0\), then
\begin{equation}
 \vartheta_t^{\rm pol}(k)
 =|\partial_t\log b_t(k)|.
 \label{eq:v6v27-scalar-clock}
\end{equation}
\end{lemma}

\begin{proof}
On the transverse range,
\(C_t=(\widehat k^2b_t)^{-1}I\) and
\(\dot G_t=\widehat k^2\dot b_tI\).  Their normalized product is
\((\dot b_t/b_t)I\).
\end{proof}

For nonzero \(R_t\), Definition~\ref{def:v6v27-relative-clock} remains the
exact matrix clock and automatically measures scalar and hypercubic
variations in the correct covariance metric.

\subsection{Renormalized scalar matching}

Choose one nonzero reference momentum \(k_*\) in the matching window and
write
\begin{equation}
 \beta_n^{\rm ren}:=b_n(k_*),
 \qquad
 \widetilde b_n(k):=b_n(k)-b_n(k_*).
 \label{eq:v6v27-renormalized-beta}
\end{equation}
The adjacent matching condition assigns the constant transverse quadratic
coefficient to \(\beta_n^{\rm ren}\); only
\(\widetilde b_n\) belongs to the irrelevant residual score.

\begin{proposition}[Exact split of the scalar score]
\label{prop:v6v27-score-split}
For every differentiable interpolation,
\begin{equation}
 \partial_tb_t(k)
 =\partial_t\beta_t^{\rm ren}
 +\partial_t\widetilde b_t(k),
 \qquad
 \widetilde b_t(k_*)=0.
 \label{eq:v6v27-score-split}
\end{equation}
The first term is the marginal charge-running coordinate.  The second is a
momentum-dependent residual and cannot be absorbed into the charge without
changing the matching prescription.
\end{proposition}

\begin{proof}
This is the definition \eqref{eq:v6v27-renormalized-beta} differentiated
in \(t\).  Evaluation at \(k_*\) proves the last identity.
\end{proof}

The total-variation clock for the marginal part is
\begin{equation}
 \sum_n
 |\log\beta_{n+1}^{\rm ren}-\log\beta_n^{\rm ren}|,
 \label{eq:v6v27-marginal-clock}
\end{equation}
provided the full transverse kernel stays positive.

\subsection{Gapped charged modes give a local residual}

Assume an integrated charged block has an exponentially decaying kernel
uniformly at the step under consideration.  Its quadratic polarization is
then analytic in a fixed complex momentum strip.  Suppose also that the
lattice symmetries include inversion and the hypercubic group.

\begin{theorem}[Analytic polarization matching]
\label{thm:v6v27-analytic-polarization}
For a gauge-invariant exponentially local gapped charged block, its
quadratic polarization has an expansion near \(k=0\) of the form
\begin{equation}
 \Pi_{\mu\nu}(k)
 =\widehat k^2\pi(0)P^T_{\mu\nu}(k)
 +O(|\widehat k|^4)
 \label{eq:v6v27-polarization-expansion}
\end{equation}
as an analytic tensor after the transverse expression is assembled.
After \(\pi(0)\) is absorbed into \(\beta^{\rm ren}\), the relative
quadratic score on a shell of wavelength \(L_j\) is
\begin{equation}
 O(L_j^{-2}).
 \label{eq:v6v27-residual-clock}
\end{equation}
Hence its infrared shell sum is absolutely summable.
\end{theorem}

\begin{proof}
Exponential locality gives an analytic Fourier symbol.  Gauge invariance
gives the two Ward identities
\eqref{eq:v6v27-Ward-transversality}, so a photon mass term is absent.
Inversion removes odd Taylor terms.  Hypercubic symmetry makes the unique
isotropic transverse term at quadratic order a scalar multiple of
\(\widehat k^2P^T\); every remaining analytic transverse tensor starts at
order four.  This proves \eqref{eq:v6v27-polarization-expansion}.

The base inverse covariance is order \(|\widehat k|^2\).  Dividing the
matched order-four residual by it in the relative clock gives order
\(|\widehat k|^2\), which is \(O(L_j^{-2})\) on the shell.  The dyadic sum
\(\sum_jL_j^{-2}\) converges.
\end{proof}

\begin{firewall}[Massless charged fields are not a local polarization
owner]
\label{fw:v6v27-massless-polarization}
Theorem~\ref{thm:v6v27-analytic-polarization} requires a gapped integrated
block.  Retained massless or light charged fields can produce nonanalytic
momentum dependence and logarithmic running.  They must stay in the
interacting massless skeleton; their contribution belongs to the marginal
kernel \(b_n(k)\), not to the analytic local residual.
\end{firewall}

\subsection{A conditional polarization score card}

Let \(G_{n,t}(k)\) be the full positive transverse quadratic kernel on a
shell, including \(R_{n,t}\), and put
\begin{equation}
 \eta_n^{\rm pol}
 :=K\int_0^1
 \sup_{j,k\in\mathcal A_j}
 \left\|C_{n,t}(k)^{1/2}
 \partial_tG_{n,t}(k)
 C_{n,t}(k)^{1/2}\right\|dt.
 \label{eq:v6v27-polarization-eta}
\end{equation}

\begin{theorem}[Quadratic polarization contribution to IMSSC]
\label{thm:v6v27-polarization-IMSSC}
Assume:
\begin{enumerate}
\item \(G_{n,t}\) is positive on every nonzero transverse shell;
\item the complete charged-field effective action is gauge invariant;
\item the quadratic score is inserted by
      Theorem~\ref{thm:v6v27-kernel-score} before graphwise norms; and
\item the nonquadratic connected score rows obey the IMS endpoint bounds.
\end{enumerate}
Then the quadratic polarization part of the V25 IMSSC is bounded by
\(\eta_n^{\rm pol}\).  It is summable if the marginal logarithmic clock
\eqref{eq:v6v27-marginal-clock}, the relative anisotropic variation, and
the analytic residual clocks are all summable.
\end{theorem}

\begin{proof}
The quadratic score becomes one marked covariance-variation edge by
Theorem~\ref{thm:v6v27-kernel-score}.  Normalize that edge by the retained
photon carrier in the IMS denominator.  The remaining operator is exactly
the relative clock in \eqref{eq:v6v27-polarization-eta}.  Gauge invariance
applies the Ward projection before the norm.  The scalar marginal part is
Lemma~\ref{lem:v6v27-scalar-clock}; the gapped analytic residual is
summable by Theorem~\ref{thm:v6v27-analytic-polarization}; the traceless
remainder is measured directly by the same relative matrix norm.
\end{proof}

\begin{firewall}[No sign or beta function has been proved]
\label{fw:v6v27-no-beta-sign}
Nothing above proves the sign of the charged-field contribution to
\(b_n\), the existence of a nonperturbative hypercharge trajectory, or
finite total variation of \eqref{eq:v6v27-marginal-clock}.  Positivity of
the full transverse kernel and summability are explicit hypotheses.
\end{firewall}

\begin{firewall}[The chiral phase is outside the positive Hessian]
\label{fw:v6v27-phase-separate}
The quadratic kernel in this note belongs to the positive modulus backbone.
The chiral determinant phase must satisfy the V18 complex polymer and V19
determinant-line conditions.  It cannot be inserted into
\eqref{eq:v6v27-general-Gaussian} as though it were a positive quadratic
form.
\end{firewall}

\begin{theorem}[Sixth-series V27 result]
\label{thm:v6v27-status}
Gauge invariance makes the exact photon Hessian transverse.  It decomposes
uniquely into a scalar running kernel \(\widehat k^2b_n(k)P^T\) and a
transverse traceless lattice remainder.  For any positive interpolating
kernel, the extensive centered quadratic score is exactly one marked edge
\(\dot C=-C\dot GC\); its intrinsic clock is the relative operator
\(\|C^{1/2}\dot GC^{1/2}\|\).  Charge matching separates the marginal
variation of \(\beta^{\rm ren}\) from momentum-dependent residuals.  A
gapped local charged block has an order-four matched residual and therefore
an \(O(L^{-2})\) summable shell clock.  Retained light charged fields and
the marginal logarithmic clock remain inside the massless skeleton, with
no sign or ultraviolet trajectory claimed.
\end{theorem}

\begin{proof}
Use Lemma~\ref{lem:v6v27-Ward-transversality},
Proposition~\ref{prop:v6v27-decomposition},
Theorem~\ref{thm:v6v27-kernel-score},
Proposition~\ref{prop:v6v27-score-split},
Theorem~\ref{thm:v6v27-analytic-polarization}, and
Theorem~\ref{thm:v6v27-polarization-IMSSC}, subject to the three firewalls.
\end{proof}

V28 will formulate the retained-light-field contribution as a shellwise
current--current connected score.  It will prove the exact Ward cancellation
of the longitudinal photon sector and identify the additional bound needed
to control the scalar marginal clock without assuming a perturbative beta
function.


\clearpage
\section{V28: retained charged currents and the marginal susceptibility}
\label{sec:v6v28}

\subsection{Exact matter Hessian}

V27 treats a gapped charged block by an analytic polarization remainder.
Retained light charged fields are different: their connected current
correlations can be long range and must remain inside the massless
skeleton.  The exact quadratic photon insertion contains both a connected
current covariance and a contact, or seagull, term.

Let \(\psi\) denote finite-cutoff matter variables with positive modulus
measure at fixed background photon field \(A\), and write
\begin{equation}
 Z_m[A]=\int e^{-S_m[A,\psi]},d\lambda_m(\psi),
 \qquad
 W_m[A]:=-\log Z_m[A].
 \label{eq:v6v28-matter-effective-action}
\end{equation}
Define
\begin{align}
 J_\mu(x)&:=\frac{\delta S_m}{\delta A_\mu(x)},
 \label{eq:v6v28-current}\\
 K_{\mu\nu}(x,y)&:=
 \frac{\delta^2S_m}
 {\delta A_\mu(x)\delta A_\nu(y)}.
 \label{eq:v6v28-seagull}
\end{align}

\begin{theorem}[Exact seagull--current formula]
\label{thm:v6v28-seagull-current}
Whenever differentiation under the finite-cutoff integral is justified,
\begin{align}
 \frac{\delta W_m}{\delta A_\mu(x)}
 &=\langle J_\mu(x)\rangle_A,
 \label{eq:v6v28-first-derivative}\\
 \Pi_{\mu\nu}(x,y)
 :=\frac{\delta^2 W_m}
 {\delta A_\mu(x)\delta A_\nu(y)}
 &=\langle K_{\mu\nu}(x,y)\rangle_A
 -\langle J_\mu(x);J_\nu(y)\rangle_{A,c}.
 \label{eq:v6v28-complete-polarization}
\end{align}
\end{theorem}

\begin{proof}
Differentiate \(-\log Z_m\).  The first derivative is the normalized
expectation of \(\delta S_m/\delta A\).  Differentiating that expectation
produces the expectation of \(K\) and minus the covariance with the
derivative of the action.  This gives
\eqref{eq:v6v28-complete-polarization}.
\end{proof}

Neither term on the right is the polarization by itself.  In particular,
the contact term is required by the lattice Ward identity.

\subsection{Ward cancellation before norms}

Assume the regularized matter partition is gauge invariant,
\begin{equation}
 W_m[A+\nabla\lambda]=W_m[A].
 \label{eq:v6v28-matter-gauge-invariance}
\end{equation}

\begin{theorem}[Complete current Ward identity]
\label{thm:v6v28-current-Ward}
At a gauge-invariant background,
\begin{align}
 \nabla_\mu^*\langle J_\mu(x)\rangle_A&=0,
 \label{eq:v6v28-one-current-Ward}\\
 \nabla_\mu^*\Pi_{\mu\nu}(x,y)&=0.
 \label{eq:v6v28-two-current-Ward}
\end{align}
Equivalently in momentum space,
\begin{equation}
 \widehat k_\mu
 \left[widehat{\langle K_{\mu\nu}\rangle}
 -\widehat{\langle J_\mu;J_\nu\rangle_c}\right](k)=0.
 \label{eq:v6v28-Ward-cancellation}
\end{equation}
\end{theorem}

\begin{proof}
Differentiate \eqref{eq:v6v28-matter-gauge-invariance} with respect to the
gauge parameter to obtain \eqref{eq:v6v28-one-current-Ward}.  Differentiate
that identity with respect to \(A_\nu(y)\).  The derivative of the expected
current is exactly \eqref{eq:v6v28-complete-polarization}, proving
\eqref{eq:v6v28-two-current-Ward}.  Fourier transformation gives the last
formula.
\end{proof}

\begin{corollary}[Longitudinal photon cancellation]
\label{cor:v6v28-longitudinal-cancellation}
When the complete tensor \(\Pi\) is inserted between photon propagators,
all longitudinal projector contributions vanish.  This conclusion need
not hold for \(\langle K\rangle\) and
\(\langle J;J\rangle_c\) separately.
\end{corollary}

\begin{proof}
Contract either side with \(\widehat k\) and use
\eqref{eq:v6v28-Ward-cancellation}.  Separate terms are not transverse
unless an additional identity is available.
\end{proof}

\begin{firewall}[The Ward identity belongs to the complete regulated
family]
\label{fw:v6v28-complete-regulator}
If a cutoff, chiral phase convention, or mode split is not gauge invariant,
its missing counterterm and cross-shell pieces must be included before
Theorem~\ref{thm:v6v28-current-Ward} is used.  Imposing transversality on
one selected diagram after taking norms is not an exact Ward proof.
\end{firewall}

\subsection{The scalar marginal susceptibility}

At \(A=0\), apply the transverse scalar projection from V27 and define
\begin{equation}
 \chi_n(k):=
 \frac{1}{(D-1)\widehat k^2}
 \operatorname{Tr}_T\widehat\Pi_n(k),
 \qquad k\ne0.
 \label{eq:v6v28-susceptibility}
\end{equation}
This definition is made after the seagull and connected covariance are
combined.

Let \(\Delta_n\Pi\) be the contribution of one exact adjacent retained-field
step and let \(k_{*,n}\) be its charge-matching momentum.  Define
\begin{equation}
 \delta\beta_n^{\rm light}
 :=\frac{1}{(D-1)\widehat k_{*,n}^2}
 \operatorname{Tr}_T\widehat{\Delta_n\Pi}(k_{*,n}).
 \label{eq:v6v28-beta-increment}
\end{equation}

\begin{proposition}[Exact charge-coordinate increment]
\label{prop:v6v28-charge-increment}
If the photon inverse kernel before the step has scalar coefficient
\(b_n(k)\), then matching the transverse quadratic form at \(k_{*,n}\)
gives
\begin{equation}
 \beta_{n+1}^{\rm ren}
 =\beta_n^{\rm ren}+\delta\beta_n^{\rm light}
 +\delta\beta_n^{\rm high}+\delta\beta_n^{\rm ct},
 \label{eq:v6v28-charge-recursion}
\end{equation}
where the last two terms are respectively the already-integrated gapped
block and the chosen counterterm contribution.  Equation
\eqref{eq:v6v28-charge-recursion} is an identity once the three pieces and
matching prescription are fixed; it asserts no sign.
\end{proposition}

\begin{proof}
Evaluate the scalar transverse decomposition of the complete adjacent
quadratic kernel at \(k_{*,n}\).  The light, high, and counterterm Hessians
add linearly.  By definition their scalar projections are the three
increments in \eqref{eq:v6v28-charge-recursion}.
\end{proof}

\subsection{A nonperturbative susceptibility card}

Define the Ward-first marginal magnitude
\begin{equation}
 \mathfrak q_n
 :=\sup_{t\in[0,1]}
 \frac{1}{(D-1)\widehat k_{*,n}^2}
 \left|
 \operatorname{Tr}_T
 \mathcal W_2
 \left(
 \widehat{\langle K\rangle}_{n,t}
 -\widehat{\langle J;J\rangle}_{n,t,c}
 \right)(k_{*,n})
 \right|,
 \label{eq:v6v28-q-card}
\end{equation}
where \(\mathcal W_2\) means that all regulator, counterterm, and
gauge-related contributions needed by the two-point Ward identity are
assembled first.

\begin{lemma}[Susceptibility bounds the charge increment]
\label{lem:v6v28-q-bounds-beta}
For an adjacent interpolation of length one,
\begin{equation}
 |\delta\beta_n^{\rm light}|\le\mathfrak q_n
 \label{eq:v6v28-beta-bound}
\end{equation}
when \(\mathfrak q_n\) is applied to the interpolated step derivative.
\end{lemma}

\begin{proof}
The adjacent Hessian difference is the integral of its path derivative.
Project it by the linear scalar functional in
\eqref{eq:v6v28-beta-increment}, take the absolute value only after
\(\mathcal W_2\), and use the uniform bound
\eqref{eq:v6v28-q-card} over a unit interval.
\end{proof}

\begin{theorem}[Relative susceptibility summability criterion]
\label{thm:v6v28-relative-susceptibility}
Assume the complete renormalized inverse couplings remain positive and
\begin{equation}
 |\delta\beta_n^{\rm tot}|le\frac12\beta_n^{\rm ren},
 \qquad
 \sum_n\frac{|\delta\beta_n^{\rm tot}|}
 {\beta_n^{\rm ren}}<\infty,
 \label{eq:v6v28-relative-sum}
\end{equation}
where \(\delta\beta_n^{\rm tot}\) is the full right-side increment in
\eqref{eq:v6v28-charge-recursion}.  Then
\begin{equation}
 \sum_n
 |\log\beta_{n+1}^{\rm ren}-\log\beta_n^{\rm ren}|<\infty,
 \label{eq:v6v28-log-sum}
\end{equation}
and the marginal scalar part of the V27 polarization clock is summable.
\end{theorem}

\begin{proof}
For \(|x|\le1/2\), \(|\log(1+x)|\le2|x|\).  Apply this with
\(x=\delta\beta_n^{\rm tot}/\beta_n^{\rm ren}\), use
\eqref{eq:v6v28-charge-recursion}, and sum.  The scalar-clock identity is
Lemma~\ref{lem:v6v27-scalar-clock}.
\end{proof}

\begin{corollary}[A sufficient current card]
\label{cor:v6v28-current-card}
If the high and counterterm increments have finite relative variation and
\begin{equation}
 \sum_n\frac{\mathfrak q_n}{\beta_n^{\rm ren}}<\infty,
 \label{eq:v6v28-current-sum}
\end{equation}
then the light-current marginal clock is summable, provided the small-step
condition in \eqref{eq:v6v28-relative-sum} holds after finitely many initial
steps.
\end{corollary}

\begin{proof}
Use Lemma~\ref{lem:v6v28-q-bounds-beta}, add the two other finite-variation
sequences, and apply Theorem~\ref{thm:v6v28-relative-susceptibility} to the
tail.  Finitely many initial positive steps cost a finite factor.
\end{proof}

This is the precise nonperturbative scalar bound needed from the retained
charged sector.  It is stronger than boundedness of one current two-point
function: it asks for summable \emph{relative adjacent variation} after
matching.

\subsection{Why positivity does not fix the sign}

The covariance tensor
\(\langle J;J\rangle_c\) is positive semidefinite as a quadratic form in
real test sources.  In \eqref{eq:v6v28-complete-polarization} it is
subtracted from the expected seagull.

\begin{proposition}[No sign from covariance positivity alone]
\label{prop:v6v28-no-sign}
Without an additional convexity or spectral theorem for the complete
matter effective action, positivity of
\(\langle J;J\rangle_c\) does not determine the sign of \(\Pi\) or of
\(\delta\beta_n^{\rm light}\).
\end{proposition}

\begin{proof}
Equation \eqref{eq:v6v28-complete-polarization} is the difference of the
seagull expectation and a positive covariance.  Either term may dominate
unless a relation between them is proved.  The transverse scalar projection
is linear and therefore does not restore a sign.
\end{proof}

\subsection{Massless nonanalyticity remains in the skeleton}

For a retained light field, \(\chi_n(k)\) need not be analytic at
\(k=0\).  The Ward identity supplies the factor \(\widehat k^2P^T\) but
does not make the scalar coefficient analytic or bounded uniformly over all
cutoff steps.

\begin{firewall}[No local Taylor remainder for the retained current]
\label{fw:v6v28-no-light-Taylor}
The order-four analytic residual theorem of V27 applies only to a gapped
integrated charged block.  It cannot be applied to
\eqref{eq:v6v28-susceptibility} for retained light fields.  Logarithmic or
other nonanalytic scale dependence belongs to the IMS skeleton and the
charge recursion.
\end{firewall}

\begin{firewall}[The chiral measure is still split]
\label{fw:v6v28-chiral-split}
The positive expectations in this note refer to the modulus backbone or a
vectorlike retained representation.  The full chiral Standard Model
functional additionally requires the V18 complex phase expansion and the
V19 locality-compatible determinant-line trivialization.  Ward invariance
of the modulus alone does not prove the full chiral gauge identity.
\end{firewall}

\begin{theorem}[Sixth-series V28 result]
\label{thm:v6v28-status}
The exact retained-matter photon Hessian is the expected seagull minus the
connected current covariance.  Gauge invariance makes this complete
combination transverse; its two terms may not be divided by
\(\widehat k^2\) or normed separately.  The transverse scalar projection at
the matching momentum gives the exact light-field increment of the
renormalized inverse coupling.  A Ward-first susceptibility card with
\(\sum_n\mathfrak q_n/\beta_n^{\rm ren}<\infty\), together with finite
relative high and counterterm variation, is sufficient for the marginal
photon score clock.  No sign, beta function, massless analyticity, or chiral
phase construction has been assumed.
\end{theorem}

\begin{proof}
Use Theorems~\ref{thm:v6v28-seagull-current} and
\ref{thm:v6v28-current-Ward},
Proposition~\ref{prop:v6v28-charge-increment},
Lemma~\ref{lem:v6v28-q-bounds-beta},
Theorem~\ref{thm:v6v28-relative-susceptibility}, and
Proposition~\ref{prop:v6v28-no-sign}, subject to the three firewalls.
\end{proof}

V29 will go upstream at the susceptibility card.  It will express the
adjacent current covariance through a shell-resolvent identity, determine
the exact gap or decay estimate required to bound \(\mathfrak q_n\), and
test whether the proposed bound is compatible with a genuinely massless
charged sector rather than smuggling in a fermion mass gap.


\clearpage
\section{V29: canonical photon dilation and the additive charge clock}
\label{sec:v6v29}

\subsection{Why the fixed-field logarithmic clock is not intrinsic}

V26--V28 use a fixed photon coordinate \(A\), for which varying the inverse
coupling \(\beta\) changes the Gaussian covariance and produces the clock
\(|\Delta\log\beta|\).  That is a valid sufficient estimate in that
coordinate, but it is too strong for a trajectory on which the canonically
normalized charge tends to zero.  An exact field dilation removes the free
quadratic variation and transfers it to the interaction vertex.

On one finite transverse shell, let
\begin{equation}
 d\gamma_\beta(A)=Z_\beta^{-1}
 \exp\left[-\frac\beta2\langle A,KA\rangle\right]dA,
 \qquad K>0.
 \label{eq:v6v29-beta-Gaussian}
\end{equation}
Define the canonical photon coordinate and charge by
\begin{equation}
 \mathcal A:=\sqrt\beta\,A,
 \qquad
 e:=\beta^{-1/2}.
 \label{eq:v6v29-canonical-dilation}
\end{equation}

\begin{theorem}[Exact Gaussian dilation]
\label{thm:v6v29-Gaussian-dilation}
For every integrable observable \(F\),
\begin{equation}
 \int F(A)\,d\gamma_\beta(A)
 =\int F(e\mathcal A)\,d\gamma_1(\mathcal A),
 \label{eq:v6v29-Gaussian-pushforward}
\end{equation}
where \(d\gamma_1\) has covariance \(K^{-1}\).  All determinant powers of
\(\beta\) cancel in the normalized pushforward.
\end{theorem}

\begin{proof}
Make the linear change of variables
\(A=\beta^{-1/2}\mathcal A\).  Its Jacobian is
\(\beta^{-N/2}\), while
\(Z_\beta=\beta^{-N/2}Z_1\) in the \(N\)-dimensional shell.  The normalized
Jacobian therefore cancels and the quadratic form becomes
\(\langle\mathcal A,K\mathcal A\rangle/2\).
\end{proof}

\subsection{Equivalence of score and moving-observable forms}

Let \(\beta_t>0\), \(e_t=\beta_t^{-1/2}\), and put
\(G_t(\mathcal A)=F(e_t\mathcal A)\).

\begin{lemma}[Dilation derivative]
\label{lem:v6v29-dilation-derivative}
\begin{equation}
 \partial_tG_t(\mathcal A)
 =-\frac12(\partial_t\log\beta_t)
 \mathcal A\cdot\nabla_{\mathcal A}G_t(\mathcal A)
 =\dot e_t\,\mathcal A\cdot(\nabla_AF)(e_t\mathcal A).
 \label{eq:v6v29-dilation-derivative}
\end{equation}
\end{lemma}

\begin{proof}
Differentiate \(e_t=\beta_t^{-1/2}\) and apply the chain rule.
\end{proof}

\begin{theorem}[Free score cancellation by canonical dilation]
\label{thm:v6v29-free-score-cancellation}
The fixed-field quadratic score identity of V26 and the moving-observable
identity obtained from \eqref{eq:v6v29-Gaussian-pushforward} agree exactly:
\begin{equation}
 \mathbb E_{\gamma_{\beta_t}}
 [F\mathcal S_t^\gamma]
 =\mathbb E_{\gamma_1}[\partial_tG_t].
 \label{eq:v6v29-score-dilation-equivalence}
\end{equation}
If observables and interactions are expressed directly in the canonical
coordinate \(\mathcal A\), the free Gaussian law has no coupling score.
\end{theorem}

\begin{proof}
The standard Gaussian Stein identity gives
\begin{equation}
 \mathbb E_{\gamma_1}
 [\mathcal A\cdot\nabla_{\mathcal A}G]
 =\operatorname{Tr}
 (K^{-1}\mathbb E_{\gamma_1}D^2G).
 \label{eq:v6v29-Euler-Stein}
\end{equation}
Insert Lemma~\ref{lem:v6v29-dilation-derivative}.  Under the change of
variables, \(D_{\mathcal A}^2G=e_t^2D_A^2F\) and
\(C_{\beta_t}=e_t^2K^{-1}\).  The result is exactly the covariance-variation
formula \eqref{eq:v6v26-score-edge}, proving
\eqref{eq:v6v29-score-dilation-equivalence}.  If all quantities are defined
in \(\mathcal A\), the reference density is \(\gamma_1\) and its derivative
vanishes.
\end{proof}

\begin{assessment}[Scope correction to the V26 clock]
\label{ass:v6v29-V26-scope}
The V26 logarithmic clock is correct for transport at fixed unnormalized
field \(A\).  It is not a coordinate-independent necessity.  A canonical
dilation replaces it by moving interaction and observable terms.  The
statement that \(\beta\to0\) is at infinite distance survives because then
\(e\to\infty\), but \(\beta\to\infty\) can have finite additive
\(e\)-variation even though \(|\log\beta|\) diverges.
\end{assessment}

\subsection{Where the running charge appears}

In the convention with unit charge multiplying the unnormalized connection,
the matter covariant derivative has the form
\begin{equation}
 D(A)=\partial+iY A.
 \label{eq:v6v29-unnormalized-covariant-derivative}
\end{equation}
After \eqref{eq:v6v29-canonical-dilation},
\begin{equation}
 D(\mathcal A)=\partial+ieY\mathcal A.
 \label{eq:v6v29-canonical-covariant-derivative}
\end{equation}
Thus the charge becomes an interaction coordinate.

Let a positive canonical-coordinate action be
\begin{equation}
 S_t(\mathcal A,\psi)
 =\frac12\langle\mathcal A,K\mathcal A\rangle
 +S_m(e_t\mathcal A,\psi)
 +S_{{\rm ct},t}(\mathcal A,\psi).
 \label{eq:v6v29-canonical-action}
\end{equation}

\begin{proposition}[Canonical matched score]
\label{prop:v6v29-canonical-score}
The centered action score is
\begin{equation}
 \mathcal S_t^{\rm can}
 =-\dot e_t
 \left(\partial_eS_m-
 \mathbb E_t\partial_eS_m\right)
 -\left(\partial_tS_{{\rm ct},t}
 -\mathbb E_t\partial_tS_{{\rm ct},t}\right).
 \label{eq:v6v29-canonical-score}
\end{equation}
There is no extensive free Maxwell quadratic score.
\end{proposition}

\begin{proof}
Differentiate \eqref{eq:v6v29-canonical-action}; the free quadratic form is
independent of \(t\).  Apply the centered-action formula
\eqref{eq:v6v25-centered-action-score}.
\end{proof}

At small \(e\), \(\partial_eS_m\) begins with the complete
photon--current vertex and its gauge-completion terms.  The Ward-first IMS
norm must be applied to their complete connected insertion.

\subsection{The additive charge card}

\begin{definition}[Connected charge-derivative card]
\label{def:v6v29-charge-card}
On a canonical path interval \(I\), the CCD card holds with constant
\(M_I\) if
\begin{equation}
 \left\|
 \kappa_{r+1}^{e}
 (F_1,\ldots,F_r,
  \partial_eS_m-\mathbb E_e\partial_eS_m)
 \right\|_{\rm IMS}
 \le M_I\prod_i\|F_i\|_{\rm end}
 \label{eq:v6v29-charge-card}
\end{equation}
uniformly in \(e\in I\), cutoff, volume, and \(r\).
\end{definition}

\begin{theorem}[Additive-charge transport]
\label{thm:v6v29-additive-transport}
Assume the CCD card along every adjacent interval, with a common bound
\(M\), and assume the matched counterterm score has summable IMS norm.
Then the canonical photon--matter contribution to the V25 score clock is
\begin{equation}
 \eta_n^{e}
 \le M|e_{n+1}-e_n|+\eta_n^{\rm ct}.
 \label{eq:v6v29-additive-clock}
\end{equation}
It is summable whenever
\begin{equation}
 \sum_n|e_{n+1}-e_n|<\infty,
 \qquad
 \sum_n\eta_n^{\rm ct}<\infty.
 \label{eq:v6v29-additive-variation}
\end{equation}
\end{theorem}

\begin{proof}
Use the connected scale-score identity of V25 with
Proposition~\ref{prop:v6v29-canonical-score}.  Parametrize the adjacent
path monotonically in \(e\).  The CCD card bounds its matter insertion by
\(M|de|\); integrate to obtain the first term of
\eqref{eq:v6v29-additive-clock}.  Add the counterterm clock and sum.
\end{proof}

\begin{corollary}[Asymptotically vanishing charge has finite additive
length]
\label{cor:v6v29-zero-charge-length}
If \(e_n\downarrow0\), then
\begin{equation}
 \sum_n|e_{n+1}-e_n|=e_0<\infty.
 \label{eq:v6v29-zero-charge-variation}
\end{equation}
Thus a uniformly proved CCD card would permit a vanishing-charge tail even
though the fixed-field logarithmic clock diverges.
\end{corollary}

\begin{proof}
The monotone variation telescopes.
\end{proof}

This is the coordinate appropriate to a candidate asymptotically free
sector.  It does not prove that hypercharge or either non-Abelian coupling
follows such a trajectory.

\begin{corollary}[Diverging charge has infinite additive length]
\label{cor:v6v29-diverging-charge}
If \(e_n\to\infty\), then
\begin{equation}
 \sum_n|e_{n+1}-e_n|=\infty.
 \label{eq:v6v29-diverging-variation}
\end{equation}
Consequently a CCD card with a nondecaying uniform constant cannot close
that trajectory by additive transport.
\end{corollary}

\begin{proof}
The total variation dominates \(|e_N-e_0|\), which tends to infinity.
\end{proof}

The hypercharge obstruction therefore remains, but for the intrinsic
reason that its canonical interaction coordinate leaves every bounded
interval, not because the free Gaussian normalization was held fixed.

\subsection{Relation to the susceptibility recursion}

Since \(e=\beta^{-1/2}\), an adjacent inverse-coupling increment obeys
\begin{equation}
 e_{n+1}-e_n
 =e_n\left[
 \left(1+\frac{\delta\beta_n}{\beta_n}\right)^{-1/2}-1
 \right].
 \label{eq:v6v29-e-beta-relation}
\end{equation}

\begin{lemma}[Susceptibility-to-charge conversion]
\label{lem:v6v29-susceptibility-conversion}
If \(|\delta\beta_n|\le\beta_n/2\), then
\begin{equation}
 |e_{n+1}-e_n|
 \le C e_n\frac{|\delta\beta_n|}{\beta_n}.
 \label{eq:v6v29-e-beta-bound}
\end{equation}
Hence the V28 Ward-first susceptibility card is sufficient for canonical
additive transport under the weaker weighted condition
\begin{equation}
 \sum_n e_n\frac{\mathfrak q_n}{\beta_n}<\infty,
 \label{eq:v6v29-weighted-susceptibility}
\end{equation}
plus the corresponding high and counterterm terms.
\end{lemma}

\begin{proof}
On \(|x|\le1/2\), the derivative of \((1+x)^{-1/2}\) is uniformly bounded.
Apply the mean-value theorem to
\eqref{eq:v6v29-e-beta-relation}, then use
\(|\delta\beta_n^{\rm light}|\le\mathfrak q_n\) from V28.
\end{proof}

When \(e_n\to0\), the extra factor \(e_n\) makes
\eqref{eq:v6v29-weighted-susceptibility} strictly weaker than the fixed-field
relative sum \eqref{eq:v6v28-current-sum}.

\subsection{Observable normalization}

Physical correlation functions must specify whether their photon endpoint
is \(A\), \(\mathcal A\), or the field strength.  Under the dilation,
\begin{equation}
 A=e\mathcal A,
 \qquad
 F(A)=eF(\mathcal A)
 \label{eq:v6v29-endpoint-dilation}
\end{equation}
for a linear Abelian field strength.

\begin{proposition}[Endpoint factors are part of matching]
\label{prop:v6v29-endpoint-matching}
Transport in canonical variables yields uniform bounds for canonical
endpoints.  Bounds for unnormalized photon endpoints acquire the explicit
powers of \(e_n\) dictated by \eqref{eq:v6v29-endpoint-dilation}.  These
powers must be included in the V14 interface and continuum normalization;
they cannot be hidden in the transport constant.
\end{proposition}

\begin{proof}
Substitute \(A=e\mathcal A\) in each multilinear observable.  Every linear
photon or field-strength endpoint contributes one factor of \(e\).  This is
an exact algebraic change of variables.
\end{proof}

\begin{firewall}[The CCD card is unproved]
\label{fw:v6v29-CCD-open}
Theorem~\ref{thm:v6v29-additive-transport} is conditional on a uniform
connected charge-derivative card.  Such a card has not been proved for the
retained chiral Standard Model fields, and it may fail at strong charge or
at determinant-line singularities.
\end{firewall}

\begin{firewall}[Canonical dilation does not remove the Landau problem]
\label{fw:v6v29-no-Landau-cure}
Changing variables removes the free Gaussian normalization clock.  It does
not prevent \(e\) from growing, construct a bounded ultraviolet trajectory,
or improve the bad-block rate.  A trajectory with \(e\to\infty\) still has
infinite additive length and leaves every interval on which a uniform CCD
card could be established perturbatively.
\end{firewall}

\begin{theorem}[Sixth-series V29 result]
\label{thm:v6v29-status}
The fixed-field photon clock \(|\Delta\log\beta|\) is a sufficient but
coordinate-dependent transport cost.  The exact canonical dilation
\(\mathcal A=\sqrt\beta A\) fixes the free Gaussian law and moves the
running into \(e=\beta^{-1/2}\).  The complete matched score is then one
centered charge derivative plus counterterms, and a uniform CCD card yields
the additive clock \(|\Delta e|\).  A monotone trajectory with \(e\to0\)
has finite additive length, while \(e\to\infty\) has infinite length.  The
V28 susceptibility controls this clock through the weighted sum
\(\sum e_n\mathfrak q_n/\beta_n\).  The hypercharge obstruction remains,
but no longer rests on an avoidable field normalization.
\end{theorem}

\begin{proof}
Use Theorems~\ref{thm:v6v29-Gaussian-dilation},
\ref{thm:v6v29-free-score-cancellation}, and
\ref{thm:v6v29-additive-transport},
Corollaries~\ref{cor:v6v29-zero-charge-length} and
\ref{cor:v6v29-diverging-charge},
Lemma~\ref{lem:v6v29-susceptibility-conversion}, and
Proposition~\ref{prop:v6v29-endpoint-matching}, subject to both firewalls.
\end{proof}

V30 will perform the next compulsory companion audit.  It will then test
the CCD card on a finite-volume retained charged shell by differentiating
the complete current cumulant in the canonical coupling, keeping the
seagull, phase, and Ward-completion terms assembled before norms.


\clearpage
\section{V30: compulsory audit and the covariant charge-score leaf}
\label{sec:v6v30}

\subsection{Read-only companion audit}

The active companion files inspected at this checkpoint were
\begin{center}
\begin{tabular}{p{0.62\textwidth}r}
\hline
file & bytes\\
\hline
\texttt{Renewal\_Yang\_Mills\_Mass\_Gap\_Research\_Notes\_V23.tex}
& \(343\,899\)\\
\texttt{Renewal\_NS\_Stability\_Research\_Notes\_V31.tex}
& \(887\,537\)\\
\hline
\end{tabular}
\end{center}
Their SHA--256 digests were
\begin{align}
 &\texttt{138b0be3c3ed5f1eb3bfeb9f3acfc95a7d240c3599eedbfe8d29c4bd8f2a04e5},
 \label{eq:v6v30-YM-hash}\\
 &\texttt{f1121b62238ad5491bb7cf03635ea9934942edf573ed8faaa9ee79e1d38a6696}.
 \label{eq:v6v30-NS-hash}
\end{align}
In the second digest the visual space is typographical; the exact recorded
digest is the uninterrupted string
\texttt{f1121b62238ad5491bb7cf03635ea9934942edf573ed8faaa9ee79e1d38a6696}.
Neither file was modified.

Yang--Mills V20--V23 proves a uniform spectral gap along an intrinsic
one-well interpolation and obtains small value and gradient score rows.  It
also proves that differentiating a principal-chart cutoff creates boundary
distributions whose translated-Gaussian derivative rows grow
superexponentially.  That chart route is abandoned in favor of an
intrinsic score leaf with a negative Sobolev norm.

The nonzero transverse photon shell in V26 is a linear vector space with no
radial chart boundary, so the Yang--Mills surface obstruction is absent
there.  Compact torons, large-gauge fundamental regions, and group-valued
non-Abelian links do have global boundaries and must remain intrinsic.  The
NS V31 endpoint is unchanged and continues to support the same order of
operations: form the complete correlated object before absolute norms.

\begin{firewall}[No gap or score row is transferred]
\label{fw:v6v30-no-transfer}
The Yang--Mills gap concerns a compact one-link law.  It does not give a gap
for massless charged matter, a photon--fermion shell, a gravitational
constraint quotient, or the Standard Model determinant phase.  The
negative-norm method is re-derived below with independent hypotheses.
\end{firewall}

\subsection{Canonical dilation moves the gauge quotient}

V29 uses \(\mathcal A=\sqrt\beta A\) and \(e=\beta^{-1/2}\).  If the
unnormalized connection transforms as
\begin{equation}
 A\longmapsto A-\nabla\omega,
 \label{eq:v6v30-A-gauge-map}
\end{equation}
then the canonical field transforms as
\begin{equation}
 \mathcal A\longmapsto
 \mathcal A-e^{-1}\nabla\omega.
 \label{eq:v6v30-canonical-gauge-map}
\end{equation}
Thus the coordinate description of a gauge orbit depends on \(e\).

Let \(X_e\) be the reduced configuration space at charge \(e\), and choose
a family of measurable differentiable identifications
\begin{equation}
 T_e:X_0\longrightarrow X_e
 \label{eq:v6v30-quotient-identification}
\end{equation}
on one constant-rank chart.  Pull every density and observable back to the
fixed space \(X_0\).

\begin{theorem}[Covariant moving-law cumulant identity]
\label{thm:v6v30-covariant-score}
Let \(\widetilde\nu_e:=T_e^*\nu_e\),
\(\widetilde F_{i,e}:=F_{i,e}\circ T_e\), and let
\begin{equation}
 \widetilde{\mathcal S}_e
 :=\partial_e\log\frac{d\widetilde\nu_e}{d\lambda_0}.
 \label{eq:v6v30-pulled-score}
\end{equation}
Then
\begin{align}
 \frac d{de}\kappa_r^{\nu_e}(F_{1,e},\ldots,F_{r,e})
 &=\kappa_{r+1}^{\widetilde\nu_e}
 (\widetilde F_{1,e},\ldots,\widetilde F_{r,e},
  \widetilde{\mathcal S}_e)
 \notag\\
 &\quad+
 \sum_{i=1}^r
 \kappa_r^{\widetilde\nu_e}
 (\widetilde F_{1,e},\ldots,
  \partial_e\widetilde F_{i,e},\ldots,
  \widetilde F_{r,e}).
 \label{eq:v6v30-covariant-score}
\end{align}
\end{theorem}

\begin{proof}
The pullback turns the problem into a differentiable family of normalized
laws and moving observables on one fixed measure space.  Apply the moving
version of Theorem~\ref{thm:v6v25-score-cumulant}.
\end{proof}

\begin{definition}[Horizontal physical observable]
\label{def:v6v30-horizontal-observable}
A family \(F_e\) is horizontal for \(T_e\) if
\begin{equation}
 \partial_e(F_e\circ T_e)=0.
 \label{eq:v6v30-horizontal-observable}
\end{equation}
\end{definition}

Only for horizontal endpoints do the moving-observable terms in
\eqref{eq:v6v30-covariant-score} vanish.

\begin{proposition}[Moving Ward projector]
\label{prop:v6v30-moving-Ward}
If \(\mathcal P_e\) is the Ward or physical-state projector in canonical
coordinates, then
\begin{equation}
 \partial_e(\mathcal P_eA_e)
 =\mathcal P_e\partial_eA_e+dot{\mathcal P}_eA_e.
 \label{eq:v6v30-moving-Ward}
\end{equation}
The second term is absent only if the quotient identification conjugates
\(\mathcal P_e\) to a fixed projector on \(X_0\).
\end{proposition}

\begin{proof}
This is the product rule.  If
\(T_e^{-1}\mathcal P_eT_e\) is constant, its derivative vanishes after
pullback; conversely a varying conjugated projector contributes the stated
term.
\end{proof}

\begin{assessment}[Append-only qualification of V29]
\label{ass:v6v30-V29-qualification}
The additive charge theorem V29 is valid for a fixed transverse shell with
horizontally identified endpoints and Ward projector.  In the full gauge
theory, the quotient-transport and
\(\dot{\mathcal P}_e\) terms must be included in the CCD card.  Canonical
dilation does not make them vanish automatically.
\end{assessment}

\subsection{The complete cumulant Wick polynomial}

Let \(F_I:=\prod_{i\in I}F_i\) in a fixed canonical order, and let
\(\Pi(E)\) be the set of partitions of a finite set \(E\), including the
empty partition of \(\varnothing\).  Define
\begin{equation}
 \mathfrak W_r^\nu(F_1,\ldots,F_r)
 :=\sum_{\varnothing\ne B\subseteq[r]}
 \sum_{\pi\in\Pi([r]\setminus B)}
 |\pi|!(-1)^{|\pi|}
 F_B\prod_{C\in\pi}\mathbb E_\nu F_C.
 \label{eq:v6v30-Wick-polynomial}
\end{equation}

\begin{theorem}[Score--Wick identity]
\label{thm:v6v30-score-Wick}
For every centered scalar score \(S\),
\begin{equation}
 \boxed{\qquad
 \kappa_{r+1}^\nu(F_1,\ldots,F_r,S)
 =\mathbb E_\nu
 [S\,\mathfrak W_r^\nu(F_1,\ldots,F_r)].
 \qquad}
 \label{eq:v6v30-score-Wick}
\end{equation}
\end{theorem}

\begin{proof}
Expand the cumulant by partitions of \([r]\cup\{*\}\), where \(*\)
labels \(S\).  Partitions in which \(\{*\}\) is a singleton vanish because
\(\mathbb ES=0\).  In every remaining partition, the score block is
uniquely \(B\cup\{*\}\) with nonempty \(B\subseteq[r]\); the remaining
blocks form \(\pi\in\Pi([r]\setminus B)\).  The total number of blocks is
\(|\pi|+1\), whose cumulant coefficient is
\(|\pi|!(-1)^{|\pi|}\).  Factoring the score expectation gives exactly
\eqref{eq:v6v30-Wick-polynomial}.
\end{proof}

The proper-partition cancellations are contained inside one polynomial.
They must not be estimated term by term.

\subsection{One negative-norm score leaf}

Let \(L\ge0\) be a self-adjoint reversible generator on
\(L^2(\nu)\), with constants as its kernel.  Let \(Q\) project orthogonally
off constants and set \(R=(L|_{Q})^{-1}\) on its domain.  Define
\begin{equation}
 \|S\|_{-1,L}^2:=\langle S,RS\rangle
 =\|dR S\|_2^2,
 \label{eq:v6v30-negative-norm}
\end{equation}
where the last equality uses the Dirichlet differential of \(L=d^*d\).

\begin{theorem}[Negative-norm score-leaf bound]
\label{thm:v6v30-score-leaf}
For a centered score \(S\) in the domain of \(R^{1/2}\),
\begin{equation}
 |\kappa_{r+1}^\nu(F_1,\ldots,F_r,S)|
 \le
 \|S\|_{-1,L}
 \|d\mathfrak W_r^\nu(F_1,\ldots,F_r)\|_2.
 \label{eq:v6v30-score-leaf}
\end{equation}
If \(L\) has gap \(\lambda>0\), then
\begin{equation}
 \|S\|_{-1,L}\le\lambda^{-1/2}\|S\|_2.
 \label{eq:v6v30-gap-negative}
\end{equation}
\end{theorem}

\begin{proof}
The score--Wick identity and centering give
\[
 \mathbb E[S\mathfrak W_r]
 =\langle LR S,\mathfrak W_r\rangle
 =\langle dR S,d\mathfrak W_r\rangle.
\]
Cauchy--Schwarz proves \eqref{eq:v6v30-score-leaf}.  The spectral theorem
gives \(R\le\lambda^{-1}Q\), proving
\eqref{eq:v6v30-gap-negative}.
\end{proof}

This is noncircular: it asks for a derivative of the complete Wick
polynomial, not the desired cumulant bound at arity \(r+1\).

\subsection{The charge-score leaf}

In canonical variables let the full pulled-back matched charge score be
\begin{equation}
 S_e^{\rm ch}:=
 Q\left[
 -\partial_eS_m-partial_eS_{\rm ct}
 +S_e^{\rm quot}
 \right],
 \label{eq:v6v30-full-charge-score}
\end{equation}
where \(S_e^{\rm quot}\) contains the Jacobian and action terms generated by
the quotient identification \(T_e\).  Moving endpoint and Ward-projector
terms remain those in
\eqref{eq:v6v30-covariant-score}--
\eqref{eq:v6v30-moving-Ward}.

\begin{definition}[Charge-score negative-norm card]
\label{def:v6v30-charge-negative-card}
The CSN card holds on a coupling interval \(I\) if there is \(a_I<\infty\)
such that
\begin{equation}
 \sup_{e\in I}\|S_e^{\rm ch}\|_{-1,L_e}\le a_I
 \label{eq:v6v30-charge-negative-card}
\end{equation}
for a shellwise generator whose zero and physical massless modes are
retained rather than inverted.
\end{definition}

\begin{definition}[Wick derivative card]
\label{def:v6v30-Wick-derivative-card}
The WDC holds if
\begin{equation}
 \|d\mathfrak W_r^{\nu_e}(F_1,ldots,F_r)\|_2
 \le r!K^r\mathcal T_r^{\rm hyb}
 \prod_i\|F_i\|_{\rm end}
 \label{eq:v6v30-Wick-derivative-card}
\end{equation}
after complete Ward projection, uniformly in \(e\in I\).
\end{definition}

\begin{theorem}[CSN plus WDC implies the CCD card]
\label{thm:v6v30-CSN-WDC}
Assume the CSN and WDC, horizontal endpoints, and a horizontally constant
Ward projector.  Then the connected charge-derivative card of V29 holds
with constant \(M_I\le a_I\) times the WDC base.  Consequently the additive
charge transport theorem applies on \(I\).
\end{theorem}

\begin{proof}
Apply Theorem~\ref{thm:v6v30-score-leaf} with
\(S=S_e^{\rm ch}\), then use Definitions
\ref{def:v6v30-charge-negative-card} and
\ref{def:v6v30-Wick-derivative-card}.  Horizontal transport removes the
moving endpoint and projector terms.  Divide by the IMS denominator to
obtain the CCD bound, and invoke
Theorem~\ref{thm:v6v29-additive-transport}.
\end{proof}

\subsection{Compatibility with a massless charged shell}

A full volume-uniform gap is neither expected nor assumed.  Decompose the
matter and photon generator into scale shells and invert only on the
centered complement of the retained low modes.

\begin{proposition}[Shellwise negative-norm scaling]
\label{prop:v6v30-shell-negative}
If the shell generator has gap \(\lambda_j\ge cL_j^{-z}\) and the centered
charge score satisfies
\begin{equation}
 \|S_{e,j}^{\rm ch}\|_2\le C L_j^{-s},
 \label{eq:v6v30-score-L2-scale}
\end{equation}
then
\begin{equation}
 \|S_{e,j}^{\rm ch}\|_{-1}
 \le C L_j^{z/2-s}.
 \label{eq:v6v30-score-negative-scale}
\end{equation}
Its shell sum is absolutely convergent if \(s>z/2\).
\end{proposition}

\begin{proof}
Use \eqref{eq:v6v30-gap-negative}.  The dyadic series
\(\sum_jL_j^{z/2-s}\) converges exactly when its exponent is negative.
\end{proof}

At the borderline \(s=z/2\), a fixed gap estimate gives one constant per
shell and cannot prove the CCD card.  An additional Ward cancellation,
running clock, or exact telescoping identity is required.

\begin{firewall}[Grassmann fields do not supply a reversible probability
generator]
\label{fw:v6v30-Grassmann}
The negative-norm argument is literal for a positive modulus backbone or a
reflection-positive transfer representation.  Retained chiral Grassmann
fields are not classical random variables with a reversible Markov
generator.  Applying Theorem~\ref{thm:v6v30-score-leaf} to them requires an
operator/transfer analogue and the determinant-line conditions of V19.
\end{firewall}

\begin{firewall}[WDC is the live all-order gate]
\label{fw:v6v30-WDC-open}
The complete Wick polynomial is exact, but no uniform derivative bound
\eqref{eq:v6v30-Wick-derivative-card} has been proved for the coupled
massless SM--Einstein skeleton.  Expanding its partition terms separately
would reintroduce the factorial leakage identified by the companion
program and is not an admissible proof.
\end{firewall}

\begin{theorem}[Sixth-series V30 result]
\label{thm:v6v30-status}
The compulsory audit confirms that boundary-differentiated compact charts
cannot provide all-order score rows; compact zero modes must remain
intrinsic.  Canonical photon dilation also moves the gauge quotient, so the
exact coupling derivative contains quotient, endpoint, and Ward-projector
transport terms unless a horizontal identification is constructed.  After
that correction, the complete charge-score cumulant equals one centered
score paired with an exact cumulant Wick polynomial.  A single
negative-norm score leaf plus a Ward-first Wick derivative card implies the
V29 CCD card without recursive use of the desired arity bound.  The WDC,
the shellwise massless score gain, and a transfer version for retained
chiral fermions remain open.
\end{theorem}

\begin{proof}
Use Theorem~\ref{thm:v6v30-covariant-score},
Proposition~\ref{prop:v6v30-moving-Ward},
Theorems~\ref{thm:v6v30-score-Wick} and
\ref{thm:v6v30-score-leaf},
Theorem~\ref{thm:v6v30-CSN-WDC}, and
Proposition~\ref{prop:v6v30-shell-negative}, subject to the three
firewalls.
\end{proof}

V31 will differentiate the complete cumulant Wick polynomial before
expanding its partition coefficients.  It will derive an exact rooted
subset recursion for \(d\mathfrak W_r\), compute its coefficient census on
stars and balanced trees, and decide whether WDC has exponential or
factorial grade.


\clearpage
\section{V31: the Wick derivative census and the failure of an unpaired norm}
\label{sec:v6v31}

\subsection{Exact rooted derivative polynomial}

V30 reduces a score cumulant to a negative-norm score leaf paired with the
derivative of the complete cumulant Wick polynomial.  We now differentiate
that polynomial exactly before taking norms.

For a finite label set \(I\), define the Appell--cumulant polynomial
\begin{equation}
 \mathfrak R_I^\nu(F)
 :=\sum_{A\subseteq I}
 \sum_{\pi\in\Pi(I\setminus A)}
 |\pi|!(-1)^{|\pi|}
 F_A\prod_{C\in\pi}\mathbb E_\nu F_C,
 \label{eq:v6v31-rooted-polynomial}
\end{equation}
with \(F_\varnothing=1\) and \(\mathfrak R_\varnothing=1\).

\begin{theorem}[Exact derivative of the cumulant Wick polynomial]
\label{thm:v6v31-Wick-derivative}
For every derivation \(d\) acting on the configuration dependence of the
observables and leaving their scalar expectations fixed,
\begin{equation}
 \boxed{\qquad
 d\mathfrak W_r^\nu(F_1,\ldots,F_r)
 =\sum_{i=1}^r
 (dF_i)\,
 \mathfrak R_{[r]\setminus\{i\}}^\nu(F).
 \qquad}
 \label{eq:v6v31-Wick-derivative}
\end{equation}
\end{theorem}

\begin{proof}
Differentiate \eqref{eq:v6v30-Wick-polynomial}.  For each nonempty
unaveraged block \(B\), the product rule chooses one \(i\in B\) and leaves
\((dF_i)F_{B\setminus\{i\}}\).  Fix \(i\) and set
\(A=B\setminus\{i\}\).  Then \(A\) runs over every subset of
\([r]\setminus\{i\}\), while the averaged partition is any partition of
its complement.  The coefficient remains
\(|\pi|!(-1)^{|\pi|}\), giving
\eqref{eq:v6v31-rooted-polynomial}.
\end{proof}

\begin{lemma}[Normalized exponential generating function]
\label{lem:v6v31-generating-function}
For commuting scalar inputs,
\begin{equation}
 \mathfrak R_I^\nu(F)
 =\left.
 \partial_{t_I}
 \frac{\exp(\sum_{i\in I}t_iF_i)}
 {\mathbb E_\nu\exp(\sum_{i\in I}t_iF_i)}
 \right|_{t=0},
 \label{eq:v6v31-generating-function}
\end{equation}
where \(\partial_{t_I}=\prod_{i\in I}\partial_{t_i}\).
In particular, \(\mathbb E\mathfrak R_I=0\) for nonempty \(I\).
\end{lemma}

\begin{proof}
Expand the numerator by the subset of labels differentiated there.  The
multivariate reciprocal identity
\[
 \left(\mathbb Ee^{\sum t_iF_i}\right)^{-1}
 =\sum_{E}
 \sum_{\pi\in\Pi(E)}
 |\pi|!(-1)^{|\pi|}
 \prod_{C\in\pi}\mathbb EF_C
 \prod_{i\in E}t_i
\]
at square-free coefficient level supplies the complement partition and
the coefficient in \eqref{eq:v6v31-rooted-polynomial}.  The normalized
exponential has expectation one, so every nonconstant square-free
coefficient is centered.
\end{proof}

\subsection{Exact subset recursion}

Let \(\kappa(F_i,F_B)\) denote the joint cumulant of \(F_i\) with the
variables labelled by \(B\), with
\(\kappa(F_i,F_\varnothing)=\mathbb EF_i\).

\begin{theorem}[Appell subset recursion]
\label{thm:v6v31-Appell-recursion}
For \(i\notin I\),
\begin{equation}
 \mathfrak R_{I\cup\{i\}}
 =F_i\mathfrak R_I
 -\sum_{B\subseteq I}
 \kappa(F_i,F_B)\,
 \mathfrak R_{I\setminus B}.
 \label{eq:v6v31-Appell-recursion}
\end{equation}
\end{theorem}

\begin{proof}
Write the normalized exponential as
\(G(t)=e^{t\cdot F-K(t)}\), where
\(K(t)=\log\mathbb Ee^{t\cdot F}\) is the cumulant generating function.
Then
\[
 \partial_{t_i}G=(F_i-\partial_{t_i}K)G.
\]
Apply \(\partial_{t_I}\), use the product rule, and evaluate at zero.  The
derivative \(\partial_{t_B}\partial_{t_i}K(0)\) is
\(\kappa(F_i,F_B)\), while the complementary derivative of \(G\) is
\(\mathfrak R_{I\setminus B}\).  This proves the recursion.
\end{proof}

The recursion uses cumulants of strictly smaller total arity than the score
cumulant being proved.  It is therefore compatible with strong induction,
provided its subset sum is controlled before spatial norms.

\subsection{The coefficient census is only factorial-exponential}

Let
\begin{equation}
 a_m:=\sum_{\pi\in\Pi([m])}|\pi|!
 =\sum_{k=0}^m k!S(m,k)
 \label{eq:v6v31-ordered-Bell}
\end{equation}
be the ordered Bell number.

\begin{lemma}[Ordered-partition bound]
\label{lem:v6v31-ordered-Bell-bound}
There is a numerical \(C\) such that
\begin{equation}
 a_m\le C^m m!
 \label{eq:v6v31-ordered-Bell-bound}
\end{equation}
for every \(m\ge0\).
\end{lemma}

\begin{proof}
The exponential generating function is
\begin{equation}
 \sum_{m\ge0}\frac{a_m}{m!}z^m
 =\frac1{2-e^z}.
 \label{eq:v6v31-ordered-Bell-EGF}
\end{equation}
Indeed, an ordered partition is a finite ordered sequence of nonempty
labelled sets, whose exponential generating function is
\(1/[1-(e^z-1)]\).  This function is analytic on every disk
\(|z|<\log2\).  Cauchy's coefficient estimate on a fixed smaller disk gives
\eqref{eq:v6v31-ordered-Bell-bound}.
\end{proof}

\begin{proposition}[Universal pointwise factorial grade]
\label{prop:v6v31-universal-grade}
If \(|F_i|\le1\), then
\begin{equation}
 |\mathfrak R_I(F)|\le C^{|I|}|I|!.
 \label{eq:v6v31-R-bound}
\end{equation}
If additionally \(|dF_i|\le b_i\), then
\begin{equation}
 |d\mathfrak W_r(F)|
 \le C^rr!\max_i b_i.
 \label{eq:v6v31-dW-bound}
\end{equation}
\end{proposition}

\begin{proof}
For \(|I|=n\), take absolute values in
\eqref{eq:v6v31-rooted-polynomial}:
\[
 |\mathfrak R_I|
 \le\sum_{a=0}^n\binom na a_{n-a}.
\]
Use Lemma~\ref{lem:v6v31-ordered-Bell-bound} and
\(\binom na(n-a)!=n!/a!\).  The remaining sum over \(a\) is bounded by an
exponential base.  Insert this bound into
\eqref{eq:v6v31-Wick-derivative}; the factor \(r(r-1)!=r!\) proves the last
claim.
\end{proof}

Thus no extra factorial beyond the allowed observable \(r!\) is forced by
the algebraic coefficient census.  The live issue is spatial connectedness,
not arity counting.

\subsection{Counterexample to the unpaired spatial WDC}

Consider a product probability space
\(\Omega=\prod_{i=1}^r\Omega_i\).  Let \(F_i\) be a centered nonconstant
bounded function depending only on coordinate \(i\), with
\(\mathbb EF_i^2=1\) and \(\mathbb E|d_iF_i|^2>0\).  Place their anchors at
arbitrarily separated spatial sites.

\begin{lemma}[Independent Wick product]
\label{lem:v6v31-independent-Wick}
For the independent centered family,
\begin{equation}
 \mathfrak R_{[r]}(F_1,\ldots,F_r)=\prod_{i=1}^rF_i.
 \label{eq:v6v31-independent-Wick}
\end{equation}
\end{lemma}

\begin{proof}
Independence gives
\(\mathbb Ee^{\sum t_iF_i}=\prod_i\mathbb Ee^{t_iF_i}\).  The normalized
generating function in \eqref{eq:v6v31-generating-function} factors into
\(\prod_i[e^{t_iF_i}/\mathbb Ee^{t_iF_i}]\).  Its square-free coefficient
is \(\prod_i(F_i-\mathbb EF_i)=\prod_iF_i\).
\end{proof}

\begin{theorem}[The unpaired WDC is false in a decaying tree norm]
\label{thm:v6v31-WDC-false}
For the family above,
\begin{equation}
 \|d\mathfrak R_{[r]}\|_2^2
 =\sum_{i=1}^r\mathbb E|d_iF_i|^2>0,
 \label{eq:v6v31-independent-gradient}
\end{equation}
independently of the distances between anchors.  Therefore no estimate of
the form
\begin{equation}
 \|d\mathfrak W_r\|_2
 \le r!K^r\mathcal T_r^{\rm hyb}(x_1,\ldots,x_r)
 \label{eq:v6v31-false-spatial-WDC}
\end{equation}
can hold for every such family when the hybrid connected tree weight tends
to zero as the anchors split into distant groups.
\end{theorem}

\begin{proof}
Different coordinate gradients are orthogonal in the product Dirichlet
form.  Differentiate \eqref{eq:v6v31-independent-Wick}; independence and
\(\mathbb EF_j^2=1\) give
\eqref{eq:v6v31-independent-gradient}.  The right side of
\eqref{eq:v6v31-false-spatial-WDC} tends to zero when all connecting
massive or power-law root weights tend to zero with separation, while the
left side is fixed and positive.
\end{proof}

\begin{assessment}[Append-only correction to V30]
\label{ass:v6v31-V30-correction}
The score--Wick identity and negative-norm pairing in V30 remain exact.
The unpaired spatial WDC
\eqref{eq:v6v30-Wick-derivative-card} is false as a general hypothesis for
decaying connected tree weights.  Consequently Theorem
\ref{thm:v6v30-CSN-WDC} remains a valid implication from an overstrong
hypothesis but is not the correct acquisition target.
\end{assessment}

The counterexample does not contradict connected score decay.  If a local
score is independent of one of the separated factors, the complete
cumulant on the left of \eqref{eq:v6v30-score-Wick} vanishes.  The global
Cauchy--Schwarz step loses that cancellation by norming the Wick gradient
before pairing it with the score current.

\subsection{The paired local score-current card}

Let \(u_S=R S\) and define its Dirichlet current
\begin{equation}
 j_S(x):=d_xu_S.
 \label{eq:v6v31-score-current}
\end{equation}
Then the exact negative-norm identity is spatially resolved:
\begin{equation}
 \kappa_{r+1}(F_1,ldots,F_r,S)
 =\sum_x\mathbb E
 \left[j_S(x)\cdot d_x\mathfrak W_r(F)\right].
 \label{eq:v6v31-local-pairing}
\end{equation}

\begin{definition}[Paired local Wick card]
\label{def:v6v31-paired-card}
For a score anchored at a support \(A\), the PLWC holds with activity
\(a_S\) if
\begin{equation}
 \left|
 \sum_x\mathbb E
 [j_S(x)\cdot d_x\mathfrak W_r(F)]
 \right|
 \le a_S r!K^r
 \mathcal T_{r,A}^{\rm hyb}(x_1,ldots,x_r)
 \prod_i\|F_i\|_{\rm end},
 \label{eq:v6v31-paired-card}
\end{equation}
where \(\mathcal T_{r,A}^{\rm hyb}\) is the hybrid tree sum on the
observable roots together with one distinguished score root.  Ward and
proper-partition sums are completed inside the expectation.
\end{definition}

\begin{theorem}[PLWC gives the connected charge-derivative card]
\label{thm:v6v31-PLWC-CCD}
If the complete canonical charge score decomposes into rooted score
supports with summable activities \(a_{S,A}\), and each satisfies the PLWC,
then the V29 CCD card holds with
\begin{equation}
 M\le K\sup_x\sum_{A:a(A)=x}a_{S,A}e^{\tau|A|},
 \label{eq:v6v31-PLWC-activity}
\end{equation}
after the distinguished score root is summed by the V23 mixed tree
compiler wherever its incident edge meets the strict scale criterion.
\end{theorem}

\begin{proof}
Equation \eqref{eq:v6v31-local-pairing} and the PLWC bound each rooted score
contribution before any global gradient norm.  Sum its support and anchor
with the stated local activity.  The mixed tree compiler removes the score
root when its massive or Ward-dressed massless attachment is summable; if
it is a physical nonsummable carrier, retain that edge in the IMS skeleton.
The resulting bound is exactly the CCD normalization of V29.
\end{proof}

\begin{firewall}[Separate negative and gradient norms lose connectedness]
\label{fw:v6v31-no-global-Cauchy}
One may use Cauchy--Schwarz locally inside a connected conditional or shell
decomposition, but not as the global product
\(\|S\|_{-1}\|d\mathfrak W_r\|_2\) when spatial decay is required.  The
independent-product counterexample proves that this separation discards the
needed cancellation.
\end{firewall}

\begin{firewall}[PLWC remains an analytic gate]
\label{fw:v6v31-PLWC-open}
Definition~\ref{def:v6v31-paired-card} is weaker and correctly connected,
but it has not been proved for the interacting massless Standard Model
current or gravitational score.  Its proof must use locality of the score
current together with the exact Appell recursion, without expanding
proper-partition terms separately.
\end{firewall}

\begin{theorem}[Sixth-series V31 result]
\label{thm:v6v31-status}
The derivative of the complete score Wick polynomial is the exact sum
\(\sum_i(dF_i)\mathfrak R_{[r]\setminus i}\).  Its ordered-partition
coefficient mass is only \(C^rr!\), so there is no new arity factorial.
However, a product of independent centered variables has a Wick gradient
whose global \(L^2\) norm does not decay with spatial separation.  The V30
unpaired WDC is therefore false in any decaying connected tree norm.  The
correct target is the paired local score-current card
\eqref{eq:v6v31-paired-card}, which retains the expectation and score root
until connected cancellation has occurred.
\end{theorem}

\begin{proof}
Use Theorems~\ref{thm:v6v31-Wick-derivative} and
\ref{thm:v6v31-Appell-recursion},
Lemma~\ref{lem:v6v31-ordered-Bell-bound},
Proposition~\ref{prop:v6v31-universal-grade},
Theorem~\ref{thm:v6v31-WDC-false}, and
Theorem~\ref{thm:v6v31-PLWC-CCD}, subject to both firewalls.
\end{proof}

V32 will attack the PLWC by conditioning on the score support.  It will
prove the exact cancellation of every observable component independent of
that support, derive the resulting connected support hypergraph, and test
whether its leaf pruning has one factorial or two.


\clearpage
\section{V32: exact score-support cancellation and the rooted support graph}
\label{sec:v6v32}

V31 showed that the score current must remain paired with the Wick
derivative until spatial connectedness has been extracted.  This note
proves that extraction in an exact finite-range product model.  The result
is not yet the interacting Standard Model estimate, but it isolates the
precise cancellation that an approximate mixing argument must preserve.

We first correct a harmless typographical omission in
\eqref{eq:v6v31-local-pairing}: its left side is
\(\kappa_{r+1}(F_1,\ldots,F_r,S)\).

\subsection{Local product model and Appell factorization}

Let
\begin{equation}
 (\Omega,\nu)=\prod_{y\in\Lambda}(\Omega_y,\nu_y),
 \qquad
 \mathcal F_B:=\sigma(\omega_y:y\in B),
 \label{eq:v6v32-product-space}
\end{equation}
where \(\Lambda\) is finite.  A random variable has support in \(B\) if it
is \(\mathcal F_B\)-measurable.  The coordinate derivation \(d_x\) obeys
the Leibniz rule and annihilates every variable measurable with respect to
\(\mathcal F_{\Lambda\setminus\{x\}}\).

\begin{lemma}[Factorization of the Appell--cumulant polynomial]
\label{lem:v6v32-Appell-factorization}
Suppose \(I=I_1\sqcup\cdots\sqcup I_q\) and the families
\((F_i)_{i\in I_a}\) are mutually independent.  Then
\begin{equation}
 \mathfrak R_I(F)=\prod_{a=1}^q\mathfrak R_{I_a}(F).
 \label{eq:v6v32-Appell-factorization}
\end{equation}
If \(I_a\ne\varnothing\), then
\(\mathbb E\mathfrak R_{I_a}=0\).
\end{lemma}

\begin{proof}
The joint moment generating function factors over the independent
families.  Hence the normalized exponential in
\eqref{eq:v6v31-generating-function} is the product of the corresponding
normalized exponentials.  Taking the square-free coefficient indexed by
\(I\) gives \eqref{eq:v6v32-Appell-factorization}.  Centering is the last
assertion of Lemma~\ref{lem:v6v31-generating-function}.
\end{proof}

For a fixed summand of the local score pairing, let \(J_x\) be a bounded
score current with support \(B_0\), where \(x\in B_0\).  Let \(F_i\) have
support \(B_i\), \(1\le i\le r\).  Define the rooted intersection graph
\(G_x\) on \(\{0,1,\ldots,r\}\) by
\begin{equation}
 p\sim q\quad\Longleftrightarrow\quad B_p\cap B_q\ne\varnothing.
 \label{eq:v6v32-intersection-graph}
\end{equation}

\begin{theorem}[Exact rooted support cancellation]
\label{thm:v6v32-rooted-cancellation}
In the product model \eqref{eq:v6v32-product-space},
\begin{equation}
 \boxed{\qquad
 G_x\text{ disconnected}
 \quad\Longrightarrow\quad
 \mathbb E\!left[J_x\,d_x\mathfrak W_r(F_1,\ldots,F_r)\right]=0.
 \qquad}
 \label{eq:v6v32-rooted-cancellation}
\end{equation}
Thus every nonzero local score term has a support graph joining the score
root to every observable root.
\end{theorem}

\begin{proof}
Use Theorem~\ref{thm:v6v31-Wick-derivative} and fix its summand indexed by
\(i\).  Let \(C\) be a connected component of \(G_x\) not containing the
root \(0\).

If \(i\in C\), then \(x\in B_0\) and \(B_i\cap B_0=\varnothing\), so
\(x\notin B_i\) and \(d_xF_i=0\).  If \(i\notin C\), the coordinate union
\(\bigcup_{k\in C}B_k\) is disjoint from the coordinate unions belonging
to all other graph components.  Lemma
\ref{lem:v6v32-Appell-factorization} therefore factors
\(\mathfrak R_{[r]\setminus\{i\}}\) into one factor
\(\mathfrak R_C\) and factors from the other components.  The first has
zero expectation, while it is independent of \(J_xd_xF_i\) and all the
other factors.  The expectation of this summand is zero.  Every derivative
summand vanishes, proving the claim.
\end{proof}

\begin{corollary}[Conditional projection form]
\label{cor:v6v32-conditional-projection}
For every \(x\),
\begin{equation}
 \mathbb E[J_xd_x\mathfrak W_r]
 =\mathbb E\!left[
 J_x\,\mathbb E(d_x\mathfrak W_r\mid\mathcal F_{B_0})
 \right].
 \label{eq:v6v32-conditional-projection}
\end{equation}
In the product model the conditional expectation on the right deletes
every support component not joined to \(B_0\).
\end{corollary}

\begin{proof}
The identity is the defining property of conditional expectation because
\(J_x\) is \(\mathcal F_{B_0}\)-measurable.  The deletion statement is
Theorem~\ref{thm:v6v32-rooted-cancellation} applied component by component.
\end{proof}

\subsection{From connected supports to rooted trees}

Let \(\mathbb T_{r+1}\) be the set of labelled spanning trees on
\(\{0,\ldots,r\}\).

\begin{lemma}[Rooted tree extraction]
\label{lem:v6v32-tree-extraction}
For arbitrary supports,
\begin{equation}
 \mathbf 1_{\{G_x\ {
m connected}\}}
 \le
 \sum_{T\in\mathbb T_{r+1}}
 \prod_{\{p,q\}\in E(T)}
 \mathbf 1_{\{B_p\cap B_q\ne\varnothing\}}.
 \label{eq:v6v32-tree-extraction}
\end{equation}
Moreover
\begin{equation}
 |\mathbb T_{r+1}|=(r+1)^{r-1}\le C^r r!.
 \label{eq:v6v32-Cayley-grade}
\end{equation}
\end{lemma}

\begin{proof}
Every finite connected graph contains a spanning tree, which proves the
indicator inequality.  Cayley's formula gives the equality.  The
elementary bound \(r!\ge(r/e)^r\), obtained by integrating \(\log t\),
and \((1+1/r)^r\le e\) give the last inequality after enlarging the
constant for the first two values of \(r\).
\end{proof}

The estimate \eqref{eq:v6v32-Cayley-grade} says that the connected support
census itself has the permitted factorial-exponential grade.  It must not
be multiplied blindly by the absolute coefficient census of V31.  The
Appell recursion performs the coefficient and connectivity pruning in one
subset convolution.

\begin{lemma}[One-factorial subset convolution]
\label{lem:v6v32-one-factorial-convolution}
For \(m\ge0\) and \(K\ge1\),
\begin{equation}
 \sum_{B\subseteq[m]}
 (|B|+1)!K^{|B|+1}(m-|B|)!K^{m-|B|}
 \le m!(4K)^{m+1}.
 \label{eq:v6v32-one-factorial-convolution}
\end{equation}
Consequently a leaf-pruning recursion in which a rooted connected block of
size \(|B|+1\) and a complementary Appell remainder of size \(m-|B|\)
each have factorial-exponential bounds still has only one factorial.
\end{lemma}

\begin{proof}
Group subsets by \(b=|B|\).  Since
\begin{equation}
 \binom mb(b+1)!(m-b)!=m!(b+1),
 \label{eq:v6v32-binomial-factorial-cancel}
\end{equation}
the left side is
\[
 K^{m+1}m!\sum_{b=0}^m(b+1)
 =K^{m+1}m!\frac{(m+1)(m+2)}2.
\]
The polynomial factor is at most \(4^{m+1}\), proving the stated bound.
The final assertion is precisely the same calculation with the numerical
constants of the two recursive factors included in \(K\).
\end{proof}

\begin{proposition}[Exact product-model PLWC]
\label{prop:v6v32-product-PLWC}
Assume in addition that \(|J_x|\le a_x\),
\(|d_xF_i|\le b_{i,x}\), and \(|F_i|\le1\).  Then
\begin{equation}
 \left|\mathbb E[J_xd_x\mathfrak W_r(F)]\right|
 \le a_x C^rr!\max_i b_{i,x}
 \mathbf 1_{\{G_x\ {
m connected}\}}.
 \label{eq:v6v32-product-PLWC}
\end{equation}
Combining this with \eqref{eq:v6v32-tree-extraction} gives a rooted support
tree bound.  If an overlap edge is dominated by the required massive or
Ward-dressed massless kernel, summing \(x\) gives the PLWC of V31.
\end{proposition}

\begin{proof}
Theorem~\ref{thm:v6v32-rooted-cancellation} inserts the connectedness
indicator.  On its support, use
\eqref{eq:v6v31-Wick-derivative} and
Proposition~\ref{prop:v6v31-universal-grade}.  This gives the first
inequality.  Lemma~\ref{lem:v6v32-tree-extraction} gives the tree form.
The last statement follows when each overlap indicator and the score
activity have a summable kernel majorant of the type required in
Definition~\ref{def:v6v31-paired-card}.
\end{proof}

\begin{firewall}[Product support is not interacting mixing]
\label{fw:v6v32-product-only}
Theorem~\ref{thm:v6v32-rooted-cancellation} uses exact independence of
disjoint coordinate sets.  In a gauge, fermion, or gravitational Gibbs
law, disjoint supports are generally correlated.  Replacing the zero in
\eqref{eq:v6v32-rooted-cancellation} by a decaying bound requires a
conditional mixing estimate uniform in the gauge quotient and in the
retained massless modes.
\end{firewall}

\begin{firewall}[Do not multiply the two censuses]
\label{fw:v6v32-no-double-census}
Bounding the complete Appell polynomial by \(C^rr!\) and then summing all
labelled trees by another \(C^rr!\) produces a false two-factorial budget.
The subset identity
\eqref{eq:v6v32-binomial-factorial-cancel} shows how a valid leaf-pruning
proof must combine the connected block and its complement before taking
the arity supremum.
\end{firewall}

\begin{theorem}[Sixth-series V32 result]
\label{thm:v6v32-status}
For a local product law, the paired score-current term vanishes exactly
unless the score support and every observable support lie in one rooted
intersection component.  A connected term admits a rooted spanning-tree
extraction, and the Appell subset recursion has a one-factorial convolution
census.  Hence neither exact support cancellation nor recursive leaf
pruning forces a second factorial.  The remaining analytic problem is to
replace exact product cancellation by a uniform conditional mixing kernel
for the interacting SM--Einstein law.
\end{theorem}

\begin{proof}
Combine Theorem~\ref{thm:v6v32-rooted-cancellation},
Lemma~\ref{lem:v6v32-tree-extraction},
Lemma~\ref{lem:v6v32-one-factorial-convolution}, and
Proposition~\ref{prop:v6v32-product-PLWC}, subject to both firewalls.
\end{proof}

V33 will formulate the needed replacement as a conditional oscillation
estimate.  It will derive a rooted martingale telescoping formula, identify
the exact operator norm that must decay along one support edge, and decide
whether this estimate follows from the existing one-well transfer gap or
requires a stronger spatial mixing input.


\clearpage
\section{V33: conditional projection, edge influence, and the gap obstruction}
\label{sec:v6v33}

The product cancellation in V32 has an exact formulation that survives for
an interacting law.  The score must first be projected onto its local
sigma algebra; only then should one use the generator and its negative
norm.  This order exposes a two-support conditional projection operator.

\subsection{Projection before the negative norm}

Let \(P_AH:=\mathbb E(H\mid\mathcal F_A)\) and
\(P_A^\circ H:=P_AH-\mathbb EH\).  These are orthogonal projections in
\(L^2(\nu)\), with \(P_A^\circ\) restricted to the centered subspace.

\begin{theorem}[Exact rooted projection identity]
\label{thm:v6v33-rooted-projection}
Let \(S_A\) be a centered score measurable with respect to
\(\mathcal F_A\).  For every \(H\in L^2(\nu)\),
\begin{equation}
 \langle S_A,H\rangle
 =\langle S_A,P_A^\circ H\rangle.
 \label{eq:v6v33-rooted-projection}
\end{equation}
If \(-\mathcal L\) is a nonnegative reversible generator, \(R\) is its
inverse on centered functions, and the required domains are satisfied,
then
\begin{equation}
 \langle S_A,H\rangle
 =\mathcal E(RS_A,P_A^\circ H).
 \label{eq:v6v33-projected-energy}
\end{equation}
In particular, with \(H=\mathfrak W_r(F)\), the score cumulant is obtained
by differentiating the rooted projection \(P_A^\circ\mathfrak W_r\), not
by first taking a global norm of \(d\mathfrak W_r\).
\end{theorem}

\begin{proof}
Because \(S_A\) is \(\mathcal F_A\)-measurable,
\(\langle S_A,H\rangle=\langle S_A,P_AH\rangle\).  Its mean is zero, so
the constant part of \(P_AH\) vanishes in the pairing.  This proves
\eqref{eq:v6v33-rooted-projection}.  On the centered subspace,
\(-\mathcal L RS_A=S_A\); reversibility and the definition of the
Dirichlet form give
\(\langle S_A,Q\rangle=\mathcal E(RS_A,Q)\) for centered \(Q\) in the
form domain.  Take \(Q=P_A^\circ H\).
\end{proof}

When the form has a local carré du champ,
\eqref{eq:v6v33-projected-energy} reads
\begin{equation}
 \langle S_A,H\rangle
 =\sum_x\mathbb E\!left[
 d_xRS_A\cdot d_xP_A^\circ H
 \right].
 \label{eq:v6v33-projected-current}
\end{equation}
The conditional projection now carries the cancellation lost in the
global Cauchy step of V30.

\subsection{The exact one-edge norm}

For two supports define the maximal conditional influence
\begin{equation}
 \rho_2(A,B)
 :=\|P_A^\circ P_B^\circ\|_{L^2_0(\nu)\to L^2_0(\nu)}.
 \label{eq:v6v33-rho2}
\end{equation}
It lies in \([0,1]\).

\begin{lemma}[One-edge covariance bound]
\label{lem:v6v33-one-edge}
If \(f\) is centered and \(\mathcal F_A\)-measurable and \(g\) is
centered and \(\mathcal F_B\)-measurable, then
\begin{equation}
 |\operatorname{Cov}_\nu(f,g)|
 \le \rho_2(A,B)\|f\|_2\|g\|_2.
 \label{eq:v6v33-one-edge}
\end{equation}
For a product law and disjoint \(A,B\), \(\rho_2(A,B)=0\).
\end{lemma}

\begin{proof}
Since \(f=P_A^\circ f\) and \(g=P_B^\circ g\), self-adjointness of the
conditional projections gives
\[
 \langle f,g\rangle
 =\langle f,P_A^\circ P_B^\circ g\rangle.
\]
Cauchy--Schwarz and the operator norm prove the estimate.  Under product
independence, conditioning a centered \(B\)-variable on the disjoint
coordinates in \(A\) gives zero.
\end{proof}

The energy version actually needed after
\eqref{eq:v6v33-projected-current} is
\begin{equation}
 \rho_{\mathcal E}(A,B)
 :=\sup_{\substack{g\in L^2_0(\mathcal F_B)\\g\ne0}}
 \frac{\|dP_A^\circ g\|_{L^2(\ell^2)}}{\|g\|_{\rm end}}.
 \label{eq:v6v33-energy-influence}
\end{equation}
The spatial acquisition target is a kernel bound
\begin{equation}
 \rho_{\mathcal E}(A,B)
 \le C|A|^{c}|B|^{c}
 \exp[-m\,d(A,B)]
 \label{eq:v6v33-energy-influence-decay}
\end{equation}
in a massive channel, with the corresponding Ward-dressed power kernel in
a physical massless channel.  Endpoint norms may replace the displayed
polynomial support factors if they already include them.

\subsection{Doob telescoping and separator chains}

Let \(A=U_0\subset U_1\subset\cdots\subset U_m\), and set
\(P_k=P_{U_k}\).  If \(H\) is \(\mathcal F_{U_m}\)-measurable, define
\(D_0H=P_0H-\mathbb EH\) and
\(D_kH=(P_k-P_{k-1})H\) for \(k\ge1\).

\begin{lemma}[Rooted Doob decomposition]
\label{lem:v6v33-Doob}
One has the exact orthogonal decomposition
\begin{equation}
 H-\mathbb EH=\sum_{k=0}^mD_kH,
 \qquad
 P_A^\circ H=D_0H,
 \label{eq:v6v33-Doob}
\end{equation}
and
\begin{equation}
 \|H-\mathbb EH\|_2^2=\sum_{k=0}^m\|D_kH\|_2^2.
 \label{eq:v6v33-Doob-orthogonal}
\end{equation}
Thus the connected score contribution is exactly the first martingale
increment in an exploration grown from the score support.
\end{lemma}

\begin{proof}
The first equality telescopes because \(P_mH=H\).  Conditional expectation
onto a nested sequence of sigma algebras is an orthogonal martingale
projection, so its successive differences are mutually orthogonal.  The
second equality and the norm identity follow.
\end{proof}

For a genuine Markov separator chain \(A_0,A_1,\ldots,A_q\), assume that
the conditional law makes the past and future independent given each
intermediate separator.  On centered functions transported from
\(A_q\) to \(A_0\), conditional expectation then factors through the
separators.

\begin{proposition}[Multiplication along a separator path]
\label{prop:v6v33-separator-product}
Under the preceding Markov separator assumption,
\begin{equation}
 \|P_{A_0}^\circ P_{A_q}^\circ\|_{2\to2}
 \le\prod_{k=0}^{q-1}\rho_2(A_k,A_{k+1}).
 \label{eq:v6v33-separator-product}
\end{equation}
Hence a uniform edge contraction \(\rho_2(A_k,A_{k+1})\le\theta<1\)
gives exponential decay in the number of separating layers.
\end{proposition}

\begin{proof}
The Markov property identifies the conditional expectation from the last
separator to the first with the composition of the successive conditional
expectation operators.  Remove constants at every step.  Submultiplicativity
of operator norms gives the displayed product.
\end{proof}

This proposition is the operator form of rooted tree pruning: each leaf is
conditioned through its parent separator before another leaf is exposed.
Lemma~\ref{lem:v6v32-one-factorial-convolution} then controls the subset
census without a second factorial.

\subsection{A spectral gap alone does not give edge decay}

\begin{theorem}[Gap-only obstruction]
\label{thm:v6v33-gap-obstruction}
A positive reversible spectral gap, without a locality or separator
hypothesis, does not imply any decay of \(\rho_2(A,B)\) with a declared
distance between \(A\) and \(B\).
\end{theorem}

\begin{proof}
Let \(u,v\in\{-1,1\}\) have law
\begin{equation}
 \nu_\varepsilon(u,v)=Z_\varepsilon^{-1}e^{\varepsilon uv},
 \qquad \varepsilon>0.
 \label{eq:v6v33-two-spin-law}
\end{equation}
Use the reversible two-site heat-bath generator that updates either spin
from its conditional law.  This is an irreducible generator on four
states, so its spectral gap \(\lambda_\varepsilon\) is strictly positive.
Neither the law nor the generator depends on a distance assigned to the
two coordinate labels.  Yet
\begin{equation}
 \mathbb Eu=\mathbb Ev=0,
 \qquad
 \mathbb E(uv)=\tanh\varepsilon,
 \label{eq:v6v33-two-spin-correlation}
\end{equation}
and Lemma~\ref{lem:v6v33-one-edge} implies
\(\rho_2(\{u\},\{v\})\ge\tanh\varepsilon\).  Assigning the two labels
arbitrarily large geometric separation leaves the gap and this lower bound
unchanged.  Therefore the gap alone contains no spatial decay statement.
\end{proof}

\begin{assessment}[What the YM one-well gap supplies]
\label{ass:v6v33-YM-gap}
The intrinsic one-well path gap audited in V30 controls the inverse
generator and therefore the size of a negative-norm score leaf.  By
Theorem~\ref{thm:v6v33-gap-obstruction}, it cannot by itself prove
\eqref{eq:v6v33-energy-influence-decay}.  A local propagation estimate for
the semigroup, or a uniform Markov-separator contraction, is an additional
input.
\end{assessment}

\begin{firewall}[Conditioning must respect the quotient]
\label{fw:v6v33-quotient-conditioning}
For gauge and diffeomorphism variables, \(\mathcal F_A\) cannot be chosen
by fixing raw compact coordinates independently across a cut.  It must be
defined on the physical quotient, or with boundary flux/edge variables
retained so that gluing is exact.  Otherwise
\(P_A\) breaks the Ward identity and the influence norm does not measure a
physical channel.
\end{firewall}

\begin{firewall}[Massless edges have no exponential target]
\label{fw:v6v33-massless-influence}
Equation~\eqref{eq:v6v33-energy-influence-decay} is a massive-channel
target.  A retained photon, graviton, or other physical massless carrier
must remain as a power-law edge in the IMS skeleton.  Only Ward-projected
gauge artifacts may be removed or improved by quotient conditioning.
\end{firewall}

\begin{theorem}[Sixth-series V33 result]
\label{thm:v6v33-status}
The exact interacting replacement for product cancellation is to project
the Wick polynomial onto the centered sigma algebra of the score support
before applying the negative-norm identity.  The required one-edge object
is the off-diagonal conditional influence
\(P_A^\circ P_B^\circ\), together with its energy version
\eqref{eq:v6v33-energy-influence}.  Doob exploration identifies the rooted
part as the first martingale increment, and Markov separators multiply
edge contractions.  A reversible spectral gap alone cannot supply spatial
decay; it must be paired with a locality or separator estimate.
\end{theorem}

\begin{proof}
Use Theorem~\ref{thm:v6v33-rooted-projection},
Lemma~\ref{lem:v6v33-one-edge},
Lemma~\ref{lem:v6v33-Doob},
Proposition~\ref{prop:v6v33-separator-product}, and
Theorem~\ref{thm:v6v33-gap-obstruction}, subject to both firewalls.
\end{proof}

V34 will combine the one-well spectral gap with a quantitative finite-speed
semigroup estimate.  By splitting the covariance integral before and after
the propagation time, it will derive the precise exponential clustering
rate and expose the constants that must remain uniform under RG blocking.


\clearpage
\section{V34: gap plus finite propagation gives conditional clustering}
\label{sec:v6v34}

V33 proved that a spectral gap has no spatial content by itself.  We now
add the missing local propagation estimate and derive the resulting decay
rate without hiding its dependence on the gap or propagation velocity.

Let \(P_t=e^{t\mathcal L}\) be a reversible Markov semigroup with
Dirichlet form
\begin{equation}
 \mathcal E(f,g)=\sum_x\mathbb E[d_xf\cdot d_xg].
 \label{eq:v6v34-local-form}
\end{equation}
For \(A\subseteq\Lambda\), write
\(\|d_Af\|_2^2=\sum_{x\in A}\mathbb E|d_xf|^2\).

\begin{hypothesis}[Finite propagation and energy decay]
\label{hyp:v6v34-FPED}
There are constants \(C_{\rm fp},C_{\rm en},\mu,v,\lambda>0\), uniform in
the volume under consideration, such that every centered observable \(g\)
supported in \(B\) satisfies
\begin{align}
 \|d_AP_tg\|_2
 &\le C_{\rm fp}e^{vt-\mu d(A,B)}\|g\|_{\rm end},
 \label{eq:v6v34-finite-propagation}\\
 \|dP_tg\|_2
 &\le C_{\rm en}e^{-\lambda t}\|g\|_{\rm end}.
 \label{eq:v6v34-energy-decay}
\end{align}
The second row follows from a spectral gap when the endpoint norm controls
the initial Dirichlet norm.  The first row is an independent locality
statement.
\end{hypothesis}

\begin{lemma}[Covariance integral]
\label{lem:v6v34-covariance-integral}
If \(f,g\) are centered and the semigroup converges to equilibrium, then
\begin{equation}
 \operatorname{Cov}_\nu(f,g)
 =\int_0^\infty\mathcal E(f,P_tg)\,dt.
 \label{eq:v6v34-covariance-integral}
\end{equation}
\end{lemma}

\begin{proof}
Reversibility gives
\[
 \frac d{dt}\langle f,P_tg\rangle
 =\langle f,\mathcal LP_tg\rangle
 =-\mathcal E(f,P_tg).
\]
Integrate from zero to infinity.  The terminal inner product is zero
because both variables are centered and \(P_tg\) converges to zero.
\end{proof}

\begin{theorem}[Quantitative exponential clustering]
\label{thm:v6v34-exponential-clustering}
Assume Hypothesis~\ref{hyp:v6v34-FPED}.  If centered \(f\) is supported in
\(A\) and centered \(g\) is supported in \(B\), then
\begin{equation}
 |\operatorname{Cov}_\nu(f,g)|
 \le C_*\|d_Af\|_2\|g\|_{\rm end}
 \exp[-m_*d(A,B)],
 \label{eq:v6v34-exponential-clustering}
\end{equation}
where
\begin{equation}
 m_*:=\frac{\mu\lambda}{v+\lambda},
 \qquad
 C_*:=\frac{C_{\rm fp}}v+\frac{C_{\rm en}}\lambda.
 \label{eq:v6v34-effective-mass}
\end{equation}
The first fraction is interpreted by its continuous \(v\downarrow0\)
limit, with an arbitrarily small loss in \(m_*\).
\end{theorem}

\begin{proof}
Set \(D=d(A,B)\) and
\begin{equation}
 t_*:=\frac{\mu D}{v+\lambda}.
 \label{eq:v6v34-splitting-time}
\end{equation}
Because \(d_xf=0\) for \(x\notin A\), Cauchy--Schwarz in
\eqref{eq:v6v34-covariance-integral} gives
\[
 |\mathcal E(f,P_tg)|
 \le\|d_Af\|_2\|d_AP_tg\|_2.
\]
On \([0,t_*]\), insert
\eqref{eq:v6v34-finite-propagation}; on \([t_*,\infty)\), insert
\eqref{eq:v6v34-energy-decay}.  The two integrals are at most
\begin{align*}
 \frac{C_{\rm fp}}v
 e^{-\mu D+vt_*}\|d_Af\|_2\|g\|_{\rm end},
 \qquad
 \frac{C_{\rm en}}\lambda
 e^{-\lambda t_*}\|d_Af\|_2\|g\|_{\rm end}.
\end{align*}
Both exponents equal \(-m_*D\).  Their sum proves the claim.  When
\(v=0\), the early integral contributes a factor \(t_*e^{-\mu D}\),
which is absorbed by replacing \(\mu\) with any smaller positive rate.
\end{proof}

\subsection{Conditional influence consequence}

The covariance theorem gives the conditional operator of V33 once local
test functions have a controlled upper energy.

\begin{hypothesis}[Root energy regularity]
\label{hyp:v6v34-root-energy}
For the root observable class on \(A\),
\begin{equation}
 \|d_Af\|_2\le\Lambda_A^{1/2}\|f\|_2
 \qquad
 (f\in L^2_0(\mathcal F_A)).
 \label{eq:v6v34-root-energy}
\end{equation}
At fixed lattice cutoff this is an upper spectral bound for the local
form.  In a continuum limit it must be replaced by a Sobolev endpoint
norm with uniform constants.
\end{hypothesis}

\begin{corollary}[Off-diagonal conditional projection]
\label{cor:v6v34-rho-decay}
Under Hypotheses~\ref{hyp:v6v34-FPED} and
\ref{hyp:v6v34-root-energy},
\begin{equation}
 \rho_2(A,B)
 \le C_*\Lambda_A^{1/2}e^{-m_*d(A,B)}
 \label{eq:v6v34-rho-decay}
\end{equation}
when the endpoint norm on \(B\) is the \(L^2\) norm used in the definition
of \(\rho_2\).  With a stronger endpoint norm, the same proof gives the
corresponding mixed operator bound.
\end{corollary}

\begin{proof}
By duality,
\begin{equation}
 \|P_A^\circ g\|_2
 =\sup_{\substack{f\in L^2_0(\mathcal F_A)\\\|f\|_2=1}}
 |\operatorname{Cov}(f,g)|.
 \label{eq:v6v34-projection-duality}
\end{equation}
Apply Theorem~\ref{thm:v6v34-exponential-clustering} and then
\eqref{eq:v6v34-root-energy}.
\end{proof}

For the projected energy norm in
\eqref{eq:v6v33-energy-influence}, one needs the differentiated analogue
of Corollary~\ref{cor:v6v34-rho-decay}:
\begin{equation}
 \|dP_A^\circ g\|_2
 \le C_{\rm pr}\Lambda_A^{1/2}
 e^{-m_*d(A,B)}\|g\|_{\rm end}.
 \label{eq:v6v34-projected-energy-decay}
\end{equation}
This does not follow formally from an \(L^2\) projection bound for an
unbounded continuum gradient.  It is a separate regularity row.

\begin{proposition}[Conditional route to the PLWC]
\label{prop:v6v34-conditional-PLWC}
Suppose the differentiated projection estimate
\eqref{eq:v6v34-projected-energy-decay} is stable under conditioning on
previously exposed rooted blocks, and suppose each connected cumulant
block obeys the existing IMS endpoint estimate.  Then the Appell recursion
\eqref{eq:v6v31-Appell-recursion}, applied from the leaves toward the score
root, proves the PLWC with one factorial and edge rate \(m_*\).
\end{proposition}

\begin{proof}
For each leaf block, condition its Appell remainder toward the adjacent
rooted block and apply \eqref{eq:v6v34-projected-energy-decay}; this supplies
one decaying edge before the leaf variables are summed.  Remove the leaf
and repeat.  At a branching vertex, sum the subset of labels assigned to
the connected block together with its complementary Appell remainder.
Lemma~\ref{lem:v6v32-one-factorial-convolution} bounds that subset sum by
\(C^rr!\), rather than multiplying a partition factorial by a tree
factorial.  The IMS compiler handles the resulting hybrid edge product.
Uniform stability under the preceding conditionings permits the induction
at every vertex.
\end{proof}

\subsection{RG uniformity and the closing-gap diagnostic}

At RG scale \(j\), let the four constants be
\(\mu_j,v_j,\lambda_j,C_{*,j}\).  The induced clustering rate is
\begin{equation}
 m_{*,j}=\frac{\mu_j\lambda_j}{v_j+\lambda_j}.
 \label{eq:v6v34-scale-effective-mass}
\end{equation}

\begin{lemma}[Uniform-rate criterion]
\label{lem:v6v34-uniform-rate}
If \(\inf_j\mu_j>0\), \(\inf_j\lambda_j>0\),
\(\sup_jv_j<\infty\), and the prefactors and conditional regularity
constants are uniformly bounded, then \(\inf_jm_{*,j}>0\).  If instead
\(\lambda_j\downarrow0\) while \(v_j\) stays bounded below away from zero,
then \(m_{*,j}\downarrow0\).
\end{lemma}

\begin{proof}
Both assertions follow directly from
\eqref{eq:v6v34-scale-effective-mass}.
\end{proof}

Thus a closing massless gap cannot be hidden in a massive exponential
tree edge.  Its channel must either acquire a Ward improvement strong
enough for the required power-law summation or remain explicitly in the
physical IMS skeleton.

\begin{firewall}[Finite propagation is an acquisition target]
\label{fw:v6v34-FP-open}
No uniform estimate of the form
\eqref{eq:v6v34-finite-propagation} has yet been proved for the coupled
gauge--fermion--gravity transfer generator.  In particular, quotient
charts and determinant-line transport may change the commutator constants
entering \(v\) and \(\mu\).
\end{firewall}

\begin{firewall}[Projection regularity is not a consequence of the gap]
\label{fw:v6v34-regularity-open}
The gap controls a lower energy bound.  The root regularity
\eqref{eq:v6v34-root-energy} and its differentiated conditional version
\eqref{eq:v6v34-projected-energy-decay} are upper regularity statements.
They require a local elliptic, heat-kernel, or finite-state argument and
must remain uniform as the ultraviolet cutoff is removed.
\end{firewall}

\begin{theorem}[Sixth-series V34 result]
\label{thm:v6v34-status}
A reversible gap \(\lambda\) combined with finite propagation parameters
\((\mu,v)\) gives exponential clustering at the explicit rate
\(m_*=\mu\lambda/(v+\lambda)\).  With root energy regularity this yields
decay of the conditional influence operator; with differentiated
conditional regularity it supplies the edge estimate needed by the PLWC.
The formula also proves that a closing massless gap destroys a uniform
massive edge and must be treated by Ward power counting or as a retained
physical carrier.
\end{theorem}

\begin{proof}
Combine Theorem~\ref{thm:v6v34-exponential-clustering},
Corollary~\ref{cor:v6v34-rho-decay},
Proposition~\ref{prop:v6v34-conditional-PLWC}, and
Lemma~\ref{lem:v6v34-uniform-rate}, subject to both firewalls.
\end{proof}

V35 will perform the compulsory YM/NS audit.  It will compare their newest
propagation, quotient, and negative-norm estimates with the two open rows
\eqref{eq:v6v34-finite-propagation} and
\eqref{eq:v6v34-projected-energy-decay}, then redirect the SM proof toward
the strongest common acquisition target.


\clearpage
\section{V35: compulsory companion audit and the payer-complete root card}
\label{sec:v6v35}

\subsection{Read-only YM/NS audit}

Before choosing the next SM step, the newest active companion notes were
read without modification:
\begin{center}
\begin{tabular}{p{0.59\textwidth}r r}
\hline
file & bytes & logical lines\\
\hline
\texttt{Renewal\_Yang\_Mills\_Mass\_Gap\_Research\_Notes\_V25.tex}
& \(376\,620\) & \(11\,730\)\\
\texttt{Renewal\_NS\_Stability\_Research\_Notes\_V31.tex}
& \(887\,537\) & \(23\,052\)\\
\hline
\end{tabular}
\end{center}
Their SHA--256 digests are
\begin{align}
 &\texttt{048b37af3b11bb3f0b8249423f34f5b1724309e0a807c213bf44f65528521ce8},
 \label{eq:v6v35-YM-hash}\\
 &\texttt{f1121b62238ad5491bb7cf03635ea9934942edf573ed8faaa9ee79e1d38a6696}.
 \label{eq:v6v35-NS-hash}
\end{align}

YM V25 computes the complete quartic Gaussian score graph and transfers its
nonzero arity-four coefficient to the actual \(SU(2)\) geodesic--Wilson
path.  The score supplies exactly the ordinary tree homogeneity and no
additional inverse-coupling factor for a separately estimated final payer.
Its V24 score--Wick pairing card is therefore false; the corrected target
applies the exact output payer to the complete score cumulant before taking
the operator norm.

NS V31 reconstructs unweighted critical coarea as the exact first moment of
flat dyadic probability marks.  Heat formation supplies that moment, but
with the critical source norm; restoring it from the energy-paid weak source
returns the continuation-class stock.  Thus a restoring operation performed
after a separate norm can spend exactly the gain it was meant to preserve.

\begin{firewall}[No companion conclusion is imported]
\label{fw:v6v35-no-import}
The YM payer is a Peter--Weyl mass--resolvent operation, while the NS first
moment is a Fourier coarea weight.  Neither proves a Standard Model Ward
projection, a conditional gauge estimate, a fermion determinant bound, or
an Einstein continuum limit.  The transferable result is the order of
operations: the complete correlated insertion and its final restoring
operation must be assembled before taking norms.
\end{firewall}

\subsection{A sharp upper-regularity obstruction}

V34 left the differentiated conditional estimate
\eqref{eq:v6v34-projected-energy-decay} as a separate regularity row.  The
following elementary example proves that neither a gap nor perfect locality
can make its constant uniform over an unrestricted ultraviolet observable
class.

\begin{theorem}[The gap does not bound differentiated projection]
\label{thm:v6v35-upper-regularity-obstruction}
Let \(\gamma\) be standard Gaussian measure on \(\mathbb R\), and let
\(\mathcal L=\partial_x^2-x\partial_x\) be the Ornstein--Uhlenbeck
generator.  Its spectral gap is one.  If \(h_n\) is the \(L^2(\gamma)\)-
normalized probabilists' Hermite polynomial of degree \(n\), then
\begin{equation}
 \|h_n\|_2=1,
 \qquad
 \|dh_n\|_2=\sqrt n.
 \label{eq:v6v35-Hermite-energy}
\end{equation}
For the full one-site sigma algebra \(A\),
\(P_A^\circ h_n=h_n\) for \(n\ge1\).  Consequently no estimate
\begin{equation}
 \|dP_A^\circ g\|_2\le C\|g\|_2
 \label{eq:v6v35-false-uniform-upper}
\end{equation}
can follow from the unit gap and locality with a constant uniform over
ultraviolet degree.
\end{theorem}

\begin{proof}
The Hermite eigenvalue equation is
\(-\mathcal Lh_n=nh_n\).  Gaussian integration by parts gives
\[
 \|dh_n\|_2^2=\mathcal E(h_n,h_n)
 =\langle h_n,-\mathcal Lh_n\rangle=n.
\]
Every nonconstant Hermite polynomial is centered and measurable on the sole
site, so its centered conditional projection is itself.  Letting
\(n\to\infty\) contradicts any uniform constant in
\eqref{eq:v6v35-false-uniform-upper}.
\end{proof}

The obstruction is ultraviolet rather than spatial.  It shows that the
root factor \(\Lambda_A^{1/2}\) in V34 must either be included in the
endpoint clock or cancelled by the exact final physical projection.  It
cannot be restored after a scale-free conditional norm.

\subsection{The payer-complete rooted influence}

Let \(\mathcal P_{\rm phys}\) denote the complete output operation at the
current RG step: horizontal gauge projection, Ward subtraction, selected
Schur complement or resolvent, and the prescribed canonical endpoint
normalization.  It acts on the output tensor after the score cumulant is
assembled.  Define the payer-complete rooted influence by
\begin{equation}
 \mathfrak I_{\rm pc}(A;F_1,\ldots,F_r)
 :=\mathcal P_{\rm phys}
 \left\langle
 S_A,P_A^\circ\mathfrak W_r(F_1,\ldots,F_r)
 \right\rangle.
 \label{eq:v6v35-payer-complete-influence}
\end{equation}
Using V33, the same object is
\begin{equation}
 \mathfrak I_{\rm pc}
 =\mathcal P_{\rm phys}
 \mathcal E\!left(
 RS_A,P_A^\circ\mathfrak W_r(F)
 \right).
 \label{eq:v6v35-payer-complete-energy}
\end{equation}
Both equalities are taken before an output norm.

\begin{definition}[Payer-complete rooted card]
\label{def:v6v35-PCRIC}
The PCRIC holds on a shell if
\begin{equation}
 \|\mathfrak I_{\rm pc}(A;F_1,\ldots,F_r)\|_{\rm out}
 \le a_A r!K^r
 \mathcal T_{r,A}^{\rm hyb}
 \prod_{i=1}^r\|F_i\|_{\rm end},
 \label{eq:v6v35-PCRIC}
\end{equation}
with the score activity, physical output payer, derivative restoration, and
Ward projector all included in \(a_A\) and the displayed edge weights.
No later factor growing with the cutoff or inverse coupling is permitted.
\end{definition}

\begin{proposition}[PCRIC is the scale-matched replacement for PLWC]
\label{prop:v6v35-PCRIC-CCD}
If the complete canonical charge derivative decomposes into rooted scores
with summable activities and each rooted contribution satisfies PCRIC,
then the V29 connected charge-derivative card holds.  The one-factorial
Appell leaf pruning of V32 remains valid because
\(\mathcal P_{\rm phys}\) is applied linearly to each complete rooted
contribution before its norm.
\end{proposition}

\begin{proof}
Sum \eqref{eq:v6v35-PCRIC} over rooted supports and use the mixed tree
compiler exactly as in Theorem~\ref{thm:v6v31-PLWC-CCD}.  There is no
post-composition payer, so no additional scale factor appears after the
root is summed.  At every Appell branching, combine the rooted block and
its complement by Lemma~\ref{lem:v6v32-one-factorial-convolution}; linearity
allows the complete physical output operation to remain outside the
internal subset sum but inside the final norm.
\end{proof}

\subsection{Why a generic extra score clock is false}

The YM audit suggests a model-independent Gaussian test.  Let \(X\) be
standard Gaussian and \(H_4=X^4-6X^2+3\).

\begin{lemma}[Quartic-root saturation]
\label{lem:v6v35-quartic-saturation}
For \(t_i\to0\),
\begin{equation}
 \kappa_5(e^{it_1X},e^{it_2X},e^{it_3X},e^{it_4X},H_4)
 =24t_1t_2t_3t_4+O_{\rm an}(|t|^6).
 \label{eq:v6v35-quartic-saturation}
\end{equation}
If \(t_i\asymp b\) and the score coefficient has size \(b^2\), the score
cumulant is generically of order \(b^6\), exactly the four-input tree
homogeneity.  There is no generic further factor \(b^2\).
\end{lemma}

\begin{proof}
Gaussian normal ordering gives \(H_4=:X^4:\).  At total input degree four,
the only connected contraction joins one of the four root legs to each
exponential input.  Its multiplicity is \(4!=24\); every other connected
graph contains an additional input--input contraction and has degree at
least six.  Multiplication by a score coefficient \(b^2\) gives the stated
homogeneity.
\end{proof}

\begin{corollary}[No separated large-payer argument]
\label{cor:v6v35-no-separated-payer}
A universal proof that first demands an extra score factor \(b^2\) beyond
tree homogeneity and then applies an output operator of norm
\(O(b^{-2})\) is false in the quartic Gaussian tangent model.  A valid
argument must show either that the physical payer is \(O(1)\) on the
quartic score range or that its exact action cancels the saturating output
sector before norms.
\end{corollary}

\begin{proof}
Lemma~\ref{lem:v6v35-quartic-saturation} gives a nonzero order-\(b^6\)
coefficient.  It cannot obey an order-\(b^8\) bound with a constant uniform
as \(b\downarrow0\).  A subsequent norm estimate of size \(b^{-2}\) cannot
restore the lost assertion.  Only information about the payer's action on
the actual coefficient can change this conclusion.
\end{proof}

\begin{assessment}[Audit-induced redirection]
\label{ass:v6v35-redirection}
The separate acquisition targets in V34 are too coarse when a derivative,
Ward quotient, or resolvent normalization grows with scale.  The correct
target is PCRIC: condition on the physical score support, retain the
complete Appell polynomial, apply the complete physical output operation,
and only then take the output norm.  Finite propagation and the gap remain
inputs to this card, but they are not estimated independently of the final
scale-restoring operation.
\end{assessment}

\begin{firewall}[PCRIC remains unproved]
\label{fw:v6v35-PCRIC-open}
No PCRIC has been established for a chiral Standard Model shell coupled to
Einstein variables.  In particular, the physical conditional sigma algebra,
the horizontal Ward projector, the determinant-line transport, and the
output-sector action of the Schur complement have not yet been assembled in
one finite-cutoff theorem.
\end{firewall}

\begin{theorem}[Sixth-series V35 result]
\label{thm:v6v35-status}
The compulsory audit records YM V25 and NS V31 and changes the SM
acquisition order.  A unit spectral gap does not control differentiated
conditional projection uniformly over ultraviolet degree.  A quartic
Gaussian score saturates ordinary tree homogeneity and has no spare generic
clock for a separately bounded large payer.  The scale-correct sufficient
target is therefore the payer-complete rooted card
\eqref{eq:v6v35-PCRIC}, which places conditioning, negative-norm pairing,
Ward projection, derivative restoration, and output normalization inside
one rooted estimate before norms.
\end{theorem}

\begin{proof}
Use Theorem~\ref{thm:v6v35-upper-regularity-obstruction},
Proposition~\ref{prop:v6v35-PCRIC-CCD},
Lemma~\ref{lem:v6v35-quartic-saturation}, and
Corollary~\ref{cor:v6v35-no-separated-payer}, together with the read-only
audit and its transfer firewall.
\end{proof}

V36 will construct the finite-cutoff physical conditional algebra on a
gauge region with boundary flux variables retained.  It will prove the
associated gluing and Ward-intertwining identities, then determine whether
the horizontal conditional projection commutes with the final photon and
Einstein output projectors or produces a new boundary-current term.


\clearpage
\section{V36: physical regional conditioning and boundary-sourced projection defects}
\label{sec:v6v36}

The PCRIC of V35 conditions on a score region and then applies the complete
physical output projection.  Raw link coordinates do not factor across a
gauge cut.  This note constructs the finite-graph cut with boundary edge
modes, proves exact commutation with gauge averaging, and isolates the
distinct defect caused by a background-dependent horizontal projector.

\subsection{Extended gauge cut}

Let \(\Gamma=(V,E)\) be a finite oriented graph and \(G\) a compact group
with normalized Haar measure.  A link configuration is
\(U=(U_e)_{e\in E}\in G^E\), with \(U_{\bar e}=U_e^{-1}\).  The vertex
gauge group \(\mathcal G=G^V\) acts by
\begin{equation}
 (g\cdot U)_e=g_{s(e)}U_eg_{t(e)}^{-1}.
 \label{eq:v6v36-gauge-action}
\end{equation}

Choose a vertex region \(A\subset V\).  Split every edge crossing its
boundary at a new vertex \(b_e\), and write
\begin{equation}
 U_e=u_e^{A}u_e^{A^c}
 \label{eq:v6v36-cut-product}
\end{equation}
after orienting the crossing edge from \(A\) to \(A^c\).  The boundary edge
group \(\mathcal H_\partial=G^{E_\partial}\) acts by
\begin{equation}
 (u_e^A,u_e^{A^c})longmapsto
 (u_e^Ah_e^{-1},h_eu_e^{A^c}).
 \label{eq:v6v36-boundary-edge-action}
\end{equation}
The glued variable \(U_e\) is invariant under this action.

\begin{theorem}[Exact compact-group gluing]
\label{thm:v6v36-compact-gluing}
For each cut edge, multiplication induces a measure-preserving bijection
of quotient probability spaces
\begin{equation}
 (G\times G)/G_{\rm diag}\simeq G,
 \qquad [(a,b)]\longmapsto ab,
 \label{eq:v6v36-one-edge-quotient}
\end{equation}
where \(h\in G_{\rm diag}\) acts as \((a,b)\mapsto(ah^{-1},hb)\).
Consequently
\begin{equation}
 L^2(G^E)
 \simeq
 \left(
 L^2(G^{E_A^{\rm ext}})\widehat\otimes
 L^2(G^{E_{A^c}^{\rm ext}})
 \right)^{\mathcal H_\partial},
 \label{eq:v6v36-Hilbert-gluing}
\end{equation}
before imposing the remaining vertex gauge constraints.  Imposing the
interior vertex groups on both factors and then the boundary edge group is
equivalent to imposing the original vertex gauge group on the glued graph.
\end{theorem}

\begin{proof}
The product \(ab\) is constant on diagonal orbits.  If \(ab=a'b'\), then
with \(h=a'^{-1}a\) one has \(a'=ah^{-1}\) and \(b'=hb\), so the induced
map is bijective.  For every integrable \(F\), Haar invariance gives
\begin{equation}
 \int_{G\times G}F(ab)\,da\,db
 =\int_G F(c)\,dc.
 \label{eq:v6v36-Haar-product}
\end{equation}
Thus the quotient pushforward is Haar measure and pullback is an isometry
onto the diagonal-invariant subspace.  Tensor this construction over the
cut edges and leave uncut edges unchanged to obtain
\eqref{eq:v6v36-Hilbert-gluing}.  Interior gauge actions commute with the
boundary edge actions and reconstruct the vertex transformations in
\eqref{eq:v6v36-gauge-action}; taking the commuting invariant subspaces in
either order gives the last assertion.
\end{proof}

The boundary variables in the extended factors are not disposable gauge
coordinates.  Their left/right representation labels carry the electric
flux needed to glue Gauss-law sectors.

\begin{corollary}[Peter--Weyl boundary-sector decomposition]
\label{cor:v6v36-boundary-PW}
For one cut edge,
\begin{equation}
 L^2(G)\simeq\widehat\bigoplus_{\pi\in\widehat G}
 V_\pi\otimes V_\pi^*.
 \label{eq:v6v36-PW-edge}
\end{equation}
Across a general cut, gluing contracts equal boundary representation labels
and their matrix indices to the trivial \(\mathcal H_\partial\) sector.
Hence a physical regional algebra must retain the boundary representation
and flux data until the two sides are glued.
\end{corollary}

\begin{proof}
This is the Peter--Weyl theorem on every boundary copy of \(G\).  Haar
averaging the diagonal boundary action is the orthogonal projection onto
its invariant tensors, which pairs a representation with its dual and
contracts the corresponding indices.
\end{proof}

\subsection{Gauge-covariant conditional expectation}

Let \(\nu\) be a gauge-invariant finite-cutoff probability law, including
any retained matter variables, and let \(\mathcal F_A^{\rm ext}\) be the
sigma algebra generated by the extended \(A\)-variables and their boundary
edge modes.  It is invariant under interior and boundary gauge actions.
Write \(P_A^{\rm ext}\) for conditional expectation onto this algebra and
\begin{equation}
 \mathsf Q_{\rm gauge}:=int_{\mathcal G\times\mathcal H_\partial}
 \mathsf U_g\,dg
 \label{eq:v6v36-gauge-average}
\end{equation}
for orthogonal gauge averaging.

\begin{theorem}[Ward intertwining of physical conditioning]
\label{thm:v6v36-Ward-intertwining}
The extended conditional expectation commutes with every gauge unitary and
with gauge averaging:
\begin{equation}
 P_A^{\rm ext}\mathsf U_g=\mathsf U_gP_A^{\rm ext},
 \qquad
 P_A^{\rm ext}\mathsf Q_{\rm gauge}
 =\mathsf Q_{\rm gauge}P_A^{\rm ext}.
 \label{eq:v6v36-Ward-intertwining}
\end{equation}
For a strongly continuous one-parameter gauge subgroup with generator
\(D_\xi\),
\begin{equation}
 D_\xi P_A^{\rm ext}F=P_A^{\rm ext}D_\xi F
 \label{eq:v6v36-infinitesimal-Ward}
\end{equation}
on the common generator domain.  In particular, physical inputs remain
physical after conditioning.
\end{theorem}

\begin{proof}
Because \(\nu\) and \(\mathcal F_A^{\rm ext}\) are invariant, the unitary
image of the closed subspace
\(L^2(\mathcal F_A^{\rm ext})\) is itself.  Its orthogonal projector
therefore commutes with every \(\mathsf U_g\).  Haar integration gives the
second equality.  Differentiate the unitary intertwining identity at the
identity on the common generator domain to obtain
\eqref{eq:v6v36-infinitesimal-Ward}.
\end{proof}

\begin{assessment}[Why boundary modes are compulsory]
\label{ass:v6v36-boundary-modes}
If crossing links are assigned wholly to one side and their boundary action
is discarded, the resulting raw coordinate sigma algebra is not invariant
under gauge transformations at the cut.  The proof of
Theorem~\ref{thm:v6v36-Ward-intertwining} then fails, and conditioning can
create a spurious Ward source.  The extended boundary modes are the minimal
data that make regional conditioning gauge covariant.
\end{assessment}

\subsection{Background-dependent horizontal projection}

Gauge averaging and horizontal output projection are different operations.
Let \(\mathscr H_{\rm out}\) be the finite-cutoff output space, and let
\(Q(U)\) be a measurable orthogonal projector on it, such as a
background-dependent photon or Einstein horizontal projector.  For an
\(\mathscr H_{\rm out}\)-valued random variable \(F\), define the
operator-valued conditional mean
\begin{equation}
 \overline Q_A:=P_A^{\rm ext}Q.
 \label{eq:v6v36-conditional-projector-mean}
\end{equation}

\begin{theorem}[Exact conditional projector defect]
\label{thm:v6v36-projector-defect}
One has
\begin{equation}
 \boxed{\quad
 P_A^{\rm ext}(QF)
 -\overline Q_A,P_A^{\rm ext}F
 =P_A^{\rm ext}\!left[
 (Q-\overline Q_A)(F-P_A^{\rm ext}F)
 \right].
 \quad}
 \label{eq:v6v36-projector-defect}
\end{equation}
Thus conditioning commutes with the output projector when \(Q\) is
\(\mathcal F_A^{\rm ext}\)-measurable; in general the defect is an exact
conditional covariance between projector motion and the exterior
fluctuation of the output.
\end{theorem}

\begin{proof}
Write \(Q=\overline Q_A+(Q-\overline Q_A)\) and
\(F=P_A^{\rm ext}F+(F-P_A^{\rm ext}F)\).  Conditional expectation kills
each centered factor separately.  The product of the two centered factors
is the only remainder, giving the displayed identity.
\end{proof}

The right side of \eqref{eq:v6v36-projector-defect} belongs inside the
PCRIC.  Bounding \(Q-\overline Q_A\) and
\(F-P_A^{\rm ext}F\) separately would again discard their correlated
boundary cancellation.

\subsection{Boundary source for the Hodge commutator}

There is also a deterministic commutator caused by cutting a horizontal
field.  Let \(K:\mathscr H_0\to\mathscr H_1\) be the linearized gauge
differential, let \(M=K^*K\) be invertible on the gauge-orthogonal scalar
space, and set
\begin{equation}
 Q_H:=I-KM^{-1}K^*.
 \label{eq:v6v36-horizontal-projector}
\end{equation}
Let \(C_A\) be multiplication by a regional cutoff on one-forms.

\begin{proposition}[Exact boundary commutator]
\label{prop:v6v36-boundary-commutator}
On the common domain,
\begin{align}
 [C_A,Q_H]
 ={}&-[C_A,K]M^{-1}K^*
      -K[C_A,M^{-1}]K^*
      -KM^{-1}[C_A,K^*],
 \label{eq:v6v36-boundary-commutator}\\
 [C_A,M^{-1}]
 ={}&-M^{-1}[C_A,M]M^{-1}.
 \label{eq:v6v36-resolvent-commutator}
\end{align}
For a local lattice differential, every primitive commutator
\([C_A,K]\), \([C_A,K^*]\), and \([C_A,M]\) is supported on links or
vertices meeting the cut.  The defect is therefore boundary-sourced and
then propagated by the gauge Green operator \(M^{-1}\).
\end{proposition}

\begin{proof}
Expand \([C_A,KM^{-1}K^*]\) by the Leibniz rule for commutators.  The
inverse identity follows from
\([C_A,MM^{-1}]=0\) after multiplication by \(M^{-1}\).  A local
difference operator commutes with a constant regional multiplier away from
edges on which that multiplier changes, proving the support statement.
\end{proof}

\begin{lemma}[The regional and horizontal projectors need not commute]
\label{lem:v6v36-noncommuting-example}
On \(\mathbb R^2\), let the gauge direction be
\(K1=(1,1)\), let \(Q_H\) project onto \({\rm span}(1,-1)\), and let
\(C_A={\rm diag}(1,0)\).  Then
\begin{equation}
 [C_A,Q_H]=\frac12
 \begin{pmatrix}0&-1\\1&0\end{pmatrix}\ne0.
 \label{eq:v6v36-two-edge-commutator}
\end{equation}
\end{lemma}

\begin{proof}
Here
\(Q_H=\frac12\left(\begin{smallmatrix}1&-1\\-1&1\end{smallmatrix}\right)\).
Direct matrix multiplication gives the displayed commutator.
\end{proof}

The same algebra applies to the linearized Einstein projector after
replacing \(K\) by the infinitesimal diffeomorphism operator and \(M\) by
its gauge-fixed normal operator.  Uniform estimates still depend on the
corresponding Green kernel and boundary trace norms.

\begin{firewall}[Gluing is finite-cutoff kinematics]
\label{fw:v6v36-finite-cutoff}
Theorem~\ref{thm:v6v36-compact-gluing} and the Ward intertwining theorem are
exact on a finite compact lattice.  They do not construct the interacting
chiral fermion measure, control an anomaly, remove the lattice cutoff, or
establish an Einstein measure.
\end{firewall}

\begin{firewall}[The boundary defect still needs a rooted estimate]
\label{fw:v6v36-boundary-estimate-open}
Equations \eqref{eq:v6v36-projector-defect} and
\eqref{eq:v6v36-boundary-commutator} identify the defect but do not bound
it.  A PCRIC proof must show that its boundary source carries a summable
massive kernel or the correct Ward-dressed physical massless edge, without
applying a Green-operator norm after the correlated product has been split.
\end{firewall}

\begin{theorem}[Sixth-series V36 result]
\label{thm:v6v36-status}
Finite-cutoff regional gauge conditioning is well defined after every cut
link is split and its boundary representation/flux mode is retained.
Conditional expectation then intertwines gauge averaging and the
infinitesimal Ward generators exactly.  A background-dependent horizontal
output projector has the exact conditional covariance defect
\eqref{eq:v6v36-projector-defect}, while a deterministic regional cut has
the boundary-sourced Hodge commutator
\eqref{eq:v6v36-boundary-commutator}.  Thus physical conditioning creates no
spurious gauge charge, but the photon and Einstein horizontal projectors
generically contribute a real boundary-current row to PCRIC.
\end{theorem}

\begin{proof}
Use Theorem~\ref{thm:v6v36-compact-gluing},
Corollary~\ref{cor:v6v36-boundary-PW},
Theorem~\ref{thm:v6v36-Ward-intertwining},
Theorem~\ref{thm:v6v36-projector-defect}, and
Proposition~\ref{prop:v6v36-boundary-commutator}; noncommutation is witnessed
by Lemma~\ref{lem:v6v36-noncommuting-example}.
\end{proof}

V37 will estimate the boundary-sourced Hodge defect without taking a global
Green-operator norm.  It will derive a boundary-to-bulk kernel formula,
separate massive and physical massless channels, and test the shell sum at
the exact codimension-one threshold.


\clearpage
\section{V37: boundary Green kernels and the codimension-one Ward threshold}
\label{sec:v6v37}

V36 expresses the failure of a regional cutoff to commute with the
horizontal gauge projector as a boundary source propagated by the gauge
Green operator.  Taking a global norm of that inverse would lose both
locality and the physical charge split.  We therefore keep its kernel and
perform the boundary sum before any operator norm.

\subsection{Exact boundary-to-bulk representation}

Retain the notation
\(Q_H=I-KM^{-1}K^*\), \(M=K^*K\), and let \(C_A\) be a regional
multiplier.  Put
\begin{equation}
 B_1:=[C_A,K],\qquad
 B_0:=[C_A,M],\qquad
 B_{-1}:=[C_A,K^*].
 \label{eq:v6v37-boundary-operators}
\end{equation}
Each \(B_q\) is supported on the cut for a local lattice differential.

\begin{proposition}[Three-term boundary Green formula]
\label{prop:v6v37-three-term-formula}
On the common domain,
\begin{equation}
 [C_A,Q_H]f
 =-B_1M^{-1}K^*f
   +KM^{-1}B_0M^{-1}K^*f
   -KM^{-1}B_{-1}f.
 \label{eq:v6v37-three-term-formula}
\end{equation}
Every term is obtained by applying a boundary-supported operator before at
least one occurrence of the Green kernel \(M^{-1}\).  Thus the commutator
can be estimated as a sum of boundary-to-bulk kernels without using
\(\|M^{-1}\|\).
\end{proposition}

\begin{proof}
Insert the resolvent commutator
\eqref{eq:v6v36-resolvent-commutator} into
\eqref{eq:v6v36-boundary-commutator} and use the definitions in
\eqref{eq:v6v37-boundary-operators}.  Local support of the primitive
commutators was proved in
Proposition~\ref{prop:v6v36-boundary-commutator}.
\end{proof}

For example, if \(G=M^{-1}\) has kernel \(G(x,y)\), the last term has the
schematic form
\begin{equation}
 (KGB_{-1}f)(x)
 =\sum_{y\in\partial A}K_xG(x,y)(B_{-1}f)(y).
 \label{eq:v6v37-boundary-kernel}
\end{equation}
The middle term is a composition of two such propagations joined at a
boundary insertion.  Its two Green factors must remain convolved through
that insertion before a norm is taken.

\subsection{The codimension-one census}

Let \(d\ge3\), and assume a discrete boundary \(\Sigma\subset\mathbb Z^d\)
satisfies the uniform shell count
\begin{equation}
 \#\{y\in\Sigma:n\le |x-y|<n+1\}
 \le C_\Sigma(1+n)^{d-2}
 \label{eq:v6v37-boundary-shell-count}
\end{equation}
for every \(x\) and \(n\ge0\).

\begin{lemma}[Sharp boundary summability threshold]
\label{lem:v6v37-boundary-threshold}
If a kernel obeys
\begin{equation}
 |H_\alpha(x,y)|\le C_H(1+|x-y|)^{-\alpha},
 \label{eq:v6v37-power-kernel}
\end{equation}
then its absolute boundary shell majorant is summable at infinity when
\(\alpha>d-1\).  At \(\alpha=d-1\), the majorant is harmonic and diverges
logarithmically.  More precisely,
\begin{align}
 \sum_{\substack{y\in\Sigma\\|x-y|\ge R}}
 |H_\alpha(x,y)|
 &\le C\sum_{n\ge R}(1+n)^{d-2-\alpha},
 \label{eq:v6v37-boundary-series}\
 \sum_{n=R}^{N}(1+n)^{-1}
 &\asymp\log(N/R)qquad(1\le R<N).
 \label{eq:v6v37-boundary-log}
\end{align}
\end{lemma}

\begin{proof}
Partition \(\Sigma\) into the unit shells in
\eqref{eq:v6v37-boundary-shell-count} and use
\eqref{eq:v6v37-power-kernel}.  The resulting numerical series has exponent
\(d-2-\alpha\), so it converges exactly when that exponent is less than
\(-1\).  At equality it is the harmonic series, whose partial sums are
comparable to the displayed logarithm by the integral test.
\end{proof}

For a massless Laplace-type gauge normal operator in \(d\) dimensions,
\(G(x,y)\) has degree \(2-d\) and one differentiated Green kernel has
degree \(1-d\).  Lemma~\ref{lem:v6v37-boundary-threshold} therefore places
the naive boundary Hodge current exactly at the logarithmic threshold.

\begin{corollary}[Massive boundary propagation]
\label{cor:v6v37-massive-boundary}
If
\begin{equation}
 |H_m(x,y)|\le C_H e^{-m|x-y|}(1+|x-y|)^{-(d-1)},
 \label{eq:v6v37-massive-kernel}
\end{equation}
then its boundary sum is finite for every \(m>0\).  For small \(m\), the
far-field majorant can grow like \(1+|\log m|\); hence fixed positivity of
the mass, or an additional cancellation, is needed for a scale-uniform
constant as \(m\downarrow0\).
\end{corollary}

\begin{proof}
The shell sum is bounded by
\(C\sum_{n\ge0}e^{-mn}/(1+n)\).  This converges for \(m>0\) and is bounded
by \(C(1+|\log m|)\) by splitting at \(n\simeq m^{-1}\).
\end{proof}

\subsection{Ward neutralization gains the missing power}

Partition \(\Sigma\) into finite boundary blocks \(D\), choose a base point
\(y_D\in D\), and let \(\sigma_D(y)\) be a scalar boundary source with
\begin{equation}
 \sum_{y\in D}\sigma_D(y)=0.
 \label{eq:v6v37-neutral-source}
\end{equation}

\begin{theorem}[Neutral boundary block gains one derivative]
\label{thm:v6v37-neutral-gain}
Suppose \(H(x,y)\) is differentiable in its second variable along lattice
edges and
\begin{equation}
 |\nabla_yH(x,y)|
 \le C_H(1+|x-y|)^{-d}
 \label{eq:v6v37-kernel-gradient}
\end{equation}
whenever \(d(x,D)\ge2\operatorname{diam}D\).  Then
\begin{equation}
 \left|
 \sum_{y\in D}H(x,y)\sigma_D(y)
 \right|
 \le C_H\operatorname{diam}(D)
 (1+d(x,D))^{-d}
 \sum_{y\in D}|\sigma_D(y)|.
 \label{eq:v6v37-neutral-gain}
\end{equation}
The resulting far boundary majorant is absolutely summable by
Lemma~\ref{lem:v6v37-boundary-threshold}.
\end{theorem}

\begin{proof}
Use \eqref{eq:v6v37-neutral-source} to subtract a constant kernel value:
\[
 \sum_{y\in D}H(x,y)\sigma_D(y)
 =\sum_{y\in D}[H(x,y)-H(x,y_D)]\sigma_D(y).
\]
Join \(y_D\) to \(y\) by a path in a fixed enlargement of \(D\) and
sum the discrete gradient of \(H\) along that path.  The distance
assumption keeps every intermediate point comparable to \(d(x,D)\), giving
\[
 |H(x,y)-H(x,y_D)|
 \le C_H\operatorname{diam}(D)(1+d(x,D))^{-d}.
\]
Sum against \(|\sigma_D(y)|\).  Since \(d>d-1\), boundary summability is
Lemma~\ref{lem:v6v37-boundary-threshold}.
\end{proof}

For a nonabelian linearized source, choose parallel transports
\(\tau_{y\to y_D}\) inside the boundary block and replace
\eqref{eq:v6v37-neutral-source} by the covariant neutrality condition
\begin{equation}
 \sum_{y\in D}\tau_{y\to y_D}\sigma_D(y)=0.
 \label{eq:v6v37-covariant-neutrality}
\end{equation}
The same proof applies to the covariantly transported kernel, provided its
first covariant difference satisfies \eqref{eq:v6v37-kernel-gradient}.

\subsection{Physical charge split}

Let \(\Pi_D^0\) denote the blockwise covariant constant mode and
\(\Pi_D^\perp=I-\Pi_D^0\).  Decompose every boundary current as
\begin{equation}
 \sigma_D=\sigma_D^{\rm ch}+\sigma_D^{\rm neu},
 \qquad
 \sigma_D^{\rm ch}:=\Pi_D^0\sigma_D,
 \qquad
 \sigma_D^{\rm neu}:=\Pi_D^\perp\sigma_D.
 \label{eq:v6v37-charge-neutral-split}
\end{equation}

\begin{proposition}[Only the neutral row is automatically summable]
\label{prop:v6v37-charge-split}
Under the massless kernel hypotheses above,
\(\sigma_D^{\rm neu}\) gains the summable power in
\eqref{eq:v6v37-neutral-gain}.  A nonzero charge mode
\(\sigma_D^{\rm ch}\) retains the degree \(1-d\) boundary kernel and is at
the logarithmic threshold.  It may be removed only if an exact Gauss/Ward
identity proves that its coefficient is zero.  Otherwise it must remain as
a physical massless edge in the IMS skeleton.
\end{proposition}

\begin{proof}
The neutral projection satisfies the covariant zero-sum identity, so apply
Theorem~\ref{thm:v6v37-neutral-gain}.  The constant projection has nonzero
block sum in general, so the subtraction used in that proof is unavailable.
Lemma~\ref{lem:v6v37-boundary-threshold} then gives the logarithmic
threshold for a degree \(1-d\) kernel.  Gauss law identifies this constant
boundary flux with enclosed physical charge; absent its vanishing, deleting
it would delete a physical channel.
\end{proof}

The gravitational analogue splits boundary momentum and Hamiltonian
constraint charges from neutral diffeomorphism modes.  The latter can gain
a derivative after the linearized Bianchi identity; ADM-type charge modes
must be retained unless the chosen sector fixes them to zero.

\subsection{A boundary-rooted PCRIC reduction}

Let \(a_D^{\rm neu}\) and \(a_D^{\rm ch}\) be the activities of the two
rows in \eqref{eq:v6v37-charge-neutral-split} after the complete physical
output operation.

\begin{theorem}[Boundary defect compilation]
\label{thm:v6v37-boundary-compilation}
Assume:
\begin{enumerate}[label=\textup{(\roman*)},leftmargin=3.7em]
\item the primitive boundary commutators in
      \eqref{eq:v6v37-boundary-operators} have uniformly local endpoint
      bounds;
\item the massive Green rows obey
      \eqref{eq:v6v37-massive-kernel} with a scale-uniform positive mass;
\item every gauge-artifact massless row is covariantly neutral and obeys
      \eqref{eq:v6v37-kernel-gradient}; and
\item every nonzero physical charge row is retained as an IMS massless edge
      and satisfies its prescribed endpoint bound.
\end{enumerate}
Then the full boundary commutator in
\eqref{eq:v6v37-three-term-formula} reduces to rooted activities whose
incident kernels are either absolutely summable or explicitly retained
physical massless edges.  If the remaining score--Appell bulk pairing obeys
PCRIC, adjoining the boundary defect preserves PCRIC with one additional
distinguished boundary root and no logarithmic leakage.
\end{theorem}

\begin{proof}
Decompose every primitive boundary source block by
\eqref{eq:v6v37-charge-neutral-split}.  Massive rows are summable by
Corollary~\ref{cor:v6v37-massive-boundary}.  Neutral massless rows are
summable by Theorem~\ref{thm:v6v37-neutral-gain}.  Retain each physical
charge row rather than summing it.  Apply this decomposition inside each
Green convolution of Proposition~\ref{prop:v6v37-three-term-formula}; the
middle term is compiled one boundary insertion at a time.  The mixed tree
compiler then attaches the resulting distinguished boundary root.  Since
the decomposition is made before the output norm, the PCRIC order of
operations is preserved.
\end{proof}

\begin{firewall}[Ward neutrality is a sector statement]
\label{fw:v6v37-neutrality-sector}
Covariant neutrality cannot be imposed on a block carrying electric,
hypercharge, color, energy, or momentum flux.  It must follow from the
actual Ward/Bianchi identity and the selected physical sector.  Declaring a
physical charge block neutral would manufacture the missing power by
removing a real massless channel.
\end{firewall}

\begin{firewall}[Kernel bounds remain open for the coupled background]
\label{fw:v6v37-kernel-open}
The power census is exact once the Green kernel estimates are assumed.  No
uniform covariant bound \eqref{eq:v6v37-kernel-gradient} has yet been proved
for the random coupled SM--Einstein background, and the middle two-Green
term may require a weighted convolution estimate near rough boundaries.
\end{firewall}

\begin{theorem}[Sixth-series V37 result]
\label{thm:v6v37-status}
The horizontal regional commutator is a sum of boundary-supported sources
propagated by one or two gauge Green kernels.  On a codimension-one boundary,
a massless differentiated Green kernel of degree \(1-d\) is exactly
logarithmically nonsummable.  Covariant Ward neutrality permits subtraction
of the block monopole and gains one derivative, producing the summable
degree \(-d\).  Nonzero physical charge modes cannot receive this gain and
must remain as explicit massless IMS edges.  Under uniform kernel bounds,
this split compiles the boundary defect into PCRIC without logarithmic
leakage.
\end{theorem}

\begin{proof}
Use Proposition~\ref{prop:v6v37-three-term-formula},
Lemma~\ref{lem:v6v37-boundary-threshold},
Theorem~\ref{thm:v6v37-neutral-gain},
Proposition~\ref{prop:v6v37-charge-split}, and
Theorem~\ref{thm:v6v37-boundary-compilation}, subject to both firewalls.
\end{proof}

V38 will turn to the retained chiral fermions.  It will formulate regional
conditioning at the level of determinant lines, prove the finite-matrix
Schur gluing identity and its phase cocycle, and identify the precise
anomaly-cancellation condition needed for the boundary-rooted Ward
intertwiner.


\clearpage
\section{V38: chiral Schur gluing, determinant cocycles, and SM anomaly arithmetic}
\label{sec:v6v38}

The gauge cut of V36 is a statement about positive compact variables.
Retained chiral fermions instead glue through Berezin integration and a
determinant line.  This note proves the finite-matrix gluing identity,
locates its boundary Green operator, and computes the local anomaly
coefficients of one Standard Model generation.  The calculation identifies
the exact remaining gap between anomaly arithmetic and a uniform
finite-cutoff determinant-line PCRIC.

\subsection{Finite-dimensional fermion gluing}

Let \(H_+=H_{+,A}\oplus H_{+,B}\) and
\(H_-=H_{-,A}\oplus H_{-,B}\) be finite-dimensional chiral spaces of equal
total dimension.  Write a chiral Dirac matrix as
\begin{equation}
 D_+=
 \begin{pmatrix}
 D_A&B_{AB}\\
 B_{BA}&D_B
 \end{pmatrix}:
 H_+\longrightarrow H_-.
 \label{eq:v6v38-block-Dirac}
\end{equation}
The off-diagonal blocks are supported at the regional cut for a local
lattice Dirac operator.

\begin{theorem}[Exact chiral Schur gluing]
\label{thm:v6v38-Schur-gluing}
If \(D_B\) is invertible, define
\begin{equation}
 S_A:=D_A-B_{AB}D_B^{-1}B_{BA}.
 \label{eq:v6v38-Schur-complement}
\end{equation}
Then
\begin{equation}
 \det D_+=\det D_B\det S_A.
 \label{eq:v6v38-determinant-gluing}
\end{equation}
With one fixed Berezin orientation convention, integrating the exterior
fermions gives
\begin{align}
 &\int d\overline\psi_B,d\psi_B
 \exp[-\overline\psi D_+\psi]
 \notag\\
 &\qquad=
 \det D_B\,
 \exp[-\overline\psi_A S_A\psi_A].
 \label{eq:v6v38-Berezin-gluing}
\end{align}
Thus the exterior determinant and the boundary-to-boundary propagator
\(B_{AB}D_B^{-1}B_{BA}\) are one gluing object and must be estimated before
taking separate norms.
\end{theorem}

\begin{proof}
Multiply \(D_+\) on the left by the unit-determinant block matrix
\[
 L=\begin{pmatrix}I&-B_{AB}D_B^{-1}\\0&I\end{pmatrix}.
\]
Then
\[
 LD_+=
 \begin{pmatrix}S_A&0\\B_{BA}&D_B\end{pmatrix},
\]
whose determinant is \(\det S_A\det D_B\).  This proves
\eqref{eq:v6v38-determinant-gluing}.  Complete the square by the Grassmann
translations
\[
 \psi_B\mapsto\psi_B+D_B^{-1}B_{BA}\psi_A,
 \qquad
 \overline\psi_B\mapsto
 \overline\psi_B+\overline\psi_AB_{AB}D_B^{-1}.
\]
Berezin translation invariance separates a Gaussian exterior integral
equal to \(\det D_B\) and leaves the quadratic form
\(\overline\psi_AS_A\psi_A\), proving
\eqref{eq:v6v38-Berezin-gluing}.
\end{proof}

\begin{proposition}[Gauge covariance of the Schur complement]
\label{prop:v6v38-Schur-covariance}
Suppose a gauge transformation respects the regional decomposition and
acts by
\begin{equation}
 D_+(gU)=G_-(g)D_+(U)G_+(g)^{-1},
 \label{eq:v6v38-Dirac-covariance}
\end{equation}
with block diagonal \(G_\pm=G_{\pm,A}\oplus G_{\pm,B}\).  Then
\begin{equation}
 S_A(gU)=G_{-,A}(g)S_A(U)G_{+,A}(g)^{-1}.
 \label{eq:v6v38-Schur-covariance}
\end{equation}
The exterior inverse appearing in the Schur complement is therefore a
gauge-covariant boundary propagator.
\end{proposition}

\begin{proof}
Transform all four blocks in \eqref{eq:v6v38-block-Dirac}.  The inverse
transforms as
\(D_B(gU)^{-1}=G_{+,B}D_B(U)^{-1}G_{-,B}^{-1}\).  The intermediate
exterior factors cancel in
\(B_{AB}D_B^{-1}B_{BA}\), leaving exactly
\eqref{eq:v6v38-Schur-covariance}.
\end{proof}

\subsection{Determinant line and phase cocycle}

For a chiral map \(D_+:H_+\to H_-\), its determinant belongs naturally to
\begin{equation}
 \mathcal L_D:=\det H_-\otimes(\det H_+)^*.
 \label{eq:v6v38-determinant-line}
\end{equation}
Choosing bases trivializes this line and represents the determinant by a
complex number; changing the trivialization changes its phase.

On a patch where the determinant is nonzero, define the gauge phase
\begin{equation}
 c(g;U):=\frac{\det D_+(gU)}{\det D_+(U)}.
 \label{eq:v6v38-phase-cocycle}
\end{equation}

\begin{lemma}[Exact cocycle identity]
\label{lem:v6v38-cocycle}
For composable gauge transformations,
\begin{equation}
 c(g_1g_2;U)=c(g_1;g_2U)c(g_2;U).
 \label{eq:v6v38-cocycle}
\end{equation}
A gauge-invariant scalar fermion measure exists on the patch exactly when
the product cocycle of all chiral species is a coboundary removable by the
allowed local counterterm/trivialization.
\end{lemma}

\begin{proof}
Insert \(\det D_+(g_2U)\) between the numerator and denominator of
\eqref{eq:v6v38-phase-cocycle}.  A change of determinant-line
trivialization by a nonzero function \(b(U)\) replaces the cocycle by
\(b(gU)c(g;U)b(U)^{-1}\).  It becomes one precisely when the original
cocycle is the corresponding coboundary.
\end{proof}

The determinant identity in Theorem~\ref{thm:v6v38-Schur-gluing} is an
identity of determinant lines.  Treating \(\det D_B\) and \(\det S_A\) as
independently phased scalar weights before gluing can violate
\eqref{eq:v6v38-cocycle} even though their tensor product is well defined.

\subsection{One-generation Standard Model anomaly coefficients}

Write every fermion as a left-handed Weyl field.  One generation, without
a sterile right-handed neutrino, is
\begin{equation}
 \begin{array}{c|c}
 \text{field}&SU(3)_c\times SU(2)_L\times U(1)_Y\\
 \hline
 Q&(\mathbf3,\mathbf2)_{1/6}\\
 u^c&(\overline{\mathbf3},\mathbf1)_{-2/3}\\
 d^c&(\overline{\mathbf3},\mathbf1)_{1/3}\\
 L&(\mathbf1,\mathbf2)_{-1/2}\\
 e^c&(\mathbf1,\mathbf1)_{1}.
 \end{array}
 \label{eq:v6v38-SM-generation}
\end{equation}

\begin{theorem}[Perturbative anomaly arithmetic]
\label{thm:v6v38-SM-anomaly-cancellation}
For the multiplets in \eqref{eq:v6v38-SM-generation}, the coefficients of
the local
\(SU(3)^3\), \(SU(3)^2U(1)\), \(SU(2)^2U(1)\), \(U(1)^3\), and
mixed gravitational--\(U(1)\) anomalies all vanish.  The number of
left-handed \(SU(2)\) doublets is four, so the ordinary mod-two Witten
anomaly is absent for one generation.
\end{theorem}

\begin{proof}
Normalize the quadratic index of a color or weak fundamental to
\(T(\mathbf3)=T(\mathbf2)=1/2\), and the cubic color anomaly so that a
fundamental and antifundamental contribute opposite signs.  For
\(SU(3)^3\), the weak doublet \(Q\) gives two color fundamentals while
\(u^c,d^c\) give two antifundamentals, hence \(2-1-1=0\).

The \(SU(3)^2U(1)\) coefficient is
\begin{equation}
 2\left(\frac16\right)\frac12
 +\left(-\frac23\right)\frac12
 +\left(\frac13\right)\frac12=0.
 \label{eq:v6v38-SU3SU3Y}
\end{equation}
The \(SU(2)^2U(1)\) coefficient is
\begin{equation}
 3\left(\frac16\right)\frac12
 +\left(-\frac12\right)\frac12=0.
 \label{eq:v6v38-SU2SU2Y}
\end{equation}
The cubic hypercharge coefficient, including color and weak
multiplicities, is
\begin{align}
 &6\left(\frac16\right)^3
 +3\left(-\frac23\right)^3
 +3\left(\frac13\right)^3
 +2\left(-\frac12\right)^3+1^3
 \notag\\
 &\qquad=\frac1{36}-\frac89+\frac19-\frac14+1=0.
 \label{eq:v6v38-Y-cubic}
\end{align}
The mixed gravitational coefficient is the hypercharge trace
\begin{equation}
 6\left(\frac16\right)
 +3\left(-\frac23\right)
 +3\left(\frac13\right)
 +2\left(-\frac12\right)+1=0.
 \label{eq:v6v38-grav-Y}
\end{equation}
There is no perturbative \(SU(2)^3\) coefficient because its representations
are pseudoreal.  Counting color copies, \(Q\) supplies three weak doublets
and \(L\) one, for an even total of four; this removes the standard mod-two
global \(SU(2)\) anomaly.
\end{proof}

Adding a sterile left-handed conjugate neutrino
\((\mathbf1,\mathbf1)_0\) changes none of these coefficients.

\begin{corollary}[Local determinant-line curvature cancels]
\label{cor:v6v38-local-curvature}
In any regulator whose determinant-line curvature is given by the standard
four-dimensional local anomaly polynomial, the sum of gauge and mixed
gauge--gravitational curvature contributions from one Standard Model
generation vanishes.
\end{corollary}

\begin{proof}
The degree-six anomaly polynomial is a linear combination of the group
invariants computed in
Theorem~\ref{thm:v6v38-SM-anomaly-cancellation}.  Each coefficient is zero.
\end{proof}

\subsection{Determinant-complete regional card}

Let \(\mathfrak D_B:=\det D_B\) and let all boundary fermion output legs be
propagated with \(S_A^{-1}\) or the prescribed Schur output before taking a
norm.

\begin{definition}[Determinant-complete fermion root]
\label{def:v6v38-DCF-root}
The regional fermion contribution is determinant-complete if it is kept in
the form
\begin{equation}
 \mathfrak I_A^{\rm f}
 :=\mathcal P_{\rm phys}
 \left[
 \mathfrak D_B\,
 \mathcal B_A(S_A;F_1,\ldots,F_r,S_A^{\rm score})
 \right],
 \label{eq:v6v38-DCF-root}
\end{equation}
where \(\mathcal B_A\) is the complete Berezin/Schur boundary tensor and
the product determinant line is trivialized only after all Standard Model
species and allowed counterterms are assembled.
\end{definition}

\begin{proposition}[Finite-matrix Ward consistency]
\label{prop:v6v38-finite-Ward-consistency}
At finite dimension, if the total determinant cocycle is trivialized and
the boundary gauge transformations respect the block decomposition, then
\eqref{eq:v6v38-DCF-root} intertwines the boundary gauge action.  No extra
Ward source is produced by Schur gluing.
\end{proposition}

\begin{proof}
Proposition~\ref{prop:v6v38-Schur-covariance} gives covariance of the
effective boundary matrix.  Theorem~\ref{thm:v6v38-Schur-gluing} identifies
the product of exterior and Schur determinants with the original
determinant line.  Triviality of the total cocycle makes that line scalar
and gauge invariant.  Applying the covariant physical output operation
therefore preserves the boundary Ward action.
\end{proof}

\begin{firewall}[Local arithmetic is not global trivialization]
\label{fw:v6v38-global-anomaly-open}
Theorem~\ref{thm:v6v38-SM-anomaly-cancellation} proves the usual local
coefficient cancellations and the ordinary \(SU(2)\) mod-two check.  It
does not construct a nonperturbative chiral lattice regulator, prove
triviality of the determinant line for every topology and every quotient
form of the Standard Model gauge group, or control eta-invariant phases on
general Einstein backgrounds.
\end{firewall}

\begin{firewall}[Zero modes require an exact-sequence gluing theorem]
\label{fw:v6v38-zero-modes}
The Schur formula assumes \(D_B\) invertible.  In a topological or massless
sector with exterior zero modes, determinants must be replaced by
determinant lines of kernels and cokernels, with boundary zero-mode legs
retained.  Introducing a small mass and dividing it out after separate
norms is not a proof of the required uniform gluing estimate.
\end{firewall}

\begin{firewall}[No positive fermion generator]
\label{fw:v6v38-no-positive-generator}
Berezin integration is not conditional expectation in a positive
probability space.  The reversible-generator PLWC argument cannot be
applied directly to the fermion variables.  Their contribution must be
integrated into determinant-complete bosonic coefficients or treated by an
operator/transfer estimate.
\end{firewall}

\begin{theorem}[Sixth-series V38 result]
\label{thm:v6v38-status}
Finite-dimensional chiral regional integration is exactly a Schur gluing:
the exterior determinant and the boundary propagator form one
gauge-covariant determinant-line object.  One Standard Model generation
cancels every standard four-dimensional perturbative gauge and mixed
gauge--gravitational anomaly coefficient and has an even number of weak
doublets.  These facts remove the local anomaly-polynomial obstruction to a
determinant-complete regional Ward intertwiner, but global determinant-line
trivialization, zero-mode gluing, and uniform Schur kernel bounds remain
open.
\end{theorem}

\begin{proof}
Use Theorem~\ref{thm:v6v38-Schur-gluing},
Proposition~\ref{prop:v6v38-Schur-covariance},
Lemma~\ref{lem:v6v38-cocycle},
Theorem~\ref{thm:v6v38-SM-anomaly-cancellation},
Corollary~\ref{cor:v6v38-local-curvature}, and
Proposition~\ref{prop:v6v38-finite-Ward-consistency}, subject to the three
firewalls.
\end{proof}

V39 will treat the exterior zero-mode problem by determinant-line exact
sequences.  It will prove the finite-dimensional pseudodeterminant gluing
formula with explicit boundary kernel/cokernel legs and determine which
massless fermion channels must remain in the IMS skeleton.


\clearpage
\section{V39: zero-mode Schur gluing and explicit massless fermion legs}
\label{sec:v6v39}

The invertible exterior block in V38 excludes topological and protected
massless sectors.  Adding a regulator mass and estimating its determinant
and Schur inverse separately is unsafe: the vanishing determinant and the
diverging inverse can cancel only after the boundary zero-mode legs are
assembled.  We now make that cancellation exact in finite dimension.

\subsection{Reduced exterior block}

Write the full chiral matrix as
\begin{equation}
 D=
 \begin{pmatrix}A&B\\C&E\end{pmatrix}:
 X_A\oplus X_B\longrightarrow Y_A\oplus Y_B,
 \label{eq:v6v39-block-D}
\end{equation}
with equal total domain and codomain dimensions.  Let
\begin{equation}
 K:=\ker E,
 \qquad
 C_0:=\ker E^*,
 \label{eq:v6v39-kernel-cokernel}
\end{equation}
and choose orthogonal decompositions
\begin{equation}
 X_B=K\oplus K^\perp,
 \qquad
 Y_B=C_0\oplus C_0^\perp.
 \label{eq:v6v39-orthogonal-split}
\end{equation}
Let \(Q_K,Q_\perp,P_0,P_\perp\) be the corresponding projections.  The
reduced exterior operator
\begin{equation}
 E_\perp:=P_\perp E Q_\perp:K^\perp\longrightarrow C_0^\perp
 \label{eq:v6v39-reduced-E}
\end{equation}
is invertible.

\begin{definition}[Zero-mode effective boundary operator]
\label{def:v6v39-zero-effective}
Define
\begin{equation}
 S_0:=
 \begin{pmatrix}
 A-BQ_\perp E_\perp^{-1}P_\perp C&BQ_K\\
 P_0C&0
 \end{pmatrix}:
 X_A\oplus K\longrightarrow Y_A\oplus C_0.
 \label{eq:v6v39-zero-effective}
\end{equation}
The columns \(BQ_K\) and rows \(P_0C\) are the retained exterior kernel
and cokernel boundary legs.
\end{definition}

\begin{theorem}[Exact zero-mode Schur gluing]
\label{thm:v6v39-zero-Schur}
Choose determinant-line orientations for
\eqref{eq:v6v39-orthogonal-split}.  Then there is a canonical determinant-
line identity, with the corresponding fixed orientation sign in scalar
coordinates,
\begin{equation}
 \boxed{\qquad
 \det D=\det E_\perp\,\det S_0.
 \qquad}
 \label{eq:v6v39-zero-Schur}
\end{equation}
In particular, no inverse is applied to \(K\), and no determinant is
assigned to the zero eigenvalues of \(E\).
\end{theorem}

\begin{proof}
Reorder the exterior domain as \(K\oplus K^\perp\) and its codomain as
\(C_0\oplus C_0^\perp\).  Regard \(X_A\oplus K\) and
\(Y_A\oplus C_0\) as the slow domain and codomain.  Relative to the slow
and fast decomposition, \(D\) has blocks
\begin{align}
 D_{00}&=
 \begin{pmatrix}A&BQ_K\\P_0C&0\end{pmatrix},
 &
 D_{01}&=
 \begin{pmatrix}BQ_\perp\\0\end{pmatrix},
 \notag\\
 D_{10}&=
 \begin{pmatrix}P_\perp C&0\end{pmatrix},
 &
 D_{11}&=E_\perp.
 \label{eq:v6v39-slow-fast-blocks}
\end{align}
The zero entries follow from \(EQ_K=0\) and
\(P_0E=0\).  Since \(D_{11}\) is invertible, ordinary block elimination
gives
\[
 \det D=\det E_\perp
 \det(D_{00}-D_{01}E_\perp^{-1}D_{10}).
\]
The second determinant is exactly \(\det S_0\).  Reordering bases produces
only the fixed sign encoded by the chosen determinant-line orientations.
\end{proof}

\begin{corollary}[Invertibility criterion with zero modes]
\label{cor:v6v39-invertibility}
The full matrix \(D\) is invertible if and only if \(S_0\) is invertible.
Thus exterior zero modes need not force a zero of the full determinant: they
may be paired through the regional boundary maps \(BQ_K\) and \(P_0C\).
If either a kernel direction is killed by \(BQ_K\) or a cokernel direction
is missed by \(P_0C\) in a way that leaves \(S_0\) singular, the full
determinant vanishes.
\end{corollary}

\begin{proof}
The block Gaussian eliminations in the proof of
Theorem~\ref{thm:v6v39-zero-Schur} are invertible, and \(E_\perp\) is
invertible.  Hence the kernel of \(D\) is isomorphic to the kernel of
\(S_0\), and the assertion follows.
\end{proof}

\subsection{Determinant-line exact sequence}

With the convention
\begin{equation}
 \operatorname{Det}(E)
 :=\det C_0\otimes(\det K)^*,
 \label{eq:v6v39-Fredholm-line}
\end{equation}
the nonzero block \(E_\perp\) canonically identifies the full matrix
determinant line with its zero-mode line and the effective boundary line.

\begin{proposition}[Zero-mode line is carried by the boundary legs]
\label{prop:v6v39-zero-line}
The decomposition used in Theorem~\ref{thm:v6v39-zero-Schur} gives a
canonical isomorphism
\begin{equation}
 \det(Y_A\oplus Y_B)\otimes
 \det(X_A\oplus X_B)^*
 \simeq
 \det(C_0^\perp)\otimes\det(K^\perp)^*
 \otimes
 \det(Y_A\oplus C_0)\otimes\det(X_A\oplus K)^*.
 \label{eq:v6v39-line-factorization}
\end{equation}
Under this isomorphism, the determinant section of \(D\) is the tensor
product of the reduced determinant section of \(E_\perp\) and the
determinant section of \(S_0\).  Gauge phases on \(K\) and \(C_0\) are
therefore transported by the boundary legs inside \(S_0\), not by an
independent scalar pseudodeterminant.
\end{proposition}

\begin{proof}
Apply multiplicativity of top exterior powers to the direct sums in
\eqref{eq:v6v39-orthogonal-split} and group the fast determinant factors
with \(E_\perp\).  The remaining factors are exactly the determinant line
of the slow map \(S_0:X_A\oplus K\to Y_A\oplus C_0\).  The block
elimination in Theorem~\ref{thm:v6v39-zero-Schur} has determinant one, so
it preserves the resulting determinant section.
\end{proof}

\subsection{The regulator cancellation cannot be separated}

\begin{lemma}[One-zero-mode saturation]
\label{lem:v6v39-regulator-saturation}
Let
\begin{equation}
 D_m=\begin{pmatrix}a&b\\c&m\end{pmatrix},
 \qquad bc\ne0.
 \label{eq:v6v39-two-by-two-regulator}
\end{equation}
For \(m\ne0\), ordinary Schur gluing gives
\begin{equation}
 \det D_m
 =m\left(a-\frac{bc}{m}\right)=am-bc.
 \label{eq:v6v39-regulator-cancellation}
\end{equation}
As \(m\to0\), the exterior determinant vanishes linearly and the Schur
complement diverges linearly, while their product tends to the nonzero
value \(-bc\).  At \(m=0\), Theorem~\ref{thm:v6v39-zero-Schur} gives the
effective zero-mode matrix
\(S_0=\left(\begin{smallmatrix}a&b\\c&0\end{smallmatrix}\right)\)
directly.
\end{lemma}

\begin{proof}
The determinant calculation is immediate.  For the zero exterior block,
\(K=C_0=\mathbb C\) and there is no nonzero exterior complement, so
Definition~\ref{def:v6v39-zero-effective} is exactly the displayed
two-by-two matrix.
\end{proof}

\begin{corollary}[No separate mass-removal bounds]
\label{cor:v6v39-no-separated-mass}
There is no uniform argument that bounds the regulated exterior determinant
and regulated Schur complement separately and then removes the mass by
absolute values.  Even in dimension two their individual scales are
\(O(m)\) and \(O(m^{-1})\), while the complete determinant has a finite
nonzero limit supplied by the zero-mode boundary legs.
\end{corollary}

\begin{proof}
Use Lemma~\ref{lem:v6v39-regulator-saturation}.  Separate absolute bounds
discard the signed product \(-bc\) that survives the limit.
\end{proof}

\subsection{Gapped and massless fermion channels}

Let the singular values of \(E_\perp\) be bounded below by
\(m_{\rm f}>0\) on a given exterior sector.  Then
\(E_\perp^{-1}\) is the gapped boundary propagator used in
\eqref{eq:v6v39-zero-effective}.  In a local lattice theory, its kernel is
expected to obey a massive off-diagonal bound.  The spaces \(K\) and
\(C_0\), by contrast, have no such inverse.

\begin{theorem}[Finite-cutoff fermion channel split]
\label{thm:v6v39-fermion-channel-split}
At finite dimension, the regional fermion gluing decomposes exactly into:
\begin{enumerate}[label=\textup{(\roman*)},leftmargin=3.7em]
\item the reduced determinant and propagator of \(E_\perp\), which may be
      compiled as a massive channel when a uniform singular-value and
      kernel bound is proved; and
\item the kernel/cokernel legs \(BQ_K\) and \(P_0C\) inside \(S_0\), which
      must remain as explicit finite-rank massless/topological output legs.
\end{enumerate}
No regulator mass is required for this split.
\end{theorem}

\begin{proof}
Theorem~\ref{thm:v6v39-zero-Schur} contains every nonzero exterior mode in
\(E_\perp\) and every zero mode in the domain or codomain of \(S_0\).
The former has an inverse; the latter does not.  These subspaces are
orthogonal and exhaustive by \eqref{eq:v6v39-orthogonal-split}, proving the
exact dichotomy.
\end{proof}

For the Standard Model, genuinely massless chiral modes before symmetry
breaking and topological zero modes after it cannot be assigned an
artificial massive tree edge.  Their boundary legs must be retained until
gauge-invariant fermion insertions, Yukawa pairings, or anomaly-cancelled
species products close them.

\subsection{Zero-mode-complete PCRIC row}

\begin{definition}[Zero-mode-complete fermion root]
\label{def:v6v39-ZMC-root}
Replace the invertible Schur object in
\eqref{eq:v6v38-DCF-root} by
\begin{equation}
 \mathfrak I_A^{\rm f,0}
 :=\mathcal P_{\rm phys}
 \left[
 \det E_\perp\,
 \mathcal B_A(S_0;F_1,\ldots,F_r,S_A^{\rm score})
 \right],
 \label{eq:v6v39-ZMC-root}
\end{equation}
where all \(K\) and \(C_0\) legs remain among the boundary outputs of
\(\mathcal B_A\).  The norm is taken only after the anomaly-cancelled
species product and physical output pairing have saturated those legs.
\end{definition}

\begin{proposition}[Conditional sufficiency]
\label{prop:v6v39-ZMC-sufficiency}
Suppose the reduced exterior kernels have uniform massive bounds, the
dimensions and endpoint norms of the retained zero-mode spaces have a
uniform topological-sector bound, and the complete physical boundary
pairing of all zero-mode legs obeys the one-factorial rooted tree estimate.
Then the fermion part of PCRIC follows without a regulator-mass loss.
\end{proposition}

\begin{proof}
Compile the \(E_\perp^{-1}\) rows as massive edges and leave the finite-rank
\(K,C_0\) rows in the rooted skeleton.  Apply the assumed physical pairing
before the output norm.  The determinant identity
\eqref{eq:v6v39-zero-Schur} prevents any separate vanishing or diverging
mass factor, while the Appell subset convolution of V32 supplies the stated
arity grade.
\end{proof}

\begin{firewall}[Uniform zero-mode dimension is not established]
\label{fw:v6v39-zero-dimension-open}
Topological charge can vary with the gauge background, and small singular
values can approach zero without belonging to a fixed kernel bundle.  No
uniform bound on the dimension, conditioning, or determinant-line transport
of these spaces has yet been proved over the full SM--Einstein configuration
class.
\end{firewall}

\begin{firewall}[Reduced determinant needs phase control]
\label{fw:v6v39-reduced-phase}
The scalar product in \eqref{eq:v6v39-ZMC-root} is shorthand for the
determinant-line tensor product of V38 and
Proposition~\ref{prop:v6v39-zero-line}.  Taking an absolute
pseudodeterminant before anomaly-cancelled gluing would erase the phase
cocycle and is not an admissible construction of a chiral measure.
\end{firewall}

\begin{theorem}[Sixth-series V39 result]
\label{thm:v6v39-status}
Exterior fermion zero modes admit an exact Schur gluing after the exterior
space is split into kernel, cokernel, and invertible complement.  The full
determinant is \(\det E_\perp\det S_0\); all zero-mode information appears
as explicit boundary legs in \(S_0\).  A two-dimensional example proves
that regulator determinant and inverse divergences cancel only in this
complete object.  Hence reduced nonzero modes may become massive tree
edges, while kernel and cokernel modes remain explicit physical or
topological IMS legs until their complete SM boundary pairing is formed.
\end{theorem}

\begin{proof}
Use Theorem~\ref{thm:v6v39-zero-Schur},
Proposition~\ref{prop:v6v39-zero-line},
Lemma~\ref{lem:v6v39-regulator-saturation},
Theorem~\ref{thm:v6v39-fermion-channel-split}, and
Proposition~\ref{prop:v6v39-ZMC-sufficiency}, subject to both firewalls.
\end{proof}

V40 will perform the compulsory YM/NS audit.  It will compare their newest
complete-payer and formation-before-unmarking results with the boundary and
zero-mode-complete PCRIC, then choose whether the next common bottleneck is
a covariant Green-kernel bound, determinant-line transport, or the
Einstein constraint-sector analogue.


\clearpage
\section{V40: compulsory audit, normalization correction, and compact SM heat forests}
\label{sec:v6v40}

\subsection{Read-only companion audit}

The newest active companion notes inspected before this append were
\begin{center}
\begin{tabular}{p{0.59\textwidth}r r}
\hline
file & bytes & logical lines\\
\hline
\texttt{Renewal\_Yang\_Mills\_Mass\_Gap\_Research\_Notes\_V29.tex}
& \(433\,724\) & \(13\,605\)\\
\texttt{Renewal\_NS\_Stability\_Research\_Notes\_V31.tex}
& \(887\,537\) & \(23\,052\)\\
\hline
\end{tabular}
\end{center}
Their SHA--256 digests were
\begin{align}
 &\texttt{533daa54c78a98e56657c62a0b26b15150383cad003a1e76319dd333f2dc1782},
 \label{eq:v6v40-YM-hash}\\
 &\texttt{f1121b62238ad5491bb7cf03635ea9934942edf573ed8faaa9ee79e1d38a6696}.
 \label{eq:v6v40-NS-hash}
\end{align}
Neither file was modified.

YM V26 proves that the inherited \(O(\alpha)\) first-connectivity payer
acts with that large eigenvalue on nonzero spin-one and spin-two pieces of
the first quartic score tensor.  The locally paid score card is therefore
false as an exact sector statement, not merely because a norm was taken too
early.  YM V27--V29 then resums the raw scalar score cycles, proves
noncommutative subset Casimir stability, retains Duhamel time ordering, and
derives an all-order strong heat-kernel forest card without a factorial.

NS V31 remains unchanged.  Its exact dyadic unmarking restores one first
moment, and heat formation pays that moment only at critical source
regularity.  Applying a scale-restoring operation after a locally normalized
bound therefore returns the continuation-class stock.

\begin{firewall}[Audit transfer boundary]
\label{fw:v6v40-audit-transfer}
No YM Wilson bridge or NS continuation estimate is imported.  The valid
common lessons are: keep the local comparison at its raw normalization;
apply a large first-connectivity payer only together with the numerator that
pays it; and retain noncommutative time ordering through the forest
derivative.
\end{firewall}

\subsection{Append-only correction to the V35 card}

Separate the final operations into two classes.  Let
\(\mathcal Q_{\rm kin}\) contain the norm-nonincreasing or scale-accounted
kinematic operations required to define a physical local output: gauge
averaging, horizontal/Ward projection with its boundary row, determinant-
line gluing, and fixed endpoint normalization.  Let
\(\mathcal R_{\rm join}\) denote a potentially large reduced resolvent or
first-connectivity payer used only after a genuinely connecting covariance
or numerator \(\mathcal N_{\rm conn}\) has been formed.

\begin{assessment}[PCRIC was overbundled]
\label{ass:v6v40-PCRIC-correction}
Definition~\ref{def:v6v35-PCRIC} and
Proposition~\ref{prop:v6v35-PCRIC-CCD} remain logically valid implications
from their stated hypothesis.  However, including a later large
first-connectivity payer inside every local score root makes that hypothesis
generically overstrong and can make it false, as the audited YM sector test
demonstrates.  From V40 onward PCRIC is replaced by the two-stage raw-root
and complete-junction cards below.
\end{assessment}

\begin{definition}[Covariant raw rooted card]
\label{def:v6v40-CRRC}
The CRRC is
\begin{equation}
 \left\|
 \mathcal Q_{\rm kin}
 \left\langle S_A,P_A^\circ
 \mathfrak W_r(F_1,\ldots,F_r)\right\rangle
 \right\|_{\rm out}
 \le a_A r!K^r\mathcal T_{r,A}^{\rm hyb}
 \prod_i\|F_i\|_{\rm end}.
 \label{eq:v6v40-CRRC}
\end{equation}
Every boundary, determinant, and zero-mode leg needed for covariance is
retained inside \(\mathcal Q_{\rm kin}\), but no large later junction
resolvent is applied.
\end{definition}

\begin{definition}[Complete junction card]
\label{def:v6v40-CJC}
For a first connection between rooted components, the CJC is the estimate
\begin{equation}
 \left\|
 \mathcal R_{\rm join}\mathcal N_{\rm conn}
 \sum_A\mathcal Q_{\rm kin}
 \left\langle S_A,P_A^\circ\mathfrak W_r(F)\right\rangle
 \right\|_{\rm out}
 \le r!K^r\mathcal T_r^{\rm hyb}
 \prod_i\|F_i\|_{\rm end},
 \label{eq:v6v40-CJC}
\end{equation}
where \(\mathcal R_{\rm join}\mathcal N_{\rm conn}\) is bounded as one
operator on the selected connected sector.  A norm estimate for
\(\mathcal R_{\rm join}\) alone is not used.
\end{definition}

\begin{proposition}[Two-stage sufficiency]
\label{prop:v6v40-two-stage}
If all local roots obey CRRC, the root activities are summable, and every
first connection obeys CJC, then the connected charge-derivative card of V29
holds with one factorial.  No local root is charged by the large junction
payer.
\end{proposition}

\begin{proof}
Use CRRC and the Appell convolution of V32 to compile each rooted component
at raw normalization.  At the first edge joining components, keep the
connecting numerator and apply CJC to the complete sum.  Subsequent
summable edges are handled by the mixed tree compiler.  Every label is
assigned once to a rooted block and every component pays the junction at
most once, so the subset identity
\eqref{eq:v6v32-binomial-factorial-cancel} preserves the one-factorial
grade.
\end{proof}

\subsection{An abstract time-ordered heat forest}

We now prove the transferable part of the YM heat endpoint in an abstract
operator form.  Let \(b_i>0\), and let \(V_{ij}=V_{ij}^*\) act on a finite
tensor product.  For \(0\le x_{ij}\le1\), set
\begin{equation}
 H(\boldsymbol x):=\sum_{i<j}x_{ij}V_{ij},
 \qquad
 U(\boldsymbol x):=e^{-H(\boldsymbol x)}.
 \label{eq:v6v40-operator-heat}
\end{equation}

For a labelled tree \(T\) and \(u_e\in[0,1]\), define the BKAR minimum
matrix
\begin{equation}
 x_{ij}^T(\boldsymbol u)
 :=\min_{e\in\operatorname{path}_T(i,j)}u_e.
 \label{eq:v6v40-BKAR-minimum}
\end{equation}

\begin{hypothesis}[Operator forest stability]
\label{hyp:v6v40-operator-stability}
There are \(B,c\ge0\) such that at every tree or forest minimum matrix,
\begin{equation}
 H(\boldsymbol x^F(\boldsymbol u))\ge-BmI,
 \qquad
 \|V_{ij}\|\le c b_ib_j.
 \label{eq:v6v40-operator-stability}
\end{equation}
\end{hypothesis}

\begin{lemma}[Factorial-free ordered derivative]
\label{lem:v6v40-ordered-derivative}
Under Hypothesis~\ref{hyp:v6v40-operator-stability}, for distinct edges
\(e_1,\ldots,e_k\),
\begin{equation}
 \left\|
 \partial_{e_1}\cdots\partial_{e_k}
 U(\boldsymbol x^F(\boldsymbol u))
 \right\|
 \le e^{Bm}c^k\prod_{ij\in\{e_1,\ldots,e_k\}}b_ib_j.
 \label{eq:v6v40-ordered-derivative}
\end{equation}
\end{lemma}

\begin{proof}
Iterated Duhamel differentiation gives
\begin{align*}
 \partial_{e_1}\cdots\partial_{e_k}e^{-H}
 =(-1)^k\sum_{\pi\in S_k}\int_{\Delta_k}
 e^{-\tau_0H}V_{e_{\pi(1)}}e^{-\tau_1H}\cdots
 V_{e_{\pi(k)}}e^{-\tau_kH}\,d\boldsymbol\tau,
\end{align*}
where \(\tau_j\ge0\), \(\sum_j\tau_j=1\), and
\(\operatorname{vol}\Delta_k=1/k!\).  Stability gives
\(\|e^{-\tau H}\|\le e^{\tau Bm}\), so the exponential pieces in every
ordered word contribute at most \(e^{Bm}\).  Insert the pair-row bounds.
The \(k!\) words cancel the simplex volume exactly.
\end{proof}

\begin{theorem}[Abstract strong heat-forest card]
\label{thm:v6v40-abstract-heat-forest}
Suppose an operator-valued moment family has a BKAR representation whose
connected cumulant is
\begin{equation}
 \kappa_m
 =\sum_{T\in\mathbb T_m}
 \int_{[0,1]^{m-1}}
 \partial_{E(T)}[A_0U](\boldsymbol x^T(\boldsymbol u))
 \,d\boldsymbol u,
 \label{eq:v6v40-BKAR-cumulant}
\end{equation}
where \(\|A_0\|\le1\) and \(A_0\) commutes with every pair interaction.
Under Hypothesis~\ref{hyp:v6v40-operator-stability},
\begin{align}
 \|\kappa_m\|
 &\le e^{Bm}c^{m-1}
 \sum_{T\in\mathbb T_m}\prod_i b_i^{d_i(T)}
 \notag\\
 &=e^{Bm}c^{m-1}
 \left(\prod_i b_i\right)
 \left(\sum_i b_i\right)^{m-2}.
 \label{eq:v6v40-abstract-heat-card}
\end{align}
No commutator expansion and no factorial occur.
\end{theorem}

\begin{proof}
Apply Lemma~\ref{lem:v6v40-ordered-derivative} with \(k=m-1\) to every
tree integral; the parameter cube has volume one.  Each tree edge
\(ij\) contributes \(b_ib_j\), so its product is
\(\prod_i b_i^{d_i(T)}\).  Sum the trees and use the weighted Pruefer
identity
\begin{equation}
 \sum_{T\in\mathbb T_m}\prod_i b_i^{d_i(T)}
 =\left(\prod_i b_i\right)\left(\sum_i b_i\right)^{m-2}.
 \label{eq:v6v40-weighted-Pruefer}
\end{equation}
\end{proof}

\subsection{Application to the compact Standard Model gauge product}

Let
\begin{equation}
 G_{\rm cSM}=SU(3)_c\times SU(2)_L\times U(1)_Y
 \label{eq:v6v40-compact-SM-group}
\end{equation}
at finite cutoff, with independent central heat times \(s_a>0\) for its
Lie-algebra factors.  Quotienting by the common finite central subgroup does
not change the following calculation on representations that descend to
the quotient.

For input representation \(R_i\), let \(J_{a,i}^\ell\) be Hermitian
generators and
\begin{equation}
 C_{a,i}:=\sum_\ell(J_{a,i}^\ell)^2,
 \qquad
 b_i^2:=\sum_a s_a C_{a,i}.
 \label{eq:v6v40-SM-clock}
\end{equation}
On reducible inputs, read the estimates blockwise and take the supremum of
the displayed scalar Casimirs.  Set
\begin{equation}
 V_{ij}:=2\sum_{a,\ell}s_aJ_{a,i}^\ell J_{a,j}^\ell.
 \label{eq:v6v40-SM-pair}
\end{equation}

\begin{lemma}[SM subset stability and pair row]
\label{lem:v6v40-SM-stability}
For every label subset \(I\),
\begin{equation}
 \sum_{\substack{i<j\\i,j\in I}}V_{ij}
 =\sum_a s_a\left[
 C_{a,I}^{\rm tot}-\sum_{i\in I}C_{a,i}
 \right]
 \ge-\sum_{i\in I}b_i^2I,
 \label{eq:v6v40-SM-subset-stability}
\end{equation}
and
\begin{equation}
 \|V_{ij}\|\le2b_ib_j.
 \label{eq:v6v40-SM-pair-bound}
\end{equation}
Every BKAR minimum matrix therefore satisfies
\(H\ge-\sum_i b_i^2I\).
\end{lemma}

\begin{proof}
For each compact factor,
\begin{equation}
 C_{a,I}^{\rm tot}
 :=\sum_\ell\left(\sum_{i\in I}J_{a,i}^\ell\right)^2
 =\sum_{i\in I}C_{a,i}
 +2\sum_{i<j\in I}\sum_\ell J_{a,i}^\ell J_{a,j}^\ell.
\end{equation}
The total Casimir is a sum of squares, proving the subset lower bound.
For one pair, row--column Cauchy--Schwarz gives
\[
 \left\|\sum_\ell J_{a,i}^\ell J_{a,j}^\ell\right\|
 \le\sqrt{C_{a,i}C_{a,j}}.
\]
Sum in \(a\) and apply scalar Cauchy--Schwarz to
\(\sqrt{s_aC_{a,i}}\sqrt{s_aC_{a,j}}\), proving the pair row.  Finally,
every BKAR minimum matrix is the integral of disjoint threshold-block
partitions; apply the subset bound on each block and integrate.
\end{proof}

\begin{theorem}[Raw compact-SM heat endpoint]
\label{thm:v6v40-SM-heat-endpoint}
Let \(M_I\) be the tensor moment of the product central heat kernel on
\(G_{\rm cSM}\).  Its complete cumulant obeys
\begin{equation}
 \boxed{\quad
 \|\kappa_m^{\rm heat}(R_1,\ldots,R_m)\|
 \le (2e)^m
 \left(\prod_i b_i\right)
 \left(\sum_i b_i\right)^{m-2}
 \quad}
 \label{eq:v6v40-SM-heat-endpoint}
\end{equation}
whenever \(b_i\le1\).  The constant is independent of the arity and
representation dimensions.
\end{theorem}

\begin{proof}
The heat Fourier transform on the tensor product is
\begin{equation}
 M_I=\exp\left[-\sum_as_aC_{a,I}^{\rm tot}\right].
 \label{eq:v6v40-SM-heat-moment}
\end{equation}
Separate its one-label Casimirs as the commuting contraction \(A_0\) and
interpolate the pair terms \eqref{eq:v6v40-SM-pair}.  At a partition
endpoint the exponential factors into the product of the block moments.
The multivariate BKAR forest formula followed by the partition Moebius sum
therefore gives \eqref{eq:v6v40-BKAR-cumulant}.  Lemma
\ref{lem:v6v40-SM-stability} verifies Hypothesis
\ref{hyp:v6v40-operator-stability} with \(B=1\) and \(c=2\), because
\(\sum_i b_i^2\le m\).  Apply
Theorem~\ref{thm:v6v40-abstract-heat-forest} and enlarge the harmless
constant.
\end{proof}

\begin{assessment}[What the audit acquires for SM]
\label{ass:v6v40-SM-acquisition}
The compact Standard Model gauge product now has a proved all-order raw heat
endpoint for representation-valued link inputs.  Noncommutativity is not a
factorial obstruction: the Duhamel simplex retains insertion order, while
subset Casimir positivity controls every intervening exponential.  The live
gauge gate is a raw, gapped comparison from the interacting Wilson/matter
law to this endpoint, together with the determinant-complete chiral rows of
V38--V39.
\end{assessment}

\begin{firewall}[The heat endpoint is not the interacting SM law]
\label{fw:v6v40-heat-not-SM}
Theorem~\ref{thm:v6v40-SM-heat-endpoint} concerns a finite-cutoff product
central heat kernel.  It does not include plaquette interactions, Higgs and
Yukawa terms, chiral determinants, Einstein variables, a comparison-path
gap, cutoff removal, or reconstruction of a continuum quantum field theory.
\end{firewall}

\begin{firewall}[Large junction payer remains open]
\label{fw:v6v40-junction-open}
CRRC and the heat endpoint are raw local estimates.  No CJC has yet been
proved for the SM--Einstein first-connectivity junction.  Attaching a large
reduced resolvent directly to
\eqref{eq:v6v40-SM-heat-endpoint} would repeat the normalization error
corrected in this version.
\end{firewall}

\begin{theorem}[Sixth-series V40 result]
\label{thm:v6v40-status}
The compulsory audit corrects the V35 acquisition order: local physical
roots use the covariant raw rooted card, while any large reduced resolvent
is applied only in a complete first-connectivity junction with its paying
numerator.  Independently, a time-ordered operator forest theorem shows
that subset stability and pair-row bounds yield a factorial-free strong
heat card.  Compact Casimir positivity verifies these hypotheses for the
finite-cutoff \(SU(3)\times SU(2)\times U(1)\) heat endpoint, giving the
all-order bound \eqref{eq:v6v40-SM-heat-endpoint}.  The interacting raw
comparison, chiral determinant rows, and global junction card remain open.
\end{theorem}

\begin{proof}
Use Proposition~\ref{prop:v6v40-two-stage},
Lemma~\ref{lem:v6v40-ordered-derivative},
Theorem~\ref{thm:v6v40-abstract-heat-forest},
Lemma~\ref{lem:v6v40-SM-stability}, and
Theorem~\ref{thm:v6v40-SM-heat-endpoint}, subject to the three firewalls.
\end{proof}

V41 will construct a raw compact-gauge comparison path to the heat endpoint
with no chart-boundary differentiation.  It will derive the exact score in
Fourier/Peter--Weyl variables, state the intrinsic path-gap card needed for
the score leaf, and test that card on the abelian and one-plaquette sectors
before coupling the chiral determinant.


\clearpage
\section{V41: intrinsic heat--interaction path and the global-comparison no-go}
\label{sec:v6v41}

V40 supplies a raw all-order heat endpoint for the compact Standard Model
gauge product.  To compare it with an interacting gauge law, the path score
must remain a sum of local physical insertions and must not differentiate a
principal-coordinate boundary.  A geometric interpolation of positive
densities has exactly these properties.

\subsection{Coordinate-free geometric interpolation}

Let \(X\) be a compact finite-cutoff configuration space with reference
measure \(m\).  Let \(h,t:X\to(0,\infty)\) be normalized densities, where
\(h\) is the heat endpoint and \(t\) is the target bosonic density.  Put
\begin{equation}
 \Phi:=\log t-\log h,
 \qquad
 d\nu_u:=Z_u^{-1}e^{u\Phi}h\,dm,
 \qquad 0\le u\le1.
 \label{eq:v6v41-geometric-path}
\end{equation}

\begin{theorem}[Exact intrinsic score]
\label{thm:v6v41-intrinsic-score}
For every integrable \(F\) independent of \(u\),
\begin{equation}
 \frac d{du}\mathbb E_uF
 =\mathbb E_u[F S_u],
 \qquad
 S_u:=\Phi-\mathbb E_u\Phi.
 \label{eq:v6v41-score-identity}
\end{equation}
For fixed observables \(F_1,\ldots,F_r\),
\begin{equation}
 \frac d{du}\kappa_r^{\nu_u}(F_1,\ldots,F_r)
 =\kappa_{r+1}^{\nu_u}(F_1,\ldots,F_r,S_u).
 \label{eq:v6v41-cumulant-derivative}
\end{equation}
If \(h,t\) are gauge invariant, then \(S_u\) is gauge invariant.  No
coordinate chart or chart-boundary distribution appears.
\end{theorem}

\begin{proof}
Differentiate \(Z_u^{-1}e^{u\Phi}\) under the finite measure integral.
The derivative of \(\log Z_u\) is \(\mathbb E_u\Phi\), proving the first
identity.  Apply it to the moment--cumulant partition formula.  Terms in
which the score belongs to a block by itself vanish because
\(\mathbb E_uS_u=0\); the remaining partition sum is exactly the
\((r+1)\)-point cumulant.  Gauge invariance follows because logarithms and
normalization preserve invariant scalar functions.
\end{proof}

\begin{corollary}[Raw endpoint comparison]
\label{cor:v6v41-raw-comparison}
\begin{equation}
 \kappa_r^{t}(F_1,\ldots,F_r)
 =\kappa_r^{h}(F_1,\ldots,F_r)
 +\int_0^1
 \kappa_{r+1}^{\nu_u}(F_1,\ldots,F_r,S_u)\,du.
 \label{eq:v6v41-raw-comparison}
\end{equation}
Thus the heat card of V40 and a uniform raw score card along
\(\nu_u\) imply the target raw card.  No first-connectivity payer occurs in
this identity.
\end{corollary}

\begin{proof}
Integrate \eqref{eq:v6v41-cumulant-derivative} from zero to one.
\end{proof}

\subsection{Local score decomposition on a gauge lattice}

Let \(E\) and \(P\) be the links and plaquettes of a finite lattice.  Take
the product central heat density
\begin{equation}
 h(U)=\prod_{e\in E}h_{s_e}(U_e)
 \label{eq:v6v41-product-heat-density}
\end{equation}
and a positive bosonic gauge--Higgs target
\begin{equation}
 t(U,\varphi)=Z_t^{-1}
 \exp[-S_{\rm g}(U)-S_{\rm H}(U,\varphi)],
 \label{eq:v6v41-bosonic-target}
\end{equation}
written relative to a fixed product reference density for the Higgs field.
For the moment, the chiral determinant is excluded from \(t\); its phase
and zero modes require V38--V39.

When \(S_{\rm g}\) is a plaquette sum and \(S_{\rm H}\) is a finite-range
site/link sum, the log-density difference has a local decomposition
\begin{equation}
 \Phi=\sum_{a\in\mathcal A}\phi_a,
 \label{eq:v6v41-local-log-ratio}
\end{equation}
where every support \(A_a\) has uniformly bounded lattice diameter.

\begin{proposition}[Centered local score roots]
\label{prop:v6v41-local-score-roots}
Choose deterministic weights \(w_a>0\) with \(\sum_aw_a=1\), and set
\begin{equation}
 S_{a,u}:=\phi_a-\mathbb E_u\phi_a,
 \qquad
 S_u=\sum_{a\in\mathcal A}S_{a,u}.
 \label{eq:v6v41-local-score-roots}
\end{equation}
Each \(S_{a,u}\) is centered, gauge invariant when the local term is, and
supported on \(A_a\) in the extended boundary sense of V36.  Consequently
\begin{equation}
 \kappa_{r+1}^{\nu_u}(F_1,\ldots,F_r,S_u)
 =\sum_a
 \left\langle
 S_{a,u},P_{A_a}^{\circ}\mathfrak W_r^{\nu_u}(F)
 \right\rangle.
 \label{eq:v6v41-rooted-score-sum}
\end{equation}
\end{proposition}

\begin{proof}
Equation \eqref{eq:v6v41-local-score-roots} follows by summing
\(\phi_a-\mathbb E_u\phi_a\).  The support and invariance statements follow
from the local action.  Apply the score--Wick identity of V30 to every
summand and then the rooted projection identity
\eqref{eq:v6v33-rooted-projection}.
\end{proof}

The irrelevant weights \(w_a\) are not needed in
\eqref{eq:v6v41-local-score-roots}; they become useful only when assigning
an anchor to overlapping representations of the same local action term.
No extensive global score norm is taken.

\subsection{The intrinsic path-gap card}

Let \(\mathcal E_u\) be a gauge-covariant local Dirichlet form on the
physical quotient of \(\nu_u\).

\begin{definition}[Uniform physical path gap]
\label{def:v6v41-UPG}
The UPG holds if there is \(\lambda_*>0\), independent of the volume,
lattice scale, and \(u\in[0,1]\), such that
\begin{equation}
 \operatorname{Var}_{\nu_u}(f)
 \le\lambda_*^{-1}\mathcal E_u(f,f)
 \label{eq:v6v41-UPG}
\end{equation}
for every centered physical \(f\) orthogonal to the retained massless and
topological sectors.
\end{definition}

If UPG holds, every local score has
\begin{equation}
 \|S_{a,u}\|_{-1,u}
 \le\lambda_*^{-1/2}\|S_{a,u}\|_2.
 \label{eq:v6v41-score-negative-gap}
\end{equation}
This is only the size row.  The CRRC still requires the conditional
connected pairing and the boundary/zero-mode splits of V36--V39.

\subsection{Why global density comparison is the wrong gap proof}

For bounded \(\Phi\), the Holley--Stroock comparison gives a lower gap
bound degraded by \(e^{-\operatorname{osc}\Phi}\).  Even on one abelian
plaquette this degradation is exponentially bad at weak coupling.

Let \(G=U(1)\), \(-\pi\le\theta<\pi\), and use the heat kernel with Fourier
multiplier \(e^{-sn^2/2}\):
\begin{equation}
 h_s(\theta)=\frac1{\sqrt{2\pi s}}
 \sum_{k\in\mathbb Z}
 \exp\left[-\frac{(\theta+2\pi k)^2}{2s}\right].
 \label{eq:v6v41-circle-heat}
\end{equation}
The Wilson density is
\begin{equation}
 w_\beta(\theta)=\frac{e^{\beta\cos\theta}}{2\pi I_0(\beta)}.
 \label{eq:v6v41-von-Mises}
\end{equation}
Calibrate \(s=\beta^{-1}\).

\begin{theorem}[Abelian oscillation obstruction]
\label{thm:v6v41-abelian-oscillation}
There are constants \(c,C>0\) such that, for all sufficiently large
\(\beta\),
\begin{equation}
 \operatorname{osc}_{\theta\in[-\pi,\pi)}
 \log\frac{w_\beta(\theta)}{h_{1/\beta}(\theta)}
 \ge\left(\frac{\pi^2}{2}-2\right)\beta-C
 \ge c\beta.
 \label{eq:v6v41-abelian-oscillation}
\end{equation}
Hence a global bounded-density perturbation theorem yields at best an
exponentially small comparison factor, even though both endpoint measures
are concentrated in the same quadratic well.
\end{theorem}

\begin{proof}
At \(\theta=0\), the periodic Gaussian formula gives
\begin{equation}
 h_{1/\beta}(0)
 =\sqrt{\frac\beta{2\pi}}[1+O(e^{-2\pi^2\beta})].
 \label{eq:v6v41-heat-zero}
\end{equation}
At \(\theta=\pi\), the two nearest images contribute equally and all other
images are exponentially smaller:
\begin{equation}
 h_{1/\beta}(\pi)
 =2\sqrt{\frac\beta{2\pi}}
 e^{-\beta\pi^2/2}[1+O(e^{-4\pi^2\beta})].
 \label{eq:v6v41-heat-pi}
\end{equation}
The Wilson normalization cancels in the ratio of its values, and
\begin{equation}
 \log\frac{w_\beta(\pi)}{w_\beta(0)}=-2\beta.
 \label{eq:v6v41-Wilson-ratio}
\end{equation}
Subtract the logarithm of
\eqref{eq:v6v41-heat-pi}/\eqref{eq:v6v41-heat-zero}.  This gives
\[
 \log\frac{w_\beta(\pi)/h_{1/\beta}(\pi)}
 {w_\beta(0)/h_{1/\beta}(0)}
 =\left(\frac{\pi^2}{2}-2\right)\beta-\log2+o(1).
\]
The oscillation is at least the absolute difference of these two values,
proving the claim.
\end{proof}

\begin{corollary}[Holley--Stroock cannot acquire UPG]
\label{cor:v6v41-HS-no-go}
On the calibrated abelian one-plaquette path, a gap proof based only on the
global oscillation of \(\log(w_\beta/h_{1/\beta})\) loses a factor
\(e^{-c\beta}\).  It cannot yield the scale-uniform or polynomially scaled
UPG required by the raw comparison programme.
\end{corollary}

\begin{proof}
The bounded perturbation comparison multiplies the reference Poincare
constant by at least the exponential of the potential oscillation.  Apply
Theorem~\ref{thm:v6v41-abelian-oscillation}.
\end{proof}

The obstruction is created at the antipodal tail, not near the common
quadratic well.  A valid proof must use an intrinsic one-well Hardy or
conditional block argument that pays the exponentially small tail by its
small mass rather than by a supremum density ratio.

\subsection{One-plaquette quotient test}

On a simply connected one-plaquette graph, a spanning-tree gauge fixing
reduces every gauge-invariant observable to the conjugacy class of the
plaquette holonomy.  The geometric path
\eqref{eq:v6v41-geometric-path} therefore descends to the class space, and
its score is the centered class function
\begin{equation}
 S_u(U_p)=log w(U_p)-\log h(U_p)
 -\mathbb E_u[\log w(U_p)-\log h(U_p)].
 \label{eq:v6v41-one-plaquette-score}
\end{equation}

\begin{proposition}[The one-plaquette test has no gauge-zero-mode defect]
\label{prop:v6v41-one-plaquette-quotient}
After spanning-tree gauge fixing, the physical Dirichlet form and score in
the one-plaquette sector act only on the conjugacy-class variable.  Gauge
orbit directions are absent, and no chart boundary derivative occurs.
Thus failure or success of the path gap in this sector is a genuine
one-well question rather than a gauge-volume artefact.
\end{proposition}

\begin{proof}
Tree gauge transformations set all tree links to the identity.  The sole
remaining loop variable is the plaquette holonomy, defined up to conjugacy
at the root.  Gauge-invariant functions and the quotient Dirichlet form are
therefore class functions of that variable.  Equation
\eqref{eq:v6v41-one-plaquette-score} is also a class function and was
defined from global positive densities, so no coordinate cut was used.
\end{proof}

\begin{firewall}[The path is bosonic until determinant completion]
\label{fw:v6v41-bosonic-path}
A chiral determinant can be complex or vanish.  Its logarithm cannot be
inserted into the positive geometric path before the anomaly-cancelled
determinant line and the zero-mode boundary legs of V38--V39 are assembled.
The path in this version covers only the positive compact gauge--Higgs
backbone.
\end{firewall}

\begin{firewall}[UPG is not proved]
\label{fw:v6v41-UPG-open}
The intrinsic score and its local decomposition are exact, but no
volume- and scale-uniform physical path gap has been proved for the
interacting lattice.  The abelian calculation rules out only the coarse
global density-ratio proof; it does not disprove an intrinsic one-well or
conditional block gap.
\end{firewall}

\begin{theorem}[Sixth-series V41 result]
\label{thm:v6v41-status}
The heat-to-target geometric interpolation gives an exact coordinate-free,
gauge-invariant score and a raw cumulant comparison.  For finite-range
gauge--Higgs actions the score splits into centered link, plaquette, and
site roots, so the CRRC can be applied locally.  On the calibrated abelian
one-plaquette model, the global heat/Wilson log-density ratio has oscillation
at least \(c\beta\); a Holley--Stroock comparison therefore cannot prove the
required path gap.  The next gap argument must be intrinsic and local,
paying remote compact tails by their probability mass.
\end{theorem}

\begin{proof}
Use Theorem~\ref{thm:v6v41-intrinsic-score},
Corollary~\ref{cor:v6v41-raw-comparison},
Proposition~\ref{prop:v6v41-local-score-roots},
Theorem~\ref{thm:v6v41-abelian-oscillation}, and
Proposition~\ref{prop:v6v41-one-plaquette-quotient}, subject to both
firewalls.
\end{proof}

V42 will prove a block tensorization theorem for the physical path gap from
uniform conditional one-block gaps and a strict interblock influence
matrix.  It will make the volume dependence explicit and test whether the
plaquette incidence matrix is contractive in the ultraviolet normalization
or whether a multiscale block gap is required.


\clearpage
\section{V42: block-gap tensorization and the plaquette-incidence obstruction}
\label{sec:v6v42}

V41 reduces the compact bosonic comparison to a uniform physical path gap
and rules out a global density-ratio proof.  This note gives a precise
block tensorization theorem.  It then tests the most direct link-block
influence matrix on the abelian plaquette tangent theory and finds that
absolute incidence is noncontractive by a fixed factor.  The correction is
to retain signed Hodge structure and separate the physical massless shell.

\subsection{A quantitative martingale-influence theorem}

Let \(\Lambda=\{1,\ldots,N\}\) index finite blocks.  Let
\(P_i f=\mathbb E(f\mid\mathcal F_{i^c})\) be conditional expectation with
block \(i\) resampled, and put
\begin{equation}
 v_i(f):=\|(I-P_i)f\|_2
 =\left(\mathbb E\operatorname{Var}_i(f)\right)^{1/2}.
 \label{eq:v6v42-local-variance-vector}
\end{equation}
Choose an ordering of the blocks, let \(\mathcal G_i\) reveal the first
\(i\) blocks, and define the Doob differences
\begin{equation}
 \Delta_i f:=\mathbb E(f\mid\mathcal G_i)
 -\mathbb E(f\mid\mathcal G_{i-1}),
 \qquad
 m_i(f):=\|\Delta_i f\|_2.
 \label{eq:v6v42-Doob-vector}
\end{equation}

\begin{definition}[Block martingale influence matrix]
\label{def:v6v42-BMI}
A nonnegative matrix \(C=(c_{ij})\), with zero diagonal, is a BMI matrix if
for every \(f\in L^2\),
\begin{equation}
 m_i(f)\le v_i(f)+\sum_{j=1}^Nc_{ij}m_j(f)
 \qquad(1\le i\le N).
 \label{eq:v6v42-BMI}
\end{equation}
This is an operator statement about propagation of one revealed block; the
entries are not defined by absolute action coefficients unless that bound
has actually been proved.
\end{definition}

\begin{theorem}[Block-gap tensorization]
\label{thm:v6v42-block-tensorization}
Assume \(C\) is a BMI matrix with
\begin{equation}
 \|C\|_{\ell^2\to\ell^2}\le\eta<1.
 \label{eq:v6v42-strict-influence}
\end{equation}
Assume also uniform conditional block gaps
\begin{equation}
 \mathbb E\operatorname{Var}_i(f)
 \le\lambda_0^{-1}\mathcal E_i(f,f),
 \qquad
 \mathcal E=\sum_i\mathcal E_i.
 \label{eq:v6v42-conditional-gap}
\end{equation}
Then the global law has the volume-independent Poincare inequality
\begin{equation}
 \operatorname{Var}(f)
 \le\frac1{\lambda_0(1-\eta)^2}\mathcal E(f,f).
 \label{eq:v6v42-global-gap}
\end{equation}
Thus its physical gap is at least \(\lambda_0(1-\eta)^2\).
\end{theorem}

\begin{proof}
The Doob differences are orthogonal and telescope, so
\begin{equation}
 \operatorname{Var}(f)=\sum_i m_i(f)^2=\|m(f)\|_{\ell^2}^2.
 \label{eq:v6v42-Doob-isometry}
\end{equation}
The componentwise BMI inequality gives
\(m\le v+Cm\).  Since \(C\) is nonnegative, iteration yields
\begin{equation}
 m\le\sum_{k=0}^\infty C^kv=(I-C)^{-1}v.
 \label{eq:v6v42-Neumann-influence}
\end{equation}
Taking the \(\ell^2\) norm and using
\(\|(I-C)^{-1}\|\le(1-\eta)^{-1}\) gives
\[
 \operatorname{Var}(f)
 \le(1-\eta)^{-2}\sum_i v_i(f)^2.
\]
Apply \eqref{eq:v6v42-conditional-gap} and sum in \(i\).
\end{proof}

\begin{remark}[The theorem isolates the real acquisition]
\label{rem:v6v42-BMI-acquisition}
The conditional one-block gap controls the cost of resampling a block.  The
BMI inequality controls how that resampling changes the rooted martingale.
Neither follows from the other.  A proof of
\eqref{eq:v6v42-BMI} must retain the sign, gauge quotient, and boundary
variables of V36 when an absolute influence estimate is not contractive.
\end{remark}

\subsection{Gaussian tangent verification}

For a centered Gaussian law on \(\mathbb R^N\) with positive precision
matrix \(H\), write
\begin{equation}
 H=D-R,
 \qquad
 D=\operatorname{diag}(H_{11},\ldots,H_{NN}),
 \qquad
 C_H:=D^{-1/2}RD^{-1/2}.
 \label{eq:v6v42-Gaussian-split}
\end{equation}

\begin{proposition}[Signed Gaussian block criterion]
\label{prop:v6v42-Gaussian-criterion}
If \(D>0\) and \(\|C_H\|<1\), then
\begin{equation}
 H\ge(1-\|C_H\|)D
 \label{eq:v6v42-Gaussian-coercivity}
\end{equation}
and the Gaussian Poincare constant is bounded by
\begin{equation}
 \lambda_{\min}(H)^{-1}
 \le[(1-\|C_H\|)\lambda_{\min}(D)]^{-1}.
 \label{eq:v6v42-Gaussian-gap}
\end{equation}
Replacing \(C_H\) by its entrywise absolute value is sufficient but may be
strictly weaker.
\end{proposition}

\begin{proof}
Factor
\(H=D^{1/2}(I-C_H)D^{1/2}\).  Since \(C_H\) is self-adjoint,
\(I-C_H\ge(1-\|C_H\|)I\), proving the form bound.  The Gaussian Poincare
constant is the reciprocal of the smallest precision eigenvalue.  Finally
\(\|C_H\|\le\||C_H|\|\), so entrywise absolute control is only a
sufficient majorant.
\end{proof}

This tangent criterion is the Hessian analogue of Theorem
\ref{thm:v6v42-block-tensorization}.  It already detects why raw link
blocks are poorly adapted to a gauge plaquette action.

\subsection{Exact abelian plaquette incidence}

Consider a periodic \(d\)-dimensional hypercubic lattice, \(d\ge2\), and
the quadratic abelian plaquette action
\begin{equation}
 S_2(A)=\frac\beta2\sum_p(dA)_p^2
 =\frac12\langle A,HA\rangle,
 \qquad H=\beta d^*d.
 \label{eq:v6v42-abelian-action}
\end{equation}
Every oriented link belongs to
\begin{equation}
 n_p=2(d-1)
 \label{eq:v6v42-link-plaquette-degree}
\end{equation}
plaquettes.  In each such plaquette it couples to the other three boundary
links with coefficient of magnitude \(\beta\).

\begin{theorem}[Absolute link incidence is noncontractive]
\label{thm:v6v42-incidence-no-go}
For the link-coordinate Hessian of \eqref{eq:v6v42-abelian-action},
\begin{equation}
 H_{ee}=\beta n_p,
 \qquad
 \sum_{e'\ne e}|H_{ee'}|=3\beta n_p
 \label{eq:v6v42-absolute-row}
\end{equation}
on a lattice large enough to avoid degenerate identifications.  Hence the
entrywise absolute, diagonally normalized influence matrix has row sum
\begin{equation}
 \sum_{e'\ne e}
 \frac{|H_{ee'}|}{\sqrt{H_{ee}H_{e'e'}}}=3.
 \label{eq:v6v42-row-sum-three}
\end{equation}
It is not a strict Dobrushin/BMI contraction at any value of \(\beta\).
\end{theorem}

\begin{proof}
Expanding one plaquette square gives diagonal coefficient \(\beta\) for
each of its four links and three off-diagonal coefficients of magnitude
\(\beta\) in the Hessian row of a fixed link.  A fixed link lies in
\(n_p\) plaquettes.  On a nondegenerate periodic lattice, the resulting
neighbor pairs are distinct, giving the two identities in
\eqref{eq:v6v42-absolute-row}.  All link diagonal entries are equal, so
normalization divides the absolute off-diagonal row by \(\beta n_p\),
giving three.  The common factor \(\beta\) cancels.
\end{proof}

The failure is not instability of the action.  The signed Hessian is the
nonnegative square \(\beta d^*d\).  Absolute link incidence has discarded
the orientation cancellations and the gauge complex.

\begin{corollary}[Weak coupling does not improve absolute incidence]
\label{cor:v6v42-beta-no-help}
Letting \(\beta\to\infty\) narrows the one-plaquette conditional law but
does not make the normalized absolute influence matrix contractive.  A
proof of UPG cannot use inverse temperature alone to turn
\eqref{eq:v6v42-row-sum-three} into a number below one.
\end{corollary}

\begin{proof}
The normalized row sum in
\eqref{eq:v6v42-row-sum-three} is independent of \(\beta\).
\end{proof}

\subsection{Hodge and shell correction}

On the gauge-orthogonal one-form space, the Fourier symbol of the abelian
Hessian is
\begin{equation}
 H(k)=\beta\widehat k^{,2}P_T(k),
 \label{eq:v6v42-transverse-symbol}
\end{equation}
consistent with the exact Ward decomposition in V27.  The zero and small
momenta are physical photon channels; no volume-independent gap exists on
the full transverse space.

\begin{proposition}[Shellwise physical gap]
\label{prop:v6v42-shell-gap}
On a transverse momentum shell \(\mathcal S_j\) satisfying
\(\widehat k^{,2}\ge\kappa_j^2>0\),
\begin{equation}
 \langle A,HA\rangle
 \ge\beta\kappa_j^2\|A\|_2^2
 \qquad(A\in\operatorname{Ran}P_T\mathbf1_{\mathcal S_j}).
 \label{eq:v6v42-shell-gap}
\end{equation}
The low-momentum complement must remain a physical massless IMS edge.  It
cannot be included in the UPG of Definition~\ref{def:v6v41-UPG}.
\end{proposition}

\begin{proof}
Functional calculus of \eqref{eq:v6v42-transverse-symbol} gives the lower
bound on the indicated shell.  At \(k=0\), and in the infinite-volume limit
as \(k\to0\), the eigenvalue vanishes, proving the necessity of the
retained complement.
\end{proof}

\begin{assessment}[Correct block choice]
\label{ass:v6v42-correct-blocks}
The linkwise absolute influence route is ruled out already by the free
abelian tangent.  A viable BMI proof must use blocks adapted to the signed
de Rham/gauge complex, or a multiscale shell decomposition on the physical
quotient.  Massive Higgsed gauge components may use a true block gap;
unbroken photon and graviton components require shellwise coercivity plus
explicit massless edges.
\end{assessment}

\begin{firewall}[The BMI inequality is not acquired]
\label{fw:v6v42-BMI-open}
Theorem~\ref{thm:v6v42-block-tensorization} is an exact implication from
\eqref{eq:v6v42-BMI}.  No strict BMI matrix has yet been constructed for
the nonlinear compact gauge--Higgs path.  The plaquette calculation proves
that an entrywise absolute link matrix is the wrong candidate.
\end{firewall}

\begin{firewall}[Shell coercivity is not conditional mixing]
\label{fw:v6v42-shell-not-mixing}
Proposition~\ref{prop:v6v42-shell-gap} is a quadratic Fourier statement.
It does not control compact large-field tails, nonlinear plaquette
interactions, conditional boundary laws, or the chiral determinant.  Those
rows must be added without destroying the signed shell gap.
\end{firewall}

\begin{theorem}[Sixth-series V42 result]
\label{thm:v6v42-status}
Uniform conditional block gaps plus a strict martingale influence matrix
give a volume-independent physical gap
\(\lambda_0(1-\eta)^2\).  The abelian plaquette tangent proves that the
obvious absolute link influence matrix has normalized row sum three,
independent of weak coupling, even though the signed Hessian is
nonnegative.  The correct gap acquisition must therefore retain the Hodge
sign structure and work shellwise on the physical quotient, leaving the
unbroken photon and graviton low modes as explicit massless IMS channels.
\end{theorem}

\begin{proof}
Use Theorem~\ref{thm:v6v42-block-tensorization},
Proposition~\ref{prop:v6v42-Gaussian-criterion},
Theorem~\ref{thm:v6v42-incidence-no-go}, and
Proposition~\ref{prop:v6v42-shell-gap}, subject to both firewalls.
\end{proof}

V43 will construct a signed multiscale BMI on transverse shells.  It will
derive the Schur complement for one high shell in the nonlinear bosonic
Hessian, prove a quantitative gap under a small off-shell coupling norm,
and expose the exact large-field tail estimate needed to make that norm
uniform along the intrinsic path.


\clearpage
\section{V43: signed shell Schur gaps and the large-field capacity gate}
\label{sec:v6v43}

V42 shows that absolute link incidence discards the gauge-complex signs.
We now use the signed operator Hessian on physical shells.  The small-field
part is controlled by a normalized Schur coupling.  Extending that control
to the compact large-field region requires a capacity estimate; small tail
probability by itself is insufficient.

\subsection{Signed two-block Schur coercivity}

Let \(X,Y\) be Hilbert spaces and
\begin{equation}
 \mathcal H=
 \begin{pmatrix}A&B\\B^*&C\end{pmatrix},
 \qquad A\ge aI>0,
 \qquad C\ge cI>0.
 \label{eq:v6v43-block-Hessian}
\end{equation}
Define the normalized off-shell coupling
\begin{equation}
 \delta:=\|A^{-1/2}BC^{-1/2}\|.
 \label{eq:v6v43-normalized-coupling}
\end{equation}

\begin{theorem}[Signed Schur gap]
\label{thm:v6v43-signed-Schur-gap}
If \(\delta<1\), then
\begin{align}
 \mathcal H
 &\ge(1-\delta)(A\oplus C),
 \label{eq:v6v43-full-coercivity}\\
 S_A:=A-BC^{-1}B^*
 &\ge(1-\delta^2)A,
 \label{eq:v6v43-Schur-coercivity}\\
 S_C:=C-B^*A^{-1}B
 &\ge(1-\delta^2)C.
 \label{eq:v6v43-Schur-coercivity-C}
\end{align}
In particular, the full gap is at least
\((1-\delta)\min\{a,c\}\), while the effective \(X\)-shell gap is at
least \((1-\delta^2)a\).
\end{theorem}

\begin{proof}
Put \(u=A^{1/2}x\), \(v=C^{1/2}y\), and
\(T=A^{-1/2}BC^{-1/2}\).  Then
\begin{align*}
 \langle(x,y),\mathcal H(x,y)\rangle
 &=\|u\|^2+\|v\|^2+2\operatorname{Re}\langle u,Tv\rangle\\
 &\ge(1-\delta)(\|u\|^2+\|v\|^2),
\end{align*}
which proves \eqref{eq:v6v43-full-coercivity}.  Moreover
\[
 BC^{-1}B^*=A^{1/2}TT^*A^{1/2}le\delta^2A,
\]
giving \eqref{eq:v6v43-Schur-coercivity}.  Exchange the two blocks for the
last estimate.
\end{proof}

Unlike an entrywise absolute influence matrix, \(\delta\) retains all
orientation and representation signs inside \(B\) before taking one
operator norm.

\subsection{Physical shell version with retained kernels}

Let \(P_j\) be a transverse physical shell projector and
\(P_{\rm r}=I-P_j\).  Split a bosonic Hessian on the gauge quotient as
\begin{equation}
 \mathcal H_u=
 \begin{pmatrix}
 A_{j,u}&B_{j,u}\\B_{j,u}^*&C_{j,u}
 \end{pmatrix}
 \label{eq:v6v43-shell-Hessian}
\end{equation}
on \(P_j\mathscr H\oplus P_{\rm r}\mathscr H\).  If
\(C_{j,u}\) has a kernel \(K_{j,u}\) consisting of retained physical
massless/topological modes, let \(Q_{j,u}\) project onto its orthogonal
complement.

\begin{proposition}[Reduced shell Schur estimate]
\label{prop:v6v43-reduced-shell}
Assume
\begin{equation}
 A_{j,u}\ge a_jI,
 \qquad
 Q_{j,u}C_{j,u}Q_{j,u}\ge c_jQ_{j,u},
 \qquad
 P_{K_{j,u}}B_{j,u}^*=0.
 \label{eq:v6v43-reduced-shell-assumptions}
\end{equation}
Define
\begin{equation}
 \delta_{j,u}:=left\|
 A_{j,u}^{-1/2}B_{j,u}Q_{j,u}
 (Q_{j,u}C_{j,u}Q_{j,u})^{-1/2}
 \right\|.
 \label{eq:v6v43-reduced-delta}
\end{equation}
If \(\delta_{j,u}<1\), elimination of the gapped complement gives a shell
Hessian at least \((1-\delta_{j,u}^2)a_j\).  If
\(P_{K_{j,u}}B_{j,u}^*\ne0\), those couplings cannot be Schur-eliminated
and must remain as explicit massless IMS legs.
\end{proposition}

\begin{proof}
Apply Theorem~\ref{thm:v6v43-signed-Schur-gap} to
\(P_j\mathscr H\oplus Q_{j,u}P_{\rm r}\mathscr H\).  The third assumption
makes the retained kernel an uncoupled output for this elimination.  If it
fails, the inverse of \(C_{j,u}\) does not exist on the range of the
coupling, so no Schur estimate can remove that row.
\end{proof}

This is the bosonic analogue of the zero-mode-complete fermion split in
V39.

\subsection{Small-field nonlinear perturbation}

Write the physical Hessian on a good one-well region \(\mathcal G_j\) as
\begin{equation}
 \mathcal H_u(\omega)=\mathcal H_{0,j}+W_{j,u}(\omega),
 \qquad
 \mathcal H_{0,j}>0
 \label{eq:v6v43-Hessian-perturbation}
\end{equation}
after removing retained massless directions.

\begin{lemma}[Normalized perturbation stability]
\label{lem:v6v43-normalized-perturbation}
If
\begin{equation}
 \sup_{\omega\in\mathcal G_j,\,u\in[0,1]}
 \left\|
 \mathcal H_{0,j}^{-1/2}W_{j,u}(\omega)
 \mathcal H_{0,j}^{-1/2}
 \right\|\le\varepsilon_j<1,
 \label{eq:v6v43-good-set-smallness}
\end{equation}
then
\begin{equation}
 \mathcal H_u(\omega)
 \ge(1-\varepsilon_j)\mathcal H_{0,j}
 \qquad(\omega\in\mathcal G_j).
 \label{eq:v6v43-good-set-gap}
\end{equation}
The same conclusion holds for any shell Schur complement whose normalized
off-shell coupling remains below one.
\end{lemma}

\begin{proof}
Factor the quadratic form as
\[
 \mathcal H_u
 =\mathcal H_{0,j}^{1/2}
 \left[I+mathcal H_{0,j}^{-1/2}W_{j,u}
 \mathcal H_{0,j}^{-1/2}\right]
 \mathcal H_{0,j}^{1/2}.
\]
The middle self-adjoint factor is at least
\((1-\varepsilon_j)I\).  For a shell Schur complement, combine this bound
with Theorem~\ref{thm:v6v43-signed-Schur-gap}.
\end{proof}

The good-set Hessian is not yet a global Poincare inequality.  The
configuration can leave \(\mathcal G_j\), and a rare region can still form
a narrow conductance neck.

\subsection{Small bad mass is not enough}

\begin{theorem}[Rare-neck counterexample]
\label{thm:v6v43-rare-neck}
There is a family of smooth positive probability densities \(\rho_\epsilon\)
on \([-2,2]\) and fixed sets
\(\mathcal B=[-1/2,1/2]\) such that
\begin{equation}
 \nu_\epsilon(\mathcal B)\longrightarrow0,
 \label{eq:v6v43-small-bad-mass}
\end{equation}
while the spectral gap of the reversible one-dimensional Dirichlet form
\begin{equation}
 \mathcal E_\epsilon(f)=
 \int_{-2}^2|f'(x)|^2\rho_\epsilon(x)\,dx
 \label{eq:v6v43-neck-form}
\end{equation}
tends to zero.  Therefore exponentially small large-field probability does
not by itself extend a good-set Hessian gap to a global path gap.
\end{theorem}

\begin{proof}
Choose a smooth even density that is comparable to one on
\([-2,-1]\cup[1,2]\), comparable to \(\epsilon^2\) on
\([-1/2,1/2]\), and smoothly interpolated on the remaining fixed
intervals, then normalize it.  Its mass on \(\mathcal B\) is
\(O(\epsilon^2)\).

Let \(f_\epsilon\) equal \(-1\) on \([-2,-1/2]\), \(+1\) on
\([1/2,2]\), and interpolate linearly across \(\mathcal B\).  Evenness
gives \(\mathbb E f_\epsilon=0\), while
\(\operatorname{Var}(f_\epsilon)\ge c>0\) because both outer wells have
fixed mass.  The derivative is supported in \(\mathcal B\) and uniformly
bounded, so
\(\mathcal E_\epsilon(f_\epsilon)=O(\epsilon^2)\).  The Rayleigh quotient
upper-bounds the spectral gap by \(C\epsilon^2\), proving the claim.
\end{proof}

The correct tail quantity couples its mass to the inverse density across
the return path.

\subsection{One-dimensional capacity/Hardy card}

Let \(d\nu(r)=\rho(r)\,dr\) on \([0,R]\), with \(\rho>0\), and choose a
one-well median \(r_0\).  Define
\begin{align}
 B_+&:=\sup_{r_0<r<R}
 \nu([r,R])\int_{r_0}^r\rho(s)^{-1}\,ds,
 \label{eq:v6v43-Hardy-plus}\\
 B_-&:=\sup_{0<r<r_0}
 \nu([0,r])\int_r^{r_0}\rho(s)^{-1}\,ds.
 \label{eq:v6v43-Hardy-minus}
\end{align}

\begin{theorem}[One-well Hardy sufficiency]
\label{thm:v6v43-Hardy}
If \(B:=\max\{B_+,B_-\}<\infty\), then
\begin{equation}
 \operatorname{Var}_\nu(f)
 \le4B\int_0^R|f'(r)|^2\rho(r)\,dr
 \label{eq:v6v43-Hardy-Poincare}
\end{equation}
for every absolutely continuous \(f\).
\end{theorem}

\begin{proof}
Subtract \(f(r_0)\).  On the right half-interval,
\begin{equation}
 |f(r)-f(r_0)|^2
 \le\left(\int_{r_0}^r|f'(s)|^2\rho(s)\,ds\right)
 \left(\int_{r_0}^r\rho(s)^{-1}\,ds\right)
 \label{eq:v6v43-point-Hardy}
\end{equation}
by Cauchy--Schwarz.  The standard integration-by-parts form of the weighted
Hardy inequality applies the tail factor
\(\nu([r,R])\) before the supremum and gives
\begin{equation}
 \int_{r_0}^R|f-f(r_0)|^2\,d\nu
 \le4B_+\int_{r_0}^R|f'|^2\,d\nu.
\end{equation}
The left half is identical with \(B_-\).  Finally, the mean minimizes \(c\mapsto\int|f-c|^2d\nu\), so
\(\operatorname{Var}_\nu(f)\le\int|f-f(r_0)|^2d\nu\).  This proves
\eqref{eq:v6v43-Hardy-Poincare}.
\end{proof}

\begin{definition}[Uniform conditional tail-capacity card]
\label{def:v6v43-UTC}
The UTC holds on shell \(j\) if every one-block radial conditional density
\(\rho_{j,u}^\xi\), uniformly in path parameter \(u\), physical boundary
condition \(\xi\), and retained sector, satisfies
\begin{equation}
 \max\{B_+(\rho_{j,u}^\xi),B_-(\rho_{j,u}^\xi)\}
 \le C a_j^{-1},
 \label{eq:v6v43-UTC}
\end{equation}
where \(a_j\) is the quadratic shell scale.  Angular conditional modes must
obey the corresponding physical quotient gap.
\end{definition}

\begin{proposition}[What UTC would complete]
\label{prop:v6v43-UTC-sufficiency}
Assume the good-set normalized Hessian bound
\eqref{eq:v6v43-good-set-smallness}, the signed reduced shell coupling
\(\sup_u\delta_{j,u}<1\), UTC, and a uniform angular quotient gap.  Then
the one-shell conditional Poincare constant has the correct scale
\(O(a_j^{-1})\).  If the shell martingale influence matrix of V42 is also
strict, these conditional gaps tensorize to UPG on the gapped physical
shells.
\end{proposition}

\begin{proof}
Lemma~\ref{lem:v6v43-normalized-perturbation} and
Proposition~\ref{prop:v6v43-reduced-shell} control the quadratic core and
off-shell elimination.  Theorem~\ref{thm:v6v43-Hardy} controls radial
escape and return with constant \(O(a_j^{-1})\); the angular hypothesis
controls the remaining quotient directions.  Apply
Theorem~\ref{thm:v6v42-block-tensorization} to the resulting conditional
shell gaps.
\end{proof}

\begin{firewall}[UTC is stronger than a tail bound]
\label{fw:v6v43-UTC-open}
No UTC has yet been proved for the coupled compact gauge--Higgs path with
arbitrary physical boundary data.  The rare-neck counterexample shows that
an estimate of \(\nu(\mathcal G_j^c)\) alone cannot replace it.  One must
control the inverse-density return integral or an equivalent conductance
capacity.
\end{firewall}

\begin{firewall}[Signed Schur smallness remains shell dependent]
\label{fw:v6v43-delta-open}
The quadratic Hodge operator has a signed shell gap, but no uniform bound
\(\delta_{j,u}<1\) has been proved after nonlinear Higgs, boundary,
determinant, and Einstein couplings are inserted.  Couplings to retained
massless kernels must stay explicit rather than being hidden in a
pseudoinverse.
\end{firewall}

\begin{theorem}[Sixth-series V43 result]
\label{thm:v6v43-status}
A normalized signed off-shell coupling \(\delta<1\) gives exact Schur
coercivity \((1-\delta^2)A\) on a physical shell and retains all massless
kernel legs explicitly.  A normalized nonlinear Hessian perturbation below
one preserves the small-field shell gap.  Extending that gap globally
cannot follow from small large-field probability: a rare-neck family has
vanishing bad mass and vanishing spectral gap.  The correct acquisition is
the uniform conditional tail-capacity card
\eqref{eq:v6v43-UTC}, combined with signed shell influence rather than
absolute link incidence.
\end{theorem}

\begin{proof}
Use Theorem~\ref{thm:v6v43-signed-Schur-gap},
Proposition~\ref{prop:v6v43-reduced-shell},
Lemma~\ref{lem:v6v43-normalized-perturbation},
Theorem~\ref{thm:v6v43-rare-neck},
Theorem~\ref{thm:v6v43-Hardy}, and
Proposition~\ref{prop:v6v43-UTC-sufficiency}, subject to both firewalls.
\end{proof}

V44 will attack UTC for a calibrated compact one-well conditional density.
It will prove explicit upper and lower radial density bounds, evaluate both
Hardy products without a global oscillation factor, and identify which
boundary or matter terms can create a second well.


\clearpage
\section{V44: a uniform abelian path capacity and the staple-cancellation no-go}
\label{sec:v6v44}

V43 replaces large-field tail probability by the one-well Hardy products.
They can be evaluated without a global density-ratio oscillation.  We first
prove a general monotone one-well estimate on a compact interval and apply
it to the entire calibrated \(U(1)\) heat--Wilson path.  We then show that
uniformity in arbitrary multi-plaquette boundary data fails for one-link
blocks because the conditional staple can cancel exactly.

\subsection{A compact monotone one-well Hardy theorem}

Let \(V\in C^1([0,R])\) be increasing, with \(V(0)=0\).  Assume there are
positive constants \(c_0,C_0,c_1,r_c\), independent of \(\alpha\ge1\),
such that
\begin{align}
 c_0r^2\le V(r)&\le C_0r^2
 &&(0\le r\le r_c),
 \label{eq:v6v44-core-quadratic}\\
 V'(r)&\ge c_1\min\{r,R-r\}
 &&(0<r<R).
 \label{eq:v6v44-one-well-drift}
\end{align}
Let \(\nu_\alpha\) have density
\begin{equation}
 \rho_\alpha(r)=Z_\alpha^{-1}e^{-\alpha V(r)}
 \label{eq:v6v44-one-well-density}
\end{equation}
on \([0,R]\), and choose any median \(r_\alpha\).

\begin{theorem}[Uniform compact one-well capacity]
\label{thm:v6v44-one-well-Hardy}
There is \(C\), depending only on the constants in
\eqref{eq:v6v44-core-quadratic}--\eqref{eq:v6v44-one-well-drift} and on
\(R\), such that
\begin{equation}
 r_\alpha\asymp\alpha^{-1/2}
 \label{eq:v6v44-median-scale}
\end{equation}
and
\begin{align}
 \sup_{r>r_\alpha}
 \nu_\alpha([r,R])
 \int_{r_\alpha}^r\rho_\alpha(s)^{-1}\,ds
 &\le\frac C\alpha,
 \label{eq:v6v44-Hardy-right}\\
 \sup_{r<r_\alpha}
 \nu_\alpha([0,r])
 \int_r^{r_\alpha}\rho_\alpha(s)^{-1}\,ds
 &\le\frac C\alpha.
 \label{eq:v6v44-Hardy-left}
\end{align}
Consequently the radial spectral gap is at least \(c\alpha\).
\end{theorem}

\begin{proof}
The normalization cancels from each Hardy product, so work with
\(q_\alpha=e^{-\alpha V}\).  On
\([0,r_c]\), \eqref{eq:v6v44-core-quadratic} compares the mass with a
Gaussian at scale \(\alpha^{-1/2}\).  On
\([r_c,R]\), monotonicity and \(V(r_c)>0\) make the mass exponentially
smaller.  Hence fixed positive fractions of the total mass lie between
\(c\alpha^{-1/2}\) and \(C\alpha^{-1/2}\), proving
\eqref{eq:v6v44-median-scale}.

For \(r<r_\alpha\), the quadratic comparison and the changes of variables
\(y=\sqrt\alpha r\), \(z=\sqrt\alpha s\) give
\begin{align*}
 \left(\int_0^rq_\alpha(s)\,ds\right)
 \left(\int_r^{r_\alpha}q_\alpha(s)^{-1}\,ds\right)
 \le\frac C\alpha.
\end{align*}
All scaled integration endpoints remain in a fixed interval, and the bound
also holds as \(r\downarrow0\) by direct Gaussian integration.  This proves
\eqref{eq:v6v44-Hardy-left}.

For the right product, first take
\(r\le R/2\) with \(r\ge r_\alpha\).  On the interval from \(r\) to
\(R/2\), the drift is at least \(c\alpha r\), so integration of
\((\log q_\alpha)'=-\alpha V'\) gives
\begin{equation}
 \int_r^{R/2}q_\alpha(s)\,ds
 \le\frac{Cq_\alpha(r)}{\alpha r}.
 \label{eq:v6v44-right-tail-core}
\end{equation}
The remaining interval \([R/2,R]\) is exponentially smaller than the
right side after constants are enlarged.  Reversing the same drift estimate
gives
\begin{equation}
 \int_{r_\alpha}^rq_\alpha(s)^{-1}\,ds
 \le\frac{Cq_\alpha(r)^{-1}}{\alpha r}
 \label{eq:v6v44-inverse-core}
\end{equation}
when \(r\) is a fixed multiple of \(\alpha^{-1/2}\) or larger; the bounded
scaled core covers the remaining values.  Their product is at most
\(C/(\alpha^2r^2)\le C/\alpha\).

Finally let \(r=R-z\ge R/2\).  Monotonicity gives
\begin{equation}
 \int_r^Rq_\alpha(s)\,ds\le zq_\alpha(r).
 \label{eq:v6v44-end-tail}
\end{equation}
On \([R/2,r]\), the drift lower bound is
\(c\alpha(R-s)\ge c\alpha z\); hence
\begin{equation}
 \int_{r_\alpha}^rq_\alpha(s)^{-1}\,ds
 \le\frac{Cq_\alpha(r)^{-1}}{\alpha z},
 \label{eq:v6v44-end-inverse}
\end{equation}
with the part below \(R/2\) absorbed by the same exponential comparison.
The product of \eqref{eq:v6v44-end-tail} and
\eqref{eq:v6v44-end-inverse} is at most \(C/\alpha\).  This proves the
right Hardy bound.  Apply Theorem~\ref{thm:v6v43-Hardy} for the gap.
\end{proof}

\begin{lemma}[Stability under bounded density distortion]
\label{lem:v6v44-bounded-distortion}
If two normalized densities satisfy
\(c\rho\le\widetilde\rho\le C\rho\), then each Hardy product for
\(\widetilde\rho\) is at most \(C/c\) times the corresponding product for
\(\rho\), after moving the median within a fixed quantile interval.  Thus a
uniform multiplicative density comparison preserves the \(O(\alpha^{-1})\)
capacity scale even when the global logarithmic oscillation between two
different endpoints is order \(\alpha\).
\end{lemma}

\begin{proof}
The tail mass is enlarged by at most \(C\), while the inverse-density
integral is enlarged by at most \(c^{-1}\).  Density comparability moves a
median only between fixed quantiles of the reference law; the proof of
Theorem~\ref{thm:v6v44-one-well-Hardy} is uniform on that quantile interval.
\end{proof}

\subsection{The full calibrated \(U(1)\) heat--Wilson path}

Retain the circle heat kernel \eqref{eq:v6v41-circle-heat} and Wilson density
\eqref{eq:v6v41-von-Mises}.  On the class interval \(0\le r\le\pi\), set
\begin{equation}
 \rho_{\beta,u}(r)=Z_{\beta,u}^{-1}
 h_{1/\beta}(r)^{1-u}e^{u\beta\cos r},
 \qquad 0\le u\le1.
 \label{eq:v6v44-abelian-path-density}
\end{equation}

\begin{lemma}[Uniform periodic-Gaussian envelope]
\label{lem:v6v44-heat-envelope}
For \(\beta\ge1\) and \(0\le r\le\pi\),
\begin{equation}
 c\sqrt\beta e^{-\beta r^2/2}
 \le h_{1/\beta}(r)
 \le C\sqrt\beta e^{-\beta r^2/2},
 \label{eq:v6v44-heat-envelope}
\end{equation}
with numerical constants independent of \(\beta,r\).
\end{lemma}

\begin{proof}
The \(k=0\) term in the periodic Gaussian sum gives the lower bound.  For
\(0\le r\le\pi\), the \(k=-1\) term divided by the \(k=0\) term is
\begin{equation}
 \exp[-2\pi\beta(\pi-r)]\le1.
 \label{eq:v6v44-second-image}
\end{equation}
All remaining image ratios are bounded by two geometric Gaussian tails,
uniformly for \(\beta\ge1\).  Their sum is bounded by a numerical constant,
proving the upper estimate.
\end{proof}

Define
\begin{equation}
 V_u(r):=(1-u)\frac{r^2}{2}+u(1-\cos r).
 \label{eq:v6v44-effective-path-potential}
\end{equation}

\begin{lemma}[Uniform one-well geometry]
\label{lem:v6v44-uniform-one-well}
Uniformly in \(u\in[0,1]\),
\begin{align}
 c r^2\le V_u(r)&\le Cr^2
 &&(0\le r\le\pi/2),
 \label{eq:v6v44-path-core}\\
 V_u'(r)=(1-u)r+u\sin r
 &\ge c\min\{r,\pi-r\}
 &&(0<r<\pi).
 \label{eq:v6v44-path-drift}
\end{align}
\end{lemma}

\begin{proof}
Use the elementary quadratic bounds for \(1-\cos r\) on
\([0,\pi/2]\).  Also
\(\sin r\ge(2/\pi)\min\{r,\pi-r\}\).  If \(u\ge1/2\), the sine term
gives the drift bound; if \(u<1/2\), the term \((1-u)r\) does.
\end{proof}

\begin{theorem}[Uniform abelian path UTC]
\label{thm:v6v44-abelian-UTC}
There is a numerical \(C\) such that the two Hardy products of
\(\rho_{\beta,u}\) satisfy
\begin{equation}
 \max\{B_+(\rho_{\beta,u}),B_-(\rho_{\beta,u})\}
 \le\frac C\beta
 \label{eq:v6v44-abelian-UTC}
\end{equation}
for every \(\beta\ge1\) and \(u\in[0,1]\).  Hence the intrinsic radial
path gap is at least \(c\beta\), uniformly along the entire heat--Wilson
interpolation.
\end{theorem}

\begin{proof}
Lemma~\ref{lem:v6v44-heat-envelope} shows that
\(\rho_{\beta,u}\) is uniformly multiplicatively comparable to the
normalized density
\begin{equation}
 \widetilde\rho_{\beta,u}(r)
 \propto e^{-\beta V_u(r)}.
 \label{eq:v6v44-comparison-density}
\end{equation}
The constants remain uniform after raising the heat envelope to the power
\(1-u\).  Lemma~\ref{lem:v6v44-uniform-one-well} verifies the hypotheses of
Theorem~\ref{thm:v6v44-one-well-Hardy} with \(\alpha=\beta\), uniformly in
\(u\).  Apply Lemma~\ref{lem:v6v44-bounded-distortion} and then the Hardy
Poincare theorem.
\end{proof}

This result explains why the order-\(\beta\) log-density oscillation in V41
was not a genuine gap obstruction: the Hardy product pairs each small tail
with the inverse-density cost of reaching that same tail.

\subsection{Arbitrary boundary data can cancel the staple}

Now place a link angle \(\theta_e\) in a multi-plaquette \(U(1)\) Wilson
action and fix every other link.  Each incident plaquette contributes
\(\beta\cos(\sigma_p\theta_e+\phi_p)\), with
\(\sigma_p=\pm1\).  After conjugating phases, define the complex staple
\begin{equation}
 H_e(\xi):=\sum_{p\ni e}e^{i\psi_p(\xi)}.
 \label{eq:v6v44-staple}
\end{equation}
Then the conditional density of \(\theta_e\) is
\begin{equation}
 d\nu_e^\xi(\theta)
 =\frac{
 e^{\beta\operatorname{Re}(e^{i\theta}H_e(\xi))}
 }{2\pi I_0(\beta|H_e(\xi)|)}\,d\theta.
 \label{eq:v6v44-conditional-von-Mises}
\end{equation}

\begin{theorem}[Staple-cancellation obstruction]
\label{thm:v6v44-staple-cancellation}
There are admissible fixed boundary configurations \(\xi\) for which
\begin{equation}
 H_e(\xi)=0.
 \label{eq:v6v44-zero-staple}
\end{equation}
For such data the conditional link law is uniform, so its spectral gap in
the standard angular Dirichlet form is order one rather than order
\(\beta\).  Therefore the arbitrary-boundary, one-link UTC with target
constant \(C/\beta\) is false on a multi-plaquette lattice.
\end{theorem}

\begin{proof}
At least two plaquettes meet an interior link.  Choose their fixed staple
phases in opposite pairs; on a hypercubic lattice the remaining incident
phases can likewise be paired, giving
\(\sum_{p\ni e}e^{i\psi_p}=0\).  These phases are realized by assigning the
other link angles on the finite star of \(e\), with any consistency
constraints absorbed outside that star.  Equation
\eqref{eq:v6v44-conditional-von-Mises} then becomes normalized Haar measure
\(d\theta/(2\pi)\).  Its first nonconstant Fourier modes have gap one for
the form \(\int|f'|^2d\theta/(2\pi)\), independent of \(\beta\).  It cannot
satisfy a lower gap bound \(c\beta\) uniformly as \(\beta\to\infty\).
\end{proof}

\begin{assessment}[Correction to the UTC quantifier]
\label{ass:v6v44-UTC-correction}
Definition~\ref{def:v6v43-UTC} is a useful sufficient card, but its demand
for the shell-scale gap under every one-link boundary condition is too
strong.  The correct alternatives are:
\begin{enumerate}[label=\textup{(\roman*)},leftmargin=3.7em]
\item use blocks large enough that internal plaquette curvature cannot be
      canceled by arbitrary exterior staples;
\item prove a good-boundary UTC and pay the bad-boundary set by a joint
      boundary capacity, not by its probability alone; or
\item retain the signed multiscale Hodge Hessian and avoid conditioning down
      to one raw link.
\end{enumerate}
\end{assessment}

\begin{firewall}[The proved UTC is one-plaquette abelian]
\label{fw:v6v44-scope}
Theorem~\ref{thm:v6v44-abelian-UTC} proves the complete calibrated radial
path gap for the abelian one-plaquette quotient.  It does not prove a
nonabelian angular gap, a volume-uniform interacting path gap, or stability
under chiral determinant and Einstein boundary terms.
\end{firewall}

\begin{firewall}[Bad staples require capacity, not exclusion]
\label{fw:v6v44-bad-staples}
The configurations with small \(|H_e|\) cannot simply be discarded as rare.
The rare-neck theorem of V43 requires a conductance/capacity estimate for
their connection to the good-staple region.  No such joint boundary
capacity has yet been proved.
\end{firewall}

\begin{theorem}[Sixth-series V44 result]
\label{thm:v6v44-status}
Monotone compact one-well densities with quadratic core and a two-sided
endpoint drift have Hardy products \(O(\alpha^{-1})\).  The full calibrated
\(U(1)\) heat--Wilson path satisfies these hypotheses uniformly and has
radial gap at least \(c\beta\), despite its order-\(\beta\) global
log-density oscillation.  On a multi-plaquette lattice, however, arbitrary
boundary staples can cancel and reduce a one-link conditional gap to order
one.  The full path gap must therefore use larger physical blocks, a joint
bad-boundary capacity, or signed multiscale Hodge conditioning.
\end{theorem}

\begin{proof}
Use Theorem~\ref{thm:v6v44-one-well-Hardy},
Lemma~\ref{lem:v6v44-heat-envelope},
Lemma~\ref{lem:v6v44-uniform-one-well},
Theorem~\ref{thm:v6v44-abelian-UTC}, and
Theorem~\ref{thm:v6v44-staple-cancellation}, subject to both firewalls.
\end{proof}

V45 will perform the compulsory YM/NS audit.  It will compare the newest
Wilson bridge and NS formation-capacity results with the abelian UTC and
staple obstruction, then select the correct nonabelian block or capacity
acquisition before continuing the raw SM comparison.


\clearpage
\section{V45: compulsory audit and the joint staple-capacity direction}
\label{sec:v6v45}

\subsection{Read-only companion audit}

The newest active companion notes inspected before this append were
\begin{center}
\begin{tabular}{p{0.59\textwidth}r r}
\hline
file & bytes & logical lines\\
\hline
\texttt{Renewal\_Yang\_Mills\_Mass\_Gap\_Research\_Notes\_V32.tex}
& \(471\,910\) & \(14\,849\)\\
\texttt{Renewal\_NS\_Stability\_Research\_Notes\_V31.tex}
& \(887\,537\) & \(23\,052\)\\
\hline
\end{tabular}
\end{center}
Their SHA--256 digests were
\begin{align}
 &\texttt{a98720fac11d309a3dca050b2613a2a61f0411329f4ee61253eff7558719a47b},
 \label{eq:v6v45-YM-hash}\\
 &\texttt{f1121b62238ad5491bb7cf03635ea9934942edf573ed8faaa9ee79e1d38a6696}.
 \label{eq:v6v45-NS-hash}
\end{align}
Neither file was modified.

YM V30--V31 proves uniform overlap, concentration, paired heat-image slope
bounds, and a spectral gap of order \(\alpha\) along the entire calibrated
\(SU(2)\) heat--Wilson bridge.  YM V32 chooses the second-order heat time
that cancels the quadratic bridge score, proves a global quartic-value and
cubic-gradient profile, and obtains an intrinsic negative score norm
\(O(\alpha^{-3/2})\).  Its remaining gate is the all-arity raw bridge-score
forest, not the one-link gap.

NS V31 is unchanged.  Its sharp formation and unmarking tests continue to
warn that a small normalized exceptional set does not pay a restoring
operation unless its physical ancestry or capacity is retained.

\begin{firewall}[No duplicate transfer]
\label{fw:v6v45-no-duplicate}
The \(SU(2)\) bridge gap and score rows are not rederived here.  They are
recorded as proved companion inputs at the exact checkpoint above.  They do
not imply an \(SU(3)\) alcove gap, a multi-plaquette path gap, a chiral
determinant estimate, or an Einstein continuum theorem.
\end{firewall}

\subsection{Product and finite-central-quotient gap reduction}

Let \((X_a,\nu_a,\mathcal E_a)\), \(1\le a\le q\), be reversible factors,
and equip their product with
\begin{equation}
 \nu=\bigotimes_{a=1}^q\nu_a,
 \qquad
 \mathcal E(f,f)=\sum_{a=1}^q
 \int\mathcal E_a(f(\cdot,x_{\ne a}),f(\cdot,x_{\ne a}))
 \,d\nu_{\ne a}.
 \label{eq:v6v45-product-form}
\end{equation}

\begin{theorem}[Exact product tensorization of the gap]
\label{thm:v6v45-product-gap}
If factor \(a\) has Poincare gap \(\lambda_a>0\), then the product gap is
\begin{equation}
 \lambda_{\rm prod}=\min_a\lambda_a.
 \label{eq:v6v45-product-gap}
\end{equation}
\end{theorem}

\begin{proof}
Iterated conditional variance gives
\begin{equation}
 \operatorname{Var}_\nu(f)
 \le\sum_a\int
 \operatorname{Var}_{\nu_a}(f(\cdot,x_{\ne a})),d\nu_{\ne a}.
 \label{eq:v6v45-Efron-Stein}
\end{equation}
Apply the factor Poincare inequalities and bound
\(\lambda_a^{-1}\le(\min_b\lambda_b)^{-1}\).  This proves the lower gap
bound.  Conversely, a centered test function depending only on a factor
whose gap approaches the minimum has the same Rayleigh quotient in the
product, proving equality.
\end{proof}

\begin{proposition}[Finite quotient cannot lower the physical gap]
\label{prop:v6v45-finite-quotient-gap}
Let a finite group \(\Gamma\) act measure-preservingly on the product and
commute with its generator.  The quotient law, identified with the
\(\Gamma\)-invariant subspace of \(L^2(\nu)\), has gap at least
\(\lambda_{\rm prod}\).
\end{proposition}

\begin{proof}
The quotient Rayleigh infimum is taken over the centered invariant
subspace, a subset of all centered product functions.  Restricting the
infimum cannot decrease it.
\end{proof}

\begin{corollary}[One-link compact-SM reduction]
\label{cor:v6v45-one-link-reduction}
For a factorized one-link bridge on the covering group
\(SU(3)\times SU(2)\times U(1)\), followed by the finite central quotient
on representations that descend, the physical gap is bounded below by the
minimum of the three factor gaps.  V44 supplies the calibrated \(U(1)\)
factor and the audited YM checkpoint supplies the \(SU(2)\) factor.  The
unacquired compact one-link factor is therefore the \(SU(3)\) Weyl-alcove
gap.
\end{corollary}

\begin{proof}
Apply Theorem~\ref{thm:v6v45-product-gap} and
Proposition~\ref{prop:v6v45-finite-quotient-gap}.  The final status statement
uses only the two explicitly recorded factor results.
\end{proof}

Matter and plaquette interactions destroy this product factorization; the
corollary is a one-link endpoint reduction, not the lattice path gap.

\subsection{Geometry of the zero-staple set}

An interior link in a \(d\)-dimensional hypercubic lattice has
\begin{equation}
 n=2(d-1)
 \label{eq:v6v45-staple-number}
\end{equation}
incident plaquette staples.  Model their abelian phases by
\(\phi=(\phi_1,\ldots,\phi_n)\in\mathbb T^n\) and define
\begin{equation}
 H(\phi):=\sum_{j=1}^ne^{i\phi_j}\in\mathbb C\simeq\mathbb R^2.
 \label{eq:v6v45-staple-map}
\end{equation}

\begin{lemma}[Rank of the staple map]
\label{lem:v6v45-staple-rank}
At \(H(\phi)=0\), the differential \(dH_\phi\) has rank two unless all
vectors \(e^{i\phi_j}\) lie on one unoriented line.  Hence the regular
zero-staple set is a smooth submanifold of codimension two.  The singular
zero set consists of configurations
\begin{equation}
 e^{i\phi_j}\in\{e^{i\psi},-e^{i\psi}\}
 \label{eq:v6v45-collinear-staples}
\end{equation}
with equal numbers of the two signs; each such stratum has dimension one.
\end{lemma}

\begin{proof}
The derivative in coordinate \(j\) is the real two-vector
\(ie^{i\phi_j}\).  These columns span \(\mathbb R^2\) unless they are all
collinear, which is equivalent to
\eqref{eq:v6v45-collinear-staples}.  At a zero of their sum, a collinear
family has the same number of positive and negative vectors.  Apart from
the common angle \(\psi\), only a finite sign assignment remains, so every
singular stratum is one-dimensional.  The regular-value theorem gives the
codimension-two assertion.
\end{proof}

In four Euclidean lattice dimensions, \(n=6\).  The singular strata then
have codimension five and the regular cancellation set has codimension two.
Thus the exact boundary values found in V44 exist, but they are not a
generic hypersurface separating the six-phase staple block.

\subsection{Reference small-staple mass in four dimensions}

Let the six phases be independent and Haar-uniform.  Let \(J_0\) be the
order-zero Bessel function.

\begin{theorem}[Quadratic small-staple volume]
\label{thm:v6v45-small-staple-volume}
For \(n=6\), the random sum
\(H=\sum_{j=1}^6e^{i\phi_j}\) has a bounded continuous density
\(p_6\) on \(\mathbb R^2\), with \(p_6(0)>0\).  Consequently there are
\(c,C,\epsilon_0>0\) such that
\begin{equation}
 c\epsilon^2
 \le\mathbb P(|H|\le\epsilon)
 \le C\epsilon^2
 \qquad(0<\epsilon\le\epsilon_0).
 \label{eq:v6v45-small-staple-volume}
\end{equation}
\end{theorem}

\begin{proof}
The characteristic function of one uniform unit vector is
\(J_0(|\xi|)\); hence that of the six-fold sum is
\(J_0(|\xi|)^6\).  The standard oscillatory integral representation
\begin{equation}
 J_0(r)=\frac1{2\pi}\int_0^{2\pi}e^{ir\cos\theta}\,d\theta
 \label{eq:v6v45-Bessel-integral}
\end{equation}
and a split into neighborhoods of the two stationary points and their
complement gives
\(|J_0(r)|\le C(1+r)^{-1/2}\).  Therefore
\begin{equation}
 \int_{\mathbb R^2}|J_0(|\xi|)|^6,d\xi
 \le C\int_0^\infty r(1+r)^{-3},dr<\infty.
 \label{eq:v6v45-Bessel-L1}
\end{equation}
Fourier inversion yields a bounded continuous density.

Let \(p_3\) be the density of a three-vector sum.  Plancherel applies
because \(J_0^3\in L^2(\mathbb R^2)\), and convolution gives
\begin{equation}
 p_6(0)=\int_{\mathbb R^2}p_3(z)p_3(-z)\,dz
 =\int_{\mathbb R^2}p_3(z)^2\,dz>0.
 \label{eq:v6v45-density-positive}
\end{equation}
Continuity makes \(p_6\) bounded above and below by positive constants on a
small disk around zero.  Integrating over the disk of radius \(\epsilon\)
proves \eqref{eq:v6v45-small-staple-volume}.
\end{proof}

\subsection{Reference capacity of the bad-staple tube}

On \(\mathbb T^6\), define the one-capacity
\begin{equation}
 \operatorname{Cap}_1(A)
 :=\inf_{f\ge1\ {
m near}\ A}
 \int_{\mathbb T^6}(|\nabla f|^2+|f|^2)\,d\phi.
 \label{eq:v6v45-one-capacity}
\end{equation}

\begin{lemma}[Six-dimensional isocapacitary lower bound]
\label{lem:v6v45-isocapacity}
There is \(c>0\) such that every measurable
\(A\subset\mathbb T^6\) satisfies
\begin{equation}
 \operatorname{Cap}_1(A)
 \ge c|A|^{2/3}.
 \label{eq:v6v45-isocapacity}
\end{equation}
\end{lemma}

\begin{proof}
The Sobolev inequality on the compact six-torus is
\begin{equation}
 \|f\|_{L^3}^2
 \le C\int_{\mathbb T^6}(|\nabla f|^2+|f|^2)\,d\phi.
 \label{eq:v6v45-Sobolev-six}
\end{equation}
If \(f\ge1\) on \(A\), then
\(\|f\|_{L^3}^2\ge|A|^{2/3}\).  Take the infimum over admissible
\(f\).
\end{proof}

\begin{corollary}[The reference zero-staple tube is not a low-capacity
neck]
\label{cor:v6v45-reference-capacity}
Let \(A_\epsilon=\{\phi:|H(\phi)|\le\epsilon\}\).  For small \(\epsilon\),
\begin{equation}
 |A_\epsilon|\asymp\epsilon^2,
 \qquad
 \operatorname{Cap}_1(A_\epsilon)\ge c\epsilon^{4/3},
 \qquad
 \frac{|A_\epsilon|}{\operatorname{Cap}_1(A_\epsilon)}
 \le C\epsilon^{2/3}.
 \label{eq:v6v45-reference-capacity}
\end{equation}
Thus the bad-staple tube has more than enough reference capacity relative
to its mass; the one-link failure in V44 is not a bottleneck of the full
six-phase product block.
\end{corollary}

\begin{proof}
The volume statement is Theorem
\ref{thm:v6v45-small-staple-volume}.  Insert it into
Lemma~\ref{lem:v6v45-isocapacity} and divide.
\end{proof}

\begin{assessment}[Audit-induced block choice]
\label{ass:v6v45-block-choice}
The companion audit removes the need to revisit the isolated \(SU(2)\)
bridge gap.  V44's zero-staple boundary values should be handled by a
joint plaquette-star block: in four dimensions their regular set has
codimension two, their product-reference tube has mass \(\asymp\epsilon^2\),
and its capacity dominates that mass.  The next analytic task is to retain
this favorable ratio under the interacting geometric path and the gauge
quotient, rather than imposing a false pointwise boundary gap.
\end{assessment}

\begin{firewall}[Reference capacity is not interacting capacity]
\label{fw:v6v45-interacting-capacity}
Theorem~\ref{thm:v6v45-small-staple-volume} uses six independent uniform
phases, and Lemma~\ref{lem:v6v45-isocapacity} uses the unweighted elliptic
form.  The interacting weak-coupling density is not uniformly comparable
to Haar on the whole star.  A weighted capacity proof must exploit its
one-well action and preserve boundary gauge constraints.
\end{firewall}

\begin{firewall}[The \(SU(3)\) factor remains open]
\label{fw:v6v45-SU3-open}
Product tensorization reduces the isolated compact-SM bridge to its factor
gaps but does not supply the \(SU(3)\) Weyl-alcove estimate.  That factor
and the interacting star capacity are distinct live rows.
\end{firewall}

\begin{theorem}[Sixth-series V45 result]
\label{thm:v6v45-status}
The compulsory audit records the proved companion \(SU(2)\) bridge gap and
negative score row, so the isolated compact-SM one-link gap reduces to the
\(SU(3)\) factor by product tensorization and finite-quotient monotonicity.
For the four-dimensional lattice interaction, the six-phase zero-staple
set is generically codimension two, its product-reference tube has mass
\(\asymp\epsilon^2\), and six-dimensional Sobolev capacity dominates that
mass by a factor \(\epsilon^{-2/3}\).  This selects a joint plaquette-star
capacity, rather than a one-link arbitrary-boundary gap, as the next
interacting acquisition.
\end{theorem}

\begin{proof}
Use Theorem~\ref{thm:v6v45-product-gap},
Proposition~\ref{prop:v6v45-finite-quotient-gap},
Lemma~\ref{lem:v6v45-staple-rank},
Theorem~\ref{thm:v6v45-small-staple-volume}, and
Corollary~\ref{cor:v6v45-reference-capacity}, subject to the three
firewalls.
\end{proof}

V46 will prove a weighted star-capacity comparison on the one-well part of
the geometric path.  It will localize the six-phase action around the
aligned staple manifold, derive a scale-covariant Sobolev inequality, and
test whether the small-staple tube retains a favorable mass-to-capacity
ratio after weak-coupling concentration.


\clearpage
\section{V46: the integrated star well, exponential bad-staple mass, and the capacity gap}
\label{sec:v6v46}

V45 identifies the six-phase plaquette star as the correct block for the
four-dimensional abelian staple cancellation.  We now integrate the central
link exactly.  The resulting block has a unique aligned minimum modulo the
common phase and gives exponential suppression of a small resultant.  A
Gaussian reference has the required scale-covariant capacity, but the
transfer of that capacity to the full star weight remains a distinct step.

\subsection{Relative star coordinates}

For \(n\ge3\), let
\begin{equation}
 H(\phi)=\sum_{j=1}^ne^{i\phi_j},
 \qquad \phi\in\mathbb T^n.
 \label{eq:v6v46-star-resultant}
\end{equation}
The common rotation
\(\phi_j\mapsto\phi_j+\psi\) changes only the phase of \(H\), not
\(|H|\).  Near the aligned manifold, write
\begin{equation}
 \phi_j=\psi+x_j,
 \qquad
 \sum_{j=1}^nx_j=0.
 \label{eq:v6v46-relative-coordinates}
\end{equation}
The relative space has dimension \(m=n-1\).

\begin{lemma}[Small resultant is outside the aligned core]
\label{lem:v6v46-resultant-distance}
For every real lift \(x_j\) satisfying \(\sum_jx_j=0\),
\begin{equation}
 |H|\ge\operatorname{Re}(e^{-i\psi}H)
 =\sum_{j=1}^n\cos x_j
 \ge n-\frac12\sum_{j=1}^nx_j^2.
 \label{eq:v6v46-resultant-distance}
\end{equation}
Consequently
\begin{equation}
 |H|\le\epsilon<n
 \quad\Longrightarrow\quad
 \sum_jx_j^2\ge2(n-\epsilon)
 \label{eq:v6v46-bad-outside-core}
\end{equation}
whenever one common lift chart contains the configuration.
\end{lemma}

\begin{proof}
The first inequality chooses the real component in the common-phase
direction.  The identity is the definition of \(H\), and
\(\cos x\ge1-x^2/2\) gives the last inequality.  Rearrangement proves the
consequence.
\end{proof}

Thus the zero-staple tube is a fixed relative distance from the aligned
well; it is not a small fluctuation whose conditional gap merely loses a
factor \(|H|\).

\subsection{Exact integration of the central link}

Fix the staple phases and integrate the central link angle \(\theta\).  The
Wilson factor is
\begin{equation}
 \exp\{\beta\operatorname{Re}(e^{i\theta}H(\phi))\}.
 \label{eq:v6v46-central-link-factor}
\end{equation}

\begin{proposition}[Bessel star weight]
\label{prop:v6v46-Bessel-star}
Normalized Haar integration of the central link gives
\begin{equation}
 \int_0^{2\pi}
 e^{\beta\operatorname{Re}(e^{i\theta}H)}\frac{d\theta}{2\pi}
 =I_0(\beta|H|).
 \label{eq:v6v46-Bessel-star}
\end{equation}
Hence the relative star law induced from independent staple phases has
density proportional to \(I_0(\beta|H(\phi)|)\).
\end{proposition}

\begin{proof}
Write \(H=|H|e^{-i\theta_0}\), shift
\(\theta\mapsto\theta+\theta_0\), and use
\begin{equation}
 I_0(z)=\frac1{2\pi}\int_0^{2\pi}e^{z\cos\theta}\,d\theta.
 \label{eq:v6v46-I0-integral}
\end{equation}
\end{proof}

The effective one-well potential at exponential scale is
\begin{equation}
 V_*(\phi):=n-|H(\phi)|.
 \label{eq:v6v46-star-potential}
\end{equation}

\begin{theorem}[Critical-point classification of the star well]
\label{thm:v6v46-star-critical}
Modulo common rotation, the only local minima of \(V_*\) are the aligned
configurations \(e^{i\phi_1}=\cdots=e^{i\phi_n}\), where \(V_*=0\).
Every critical configuration with \(H\ne0\) has each staple vector parallel
or antiparallel to \(H\).  A mixed-sign such configuration is not a local
minimum of \(V_*\).  Points with \(H=0\) are global maxima of \(V_*\), not
additional wells.
\end{theorem}

\begin{proof}
Differentiate \(|H|^2\):
\begin{equation}
 \partial_{\phi_j}|H|^2
 =2\operatorname{Im}(ie^{i\phi_j}\overline H)
 =-2\operatorname{Im}(e^{i\phi_j}\overline H).
 \label{eq:v6v46-star-derivative}
\end{equation}
At a critical point with \(H\ne0\), every
\(e^{i\phi_j}\overline H\) is real, so the vector is parallel or
antiparallel to \(H\).

Rotate so these vectors are signs \(s_j\in\{1,-1\}\), and suppose
\(M=\sum_js_j>0\).  For small perturbations \(\delta_j\),
\begin{equation}
 \left|\sum_js_je^{i\delta_j}\right|
 =M-\frac12\sum_js_j\delta_j^2
 +\frac1{2M}\left(\sum_js_j\delta_j\right)^2
 +O(|\delta|^3).
 \label{eq:v6v46-signed-expansion}
\end{equation}
If all signs are positive, restrict away from the common rotation by
\(\sum_j\delta_j=0\); the quadratic term is
\(-\frac12\sum_j\delta_j^2\), so \(|H|\) has a strict maximum and \(V_*\)
has a strict minimum.  If a negative sign is present, perturb only that vector by a small amount
\(a\).  Formula \eqref{eq:v6v46-signed-expansion} gives the positive
quadratic change \(a^2/2+a^2/(2M)\).  Subtracting the common component of
this perturbation is a common rotation and leaves \(|H|\) unchanged.  Thus
the mixed point is not a local maximum of
\(|H|\), hence not a local minimum of \(V_*\).  Finally
\(0\le|H|\le n\), so \(H=0\) gives the largest value \(V_*=n\).
\end{proof}

\subsection{Exponential suppression of small staples}

Let \(\mu_{\beta,n}\) be the probability measure on \(\mathbb T^n\) with
density proportional to \(I_0(\beta|H|)\).  Use the elementary bounds
\begin{equation}
 I_0(z)\le e^z,
 \qquad
 I_0(z)\ge c\frac{e^z}{\sqrt{1+z}}quad(z\ge0).
 \label{eq:v6v46-I0-bounds}
\end{equation}
The lower bound follows by integrating
\eqref{eq:v6v46-I0-integral} on \(|\theta|\le(1+z)^{-1/2}\) and using
\(\cos\theta\ge1-\theta^2/2\).

\begin{theorem}[Bad-staple mass has the full action barrier]
\label{thm:v6v46-bad-mass}
For every fixed \(0\le\epsilon<n\), there are \(C,c>0\) such that
\begin{equation}
 \mu_{\beta,n}(|H|\le\epsilon)
 \le C\beta^{n/2}e^{-\beta(n-\epsilon-c/\beta)}
 \label{eq:v6v46-bad-mass}
\end{equation}
for all \(\beta\ge1\).  In particular, for the four-dimensional star
\(n=6\), this is at most \(C\beta^3e^{-\beta(6-\epsilon)+C}\).
\end{theorem}

\begin{proof}
The numerator of the bad event is at most
\((2\pi)^ne^{\beta\epsilon}\) by the upper bound in
\eqref{eq:v6v46-I0-bounds}.  For the denominator, use relative coordinates
with \(|x|\le c_0\beta^{-1/2}\).  Lemma
\ref{lem:v6v46-resultant-distance} gives
\(|H|\ge n-C|x|^2\ge n-C/\beta\).  The lower Bessel bound is therefore
\begin{equation}
 I_0(\beta|H|)
 \ge c\beta^{-1/2}e^{\beta n-C}.
 \label{eq:v6v46-core-I0}
\end{equation}
The relative ball has dimension \(n-1\) and volume
\(c\beta^{-(n-1)/2}\); the common phase has fixed volume.  Hence the
partition function is at least
\(c\beta^{-n/2}e^{\beta n-C}\).  Divide the numerator by this lower bound.
\end{proof}

This theorem proves concentration but, by V43, not the required capacity.

\subsection{The scale-covariant Gaussian reference capacity}

Let \(\gamma_{\beta,m}\) be the centered Gaussian law on
\(\mathbb R^m\) with density proportional to
\(e^{-\beta|x|^2/2}\) and Dirichlet form
\begin{equation}
 \mathcal E_\beta(f)=\int|\nabla f|^2d\gamma_{\beta,m}.
 \label{eq:v6v46-Gaussian-form}
\end{equation}
Its Poincare gap is \(\beta\).

For disjoint measurable sets \(A,B\), define the condenser capacity
\begin{equation}
 \operatorname{Cap}_\beta(A,B)
 :=\inf\{\mathcal E_\beta(f):0\le f\le1,
 f|_A=1, f|_B=0\}.
 \label{eq:v6v46-condenser-capacity}
\end{equation}

\begin{lemma}[Gap gives a set-capacity lower bound]
\label{lem:v6v46-gap-capacity}
If a probability law has Poincare gap \(\lambda\), then
\begin{equation}
 \operatorname{Cap}(A,B)
 \ge\lambda\,\nu(A)\nu(B).
 \label{eq:v6v46-gap-capacity}
\end{equation}
\end{lemma}

\begin{proof}
For every admissible \(f\), the independent-copy formula gives
\begin{equation}
 \operatorname{Var}(f)
 =\frac12\iint|f(x)-f(y)|^2d\nu(x)d\nu(y)
 \ge\nu(A)\nu(B).
\end{equation}
The Poincare inequality gives
\(\mathcal E(f)\ge\lambda\operatorname{Var}(f)\).  Take the infimum.
\end{proof}

\begin{corollary}[Gaussian bad-set capacity is favorable]
\label{cor:v6v46-Gaussian-capacity}
Let \(B\) be a Gaussian core of mass at least one half and let \(A\) be
any disjoint bad set.  Then
\begin{equation}
 \operatorname{Cap}_\beta(A,B)
 \ge\frac\beta2\gamma_{\beta,m}(A),
 \qquad
 \frac{\gamma_{\beta,m}(A)}
 {\operatorname{Cap}_\beta(A,B)}\le\frac2\beta.
 \label{eq:v6v46-Gaussian-capacity}
\end{equation}
The ratio has exactly the one-block scale required by UTC.
\end{corollary}

\begin{proof}
Apply Lemma~\ref{lem:v6v46-gap-capacity} with
\(\lambda=\beta\) and \(\gamma_{\beta,m}(B)\ge1/2\).
\end{proof}

\subsection{The weighted star-capacity card}

Fix an aligned core
\begin{equation}
 \mathcal K_{\beta,L}
 :=\{\operatorname{dist}(\phi,\mathcal M_{\rm align})le L/\sqrt\beta\}
 \label{eq:v6v46-aligned-core}
\end{equation}
with mass at least one half, and let
\(\mathcal A_\epsilon=\{|H|\le\epsilon\}\).

\begin{definition}[Weighted star-capacity card]
\label{def:v6v46-WSC}
The WSC holds if, uniformly along the geometric path and in admissible
physical exterior data,
\begin{equation}
 \operatorname{Cap}_{\nu_{\beta,u}^\xi}
 (\mathcal A_\epsilon,\mathcal K_{\beta,L})
 \ge c\beta\,
 \nu_{\beta,u}^\xi(\mathcal A_\epsilon)
 \label{eq:v6v46-WSC}
\end{equation}
for all \(0<\epsilon\le\epsilon_0\), with the quotient Dirichlet form and
boundary modes retained.
\end{definition}

\begin{proposition}[WSC removes the pointwise staple obstruction]
\label{prop:v6v46-WSC-use}
Assume the aligned core has its normalized Hessian gap, WSC holds on every
bad-staple tube in a star block, and the intermediate staple region has the
monotone one-well drift bound.  Then the loss of the one-link conditional
gap at \(H=0\) does not create a block Poincare bottleneck at scale
\(\beta\).  The resulting star block is eligible for the BMI tensorization
of V42.
\end{proposition}

\begin{proof}
The core gap controls fluctuations supported in
\(\mathcal K_{\beta,L}\).  The monotone region is controlled by the radial
Hardy argument of V44.  WSC gives the measure--capacity inequality with
constant \(O(\beta^{-1})\) for functions separating a bad tube from the
core.  Applying this decomposition to the positive and negative level sets
of a centered function yields a block Poincare inequality of scale
\(O(\beta^{-1})\).  This is the standard level-set proof of the
measure--capacity criterion; every crossing of the low-staple region is
charged by \eqref{eq:v6v46-WSC}, rather than by a pointwise conditional
link gap.  Theorem~\ref{thm:v6v42-block-tensorization} then applies once its
separate BMI hypothesis is verified.
\end{proof}

\begin{assessment}[What has and has not transferred]
\label{ass:v6v46-transfer}
The integrated star has no second local minimum and the bad-staple mass
contains the full barrier \(e^{-\beta(n-\epsilon)}\).  The Gaussian tangent
has the exact favorable capacity ratio \(O(\beta^{-1})\).  These facts make
WSC plausible, but they do not prove it for the nonlinear compact star:
global density comparison between that law and its Gaussian tangent again
has an order-\(\beta\) oscillation.
\end{assessment}

\begin{firewall}[Exponential mass is not WSC]
\label{fw:v6v46-mass-not-capacity}
Theorem~\ref{thm:v6v46-bad-mass} cannot be substituted for
\eqref{eq:v6v46-WSC}.  The rare-neck counterexample in V43 has the same
logical form: a small set may be separated from the core by an even smaller
conductance.  A weighted coarea or flow proof is still required.
\end{firewall}

\begin{firewall}[The star law omits shared-plaquette constraints]
\label{fw:v6v46-star-model}
The density \(I_0(\beta|H|)\) results from integrating one central abelian
link with independent staple phases.  In the full lattice, neighboring
stars share links and plaquettes, Bianchi constraints correlate their
relative phases, and chiral/matter/Einstein terms modify the block weight.
Those effects belong in the WSC and BMI hypotheses and are not controlled
by the present calculation.
\end{firewall}

\begin{theorem}[Sixth-series V46 result]
\label{thm:v6v46-status}
Integrating a central abelian link gives the exact star weight
\(I_0(\beta|H|)\).  Its exponential-scale potential
\(n-|H|\) has only the aligned minimum modulo common rotation; all mixed
critical configurations are unstable and zero staples lie at the top of
the barrier.  The bad-staple mass is exponentially small with exponent
\(\beta(n-\epsilon)\).  The Gaussian relative-star reference has capacity
at least \(\beta\) times bad mass, identifying the correct
scale-covariant WSC.  Proving that weighted capacity for the interacting
compact star, rather than estimating its mass again, is the remaining block
gap step.
\end{theorem}

\begin{proof}
Use Lemma~\ref{lem:v6v46-resultant-distance},
Proposition~\ref{prop:v6v46-Bessel-star},
Theorem~\ref{thm:v6v46-star-critical},
Theorem~\ref{thm:v6v46-bad-mass},
Corollary~\ref{cor:v6v46-Gaussian-capacity}, and
Proposition~\ref{prop:v6v46-WSC-use}, subject to both firewalls.
\end{proof}

V47 will attack WSC by constructing an explicit weighted flow from each
small-staple tube to the aligned core.  It will use the gradient ascent of
\(|H|\), classify its singular strata, and test whether the coarea Jacobian
can be bounded with one exponential action payer and no inverse-staple
singularity.


\clearpage
\section{V47: latent-link quotient tensorization closes the independent-star capacity}
\label{sec:v6v47}

V46 proposed a direct weighted flow for the nonlinear Bessel star.  There
is a stronger exact route.  The central link is a latent common group
rotation, and after it is integrated the relative star law is the diagonal
quotient of independent one-plaquette laws.  The product gap therefore
descends to the star quotient and immediately implies the weighted
capacity card.

\subsection{A compact-group latent quotient identity}

Let \(G\) be a compact group with normalized Haar measure and let
\(q:G\to(0,\infty)\) be a normalized density.  For \(n\ge2\), define the
star weight on \(G^n\) by
\begin{equation}
 W_q(U_1,\ldots,U_n)
 :=\int_G\prod_{j=1}^nq(gU_j)\,dg.
 \label{eq:v6v47-star-weight}
\end{equation}
It is invariant under diagonal left multiplication
\((U_1,\ldots,U_n)\mapsto(hU_1,\ldots,hU_n)\).

\begin{theorem}[Exact latent quotient law]
\label{thm:v6v47-latent-quotient}
Let \(\mu_q\) be the probability law with density \(W_q\) relative to
Haar measure on \(G^n\).  For every bounded diagonal-left-invariant
observable \(F\),
\begin{equation}
 \int_{G^n}F(U)W_q(U)\,dU
 =\int_{G^n}F(V)\prod_{j=1}^nq(V_j)\,dV.
 \label{eq:v6v47-latent-quotient}
\end{equation}
Thus the relative star law on \(G^n/G_{\rm diag}\) is exactly the pushforward
of \(q^{\otimes n}\) under the diagonal quotient map.
\end{theorem}

\begin{proof}
Insert \eqref{eq:v6v47-star-weight}, use Fubini, and for fixed \(g\) change
variables \(V_j=gU_j\).  Haar invariance gives \(dV=dU\), while diagonal
invariance gives
\(F(U_1,\ldots,U_n)=F(g^{-1}V_1,\ldots,g^{-1}V_n)=F(V)\).  The inner
integral is therefore independent of \(g\), and integration of normalized
Haar measure proves the identity.  Taking \(F=1\) also proves that
\(W_q\) is normalized.
\end{proof}

For an abelian Wilson star, take
\(q_\beta(e^{i\theta})=e^{\beta\cos\theta}/I_0(\beta)\).
Equation \eqref{eq:v6v47-star-weight} then gives, up to its already fixed
normalization,
\begin{equation}
 \int_0^{2\pi}prod_j e^{\beta\cos(\theta+\phi_j)}
 \frac{d\theta}{2\pi}
 =I_0(\beta|H(\phi)|),
 \label{eq:v6v47-Bessel-as-quotient}
\end{equation}
which recovers Proposition~\ref{prop:v6v46-Bessel-star}.

\subsection{Quotient gap theorem}

Let \(\mathcal E_q\) be a left-invariant reversible Dirichlet form for
\(q\,dU\).  On the product use
\begin{equation}
 \mathcal E_q^{(n)}(F,F)
 :=\sum_{j=1}^n\int|\nabla_jF|^2\prod_kq(U_k)\,dU.
 \label{eq:v6v47-product-energy}
\end{equation}
The quotient form is the restriction of this form to diagonal-invariant
functions.  Equivalently, it is the horizontal form obtained after removing
the common left-rotation direction.

\begin{theorem}[One-plaquette gap descends to the star]
\label{thm:v6v47-star-gap}
If \(q\,dU\) has Poincare gap \(\lambda_q\), then the relative star law
\(\mu_q\) on \(G^n/G_{\rm diag}\) has gap at least \(\lambda_q\):
\begin{equation}
 \operatorname{Var}_{\mu_q}(F)
 \le\lambda_q^{-1}\mathcal E_{\rm star}(F,F)
 \label{eq:v6v47-star-gap}
\end{equation}
for every centered physical relative observable \(F\).
\end{theorem}

\begin{proof}
By Theorem~\ref{thm:v6v47-latent-quotient}, the variance of a relative
observable under the star law equals its variance under \(q^{\otimes n}\).
The product tensorization theorem
\ref{thm:v6v45-product-gap} bounds that variance by
\(\lambda_q^{-1}\mathcal E_q^{(n)}(F,F)\).  Since \(F\) is invariant under
the common rotation, its product gradient is horizontal and its product
energy is exactly the quotient star energy.  This proves the claim.
\end{proof}

\begin{remark}[The common mode is deliberately excluded]
\label{rem:v6v47-common-mode}
On the unreduced star, the common group rotation may have a gap of order one
or may be a gauge zero mode.  Theorem~\ref{thm:v6v47-star-gap} concerns the
physical relative quotient.  It does not assign the one-plaquette
\(O(\beta)\) gap to an unphysical common direction.
\end{remark}

\subsection{The entire abelian star path}

Let \(q_{\beta,u}\) be the normalized one-plaquette density of the
calibrated \(U(1)\) heat--Wilson geometric path in
\eqref{eq:v6v44-abelian-path-density}, transported from the class interval
to the circle.  Define
\begin{equation}
 W_{\beta,u}(U_1,ldots,U_n)
 :=\int_{U(1)}\prod_{j=1}^nq_{\beta,u}(gU_j)\,dg.
 \label{eq:v6v47-abelian-star-path}
\end{equation}

\begin{corollary}[Uniform abelian star-path gap]
\label{cor:v6v47-abelian-star-gap}
For every fixed \(n\), the relative star path
\eqref{eq:v6v47-abelian-star-path} has gap
\begin{equation}
 \lambda_{\rm star}(\beta,u)\ge c\beta
 \label{eq:v6v47-abelian-star-gap}
\end{equation}
uniformly in \(u\in[0,1]\) and \(\beta\ge1\).
\end{corollary}

\begin{proof}
Theorem~\ref{thm:v6v44-abelian-UTC} gives the one-plaquette gap
\(\lambda_{q_{\beta,u}}\ge c\beta\).  Apply
Theorem~\ref{thm:v6v47-star-gap}.
\end{proof}

This result includes the zero-staple configurations of V44.  Their
one-link conditional law may be uniform, but a relative star fluctuation
cannot localize there without paying the product one-plaquette energy in
the latent lift.

\subsection{Capacity consequence}

Let \(\mathcal K_{\beta,L}\) be the aligned relative core in
\eqref{eq:v6v46-aligned-core}.  For fixed \(n\), choose \(L\) so that its
star-path mass is at least one half, uniformly in \(u\) and large \(\beta\).

\begin{lemma}[Uniform aligned-core mass]
\label{lem:v6v47-core-mass}
Such an \(L\) exists for the abelian path.
\end{lemma}

\begin{proof}
The density envelopes used in
Theorem~\ref{thm:v6v44-abelian-UTC} give
\begin{equation}
 \mathbb E_{q_{\beta,u}}d(U,1)^2\le C\beta^{-1}
 \label{eq:v6v47-factor-second-moment}
\end{equation}
uniformly in \(u\).  Under the latent product representation, the squared
distance to the aligned quotient is bounded by the sum of the \(n\) factor
distances to the identity.  Markov's inequality gives
\[
 \mu_{\beta,u}^{\rm star}
 (\mathcal K_{\beta,L}^c)
 \le Cn/L^2.
\]
Choose \(L^2\ge2Cn\).
\end{proof}

\begin{theorem}[WSC for the independent abelian star path]
\label{thm:v6v47-WSC}
For every measurable relative bad set \(A\) disjoint from
\(\mathcal K_{\beta,L}\),
\begin{equation}
 \operatorname{Cap}_{\mu_{\beta,u}^{\rm star}}
 (A,\mathcal K_{\beta,L})
 \ge\frac{c\beta}{2}\mu_{\beta,u}^{\rm star}(A).
 \label{eq:v6v47-WSC}
\end{equation}
In particular, Definition~\ref{def:v6v46-WSC} holds for every
small-staple tube \(A=\{|H|\le\epsilon\}\), uniformly along the complete
independent abelian star path.
\end{theorem}

\begin{proof}
Corollary~\ref{cor:v6v47-abelian-star-gap} gives gap \(c\beta\), and
Lemma~\ref{lem:v6v47-core-mass} gives core mass at least one half.  Apply
the gap-capacity inequality \eqref{eq:v6v46-gap-capacity}.
\end{proof}

The weighted flow proposed at the end of V46 is therefore unnecessary for
the independent star.  The latent quotient proves the capacity with its
correct scale and includes all singular staple strata automatically.

\subsection{Nonabelian extension}

Theorems~\ref{thm:v6v47-latent-quotient} and
\ref{thm:v6v47-star-gap} are group-independent.  For central
one-plaquette densities, orientation reversals and left/right placements of
the central link are absorbed by inversion and Haar changes of variables.

\begin{corollary}[Conditional nonabelian star reduction]
\label{cor:v6v47-nonabelian-star}
For a compact group \(G\), any independent star path of the form
\eqref{eq:v6v47-star-weight} has a physical relative gap at least equal to
its one-plaquette factor gap.  Thus the audited \(SU(2)\) bridge gap would
give the corresponding independent \(SU(2)\) star gap without a new staple
analysis.  The \(SU(3)\) star remains reduced to the unacquired
one-plaquette Weyl-alcove gap.
\end{corollary}

\begin{proof}
Apply Theorem~\ref{thm:v6v47-star-gap}; the last sentence is the V45 audit
status.
\end{proof}

\subsection{The remaining lattice correction}

An actual hypercubic block need not have the exact density
\eqref{eq:v6v47-star-weight}.  Neighboring plaquette holonomies are linked
by Bianchi relations and by plaquettes leaving the star; matter and
determinant terms add further couplings.  Write the true conditional block
density relative to the latent star law as
\begin{equation}
 d\widetilde\mu_{\beta,u}^\xi
 =\frac{e^{-R_{\beta,u}^\xi}}{\widetilde Z_{\beta,u}^\xi}
 d\mu_{\beta,u}^{\rm star}.
 \label{eq:v6v47-star-remainder}
\end{equation}

\begin{definition}[Star remainder card]
\label{def:v6v47-SRC}
The SRC holds if the remainder is controlled in a way that preserves the
star gap or capacity uniformly in physical exterior data.  Two sufficient
forms are:
\begin{align}
 \operatorname{osc}R_{\beta,u}^\xi&\le C,
 \label{eq:v6v47-SRC-osc}\\
 \left\|
 (\mathcal H_{\rm star})^{-1/2}
 \nabla^2R_{\beta,u}^\xi
 (\mathcal H_{\rm star})^{-1/2}
 \right\|&\le1-\delta
 \label{eq:v6v47-SRC-Hessian}
\end{align}
on a core, together with the WSC analogue on its complement.
\end{definition}

\begin{proposition}[Bounded remainder preserves the star gap]
\label{prop:v6v47-bounded-remainder}
If \eqref{eq:v6v47-SRC-osc} holds, then
\begin{equation}
 \operatorname{gap}(\widetilde\mu_{\beta,u}^\xi)
 \ge e^{-C}\operatorname{gap}(\mu_{\beta,u}^{\rm star})
 \ge ce^{-C}\beta.
 \label{eq:v6v47-bounded-remainder-gap}
\end{equation}
\end{proposition}

\begin{proof}
The density ratio between the two normalized laws has oscillation at most
\(C\).  Apply the bounded perturbation comparison to the variance and
Dirichlet integrals, then use
Corollary~\ref{cor:v6v47-abelian-star-gap} or the corresponding nonabelian
factor theorem.
\end{proof}

\begin{firewall}[The true remainder is not known to be bounded]
\label{fw:v6v47-SRC-open}
At weak coupling, a plaquette or matter term crossing the star boundary can
have oscillation of order \(\beta\).  Therefore
\eqref{eq:v6v47-SRC-osc} is not assumed for the full lattice.  The
normalized Hessian/core plus weighted-capacity alternative in
\eqref{eq:v6v47-SRC-Hessian} is the live SRC route.
\end{firewall}

\begin{firewall}[Bianchi constraints are not product variables]
\label{fw:v6v47-Bianchi-open}
The latent quotient identity is exact for a star whose plaquette factors
are independent before quotienting the central link.  A four-dimensional
Bianchi constraint reduces their allowed holonomy manifold and can change
the quotient Dirichlet form.  A physical coordinate chart or exact
cochain-complex disintegration is required before applying the star theorem
to a hypercubic block.
\end{firewall}

\begin{theorem}[Sixth-series V47 result]
\label{thm:v6v47-status}
The independent compact-group star law is exactly the diagonal quotient of
a product of one-plaquette laws.  Its physical relative gap is therefore no
smaller than the one-plaquette gap.  Applied to the complete calibrated
abelian path, this gives a uniform \(c\beta\) star gap and proves WSC for
every small-staple tube, including all singular strata.  The construction
extends conditionally to \(SU(2)\) and \(SU(3)\) once their factor gaps are
available.  The remaining lattice problem is the SRC: control the Bianchi,
cross-boundary, matter, and determinant remainder relative to the latent
star without an order-\(\beta\) oscillation loss.
\end{theorem}

\begin{proof}
Use Theorem~\ref{thm:v6v47-latent-quotient},
Theorem~\ref{thm:v6v47-star-gap},
Corollary~\ref{cor:v6v47-abelian-star-gap},
Theorem~\ref{thm:v6v47-WSC}, and
Proposition~\ref{prop:v6v47-bounded-remainder}, subject to both firewalls.
\end{proof}

V48 will construct the exact abelian cochain disintegration on a contractible
hypercubic block.  It will separate independent plaquette coordinates,
Bianchi constraints, boundary flux, and harmonic modes, and determine the
precise SRC produced by gluing the star to its neighboring cubes.


\clearpage
\section{V48: abelian cochain disintegration, Bianchi restriction, and the photon reconstruction cost}
\label{sec:v6v48}

The SRC in V47 grouped Bianchi relations with ordinary density remainders.
They are different.  A Bianchi identity restricts the support of the
plaquette law to a closed cochain subspace; it is not a bounded potential.
On a contractible block the link quotient is exactly isomorphic to that
subspace.  The Gaussian restriction preserves the curvature gap, while
reconstruction of a link field costs one inverse derivative and recovers
the physical photon shell scale.

\subsection{Exact finite cochain complex}

Let \(K\) be a finite contractible cubical complex with real cochain spaces
\(C^k(K)\) and coboundaries
\begin{equation}
 C^0(K)\xrightarrow{d_0}C^1(K)
 \xrightarrow{d_1}C^2(K)
 \xrightarrow{d_2}C^3(K),
 \qquad d_1d_0=0,quad d_2d_1=0.
 \label{eq:v6v48-cochain-complex}
\end{equation}
Use the standard cell inner products.  Contractibility gives
\begin{equation}
 \ker d_1=\operatorname{im}d_0,
 \qquad
 \ker d_2=\operatorname{im}d_1.
 \label{eq:v6v48-exactness}
\end{equation}
Constants are the kernel of \(d_0\).

\begin{theorem}[Link quotient equals closed curvature]
\label{thm:v6v48-link-curvature-isomorphism}
The curvature map induces a linear isomorphism
\begin{equation}
 \overline d_1:
 C^1(K)/\operatorname{im}d_0
 \xrightarrow{\ \simeq\ }
 Z^2(K):=\ker d_2,
 \qquad [A]\longmapsto d_1A.
 \label{eq:v6v48-link-curvature-isomorphism}
\end{equation}
The unique minimum-norm representative of the gauge class with curvature
\(F\in Z^2(K)\) is
\begin{equation}
 A_F=d_1^*(d_1d_1^*)^{-1}_{Z^2}F,
 \label{eq:v6v48-minimal-reconstruction}
\end{equation}
where the inverse is taken on \(Z^2=\operatorname{im}d_1\).
\end{theorem}

\begin{proof}
The kernel of \(A\mapsto d_1A\) is
\(\ker d_1=\operatorname{im}d_0\), so the induced quotient map is
injective.  Its range is \(\operatorname{im}d_1=\ker d_2\), so it is
surjective.  The orthogonal complement of the gauge directions is
\begin{equation}
 (\ker d_1)^\perp=\operatorname{im}d_1^*.
 \label{eq:v6v48-Coulomb-space}
\end{equation}
On this space \(d_1\) is bijective onto \(Z^2\).  For
\(F\in Z^2\), formula \eqref{eq:v6v48-minimal-reconstruction} lies in
\(\operatorname{im}d_1^*\) and
\[
 d_1A_F=d_1d_1^*(d_1d_1^*)^{-1}_{Z^2}F=F.
\]
It is orthogonal to every gauge direction and is therefore the unique
minimum-norm representative.
\end{proof}

\begin{corollary}[Bianchi identity is the complete local constraint]
\label{cor:v6v48-Bianchi-complete}
On a contractible real block, a two-cochain is the curvature of a link
field exactly when
\begin{equation}
 d_2F=0.
 \label{eq:v6v48-Bianchi}
\end{equation}
There is no additional harmonic two-form.  On a noncontractible or compact
angle block, harmonic and integer-flux sectors must be added explicitly.
\end{corollary}

\subsection{Gauge-fixed measure disintegration}

Restrict \(d_1\) to the Coulomb subspace
\(\mathscr C^1=(\operatorname{im}d_0)^\perp\), and let
\begin{equation}
 J_K:=\left|det(d_1|_{\mathscr C^1}:\mathscr C^1\to Z^2)\right|^{-1}.
 \label{eq:v6v48-cochain-Jacobian}
\end{equation}

\begin{proposition}[Constant-Jacobian curvature disintegration]
\label{prop:v6v48-measure-disintegration}
For every integrable gauge-invariant function \(G(d_1A)\),
\begin{equation}
 \int_{\mathscr C^1}G(d_1A)\,dA
 =J_K\int_{Z^2}G(F)\,dF.
 \label{eq:v6v48-measure-disintegration}
\end{equation}
Thus a plaquette product density becomes its restriction to the linear
Bianchi subspace with a field-independent Jacobian.  No oscillatory SRC
potential is created by the real linear change of variables.
\end{proposition}

\begin{proof}
Theorem~\ref{thm:v6v48-link-curvature-isomorphism} makes
\(d_1|_{\mathscr C^1}\) a linear bijection.  Apply the finite-dimensional
change-of-variables formula; its determinant is constant.
\end{proof}

With fixed boundary links, use the relative cochain complex
\(C^\bullet(K,\partial K)\).  For a four-dimensional topological ball,
its first and second relative cohomology vanish, so the same argument gives
an affine Bianchi subspace
\begin{equation}
 F=F_\xi+d_1^{\rm rel}A_{\rm int},
 \qquad
 d_2F=0,
 \label{eq:v6v48-affine-boundary-curvature}
\end{equation}
provided the prescribed boundary flux is compatible.  Incompatible or
unfixed flux belongs to the extended boundary sectors of V36.

\subsection{Compact \(U(1)\) lifts and flux sectors}

For angle-valued links, choose real lifts on a local chart.  Their
plaquette curvatures have the form
\begin{equation}
 \widetilde F=d_1\widetilde A+2\pi n,
 \qquad n\in C^2(K;\mathbb Z),
 \label{eq:v6v48-compact-curvature}
\end{equation}
with the compact Bianchi condition
\begin{equation}
 d_2\widetilde F=2\pi d_2n.
 \label{eq:v6v48-compact-Bianchi}
\end{equation}

\begin{assessment}[Local lift versus topological sectors]
\label{ass:v6v48-compact-sectors}
On a contractible small-field block with fixed admissible boundary flux, a
unique lift sector may be selected and the real disintegration applies.
Across chart overlaps or nontrivial topology, the integer cochain and
harmonic flux labels are part of the physical boundary data.  Summing them
after separate absolute bounds would lose the compact interference and is
not included in Proposition~\ref{prop:v6v48-measure-disintegration}.
\end{assessment}

\subsection{Gaussian Bianchi restriction}

Let \(Z\subset\mathbb R^P\) be any linear subspace and let
\(\gamma_{\beta,Z}\) be the centered Gaussian on \(Z\) with density
proportional to \(e^{-\beta\|F\|^2/2}\) relative to Lebesgue measure on
\(Z\).

\begin{theorem}[Bianchi restriction preserves the curvature gap]
\label{thm:v6v48-Gaussian-Bianchi-gap}
For every linear subspace \(Z\), including \(Z=\ker d_2\),
\begin{equation}
 \operatorname{Var}_{\gamma_{\beta,Z}}(g)
 \le\frac1\beta
 \int_Z\|\nabla_Zg\|^2d\gamma_{\beta,Z}.
 \label{eq:v6v48-Gaussian-Bianchi-gap}
\end{equation}
The constant is independent of the codimension and block volume.
\end{theorem}

\begin{proof}
Choose an orthonormal basis of \(Z\).  In those coordinates
\(\gamma_{\beta,Z}\) is a product of standard one-dimensional Gaussians
of precision \(\beta\).  Tensorization gives gap \(\beta\).
\end{proof}

Thus linear Bianchi constraints do not degrade the Gaussian curvature gap.
The inverse derivative appears only when one asks for the link-field metric.

\subsection{Exact photon reconstruction scale}

On the Coulomb space, the quadratic link action is
\begin{equation}
 S_2(A)=\frac\beta2\|d_1A\|^2
 =\frac12\langle A,\beta d_1^*d_1A\rangle.
 \label{eq:v6v48-link-Gaussian-action}
\end{equation}
Let
\(\sigma_{\min}^+(d_1)\) be the least nonzero singular value of
\(d_1\).

\begin{theorem}[Curvature-to-link gap conversion]
\label{thm:v6v48-link-gap}
The Gaussian law in link Coulomb coordinates has exact spectral gap
\begin{equation}
 \lambda_{A,K}
 =\beta[\sigma_{\min}^+(d_1)]^2
 =\beta\lambda_{\min}^+(d_1^*d_1).
 \label{eq:v6v48-link-gap}
\end{equation}
Equivalently, the reconstruction map satisfies
\begin{equation}
 \|A_F\|
 \le[\sigma_{\min}^+(d_1)]^{-1}\|F\|.
 \label{eq:v6v48-reconstruction-cost}
\end{equation}
On blocks of growing diameter \(L\), this inverse derivative can scale like
\(L\), so the link gap can scale like \(\beta L^{-2}\) even though the
curvature gap is \(\beta\).
\end{theorem}

\begin{proof}
Diagonalize the positive operator \(d_1^*d_1\) on
\(\mathscr C^1\).  A Gaussian with precision
\(\beta d_1^*d_1\) has gap equal to its smallest eigenvalue, proving the
first identity.  The operator norm of the inverse
\(d_1|_{\mathscr C^1}^{-1}\) is the reciprocal of its least singular
value, which proves \eqref{eq:v6v48-reconstruction-cost}.  Standard
long-wavelength one-form modes on a diameter-\(L\) block have derivative
size \(O(L^{-1})\), giving the stated possible scaling.
\end{proof}

\begin{corollary}[The photon cost is physical]
\label{cor:v6v48-photon-cost}
The favorable product/star curvature gap of V47 cannot be assigned
unchanged to link potentials on arbitrarily large blocks.  The inverse
derivative in \eqref{eq:v6v48-minimal-reconstruction} restores the
massless photon scale.  Fixed-diameter RG blocks may retain a uniform
multiple of \(\beta\); growing blocks require shellwise coercivity and an
explicit low-momentum photon edge as in V42.
\end{corollary}

\begin{proof}
Apply Theorem~\ref{thm:v6v48-link-gap} and compare fixed versus growing
least singular values.
\end{proof}

\subsection{The nonlinear Bianchi-restriction card}

Let \(q_{\beta,u}\) be a one-plaquette path density and consider its product
law on \(C^2(K)\).  The actual curvature law in a lift sector is
\begin{equation}
 d\nu_{K,u}^{\rm curv}(F)
 \propto
 \mathbf1_{\{F\in F_\xi+\operatorname{im}d_1^{\rm rel}\}}
 \prod_{p\subset K}q_{\beta,u}(F_p)\,dF_Z.
 \label{eq:v6v48-restricted-product}
\end{equation}
Here \(dF_Z\) is Lebesgue measure on the affine Bianchi subspace, so the
indicator is notation for support restriction, not a density with respect
to ambient Lebesgue measure.

\begin{definition}[Bianchi restriction card]
\label{def:v6v48-BRC}
The BRC holds if the physical curvature law
\eqref{eq:v6v48-restricted-product} has a gap
\begin{equation}
 \operatorname{gap}(\nu_{K,u}^{\rm curv})\ge c\beta
 \label{eq:v6v48-BRC}
\end{equation}
uniformly in the path, admissible boundary flux, lift sector, and fixed
block geometry.  For growing blocks the card is stated in the induced
shell metric before link reconstruction.
\end{definition}

\begin{proposition}[BRC plus reconstruction gives the block row]
\label{prop:v6v48-BRC-use}
If BRC holds and the quotient reconstruction singular value on a block is
\(\sigma_K\), then the corresponding link-coordinate block gap is at least
\begin{equation}
 c\beta\sigma_K^2.
 \label{eq:v6v48-BRC-link-gap}
\end{equation}
Together with a strict signed BMI between blocks, this supplies the UPG on
the gapped shells while retaining low photon modes.
\end{proposition}

\begin{proof}
The chain rule under \(F=d_1A\) transforms curvature gradients by the
adjoint of \(d_1^{-1}\).  Its largest norm is \(\sigma_K^{-1}\), so the
Poincare constant grows by at most \(\sigma_K^{-2}\).  Apply
Theorem~\ref{thm:v6v42-block-tensorization} to the resulting block gaps.
\end{proof}

\begin{assessment}[Correction to SRC]
\label{ass:v6v48-SRC-correction}
Bianchi identities should not be placed in the ordinary remainder
\(R_{\beta,u}^\xi\) of \eqref{eq:v6v47-star-remainder}.  Their exact owner
is the support restriction and constant-Jacobian cochain disintegration
\eqref{eq:v6v48-measure-disintegration}.  Cross-boundary plaquettes,
matter, determinant, and Einstein terms remain true density or operator
remainders and still belong to SRC/CRRC.
\end{assessment}

\begin{firewall}[Nonlinear BRC remains open]
\label{fw:v6v48-BRC-open}
Theorem~\ref{thm:v6v48-Gaussian-Bianchi-gap} proves BRC only for the
quadratic tangent.  Restricting a nonconvex compact one-well product law to
an affine Bianchi subspace can change its critical-point and capacity
geometry.  No uniform nonlinear BRC has yet been proved.
\end{firewall}

\begin{firewall}[Flux-sector sums remain physical]
\label{fw:v6v48-flux-open}
The real contractible disintegration does not bound sums over integer
magnetic flux, harmonic torus modes, or incompatible boundary sectors.
Those labels must be retained through the compact gauge gluing and cannot
be removed by the local Jacobian constant.
\end{firewall}

\begin{theorem}[Sixth-series V48 result]
\label{thm:v6v48-status}
On a contractible abelian block, the link quotient is exactly the closed
curvature subspace \(\ker d_2\), and gauge-fixed Lebesgue measure pushes
forward with a constant Jacobian.  Linear Bianchi restriction preserves
the Gaussian curvature gap \(\beta\).  Reconstructing link potentials
costs the inverse singular value of one lattice derivative, giving the
exact physical gap \(\beta[\sigma_{\min}^+(d_1)]^2\) and restoring the
low-momentum photon scale on large blocks.  The remaining nonlinear
constraint problem is the BRC, while cross-boundary and matter effects
remain genuine SRC rows.
\end{theorem}

\begin{proof}
Use Theorem~\ref{thm:v6v48-link-curvature-isomorphism},
Proposition~\ref{prop:v6v48-measure-disintegration},
Theorem~\ref{thm:v6v48-Gaussian-Bianchi-gap},
Theorem~\ref{thm:v6v48-link-gap}, and
Proposition~\ref{prop:v6v48-BRC-use}, subject to both firewalls.
\end{proof}

V49 will test nonlinear BRC on the smallest cube carrying a Bianchi
relation.  It will classify the constrained one-well critical points,
compute the tangent restricted Hessian, and decide whether the compact
constraint creates a new metastable well or only the already retained flux
sectors.


\clearpage
\section{V49: compact cube Bianchi wells and the all-flux BRC no-go}
\label{sec:v6v49}

V48 proves that a linear Bianchi restriction does not degrade a Gaussian
curvature gap.  Compact \(U(1)\) curvature has additional lift sectors.
On the smallest three-cell carrying a Bianchi relation, the full compact
constraint manifold already has metastable nonzero-flux minima.  Their
local Hessians are order \(\beta\), but their exponentially small
conductance collapses the global all-sector gap.

\subsection{The six-face cube model}

Orient the six plaquette angles of one cube so the compact Bianchi relation
is
\begin{equation}
 e^{i(x_1+cdots+x_6)}=1.
 \label{eq:v6v49-compact-cube-Bianchi}
\end{equation}
Changing orientations of individual faces replaces \(x_i\) by \(-x_i\)
and leaves the Wilson action below unchanged.  The constraint manifold is
\begin{equation}
 \mathcal M:=\left\{x\in\mathbb T^6:
 e^{i\sum_ix_i}=1\right\},
 \label{eq:v6v49-cube-manifold}
\end{equation}
and the one-cube action is
\begin{equation}
 S(x):=\sum_{i=1}^6(1-\cos x_i).
 \label{eq:v6v49-cube-action}
\end{equation}
Let \(\nu_\beta\) have density proportional to \(e^{-\beta S}\) on
\(\mathcal M\) with its induced Haar measure.

In a principal lift chart, the constraint reads
\begin{equation}
 \sum_{i=1}^6x_i=2\pi m,
 \qquad m\in\mathbb Z.
 \label{eq:v6v49-lift-sector}
\end{equation}
The integer \(m\) changes only when a plaquette representative crosses the
chart boundary; it is the local compact flux/monopole label.

\subsection{Critical points and constrained Hessian}

\begin{lemma}[Critical-point equation]
\label{lem:v6v49-critical-equation}
At an interior critical point of a lift sector,
\begin{equation}
 \sin x_1=\cdots=\sin x_6=\lambda
 \label{eq:v6v49-critical-equation}
\end{equation}
for one Lagrange multiplier \(\lambda\).  The constrained Hessian on
\begin{equation}
 T_x\mathcal M=\{\delta\in\mathbb R^6:\sum_i\delta_i=0\}
 \label{eq:v6v49-tangent-space}
\end{equation}
is
\begin{equation}
 Q_x(\delta)=\sum_{i=1}^6\cos x_i\,\delta_i^2.
 \label{eq:v6v49-constrained-Hessian}
\end{equation}
\end{lemma}

\begin{proof}
Apply a Lagrange multiplier to \(S\) under the lifted linear constraint.
Since \(\partial_{x_i}S=\sin x_i\), all first derivatives are equal.  The
ambient Hessian is diagonal with entries \(\cos x_i\); restrict it to the
tangent hyperplane.
\end{proof}

\begin{theorem}[Three explicit compact Bianchi minima]
\label{thm:v6v49-three-minima}
The compact constraint manifold has the critical points
\begin{equation}
 x^{(m)}=\left(\frac{\pi m}{3},\ldots,
 \frac{\pi m}{3}\right),
 \qquad m=0,\pm1.
 \label{eq:v6v49-uniform-flux-points}
\end{equation}
They satisfy
\begin{align}
 S(x^{(0)})&=0,
 &Q_{x^{(0)}}(\delta)&=\|\delta\|^2,
 \label{eq:v6v49-zero-minimum}\\
 S(x^{(\pm1)})&=3,
 &Q_{x^{(\pm1)}}(\delta)&=\frac12\|\delta\|^2
 \quad(\delta\in T\mathcal M).
 \label{eq:v6v49-flux-minima}
\end{align}
Thus \(x^{(0)}\) is the global minimum and
\(x^{(\pm1)}\) are strict metastable local minima with local quadratic gap
\(\beta/2\).
\end{theorem}

\begin{proof}
The sum of the six coordinates of \(x^{(m)}\) is \(2\pi m\), so every
point lies on \(\mathcal M\).  All sine values are equal, hence the points
are critical by Lemma~\ref{lem:v6v49-critical-equation}.  Substitute
\(0\) and \(\pm\pi/3\) into the action and Hessian.  Since their cosine
values are respectively one and one half, the Hessian is positive definite
on the tangent hyperplane.  Nonnegativity of every action summand makes
\(x^{(0)}\) global.
\end{proof}

The first saddle leaving the symmetric \(m=1\) well is visible in the same
critical equation.  Take five angles \(a=\pi/4\) and one angle
\(\pi-a=3\pi/4\).  Their sum is \(2\pi\), and their action is
\begin{equation}
 S_{\rm sad}=5(1-2^{-1/2})+(1+2^{-1/2})
 =6-2\sqrt2>3.
 \label{eq:v6v49-symmetric-saddle}
\end{equation}
The positive difference \(3-2\sqrt2\) is the symmetric escape barrier, but
the gap obstruction below needs only strict local minimality.

\subsection{Exponential collapse of the all-sector gap}

\begin{theorem}[Metastable flux well destroys an order-\(\beta\) global
gap]
\label{thm:v6v49-metastable-gap}
There are \(c,C,\delta>0\) such that the spectral gap of \(\nu_\beta\) on
the full compact constraint manifold satisfies
\begin{equation}
 \operatorname{gap}(\nu_\beta)
 \le C\beta^C e^{-\delta\beta}
 \label{eq:v6v49-metastable-gap}
\end{equation}
for all sufficiently large \(\beta\).  Hence the all-flux nonlinear BRC
\eqref{eq:v6v48-BRC} is false already on one cube.
\end{theorem}

\begin{proof}
By Theorem~\ref{thm:v6v49-three-minima}, choose nested neighborhoods
\(U_0\Subset U_1\) of \(x^{(1)}\) and \(\delta>0\) such that
\begin{equation}
 S\le3+\frac\delta4\quad\hbox{on }U_0,
 \qquad
 S\ge3+\delta\quad\hbox{on }
 \operatorname{supp}\nabla\chi,
 \label{eq:v6v49-cutoff-barrier}
\end{equation}
for a smooth cutoff \(\chi=1\) on \(U_0\), supported in \(U_1\).
This is possible because a strict local minimum has an isolating sublevel
neighborhood.

Laplace bounds in the five-dimensional tangent chart give
\begin{equation}
 \nu_\beta(U_0)ge c\beta^{-C}e^{-3\beta}/Z_\beta,
 \qquad
 \mathcal E_\beta(\chi)le C e^{-(3+\delta)\beta}/Z_\beta.
 \label{eq:v6v49-Laplace-well}
\end{equation}
The global minimum gives
\(Z_\beta\ge c\beta^{-5/2}\), so the metastable mass tends to zero.
For large \(\beta\), therefore,
\begin{equation}
 \operatorname{Var}_{\nu_\beta}(\chi)
 \ge c\nu_\beta(U_0)ge c\beta^{-C}e^{-3\beta}/Z_\beta.
 \label{eq:v6v49-cutoff-variance}
\end{equation}
The Rayleigh quotient of \(\chi\), using
\eqref{eq:v6v49-Laplace-well}, is at most a polynomial times
\(e^{-\delta\beta}\).  This proves the gap upper bound.  An order-\(\beta\)
BRC is incompatible with it.
\end{proof}

\begin{assessment}[A locally strong well can be globally slow]
\label{ass:v6v49-local-global}
The flux minima have positive Hessian \(\beta/2\), so every fixed-well
tangent calculation sees the correct strong local scale.  The global gap
is nevertheless exponentially small because changing the lift sector
requires crossing a higher-action neck.  This is the compact gauge version
of the rare-neck obstruction in V43.
\end{assessment}

\subsection{The zero-lift sector has no interior secondary minimum}

\begin{lemma}[One-negative-direction test]
\label{lem:v6v49-negative-direction}
At a critical point satisfying \eqref{eq:v6v49-critical-equation}, if at
least two values \(\cos x_i\) are negative, the constrained Hessian is not
positive.  If exactly one is negative and the remaining five arise from
the complementary sine branch, the constrained Hessian is also not
positive.
\end{lemma}

\begin{proof}
With two negative diagonal entries, choose a nonzero tangent vector
supported on those two coordinates with opposite values; its quadratic
form is negative.  With exactly one negative branch, all cosine magnitudes
are equal: five entries are \(c>0\) and one is \(-c\).  Set the negative
coordinate to \(-5a\) and every positive coordinate to \(a\).  The vector
is tangent and
\begin{equation}
 Q=5ca^2-c(25a^2)=-20ca^2<0.
 \label{eq:v6v49-one-negative-test}
\end{equation}
\end{proof}

\begin{proposition}[Unique interior well in the principal zero sector]
\label{prop:v6v49-zero-sector-well}
Inside the principal lift region with
\(\sum_ix_i=0\), the only interior local minimum of \(S\) is
\(x=0\).
\end{proposition}

\begin{proof}
At a local minimum the constrained Hessian is nonnegative.  Lemma
\ref{lem:v6v49-negative-direction} excludes every critical branch with a
negative cosine.  Hence all angles lie on the principal sine branch
\((-\pi/2,\pi/2)\), where equality of their sines implies equality of the
angles.  Their sum is zero, so every angle is zero.
\end{proof}

This proposition removes secondary interior wells but does not by itself
give a sector gap; a Hardy/capacity estimate is still needed near the chart
boundary.

\subsection{Corrected BRC hierarchy}

\begin{definition}[Sector-resolved Bianchi restriction card]
\label{def:v6v49-SBRC}
For a specified admissible compact flux sector \(m\), the SBRC is
\begin{equation}
 \operatorname{gap}(\nu_{K,u}^{\rm curv,m})
 \ge c\beta
 \label{eq:v6v49-SBRC}
\end{equation}
with boundary flux and chart transitions fixed consistently.  If sectors
are allowed to mix, their discrete transition process is retained as a
separate topological/monopole channel rather than absorbed into SBRC.
\end{definition}

\begin{proposition}[Logical alternatives for hypercharge]
\label{prop:v6v49-hypercharge-alternatives}
The cube no-go leaves three mathematically distinct routes:
\begin{enumerate}[label=\textup{(\roman*)},leftmargin=3.7em]
\item use a noncompact hypercharge potential, for which the integer lift
      wells are absent and the real cochain disintegration of V48 applies;
\item restrict to and prove stability of an admissible zero-flux compact
      sector, then establish SBRC there; or
\item retain compact monopole/flux sectors and prove their transition and
      activity estimates as explicit nonperturbative degrees of freedom.
\end{enumerate}
An all-sector order-\(\beta\) Poincare gap is not available.
\end{proposition}

\begin{proof}
Theorem~\ref{thm:v6v49-metastable-gap} rules out the all-sector BRC.  The
three alternatives respectively remove, freeze, or retain the sector label,
which exhausts the treatment of that obstruction at the finite cube level.
\end{proof}

\begin{firewall}[Sector restriction needs dynamical justification]
\label{fw:v6v49-sector-justification}
Selecting \(m=0\) is not a harmless gauge choice.  In compact \(U(1)\),
sector-changing monopole events are configurations of the theory.  Their
exclusion requires a declared noncompact formulation, an admissibility
condition whose stability and continuum equivalence are proved, or an
explicit suppression theorem.
\end{firewall}

\begin{firewall}[The cube model is not the full hypercharge lattice]
\label{fw:v6v49-cube-scope}
The one-cube action omits shared faces, charged matter determinants, Higgs
couplings, and Einstein variables.  These can alter sector activities and
transition barriers.  They cannot restore the false all-sector gap without
an identified mechanism that removes or rapidly connects the compact flux
wells.
\end{firewall}

\begin{theorem}[Sixth-series V49 result]
\label{thm:v6v49-status}
The smallest compact \(U(1)\) Bianchi manifold has a global zero-flux
minimum and strict metastable uniform-flux minima at plaquette angle
\(\pm\pi/3\), each with local Hessian \(\beta/2\) and action three.  A
cutoff localized in either metastable well proves that the full all-sector
gap is exponentially small, so nonlinear BRC is false as originally
quantified.  The principal zero-lift sector has no other interior local
minimum.  Hypercharge must therefore use a noncompact formulation, a
justified sector-resolved BRC, or explicit compact monopole channels.
\end{theorem}

\begin{proof}
Use Lemma~\ref{lem:v6v49-critical-equation},
Theorem~\ref{thm:v6v49-three-minima},
Theorem~\ref{thm:v6v49-metastable-gap},
Proposition~\ref{prop:v6v49-zero-sector-well}, and
Proposition~\ref{prop:v6v49-hypercharge-alternatives}, subject to both
firewalls.
\end{proof}

V50 will perform the compulsory YM/NS audit.  It will compare their newest
bridge-score and formation results with the compact-flux no-go, then choose
between a noncompact hypercharge trajectory and a sector-resolved compact
construction before returning to the mixed SM raw comparison.


\clearpage
\section{V50: companion audit and the noncompact hypercharge curvature endpoint}
\label{sec:v6v50}

V49 rules out an order--\(\beta\) Poincare gap for the compact
all-flux cube law.  The obstruction is topological rather than a defect of
the local Hessian: the nonzero lift sectors contain strict metastable wells.
This version performs the compulsory companion audit and then selects a
precise noncompact hypercharge branch.  The selection is provisional at the
level of the complete Standard Model, but it permits an unconditional
finite-cutoff theorem and a rigorous free continuum endpoint.

\subsection{Compulsory V50 audit}

The newest companion sources inspected for this audit were
\begin{center}
\begin{tabular}{p{0.27\linewidth}p{0.18\linewidth}p{0.43\linewidth}}
programme & source & audited conclusion\\ \hline
Yang--Mills & \texttt{V36} & exact binary heat-normalized activity and
binary WPRAC are proved; ternary and all-arity WPRAC remain open because a
separated high-spin estimate grows like \(e^{B_A^2}\)\\
Navier--Stokes & \texttt{V31} & the exact heat-formation ladder is proved,
but recovering the critical first moment from a weaker source norm restores
the missing scale factor and returns the continuation-class stock
\end{tabular}
\end{center}
At the time of reading, the Yang--Mills source had \(17\,182\) logical
lines, \(544\,062\) bytes, and SHA--256
\[
 \texttt{eebfa58f356e18179b6355f2e477662975272c85658560f11f6b3a22ced89549},
\]
while the Navier--Stokes source had \(23\,052\) logical lines,
\(887\,537\) bytes, and SHA--256
\[
 \texttt{f1121b62238ad5491bb7cf03635ea9934942edf573ed8faaa9ee79e1d38a6696}.
\]

\begin{assessment}[Direction forced by the audit]
\label{ass:v6v50-audit-direction}
The Yang--Mills binary theorem may be used as a two-leg local input, but not
as an all-order non-Abelian construction.  The Navier--Stokes sharpness test
shows why a disappearing inverse derivative must not be hidden in a generic
``heat gain.''  Accordingly, the hypercharge branch below keeps the least
positive singular value of the lattice curl in every potential-coordinate
estimate.  A volume-independent estimate is asserted only for observables
written in curvature coordinates, where the inverse derivative cancels
algebraically.
\end{assessment}

\subsection{Finite-complex noncompact Maxwell quotient}

Let
\[
 C^0(K;\mathbb R)\xrightarrow{d_0}C^1(K;\mathbb R)
 \xrightarrow{d_1}C^2(K;\mathbb R)
\]
be a finite real cochain complex with fixed Euclidean cell inner products.
Gauge directions, flat directions, and harmonic torons lie in
\(\ker d_1\).  Fix all of them and use the orthogonal representative space
\begin{equation}
 \mathcal H_K:=(\ker d_1)^\perp=\operatorname{im}d_1^*.
 \label{eq:v6v50-HK}
\end{equation}
Set
\begin{equation}
 L_K:=d_1^*d_1\big|_{\mathcal H_K},
 \qquad
 \sigma_K^2:=\lambda_{\min}(L_K)>0.
 \label{eq:v6v50-LK}
\end{equation}
The noncompact Maxwell law on the fixed quotient is
\begin{equation}
 d\gamma_{\beta,K}(A)
 :=Z_{\beta,K}^{-1}
 \exp\!\left[-{\beta\over2}\|d_1A\|^2\right]dA,
 \qquad A\in\mathcal H_K.
 \label{eq:v6v50-Maxwell-law}
\end{equation}

\begin{theorem}[Exact quotient and curvature gaps]
\label{thm:v6v50-two-gaps}
The law \(\gamma_{\beta,K}\) is the centered Gaussian with covariance
\((\beta L_K)^{-1}\).  It has Euclidean potential-coordinate Poincare gap
\begin{equation}
 \lambda_A=\beta\sigma_K^2.
 \label{eq:v6v50-potential-gap}
\end{equation}
The map \(d_1:\mathcal H_K\to\operatorname{im}d_1\) is a linear
isomorphism, and its pushforward has density
\begin{equation}
 d\widehat\gamma_{\beta,K}(F)
 =\widehat Z_{\beta,K}^{-1}e^{-\beta\|F\|^2/2}\,dF,
 \qquad F\in\operatorname{im}d_1.
 \label{eq:v6v50-curvature-law}
\end{equation}
Consequently, for every smooth \(\varphi\) on
\(\operatorname{im}d_1\),
\begin{equation}
 \operatorname{Var}_{\gamma_{\beta,K}}\!\bigl(\varphi(d_1A)\bigr)
 \le {1\over\beta}
 \int\|\nabla_F\varphi(d_1A)\|^2\,d\gamma_{\beta,K}(A).
 \label{eq:v6v50-curvature-gap}
\end{equation}
The constant in \eqref{eq:v6v50-curvature-gap} is independent of
\(\sigma_K\).
\end{theorem}

\begin{proof}
On \(\mathcal H_K\), the quadratic form in
\eqref{eq:v6v50-Maxwell-law} is \(\langle A,L_KA\rangle\), so Gaussian
diagonalization gives the covariance and
\eqref{eq:v6v50-potential-gap}.  The kernel of the restricted \(d_1\) is
zero, and its range is \(\operatorname{im}d_1\), proving bijectivity.
Under \(F=d_1A\), Lebesgue measure changes by the constant
\(\det(L_K)^{-1/2}\).  This constant is absorbed into
\(\widehat Z_{\beta,K}\), proving \eqref{eq:v6v50-curvature-law}.
The standard Gaussian Poincare inequality on the Euclidean space
\(\operatorname{im}d_1\) gives \eqref{eq:v6v50-curvature-gap}.
\end{proof}

\begin{remark}[The photon scale is retained rather than erased]
\label{rem:v6v50-photon-scale}
On a box of diameter \(L\), \(\sigma_K\) has the inverse-length scale of
the least nonzero momentum.  Thus \eqref{eq:v6v50-potential-gap} degenerates
in the thermodynamic limit, as a massless photon requires.  The curvature
estimate \eqref{eq:v6v50-curvature-gap} is uniform because applying \(d_1\)
removes that inverse derivative.  These two statements concern different
Dirichlet metrics and are not contradictory.
\end{remark}

\subsection{A relative-curvature perturbation theorem}

Let \(W:\mathcal H_K\to\mathbb R\) be \(C^2\), assume the law below is
normalizable, and suppose that for some \(0\le\kappa<\beta\),
\begin{equation}
 \nabla^2W(A)\succeq-\kappa L_K
 \quad\hbox{for every }A\in\mathcal H_K.
 \label{eq:v6v50-relative-curvature}
\end{equation}
Define
\begin{equation}
 d\mu_{\beta,W}(A)
 :=Z_{\beta,W}^{-1}
 \exp\!\left[-{\beta\over2}\|d_1A\|^2-W(A)\right]dA.
 \label{eq:v6v50-perturbed-law}
\end{equation}

\begin{theorem}[Relative Maxwell curvature card]
\label{thm:v6v50-RMC}
Under \eqref{eq:v6v50-relative-curvature}, every smooth
\(f:\mathcal H_K\to\mathbb R\) satisfies
\begin{equation}
 \operatorname{Var}_{\mu_{\beta,W}}(f)
 \le {1\over\beta-\kappa}
 \int\!\langle\nabla f,L_K^{-1}\nabla f\rangle,d\mu_{\beta,W}.
 \label{eq:v6v50-energy-Poincare}
\end{equation}
In particular,
\begin{equation}
 \operatorname{Var}_{\mu_{\beta,W}}(f)
 \le {1\over(\beta-\kappa)\sigma_K^2}
 \int\|\nabla f\|^2,d\mu_{\beta,W}.
 \label{eq:v6v50-Euclidean-Poincare}
\end{equation}
If \(f(A)=\varphi(d_1A)\), then
\begin{equation}
 \operatorname{Var}_{\mu_{\beta,W}}\!\bigl(\varphi(d_1A)\bigr)
 \le {1\over\beta-\kappa}
 \int\|P_{\operatorname{im}d_1}\nabla_F\varphi(d_1A)\|^2
 \,d\mu_{\beta,W}(A).
 \label{eq:v6v50-perturbed-curvature-gap}
\end{equation}
Thus the curvature-observable gap is again independent of
\(\sigma_K\).
\end{theorem}

\begin{proof}
The full action \(S(A)=\beta\langle A,L_KA\rangle/2+W(A)\) obeys
\[
 \nabla^2S(A)\succeq(\beta-\kappa)L_K>0.
\]
The finite-dimensional Brascamp--Lieb inequality therefore gives
\[
 \operatorname{Var}_{\mu_{\beta,W}}(f)
 \le\int\!\langle\nabla f,(\nabla^2S)^{-1}\nabla f\rangle,d\mu_{\beta,W}
 \le {1\over\beta-\kappa}
 \int\!\langle\nabla f,L_K^{-1}\nabla f\rangle,d\mu_{\beta,W},
\]
which is \eqref{eq:v6v50-energy-Poincare}.  Since
\(L_K\succeq\sigma_K^2I\), this implies
\eqref{eq:v6v50-Euclidean-Poincare}.

For \(f=\varphi\circ d_1\), put
\(\eta=\nabla_F\varphi(d_1A)\).  Then \(\nabla f=d_1^*\eta\).  A singular
value decomposition of \(d_1:\mathcal H_K\to\operatorname{im}d_1\)
gives the exact identity
\begin{equation}
 \langle d_1^*\eta,L_K^{-1}d_1^*\eta\rangle
 =\|P_{\operatorname{im}d_1}\eta\|^2.
 \label{eq:v6v50-inverse-curl-cancel}
\end{equation}
Insert this identity into \eqref{eq:v6v50-energy-Poincare}.
\end{proof}

\begin{corollary}[A curvature-local sufficient condition]
\label{cor:v6v50-curvature-local}
If \(W(A)=w(d_1A)\) and
\(\nabla_F^2w(F)\succeq-\kappa I\) on
\(\operatorname{im}d_1\), then \eqref{eq:v6v50-relative-curvature} holds
and all conclusions of Theorem~\ref{thm:v6v50-RMC} follow.
\end{corollary}

\begin{proof}
For \(h\in\mathcal H_K\),
\[
 \langle h,\nabla^2W(A)h\rangle
 =\langle d_1h,\nabla_F^2w(d_1A)d_1h\rangle
 \ge-\kappa\|d_1h\|^2
 =-\kappa\langle h,L_Kh\rangle.
\]
\end{proof}

\begin{firewall}[What must be proved for the interacting SM]
\label{fw:v6v50-interacting-curvature}
The relative bound \eqref{eq:v6v50-relative-curvature} is not automatic for
the Higgs integral or the chiral fermion determinant.  A complex determinant
does not define the real function \(W\) used above, and massless Dirac
singular values can make determinant derivatives large.  Consequently
Theorem~\ref{thm:v6v50-RMC} proves the noncompact Maxwell card and a precise
sufficient perturbative extension; it does not prove that the full
SM effective remainder satisfies that extension.
\end{firewall}

\subsection{Free continuum cylinder endpoint}

Fix the flat four-torus \(\mathbb T_L^4\).  For a reciprocal lattice vector
\(k\in(2\pi/L)\mathbb Z^4\setminus\{0\}\), let
\(P_k=I-k\otimes k/|k|^2\).  Let \(A_\Lambda\) be the real centered
Gaussian transverse one-form with Fourier covariance
\begin{equation}
 \mathbb E\bigl[\widehat A_\Lambda(k)
 \widehat A_\Lambda(k)^*\bigr]
 ={g_Y^2\over|k|^2}P_k\,mathbf 1_{\{|k|\le\Lambda\}}.
 \label{eq:v6v50-Fourier-covariance}
\end{equation}
Reality means \(\widehat A_\Lambda(-k)=overline{\widehat A_\Lambda(k)}\).
The zero mode is fixed to zero.

\begin{theorem}[Noncompact free hypercharge continuum field]
\label{thm:v6v50-free-continuum}
For every \(s>1\), the fields \(A_\Lambda\) converge in
\(L^2(\Omega;H^{-s}(\mathbb T_L^4;\mathbb R^4))\) to a centered Gaussian
distribution \(A\).  For every smooth real co-closed mean-zero one-form
\(J\),
\begin{equation}
 \mathbb E e^{i\langle A,J\rangle}
 =\exp\!\left[-{g_Y^2\over2}
 \sum_{k\ne0}{|\widehat J(k)|^2\over|k|^2}\right].
 \label{eq:v6v50-characteristic-functional}
\end{equation}
Moreover \(F=dA\) is a well-defined random two-form in
\(H^{-t}\) for every \(t>2\), and its finite-cutoff laws satisfy the exact
Bianchi identity \(dF=0\).
\end{theorem}

\begin{proof}
By \eqref{eq:v6v50-Fourier-covariance} and
\(\operatorname{tr}P_k=3\),
\[
 \mathbb E\|A_{\Lambda'}-A_\Lambda\|_{H^{-s}}^2
 \le Cg_Y^2
 \sum_{\Lambda<|k|\le\Lambda'}
 {1\over(1+|k|^2)^s|k|^2}.
\]
In four dimensions the number of lattice points in a shell of radius
\(r\) is \(O(r^3)\), so the tail is bounded by a constant times
\(\int_\Lambda^\infty r^{1-2s}\,dr\), which tends to zero exactly when
\(s>1\).  Completeness of
\(L^2(\Omega;H^{-s})\) gives the limit.  Pairing with smooth \(J\) is
continuous, and the finite-dimensional Gaussian characteristic function
passes to the limit.  Since \(k\cdot\widehat J(k)=0\),
\(P_k\widehat J(k)=\widehat J(k)\), giving
\eqref{eq:v6v50-characteristic-functional}.

The Fourier multiplier for \(d\) grows like \(|k|\).  Hence
\[
 \mathbb E\|dA\|_{H^{-t}}^2
 \le Cg_Y^2\sum_{k\ne0}(1+|k|^2)^{-t},
\]
which converges in four dimensions for \(t>2\).  Finally \(d^2=0\) at
every cutoff, and continuity of \(d:H^{-t}\to H^{-t-1}\) passes the identity
to the limit.
\end{proof}

\begin{corollary}[The endpoint remains massless]
\label{cor:v6v50-massless}
The transverse covariance multiplier in
\eqref{eq:v6v50-Fourier-covariance} is \(g_Y^2/|k|^2\).  It has no positive
mass term and its least finite-volume potential gap is proportional to
\((2\pi/L)^2\).  Thus the construction supplies the free massless
hypercharge continuum cylinder measure; it does not manufacture a photon
mass from the curvature Poincare estimate.
\end{corollary}

\begin{proof}
Both assertions follow directly by diagonalizing the covariance on the
transverse Fourier modes.
\end{proof}

\subsection{Updated route and claim boundary}

\begin{definition}[Noncompact hypercharge curvature card]
\label{def:v6v50-NHCC}
NHCC consists of the following scale-indexed statements after gauge, toron,
and boundary-flux fixing:
\begin{enumerate}[label=\textup{(NH\arabic*)},leftmargin=4.8em]
\item the effective real action is normalizable on \(\mathcal H_K\);
\item its real Hessian is bounded below by
      \((\beta_Y-\kappa_Y)L_K\), with
      \(0\le\kappa_Y<\beta_Y\);
\item its phase or sign is retained as an observable and satisfies a
      separate rooted determinant-line card; and
\item the constants survive the chosen ultraviolet matching trajectory.
\end{enumerate}
\end{definition}

\begin{proposition}[Exact status of NHCC]
\label{prop:v6v50-NHCC-status}
For pure noncompact hypercharge, NHCC\textup{(NH1)--(NH2)} hold with
\(\kappa_Y=0\), and the free continuum cylinder endpoint is
Theorem~\ref{thm:v6v50-free-continuum}.  For the coupled Standard Model,
NHCC\textup{(NH2)--(NH4)} remain open.
\end{proposition}

\begin{proof}
The pure statements are Theorems~\ref{thm:v6v50-two-gaps} and
\ref{thm:v6v50-RMC} with \(W=0\).  There is then no determinant phase.
The interacting assertions require the Higgs, chiral determinant, and
running-coupling estimates excluded by
Firewall~\ref{fw:v6v50-interacting-curvature} and were not proved here.
\end{proof}

\begin{firewall}[Continuum endpoint versus the full continuum]
\label{fw:v6v50-full-continuum}
Theorem~\ref{thm:v6v50-free-continuum} is a complete construction of the
free noncompact hypercharge field on a fixed torus.  It is not the coupled
SM--Einstein continuum: it does not remove a non-Abelian cutoff, construct
the Higgs--Yukawa measure, trivialize the anomaly-free determinant line
globally, control the hypercharge running coupling, construct the positive
Einstein measure, or prove a common Osterwalder--Schrader limit.
\end{firewall}

\begin{theorem}[Sixth-series V50 result]
\label{thm:v6v50-status}
After the compulsory companion audit, the noncompact hypercharge branch has
the following rigorous status.  On every finite cochain quotient its free
law is an exact Gaussian; the potential gap is
\(\beta\sigma_K^2\), while curvature observables have gap \(\beta\).
A real perturbation with negative Hessian at most \(\kappa L_K\) preserves
these estimates with \(\beta\) replaced by \(\beta-\kappa\).  On the
four-torus the free transverse cutoff fields converge in
\(H^{-s}\), \(s>1\), to the massless Gaussian hypercharge distribution.
The coupled matter relative-curvature estimate and determinant-phase
transport remain open.
\end{theorem}

\begin{proof}
Use Theorems~\ref{thm:v6v50-two-gaps},
\ref{thm:v6v50-RMC}, and \ref{thm:v6v50-free-continuum}, together with
Corollaries~\ref{cor:v6v50-curvature-local} and
\ref{cor:v6v50-massless}, subject to
Firewalls~\ref{fw:v6v50-interacting-curvature} and
\ref{fw:v6v50-full-continuum}.
\end{proof}

V51 will attack NHCC at its first genuinely interacting point.  It will
differentiate a finite-dimensional chiral determinant on the complement of
its zero modes, obtain the sharp inverse-singular-value dependence of its
gauge Hessian, and decide which part can be paid by anomaly cancellation and
which part requires a shell-resolved determinant card.


\clearpage
\section{V51: reduced determinant Hessians, inverse-gap loss, and the anomaly mismatch}
\label{sec:v6v51}

V50 proves the noncompact Maxwell curvature card for a real perturbation
whose negative Hessian is smaller than the Maxwell form.  The first coupled
question is therefore quantitative: what lower bound is supplied by a
finite-dimensional chiral determinant after its zero modes have been split
off?  This version computes the exact answer.  The calculation also
separates perturbative anomaly cancellation from coercivity of the real
determinant modulus.

\subsection{Jacobi calculus on an invertible determinant chart}

Let \(E_+\) and \(E_-\) be complex Hilbert spaces of the same finite
dimension \(r\).  On an open set \(\mathcal U\) in a real parameter space,
let
\[
 D(a):E_+\longrightarrow E_-
\]
be \(C^2\) and invertible.  Write \(R(a)=D(a)^{-1}\).  A local branch of
the logarithm defines
\begin{equation}
 \Psi(a):=-\log\det D(a),
 \qquad
 W(a):=\operatorname{Re}\Psi(a)=-\log|\det D(a)|,
 \qquad
 \Theta(a):=\operatorname{Im}\Psi(a).
 \label{eq:v6v51-logdet-fields}
\end{equation}
Only derivatives of \(\Psi\) are branch independent.

\begin{theorem}[Exact determinant score and Hessian]
\label{thm:v6v51-Jacobi}
For real tangent directions \(X,Y\),
\begin{align}
 d\Psi[X]
 &=-\operatorname{Tr}(R D_X),
 \label{eq:v6v51-first-Jacobi}\\
 d^2\Psi[X,Y]
 &=\operatorname{Tr}(R D_Y R D_X)
   -\operatorname{Tr}(R D_{XY}).
 \label{eq:v6v51-second-Jacobi}
\end{align}
If \(s_{\min}(D)\ge\rho>0\), then
\begin{align}
 |dW[X]|+|d\Theta[X]|
 &\le 2r\rho^{-1}\|D_X\|_{\rm op},
 \label{eq:v6v51-score-bound}\\
 |d^2W[X,Y]|+|d^2\Theta[X,Y]|
 &\le 2r\left(
 \rho^{-2}\|D_X\|_{\rm op}\|D_Y\|_{\rm op}
 +\rho^{-1}\|D_{XY}\|_{\rm op}\right).
 \label{eq:v6v51-Hessian-bound}
\end{align}
In particular,
\begin{equation}
 d^2W[X,X]
 \ge-r\left(
 \rho^{-2}\|D_X\|_{\rm op}^2
 +\rho^{-1}\|D_{XX}\|_{\rm op}\right).
 \label{eq:v6v51-lower-Hessian}
\end{equation}
\end{theorem}

\begin{proof}
Jacobi's identity gives
\(d\log\det D[X]=\operatorname{Tr}(RD_X)\).  Differentiating
\(RD=I\) gives \(dR[Y]=-RD_YR\), and a second differentiation proves
\eqref{eq:v6v51-second-Jacobi}.  Since
\(\|R\|_{\rm op}\le\rho^{-1}\) and
\(|\operatorname{Tr}T|\le r\|T\|_{\rm op}\), the separate bounds
\[
 |d\Psi[X]|\le r\rho^{-1}\|D_X\|_{\rm op}
\]
and
\[
 |d^2\Psi[X,Y]|
 \le r\rho^{-2}\|D_X\|_{\rm op}\|D_Y\|_{\rm op}
     +r\rho^{-1}\|D_{XY}\|_{\rm op}
\]
follow.  Real and imaginary parts are each bounded by the modulus, which
proves the displayed estimates; retaining only the real part gives
\eqref{eq:v6v51-lower-Hessian}.
\end{proof}

\begin{remark}[A sharper harmless factor]
\label{rem:v6v51-factor-two}
The factor two in \eqref{eq:v6v51-score-bound} and
\eqref{eq:v6v51-Hessian-bound} only packages the real and phase estimates in
one line.  Each of \(|dW|\), \(|d\Theta|\), \(|d^2W|\), and
\(|d^2\Theta|\) separately obeys the corresponding right side without that
factor.
\end{remark}

\subsection{Compatibility with the V39 zero-mode split}

Suppose now that \(D(a)\) has constant kernel and cokernel dimensions on a
chart.  Let \(P_+(a)\) and \(P_-(a)\) be smooth orthogonal projections onto
chosen complements of the kernel and cokernel, and let
\begin{equation}
 D_{\rm red}(a):P_+(a)E_+\longrightarrow P_-(a)E_-
 \label{eq:v6v51-reduced-D}
\end{equation}
be the invertible reduced operator used in the V39 Schur--Berezin gluing
formula.

\begin{proposition}[Reduced determinant chart formula]
\label{prop:v6v51-reduced-chart}
Choose unitary parallel transports which identify the ranges of
\(P_+(a)\) and \(P_-(a)\) with fixed spaces.  In those frames the reduced
operator is an ordinary invertible matrix \(\widetilde D_{\rm red}(a)\), and
Theorem~\ref{thm:v6v51-Jacobi} applies to
\(-\log\det\widetilde D_{\rm red}\).  Changing either parallel frame
multiplies the reduced determinant by the determinant of a unitary
transition map.  Hence the real field
\(-\log|\det D_{\rm red}|\) and all of its derivatives are frame
independent, while the phase derivatives transform as the connection of the
determinant line.
\end{proposition}

\begin{proof}
Constant rank makes the two complement bundles smooth finite-dimensional
Hermitian vector bundles.  Parallel transport supplies local unitary
trivializations.  The usual determinant calculus then proves the first
statement.  Under a change of frames \(U_+\) and \(U_-\), the matrix changes
to \(U_-\widetilde D_{\rm red}U_+^{-1}\), so its determinant is multiplied
by \(\det U_-\det U_+^{-1}\), a unit-modulus factor.  The modulus is
unchanged.  The logarithmic phase gains the transition phase, which is
exactly the local transformation law for a determinant-line connection.
\end{proof}

\begin{firewall}[Rank-changing strata remain outside one chart]
\label{fw:v6v51-rank-change}
At a rank-changing configuration, no positive \(\rho\) exists and the
complement bundles in Proposition~\ref{prop:v6v51-reduced-chart} need not
extend smoothly.  Zero-mode boundary legs from V39 must then be retained.
Deleting the vanishing singular values from each configuration separately
does not produce a smooth effective action across the stratum.
\end{firewall}

\subsection{Sharpness of the inverse singular-value powers}

\begin{proposition}[One-mode saturation]
\label{prop:v6v51-one-mode}
For \(D_\rho(x)=\rho+ix\), \(\rho>0\),
\begin{align}
 W_\rho(x)&=-\frac12\log(\rho^2+x^2),
 & W_\rho''(0)&=-\rho^{-2},
 \label{eq:v6v51-scalar-modulus}\\
 \Theta_\rho'(0)&=-\rho^{-1}.
 \label{eq:v6v51-scalar-phase}
\end{align}
Thus the powers \(\rho^{-2}\) in the real Hessian and \(\rho^{-1}\) in the
phase score cannot be improved for general determinant families.
\end{proposition}

\begin{proof}
The modulus and argument of \(\rho+ix\) are
\((\rho^2+x^2)^{1/2}\) and \(\arctan(x/\rho)\).  Insert them into
\eqref{eq:v6v51-logdet-fields} and differentiate at zero.
\end{proof}

\begin{corollary}[A spectral gap is logically necessary for a raw Hessian
card]
\label{cor:v6v51-gap-necessary}
No bound on the negative Hessian of \(-\log|\det D|\) which is uniform over
all invertible finite matrices can depend only on uniform upper bounds for
\(D_X\) and \(D_{XX}\).  Some lower-singular-value information or a
cancellation prior to absolute values is necessary.
\end{corollary}

\begin{proof}
In Proposition~\ref{prop:v6v51-one-mode}, \(D_X=i\) and \(D_{XX}=0\) are
uniform while \(-W_\rho''(0)=\rho^{-2}\to\infty\).
\end{proof}

\subsection{Anomaly cancellation does not cancel the real square trace}

Use the all-left-handed one-generation hypercharges
\begin{equation}
 \begin{array}{c|ccccc}
 \text{multiplet}&Q_L&u^c&d^c&L&e^c\\ \hline
 \text{multiplicity}&6&3&3&2&1\\
 Y&1/6&-2/3&1/3&-1/2&1
 \end{array}
 \label{eq:v6v51-hypercharge-table}
\end{equation}
already used in the V38 anomaly audit.

\begin{lemma}[Linear, cubic, and square sums]
\label{lem:v6v51-charge-sums}
For \eqref{eq:v6v51-hypercharge-table},
\begin{equation}
 \sum_f n_fY_f=0,
 \qquad
 \sum_f n_fY_f^3=0,
 \qquad
 \sum_f n_fY_f^2={10\over3}>0.
 \label{eq:v6v51-charge-sums}
\end{equation}
\end{lemma}

\begin{proof}
Directly,
\begin{align*}
 \sum_fn_fY_f
 &=1-2+1-1+1=0,\\
 \sum_fn_fY_f^3
 &={1\over36}-{8\over9}+{1\over9}-{1\over4}+1=0,\\
 \sum_fn_fY_f^2
 &={1\over6}+{4\over3}+{1\over3}+{1\over2}+1
 ={10\over3}.
\end{align*}
\end{proof}

\begin{proposition}[Anomaly-free toy family with uncancelled negative
curvature]
\label{prop:v6v51-anomaly-mismatch}
Give every Weyl component the scalar local model
\(D_f(x)=\rho+iY_fx\), with a common \(\rho>0\), and multiply all
determinants with their multiplicities.  The total phase has vanishing
linear charge trace and the cubic charge trace vanishes, but the total real
effective action satisfies
\begin{equation}
 W_{\rm tot}''(0)
 =-{1\over\rho^2}\sum_fn_fY_f^2
 =-{10\over3\rho^2}.
 \label{eq:v6v51-uncancelled-square}
\end{equation}
\end{proposition}

\begin{proof}
The first two cancellations are Lemma~\ref{lem:v6v51-charge-sums}.  Apply
Proposition~\ref{prop:v6v51-one-mode} with \(x\) replaced by \(Y_fx\),
sum with multiplicities, and use the square sum in the same lemma.
\end{proof}

\begin{assessment}[Scope of anomaly cancellation]
\label{ass:v6v51-anomaly-scope}
The V38 arithmetic is essential for local gauge consistency of the phase
and for determinant-line gluing.  Proposition~\ref{prop:v6v51-anomaly-mismatch}
shows that it is not a lower-Hessian estimate for the determinant modulus.
The surviving square trace is the same charge invariant which appears in
Abelian vacuum polarization.  Therefore the hypercharge running obstruction
and the NHCC real-curvature obstruction are related rather than mutually
cancelling.
\end{assessment}

\subsection{Exact singular-shell resolution}

At a fixed configuration choose a singular value decomposition
\(D=U\Sigma V^*\), with singular values \(\sigma_i>0\), and define
\begin{equation}
 B_X:=U^*D_XV,
 \qquad B_{XY}:=U^*D_{XY}V.
 \label{eq:v6v51-BX}
\end{equation}
Let \(\{I_a\}_a\) be any partition of the singular-value labels and let
\(P_a\) be the corresponding coordinate projections.  Put
\begin{equation}
 m_a:=\max_{i\in I_a}\sigma_i^{-1}.
 \label{eq:v6v51-ma}
\end{equation}

\begin{theorem}[Exact shell formula and absolute shell ledger]
\label{thm:v6v51-shell-ledger}
The affine part of the determinant Hessian has the exact expansion
\begin{equation}
 \operatorname{Tr}(RD_YRD_X)
 =\sum_{a,b}\operatorname{Tr}\!\left(
 P_a\Sigma^{-1}B_YP_b\Sigma^{-1}B_XP_a\right).
 \label{eq:v6v51-exact-shell-Hessian}
\end{equation}
Consequently,
\begin{align}
 |d^2\Psi[X,Y]|
 &\le\sum_{a,b}m_am_b
 \|P_aB_YP_b\|_{\rm HS}
 \|P_bB_XP_a\|_{\rm HS}
 \notag\\
 &\quad+sum_a m_a\|P_aB_{XY}P_a\|_{\rm tr}
       +\sum_{a\ne b}m_a\|P_aB_{XY}P_b\|_{\rm tr}.
 \label{eq:v6v51-shell-bound}
\end{align}
For the trace term the off-diagonal summands in the last sum are actually
zero; they are displayed only to make clear that a termwise trace-norm
majorant may discard this exact diagonal cancellation.
\end{theorem}

\begin{proof}
Cyclicity of trace and the SVD give
\[
 \operatorname{Tr}(RD_YRD_X)
 =\operatorname{Tr}(\Sigma^{-1}B_Y\Sigma^{-1}B_X).
\]
Insert \(I=\sum_aP_a\) on both sides of \(B_Y\) and use cyclicity once
more to obtain \eqref{eq:v6v51-exact-shell-Hessian}.  On the
\((a,b)\)-block, the two inverse singular-value factors have operator norms
at most \(m_a\) and \(m_b\).  The inequality
\(|\operatorname{Tr}(CD)|\le\|C\|_{\rm HS}\|D\|_{\rm HS}\) gives the
double sum.  The second Jacobi term is
\(\operatorname{Tr}(\Sigma^{-1}B_{XY})\); insertion of the shell
projections shows that only its diagonal blocks contribute.  Bounding them
by trace norm proves \eqref{eq:v6v51-shell-bound} and the final assertion.
\end{proof}

\begin{definition}[Determinant singular-shell card]
\label{def:v6v51-DSSC}
DSSC is a uniform, scale-indexed estimate of the complete shell expression
\eqref{eq:v6v51-exact-shell-Hessian}, combined across all SM multiplets
before absolute values, which proves
\begin{equation}
 d^2W_{\rm ferm}[X,X]\ge-\kappa_{\rm f}\|d_1X\|^2,
 \qquad \kappa_{\rm f}<\beta_Y-kappa_{\rm Higgs},
 \label{eq:v6v51-DSSC}
\end{equation}
on each retained determinant chart.  The phase counterpart keeps the
determinant-line connection and its rooted cumulants rather than inserting
them into the positive reference law.
\end{definition}

\begin{proposition}[A finite-volume sufficient estimate]
\label{prop:v6v51-finite-sufficient}
Suppose on a fixed determinant chart that
\begin{equation}
 \|D_X\|_{\rm op}\le a_1\|d_1X\|,
 \qquad
 \|D_{XX}\|_{\rm op}\le a_2\|d_1X\|^2,
 \qquad s_{\min}(D)\ge\rho.
 \label{eq:v6v51-derivative-assumptions}
\end{equation}
Then
\begin{equation}
 d^2W[X,X]\ge
 -r\left(a_1^2\rho^{-2}+a_2\rho^{-1}\right)\|d_1X\|^2.
 \label{eq:v6v51-finite-kappa}
\end{equation}
Thus Theorem~\ref{thm:v6v50-RMC} applies whenever the sum of this constant
and the Higgs negative-curvature constant is strictly smaller than
\(\beta_Y\).
\end{proposition}

\begin{proof}
Insert \eqref{eq:v6v51-derivative-assumptions} into
\eqref{eq:v6v51-lower-Hessian}; then add the Higgs lower bound and apply
Theorem~\ref{thm:v6v50-RMC}.
\end{proof}

\begin{firewall}[Why the finite sufficient estimate is not yet uniform]
\label{fw:v6v51-finite-not-uniform}
For ordinary linkwise bounds one typically controls \(D_X\) by
\(\|X\|\), not by \(\|d_1X\|\).  On the fixed quotient this conversion
costs \(\sigma_K^{-1}\), precisely the massless inverse derivative retained
in V50.  In addition, massless fermion charts need not have a uniform
\(\rho\).  Therefore Proposition~\ref{prop:v6v51-finite-sufficient} is a
correct finite-chart criterion, but it is not DSSC and must not be presented
as a volume- or cutoff-uniform SM estimate.
\end{firewall}

\begin{theorem}[Sixth-series V51 result]
\label{thm:v6v51-status}
On every constant-rank determinant chart, the reduced chiral determinant
has the exact score and Hessian formulas
\eqref{eq:v6v51-first-Jacobi}--\eqref{eq:v6v51-second-Jacobi}.  Their raw
uniform bounds lose one and two inverse powers of the least nonzero singular
value, and the scalar family \(\rho+ix\) proves those losses sharp.  Standard
Model linear and cubic hypercharge anomaly sums vanish, but the square sum
equals \(10/3\); consequently anomaly cancellation does not supply the
negative-Hessian bound needed by NHCC.  The exact singular-shell formula
\eqref{eq:v6v51-exact-shell-Hessian} identifies DSSC as the remaining
nonperturbative determinant input.
\end{theorem}

\begin{proof}
Use Theorem~\ref{thm:v6v51-Jacobi},
Propositions~\ref{prop:v6v51-reduced-chart},
\ref{prop:v6v51-one-mode}, and
\ref{prop:v6v51-anomaly-mismatch}, Lemma~\ref{lem:v6v51-charge-sums}, and
Theorem~\ref{thm:v6v51-shell-ledger}, subject to
Firewalls~\ref{fw:v6v51-rank-change} and
\ref{fw:v6v51-finite-not-uniform}.
\end{proof}

V52 will go upstream from DSSC rather than estimate
\eqref{eq:v6v51-shell-bound} term by term.  It will use gauge covariance to
derive the exact determinant Ward identities, separate pure-gauge from
transverse directions, and test whether the transverse quadratic response
can be represented by a curvature kernel without paying
\(\sigma_K^{-2}\).


\clearpage
\section{V52: determinant Ward identities and a curvature-kernel criterion}
\label{sec:v6v52}

V51 shows that a raw determinant Hessian loses two inverse powers of the
least Dirac singular value.  Before applying that bound, gauge covariance
must be used exactly.  This version proves the resulting Ward identities,
shows that they descend the Hessian to curvature space, and identifies the
additional locality moment which makes that descent uniform at small
momentum.

\subsection{Finite gauge-covariant determinant identities}

Let \(A\in C^1(K;\mathbb R)\) and \(\phi\in C^0(K;\mathbb R)\).  Suppose a
finite chiral matrix family obeys
\begin{equation}
 D(A+d_0\phi)
 =G_-(\phi)D(A)G_+(\phi)^{-1},
 \label{eq:v6v52-gauge-covariance}
\end{equation}
where \(G_\pm(\phi)\) are unitary.  On an invertible chart write
\(W(A)=-\log|\det D(A)|\) and choose a local phase
\(\Theta(A)=-\operatorname{Im}\log\det D(A)\).

\begin{theorem}[Exact modulus Ward identities]
\label{thm:v6v52-modulus-Ward}
For every \(A,X,\phi\) in the chart,
\begin{align}
 W(A+d_0\phi)&=W(A),
 \label{eq:v6v52-modulus-invariance}\\
 dW_A[d_0\phi]&=0,
 \label{eq:v6v52-first-Ward}\\
 d^2W_A[X,d_0\phi]&=0.
 \label{eq:v6v52-second-Ward}
\end{align}
Thus the real determinant Hessian annihilates every pure-gauge direction in
either slot.
\end{theorem}

\begin{proof}
Taking determinants in \eqref{eq:v6v52-gauge-covariance} gives
\[
 |\det D(A+d_0\phi)|
 =|\det G_-(\phi)|\,|\det D(A)|\,
  |\det G_+(\phi)|^{-1}
 =|\det D(A)|,
\]
because both gauge matrices are unitary.  This proves
\eqref{eq:v6v52-modulus-invariance}.  Differentiate it first in the gauge
parameter and then in the base point along \(X\).
\end{proof}

\begin{proposition}[Phase cocycle and the anomaly-free local Ward identity]
\label{prop:v6v52-phase-Ward}
On a fixed logarithm chart,
\begin{equation}
 \Theta(A+d_0\phi)-\Theta(A)
 =-\arg\det G_-(\phi)+\arg\det G_+(\phi)
 \pmod{2\pi}.
 \label{eq:v6v52-phase-cocycle}
\end{equation}
For a product of multiplets whose local determinant cocycle cancels, the
summed phase satisfies the analogues of
\eqref{eq:v6v52-first-Ward}--\eqref{eq:v6v52-second-Ward} on that chart.
This local statement does not trivialize the global determinant line.
\end{proposition}

\begin{proof}
Take the argument of the determinant identity following from
\eqref{eq:v6v52-gauge-covariance}.  If the product cocycle is locally one,
the right side is locally constant modulo \(2\pi\); differentiation gives
the two Ward identities.  Transition around distinct charts can retain
holonomy, so no global trivialization follows.
\end{proof}

\subsection{Exact finite-complex descent}

Fix flat and harmonic modes as in V50, so
\(d_1:\mathcal H_K\to\operatorname{im}d_1\) is an isomorphism.  Denote its
inverse by
\begin{equation}
 R_K:\operatorname{im}d_1\longrightarrow\mathcal H_K,
 \qquad R_Kd_1=I_{\mathcal H_K}.
 \label{eq:v6v52-curvature-reconstruction}
\end{equation}

\begin{theorem}[Every quotient Hessian has an exact curvature kernel]
\label{thm:v6v52-finite-factorization}
Let \(H\) be a symmetric bilinear form on \(C^1(K;\mathbb R)\) which
annihilates \(\ker d_1\) in both slots.  Define a bilinear form \(\Pi_H\)
on \(\operatorname{im}d_1\) by
\begin{equation}
 \Pi_H(F,G):=H(R_KF,R_KG).
 \label{eq:v6v52-PiH}
\end{equation}
Then
\begin{equation}
 H(X,Y)=\Pi_H(d_1X,d_1Y)
 \label{eq:v6v52-factorization}
\end{equation}
for all edge fields after their fixed flat part is removed.  Moreover,
\begin{equation}
 \|\Pi_H\|_{\rm op}
 \le\sigma_K^{-2}\|H|_{\mathcal H_K}\|_{\rm op}.
 \label{eq:v6v52-naive-factor-bound}
\end{equation}
\end{theorem}

\begin{proof}
Write \(X=X_0+X_\perp\) with \(X_0\in\ker d_1\) and
\(X_\perp\in\mathcal H_K\), and similarly for \(Y\).  The annihilation
hypothesis removes all terms containing \(X_0\) or \(Y_0\), while
\(X_\perp=R_Kd_1X\) and \(Y_\perp=R_Kd_1Y\).  This proves
\eqref{eq:v6v52-factorization}.  The singular values of \(R_K\) are the
reciprocals of those of \(d_1|_{\mathcal H_K}\), so
\(\|R_K\|=\sigma_K^{-1}\).  Apply this twice in
\eqref{eq:v6v52-PiH}.
\end{proof}

\begin{corollary}[Ward descent is exact but not uniformly bounded]
\label{cor:v6v52-Ward-not-bound}
The Hessian of the determinant modulus in
Theorem~\ref{thm:v6v52-modulus-Ward} factors through curvature.  Gauge
covariance alone yields only the bound
\eqref{eq:v6v52-naive-factor-bound}, which may lose
\(\sigma_K^{-2}\).
\end{corollary}

\begin{proof}
The Ward identities give the annihilation hypothesis of
Theorem~\ref{thm:v6v52-finite-factorization}.  Its norm estimate is then the
only estimate obtained without further structure.
\end{proof}

\begin{proposition}[Sharp gauge-invariant counterexample]
\label{prop:v6v52-Ward-counterexample}
For any \(c>0\), the quadratic functional
\begin{equation}
 W_c(A):=-{c\over2}\|P_{\mathcal H_K}A\|^2
 \label{eq:v6v52-nonlocal-quadratic}
\end{equation}
is invariant under addition of \(\ker d_1\) and satisfies both Ward
identities.  Its curvature kernel is
\begin{equation}
 \Pi_c=-cR_K^*R_K,
 \qquad \|\Pi_c\|={c\over\sigma_K^2}.
 \label{eq:v6v52-counterexample-kernel}
\end{equation}
Therefore Ward identities alone cannot prove the V50 relative-curvature
constant uniformly in volume.
\end{proposition}

\begin{proof}
The projection \(P_{\mathcal H_K}\) kills \(\ker d_1\), proving
invariance and the Ward identities.  The Hessian on \(\mathcal H_K\) is
\(-cI\).  Insert it into \eqref{eq:v6v52-PiH}; the operator representing the
result is \(-cR_K^*R_K\), whose norm is
\(c\|R_K\|^2=c\sigma_K^{-2}\).
\end{proof}

\subsection{Fourier factorization on a periodic lattice}

The counterexample is nonlocal.  We next isolate a sufficient locality
condition which removes its inverse-curl loss.  Work on a periodic
\(d\)-dimensional unit lattice and omit the zero momentum.  For
\(k\in(-\pi,\pi]^d\), set
\begin{equation}
 q_\mu(k):=e^{ik_\mu}-1,
 \qquad |q(k)|^2:=\sum_\mu|q_\mu(k)|^2.
 \label{eq:v6v52-q}
\end{equation}
The exterior derivative from one- to two-forms has symbol
\begin{equation}
 (D_1(k)a)_{\mu\nu}=q_\mu(k)a_\nu-q_\nu(k)a_\mu.
 \label{eq:v6v52-curl-symbol}
\end{equation}
On the transverse space \(q(k)^*a=0\),
\begin{equation}
 D_1(k)^*D_1(k)=|q(k)|^2I.
 \label{eq:v6v52-curl-isometry}
\end{equation}

\begin{theorem}[Uniform Fourier curvature factorization]
\label{thm:v6v52-Fourier-factorization}
Let \(H(k)\) be the Hermitian Fourier symbol of a translation-invariant real
edge-field Hessian.  Suppose for every nonzero momentum
\begin{equation}
 H(k)q(k)=0,
 \qquad
 \|H(k)\|_{\rm op}\le C_2|q(k)|^2.
 \label{eq:v6v52-transverse-quadratic}
\end{equation}
Then there is a Hermitian two-form symbol \(\Pi(k)\) such that
\begin{equation}
 H(k)=D_1(k)^*\Pi(k)D_1(k),
 \qquad
 \|\Pi(k)\|_{\rm op}\le C_2.
 \label{eq:v6v52-uniform-Pi}
\end{equation}
Consequently the quadratic form obeys
\begin{equation}
 \langle X,HX\rangle\ge-C_2\|d_1X\|^2.
 \label{eq:v6v52-relative-bound}
\end{equation}
\end{theorem}

\begin{proof}
Hermiticity and \(Hq=0\) imply that \(H\) is supported on the transverse
space.  Define
\begin{equation}
 R(k):=|q(k)|^{-2}D_1(k)^*
 \label{eq:v6v52-Rk}
\end{equation}
on \(\operatorname{im}D_1(k)\), and set
\(\Pi(k)=R(k)^*H(k)R(k)\), extended by zero on the orthogonal complement.
Equation \eqref{eq:v6v52-curl-isometry} gives
\(R(k)D_1(k)=P_{q(k)^\perp}\) and
\(\|R(k)\|=|q(k)|^{-1}\).  Hence
\[
 D_1^*\Pi D_1
 =P_{q^\perp}HP_{q^\perp}=H,
 \qquad
 \|\Pi\|\le |q|^{-2}\|H\|\le C_2.
\]
The lower form bound follows by Fourier summation.
\end{proof}

\subsection{A real-space second-moment criterion}

Let the convolution kernel of \(H\) be \(h(z)\), with the displacement
\(z\) chosen in a centered fundamental cell.  Define
\begin{equation}
 M_2(h):=\sum_z|z|^2\|h(z)\|_{\rm op}.
 \label{eq:v6v52-M2}
\end{equation}

\begin{lemma}[Second moment implies quadratic infrared vanishing]
\label{lem:v6v52-second-moment}
Assume
\begin{equation}
 \sum_zh(z)=0,
 \qquad
 \sum_zz_\alpha h(z)=0\quad(1\le\alpha\le d),
 \qquad M_2(h)<\infty.
 \label{eq:v6v52-moment-cancellations}
\end{equation}
Then
\begin{equation}
 \|H(k)\|_{\rm op}
 \le {\pi^2\over8}M_2(h)|q(k)|^2.
 \label{eq:v6v52-H-quadratic}
\end{equation}
The first-moment cancellation follows in particular when
\(h(z)=h(-z)\).
\end{lemma}

\begin{proof}
Using \eqref{eq:v6v52-moment-cancellations},
\[
 H(k)=\sum_zh(z)\bigl(e^{-ik\cdot z}-1+ik\cdot z\bigr).
\]
The scalar remainder is bounded by \(|k\cdot z|^2/2\), so
\(\|H(k)\|\le M_2(h)|k|^2/2\).  For
\(|k_\mu|\le\pi\),
\(2|\sin(k_\mu/2)|\ge2|k_\mu|/\pi\).  Thus
\(|k|^2\le(\pi^2/4)|q(k)|^2\), proving
\eqref{eq:v6v52-H-quadratic}.  Pairing \(z\) with \(-z\) proves the final
claim.
\end{proof}

\begin{corollary}[Local Ward curvature card]
\label{cor:v6v52-local-Ward-card}
If a determinant-modulus Hessian is translation invariant, obeys the Ward
identity, has zero toron response \(H(0)=0\), is even, and satisfies
\(M_2(h)\le M\) uniformly in the cutoff and volume, then
\begin{equation}
 d^2W[X,X]\ge-{\pi^2M\over8}\|d_1X\|^2.
 \label{eq:v6v52-local-Ward-relative}
\end{equation}
It therefore contributes
\(\kappa_{\rm f}\le\pi^2M/8\) to the V50 relative Maxwell curvature card.
\end{corollary}

\begin{proof}
Apply Lemma~\ref{lem:v6v52-second-moment} and
Theorem~\ref{thm:v6v52-Fourier-factorization} with
\(C_2=\pi^2M/8\).
\end{proof}

\begin{assessment}[Mass gap versus the second moment]
\label{ass:v6v52-gap-locality}
A uniform fermion singular gap can support exponential off-diagonal decay
of resolvents and hence a finite second moment, but that implication needs a
separate locality proof for the chosen lattice Dirac operator.  Massless
fermions do not provide such a gap.  The correct acquisition target is the
second moment of the \emph{combined transverse polarization kernel}, not a
termwise \(\rho^{-2}\) bound from V51.  This is the determinant analogue of
the companion YM instruction to retain the complete activity before norms.
\end{assessment}

\begin{firewall}[Torons and logarithmic running are not removed by Ward]
\label{fw:v6v52-toron-running}
The Ward identity imposes no condition at \(q=0\).  Zero toron response must
be proved or fixed as a sector condition.  Even for \(q\ne0\), transversality
allows
\(H(k)=|q(k)|^2\ell(k)P_{q(k)^\perp}\) with an unbounded scalar
\(\ell(k)\).  Thus a logarithmically growing polarization coefficient is
fully compatible with gauge invariance.  The uniform second-moment bound in
Corollary~\ref{cor:v6v52-local-Ward-card} is substantive and cannot be
replaced by the Ward identity alone.
\end{firewall}

\begin{definition}[Transverse polarization moment card]
\label{def:v6v52-TPMC}
TPMC asserts, after combining all retained matter multiplets and before
taking absolute values, that the real determinant Hessian has a
translation-covariant kernel satisfying the Ward identity, zero toron
response in the selected sector, and a uniform second moment
\begin{equation}
 \sup_{a,L}M_2(h_{a,L})<\infty.
 \label{eq:v6v52-TPMC}
\end{equation}
The phase kernel is handled separately by determinant-line transport.
\end{definition}

\begin{theorem}[Sixth-series V52 result]
\label{thm:v6v52-status}
Gauge covariance makes the determinant modulus exactly constant on pure
gauge directions and descends its Hessian to a curvature kernel.  This
algebraic descent alone may cost \(\sigma_K^{-2}\), and the invariant
quadratic example \eqref{eq:v6v52-nonlocal-quadratic} proves the loss sharp.
For a translation-invariant Hessian, Ward transversality plus quadratic
infrared vanishing yields a uniformly bounded curvature kernel.  Zero toron
response, evenness, and a uniform real-space second moment are sufficient
for that vanishing.  TPMC is therefore a precise upstream replacement for
the raw inverse-Dirac-gap estimate of V51.
\end{theorem}

\begin{proof}
Use Theorem~\ref{thm:v6v52-modulus-Ward},
Proposition~\ref{prop:v6v52-phase-Ward},
Theorem~\ref{thm:v6v52-finite-factorization},
Proposition~\ref{prop:v6v52-Ward-counterexample},
Theorem~\ref{thm:v6v52-Fourier-factorization}, and
Lemma~\ref{lem:v6v52-second-moment}, subject to
Firewall~\ref{fw:v6v52-toron-running}.
\end{proof}

V53 will test TPMC on an exactly solvable massive lattice-fermion reference.
It will derive resolvent off-diagonal decay from a finite hopping range and
a singular gap, then convert that decay into an explicit polarization
second moment.  The massless limit will be kept visible rather than absorbed
into a constant.


\clearpage
\section{V53: massive hopping resolvents and a determinant locality moment}
\label{sec:v6v53}

V52 replaces a raw inverse-Dirac-gap Hessian bound by TPMC, a second-moment
bound on the complete transverse polarization kernel.  This version proves
that card in a strong-mass finite-range reference.  The proof is elementary:
a norm-convergent hopping expansion gives exponential off-diagonal decay,
and two local gauge vertices convert that decay into a summable determinant
Hessian.  The constants expose the degeneration at the massless boundary.

\subsection{Finite-range hopping hypothesis}

Let \(\Lambda_L=(\mathbb Z/L\mathbb Z)^d\) carry the periodic graph
distance.  At every site attach a complex internal space of dimension
\(r_0\).  Consider
\begin{equation}
 D_A=m(I+T_A),
 \qquad m>0,
 \label{eq:v6v53-hopping-D}
\end{equation}
with the following assumptions, uniform in \(L\) and in the background
\(A\) inside a fixed chart:
\begin{enumerate}[label=\textup{(H\arabic*)},leftmargin=4.8em]
\item \(T_A(x,y)=0\) if \(\operatorname{dist}(x,y)>R\);
\item \(\|T_A\|_{\rm op}\le\theta<1\);
\item for an edge coordinate \(A_e\), the vertices
      \(V_e:=\partial_{A_e}T_A\) and
      \(V_{ef}:=\partial_{A_e}\partial_{A_f}T_A\) have operator norm at
      most \(c_1\) and \(c_2\), respectively, have rank at most \(r_1\),
      and are supported within distance \(R_1\) of their indexed edges;
\item \(V_{ef}=0\) when \(\operatorname{dist}(e,f)>R_2\); and
\item the family is gauge covariant as in
      \eqref{eq:v6v52-gauge-covariance}.
\end{enumerate}
Here support means that both the input and output site projections of the
vertex lie in the stated neighborhood.

\begin{lemma}[Uniform singular gap]
\label{lem:v6v53-singular-gap}
Every \(D_A\) in the chart is invertible and
\begin{equation}
 s_{\min}(D_A)\ge m(1-\theta).
 \label{eq:v6v53-gap}
\end{equation}
\end{lemma}

\begin{proof}
For every vector \(u\),
\[
 \|D_Au\|=m\|(I+T_A)u\|
 \ge m(\|u\|-\|T_Au\|)
 \ge m(1-\theta)\|u\|.
\]
The same estimate and finite dimensionality give invertibility.
\end{proof}

\subsection{Off-diagonal inverse from the hopping expansion}

For a set \(S\subset\Lambda_L\), write \(P_S\) for its site projection and
\(d(S,S')\) for the graph distance between sets.

\begin{theorem}[Explicit hopping Combes--Thomas bound]
\label{thm:v6v53-resolvent-decay}
For all site sets \(S,S'\),
\begin{equation}
 \|P_SD_A^{-1}P_{S'}\|_{\rm op}
 \le {1\over m(1-\theta)}
 \theta^{\left\lceil d(S,S')/R\right\rceil}.
 \label{eq:v6v53-resolvent-decay}
\end{equation}
\end{theorem}

\begin{proof}
The Neumann series converges in operator norm:
\[
 D_A^{-1}=m^{-1}\sum_{n=0}^\infty(-T_A)^n.
\]
Finite range implies
\(P_ST_A^nP_{S'}=0\) whenever \(nR<d(S,S')\).  If
\(N=\lceil d(S,S')/R\rceil\), the remaining tail has norm at most
\[
 m^{-1}\sum_{n=N}^\infty\theta^n
 ={\theta^N\over m(1-\theta)}.
\]
\end{proof}

\begin{remark}[Visible massless degeneration]
\label{rem:v6v53-massless-degeneration}
The decay rate is
\(\eta_\theta=R^{-1}|\log\theta|\).  Both the prefactor and the decay
length diverge as \(\theta\uparrow1\).  No estimate below is uniform at
that boundary.
\end{remark}

\subsection{Polarization kernel of the determinant modulus}

Let
\begin{equation}
 W_L(A):=-\log|\det D_A|
 \label{eq:v6v53-WL}
\end{equation}
and define its edge Hessian kernel at a translation-invariant background by
\begin{equation}
 h_L(e,f):=\partial_{A_e}\partial_{A_f}W_L(A).
 \label{eq:v6v53-hL}
\end{equation}
The V51 Jacobi formula and \(D_e=mV_e\), \(D_{ef}=mV_{ef}\) give
\begin{equation}
 h_L(e,f)
 =\operatorname{Re}\operatorname{Tr}
 \left(D_A^{-1}mV_fD_A^{-1}mV_e-D_A^{-1}mV_{ef}\right).
 \label{eq:v6v53-hessian-vertices}
\end{equation}

\begin{lemma}[Two-vertex exponential bound]
\label{lem:v6v53-two-vertex}
There are constants \(C_0<\infty\) and \(R_0<\infty\), depending only on
\(r_0,r_1,R,R_1,R_2,c_1,c_2\), such that
\begin{equation}
 |h_L(e,f)|
 \le {C_0\over(1-\theta)^2}
 \theta^{2\left\lceil
 (\operatorname{dist}(e,f)-R_0)_+/R\right\rceil}
 +{C_0\over1-\theta}
 \mathbf 1_{\{\operatorname{dist}(e,f)\le R_2\}}.
 \label{eq:v6v53-hessian-decay}
\end{equation}
The estimate is independent of \(m\), because each differentiated hopping
vertex carries the compensating factor \(m\).
\end{lemma}

\begin{proof}
Let \(S_e,S_f\) be the supports of \(V_e,V_f\).  Cyclicity of trace inserts
the support projections into the first term of
\eqref{eq:v6v53-hessian-vertices}.  The two inverse factors connect
\(S_e\) to \(S_f\) and back.  Their two factors of \(m^{-1}\) cancel the
two factors of \(m\) from the vertices.  Apply
Theorem~\ref{thm:v6v53-resolvent-decay} twice, use
\(|\operatorname{Tr}B|\le\operatorname{rank}(B)\|B\|_{\rm op}\), and
absorb the fixed support diameters and ranks into \(C_0,R_0\).  For the
second Jacobi term, Assumption \textup{(H4)} makes the expression zero
outside distance \(R_2\); one inverse bound gives the second term in
\eqref{eq:v6v53-hessian-decay}.
\end{proof}

\begin{theorem}[Uniform determinant second moment in the strong-mass
region]
\label{thm:v6v53-M2}
Fix an edge \(e_0\), identify \(f\) with its centered displacement from
\(e_0\), and put
\begin{equation}
 M_{2,L}:=\sum_f\operatorname{dist}(e_0,f)^2|h_L(e_0,f)|.
 \label{eq:v6v53-M2L}
\end{equation}
There is a dimension-dependent constant \(C_d\) such that
\begin{equation}
 \sup_LM_{2,L}
 \le C_dC_0(1+R_0+R_2)^{d+2}
 \left[
 {1\over(1-\theta)^2(1-\theta^{2/R})^{d+2}}
 +{1\over1-\theta}
 \right].
 \label{eq:v6v53-M2-bound}
\end{equation}
In particular the determinant polarization kernel has a volume-uniform
second moment for every fixed \(\theta<1\).
\end{theorem}

\begin{proof}
The number of edges at graph distance in \([n,n+1)\) from a fixed edge is
at most \(C_d(1+n)^{d-1}\), uniformly on every torus.  Insert
\eqref{eq:v6v53-hessian-decay}.  After enlarging constants to account for
\(R_0\), the nonlocal sum is bounded by
\[
 {C_dC_0(1+R_0)^{d+2}\over(1-\theta)^2}
 \sum_{n\ge0}(1+n)^{d+1}(\theta^{2/R})^n.
\]
For \(0<a<1\), comparison with derivatives of the geometric series gives
\(\sum_{n\ge0}(1+n)^{d+1}a^n\le C_d(1-a)^{-(d+2)}\).
The finite-range second-Jacobi term contributes at most
\(C_dC_0(1+R_2)^{d+2}/(1-\theta)\).  Combining the two bounds proves
\eqref{eq:v6v53-M2-bound}.
\end{proof}

\subsection{Ward cancellation and toron loops}

At finite periodic volume a constant edge field is a toron rather than a
periodic pure gauge.  It is therefore useful to separate non-winding and
winding hopping loops.  The norm-convergent identity
\begin{equation}
 \log\det D_A
 =r_0|\Lambda_L|\log m
 +\sum_{n=1}^\infty{(-1)^{n+1}\over n}
 \operatorname{Tr}(T_A^n)
 \label{eq:v6v53-loop-expansion}
\end{equation}
holds on the determinant chart.

\begin{lemma}[Only winding loops see a constant toron]
\label{lem:v6v53-toron-loops}
Assume the gauge factors in \(T_A\) are parallel transports along paths of
length at most \(R\).  Every contractible closed walk contributing to
\(\operatorname{Tr}(T_A^n)\) is unchanged by addition of a constant flat
connection.  A toron-dependent term must wind around the periodic lattice
and therefore has \(nR\ge L\).
\end{lemma}

\begin{proof}
The gauge factor around a contractible closed walk is its holonomy.  A
constant Abelian connection has zero curvature, so its holonomy on such a
walk is one.  A nontrivial torus winding has graph length at least \(L\),
while each hopping factor advances by at most \(R\); hence \(nR\ge L\).
\end{proof}

\begin{proposition}[Exponential toron Hessian remainder]
\label{prop:v6v53-toron-remainder}
Under the hypotheses above, the Hessian density in a constant toron
direction is bounded by
\begin{equation}
 C_{\rm tor}(1+L/R)^2{\theta^{\lceil L/R\rceil}\over(1-\theta)^3},
 \label{eq:v6v53-toron-remainder}
\end{equation}
where \(C_{\rm tor}\) depends only on the local hopping data.  Consequently
the toron response tends to zero exponentially as \(L\to\infty\) for fixed
\(\theta<1\).
\end{proposition}

\begin{proof}
Differentiate \eqref{eq:v6v53-loop-expansion} twice in a constant flat
direction.  For a length-\(n\) trace there are at most \(n^2\) ordered
placements of the two derivatives, while the prefactor is \(1/n\) and the
operator product is bounded by a local constant times \(\theta^n\).
Lemma~\ref{lem:v6v53-toron-loops} removes every term with
\(n<\lceil L/R\rceil\).  The elementary tail estimate
\[
 \sum_{n\ge N}n\theta^n
 \le C(1+N)^2{\theta^N\over(1-\theta)^2}
\]
proves \eqref{eq:v6v53-toron-remainder} after allowing one additional
factor \((1-\theta)^{-1}\) for the local derivative bound.  Division by
volume is implicit in the Hessian density and removes the base-point count
in the trace.
\end{proof}

\begin{theorem}[TPMC for the massive hopping reference]
\label{thm:v6v53-TPMC}
Assume in addition reflection evenness of the translation-invariant
polarization kernel.  In infinite volume, or after placing the exponentially
small finite-volume toron response in the retained flat sector, the
determinant modulus satisfies TPMC.  Its curvature-kernel norm is bounded by
\begin{equation}
 \kappa_{\rm hop}
 \le {\pi^2\over8}\sup_LM_{2,L},
 \label{eq:v6v53-kappa-hop}
\end{equation}
with the explicit right side from \eqref{eq:v6v53-M2-bound}.
\end{theorem}

\begin{proof}
Gauge covariance supplies the Ward identity by
Theorem~\ref{thm:v6v52-modulus-Ward}.  Reflection evenness supplies the
first-moment cancellation in Lemma~\ref{lem:v6v52-second-moment}.  Theorem
\ref{thm:v6v53-M2} supplies the uniform second moment.  Infinite-volume
constant flat fields are locally pure gauge, while the finite-volume
difference is the winding remainder bounded in
Proposition~\ref{prop:v6v53-toron-remainder}.  Apply
Corollary~\ref{cor:v6v52-local-Ward-card} to obtain
\eqref{eq:v6v53-kappa-hop}.
\end{proof}

\begin{corollary}[A proved interacting Maxwell window]
\label{cor:v6v53-Maxwell-window}
If the Higgs real remainder has relative negative curvature
\(\kappa_{\rm H}\), and
\begin{equation}
 \kappa_{\rm H}+\kappa_{\rm hop}<\beta_Y,
 \label{eq:v6v53-window}
\end{equation}
then the positive reference law formed from the noncompact Maxwell action,
the Higgs modulus, and the massive hopping determinant modulus satisfies the
V50 curvature-observable Poincare inequality with constant
\((\beta_Y-\kappa_{\rm H}-\kappa_{\rm hop})^{-1}\).
\end{corollary}

\begin{proof}
Theorem~\ref{thm:v6v53-TPMC} and the assumed Higgs bound add to the relative
curvature hypothesis of Theorem~\ref{thm:v6v50-RMC}.
\end{proof}

\begin{firewall}[The reference window is not the physical chiral limit]
\label{fw:v6v53-not-chiral-limit}
The condition \(\theta<1\) is a strong-mass or small-hopping condition for a
finite regulated determinant.  The constants in
\eqref{eq:v6v53-M2-bound} diverge as \(\theta\uparrow1\), and no massless
chiral continuum limit follows.  The phase and global determinant line also
remain outside the positive reference law.  The theorem proves a genuine
interacting reference window and a mechanism for TPMC; it does not prove
that the Standard Model trajectory stays in that window.
\end{firewall}

\begin{assessment}[Information imported from the YM audit]
\label{ass:v6v53-YM-lesson}
The proof keeps the two gauge vertices connected by the complete resolvent
loop until exponential decay has been extracted.  Bounding each inverse by
its global norm first would reproduce the V51 factor \(\rho^{-2}\) and lose
all separation information.  This is the determinant counterpart of the YM
V36 instruction not to triangle-bound the ternary normalized activity
before its spectral-sector cancellation is exposed.
\end{assessment}

\begin{theorem}[Sixth-series V53 result]
\label{thm:v6v53-status}
For a gauge-covariant finite-range hopping determinant with
\(\|T_A\|\le\theta<1\), the inverse has the explicit exponential bound
\eqref{eq:v6v53-resolvent-decay}.  Its two-vertex determinant Hessian has a
volume-uniform real-space second moment, and hence a uniformly bounded
curvature kernel after the Ward identity and toron separation.  This proves
TPMC and a nonempty interacting Maxwell curvature window for the massive
hopping reference.  Every constant degenerates visibly at the massless
boundary, which remains open.
\end{theorem}

\begin{proof}
Use Lemma~\ref{lem:v6v53-singular-gap},
Theorem~\ref{thm:v6v53-resolvent-decay},
Lemma~\ref{lem:v6v53-two-vertex},
Theorem~\ref{thm:v6v53-M2},
Proposition~\ref{prop:v6v53-toron-remainder}, and
Theorem~\ref{thm:v6v53-TPMC}, subject to
Firewall~\ref{fw:v6v53-not-chiral-limit}.
\end{proof}

V54 will approach the massless boundary shellwise.  It will retain the two
resolvents in the exact V51 singular-shell trace, derive a separated-shell
polarization ledger, and identify the precise weighted summability condition
which replaces the uniform hopping parameter \(\theta<1\).


\clearpage
\section{V54: shellwise polarization and the logarithmic Maxwell counterterm}
\label{sec:v6v54}

V53 proves TPMC for a finite-range determinant with a fixed strong-mass
parameter \(\theta<1\).  The bound degenerates when the fermion correlation
length diverges.  This version resolves the first massless test shell by
shell.  A proper-time polarization model shows that the remaining
four-dimensional loss is logarithmic, that every active scale contributes
an order-one amount, and that subtraction of a local Maxwell counterterm is
necessary before any cutoff-uniform curvature card can be formulated.

\subsection{A gauge-invariant transverse bubble model}

For \(\Lambda>0\), \(m\ge0\), and Euclidean external momentum
\(p\in\mathbb R^4\), define
\begin{equation}
 \mathcal B_{\Lambda,m}(p)
 :={1\over16\pi^2}\int_0^1\int_{\Lambda^{-2}}^\infty
 {e^{-sM_x(p)^2}\over s}\,ds\,dx,
 \qquad
 M_x(p)^2:=m^2+x(1-x)|p|^2.
 \label{eq:v6v54-bubble}
\end{equation}
For \(m=0\) we take \(p\ne0\).  The integral is finite: the logarithmic
endpoint singularities at \(x=0,1\) are integrable.

Let \(Q_2=\sum_fn_fY_f^2=10/3\) be the one-generation square charge from
V51.  The associated transverse quadratic functional is
\begin{equation}
 \mathcal W_{\Lambda,m}^{\rm pol}(A)
 :={Q_2\over2}\int_{\mathbb R^4}
 \mathcal B_{\Lambda,m}(p)|\widehat F(p)|^2\,dp,
 \qquad F=dA,
 \label{eq:v6v54-polarization-model}
\end{equation}
whenever \(F\) is a Schwartz two-form.  The precise tensor numerator of a
regulated chiral determinant is not asserted to equal
\eqref{eq:v6v54-polarization-model}; this model isolates its common
two-resolvent denominator and is an exact test of the absolute shell card.

\begin{lemma}[Exponential-integral representation]
\label{lem:v6v54-E1}
Let
\begin{equation}
 E_1(z):=\int_z^\infty{e^{-t}\over t}\,dt,
 \qquad z>0.
 \label{eq:v6v54-E1}
\end{equation}
Then
\begin{equation}
 \mathcal B_{\Lambda,m}(p)
 ={1\over16\pi^2}\int_0^1
 E_1\!\left({M_x(p)^2\over\Lambda^2}\right)dx.
 \label{eq:v6v54-bubble-E1}
\end{equation}
For \(0<z\le1\),
\begin{equation}
 e^{-1}\log{1\over z}
 \le E_1(z)
 \le \log{1\over z}+E_1(1).
 \label{eq:v6v54-E1-bounds}
\end{equation}
\end{lemma}

\begin{proof}
In the inner integral of \eqref{eq:v6v54-bubble}, substitute
\(t=sM_x(p)^2\).  This gives \eqref{eq:v6v54-bubble-E1}.  On
\([z,1]\), \(e^{-1}\le e^{-t}\le1\); integrate \(dt/t\) and add the
positive tail \(E_1(1)\) for the upper bound.
\end{proof}

\subsection{Every proper-time scale contributes}

Partition the proper-time integral into the intervals
\begin{equation}
 I_j(\Lambda):=
 [4^j\Lambda^{-2},4^{j+1}\Lambda^{-2}],
 \qquad j=0,1,2,\ldots.
 \label{eq:v6v54-proper-shells}
\end{equation}
For a fixed \(x\), let
\begin{equation}
 N_x(\Lambda,p,m):=
 \max\left\{0,
 \left\lfloor\log_4{\Lambda^2\over M_x(p)^2}\right\rfloor\right\}.
 \label{eq:v6v54-shell-count}
\end{equation}

\begin{theorem}[Order-one contribution per active shell]
\label{thm:v6v54-shell-contribution}
For every \(0\le j<N_x(\Lambda,p,m)\),
\begin{equation}
 e^{-1}\log4
 \le\int_{I_j(\Lambda)}{e^{-sM_x(p)^2}\over s}\,ds
 \le\log4.
 \label{eq:v6v54-one-shell}
\end{equation}
Consequently the ultraviolet part of the bubble is comparable to the number
of active logarithmic scales:
\begin{equation}
 {e^{-1}\log4\over16\pi^2}
 \int_0^1N_x\,dx
 \le\mathcal B_{\Lambda,m}(p),
 \label{eq:v6v54-shell-lower}
\end{equation}
and the sum of the active shells is at most
\begin{equation}
 {\log4\over16\pi^2}
 \int_0^1N_x\,dx.
 \label{eq:v6v54-shell-upper}
\end{equation}
\end{theorem}

\begin{proof}
If \(j<N_x\), then
\(4^{j+1}\Lambda^{-2}M_x^2\le1\).  Hence
\(e^{-1}\le e^{-sM_x^2}\le1\) throughout \(I_j\).  Since
\(\int_{I_j}ds/s=\log4\), this proves
\eqref{eq:v6v54-one-shell}.  Sum in \(j\), integrate in \(x\), and retain
or discard the positive inactive tail as stated.
\end{proof}

\begin{assessment}[The logarithm is an incidence count]
\label{ass:v6v54-incidence-count}
The two resolvents do not generate a power divergence after the transverse
curvature factors have been extracted.  They leave one constant payment per
proper-time scale.  Taking absolute values and summing those payments gives
the logarithm.  This is the determinant analogue of the NS V31 first-moment
warning: a formally small marked shell does not remain small after every
active scale is unmarked.
\end{assessment}

\subsection{Cutoff asymptotics and local subtraction}

\begin{lemma}[Small-argument limit of \(E_1\)]
\label{lem:v6v54-E1-limit}
There is a finite constant \(\gamma_E\) such that
\begin{equation}
 E_1(z)+\log z\longrightarrow-\gamma_E
 \qquad(z\downarrow0).
 \label{eq:v6v54-E1-limit}
\end{equation}
Moreover, for \(0<z\le1\),
\begin{equation}
 |E_1(z)+\log z+\gamma_E|\le z.
 \label{eq:v6v54-E1-error}
\end{equation}
\end{lemma}

\begin{proof}
Write
\[
 E_1(z)+\log z
 =\int_z^1{e^{-t}-1\over t}\,dt+\int_1^\infty{e^{-t}\over t}\,dt.
\]
The first integrand extends continuously to \(-1\) at zero, so the right
side has a finite limit; define its negative to be \(\gamma_E\).  The
difference from the limit is
\(\int_0^z(1-e^{-t})dt/t\), whose modulus is at most \(z\).
\end{proof}

\begin{theorem}[Logarithmic divergence and renormalized difference]
\label{thm:v6v54-renormalized-bubble}
If \(m>0\), then for fixed \(p\),
\begin{equation}
 \mathcal B_{\Lambda,m}(p)
 ={1\over16\pi^2}\left[
 \log\Lambda^2-\gamma_E
 -\int_0^1\log M_x(p)^2\,dx\right]+o(1)
 \label{eq:v6v54-bubble-asymptotic}
\end{equation}
as \(\Lambda\to\infty\).  If \(p,p_*\ne0\) and \(m\ge0\), then the
subtracted bubble has the finite limit
\begin{equation}
 \lim_{\Lambda\to\infty}
 \left[\mathcal B_{\Lambda,m}(p)
       -\mathcal B_{\Lambda,m}(p_*)\right]
 ={1\over16\pi^2}\int_0^1
 \log{M_x(p_*)^2\over M_x(p)^2}\,dx.
 \label{eq:v6v54-renormalized-limit}
\end{equation}
For \(m=0\), the right side equals
\((16\pi^2)^{-1}\log(|p_*|^2/|p|^2)\).
\end{theorem}

\begin{proof}
For \(m>0\), the arguments \(M_x(p)^2/\Lambda^2\) tend to zero uniformly
in \(x\).  Insert Lemma~\ref{lem:v6v54-E1-limit} into
\eqref{eq:v6v54-bubble-E1}; its uniform error is
\(O((m^2+|p|^2)/\Lambda^2)\).  This proves
\eqref{eq:v6v54-bubble-asymptotic}.

For the difference, the constant \(\log\Lambda^2-\gamma_E\) cancels.  If
\(m>0\), the same uniform argument applies.  When \(m=0\), split the
\(x\)-integral into \([\epsilon,1-\epsilon]\) and its two endpoint pieces.
On the middle interval convergence is uniform.  Lemma
\ref{lem:v6v54-E1} bounds each endpoint piece by a constant times
\(\int_0^\epsilon(1+|\log x|)dx\), uniformly after taking the difference;
this tends to zero with \(\epsilon\).  Dominated convergence therefore
proves \eqref{eq:v6v54-renormalized-limit}.  For \(m=0\), the factors
\(x(1-x)\) cancel inside the logarithmic ratio, giving the last formula.
\end{proof}

\begin{corollary}[Unrenormalized TPMC fails in the bubble test]
\label{cor:v6v54-unrenormalized-fails}
The curvature-kernel coefficient of
\(\mathcal W_{\Lambda,m}^{\rm pol}\) grows as
\begin{equation}
 {Q_2\over16\pi^2}\log\Lambda^2+O(1)
 \label{eq:v6v54-charge-log}
\end{equation}
for fixed \(m>0,p\).  Therefore no cutoff-uniform absolute TPMC bound can
hold for this unrenormalized test family.  The anomaly cancellations
\(\sum n_fY_f=\sum n_fY_f^3=0\) do not remove the logarithm because its
coefficient is \(Q_2=10/3\).
\end{corollary}

\begin{proof}
Multiply \eqref{eq:v6v54-bubble-asymptotic} by \(Q_2\) and use
Lemma~\ref{lem:v6v51-charge-sums}.
\end{proof}

\subsection{Renormalized transverse polarization card}

Fix a nonzero subtraction momentum \(p_*\).  Define the local Maxwell
counterterm coefficient
\begin{equation}
 c_{\Lambda,m}(p_*):=Q_2\mathcal B_{\Lambda,m}(p_*)
 \label{eq:v6v54-counterterm}
\end{equation}
and the subtracted curvature kernel
\begin{equation}
 \mathcal K_{\Lambda,m}^{\rm ren}(p;p_*)
 :=Q_2\left[
 \mathcal B_{\Lambda,m}(p)-\mathcal B_{\Lambda,m}(p_*)
ight].
 \label{eq:v6v54-renormalized-kernel}
\end{equation}

\begin{definition}[Renormalized TPMC]
\label{def:v6v54-rTPMC}
rTPMC asserts that the complete determinant polarization can be written as
\begin{equation}
 H_{\Lambda}^{\rm ferm}
 =c_\Lambda d_1^*d_1+d_1^*\Pi_\Lambda^{\rm ren}d_1
 +H_\Lambda^{\rm phase},
 \label{eq:v6v54-rTPMC-decomposition}
\end{equation}
where \(c_\Lambda\) is absorbed into the running Maxwell coefficient, the
real kernels \(\Pi_\Lambda^{\rm ren}\) are uniformly bounded on the chosen
renormalization window, and the phase is transported by the separate
determinant-line card.
\end{definition}

\begin{proposition}[rTPMC for the bubble model on a compact momentum window]
\label{prop:v6v54-bubble-rTPMC}
Let \(0<p_-\le|p|,|p_*|\le p_+<\infty\) and \(m\ge0\).  Then
\(\mathcal K_{\Lambda,m}^{\rm ren}(p;p_*)\) converges pointwise as
\(\Lambda\to\infty\) and is bounded uniformly in \(\Lambda\ge\Lambda_0\)
and in the stated momentum window.  Hence the transverse bubble model
satisfies rTPMC after the local counterterm
\eqref{eq:v6v54-counterterm} is extracted.
\end{proposition}

\begin{proof}
Pointwise convergence is Theorem~\ref{thm:v6v54-renormalized-bubble}.
The limiting logarithmic ratio is bounded on the compact window.  The proof
of that theorem, with \(|p|,|p_*|\in[p_-,p_+]\), gives a uniform dominated
majorant and a uniform small-argument error once \(\Lambda\ge\Lambda_0\).
The finitely bounded interval of smaller admissible cutoffs is controlled
directly from \eqref{eq:v6v54-bubble}.
\end{proof}

\begin{firewall}[Subtraction is not ultraviolet completion]
\label{fw:v6v54-not-UV-completion}
Extracting \(c_\Lambda\|F\|^2/2\) reorganizes the local logarithm into a
running inverse coupling.  It does not prove that the resulting bare
Maxwell coefficient remains positive at arbitrarily large cutoff, nor that
the full nonperturbative chiral determinant equals the bubble model.  The
hypercharge ultraviolet trajectory identified in V13--V14 therefore
remains an upstream obstruction.
\end{firewall}

\begin{firewall}[The sign needed by NHCC remains a complete-loop question]
\label{fw:v6v54-sign}
TPMC and rTPMC above are absolute kernel cards.  NHCC needs a lower bound on
the real action, for which the sign of the complete regulated fermion loop
matters.  The scalar denominator in \eqref{eq:v6v54-bubble} does not contain
the full Dirac numerator, regulator fields, Higgs dependence, or chiral
phase.  No lower-curvature sign for the physical Standard Model determinant
is inferred from this test alone.
\end{firewall}

\begin{theorem}[Sixth-series V54 result]
\label{thm:v6v54-status}
The four-dimensional two-resolvent polarization denominator has one
order-one contribution per active proper-time shell and hence a logarithmic
unrenormalized curvature coefficient.  Its coefficient is weighted by the
nonvanishing Standard Model square charge \(Q_2=10/3\), not by the cancelled
linear or cubic anomaly sums.  Subtracting its value at a fixed nonzero
momentum produces the finite limit
\eqref{eq:v6v54-renormalized-limit} and proves rTPMC for the transverse
bubble model on compact momentum windows.  Positivity of the running bare
coefficient and the complete chiral determinant card remain open.
\end{theorem}

\begin{proof}
Use Theorem~\ref{thm:v6v54-shell-contribution},
Theorem~\ref{thm:v6v54-renormalized-bubble},
Corollary~\ref{cor:v6v54-unrenormalized-fails}, and
Proposition~\ref{prop:v6v54-bubble-rTPMC}, subject to both firewalls.
\end{proof}

V55 will perform the compulsory YM/NS audit before continuing.  It will
compare their newest all-arity and stopping-time bottlenecks with the
one-payment-per-polarization-shell ledger above, then decide whether to
continue the hypercharge renormalized trajectory or move upstream to the
Higgs and gravitational positivity interfaces.


\clearpage
\section{V55: companion audit, exact subtraction cocycle, and the hypercharge horizon}
\label{sec:v6v55}

V54 finds one order-one contribution per active polarization shell and
extracts a local Maxwell counterterm in a transverse bubble model.  Before
using that extraction as part of the SM route, this version performs the
compulsory audit of the now newer YM and NS notes.  Their conclusions require
the subtraction to be made from the complete regulated loop before norms,
and they rule out treating the logarithm as a summable error.

\subsection{Compulsory V55 audit}

The audited companion sources were
\begin{center}
\begin{tabular}{p{0.25\linewidth}p{0.16\linewidth}p{0.47\linewidth}}
programme & source & newest relevant conclusion\\ \hline
Yang--Mills & \texttt{V39} & the calibrated cosine core has the required
double-factorial spectral coefficients, but the isolated Haar logarithm
fails EEMC; its complete antipodal jet is cancelled only by the paired-image
term, so the wrapped block must remain intact before all-order norms\\
Navier--Stokes & \texttt{V35} & slotwise and octave-wise absolute estimates
pay a constant per scale and cannot remove the type-I logarithm; the next
gate is the signed own-window flux, while the type-II scalar ledgers remain
jointly consistent
\end{tabular}
\end{center}
At the time of reading, the YM V39 source had \(18\,618\) logical lines,
\(588\,725\) bytes, and SHA--256
\[
 \texttt{026728ddc19aa468c78af58da8568899a375575eea2ddd35e845f0e6393747d6},
\]
while the NS V35 source had \(25\,472\) logical lines,
\(998\,459\) bytes, and SHA--256
\[
 \texttt{aea7449895cd39b2d866869048033b96a47140d037dd9e3738ebb26f3fa55c01}.
\]

\begin{theorem}[V55 transfer decision]
\label{thm:v6v55-audit-decision}
The companion audit permits exactly the following conclusions for the
present hypercharge branch.
\begin{enumerate}[label=\textup{(\roman*)},leftmargin=3.7em]
\item The V54 scalar denominator is a sharp logarithmic test but cannot be
      subtracted species by species or numerator term by numerator term.
      A counterterm is legitimate only after the complete regulated real
      polarization has been assembled and its local transverse coefficient
      identified.
\item An absolute per-shell estimate cannot by itself improve the V54
      logarithm.  Removal of the logarithm requires either exact signed
      cancellation inside the complete loop or renormalization of its local
      Maxwell part.
\item Neither companion supplies positivity of the running hypercharge
      coefficient, the global determinant line, or a nonperturbative
      ultraviolet endpoint.
\end{enumerate}
\end{theorem}

\begin{proof}
Item \textup{(i)} is the exact logical content of the YM V39 inseparability
theorem applied as a proof discipline: the full object must be assembled
before a cancellation-dependent norm is taken.  Item \textup{(ii)} is the
NS V35 constant-per-octave no-go together with the V54 proper-time shell
identity.  Item \textup{(iii)} is stated in the claim boundaries of both
companions and in the V54 firewalls.
\end{proof}

\begin{firewall}[The audit imports proof discipline, not estimates]
\label{fw:v6v55-audit-types}
YM wrapped heat kernels and NS signed energy fluxes are not determinant
polarization kernels.  No numerical constant is transferred between the
three programmes.  The audit changes the order of operations in the SM
argument and identifies a common logarithmic obstruction.
\end{firewall}

\subsection{Complete-loop subtraction}

Let \(K=d_1^*d_1\) on the fixed noncompact hypercharge quotient.  Suppose
the complete regulated real matter Hessian at cutoff \(\Lambda\), after all
species and regulator fields have been combined, admits the exact
decomposition
\begin{equation}
 H_\Lambda^{\rm mat}
 =c_\Lambda(\mu)K+d_1^*\Pi_\Lambda^{\rm ren}(\mu)d_1,
 \label{eq:v6v55-complete-subtraction}
\end{equation}
at a nonzero subtraction scale \(\mu\).  The first term is local; the second
contains the complete momentum-dependent remainder.  Define
\begin{equation}
 \beta_R(\mu):=\beta_{\rm bare}(\Lambda)+c_\Lambda(\mu).
 \label{eq:v6v55-betaR}
\end{equation}

\begin{theorem}[Exact subtraction cocycle]
\label{thm:v6v55-subtraction-cocycle}
Suppose \eqref{eq:v6v55-complete-subtraction} also holds at another
subtraction scale \(\nu\).  Then there is a scalar
\begin{equation}
 \delta c_\Lambda(\nu,\mu)
 :=c_\Lambda(\nu)-c_\Lambda(\mu)
 \label{eq:v6v55-delta-c}
\end{equation}
such that
\begin{align}
 \Pi_\Lambda^{\rm ren}(\nu)
 &=\Pi_\Lambda^{\rm ren}(\mu)-
   \delta c_\Lambda(\nu,\mu)I,
 \label{eq:v6v55-Pi-cocycle}\\
 \beta_R(\nu)
 &=\beta_R(\mu)+\delta c_\Lambda(\nu,\mu),
 \label{eq:v6v55-beta-cocycle}\\
 \delta c_\Lambda(\omega,\mu)
 &=\delta c_\Lambda(\omega,\nu)
  +\delta c_\Lambda(\nu,\mu).
 \label{eq:v6v55-cocycle}
\end{align}
The total quadratic Hessian is independent of the subtraction scale:
\begin{equation}
 \beta_R(\mu)K+d_1^*\Pi_\Lambda^{\rm ren}(\mu)d_1
 =\beta_R(\nu)K+d_1^*\Pi_\Lambda^{\rm ren}(\nu)d_1.
 \label{eq:v6v55-total-invariance}
\end{equation}
\end{theorem}

\begin{proof}
Subtract the two representations of \(H_\Lambda^{\rm mat}\).  Since the
local part is a scalar multiple of \(K=d_1^*d_1\), the difference of the
renormalized curvature kernels on \(\operatorname{im}d_1\) is the opposite
scalar multiple of the identity, which is
\eqref{eq:v6v55-Pi-cocycle}.  Definition
\eqref{eq:v6v55-betaR} gives \eqref{eq:v6v55-beta-cocycle}; telescoping
gives \eqref{eq:v6v55-cocycle}.  Adding the two transformed pieces proves
\eqref{eq:v6v55-total-invariance}.
\end{proof}

\begin{corollary}[Renormalized coercivity criterion]
\label{cor:v6v55-renormalized-coercivity}
If the Higgs remainder has relative negative curvature
\(\kappa_{\rm H}\) and
\begin{equation}
 \Pi_\Lambda^{\rm ren}(\mu)\succeq-kappa_{\rm f}(\mu)I,
 \qquad
 \beta_R(\mu)>\kappa_{\rm H}+\kappa_{\rm f}(\mu),
 \label{eq:v6v55-renormalized-margin}
\end{equation}
then the positive reference action satisfies the V50 relative Maxwell
curvature card with margin
\(\beta_R(\mu)-\kappa_{\rm H}-\kappa_{\rm f}(\mu)\).  The conclusion is
independent of the subtraction scale whenever the lower bound is transformed
together with \eqref{eq:v6v55-Pi-cocycle}.
\end{corollary}

\begin{proof}
Insert \eqref{eq:v6v55-complete-subtraction} into the total Hessian and add
the Higgs lower bound.  The result is at least the margin in
\eqref{eq:v6v55-renormalized-margin} times \(K\), so
Theorem~\ref{thm:v6v50-RMC} applies.  Equality
\eqref{eq:v6v55-total-invariance} proves scale independence of the actual
form.
\end{proof}

\subsection{Exact running in the V54 massless bubble branch}

For the massless bubble model, use the subtraction coefficient
\begin{equation}
 c_\Lambda(\mu)=Q_2\mathcal B_{\Lambda,0}(p),
 \qquad |p|=\mu>0.
 \label{eq:v6v55-bubble-counterterm}
\end{equation}

\begin{theorem}[Massless bubble running law]
\label{thm:v6v55-bubble-running}
For any \(\mu,\nu>0\), the cutoff limit of the subtraction cocycle exists
and is
\begin{equation}
 \delta c(\nu,\mu)
 :={Q_2\over16\pi^2}\log{\mu^2\over\nu^2}
 =-{Q_2\over8\pi^2}\log{\nu\over\mu}.
 \label{eq:v6v55-running-cocycle}
\end{equation}
Consequently the renormalized Maxwell coefficient in this model obeys
\begin{equation}
 \beta_R(\nu)
 =\beta_R(\mu)-{Q_2\over8\pi^2}\log{\nu\over\mu}.
 \label{eq:v6v55-beta-running}
\end{equation}
\end{theorem}

\begin{proof}
Apply \eqref{eq:v6v54-renormalized-limit} with \(m=0\),
\(|p|=\nu\), and \(|p_*|=\mu\), then multiply by \(Q_2\).  Insert the
result into \eqref{eq:v6v55-beta-cocycle}.
\end{proof}

\begin{corollary}[Finite coercive horizon of the direct bubble trajectory]
\label{cor:v6v55-horizon}
Let \(b_*>0\) be a required lower margin for the local Maxwell coefficient
at the reference scale \(\mu\).  If the running law
\eqref{eq:v6v55-beta-running} is continued without a new ultraviolet
mechanism, then
\begin{equation}
 \beta_R(\nu)>b_*
 \quad\Longleftrightarrow\quad
 \nu<\mu\exp\!\left[
 {8\pi^2\over Q_2}\bigl(\beta_R(\mu)-b_*\bigr)
 \right].
 \label{eq:v6v55-horizon}
\end{equation}
In particular, for finite \(\beta_R(\mu)\), the direct branch cannot keep a
fixed positive coercive margin at arbitrarily large \(\nu\).
\end{corollary}

\begin{proof}
Solve \(\beta_R(\mu)-(Q_2/8\pi^2)\log(\nu/\mu)>b_*\) for \(\nu\).  Since
\(Q_2=10/3>0\), the inequality has the stated finite upper endpoint.
\end{proof}

\begin{remark}[What the horizon does and does not prove]
\label{rem:v6v55-horizon-scope}
Corollary~\ref{cor:v6v55-horizon} is exact for the V54 transverse bubble
model and recovers the sign of the earlier hypercharge obstruction.  It is
not a nonperturbative theorem about the complete Standard Model beta
function.  It proves that the present direct branch cannot be declared an
all-cutoff construction merely because its subtracted two-point kernel is
finite.
\end{remark}

\subsection{Upstream decision}

\begin{proposition}[Exhaustion of the direct noncompact branch]
\label{prop:v6v55-direct-exhaustion}
Within the assumptions proved in V50--V54, exactly one of the following new
inputs is needed to pass the horizon:
\begin{enumerate}[label=\textup{(\arabic*)},leftmargin=3.5em]
\item a nonperturbative ultraviolet trajectory on which the complete
      \(\beta_R-\kappa_{\rm H}-\kappa_{\rm f}\) stays positive;
\item finite-scale matching to an enlarged gauge theory in which
      hypercharge is an embedded generator and the matched kinetic form is
      positive;
\item a different ultraviolet definition whose continuum endpoint is
      proved equivalent to the desired low-energy hypercharge sector.
\end{enumerate}
No result in V50--V55 supplies any of these inputs.
\end{proposition}

\begin{proof}
The direct noncompact construction requires the positive margin in
Corollary~\ref{cor:v6v55-renormalized-coercivity}.  The model running law
shows why subtraction alone does not preserve it.  Maintaining the same
theory, embedding it at a matching scale, or changing its ultraviolet
definition are the three logical ways to replace the failed continuation.
The claim boundaries of V50--V54 leave all three open.
\end{proof}

\begin{firewall}[No counterterm is taken from an isolated loop fragment]
\label{fw:v6v55-complete-loop}
Equations \eqref{eq:v6v55-complete-subtraction}--
\eqref{eq:v6v55-total-invariance} apply only after the complete regulated
real polarization has been assembled.  The V54 bubble may predict and test
the necessary logarithm, but its coefficient may not be inserted into the
full chiral theory without a regulator-level identification.  This is the
direct consequence of the YM V39 audit.
\end{firewall}

\begin{theorem}[Sixth-series V55 result]
\label{thm:v6v55-status}
The compulsory audit shows that the polarization logarithm must be handled
as a complete-loop local counterterm or by signed cancellation, not as a
summable shell error.  Complete-loop subtraction obeys the exact cocycle
\eqref{eq:v6v55-cocycle} and preserves the total Hessian.  In the massless
V54 bubble branch the running coefficient satisfies
\eqref{eq:v6v55-beta-running} and loses every fixed positive margin at the
finite horizon \eqref{eq:v6v55-horizon}.  The direct noncompact construction
therefore stops at a genuine ultraviolet gate and must return to the V14
matching alternatives.
\end{theorem}

\begin{proof}
Use Theorem~\ref{thm:v6v55-audit-decision},
Theorem~\ref{thm:v6v55-subtraction-cocycle},
Corollary~\ref{cor:v6v55-renormalized-coercivity},
Theorem~\ref{thm:v6v55-bubble-running}, and
Corollary~\ref{cor:v6v55-horizon}, subject to
Firewalls~\ref{fw:v6v55-audit-types} and
\ref{fw:v6v55-complete-loop}.
\end{proof}

V56 will go upstream to the finite-scale embedded-hypercharge interface of
V14.  It will extract the exact Lie-algebra kinetic matching matrix for an
embedded \(U(1)_Y\) generator, prove the Schur-complement positivity
criterion for the heavy gauge directions, and identify the threshold data
still missing from a constructive ultraviolet completion.


\clearpage
\section{V56: embedded hypercharge, heavy-direction Schur matching, and an \(SU(5)\) test}
\label{sec:v6v56}

V55 proves that the direct massless bubble trajectory cannot retain a fixed
positive Maxwell margin at arbitrarily high scales.  V14 had already listed
finite-scale embedding as an upstream alternative and proved the diagonal
tree-level matching formula.  This version treats the missing quadratic
issue: kinetic and threshold mixing with heavy gauge directions.  It derives
the exact Schur complement, proves its positivity and stability, and tests
the normalization and light representation content of the minimal
\(SU(5)\) embedding.

\subsection{Quadratic kinetic mixing at the transition}

At one curvature mode on the matching slice, write the hypercharge component
as \(f\in\mathbb R\) and the heavy components as \(G\in\mathbb R^r\).  Let
the complete real quadratic kinetic matrix, after all threshold terms at the
slice have been assembled, be
\begin{equation}
 \mathbb K=
 \begin{pmatrix}
 k&b^T\\ b&C
 \end{pmatrix},
 \qquad C=C^T.
 \label{eq:v6v56-K-block}
\end{equation}
The quadratic action is
\begin{equation}
 S^{(2)}(f,G)
 ={1\over4}\left(kf^2+2f b^TG+G^TCG\right).
 \label{eq:v6v56-mixed-action}
\end{equation}

\begin{theorem}[Exact heavy-mode marginal]
\label{thm:v6v56-Schur-marginal}
Assume \(C\succ0\).  Integrating the heavy variable against Lebesgue measure
gives
\begin{align}
 &\int_{\mathbb R^r}e^{-S^{(2)}(f,G)}\,dG
 \notag\\
 &\qquad=(4\pi)^{r/2}(\det C)^{-1/2}
 \exp\!\left[-{1\over4}k_{\rm eff}f^2\right],
 \label{eq:v6v56-Gaussian-marginal}\\
 &k_{\rm eff}:=k-b^TC^{-1}b.
 \label{eq:v6v56-Schur}
\end{align}
The full matrix \(\mathbb K\) is positive definite if and only if
\begin{equation}
 C\succ0
 \quad\hbox{and}\quad k_{\rm eff}>0.
 \label{eq:v6v56-Schur-positive}
\end{equation}
\end{theorem}

\begin{proof}
Complete the square:
\[
 kf^2+2fb^TG+G^TCG
 =(G+C^{-1}bf)^TC(G+C^{-1}bf)
 +(k-b^TC^{-1}b)f^2.
\]
Translation of \(G\) and the standard Gaussian integral prove
\eqref{eq:v6v56-Gaussian-marginal}.  The same identity shows that positivity
of \(C\) and of its Schur complement is sufficient for positivity of
\(\mathbb K\).  Conversely, restriction to \(f=0\) gives \(C\succ0\), and
evaluation at \(G=-C^{-1}bf\) gives \(k_{\rm eff}>0\).
\end{proof}

\begin{corollary}[Frozen and marginalized matching are different]
\label{cor:v6v56-frozen-versus-marginal}
If the heavy boundary curvature is fixed to \(G=0\), the hypercharge
coefficient is \(k\).  If it is integrated at the quadratic interface, the
coefficient is \(k_{\rm eff}=k-b^TC^{-1}b\le k\), with equality exactly
when \(b=0\).  Therefore the V14 diagonal formula is the marginalized
coefficient only when the complete threshold kinetic form has no
hypercharge--heavy mixing.
\end{corollary}

\begin{proof}
The frozen statement follows by substitution into
\eqref{eq:v6v56-mixed-action}; the marginalized statement is
Theorem~\ref{thm:v6v56-Schur-marginal}.  Since \(C^{-1}\succ0\),
\(b^TC^{-1}b\ge0\), with equality only for \(b=0\).
\end{proof}

\subsection{Basis-free variational form and threshold stability}

For fixed \(f\), the Schur coefficient has the variational representation
\begin{equation}
 k_{\rm eff}f^2
 =\min_{G\in\mathbb R^r}
 \begin{pmatrix}f\\G\end{pmatrix}^{\!T}
 \mathbb K
 \begin{pmatrix}f\\G\end{pmatrix}.
 \label{eq:v6v56-Schur-variational}
\end{equation}

\begin{theorem}[Relative threshold stability]
\label{thm:v6v56-threshold-stability}
Let \(\mathbb K_0\succ0\) and let \(\Delta\mathbb K\) be symmetric with
\begin{equation}
 \left\|\mathbb K_0^{-1/2}
 \Delta\mathbb K\mathbb K_0^{-1/2}\right\|_{\rm op}
 \le\varepsilon<1.
 \label{eq:v6v56-relative-threshold}
\end{equation}
If \(k_{\rm eff}^{(0)}\) and \(k_{\rm eff}\) are the Schur coefficients of
\(\mathbb K_0\) and \(\mathbb K_0+\Delta\mathbb K\), then
\begin{equation}
 (1-\varepsilon)k_{\rm eff}^{(0)}
 \le k_{\rm eff}
 \le(1+\varepsilon)k_{\rm eff}^{(0)}.
 \label{eq:v6v56-Schur-stability}
\end{equation}
In particular the complete threshold form remains positive.
\end{theorem}

\begin{proof}
Hypothesis \eqref{eq:v6v56-relative-threshold} is equivalent to the form
inequalities
\[
 (1-\varepsilon)\mathbb K_0
 \preceq\mathbb K_0+\Delta\mathbb K
 \preceq(1+\varepsilon)\mathbb K_0.
\]
Fix \(f\), minimize all three forms in \(G\), and use
\eqref{eq:v6v56-Schur-variational}.  This proves
\eqref{eq:v6v56-Schur-stability}; its lower bound is positive.
\end{proof}

\begin{corollary}[Absolute perturbation criterion]
\label{cor:v6v56-absolute-threshold}
If \(\mathbb K_0\succeq\kappa_0I\) and
\(\|\Delta\mathbb K\|_{\rm op}<\kappa_0\), then
\eqref{eq:v6v56-Schur-stability} holds with
\(\varepsilon=\|\Delta\mathbb K\|/\kappa_0<1\).
\end{corollary}

\begin{proof}
The norm of
\(\mathbb K_0^{-1/2}\Delta\mathbb K\mathbb K_0^{-1/2}\) is at most
\(\kappa_0^{-1}\|\Delta\mathbb K\|\).
\end{proof}

\begin{firewall}[Threshold matrices must be complete-loop objects]
\label{fw:v6v56-complete-threshold}
The matrix \(\Delta\mathbb K\) in
Theorem~\ref{thm:v6v56-threshold-stability} is the complete real quadratic
threshold form after heavy gauge, Higgs, fermion, ghost, and regulator
contributions appropriate to the chosen construction have been assembled.
The V55 audit forbids applying the norm criterion to isolated terms whose
large jets cancel only in the full object.
\end{firewall}

\subsection{Decoupling estimate for a massive heavy block}

Allow the heavy block to depend on momentum and suppose
\(C(p)\succeq c(p)I\), \(c(p)>0\).  Set
\(k_{\rm eff}(p)=k(p)-b(p)^TC(p)^{-1}b(p)\).

\begin{proposition}[Heavy-direction decoupling]
\label{prop:v6v56-heavy-decoupling}
For every momentum,
\begin{equation}
 0\le k(p)-k_{\rm eff}(p)
 \le{\|b(p)\|^2\over c(p)}.
 \label{eq:v6v56-decoupling}
\end{equation}
If \(c(p)\ge M^2+c_0|p|^2\) and \(\|b(p)\|\le b_0|p|\), then
\begin{equation}
 0\le k(p)-k_{\rm eff}(p)
 \le {b_0^2|p|^2\over M^2+c_0|p|^2}.
 \label{eq:v6v56-mass-decoupling}
\end{equation}
Thus the mixing correction vanishes quadratically at fixed low momentum as
the heavy mass tends to infinity.
\end{proposition}

\begin{proof}
Positivity of \(C^{-1}\) gives the lower inequality, and
\(C^{-1}\preceq c(p)^{-1}I\) gives the upper one.  Insert the two additional
bounds for \eqref{eq:v6v56-mass-decoupling}.
\end{proof}

\subsection{Concrete \(SU(5)\) normalization}

In the fundamental representation of \(SU(5)\), put
\begin{equation}
 Y=\operatorname{diag}\!\left(-{1\over3},-{1\over3},-{1\over3},
 {1\over2},{1\over2}\right).
 \label{eq:v6v56-SU5-Y}
\end{equation}
It is traceless and generates the Standard Model hypercharge subgroup.
Normalize the \(SU(5)\) generators by
\(\operatorname{Tr}_{\mathbf5}(T_aT_b)=\delta_{ab}/2\).

\begin{lemma}[The \(5/3\) kinetic normalization]
\label{lem:v6v56-five-thirds}
One has
\begin{equation}
 \operatorname{Tr}_{\mathbf5}(Y^2)={5\over6},
 \qquad
 T_Y:=\sqrt{3\over5}\,Y,
 \qquad
 \operatorname{Tr}_{\mathbf5}(T_Y^2)={1\over2}.
 \label{eq:v6v56-SU5-normalization}
\end{equation}
Consequently, with no tree-level kinetic mixing,
\begin{equation}
 g_Y=\sqrt{3\over5}\,g_5,
 \qquad
 g_Y^{-2}={5\over3}g_5^{-2}.
 \label{eq:v6v56-SU5-coupling}
\end{equation}
\end{lemma}

\begin{proof}
Equation \eqref{eq:v6v56-SU5-Y} gives
\[
 \operatorname{Tr}(Y^2)=3\cdot{1\over9}+2\cdot{1\over4}={5\over6}.
\]
Multiplication by \(3/5\) makes the squared norm \(1/2\).  In the covariant
derivative, the unified term is
\(g_5T_YA_Y=g_5\sqrt{3/5}\,Y A_Y\).  Comparing it with
\(g_YYB_Y\) in the common canonical field normalization gives
\eqref{eq:v6v56-SU5-coupling}.
\end{proof}

\begin{lemma}[One-generation representation restriction]
\label{lem:v6v56-SU5-generation}
Under \(SU(5)\supset SU(3)\times SU(2)\times U(1)_Y\),
\begin{align}
 \overline{\mathbf5}
 &\longrightarrow
 (\overline{\mathbf3},\mathbf1)_{1/3}
 \oplus(\mathbf1,\mathbf2)_{-1/2},
 \label{eq:v6v56-fivebar}\\
 \mathbf{10}
 &\longrightarrow
 (\overline{\mathbf3},\mathbf1)_{-2/3}
 \oplus(\mathbf3,\mathbf2)_{1/6}
 \oplus(\mathbf1,\mathbf1)_1.
 \label{eq:v6v56-ten}
\end{align}
Thus \(\overline{\mathbf5}\oplus\mathbf{10}\) contains exactly the
all-left-handed multiplets \(d^c,L,u^c,Q_L,e^c\) used in the V38 and V51
anomaly arithmetic.
\end{lemma}

\begin{proof}
On \(\overline{\mathbf5}\), negate the five eigenvalues of
\eqref{eq:v6v56-SU5-Y}: the first three form an
\(\overline{\mathbf3}\) with charge \(1/3\), and the last two form an
\(SU(2)\) doublet with charge \(-1/2\).  The representation
\(\mathbf{10}=\bigwedge^2\mathbf5\) has charges obtained by adding pairs of
fundamental eigenvalues.  Two color indices give charge \(-2/3\) and the
antisymmetric \(\overline{\mathbf3}\); one color and one weak index give
charge \(1/6\) and \((\mathbf3,\mathbf2)\); the two weak indices give a
singlet of charge one.
\end{proof}

\begin{corollary}[Local anomaly arithmetic is inherited correctly]
\label{cor:v6v56-SU5-anomaly}
The restricted light spectrum in
Lemma~\ref{lem:v6v56-SU5-generation} has the vanishing linear, cubic, and
mixed Standard Model anomaly coefficients already proved in V38, and it has
four \(SU(2)\) doublets per generation when color multiplicity is counted.
The local charge test and the parity test therefore pass.
\end{corollary}

\begin{proof}
The spectrum is identical, including multiplicity, to the V38 spectrum, so
its computed anomaly sums are identical.  The three colored copies of
\(Q_L\) and the single lepton doublet total four.
\end{proof}

\begin{firewall}[The \(SU(5)\) test is compatibility, not construction]
\label{fw:v6v56-SU5-scope}
The normalization and representation restriction do not construct a
reflection-positive chiral \(SU(5)\) lattice measure, prove its determinant
line globally trivial, control symmetry breaking and threshold matrices,
remove its cutoff, or establish its continuum gauge dynamics.  In
particular, the present YM companion has not proved the all-order bridge
card even for its current compact model.  The explicit \(SU(5)\) arithmetic
only proves that this embedding passes two necessary interface tests.
\end{firewall}

\begin{definition}[Embedded hypercharge kinetic card]
\label{def:v6v56-EHKC}
EHKC consists of a scale-uniform construction of the complete transition
matrix \(\mathbb K\) such that:
\begin{enumerate}[label=\textup{(EK\arabic*)},leftmargin=4.8em]
\item its heavy block is positive and its hypercharge Schur complement has a
      fixed positive lower bound;
\item its threshold correction obeys a complete-object bound such as
      \eqref{eq:v6v56-relative-threshold};
\item its light representations and determinant line match across the
      interface; and
\item its Schur coefficient agrees with the infrared renormalized
      coefficient in the exact V55 subtraction scheme.
\end{enumerate}
\end{definition}

\begin{theorem}[Sixth-series V56 result]
\label{thm:v6v56-status}
Integrating heavy gauge curvatures at a finite matching interface replaces
the frozen hypercharge coefficient by the Schur complement
\(k-b^TC^{-1}b\).  Positivity is equivalent to positivity of the heavy block
and this complement, and a relative complete-threshold perturbation of size
\(\varepsilon<1\) preserves the complement within factors
\(1\pm\varepsilon\).  The minimal \(SU(5)\) embedding has tree-level
normalization \(g_Y^{-2}=(5/3)g_5^{-2}\) and restricts
\(\overline{\mathbf5}\oplus\mathbf{10}\) to exactly one anomaly-free
Standard Model generation.  EHKC still requires a constructive positive
ultraviolet measure and its complete threshold matrix.
\end{theorem}

\begin{proof}
Use Theorems~\ref{thm:v6v56-Schur-marginal} and
\ref{thm:v6v56-threshold-stability},
Proposition~\ref{prop:v6v56-heavy-decoupling},
Lemmas~\ref{lem:v6v56-five-thirds} and
\ref{lem:v6v56-SU5-generation}, and
Corollary~\ref{cor:v6v56-SU5-anomaly}, subject to both firewalls.
\end{proof}

V57 will isolate the next EHKC gate: positivity of a complete heavy-field
Gaussian integration when the heavy block is an operator rather than a
finite matrix.  It will prove the operator Schur-complement and determinant
normalization theorem, including zero-mode projections and a quantitative
trace-class threshold criterion needed for interface amplitudes.


\clearpage
\section{V57: operator Schur complements, zero modes, and relative Gaussian normalization}
\label{sec:v6v57}

V56 treats a finite kinetic matrix at the embedded-hypercharge transition.
Actual heavy gauge variables form a field space, and their Gaussian
normalization can itself fail in a cutoff limit.  This version proves the
bounded-operator Schur theorem, gives the exact relative Fredholm determinant
for the heavy Gaussian marginal, and identifies the zero-mode alternatives.

\subsection{Bounded operator Schur theorem}

Let \(\mathcal H_Y\) and \(\mathcal H_H\) be real Hilbert spaces.  Let
\begin{equation}
 \mathbb K=
 \begin{pmatrix}
 A&B^*\\ B&C
 \end{pmatrix}
 \label{eq:v6v57-operator-block}
\end{equation}
be bounded and self-adjoint on
\(\mathcal H_Y\oplus\mathcal H_H\), with
\(C\succeq c_HI\) for some \(c_H>0\).  Define
\begin{equation}
 S_Y:=A-B^*C^{-1}B.
 \label{eq:v6v57-operator-Schur}
\end{equation}

\begin{theorem}[Operator completion of the square]
\label{thm:v6v57-operator-Schur}
For every \(f\in\mathcal H_Y\) and \(G\in\mathcal H_H\),
\begin{align}
 \left\langle
 \binom fG,\mathbb K\binom fG
 \right\rangle
 &=\langle f,S_Yf\rangle
 +\left\langle G+C^{-1}Bf,
 C(G+C^{-1}Bf)\right\rangle,
 \label{eq:v6v57-square-completion}\\
 \langle f,S_Yf\rangle
 &=\inf_{G\in\mathcal H_H}
 \left\langle
 \binom fG,\mathbb K\binom fG
 \right\rangle.
 \label{eq:v6v57-operator-variational}
\end{align}
Consequently \(\mathbb K\succ0\) if and only if
\(C\succ0\) and \(S_Y\succ0\), with uniform lower bounds understood in
the operator sense.
\end{theorem}

\begin{proof}
Bounded invertibility of \(C\) makes \(C^{-1}B\) bounded.  Expand the second
term on the right of \eqref{eq:v6v57-square-completion}; its two cross terms
are \(2\langle Bf,G\rangle\), and the added
\(\langle Bf,C^{-1}Bf\rangle\) cancels the subtraction in \(S_Y\).
The second term is minimized uniquely at \(G=-C^{-1}Bf\), proving
\eqref{eq:v6v57-operator-variational}.  Positivity follows exactly as in
the finite-dimensional Schur criterion.
\end{proof}

\begin{theorem}[Operator relative-threshold stability]
\label{thm:v6v57-operator-stability}
Let \(\mathbb K_0\succ0\) be another bounded block operator with invertible
heavy block and suppose, as quadratic forms,
\begin{equation}
 (1-\varepsilon)\mathbb K_0
 \preceq\mathbb K
 \preceq(1+\varepsilon)\mathbb K_0,
 \qquad0\le\varepsilon<1.
 \label{eq:v6v57-operator-relative}
\end{equation}
Then their hypercharge Schur operators obey
\begin{equation}
 (1-\varepsilon)S_{Y,0}
 \preceq S_Y
 \preceq(1+\varepsilon)S_{Y,0}.
 \label{eq:v6v57-operator-Schur-stability}
\end{equation}
\end{theorem}

\begin{proof}
For fixed \(f\), take the infimum over \(G\) in all three forms in
\eqref{eq:v6v57-operator-relative} and apply
\eqref{eq:v6v57-operator-variational} to each.  Since the result holds for
every \(f\), it is the operator inequality
\eqref{eq:v6v57-operator-Schur-stability}.
\end{proof}

\subsection{Relative heavy Gaussian integral}

Let \(C_0\succ0\) be self-adjoint and suppose \(C_0^{-1}\) is trace class,
so the centered Gaussian law
\begin{equation}
 \gamma_0=N(0,C_0^{-1})
 \label{eq:v6v57-gamma0}
\end{equation}
is a countably additive probability measure on \(\mathcal H_H\).  Let
\(T=T^*\) be trace class with \(\|T\|_{\rm op}<1\), and define
\begin{equation}
 C:=C_0^{1/2}(I+T)C_0^{1/2}.
 \label{eq:v6v57-relative-C}
\end{equation}
For \(j\in\mathcal H_H\), put
\begin{equation}
 \mathcal I_T(j):=
 \int\exp\!\left[
 -{1\over2}\langle G,C_0^{1/2}TC_0^{1/2}G\rangle
 -\langle j,G\rangle
 \right]d\gamma_0(G),
 \label{eq:v6v57-relative-integral}
\end{equation}
first for \(j\) in the Cameron--Martin dual so every pairing is defined.

\begin{theorem}[Exact relative determinant and source shift]
\label{thm:v6v57-relative-Gaussian}
Under the preceding hypotheses,
\begin{equation}
 \mathcal I_T(j)
 =\det(I+T)^{-1/2}
 \exp\!\left[{1\over2}\langle j,C^{-1}j\rangle\right].
 \label{eq:v6v57-relative-Gaussian}
\end{equation}
Here \(\det(I+T)\) is the Fredholm determinant.  In particular, for a
bounded mixing operator \(B:\mathcal H_Y\to\mathcal H_H\), integration of
the cross term \(\langle Bf,G\rangle\) produces the Schur action
\begin{equation}
 {1\over2}\langle f,(A-B^*C^{-1}B)f\rangle
 \label{eq:v6v57-effective-action}
\end{equation}
and the normalization \(\det(I+T)^{-1/2}\).
\end{theorem}

\begin{proof}
Choose an increasing sequence of finite-rank projections which diagonalizes
the compact self-adjoint \(T\) and exhausts \(\mathcal H_H\).  In each
finite-dimensional range, change variables by \(C_0^{1/2}\), complete the
square, and apply the ordinary Gaussian formula.  The result is
\[
 \prod_{n\le N}(1+t_n)^{-1/2}
 \exp\!\left[{1\over2}
 \langle P_Nj,C_N^{-1}P_Nj\rangle\right],
\]
where \(t_n\) are the eigenvalues of \(T\).  Trace-class summability and
\(|t_n|<1\) make \(\sum_n\log(1+t_n)\) absolutely convergent.  The source
quadratic forms converge monotonically after resolving positive and
negative spectral parts, or directly by strong convergence of the bounded
inverse \((I+T)^{-1}\).  Finite-dimensional Gaussian expectations converge
to \eqref{eq:v6v57-relative-integral} by uniform integrability, supplied by
the strict bound \(I+T\succeq(1-\|T\|)I\).  This proves
\eqref{eq:v6v57-relative-Gaussian}.  Taking \(j=Bf\) and combining its
positive shift with the original \(-\langle f,Af\rangle/2\) yields
\eqref{eq:v6v57-effective-action}.
\end{proof}

\begin{lemma}[Quantitative determinant control]
\label{lem:v6v57-det-control}
If \(T=T^*\) is trace class and \(\|T\|_{\rm op}\le\eta<1\), then
\begin{equation}
 |\log\det(I+T)|
 \le{\|T\|_1\over1-\eta}.
 \label{eq:v6v57-logdet-bound}
\end{equation}
If \(T_n\to T\) in trace norm and
\(\sup_n\|T_n\|_{\rm op}\le\eta<1\), then
\begin{equation}
 \log\det(I+T_n)\longrightarrow\log\det(I+T).
 \label{eq:v6v57-det-convergence}
\end{equation}
\end{lemma}

\begin{proof}
For every real \(|t|\le\eta\), the mean-value theorem gives
\(|\log(1+t)|\le|t|/(1-\eta)\).  Sum over the eigenvalues of \(T\) to prove
\eqref{eq:v6v57-logdet-bound}.  For convergence, use the absolutely
convergent series
\[
 \log\det(I+T)=\sum_{k=1}^\infty{(-1)^{k+1}\over k}\operatorname{Tr}(T^k).
\]
Each trace is trace-norm continuous, and the tail is uniformly bounded by a
constant times \(\eta^k\sup_n\|T_n\|_1\).  Trace-norm convergence bounds the
last supremum, so dominated convergence proves
\eqref{eq:v6v57-det-convergence}.
\end{proof}

\begin{corollary}[Interface-amplitude control]
\label{cor:v6v57-interface-amplitude}
Suppose the heavy relative perturbations \(T_n(q)\) on a common interface
obey \(\|T_n(q)\|\le\eta<1\) and
\(\|T_n(q)\|_1\le M(q)\), where the exponential of
\(M(q)/(2(1-\eta))\) is integrable against the remaining interface
amplitudes.  If \(T_n(q)\) converges in trace norm almost everywhere and the
Schur source quadratics converge, then the normalized heavy Gaussian
factors converge in \(L^1\) and may be sewn by the V14 finite-cutoff sewing
theorem.
\end{corollary}

\begin{proof}
Theorem~\ref{thm:v6v57-relative-Gaussian} gives pointwise convergence by
Lemma~\ref{lem:v6v57-det-control} and the assumed source convergence.  The
same lemma supplies the stated integrable determinant majorant.  The Schur
quadratic is nonpositive in the density exponent after the full positive
action is assembled, so its normalized factor supplies no additional
growth; otherwise its assumed convergent majorant must be included.
Dominated convergence gives \(L^1\) convergence.  Apply
Theorem~\ref{thm:v6v14-sewing}.
\end{proof}

\subsection{Zero modes are retained interface variables}

Let now \(C\succeq0\) have kernel \(Z=\ker C\), and write
\(Q=I-P_Z\).  Suppose \(C_Q:=QCQ|_{Q\mathcal H_H}\succeq c_HQ\).

\begin{proposition}[Exact zero-mode alternatives]
\label{prop:v6v57-zero-modes}
For the quadratic action
\begin{equation}
 {1\over2}\langle G,CG\rangle+\langle Bf,G\rangle,
 \label{eq:v6v57-zero-action}
\end{equation}
the following statements hold.
\begin{enumerate}[label=\textup{(\roman*)},leftmargin=3.7em]
\item Lebesgue integration over \(Z\) is never normalizable when
      \(P_ZBf=0\), because the action is constant along \(Z\).
\item If \(P_ZBf\ne0\), the action is affine and unbounded in a zero-mode
      direction, so the real Gaussian integral again does not exist.
\item After fixing the zero mode or retaining it as a separate boundary
      variable, integration over \(Q\mathcal H_H\) is governed by the
      reduced Schur operator
      \begin{equation}
       S_{Y,Q}=A-B^*QC_Q^{-1}QB.
       \label{eq:v6v57-reduced-Schur}
      \end{equation}
\end{enumerate}
\end{proposition}

\begin{proof}
Write \(G=z+G_Q\).  On \(Z\), the quadratic term vanishes.  If the linear
coefficient vanishes, integration of the constant density over a nontrivial
Euclidean zero-mode space has infinite mass.  If it does not vanish, choose
the sign of \(z\) along that coefficient to send
\eqref{eq:v6v57-zero-action} to \(-\infty\).  Removing or conditioning on
\(z\), complete the square on \(Q\mathcal H_H\) to obtain
\eqref{eq:v6v57-reduced-Schur}.
\end{proof}

\begin{assessment}[Agreement with the fermion zero-mode rule]
\label{ass:v6v57-zero-mode-agreement}
The heavy bosonic conclusion is the exact analogue of V39: a zero mode is
not repaired by deleting a vanishing determinant factor.  It becomes an
explicit boundary or sector variable, and the determinant and Schur
complement are formed only on the invertible complement.
\end{assessment}

\subsection{Why ordinary determinants fail at a field-theory limit}

\begin{proposition}[Trace-class necessity for the displayed relative
normalization]
\label{prop:v6v57-trace-necessity}
Let \(T\ge0\) be compact with eigenvalues \(t_n\ge0\).  The finite-rank
relative Gaussian normalizations
\begin{equation}
 Z_N=\prod_{n=1}^N(1+t_n)^{-1/2}
 \label{eq:v6v57-ZN}
\end{equation}
converge to a strictly positive limit if and only if
\(\sum_nt_n<\infty\), equivalently \(T\) is trace class.
\end{proposition}

\begin{proof}
For \(0\le t\le\|T\|\), the functions \(\log(1+t)\) and \(t\) are
comparable by positive constants depending only on \(\|T\|\).  Hence
\(-2\log Z_N=\sum_{n\le N}\log(1+t_n)\) has a finite limit exactly when
\(\sum_nt_n<\infty\).  The limiting product is positive exactly in that
case.
\end{proof}

\begin{firewall}[Four-dimensional field covariances need a relative or
renormalized setting]
\label{fw:v6v57-field-determinant}
For a standard second-order field operator in four dimensions, its inverse
on \(L^2\) is generally not trace class, and local threshold perturbations
need not make \(T\) trace class before counterterms are removed.  Therefore
Theorem~\ref{thm:v6v57-relative-Gaussian} is a complete theorem under its
stated Hilbert-space hypotheses, not a claim that the physical heavy gauge
field already satisfies them.  A constructive interface must choose a
distribution space and prove trace-class relative covariance, or define and
control an appropriate renormalized determinant with its local counterterms.
\end{firewall}

\begin{definition}[Operator embedded-hypercharge card]
\label{def:v6v57-OEHC}
OEHC consists of:
\begin{enumerate}[label=\textup{(OH\arabic*)},leftmargin=4.8em]
\item a positive heavy operator on the complement of explicitly retained
      zero modes;
\item a uniformly positive hypercharge Schur operator;
\item a common Gaussian reference for which the relative heavy perturbation
      has uniformly controlled trace norm after the complete counterterm
      subtraction; and
\item convergence of the relative determinant and Schur source factor in
      the V14 interface \(L^1\) norm.
\end{enumerate}
\end{definition}

\begin{theorem}[Sixth-series V57 result]
\label{thm:v6v57-status}
The heavy-field Schur complement and its relative stability extend exactly
to bounded Hilbert-space operators.  When the reference covariance is trace
class and the complete relative perturbation is trace class with operator
norm below one, the heavy Gaussian marginal has the exact Fredholm
normalization \(\det(I+T)^{-1/2}\), and trace-norm convergence supplies the
interface \(L^1\) limit.  Heavy zero modes cannot be integrated as ordinary
Gaussian variables and must be fixed or retained as boundary sectors.
OEHC records the still-unproved field-theory hypotheses.
\end{theorem}

\begin{proof}
Use Theorems~\ref{thm:v6v57-operator-Schur},
\ref{thm:v6v57-operator-stability}, and
\ref{thm:v6v57-relative-Gaussian},
Lemma~\ref{lem:v6v57-det-control},
Corollary~\ref{cor:v6v57-interface-amplitude}, and
Propositions~\ref{prop:v6v57-zero-modes} and
\ref{prop:v6v57-trace-necessity}, subject to
Firewall~\ref{fw:v6v57-field-determinant}.
\end{proof}

V58 will return to the common SM--Einstein interface.  It will combine the
gauge Schur complement with the gravitational boundary momentum form,
derive the full block positivity criterion including mixed gauge--metric
threshold terms, and test whether the Euclidean conformal direction defeats
that criterion before any probabilistic sewing is attempted.


\clearpage
\section{V58: the coupled gauge--metric quadratic form cannot cure the conformal mode}
\label{sec:v6v58}

V20 proves that the real Euclidean Einstein action is unbounded below along
bounded-amplitude high-frequency conformal packets.  V21 gives a generic
constraint and transfer Schur criterion, while V56--V57 establish the
embedded-hypercharge Schur complement.  The missing coupled statement is
whether a positive gauge sector or its threshold mixing could repair the
metric sign.  This version proves that it cannot: Gaussian elimination of a
positive gauge block only subtracts from the metric Hessian.  At the
zero-gauge background the classical mixed Hessian vanishes altogether.

\subsection{A no-rescue theorem for a negative metric direction}

Let \(\mathcal H_G\) be the already reduced gauge tangent space and
\(\mathcal H_M\) a metric tangent space on the matching slice.  Consider the
self-adjoint quadratic block
\begin{equation}
 \mathbb Q=
 \begin{pmatrix}
 K_G&J\\ J^*&K_M
 \end{pmatrix},
 \qquad K_G\succeq k_GI>0.
 \label{eq:v6v58-coupled-block}
\end{equation}

\begin{theorem}[Positive gauge integration cannot stabilize a negative
metric direction]
\label{thm:v6v58-no-rescue}
The metric Schur operator after integrating the gauge variable is
\begin{equation}
 K_M^{\rm eff}=K_M-J^*K_G^{-1}J\preceq K_M.
 \label{eq:v6v58-metric-Schur}
\end{equation}
If there is \(h\ne0\) with \(\langle h,K_Mh\rangle<0\), then
\(\mathbb Q\) is not positive and
\(\langle h,K_M^{\rm eff}h\rangle<0\).  No size or sign of \(J\) can
reverse this conclusion while \(K_G\) is integrated on a real positive
Gaussian contour.
\end{theorem}

\begin{proof}
The operator completion of the square from
Theorem~\ref{thm:v6v57-operator-Schur}, with gauge field in the eliminated
slot, gives \eqref{eq:v6v58-metric-Schur}.  Since
\(J^*K_G^{-1}J\succeq0\), the form order follows.  Independently, restriction
of \(\mathbb Q\) to vectors \((0,h)\) gives
\(\langle h,K_Mh\rangle<0\), so the full block is not positive.  The same
vector in \eqref{eq:v6v58-metric-Schur} proves strict negativity of the
effective metric form.
\end{proof}

\begin{corollary}[Order of construction at the SM--Einstein interface]
\label{cor:v6v58-order}
The gravitational constraint reduction or ultraviolet modification must
produce a positive physical metric block before the positive gauge Schur
and determinant cards can be used in a common real probability measure.
Gauge threshold mixing cannot be invoked as the mechanism which removes the
Einstein conformal instability.
\end{corollary}

\begin{proof}
If the negative conformal direction remains in the physical metric tangent
space, Theorem~\ref{thm:v6v58-no-rescue} forbids positivity of the coupled
block.  Removing it by exact constraints or changing its metric action is
therefore logically prior to the positive coupled Gaussian construction.
\end{proof}

\subsection{Classical Einstein--Maxwell Hessian at zero gauge field}

For a Riemannian four-metric \(g\) and an Abelian potential \(A\), set
\begin{equation}
 S_M(g,A)={1\over4g_Y^2}
 \int_M\sqrt g\,g^{\mu\rho}g^{\nu\sigma}
 F_{\mu\nu}(A)F_{\rho\sigma}(A),
 \qquad F(A)=dA.
 \label{eq:v6v58-Maxwell-action}
\end{equation}

\begin{lemma}[Vanishing mixed Hessian at \(A=0\)]
\label{lem:v6v58-mixed-zero}
For every background metric \(\bar g\), metric variation \(h\), and gauge
variation \(a\),
\begin{equation}
 d_gd_AS_M(\bar g,0)[h,a]=0.
 \label{eq:v6v58-mixed-Hessian-zero}
\end{equation}
The pure gauge Hessian is the positive Maxwell form
\begin{equation}
 d_A^2S_M(\bar g,0)[a,a]
 ={1\over2g_Y^2}\int_M|da|_{\bar g}^2,d\operatorname{vol}_{\bar g}.
 \label{eq:v6v58-gauge-Hessian}
\end{equation}
\end{lemma}

\begin{proof}
The map \(A\mapsto F(A)\) is linear and
\eqref{eq:v6v58-Maxwell-action} is quadratic in \(F\).  Its first
\(A\)-derivative is therefore linear in \(F(A)\) and vanishes at \(A=0\)
for every metric.  Differentiating that zero identity in the metric gives
\eqref{eq:v6v58-mixed-Hessian-zero}.  Differentiating twice in \(A\) gives
\eqref{eq:v6v58-gauge-Hessian}, with the displayed factor following from
the convention in \eqref{eq:v6v58-Maxwell-action}.
\end{proof}

\begin{proposition}[Block-diagonal classical no-go]
\label{prop:v6v58-classical-no-go}
At a background \((\bar g,A=0)\), the quadratic Einstein--Maxwell action is
block diagonal between metric and gauge variations.  Every negative
conformal direction of the Einstein Hessian from V20 remains a negative
direction of the full classical Hessian.
\end{proposition}

\begin{proof}
Lemma~\ref{lem:v6v58-mixed-zero} removes the off-diagonal block, and
\eqref{eq:v6v58-gauge-Hessian} is nonnegative after gauge and flat modes are
fixed.  Restriction to the metric conformal variation is therefore exactly
the Einstein Hessian already proved negative in V20.
\end{proof}

\subsection{Exact four-dimensional conformal invariance of Maxwell theory}

\begin{theorem}[The Maxwell kinetic action does not see the conformal
factor]
\label{thm:v6v58-Maxwell-conformal}
For \(g=e^{2\phi}\bar g\) in four dimensions,
\begin{equation}
 S_M(e^{2\phi}\bar g,A)=S_M(\bar g,A)
 \label{eq:v6v58-Maxwell-conformal}
\end{equation}
for every smooth real or complex \(\phi\) for which the expression is
defined.  Hence every derivative involving at least one conformal variation
and otherwise only gauge variations vanishes for the pure Maxwell kinetic
term.
\end{theorem}

\begin{proof}
Under the conformal change,
\(\sqrt g=e^{4\phi}\sqrt{\bar g}\) and each inverse metric contributes
\(e^{-2\phi}\).  The product in
\eqref{eq:v6v58-Maxwell-action} therefore has factor
\(e^{4\phi}e^{-2\phi}e^{-2\phi}=1\).  The curvature two-form is metric
independent.  This proves the identity and all derivative statements.
\end{proof}

\begin{corollary}[Conformal rotation leaves the Maxwell block real]
\label{cor:v6v58-rotation-Maxwell}
A formal contour rotation \(\phi=i\chi\) changes neither the pure
hypercharge Maxwell action nor its quadratic coefficient.  It therefore
does not by itself create a gauge-sector phase, but it also supplies no
positive term which could control the conformal variable.
\end{corollary}

\begin{proof}
Apply Theorem~\ref{thm:v6v58-Maxwell-conformal} with complex \(\phi\).
\end{proof}

\begin{firewall}[Masses and anomalies are outside conformal Maxwell
invariance]
\label{fw:v6v58-conformal-scope}
Higgs potentials, Yukawa masses, regulator masses, the metric dependence of
fermion determinants, and Weyl-anomaly counterterms are not conformally
invariant in the form used above.  Their complete mixed Hessian belongs to
the threshold block \(J\).  Theorem~\ref{thm:v6v58-no-rescue} still shows
that such mixing cannot repair an already negative metric principal block
inside a real positive Gaussian construction; a genuinely positive metric
term must enter the metric block itself.
\end{firewall}

\subsection{A higher-derivative scalar sign test}

As a necessary quadratic test for a modified ultraviolet metric action,
consider a conformal Fourier mode whose real quadratic symbol is
\begin{equation}
 p(s)=\alpha s^2-cs+m_0^2,
 \qquad s=|k|^2\ge\lambda_1>0,
 \qquad c>0.
 \label{eq:v6v58-conformal-symbol}
\end{equation}
The term \(-cs\) is the Einstein conformal sign; \(\alpha s^2\) models a
positive fourth-order ultraviolet term.

\begin{theorem}[Exact positivity criterion for the stabilized scalar
symbol]
\label{thm:v6v58-polynomial-criterion}
The form \(p(s)\) is strictly positive for every \(s\ge\lambda_1\) if and
only if \(\alpha>0\) and
\begin{equation}
 \begin{cases}
 m_0^2>{c^2\over4\alpha},
 &\text{if }{c\over2\alpha}\ge\lambda_1,\\
 m_0^2>c\lambda_1-\alpha\lambda_1^2,
 &\text{if }{c\over2\alpha}<\lambda_1,
 \end{cases}
 \label{eq:v6v58-polynomial-positive}
\end{equation}
apart from the special possibility that a finite ultraviolet cutoff truncates
the interval before the downward direction when \(\alpha\le0\).  No
cutoff-uniform limit with arbitrarily large \(s\) is positive when
\(\alpha<0\).
\end{theorem}

\begin{proof}
For \(\alpha>0\), the unconstrained minimum of the parabola occurs at
\(s_*=c/(2\alpha)\).  If \(s_*\ge\lambda_1\), its value is
\(m_0^2-c^2/(4\alpha)\).  If \(s_*<\lambda_1\), the parabola is increasing
on the allowed interval and its minimum is
\(p(\lambda_1)=\alpha\lambda_1^2-c\lambda_1+m_0^2\).  These give
\eqref{eq:v6v58-polynomial-positive}.  If \(\alpha<0\), then
\(p(s)\to-\infty\) as \(s\to\infty\).  The case \(\alpha=0\) also tends
to \(-\infty\) because \(c>0\).
\end{proof}

\begin{assessment}[Ultraviolet damping is not enough]
\label{ass:v6v58-UV-damping}
A positive \(\alpha\) controls high frequencies, but the exact minimum may
remain negative at \(s=c/(2\alpha)\).  Thus the phrase ``add an
\(R^2\)-term'' is not a positivity proof.  One must compute the complete
gauge-fixed and constrained scalar symbol, including its lower-order terms,
and verify \eqref{eq:v6v58-polynomial-positive} or its operator analogue.
\end{assessment}

\begin{firewall}[Euclidean polynomial positivity is not reflection
positivity]
\label{fw:v6v58-higher-derivative-OS}
Even when \eqref{eq:v6v58-polynomial-positive} holds, a fourth-order
Euclidean propagator can decompose into poles with incompatible residues.
The scalar sign test gives integrability of the quadratic density; it does
not prove a positive transfer operator, Osterwalder--Schrader positivity, or
a unitary Lorentzian theory.  Those remain separate requirements of the V20
metric-entry criterion.
\end{firewall}

\begin{definition}[Coupled gauge--metric entry card]
\label{def:v6v58-CGMEC}
CGMEC requires, on the exact constraint and gauge quotient:
\begin{enumerate}[label=\textup{(CG\arabic*)},leftmargin=4.8em]
\item a positive physical metric principal block before gauge variables are
      integrated;
\item a positive embedded-gauge block satisfying OEHC;
\item a complete mixed threshold operator whose relative norm preserves the
      two principal margins; and
\item a positive self-adjoint one-step transfer operator for the assembled
      reduced kernel.
\end{enumerate}
\end{definition}

\begin{theorem}[Sixth-series V58 result]
\label{thm:v6v58-status}
A positive gauge sector cannot repair a negative metric principal block:
its elimination changes the metric Hessian to
\(K_M-J^*K_G^{-1}J\preceq K_M\).  At zero gauge background the classical
Einstein--Maxwell mixed Hessian vanishes, and in four dimensions the Maxwell
kinetic action is exactly independent of the conformal factor.  A positive
fourth-order conformal term passes the real quadratic sign test only under
the explicit inequalities \eqref{eq:v6v58-polynomial-positive}, which do not
imply reflection positivity.  Hence a positive constrained metric block is
an independent prerequisite for the full SM--Einstein interface.
\end{theorem}

\begin{proof}
Use Theorem~\ref{thm:v6v58-no-rescue},
Lemma~\ref{lem:v6v58-mixed-zero},
Theorem~\ref{thm:v6v58-Maxwell-conformal}, and
Theorem~\ref{thm:v6v58-polynomial-criterion}, subject to both firewalls.
\end{proof}

V59 will construct the maximal positive quadratic gravity baseline which is
available without resolving the conformal problem: the transverse-traceless
linearized graviton sector on a fixed torus.  It will prove its finite-cutoff
Gaussian law, reflection-positive oscillator transfer, and distributional
continuum limit, while retaining the scalar and constraint sectors outside
the claim.


\clearpage
\section{V59: the transverse--traceless graviton Gaussian and its positive transfer}
\label{sec:v6v59}

V58 proves that the hypercharge and non-Abelian gauge sectors cannot repair
the real Euclidean conformal sign.  The maximal standard positive quadratic
gravity baseline is therefore the physical transverse--traceless sector
after linearized constraints and gauge directions have been removed.  This
version constructs that sector on a flat torus, proves its one-step transfer
operator positive, and takes its distributional ultraviolet limit.  The
construction deliberately retains homogeneous metric moduli and all
nonlinear constraints outside the Gaussian law.

\subsection{Spatial transverse--traceless modes}

Let the spatial slice be \(\mathbb T_L^3\).  For
\(k\in(2\pi/L)\mathbb Z^3\setminus\{0\}\), define
\begin{equation}
 P_{ij}(k):=\delta_{ij}-{k_ik_j\over|k|^2}
 \label{eq:v6v59-transverse-projector}
\end{equation}
and the orthogonal projector on symmetric transverse traceless tensors
\begin{equation}
 \Pi_{ij,\ell m}^{\rm TT}(k)
 :={1\over2}\left(P_{i\ell}P_{jm}+P_{im}P_{j\ell}
 -P_{ij}P_{\ell m}\right).
 \label{eq:v6v59-TT-projector}
\end{equation}

\begin{lemma}[The TT projector has rank two]
\label{lem:v6v59-TT-rank}
For every \(k\ne0\), \(\Pi^{\rm TT}(k)\) is self-adjoint and idempotent.
Its range consists exactly of symmetric tensors \(h_{ij}\) satisfying
\begin{equation}
 k_ih_{ij}=0,
 \qquad h_{ii}=0,
 \label{eq:v6v59-TT-conditions}
\end{equation}
and has real dimension two.
\end{lemma}

\begin{proof}
Choose an orthonormal frame with third axis parallel to \(k\).  Then
\(P=\operatorname{diag}(1,1,0)\).  Formula
\eqref{eq:v6v59-TT-projector} deletes every component with an index in the
third direction and subtracts the trace in the remaining two-dimensional
plane.  The surviving matrices have the form
\[
 \begin{pmatrix}a&b&0\\b&-a&0\\0&0&0\end{pmatrix},
\]
a two-dimensional real space.  In this frame the formula is visibly the
orthogonal projection onto that space, proving self-adjointness and
idempotence; the statement is frame invariant.
\end{proof}

Choose real orthonormal polarization tensors
\(e^{(\lambda)}(k)\), \(\lambda=1,2\), spanning the range, with the reality
convention \(e^{(\lambda)}(-k)=\overline{e^{(\lambda)}(k)}\).  Fix every
spatial zero mode.  At cutoff \(\Lambda\), write
\begin{equation}
 h_{ij}^{\rm TT}(\tau,x)
 =\sum_{0<|k|\le\Lambda}\sum_{\lambda=1}^2
 q_{k,\lambda}(\tau)e_{ij}^{(\lambda)}(k)e^{ik\cdot x},
 \label{eq:v6v59-mode-expansion}
\end{equation}
with the corresponding reality condition on \(q\).

\subsection{Positive quadratic action and oscillator transfer}

Let \(\kappa_g>0\) be the TT kinetic normalization.  The reduced quadratic
Euclidean action is
\begin{equation}
 S_{\Lambda}^{\rm TT}
 ={\kappa_g\over2}\sum_{k,\lambda}
 \int\left(|\dot q_{k,\lambda}(\tau)|^2
 +|k|^2|q_{k,\lambda}(\tau)|^2\right)d\tau.
 \label{eq:v6v59-TT-action}
\end{equation}
Independent real coordinates are understood in the sum.

For one real mode of frequency \(\omega>0\) and Euclidean time step
\(a>0\), define the Mehler kernel
\begin{equation}
 K_{a,\omega}(q,q')
 :=\left({\kappa_g\omega\over2\pi\sinh(a\omega)}\right)^{1/2}
 \exp\!\left[-{\kappa_g\omega\over2\sinh(a\omega)}
 \left((q^2+{q'}^2)\cosh(a\omega)-2qq'\right)\right].
 \label{eq:v6v59-Mehler}
\end{equation}

\begin{theorem}[Positive TT one-step transfer]
\label{thm:v6v59-positive-transfer}
The integral operator with kernel \(K_{a,\omega}\) is the positive
self-adjoint trace-class operator
\begin{equation}
 e^{-aH_\omega},
 \qquad
 H_\omega={1\over2}\left(-\kappa_g^{-1}{d^2\over dq^2}
 +\kappa_g\omega^2q^2\right)
 \label{eq:v6v59-oscillator-H}
\end{equation}
on \(L^2(\mathbb R,dq)\).  The finite tensor product over all TT modes with
\(0<|k|\le\Lambda\) is therefore a positive self-adjoint transfer operator.
It satisfies the V20 Osterwalder--Schrader quadratic inequality for every
finite-cutoff TT cylinder functional.
\end{theorem}

\begin{proof}
The normalized Hermite functions diagonalize \(H_\omega\) with eigenvalues
\((n+1/2)\omega\).  Mehler's summation formula gives
\[
 K_{a,\omega}(q,q')
 =\sum_{n=0}^\infty e^{-a(n+1/2)\omega}
 \psi_n(q)\psi_n(q').
\]
All eigenvalues are positive and summable, proving the first statement.
Finite tensor products preserve positivity, self-adjointness, and trace
class.  Theorem~\ref{thm:v6v20-transfer-positivity} then gives the
Osterwalder--Schrader inequality.
\end{proof}

\begin{corollary}[Nonnegative finite-cutoff Hamiltonian after vacuum shift]
\label{cor:v6v59-Hamiltonian}
Subtracting the finite zero-point constant
\(\frac12\sum_{k,\lambda}|k|\) produces
\begin{equation}
 H_\Lambda^{\rm TT}
 =\sum_{0<|k|\le\Lambda}\sum_{\lambda=1}^2
 |k|\,a_{k,\lambda}^*a_{k,\lambda}\ge0.
 \label{eq:v6v59-Fock-H}
\end{equation}
The transfer semigroup remains positive; it is multiplied only by a
positive scalar under the vacuum-energy shift.
\end{corollary}

\begin{proof}
The oscillator spectrum in the proof of
Theorem~\ref{thm:v6v59-positive-transfer} is
\((n+1/2)|k|\).  Subtract the sum of the \(1/2\)-terms.  Multiplication of
\(e^{-aH}\) by the resulting positive scalar does not change operator
positivity.
\end{proof}

\subsection{Euclidean covariance and continuum cutoff removal}

Put Euclidean time on a circle of length \(T\), with frequencies
\(p_0\in(2\pi/T)\mathbb Z\).  For \(k\ne0\), define the Gaussian TT
covariance
\begin{equation}
 \mathbb E\left[
 \widehat h_{ij}(p_0,k)
 \overline{\widehat h_{\ell m}(p_0',k')}
 \right]
 ={1\over\kappa_g}
 {\Pi_{ij,\ell m}^{\rm TT}(k)
 \delta_{p_0p_0'}\delta_{kk'}
 \over p_0^2+|k|^2},
 \label{eq:v6v59-TT-covariance}
\end{equation}
with all spatially homogeneous modes fixed.

\begin{theorem}[Distributional TT continuum field]
\label{thm:v6v59-TT-continuum}
Let \(h_\Lambda^{\rm TT}\) retain modes with
\(p_0^2+|k|^2\le\Lambda^2\) and \(k\ne0\).  For every \(s>1\), these
fields converge in
\begin{equation}
 L^2\!\left(\Omega;
 H^{-s}(S_T^1\times\mathbb T_L^3;operatorname{Sym}^2\mathbb R^3)
 \right)
 \label{eq:v6v59-TT-space}
\end{equation}
to a centered Gaussian distribution \(h^{\rm TT}\).  For every smooth real
spatial-TT test tensor \(J\) with no spatial zero mode,
\begin{equation}
 \mathbb E e^{i\langle h^{\rm TT},J\rangle}
 =\exp\!\left[-{1\over2\kappa_g}
 \sum_{p_0,k\ne0}
 {\|\widehat J(p_0,k)\|^2\over p_0^2+|k|^2}\right].
 \label{eq:v6v59-TT-characteristic}
\end{equation}
\end{theorem}

\begin{proof}
The TT projector has rank two and operator norm one by
Lemma~\ref{lem:v6v59-TT-rank}.  Therefore
\[
 \mathbb E\|h_{\Lambda'}^{\rm TT}-h_\Lambda^{\rm TT}\|_{H^{-s}}^2
 \le {2\over\kappa_g}
 \sum_{\Lambda^2<p_0^2+|k|^2\le{\Lambda'}^2}
 {(1+p_0^2+|k|^2)^{-s}\over p_0^2+|k|^2}.
\]
The four-dimensional shell count is \(O(r^3)\), so the tail is bounded by
a constant times \(\int_\Lambda^\infty r^{1-2s}dr\), which tends to zero
for \(s>1\).  Completeness gives the limit.  The finite-cutoff Gaussian
characteristic functions converge, and the TT property of \(J\) removes
the projector from its quadratic form, proving
\eqref{eq:v6v59-TT-characteristic}.
\end{proof}

\begin{theorem}[Reflection positivity survives the TT cutoff limit]
\label{thm:v6v59-TT-OS-limit}
Let \(F\) be a bounded cylinder functional of finitely many smeared TT
fields supported at positive Euclidean times.  Then the limiting Gaussian
field satisfies
\begin{equation}
 \mathbb E\bigl[\overline{F(\vartheta h^{\rm TT})}
 F(h^{\rm TT})\bigr]\ge0,
 \label{eq:v6v59-TT-OS}
\end{equation}
where \(\vartheta\) reflects Euclidean time.
\end{theorem}

\begin{proof}
Choose a sequence of cutoffs containing all spatial test modes and
approximating each smooth time smearing.  At every cutoff the expectation is
nonnegative by Theorem~\ref{thm:v6v59-positive-transfer}.  The covariance
convergence in Theorem~\ref{thm:v6v59-TT-continuum} gives convergence in law
of the finite list of smeared fields.  Since \(F\) is bounded and continuous
on that finite-dimensional list, the reflected quadratic expectations
converge.  A limit of nonnegative real numbers is nonnegative.
\end{proof}

\subsection{Masslessness and curvature observables}

\begin{proposition}[The TT endpoint is massless]
\label{prop:v6v59-TT-massless}
The covariance multiplier in \eqref{eq:v6v59-TT-covariance} has pole
\((p_0^2+|k|^2)^{-1}\) and no positive mass term.  On a finite torus the
smallest nonzero quadratic eigenvalue is
\begin{equation}
 \kappa_g\min\left\{(2\pi/T)^2,(2\pi/L)^2\right\}
 \label{eq:v6v59-TT-gap}
\end{equation}
after the fixed zero sectors are removed, and it vanishes in the large-box
limit.
\end{proposition}

\begin{proof}
Diagonalize \eqref{eq:v6v59-TT-action} in the four Fourier variables.  The
eigenvalue is \(\kappa_g(p_0^2+|k|^2)\); minimizing over the retained
nonzero modes gives the stated scale, subject to the removal of all
\(k=0\) spatial modes.
\end{proof}

\begin{remark}[Derivative observables and inverse powers]
\label{rem:v6v59-curvature-observables}
The linearized Riemann tensor contains two derivatives of \(h^{\rm TT}\).
Its covariance therefore carries two net powers of momentum rather than the
inverse quadratic pole.  As in the V50 Maxwell distinction, uniform local
estimates should be formulated in physical curvature variables while the
massless potential gap remains explicit.
\end{remark}

\begin{firewall}[TT reduction is not nonlinear Einstein gravity]
\label{fw:v6v59-not-Einstein}
The projector \eqref{eq:v6v59-TT-projector} solves only the linearized flat
constraints and is spatially nonlocal.  The construction fixes homogeneous
metric moduli, omits the conformal factor, lapse, shift, nonlinear
Hamiltonian and momentum constraints, Faddeev--Popov strata, topology
change, matter backreaction, and the nonlinear Einstein vertices.  It
therefore supplies a positive free graviton baseline and its continuum
cylinder measure, not a measure on nonlinear geometries.
\end{firewall}

\begin{firewall}[Linear reflection positivity does not close CGMEC]
\label{fw:v6v59-CGMEC-open}
Theorem~\ref{thm:v6v59-TT-OS-limit} proves reflection positivity only for
the free TT Gaussian cylinder field.  CGMEC additionally requires a local
nonlinear constrained transfer, a positive treatment of the scalar sector,
and compatibility with the embedded gauge, Higgs, and determinant interface
amplitudes.  None follows from the oscillator product alone.
\end{firewall}

\begin{definition}[Transverse--traceless gravity baseline card]
\label{def:v6v59-TTGC}
TTGC consists of the finite-cutoff positive oscillator transfer,
cutoff-uniform convergence of TT cylinder observables, explicit removal or
retention of homogeneous modes, and the distributional continuum covariance
\eqref{eq:v6v59-TT-covariance}.  Theorems
\ref{thm:v6v59-positive-transfer}--\ref{thm:v6v59-TT-OS-limit} prove TTGC on
the fixed flat torus.
\end{definition}

\begin{theorem}[Sixth-series V59 result]
\label{thm:v6v59-status}
After linearized constraint and gauge reduction on a flat spatial torus,
the two transverse--traceless graviton polarizations form positive harmonic
oscillators.  Their one-step Mehler kernels define positive self-adjoint
transfer operators, their ultraviolet cutoffs converge in
\(H^{-s}\), \(s>1\), to the massless Gaussian TT field, and reflection
positivity survives on bounded cylinder functionals.  This proves TTGC but
leaves every nonlinear, conformal, homogeneous, and constraint-sector input
of CGMEC open.
\end{theorem}

\begin{proof}
Use Lemma~\ref{lem:v6v59-TT-rank},
Theorem~\ref{thm:v6v59-positive-transfer},
Corollary~\ref{cor:v6v59-Hamiltonian},
Theorems~\ref{thm:v6v59-TT-continuum} and
\ref{thm:v6v59-TT-OS-limit}, and
Proposition~\ref{prop:v6v59-TT-massless}, subject to both firewalls.
\end{proof}

V60 will perform the compulsory YM/NS audit.  It will compare their newest
complete-object and signed-shell results with the TT projection, then decide
whether the next acquisition should attack nonlinear constraint locality,
the scalar metric sector, or the common gauge--matter--TT interface.


\clearpage
\section{V60: companion audit and an additive positive-transfer bridge}
\label{sec:v6v60}

V59 constructs the free TT graviton transfer and its reflection-positive
continuum cylinder endpoint.  This compulsory audit finds that the YM
programme has retired its exact inverse-heat geometric bridge and replaced
it by an additive positive bridge with an endpoint envelope.  The NS
programme has also proved that neither slotwise norms nor a presumed sign
remove its one-logarithm-per-scale obstruction.  These conclusions determine
how a future nonlinear metric transfer may be compared with TT: by an
additive positive-operator path which keeps the signed endpoint difference
intact.

\subsection{Compulsory V60 audit}

The newest companion files inspected were
\begin{center}
\begin{tabular}{p{0.25\linewidth}p{0.16\linewidth}p{0.47\linewidth}}
programme & source & audited mathematical endpoint\\ \hline
Yang--Mills & \texttt{V42} & at the time of reading V42 was byte-identical
to V41; V41 proves that inverse-heat normalization of the exact geometric
bridge fails already for the marked singleton, then constructs an additive
bridge with contractive complete endpoint-envelope rows and leaves its merge
cocycle EMCC open\\
Navier--Stokes & \texttt{V38} & an explicit three-shell solution has an
octave cross term of either sign, so the type-I logarithm is not removed by a
sign argument; high-total cutoff reciprocals are summable, but the type-II
regime remains open
\end{tabular}
\end{center}
The YM V42 file had \(19\,796\) logical lines, \(626\,272\) bytes, and
SHA--256
\[
 \texttt{980514090d33a8c6799bf9aef6a011234e1ba62415f5b58de5e0434f4e5af7dc}.
\]
The NS V38 file had \(26\,590\) logical lines, \(1\,051\,488\) bytes, and
SHA--256
\[
 \texttt{b2d6db9a9ba7a54da42f2419f8f1bfc8e0d1ec9baf6fc86f49caed2f195b5707}.
\]

\begin{theorem}[V60 transfer decision]
\label{thm:v6v60-audit-decision}
The audit changes the metric acquisition order as follows.
\begin{enumerate}[label=\textup{(\roman*)},leftmargin=3.7em]
\item The free TT Gaussian may not be used as an inverse-heat normalization
      for an exact nonlinear metric density merely because its finite jets
      are controlled.
\item If a positive nonlinear constrained transfer is constructed, it may
      be connected to the TT transfer by an additive operator bridge.  The
      signed endpoint difference must remain inside the complete rooted
      expression until cancellation or locality is proved.
\item No sign is assigned to mixed nonlinear metric activities on the basis
      of the TT quadratic sign.  The NS witness rules out this proof pattern
      even in its own exact parabolic setting.
\item Neither companion constructs the missing positive scalar metric
      sector or the nonlinear constraint quotient.
\end{enumerate}
\end{theorem}

\begin{proof}
Items \textup{(i)--(ii)} are the exact scope of the YM V41 singleton no-go,
additive score identity, and envelope contraction.  Item \textup{(iii)} is
the NS V38 either-sign witness interpreted only as a restriction on proof
method.  Item \textup{(iv)} follows from the claim boundaries of both files.
\end{proof}

\begin{firewall}[No companion constant is imported]
\label{fw:v6v60-audit-types}
The YM central multipliers, NS octave fluxes, and metric transfer operators
live in different spaces.  The audit imports the demonstrated failure of
two proof patterns and the additive order of operations, not a quantitative
bound.
\end{firewall}

\subsection{Additive bridge for probability densities}

Let \((X,\lambda)\) be a measure space and let \(p_0,p_1\) be strictly
positive probability densities.  Define
\begin{equation}
 p_t:=(1-t)p_0+tp_1,
 \qquad
 \sigma_t:={p_1-p_0\over p_t},
 \qquad0\le t\le1.
 \label{eq:v6v60-density-bridge}
\end{equation}

\begin{theorem}[Exact additive density transport]
\label{thm:v6v60-density-bridge}
For every bounded measurable \(F\),
\begin{align}
 p_t\sigma_t&=p_1-p_0,
 \label{eq:v6v60-signed-density}\\
 \int\sigma_tp_t\,d\lambda&=0,
 \label{eq:v6v60-centered-score}\\
 {d\over dt}\int Fp_t\,d\lambda
 &=\int F(p_1-p_0)\,d\lambda
 =\int F\sigma_tp_t\,d\lambda.
 \label{eq:v6v60-density-derivative}
\end{align}
Moreover the absolute score length is independent of \(t\):
\begin{equation}
 \int|\sigma_t|p_t\,d\lambda
 =\int|p_1-p_0|\,d\lambda
 =2\|p_1\lambda-p_0\lambda\|_{\rm TV}.
 \label{eq:v6v60-TV-length}
\end{equation}
\end{theorem}

\begin{proof}
The first identity is the definition.  Its integral is
\(\int p_1-\int p_0=0\), giving the second.  Differentiate the affine
density to prove the third.  Taking the absolute value in the first identity
and integrating proves \eqref{eq:v6v60-TV-length}.
\end{proof}

\begin{corollary}[No density-ratio lower bound is needed]
\label{cor:v6v60-no-ratio}
The integrated transport identity and total-variation length remain finite
without a uniform lower bound on \(p_1/p_0\) or \(p_0/p_1\).  Pointwise
bounds on \(\sigma_t\) may fail near a thin endpoint, but the signed density
\(p_1-p_0\) remains exact.
\end{corollary}

\begin{proof}
All right sides in Theorem~\ref{thm:v6v60-density-bridge} depend only on the
two normalized endpoints and have total variation at most two.
\end{proof}

\subsection{Positive-operator endpoint envelope}

Let \(\mathcal H_*\) be a finite-cutoff interface Hilbert space and let
\(T_0,T_1\) be positive self-adjoint operators.  For the clean finite-cutoff
statement assume
\begin{equation}
 P:=T_0+T_1\succ0
 \label{eq:v6v60-envelope-P}
\end{equation}
and define
\begin{equation}
 T_t:=(1-t)T_0+tT_1,
 \qquad
 Y_t:=P^{-1/2}T_tP^{-1/2},
 \qquad
 Z:=P^{-1/2}(T_1-T_0)P^{-1/2}.
 \label{eq:v6v60-operator-bridge}
\end{equation}

\begin{theorem}[Contractive positive-transfer bridge]
\label{thm:v6v60-operator-bridge}
For every \(t\in[0,1]\),
\begin{equation}
 0\preceq Y_t\preceq I,
 \qquad
 -I\preceq Z\preceq I,
 \qquad
 \|Y_t\|_{\rm op}\le1,
 \qquad
 \|Z\|_{\rm op}\le1.
 \label{eq:v6v60-envelope-contractions}
\end{equation}
The derivative is the constant signed endpoint row
\begin{equation}
 {dY_t\over dt}=Z.
 \label{eq:v6v60-operator-derivative}
\end{equation}
If \(T_0\) and \(T_1\) are bounded positive self-adjoint transfer
operators, then every \(T_t\) is also a bounded positive self-adjoint
transfer operator and obeys the finite-step reflection quadratic inequality.
\end{theorem}

\begin{proof}
Since \(0\preceq T_i\preceq P\), conjugation by \(P^{-1/2}\) gives
\(0\preceq P^{-1/2}T_iP^{-1/2}\preceq I\).  Convexity gives the same for
\(Y_t\).  The form inequalities
\(-P\preceq T_1-T_0\preceq P\) give \(-I\preceq Z\preceq I\).
Self-adjoint form bounds imply the two norm estimates.  Differentiation gives
\eqref{eq:v6v60-operator-derivative}.  Finally the positive cone is convex,
so \(T_t\succeq0\); Theorem~\ref{thm:v6v20-transfer-positivity} supplies the
reflection inequality.
\end{proof}

\begin{remark}[Common kernels]
\label{rem:v6v60-kernel}
If \(P\) has a kernel, positivity implies
\(\ker P=\ker T_0\cap\ker T_1\).  One may restrict every operator to
\(\overline{\operatorname{ran}P}\) and interpret the normalized sandwiches
by their quadratic forms, or regularize with \(P+\epsilon I\) and pass to
the support limit.  At finite cutoff this is exactly the zero-sector rule of
V57.
\end{remark}

\begin{corollary}[Endpoint amplitude derivative]
\label{cor:v6v60-amplitude}
For every interface vector \(f\),
\begin{equation}
 {d\over dt}\langle f,T_tf\rangle
 =\langle f,(T_1-T_0)f\rangle.
 \label{eq:v6v60-amplitude-derivative}
\end{equation}
Writing \(g=P^{1/2}f\), the right side is
\(\langle g,Zg\rangle\) and hence has modulus at most \(\|g\|^2\).
\end{corollary}

\begin{proof}
Differentiate the affine path and use
\(T_1-T_0=P^{1/2}ZP^{1/2}\), followed by
\(\|Z\|\le1\).
\end{proof}

\subsection{Tensor blocks and the merge obstruction}

For every interface block \(B\), suppose positive endpoints
\(T_{0,B},T_{1,B}\) and their envelope
\(P_B=T_{0,B}+T_{1,B}\) have been defined.  For a partition
\(\pi\) of \(B\), set
\begin{equation}
 \mathfrak C_{B,\pi}^{\rm op}
 :=P_B^{-1/2}
 \left(\bigotimes_{A\in\pi}P_A\right)P_B^{-1/2}.
 \label{eq:v6v60-operator-cocycle}
\end{equation}

\begin{proposition}[Exact normalized partition identity]
\label{prop:v6v60-partition-identity}
Let \(M_{t,A}\) be block moments or block transfer insertions and put
\(Y_{t,A}=P_A^{-1/2}M_{t,A}P_A^{-1/2}\).  Every moment--cumulant partition
term normalized by \(P_B\) has the exact form
\begin{equation}
 P_B^{-1/2}
 \left(\bigotimes_{A\in\pi}M_{t,A}\right)P_B^{-1/2}
 =\mathfrak C_{B,\pi}^{\rm op}
 \left(\bigotimes_{A\in\pi}Y_{t,A}\right),
 \label{eq:v6v60-normalized-partition}
\end{equation}
with the canonical block embeddings understood.  A derivative replaces one
\(Y_{t,A}\) by its signed normalized endpoint difference.
\end{proposition}

\begin{proof}
Insert
\(M_{t,A}=P_A^{1/2}Y_{t,A}P_A^{1/2}\) in every block and collect the
envelope factors on the outside.  The affine derivative is the product rule
together with \eqref{eq:v6v60-operator-derivative} on the differentiated
block.
\end{proof}

\begin{firewall}[Contractive blocks do not bound the merge cocycle]
\label{fw:v6v60-merge-open}
Theorem~\ref{thm:v6v60-operator-bridge} controls each complete block row,
but gives no all-arity estimate on
\(\mathfrak C_{B,\pi}^{\rm op}\).  Separate partition norms can produce
Bell-number or worse growth.  The complete connected forest or transfer
composition must retain the cocycle with the signed endpoint row, exactly as
the audited YM EMCC requires.  No nonlinear gravity merge card is proved
here.
\end{firewall}

\subsection{Application boundary for the TT baseline}

\begin{definition}[Positive operator transport card]
\label{def:v6v60-POTC}
POTC for a nonlinear constrained metric endpoint \(T_1\) relative to the TT
endpoint \(T_0\) consists of:
\begin{enumerate}[label=\textup{(PO\arabic*)},leftmargin=4.8em]
\item positive self-adjoint endpoint transfers on a common constraint and
      boundary Hilbert space;
\item a controlled endpoint envelope and signed difference as in
      Theorem~\ref{thm:v6v60-operator-bridge};
\item an all-block merge-cocycle bound applied to the complete differentiated
      expression before partition norms; and
\item cutoff convergence compatible with the V14 interface and V59 TT
      cylinder limit.
\end{enumerate}
\end{definition}

\begin{proposition}[Exact status of POTC]
\label{prop:v6v60-POTC-status}
The TT endpoint \(T_0\) and its positive transfer are proved in V59.  If a
positive nonlinear endpoint \(T_1\) on the same Hilbert space is supplied,
Theorem~\ref{thm:v6v60-operator-bridge} proves
POTC\textup{(PO1)--(PO2)} at every finite cutoff.  The nonlinear endpoint,
merge card, and common cutoff limit remain open.
\end{proposition}

\begin{proof}
Apply the cited TT and operator-bridge theorems.  The remaining statements
are precisely the hypotheses not constructed in V58--V60.
\end{proof}

\begin{firewall}[An additive path cannot repair a nonpositive endpoint]
\label{fw:v6v60-nonpositive-endpoint}
The unrotated Euclidean Einstein density is not a positive integrable
endpoint by V20.  It cannot be inserted as \(T_1\) in
Theorem~\ref{thm:v6v60-operator-bridge}.  Convex interpolation preserves
positivity only after both endpoints have independently been constructed as
positive operators.
\end{firewall}

\begin{theorem}[Sixth-series V60 result]
\label{thm:v6v60-status}
The compulsory audit retires inverse-heat comparison as a valid exact
nonlinear metric strategy and forbids removal of scale logarithms by an
unproved sign.  For any two positive finite-cutoff endpoint transfers, their
additive bridge remains positive and its endpoint-envelope normalized
unmarked and signed rows are contractions.  The exact partition identity
isolates an operator merge cocycle which must be controlled before all-order
transport.  Applied to gravity, this proves the finite-cutoff bridge
mechanism conditional on a positive nonlinear constrained endpoint; that
endpoint and its merge card remain open.
\end{theorem}

\begin{proof}
Use Theorem~\ref{thm:v6v60-audit-decision},
Theorems~\ref{thm:v6v60-density-bridge} and
\ref{thm:v6v60-operator-bridge}, and
Proposition~\ref{prop:v6v60-partition-identity}, subject to the three
firewalls.
\end{proof}

V61 will go upstream to the scalar constraint sector.  It will revisit the
linearized flat constraint symbol already isolated in V22, distinguish an
exactly removed conformal coordinate from a merely rotated one, and derive
the reduced phase-space transfer on the physical TT variables without
repeating the V59 oscillator construction.


\clearpage
\section{V61: canonical linearized constraint reduction and the TT phase space}
\label{sec:v6v61}

V22 computes the flat Hamiltonian and momentum constraint symbols and proves
that their projectors are bounded but nonlocal.  V59 then constructs the
positive TT oscillator transfer.  This version supplies the missing bridge:
it performs the complete linearized canonical reduction at every nonzero
spatial momentum.  The vacuum conformal scalar is removed by an exact
constraint and gauge quotient, not by a contour rotation.  With matter
present it is instead fixed by the Newton equation and remains a physical
long-range carrier.

\subsection{Mode coordinates for the ADM phase space}

At a fixed nonzero spatial momentum rotate coordinates so
\begin{equation}
 k=(0,0,\kappa),
 \qquad\kappa=|k|>0.
 \label{eq:v6v61-k-frame}
\end{equation}
The real canonical variables are a symmetric metric perturbation \(h_{ij}\)
and its symmetric momentum \(\pi_{ij}\), with symplectic form
\begin{equation}
 \Omega_k=\sum_{i\le j}d\pi_{ij}\wedge dh_{ij}
 \label{eq:v6v61-symplectic}
\end{equation}
after the harmless conventional rescaling of off-diagonal components which
makes the symmetric-tensor inner product Euclidean.

The V22 linearized constraints become
\begin{align}
 \mathcal H_k(h)
 &=-\kappa^2(h_{11}+h_{22})=0,
 \label{eq:v6v61-H-constraint}\\
 \mathcal M_{i,k}(\pi)
 &=-2i\kappa\pi_{i3}=0,
 \qquad i=1,2,3.
 \label{eq:v6v61-M-constraints}
\end{align}
Spatial diffeomorphisms act on configuration variables by
\begin{equation}
 \delta_\xi h_{ij}=i(k_i\xi_j+k_j\xi_i),
 \label{eq:v6v61-spatial-gauge}
\end{equation}
and the scalar constraint gauge orbit acts on momentum by
\begin{equation}
 \delta_f\pi_{ij}=f(k_i k_j-\kappa^2\delta_{ij})
 =-f\kappa^2P_{ij}(k),
 \label{eq:v6v61-scalar-gauge}
\end{equation}
up to a fixed nonzero convention factor absorbed in \(f\).

\begin{lemma}[Explicit constraint and gauge coordinates]
\label{lem:v6v61-coordinate-actions}
In the frame \eqref{eq:v6v61-k-frame}:
\begin{enumerate}[label=\textup{(\roman*)},leftmargin=3.7em]
\item the Hamiltonian constraint sets \(h_{11}+h_{22}=0\);
\item the momentum constraints set
      \(\pi_{13}=\pi_{23}=\pi_{33}=0\);
\item spatial diffeomorphisms independently remove
      \(h_{13},h_{23},h_{33}\); and
\item the scalar gauge orbit independently removes
      \(\pi_{11}+\pi_{22}\).
\end{enumerate}
\end{lemma}

\begin{proof}
Items \textup{(i)--(ii)} are
\eqref{eq:v6v61-H-constraint}--\eqref{eq:v6v61-M-constraints}.  Equation
\eqref{eq:v6v61-spatial-gauge} gives
\(\delta h_{13}=i\kappa\xi_1\),
\(\delta h_{23}=i\kappa\xi_2\), and
\(\delta h_{33}=2i\kappa\xi_3\), so the three components can be set to zero
uniquely when \(\kappa\ne0\).  Since
\(P=\operatorname{diag}(1,1,0)\),
\eqref{eq:v6v61-scalar-gauge} changes \(\pi_{11}\) and \(\pi_{22}\) by
the same arbitrary amount and leaves their difference and \(\pi_{12}\)
unchanged.  This removes their sum.
\end{proof}

\subsection{The reduced symplectic space}

Define TT coordinates
\begin{align}
 h_+&={h_{11}-h_{22}\over\sqrt2},
 &h_\times&=\sqrt2h_{12},
 \label{eq:v6v61-h-polarizations}\\
 \pi_+&={\pi_{11}-\pi_{22}\over\sqrt2},
 &\pi_\times&=\sqrt2\pi_{12}.
 \label{eq:v6v61-pi-polarizations}
\end{align}

\begin{theorem}[Exact nonzero-mode phase-space reduction]
\label{thm:v6v61-phase-reduction}
At each \(k\ne0\), the joint constraint surface
\eqref{eq:v6v61-H-constraint}--\eqref{eq:v6v61-M-constraints}, quotiented by
the gauge orbits \eqref{eq:v6v61-spatial-gauge}--
\eqref{eq:v6v61-scalar-gauge}, is the four-dimensional real phase space
\begin{equation}
 (h_+,h_\times,\pi_+,\pi_\times)\in\mathbb R^4.
 \label{eq:v6v61-reduced-phase-space}
\end{equation}
The reduced symplectic form is nondegenerate and equals
\begin{equation}
 \Omega_k^{\rm red}
 =d\pi_+\wedge dh_++d\pi_\times\wedge dh_\times.
 \label{eq:v6v61-reduced-symplectic}
\end{equation}
These are precisely the two polarizations in the range of the TT projector
\eqref{eq:v6v59-TT-projector}.
\end{theorem}

\begin{proof}
Lemma~\ref{lem:v6v61-coordinate-actions} removes four configuration
coordinates and four momentum coordinates from the twelve-dimensional
canonical mode space, leaving the four coordinates in
\eqref{eq:v6v61-reduced-phase-space}.  Substitute the constraint and gauge
conditions into \eqref{eq:v6v61-symplectic}.  The terms involving removed
coordinates vanish, while the orthonormal changes
\eqref{eq:v6v61-h-polarizations}--\eqref{eq:v6v61-pi-polarizations} give
\eqref{eq:v6v61-reduced-symplectic}.  A tensor with only
\(h_{11}=-h_{22}\) and \(h_{12}=h_{21}\) nonzero is transverse to
\(k=(0,0,\kappa)\) and traceless, and every TT tensor has this form.
\end{proof}

\begin{corollary}[The vacuum conformal scalar is removed, not rotated]
\label{cor:v6v61-conformal-removed}
For a nonzero Fourier conformal perturbation
\(h_{ij}=2\phi\delta_{ij}\), the vacuum Hamiltonian constraint gives
\begin{equation}
 -4\kappa^2\phi=0,
 \label{eq:v6v61-conformal-zero}
\end{equation}
and hence \(\phi=0\).  The variable is absent from the reduced phase space;
no complex contour is used.
\end{corollary}

\begin{proof}
The conformal tensor has \(h_{11}+h_{22}=4\phi\).  Insert it into
\eqref{eq:v6v61-H-constraint} and use \(\kappa>0\).
\end{proof}

\begin{remark}[Why this does not contradict the V20 no-go]
\label{rem:v6v61-no-contradiction}
V20 integrates the real Euclidean conformal coordinate before an exact
canonical reduction and proves that integral unbounded.  The present vacuum
calculation restricts to the linear constraint surface and quotients its
first-class gauge orbit before forming the transfer.  These are different
spaces.  The reduction is a legitimate positive baseline only where the
constraint rank and quotient are controlled.
\end{remark}

\subsection{Reduced Hamiltonian and TT transfer}

After a positive canonical rescaling, the quadratic ADM Hamiltonian on the
reduced mode takes the form
\begin{equation}
 H_k^{\rm red}
 ={1\over2}\left(
 \pi_+^2+\pi_\times^2
 +\kappa^2h_+^2+\kappa^2h_\times^2
 \right).
 \label{eq:v6v61-reduced-Hamiltonian}
\end{equation}

\begin{theorem}[Canonical derivation of the V59 TT oscillator]
\label{thm:v6v61-TT-Hamiltonian}
The reduced Hamiltonian \eqref{eq:v6v61-reduced-Hamiltonian} is the sum of
two positive oscillators of frequency \(|k|\).  Canonical quantization gives
the one-mode operators and Mehler kernels of
Theorem~\ref{thm:v6v59-positive-transfer}, after inserting the common
normalization \(\kappa_g\).  Hence the V59 transfer is the positive transfer
of the complete nonzero-mode linearized vacuum constraint quotient.
\end{theorem}

\begin{proof}
The symplectic form \eqref{eq:v6v61-reduced-symplectic} makes
\((h_+,\pi_+)\) and \((h_\times,\pi_\times)\) canonical pairs.  Each pair
has Hamiltonian \((\pi^2+\kappa^2h^2)/2\).  Rescaling
\(h\mapsto\kappa_g^{-1/2}q\),
\(\pi\mapsto\kappa_g^{1/2}p\) preserves the symplectic form and gives the
operator \eqref{eq:v6v59-oscillator-H}.  Its heat kernel is
\eqref{eq:v6v59-Mehler}.
\end{proof}

\subsection{Constraint Jacobians and the zero momentum boundary}

\begin{proposition}[Finite-cutoff reduction Jacobian]
\label{prop:v6v61-Jacobian}
At fixed \(k\ne0\), the linear maps used to fix
\((h_{13},h_{23},h_{33})\) from \((\xi_1,\xi_2,\xi_3)\) and to impose the
scalar constraint on \(h_{11}+h_{22}\) have nonzero determinants
proportional to \(\kappa^3\) and \(\kappa^2\), respectively.  Their product
is independent of the TT coordinates and is absorbed into the
finite-cutoff normalization.  It becomes singular at \(k=0\), so the
homogeneous modes cannot be included by continuity in this chart.
\end{proposition}

\begin{proof}
Lemma~\ref{lem:v6v61-coordinate-actions} gives the diagonal gauge-fixing
map \(i\kappa\operatorname{diag}(1,1,2)\), whose determinant is a nonzero
constant times \(\kappa^3\).  The scalar constraint is multiplication of
the transverse trace by \(-\kappa^2\).  Neither coefficient depends on the
TT variables.  Both vanish at \(\kappa=0\).
\end{proof}

\begin{corollary}[Homogeneous metric modes remain interface sectors]
\label{cor:v6v61-zero-modes}
The spatial zero modes, including torus shape and volume moduli, must be
fixed or retained as finite-dimensional boundary variables.  They are not
part of the V59 TT Gaussian obtained by the nonzero-mode quotient.
\end{corollary}

\begin{proof}
At \(k=0\), the constraints and gauge-fixing maps used above lose rank by
Proposition~\ref{prop:v6v61-Jacobian}.  The zero-mode rule of V57 therefore
requires a separate sector treatment.
\end{proof}

\subsection{Matter restores the scalar as a determined Newton field}

With energy-density source \(\rho(k)\), the Hamiltonian constraint is
\begin{equation}
 -\kappa^2(h_{11}+h_{22})=\kappa_N\rho(k),
 \label{eq:v6v61-sourced-H}
\end{equation}
where \(\kappa_N\) is the gravitational coupling in the chosen
normalization.

\begin{theorem}[Sourced scalar is constrained but not absent]
\label{thm:v6v61-sourced-scalar}
For \(k\ne0\), the transverse scalar trace is fixed by
\begin{equation}
 h_{11}+h_{22}=-{\kappa_N\over|k|^2}\rho(k).
 \label{eq:v6v61-Newton-trace}
\end{equation}
After quotienting the same gauge directions, the dynamical metric pairs are
still TT, but the reconstructed metric contains the nonlocal Newton carrier
\(|k|^{-2}\rho(k)\).  This carrier vanishes only when the complete matter
source is zero or exactly divisible by the lattice Laplacian in the sense
of V22.
\end{theorem}

\begin{proof}
Solve \eqref{eq:v6v61-sourced-H} for the transverse trace.  The spatial
diffeomorphism and scalar gauge actions are unchanged, so the quotient of
the freely dynamical pairs remains the TT phase space.  The inverse factor
\(|k|^{-2}\) is nonanalytic at zero and is exactly the scalar source solution
of Theorem~\ref{thm:v6v22-Newton-pole}.  V22's local inverse criterion gives
the stated divisibility alternative.
\end{proof}

\begin{firewall}[A delta constraint is not a positive nonlinear measure]
\label{fw:v6v61-delta-not-nonlinear}
At linear vacuum level, solving a constant-rank constraint and dividing by
its linear gauge orbit is exact.  In nonlinear Einstein--matter theory the
constraint surface can curve, change rank, acquire stabilizers, and couple
the scalar solution to all matter and TT modes.  The associated Jacobian,
ghost determinant, and transfer kernel must be proved positive and local in
the hybrid sense.  The linear quotient does not construct them.
\end{firewall}

\begin{firewall}[The Newton carrier remains physical]
\label{fw:v6v61-Newton-physical}
The matter-sourced scalar is not a redundant conformal integration
coordinate.  Eliminating it produces a long-range matter interaction, and
its sign and stability must be analyzed together with the matter
Hamiltonian.  Dropping \eqref{eq:v6v61-Newton-trace} would erase physical
gravity rather than fix a gauge.
\end{firewall}

\begin{definition}[Canonical constraint reduction card]
\label{def:v6v61-CCRC}
CCRC requires, at every cutoff and on every retained bounded-geometry chart:
\begin{enumerate}[label=\textup{(CR\arabic*)},leftmargin=4.8em]
\item a constant-rank Hamiltonian and momentum constraint surface with a
      controlled gauge quotient and zero-mode sectors;
\item a nondegenerate reduced symplectic form whose free tangent is the TT
      phase space;
\item a positive self-adjoint reduced transfer including the reconstructed
      Newton carrier for matter sources; and
\item cutoff-uniform Jacobian, ghost, locality, and merge bounds compatible
      with POTC and the V23 hybrid massless norm.
\end{enumerate}
\end{definition}

\begin{theorem}[Sixth-series V61 result]
\label{thm:v6v61-status}
At every nonzero spatial momentum, the complete linearized vacuum ADM
constraint surface modulo its four first-class gauge directions is exactly
the four-dimensional phase space of two TT canonical pairs.  Its Hamiltonian
is the positive two-oscillator Hamiltonian underlying the V59 transfer.  The
vacuum conformal Fourier mode is removed by the scalar constraint rather
than rotated.  With matter energy density, the scalar trace is instead fixed
to \(-\kappa_N\rho/|k|^2\) and remains the physical Newton carrier.  CCRC is
therefore proved only for the nonzero-mode linear vacuum tangent; its
nonlinear sourced extension remains open.
\end{theorem}

\begin{proof}
Use Lemma~\ref{lem:v6v61-coordinate-actions},
Theorem~\ref{thm:v6v61-phase-reduction},
Corollary~\ref{cor:v6v61-conformal-removed},
Theorem~\ref{thm:v6v61-TT-Hamiltonian},
Proposition~\ref{prop:v6v61-Jacobian}, and
Theorem~\ref{thm:v6v61-sourced-scalar}, subject to both firewalls.
\end{proof}

V62 will test stability of the reconstructed Newton sector.  It will derive
the exact quadratic completion for a positive matter density mode coupled to
the constrained scalar, identify the sign of the induced
\(|k|^{-2}\) interaction, and establish the strongest cutoff-dependent
lower bound available from a local matter kinetic energy.


\clearpage
\section{V62: the Newton completion and conditional matter stability}
\label{sec:v6v62}

V61 proves that the vacuum scalar metric mode is removed by canonical
constraint reduction, whereas a matter source reconstructs it with the
Newton multiplier \(|k|^{-2}\).  This version computes the sign produced by
eliminating a positive scalar constraint carrier and proves two sharp kinds
of stability for a scalar matter model: a finite-volume quartic criterion
and an infinite-volume lower bound at fixed \(L^2\) charge.  Neither estimate
is uniform over unrestricted Standard Model field amplitudes.

\subsection{Quadratic completion of the scalar carrier}

At a nonzero Fourier mode let \(s(k)\) be a real scalar carrier and
\(\rho(k)\) a real matter density mode.  Consider
\begin{equation}
 Q_k(s,\rho)
 ={a\over2}|k|^2|s|^2
 +b\operatorname{Re}(\overline{s}\rho),
 \qquad a>0.
 \label{eq:v6v62-scalar-quadratic}
\end{equation}

\begin{theorem}[Exact Newton completion]
\label{thm:v6v62-Newton-completion}
The minimizing constrained scalar is
\begin{equation}
 s_*(k)=-{b\over a}{\rho(k)\over|k|^2},
 \label{eq:v6v62-s-star}
\end{equation}
and the reduced matter form is
\begin{equation}
 \inf_sQ_k(s,\rho)
 =-{b^2\over2a}{|\rho(k)|^2\over|k|^2}.
 \label{eq:v6v62-Newton-mode}
\end{equation}
Thus elimination of a positive scalar carrier induces an attractive
negative Newton interaction.
\end{theorem}

\begin{proof}
Complete the square:
\[
 Q_k(s,\rho)
 ={a|k|^2\over2}
 \left|s+{b\rho\over a|k|^2}\right|^2
 -{b^2\over2a}{|\rho|^2\over|k|^2}.
\]
The square has the unique minimizer and minimum stated.
\end{proof}

\begin{remark}[Schur sign]
\label{rem:v6v62-Schur-sign}
Equation \eqref{eq:v6v62-Newton-mode} is the scalar instance of the V58
Schur order: integrating a positive block subtracts a positive quadratic
form from the retained variables.  Positivity of the carrier integration
therefore does not imply positivity of the induced matter Hamiltonian.
\end{remark}

\subsection{Finite-volume spectral control}

Let \(\mathbb T_L^3\) have volume \(L^3\), let
\(\lambda_1=(2\pi/L)^2\) be the first positive eigenvalue of
\(-\Delta\), and define for a real density \(\rho\)
\begin{equation}
 \rho_0:=\rho-L^{-3}\int_{\mathbb T_L^3}\rho,
 \qquad
 \mathcal D_L(\rho)
 :=\langle\rho_0,(-\Delta)^{-1}\rho_0\rangle.
 \label{eq:v6v62-DL}
\end{equation}

\begin{lemma}[Finite-volume Newton operator bound]
\label{lem:v6v62-torus-bound}
For every \(\rho\in L^2(\mathbb T_L^3)\),
\begin{equation}
 0\le\mathcal D_L(\rho)
 \le\lambda_1^{-1}\|\rho_0\|_2^2
 \le\lambda_1^{-1}\|\rho\|_2^2.
 \label{eq:v6v62-torus-Newton}
\end{equation}
\end{lemma}

\begin{proof}
On mean-zero Fourier modes the positive operator \((-Delta)^{-1}\) has
eigenvalues \(|k|^{-2}\le\lambda_1^{-1}\).  This gives the first two
inequalities.  Orthogonal projection to mean-zero functions is an
\(L^2\)-contraction, giving the last.
\end{proof}

Let \(\psi:\mathbb T_L^3\to\mathbb C^N\) be a scalar multiplet and put
\(\rho=|\psi|^2\).  Consider the static model energy
\begin{equation}
 \mathcal E_L(\psi)
 ={1\over2}\|\nabla\psi\|_2^2
 +{m^2\over2}\|\psi\|_2^2
 +{\lambda_H\over4}\|\psi\|_4^4
 -{\Gamma\over2}\mathcal D_L(|\psi|^2),
 \label{eq:v6v62-torus-energy}
\end{equation}
where \(\Gamma=b^2/a>0\).

\begin{theorem}[Finite-volume quartic stability margin]
\label{thm:v6v62-quartic-margin}
If
\begin{equation}
 \lambda_H>2\Gamma\lambda_1^{-1}
 ={\Gamma L^2\over2\pi^2},
 \label{eq:v6v62-quartic-condition}
\end{equation}
then
\begin{equation}
 \mathcal E_L(\psi)
 \ge {1\over2}\|\nabla\psi\|_2^2
 +{m^2\over2}\|\psi\|_2^2
 +{1\over4}
 \left(\lambda_H-2\Gamma\lambda_1^{-1}\right)
 \|\psi\|_4^4.
 \label{eq:v6v62-torus-lower}
\end{equation}
In particular the model energy is coercive in
\(H^1\cap L^4\) when \(m^2\ge0\).
\end{theorem}

\begin{proof}
Apply Lemma~\ref{lem:v6v62-torus-bound} to
\(\rho=|\psi|^2\):
\[
 \mathcal D_L(|\psi|^2)
 \le\lambda_1^{-1}\||\psi|^2\|_2^2
 =\lambda_1^{-1}\|\psi\|_4^4.
\]
Insert this in \eqref{eq:v6v62-torus-energy} and collect quartic
coefficients.  Strict positivity of all displayed leading coefficients gives
coercivity.
\end{proof}

\begin{firewall}[The spectral quartic margin is not thermodynamic]
\label{fw:v6v62-volume-margin}
Since \(\lambda_1^{-1}=L^2/(4\pi^2)\), the sufficient condition
\eqref{eq:v6v62-quartic-condition} deteriorates with the box size.  It is a
valid finite-volume cutoff criterion, not a uniform infinite-volume
Standard Model stability theorem.
\end{firewall}

\subsection{Infinite-volume form bound at fixed matter charge}

On \(\mathbb R^3\), define the Coulomb form
\begin{equation}
 \mathcal D_{\mathbb R^3}(|\psi|^2)
 :=\iint_{\mathbb R^3\times\mathbb R^3}
 { |\psi(x)|^2|\psi(y)|^2\over4\pi|x-y|}\,dx\,dy.
 \label{eq:v6v62-Coulomb}
\end{equation}
Let \(N=\|\psi\|_2^2\).

\begin{theorem}[Infinitesimal kinetic form bound at fixed \(L^2\) charge]
\label{thm:v6v62-HLS-bound}
There is a numerical Sobolev--HLS constant \(C_*<\infty\) such that for
every \(\psi\in H^1(\mathbb R^3;\mathbb C^N)\),
\begin{equation}
 \mathcal D_{\mathbb R^3}(|\psi|^2)
 \le C_*N^{3/2}\|\nabla\psi\|_2.
 \label{eq:v6v62-Coulomb-GN}
\end{equation}
Consequently, for every \(\epsilon>0\),
\begin{equation}
 {1\over2}\|\nabla\psi\|_2^2
 -{\Gamma\over2}\mathcal D_{\mathbb R^3}(|\psi|^2)
 \ge {1-\epsilon\over2}\|\nabla\psi\|_2^2
 -{\Gamma^2C_*^2\over8\epsilon}N^3.
 \label{eq:v6v62-fixed-N-lower}
\end{equation}
For every fixed \(N\), the attractive Newton form is therefore
infinitesimally form-bounded relative to the local kinetic energy.
\end{theorem}

\begin{proof}
The Hardy--Littlewood--Sobolev inequality and
\(\rho=|\psi|^2\) give
\begin{equation}
 \mathcal D_{\mathbb R^3}(|\psi|^2)
 \le C_{\rm HLS}\||\psi|^2\|_{6/5}^2
 =C_{\rm HLS}\|\psi\|_{12/5}^4.
 \label{eq:v6v62-HLS-step}
\end{equation}
Interpolate between \(L^2\) and \(L^6\):
\begin{equation}
 \|\psi\|_{12/5}
 \le\|\psi\|_2^{3/4}\|\psi\|_6^{1/4}
 \le C_S^{1/4}\|\psi\|_2^{3/4}
 \|\nabla\psi\|_2^{1/4}.
 \label{eq:v6v62-interpolation}
\end{equation}
Raise to the fourth power and combine with
\eqref{eq:v6v62-HLS-step} to obtain
\eqref{eq:v6v62-Coulomb-GN}, with \(C_*=C_{\rm HLS}C_S\).  Finally apply
Young's inequality
\[
 {\Gamma C_*\over2}N^{3/2}X
 \le{\epsilon\over2}X^2
 +{\Gamma^2C_*^2\over8\epsilon}N^3
\]
with \(X=\|\nabla\psi\|_2\).
\end{proof}

\begin{proposition}[The charge dependence is consistent with scaling]
\label{prop:v6v62-scaling}
For the fixed-charge dilation
\begin{equation}
 \psi_\ell(x)=\ell^{3/2}\psi(\ell x),
 \label{eq:v6v62-dilation}
\end{equation}
one has
\begin{equation}
 \|\psi_\ell\|_2^2=N,
 \qquad
 \|\nabla\psi_\ell\|_2^2
 =\ell^2\|\nabla\psi\|_2^2,
 \qquad
 \mathcal D_{\mathbb R^3}(|\psi_\ell|^2)
 =\ell\mathcal D_{\mathbb R^3}(|\psi|^2).
 \label{eq:v6v62-scaling-laws}
\end{equation}
Thus kinetic energy dominates arbitrarily sharp concentration at fixed
\(N\), while the optimized negative lower scale is of order
\(-\Gamma^2N^3\), consistent with
\eqref{eq:v6v62-fixed-N-lower}.
\end{proposition}

\begin{proof}
Change variables \(u=\ell x\) in each integral.  The gradient gains one
factor \(\ell\), and the Coulomb kernel gains one factor \(\ell\).  A
quadratic function \(A\ell^2-B\Gamma\ell\) has minimum of order
\(-\Gamma^2B^2/A\); the amplitude dependence of \(B^2/A\) has the
\(N^3\) homogeneity in Theorem~\ref{thm:v6v62-HLS-bound}.
\end{proof}

\subsection{Why the fixed-charge theorem does not construct a field measure}

\begin{proposition}[No uniform lower bound from the fixed-charge estimate]
\label{prop:v6v62-no-uniform-N}
The right side of \eqref{eq:v6v62-fixed-N-lower} tends to
\(-\infty\) like \(-N^3\) when the \(L^2\) charge is unrestricted.  Hence
Theorem~\ref{thm:v6v62-HLS-bound} alone does not prove normalizability of a
Euclidean scalar field measure which integrates over all \(N\).
\end{proposition}

\begin{proof}
The theorem supplies a family of lower bounds indexed by \(N\), whose
constant is \(-C\Gamma^2N^3\).  Taking an infimum over all
\(N\ge0\) gives \(-\infty\).  Additional amplitude control must be proved
from the complete Higgs potential, gauge constraints, finite-cutoff measure,
or another positive Hamiltonian mechanism.
\end{proof}

\begin{assessment}[Full energy density is harder than the mass-density
model]
\label{ass:v6v62-full-density}
The source \(\rho=|\psi|^2\) is a controlled scalar model of the Newton
completion.  The gravitational Hamiltonian constraint is sourced by the
complete energy density, including matter gradients, gauge fields,
fermions, and interactions.  Applying \((-Delta)^{-1}\) to quadratic
gradient densities can require higher Sobolev control than
\eqref{eq:v6v62-Coulomb-GN}.  No replacement of the full source by
\(|\psi|^2\) is made.
\end{assessment}

\begin{definition}[Newton-source stability card]
\label{def:v6v62-NSSC}
NSSC requires, for the complete constrained SM energy density at each
cutoff:
\begin{enumerate}[label=\textup{(NSC\arabic*)},leftmargin=5.3em]
\item an exact solution of the scalar gravitational constraint with all
      zero modes and boundary charges retained;
\item a lower bound on the induced Newton form by the positive matter and
      TT Hamiltonians, uniform in volume and cutoff;
\item an amplitude stock which is integrated by the positive reference
      measure and controls the analogue of the \(N^3\) constant; and
\item hybrid power-law locality and positive transfer compatible with CCRC
      and POTC.
\end{enumerate}
\end{definition}

\begin{firewall}[Classical form bounds do not imply Osterwalder--Schrader
positivity]
\label{fw:v6v62-OS}
Theorems~\ref{thm:v6v62-quartic-margin} and
\ref{thm:v6v62-HLS-bound} are lower-form estimates for static scalar models.
They do not construct the full constrained one-step kernel, control its
Jacobian and determinant phase, or prove reflection positivity after the
instantaneous Newton interaction is included.
\end{firewall}

\begin{theorem}[Sixth-series V62 result]
\label{thm:v6v62-status}
Eliminating a positive scalar constraint carrier induces the negative
Newton form \(-\Gamma|\rho(k)|^2/(2|k|^2)\).  On a finite three-torus it is
controlled by a local scalar quartic only under the explicit volume-dependent
margin \(\lambda_H>2\Gamma\lambda_1^{-1}\).  On \(\mathbb R^3\), for the
model density \(|\psi|^2\), it is infinitesimally form-bounded by the local
kinetic energy at fixed \(L^2\) charge, with lower constant of order
\(-\Gamma^2N^3\).  The estimates do not control unrestricted field charge
or the complete SM energy source; these are the live inputs of NSSC.
\end{theorem}

\begin{proof}
Use Theorem~\ref{thm:v6v62-Newton-completion},
Lemma~\ref{lem:v6v62-torus-bound},
Theorem~\ref{thm:v6v62-quartic-margin},
Theorem~\ref{thm:v6v62-HLS-bound}, and
Propositions~\ref{prop:v6v62-scaling} and
\ref{prop:v6v62-no-uniform-N}, subject to the two firewalls.
\end{proof}

V63 will test the amplitude stock in a finite lattice scalar reference.  It
will combine a positive onsite Higgs quartic with the discrete Newton kernel,
derive an exact matrix copositivity condition, and distinguish a
volume-uniform local stability mechanism from a bound that merely uses the
lowest Laplacian eigenvalue.


\clearpage
\section{V63: Newton--Higgs copositivity and local neutral-cell flux}
\label{sec:v6v63}

V62 controls the Newton interaction either by the lowest torus eigenvalue or
at fixed total scalar charge.  The first estimate deteriorates with volume,
and the second leaves an \(N^3\) amplitude stock.  This version gives the
exact finite-lattice amplitude criterion: a Newton--Higgs matrix must be
copositive on the cone of squared field amplitudes.  It then proves a
volume-uniform sufficient condition when the sourced density has an exact
local neutral-cell flux representation.

\subsection{Discrete Newton variational principle}

Let \(G=(V,E)\) be a finite connected graph with oriented incidence
operator
\begin{equation}
 d:C^0(V;\mathbb R)\longrightarrow C^1(E;\mathbb R),
 \qquad \Delta:=d^*d.
 \label{eq:v6v63-graph-Laplacian}
\end{equation}
Let \(P_0\) be orthogonal projection onto the mean-zero vertex functions and
let \(\Delta^+\) be the inverse of \(\Delta\) on that subspace, extended by
zero on constants.  Define
\begin{equation}
 \mathcal D_G(\rho)
 :=\langle P_0\rho,\Delta^+P_0\rho\rangle.
 \label{eq:v6v63-discrete-Newton}
\end{equation}

\begin{theorem}[Minimum-flux representation of Newton energy]
\label{thm:v6v63-minimum-flux}
For every mean-zero \(\rho\),
\begin{equation}
 \mathcal D_G(\rho)
 =\min\left\{
 \|J\|_2^2:J\in C^1(E;\mathbb R),\ d^*J=\rho
 \right\}.
 \label{eq:v6v63-flux-variational}
\end{equation}
The unique minimizer orthogonal to \(\ker d^*\) is
\begin{equation}
 J_*=d\Delta^+\rho.
 \label{eq:v6v63-minimal-flux}
\end{equation}
\end{theorem}

\begin{proof}
Put \(u=\Delta^+\rho\).  Then
\(d^*du=\Delta u=\rho\), so \(J_*=du\) is admissible, and
\[
 \|J_*\|_2^2=\langle u,\Delta u\rangle
 =\langle\Delta^+\rho,\rho\rangle=\mathcal D_G(\rho).
\]
Every other admissible flux is \(J_*+K\) with \(d^*K=0\).  Since
\(J_*\in\operatorname{im}d=(\ker d^*)^\perp\),
\(\|J_*+K\|^2=\|J_*\|^2+\|K\|^2\).  This proves minimality and uniqueness
in the stated complement.
\end{proof}

\begin{corollary}[Any explicit flux is an upper bound]
\label{cor:v6v63-flux-upper}
If \(d^*J=\rho\), then
\begin{equation}
 \mathcal D_G(\rho)\le\|J\|_2^2.
 \label{eq:v6v63-any-flux}
\end{equation}
Thus a local flux construction can control the Newton form without using
the global spectral gap of \(\Delta\).
\end{corollary}

\begin{proof}
Use \(J\) as a competitor in
\eqref{eq:v6v63-flux-variational}.
\end{proof}

\subsection{Exact copositivity criterion}

Let a scalar multiplet have squared amplitude vector
\begin{equation}
 r_v:=|\psi_v|^2\ge0,
 \qquad r\in\mathbb R_+^{V}.
 \label{eq:v6v63-amplitude-cone}
\end{equation}
The quartic Higgs--Newton part of a finite-cutoff action is
\begin{equation}
 Q_{\lambda,\Gamma}(r)
 :={\lambda\over4}\|r\|_2^2
 -{\Gamma\over2}\langle P_0r,\Delta^+P_0r\rangle
 ={1\over2}r^T\mathbb A_{\lambda,\Gamma}r,
 \label{eq:v6v63-quartic-form}
\end{equation}
where
\begin{equation}
 \mathbb A_{\lambda,\Gamma}
 :={\lambda\over2}I-\Gamma P_0\Delta^+P_0.
 \label{eq:v6v63-A-matrix}
\end{equation}

\begin{definition}[Strict copositivity]
\label{def:v6v63-copositive}
A symmetric matrix \(A\) is copositive if
\(r^TAr\ge0\) for every \(r\ge0\).  It is strictly copositive if
\begin{equation}
 \inf_{\substack{r\ge0\\\|r\|_2=1}}r^TAr>0.
 \label{eq:v6v63-strict-copositive}
\end{equation}
\end{definition}

\begin{theorem}[Exact finite-lattice amplitude criterion]
\label{thm:v6v63-copositivity}
The quartic form \(Q_{\lambda,\Gamma}\) is nonnegative for all scalar
amplitudes if and only if \(\mathbb A_{\lambda,\Gamma}\) is copositive.  It
is coercive on the amplitude cone if and only if the matrix is strictly
copositive.  If copositivity fails, then the complete action is unbounded
below along an amplitude ray whenever all its remaining terms grow at most
linearly in \(r\).
\end{theorem}

\begin{proof}
Every scalar field produces \(r\ge0\), and every vector in
\(\mathbb R_+^V\) is realized by taking one real scalar component
\(\psi_v=\sqrt{r_v}\).  The first two statements are therefore exactly
Definition~\ref{def:v6v63-copositive} applied to
\eqref{eq:v6v63-quartic-form}.  If
\(r^T\mathbb A r<0\), replace \(r\) by \(t r\), equivalently replace
\(\psi\) by \(t^{1/2}\psi\).  The quartic form is
\(t^2r^T\mathbb A r/2\), while mass, hopping, and other terms at most
linear in \(r\) are \(O(t)\).  The negative quadratic power in \(t\)
dominates.
\end{proof}

\begin{corollary}[Spectral criterion is sufficient but not necessary]
\label{cor:v6v63-spectral-sufficient}
If \(\lambda>2\Gamma\lambda_1^{-1}\), where \(\lambda_1\) is the first
positive graph-Laplacian eigenvalue, then
\(\mathbb A_{\lambda,\Gamma}\succ0\) and hence is strictly copositive.
Copositivity may hold even when the full matrix has a negative direction
outside the nonnegative cone, so this spectral condition need not be
necessary.
\end{corollary}

\begin{proof}
On all of \(C^0(V)\),
\(P_0\Delta^+P_0\preceq\lambda_1^{-1}I\).  Insert this in
\eqref{eq:v6v63-A-matrix}.  Positive definiteness implies strict
copositivity.  The last assertion is the distinction between positivity on
all vectors and positivity only on \(r\ge0\).
\end{proof}

\subsection{Local neutral-cell construction}

Partition the vertices into disjoint connected cells \(C\) of uniformly
bounded diameter.  Let \(E_C\) contain internal cell edges and let
\(\Delta_C=d_C^*d_C\) be the cell Neumann Laplacian.  Assume
\begin{equation}
 \lambda_1(\Delta_C)\ge c_{\rm cell}\ell^{-2}
 \label{eq:v6v63-cell-gap}
\end{equation}
for a common cell scale \(\ell\) and constant \(c_{\rm cell}>0\).

\begin{theorem}[Volume-uniform local neutral-cell flux]
\label{thm:v6v63-local-neutral}
Suppose a source \(\rho\) is neutral in every cell:
\begin{equation}
 \sum_{v\in C}\rho_v=0
 \qquad\text{for every cell }C.
 \label{eq:v6v63-cell-neutrality}
\end{equation}
Then there is a flux \(J\), supported only on internal cell edges, such that
\begin{equation}
 d^*J=\rho,
 \qquad
 \|J\|_2^2
 \le {\ell^2\over c_{\rm cell}}\|\rho\|_2^2.
 \label{eq:v6v63-local-flux}
\end{equation}
Consequently,
\begin{equation}
 \mathcal D_G(\rho)
 \le {\ell^2\over c_{\rm cell}}\|\rho\|_2^2,
 \label{eq:v6v63-local-Newton}
\end{equation}
with a constant independent of the number of cells and the total volume.
\end{theorem}

\begin{proof}
For each cell, neutrality makes
\(u_C:=\Delta_C^+\rho_C\) well defined.  Put
\(J_C=d_Cu_C\) and extend it by zero outside the internal cell edges.  Then
\(d^*J_C=\rho_C\) on the cell.  Disjointness of cell supports gives
\(d^*\sum_CJ_C=\rho\) and
\[
 \left\|\sum_CJ_C\right\|_2^2
 =\sum_C\langle\rho_C,\Delta_C^+\rho_C\rangle
 \le\sum_C\lambda_1(\Delta_C)^{-1}\|\rho_C\|_2^2
 \le{\ell^2\over c_{\rm cell}}\|\rho\|_2^2.
\]
Apply Corollary~\ref{cor:v6v63-flux-upper} for
\eqref{eq:v6v63-local-Newton}.
\end{proof}

\begin{corollary}[Volume-uniform quartic margin under local neutrality]
\label{cor:v6v63-local-margin}
Let \(r\ge0\) and let \(P_{\rm cell}r\) subtract the mean of \(r\) within
each cell.  If the gravitational source is
\(\rho=P_{\rm cell}r\), then
\begin{equation}
 {\lambda\over4}\|r\|_2^2
 -{\Gamma\over2}\mathcal D_G(P_{\rm cell}r)
 \ge {1\over4}
 \left(\lambda-{2\Gamma\ell^2\over c_{\rm cell}}\right)
 \|r\|_2^2.
 \label{eq:v6v63-local-quartic-margin}
\end{equation}
Thus \(\lambda>2\Gamma\ell^2/c_{\rm cell}\) is a volume-uniform sufficient
condition at fixed cell scale.
\end{corollary}

\begin{proof}
The cell projection is an \(\ell^2\)-contraction.  Apply
\eqref{eq:v6v63-local-Newton} to
\(\rho=P_{\rm cell}r\) and insert the result into the left side.
\end{proof}

\begin{assessment}[Exact divisibility replaces the global inverse]
\label{ass:v6v63-divisibility}
Theorem~\ref{thm:v6v63-local-neutral} constructs
\(\rho=d^*J\) before applying the inverse Laplacian.  Its minimum-flux
identity cancels the dangerous global \(\lambda_1^{-1}\) without declaring
the Newton field massive.  This is the scalar-constraint analogue of the V52
curvature factorization and the V22 exact-divisibility criterion.
\end{assessment}

\subsection{Why cell neutrality is new physical input}

\begin{proposition}[Global neutrality does not imply the local bound]
\label{prop:v6v63-global-not-local}
On a sequence of tori with increasing diameter, there are mean-zero sources
\(\rho_L\) of unit \(\ell^2\) norm for which
\begin{equation}
 \mathcal D_{G_L}(\rho_L)=\lambda_1(G_L)^{-1}.
 \label{eq:v6v63-low-mode-saturation}
\end{equation}
Hence global mean zero alone cannot replace
\eqref{eq:v6v63-cell-neutrality} in a volume-uniform estimate.
\end{proposition}

\begin{proof}
Choose a normalized eigenfunction of \(\Delta\) with eigenvalue
\(\lambda_1\).  It is orthogonal to constants and therefore mean zero.
Then
\(\langle\rho_L,\Delta^+\rho_L\rangle=\lambda_1^{-1}\).
\end{proof}

\begin{firewall}[Subtracting cell means changes the gravitational source]
\label{fw:v6v63-cell-source}
The physical Hamiltonian constraint couples to the complete energy density
and its global boundary charge.  Replacing that source by
\(P_{\rm cell}r\) is not a gauge choice.  A valid use of
Theorem~\ref{thm:v6v63-local-neutral} requires an exact Ward, constraint, or
multiscale decomposition in which the removed cell means are retained as
coarser source variables and are paid exactly once.  Such a decomposition is
not proved here.
\end{firewall}

\begin{definition}[Multiscale Newton flux card]
\label{def:v6v63-MNFC}
MNFC is an exact hierarchy
\begin{equation}
 \rho=\sum_{j\ge0}d^*J_j+\rho_{\rm global}
 \label{eq:v6v63-MNFC-decomposition}
\end{equation}
in which \(J_j\) is supported on cells of scale \(\ell_j\), cell means are
passed to scale \(j+1\), the global charge is retained as boundary data, and
the weighted flux stock \(\sum_j\|J_j\|_2^2\) is controlled by the positive
matter Hamiltonian without repeated payment.
\end{definition}

\begin{firewall}[Copositivity is necessary only for the stated scalar
amplitude model]
\label{fw:v6v63-copositive-scope}
The matrix \(\mathbb A_{\lambda,\Gamma}\) uses the source
\(|\psi|^2\) and an onsite quartic.  The complete SM energy density contains
derivatives, gauge energy, fermions, Yukawa terms, and constraint
counterterms.  Their full amplitude cone and induced Newton form can differ.
Theorem~\ref{thm:v6v63-copositivity} is exact for
\eqref{eq:v6v63-quartic-form}, not a claim that this matrix is the complete
SM--Einstein quartic Hessian.
\end{firewall}

\begin{theorem}[Sixth-series V63 result]
\label{thm:v6v63-status}
The discrete Newton energy is the minimum squared flux among all solutions
of \(d^*J=\rho\).  For the finite scalar amplitude model, global stability is
equivalent exactly to copositivity of
\((\lambda/2)I-\Gamma P_0\Delta^+P_0\); failure makes the action unbounded
along an amplitude ray.  A source neutral on uniformly bounded cells admits
a local flux and the volume-independent bound
\(\mathcal D_G(\rho)\le\ell^2\|\rho\|^2/c_{\rm cell}\).  Physical use of
this improvement requires MNFC, because discarded cell means are genuine
coarser gravitational sources.
\end{theorem}

\begin{proof}
Use Theorem~\ref{thm:v6v63-minimum-flux},
Theorem~\ref{thm:v6v63-copositivity},
Theorem~\ref{thm:v6v63-local-neutral},
Corollary~\ref{cor:v6v63-local-margin}, and
Proposition~\ref{prop:v6v63-global-not-local}, subject to both firewalls.
\end{proof}

V64 will construct the exact dyadic cell-mean telescope underlying MNFC on
a periodic lattice.  It will prove orthogonality of the detail sources,
identify the one-use coarse charge ledger, and compute whether the resulting
Newton flux stock is finite or logarithmic for a generic matter density.


\clearpage
\section{V64: dyadic source details and the Newton shell ledger}
\label{sec:v6v64}

V63 proves a volume-uniform Newton bound for a source which is neutral on
uniformly bounded cells, but a physical source carries nonzero cell means.
This version passes those means to coarser cells exactly once.  The resulting
martingale differences are orthogonal and locally neutral.  Their local
fluxes give a rigorous multiscale upper bound, while the exact Laplacian
spectral resolution identifies the sharp positive Newton ledger and its
logarithmic borderline.

\subsection{Dyadic conditional expectations}

Let the periodic cubic lattice have side \(2^J\).  For
\(0\le j\le J\), let \(\mathcal P_j\) be the partition into dyadic cubes of
side \(2^j\), and let
\begin{equation}
 E_j:C^0(V;\mathbb R)\longrightarrow C^0(V;\mathbb R)
 \label{eq:v6v64-Ej}
\end{equation}
be conditional expectation with respect to \(\mathcal P_j\): \(E_j\rho\)
is constant on each cube and equals its average there.  Thus \(E_0=I\),
and \(E_J\rho\) is the global mean.  Define the detail
\begin{equation}
 D_j:=E_j-E_{j+1},
 \qquad0\le j<J.
 \label{eq:v6v64-Dj}
\end{equation}

\begin{theorem}[Exact one-use cell-mean telescope]
\label{thm:v6v64-martingale}
For every \(\rho\),
\begin{equation}
 \rho-E_J\rho=\sum_{j=0}^{J-1}D_j\rho.
 \label{eq:v6v64-telescope}
\end{equation}
The detail sources are pairwise orthogonal:
\begin{equation}
 \langle D_j\rho,D_k\rho\rangle=0
 \quad(j\ne k),
 \qquad
 \|\rho-E_J\rho\|_2^2
 =\sum_{j=0}^{J-1}\|D_j\rho\|_2^2.
 \label{eq:v6v64-detail-orthogonality}
\end{equation}
Moreover \(D_j\rho\) has zero sum in every parent cube of side
\(2^{j+1}\).  Each cell mean is passed to the next scale once and appears in
exactly one difference.
\end{theorem}

\begin{proof}
The projections are nested orthogonal projections and satisfy
\(E_jE_k=E_{\max\{j,k\}}\).  Summing
\(E_j-E_{j+1}\) telescopes to \(I-E_J\).  If \(j<k\), then
\[
 D_jD_k=(E_j-E_{j+1})(E_k-E_{k+1})=E_k-E_{k+1}-E_k+E_{k+1}=0.
\]
Self-adjointness gives orthogonality and Pythagoras.  On a parent cube in
\(\mathcal P_{j+1}\), both \(E_j\rho\) and \(E_{j+1}\rho\) have the same
total sum, so their difference is neutral.  The telescoping formula proves
the one-use statement.
\end{proof}

\subsection{Local fluxes for every detail}

Assume the dyadic parent cubes have Neumann graph-Laplacian gap
\begin{equation}
 \lambda_1(\Delta_C)\ge c_d2^{-2(j+1)}
 \label{eq:v6v64-cube-gap}
\end{equation}
for a dimension constant \(c_d>0\), as holds for standard cubic cells.

\begin{theorem}[Dyadic local flux construction]
\label{thm:v6v64-detail-flux}
For every detail \(D_j\rho\), there is an edge flux \(J_j\), supported
inside the parent cubes of \(\mathcal P_{j+1}\), such that
\begin{equation}
 d^*J_j=D_j\rho,
 \qquad
 \|J_j\|_2^2
 \le C_d4^j\|D_j\rho\|_2^2.
 \label{eq:v6v64-detail-flux}
\end{equation}
The total flux \(J=\sum_{j=0}^{J-1}J_j\) satisfies
\begin{equation}
 d^*J=\rho-E_J\rho.
 \label{eq:v6v64-total-divergence}
\end{equation}
\end{theorem}

\begin{proof}
Theorem~\ref{thm:v6v64-martingale} makes \(D_j\rho\) neutral in every
parent cube.  Apply Theorem~\ref{thm:v6v63-local-neutral} to those disjoint
cubes with \(\ell=2^{j+1}\), and absorb the fixed factor four into \(C_d\).
Summing the divergence identities and using
\eqref{eq:v6v64-telescope} proves
\eqref{eq:v6v64-total-divergence}.
\end{proof}

\begin{corollary}[Rigorous local multiscale Newton bounds]
\label{cor:v6v64-local-bounds}
The discrete Newton energy satisfies both
\begin{align}
 \mathcal D_G(\rho)
 &\le C_dJ\sum_{j=0}^{J-1}4^j\|D_j\rho\|_2^2,
 \label{eq:v6v64-l2-scale-bound}\\
 \mathcal D_G(\rho)
 &\le C_d\left(
 \sum_{j=0}^{J-1}2^j\|D_j\rho\|_2
 \right)^2.
 \label{eq:v6v64-l1-scale-bound}
\end{align}
\end{corollary}

\begin{proof}
By Theorem~\ref{thm:v6v63-minimum-flux} and
\eqref{eq:v6v64-total-divergence},
\(\mathcal D_G(\rho)\le\|\sum_jJ_j\|_2^2\).  Cauchy--Schwarz in the
\(J\) scale labels gives
\(\|\sum_jJ_j\|^2\le J\sum_j\|J_j\|^2\), and
\eqref{eq:v6v64-detail-flux} proves the first bound.  The triangle
inequality followed by the same detail bound proves the second.
\end{proof}

\begin{firewall}[Orthogonal sources do not make their local fluxes
orthogonal]
\label{fw:v6v64-flux-orthogonality}
The detail sources are exactly orthogonal in vertex \(\ell^2\), but the
chosen fluxes occupy nested spatial cells and need not be orthogonal in edge
\(\ell^2\).  Dropping the factor \(J\) from
\eqref{eq:v6v64-l2-scale-bound} would require a stable multilevel flux
decomposition, not merely martingale orthogonality.
\end{firewall}

\subsection{Exact spectral Newton ledger}

Let \(Q_j\) be mutually orthogonal spectral projections of the graph
Laplacian on mean-zero functions, chosen so that on
\(\operatorname{ran}Q_j\),
\begin{equation}
 c_-L_j^{-2}I\preceq\Delta\preceq c_+L_j^{-2}I,
 \qquad L_j=2^j,
 \label{eq:v6v64-spectral-shell}
\end{equation}
and \(\sum_jQ_j=P_0\).  Empty shells are allowed.

\begin{theorem}[Positive orthogonal shell identity]
\label{thm:v6v64-spectral-ledger}
For every \(\rho\),
\begin{equation}
 \mathcal D_G(\rho)
 =\sum_j\langle Q_j\rho,\Delta^+Q_j\rho\rangle,
 \label{eq:v6v64-exact-spectral}
\end{equation}
and
\begin{equation}
 {1\over c_+}
 \sum_jL_j^2\|Q_j\rho\|_2^2
 \le\mathcal D_G(\rho)
 \le {1\over c_-}
 \sum_jL_j^2\|Q_j\rho\|_2^2.
 \label{eq:v6v64-spectral-equivalence}
\end{equation}
No cancellation between distinct spectral shells can lower the exact Newton
energy, because every term in \eqref{eq:v6v64-exact-spectral} is
nonnegative.
\end{theorem}

\begin{proof}
The spectral projections commute with \(\Delta^+\) and are mutually
orthogonal.  Insert \(P_0=\sum_jQ_j\) on both sides of
\eqref{eq:v6v63-discrete-Newton}; all cross terms vanish.  Inverting the
form bounds \eqref{eq:v6v64-spectral-shell} on each shell gives
\(c_+^{-1}L_j^2I\preceq\Delta^+\preceq c_-^{-1}L_j^2I\), proving the two
inequalities.
\end{proof}

\begin{corollary}[The critical shell profile is logarithmic]
\label{cor:v6v64-critical-log}
Suppose a family of sources has \(N\) nonempty shells on which
\begin{equation}
 a\le L_j^2\|Q_j\rho\|_2^2\le b
 \label{eq:v6v64-critical-profile}
\end{equation}
for fixed \(0<a\le b<\infty\).  Then
\begin{equation}
 {aN\over c_+}
 \le\mathcal D_G(\rho)
 \le {bN\over c_-}.
 \label{eq:v6v64-log-ledger}
\end{equation}
Since \(N\) is proportional to the logarithm of the ratio of largest to
smallest represented length scales, the Newton stock is logarithmic at this
critical profile.
\end{corollary}

\begin{proof}
Insert \eqref{eq:v6v64-critical-profile} into
\eqref{eq:v6v64-spectral-equivalence} and count the nonempty shells.
\end{proof}

\begin{assessment}[This logarithm has no signed-shell escape]
\label{ass:v6v64-positive-log}
The logarithms in the audited NS and YM programmes arose after absolute
majorization of quantities with possible cancellation.  The exact Newton
identity \eqref{eq:v6v64-exact-spectral} is already a sum of positive shell
energies.  At the critical profile, cancellation between shells cannot
remove its logarithm.  Closure requires decay stronger than the critical
profile, a positive matter stock which pays the logarithm, or an exact
source factorization which changes the shell weights.
\end{assessment}

\subsection{A scale criterion for the Newton source}

\begin{definition}[Newton spectral stock]
\label{def:v6v64-Newton-stock}
For a source family \(\rho_n\), define
\begin{equation}
 \mathfrak N(\rho_n)
 :=\sum_jL_j^2\|Q_j\rho_n\|_2^2.
 \label{eq:v6v64-Newton-stock}
\end{equation}
The Newton spectral stock card is
\begin{equation}
 \sup_n\mathfrak N(\rho_n)<\infty
 \label{eq:v6v64-NSSC-spectral}
\end{equation}
after all exact Ward and constraint decompositions of the complete SM energy
density have been made.
\end{definition}

\begin{proposition}[Power-law shell threshold]
\label{prop:v6v64-power-threshold}
If for large \(j\)
\begin{equation}
 \|Q_j\rho\|_2^2\le CL_j^{-2-\epsilon},
 \label{eq:v6v64-shell-decay}
\end{equation}
then the large-scale Newton stock is summable exactly when
\(\epsilon>0\).  At \(\epsilon=0\), the absolute and exact spectral ledger
is logarithmic unless additional shells vanish.
\end{proposition}

\begin{proof}
The shell term is at most \(CL_j^{-\epsilon}=C2^{-j\epsilon}\), whose
geometric sum converges exactly for \(\epsilon>0\).  At equality, every
nonzero comparable shell makes an order-one positive contribution by
Corollary~\ref{cor:v6v64-critical-log}.
\end{proof}

\begin{definition}[Stable multilevel Newton flux card]
\label{def:v6v64-SMNFC}
SMNFC strengthens MNFC by requiring local fluxes for the cell details with
\begin{equation}
 \left\|\sum_jJ_j\right\|_2^2
 \le C_{\rm stab}\sum_j\|J_j\|_2^2
 \label{eq:v6v64-stable-flux}
\end{equation}
uniformly in the number of scales, together with a positive matter estimate
for \(\sum_j4^j\|D_j\rho\|_2^2\).  It would remove the extra factor \(J\)
in \eqref{eq:v6v64-l2-scale-bound} without altering the physical Newton
shell stock.
\end{definition}

\begin{firewall}[Fourier orthogonality is not local owner factorization]
\label{fw:v6v64-spectral-local}
The spectral projections \(Q_j\) give the exact Newton ledger but are
nonlocal.  They may be retained as physical massless carriers in the V23
hybrid norm; they cannot automatically serve as local probability owners.
Conversely, the dyadic cell details are local but require SMNFC to obtain the
sharp \(\ell^2\) scale sum.  Equating the two decompositions without a
stable comparison would hide the live locality problem.
\end{firewall}

\begin{firewall}[The source is still the complete energy density]
\label{fw:v6v64-complete-source}
All shell and cell identities are linear in a supplied scalar source.  The
gravitational source must be the complete renormalized SM energy density,
including its boundary and zero-mode charges.  No uniform bound on its
Newton spectral stock has been proved here.
\end{firewall}

\begin{theorem}[Sixth-series V64 result]
\label{thm:v6v64-status}
Every finite-lattice source decomposes exactly into orthogonal dyadic cell
details, each neutral on its parent and paid once.  Local cell solutions give
the rigorous multiscale bounds
\eqref{eq:v6v64-l2-scale-bound}--\eqref{eq:v6v64-l1-scale-bound}; removal of
their extra scale loss requires SMNFC.  Independently, Laplacian spectral
shells give the exact positive Newton ledger
\eqref{eq:v6v64-exact-spectral}, equivalent to
\(\sum_jL_j^2\|Q_j\rho\|^2\).  An order-one contribution on every active
shell produces an unavoidable logarithm.  The complete SM energy-density
stock remains open.
\end{theorem}

\begin{proof}
Use Theorem~\ref{thm:v6v64-martingale},
Theorem~\ref{thm:v6v64-detail-flux},
Corollary~\ref{cor:v6v64-local-bounds},
Theorem~\ref{thm:v6v64-spectral-ledger}, and
Proposition~\ref{prop:v6v64-power-threshold}, subject to the three
firewalls.
\end{proof}

V65 will perform the compulsory YM/NS audit.  It will compare their newest
complete-object scale ledgers with the positive Newton shell identity and
decide whether the next attack should seek a stable local flux decomposition
or a positive matter estimate for the complete energy-density spectral
stock.


\clearpage
\section{V65: companion audit and a free composite-source Newton logarithm}
\label{sec:v6v65}

V64 identifies the exact positive Newton source stock
\(\sum_jL_j^2\|Q_j\rho\|_2^2\).  This compulsory audit shows that the
source must be renormalized as a complete composite object before an
all-scale bound is attempted.  To test that gate, this version computes the
Newton stock of the Wick square of a free massive scalar in three spatial
dimensions.  Its expected stock diverges logarithmically in the ultraviolet,
even though Wick ordering has removed the mean.

\subsection{Compulsory V65 audit}

The newest companion sources inspected were
\begin{center}
\begin{tabular}{p{0.25\linewidth}p{0.16\linewidth}p{0.47\linewidth}}
programme & source & audited mathematical endpoint\\ \hline
Yang--Mills & \texttt{V44} & a squarefree Banach-algebra source logarithm
gives exact tensor cumulants, a uniform source-invertibility simplex, and
the correct binary matrix cards; direct all-arity Cauchy bounds miss exactly
the higher Gaussian cancellation, leaving the Gaussian-renormalized matrix
source card GRMSC open\\
Navier--Stokes & \texttt{V42} & the type-I problem is reduced to a Liouville
theorem for a bounded-enstrophy ancient class; finite-energy ancient
solutions vanish, but the blow-up limit has infinite energy and the
stationary subclass contains Leray's classical stationary Liouville gate
\end{tabular}
\end{center}
At the time of reading, YM V44 had \(21\,215\) logical lines,
\(669\,788\) bytes, and SHA--256
\[
 \texttt{203ae2186b394fa08726969c85d96fde2e8497e6eb72126d7e3724129c9cdf9a},
\]
while NS V42 had \(28\,013\) logical lines, \(1\,122\,447\) bytes, and
SHA--256
\[
 \texttt{cfe70f62c7d20cf8d0f3b0d0144899f8deb082bf2689f15a46ea09647e4fc1b5}.
\]

\begin{theorem}[V65 transfer decision]
\label{thm:v6v65-audit-decision}
The audit yields the following acquisition order.
\begin{enumerate}[label=\textup{(\roman*)},leftmargin=3.7em]
\item The gravitational source must be defined as a complete renormalized
      stress-energy composite before its Newton shell stock is bounded.
      Termwise Wick monomials or diagram fragments cannot be assigned
      independent all-order owner norms when their Gaussian contractions
      cancel only in the complete source logarithm.
\item A stable local flux decomposition cannot repair an infinite exact
      spectral stock.  Source renormalization and its scale law precede the
      SMNFC locality refinement of V64.
\item If the complete source problem reaches a known open continuum gate,
      the note must identify it and move upstream rather than rename it as a
      conditional card.  This is the relevant discipline of the NS V42
      claim boundary.
\item Neither companion supplies a renormalized SM stress tensor coupled to
      nonlinear constraints.
\end{enumerate}
\end{theorem}

\begin{proof}
Item \textup{(i)} is the complete-logarithm requirement of the YM V44 GRMSC
gate applied only as an order-of-operations rule.  Item \textup{(ii)} follows
because V64's spectral identity is positive and exact.  Item \textup{(iii)}
is the explicit NS V42 reduction to a Liouville problem containing its
stationary open subclass.  Item \textup{(iv)} is the companions' stated
claim boundary.
\end{proof}

\begin{firewall}[The audit does not identify the three source theories]
\label{fw:v6v65-audit-types}
YM representation cumulants, NS ancient solutions, and gravitational
stress-energy composites are different mathematical objects.  No estimate
is transferred.  Only the demonstrated order of cancellation and the
practice of naming an upstream open gate are used.
\end{firewall}

\subsection{Massive free scalar Wick square}

Let \(\phi\) be a centered translation-invariant real Gaussian field on
\(\mathbb R^3\) with Fourier covariance
\begin{equation}
 \mathbb E\bigl[\widehat\phi(p)
 \overline{\widehat\phi(q)}\bigr]
 =(2\pi)^3\delta(p-q){1\over p^2+m^2},
 \qquad m>0.
 \label{eq:v6v65-scalar-covariance}
\end{equation}
Define the Wick square \(\rho(x)=:\!\phi(x)^2\!:\) by any standard smooth
ultraviolet approximation followed by subtraction of its mean.  Its
two-point function is independent of the chosen approximation limit as a
tempered distribution.

\begin{lemma}[Exact Wick-square spectrum]
\label{lem:v6v65-Wick-spectrum}
The Fourier covariance density of \(\rho\) is
\begin{equation}
 \mathbb E\bigl[
 \widehat\rho(k)\overline{\widehat\rho(k')}
 \bigr]
 =(2\pi)^3\delta(k-k')\,2B_m(k),
 \label{eq:v6v65-rho-covariance}
\end{equation}
where
\begin{equation}
 B_m(k):=int_{\mathbb R^3}{d^3p\over(2\pi)^3}
 {1\over(p^2+m^2)((p+k)^2+m^2)}.
 \label{eq:v6v65-Bm}
\end{equation}
\end{lemma}

\begin{proof}
For smooth cutoffs, Wick's theorem gives
\(\mathbb E[:\!\phi^2\!:(x):\!\phi^2\!:(y)]=2C_m(x-y)^2\).  Fourier
transformation turns the product into the convolution
\eqref{eq:v6v65-Bm}.  The integral is finite in three dimensions: at large
\(|p|\) its radial majorant is \(Cp^{-2}dp\), and the mass removes all
infrared singularities.  Dominated convergence removes the smooth cutoff.
\end{proof}

\subsection{Exact three-dimensional bubble}

\begin{theorem}[Closed form of the massive scalar bubble]
\label{thm:v6v65-bubble}
For \(k\ne0\),
\begin{equation}
 B_m(k)={1\over4\pi|k|}
 \arctan{|k|\over2m},
 \label{eq:v6v65-bubble-formula}
\end{equation}
and the continuous value at zero is
\begin{equation}
 B_m(0)={1\over8\pi m}.
 \label{eq:v6v65-bubble-zero}
\end{equation}
\end{theorem}

\begin{proof}
Use the Feynman parameter identity
\[
 {1\over ab}=\int_0^1{dx\over(xa+(1-x)b)^2}
\]
and shift the integration variable.  Then
\begin{align*}
 B_m(k)
 &=\int_0^1dx\int{d^3q\over(2\pi)^3}
 {1\over(q^2+m^2+x(1-x)|k|^2)^2}\\
 &={1\over8\pi}\int_0^1
 {dx\over\sqrt{m^2+x(1-x)|k|^2}}.
\end{align*}
The radial integral in the second line follows from
\(4\pi(2\pi)^{-3}\int_0^\infty r^2(r^2+M^2)^{-2}dr=1/(8\pi M)\).
Put \(x=(1+t)/2\); the remaining elementary integral is
\(2|k|^{-1}\arctan(|k|/(2m))\), proving
\eqref{eq:v6v65-bubble-formula}.  Its limit at zero is
\eqref{eq:v6v65-bubble-zero}.
\end{proof}

\begin{corollary}[Infrared regularity and ultraviolet scale]
\label{cor:v6v65-bubble-asymptotics}
For fixed \(m>0\),
\begin{equation}
 B_m(k)\longrightarrow{1\over8\pi m}
 \quad(|k|\downarrow0),
 \qquad
 |k|B_m(k)\longrightarrow{1\over8}
 \quad(|k|\to\infty).
 \label{eq:v6v65-bubble-asymptotics}
\end{equation}
\end{corollary}

\begin{proof}
Use \(\arctan z\sim z\) at zero and
\(\arctan z\to\pi/2\) at infinity in
\eqref{eq:v6v65-bubble-formula}.
\end{proof}

\subsection{Expected Newton stock}

Let \(P_{\le\Lambda}\) cut off only the external source momentum.  Define
the Newton stock per unit volume by
\begin{equation}
 \mathfrak N_{\Lambda,m}
 :=\int_{0<|k|\le\Lambda}{d^3k\over(2\pi)^3}
 {2B_m(k)\over|k|^2}.
 \label{eq:v6v65-Newton-stock}
\end{equation}
This is the translation-invariant density of
\(\mathbb E\langle P_{\le\Lambda}\rho,
(-\Delta)^{-1}P_{\le\Lambda}\rho\rangle\).

\begin{theorem}[Logarithmic composite-source Newton divergence]
\label{thm:v6v65-Newton-log}
For every \(m>0\),
\begin{equation}
 \mathfrak N_{\Lambda,m}
 ={1\over4\pi^3}
 \int_0^\Lambda{1\over r}
 \arctan{r\over2m}\,dr,
 \label{eq:v6v65-Newton-integral}
\end{equation}
and
\begin{equation}
 \mathfrak N_{\Lambda,m}
 ={1\over8\pi^2}\log{\Lambda\over m}+O(1)
 \qquad(\Lambda/m\to\infty).
 \label{eq:v6v65-Newton-log}
\end{equation}
In particular, Wick subtraction of the source mean does not make its
ultraviolet Newton spectral stock finite.
\end{theorem}

\begin{proof}
Use spherical coordinates in \eqref{eq:v6v65-Newton-stock}.  The radial
measure is \(4\pi r^2dr/(2\pi)^3\), which cancels the Newton factor
\(r^{-2}\).  Inserting \eqref{eq:v6v65-bubble-formula} gives
\eqref{eq:v6v65-Newton-integral}.  Write
\[
 \arctan{r\over2m}={\pi\over2}
 +\left(\arctan{r\over2m}-{\pi\over2}\right).
\]
The first term contributes
\((8\pi^2)^{-1}\log(\Lambda/m)\) above any fixed multiple of \(m\).  The
second term is \(O(m/r)\) for \(r\ge m\), so its product with \(dr/r\) is
integrable at infinity.  The integral below \(r=m\) is finite because
\(r^{-1}\arctan(r/(2m))\) is bounded.  This proves
\eqref{eq:v6v65-Newton-log}.  Wick ordering removes the one-point mean, while
the calculation uses the connected two-point function, so the last
statement follows.
\end{proof}

\begin{corollary}[Exact critical shell profile]
\label{cor:v6v65-critical-shell}
On every dyadic external momentum shell
\(R\le|k|\le2R\) with \(R\gg m\), the contribution to
\(\mathfrak N_{\Lambda,m}\) tends to
\begin{equation}
 {\log2\over8\pi^2}
 \label{eq:v6v65-shell-constant}
\end{equation}
as \(R/m\to\infty\).  Thus this free composite source realizes the critical
order-one-per-shell profile of Corollary~\ref{cor:v6v64-critical-log}.
\end{corollary}

\begin{proof}
Restrict \eqref{eq:v6v65-Newton-integral} to \([R,2R]\) and use dominated
convergence after scaling \(r=Ru\):
\(\arctan(Ru/(2m))\to\pi/2\) uniformly for \(1\le u\le2\).
\end{proof}

\subsection{Renormalized source card}

\begin{assessment}[The divergence is composite, not a failure of the
Newton identity]
\label{ass:v6v65-composite-divergence}
The Newton multiplier is exactly the physical \(|k|^{-2}\) carrier retained
in V22--V64.  The logarithm arises because the Wick-square covariance is
\(2B_m(k)\sim1/(4|k|)\).  A stable local flux decomposition cannot change
the positive expectation of the exact spectral stock.  The source composite
and its gravitational counterterms must be renormalized before a uniform
NSSC estimate is possible.
\end{assessment}

\begin{definition}[Renormalized Newton composite card]
\label{def:v6v65-RNCC}
RNCC requires a cutoff family of complete stress-energy sources
\(T_{00}^{(\Lambda)}\) and local gravitational counterterms such that:
\begin{enumerate}[label=\textup{(RN\arabic*)},leftmargin=4.8em]
\item the full diffeomorphism Ward identity and constraint source are defined
      before constituent Wick monomials are normed separately;
\item the local divergent part of the connected source logarithm is
      subtracted in a common regulator scheme compatible with the Einstein
      and higher-curvature couplings;
\item the renormalized Newton spectral stock of every rooted source
      insertion is uniformly bounded in the required hybrid topology; and
\item the subtraction preserves the positive constrained transfer and the
      V14 gauge--matter--metric interface.
\end{enumerate}
\end{definition}

\begin{proposition}[Mean subtraction is insufficient for RNCC]
\label{prop:v6v65-Wick-insufficient}
For the free scalar test, ordinary Wick ordering proves
\(\mathbb E\rho=0\) but RNCC\textup{(RN3)} fails by
\eqref{eq:v6v65-Newton-log}.  Any successful RNCC must therefore include a
two-source gravitational counterterm or an equivalent complete connected
subtraction, not only a one-point cosmological-constant subtraction.
\end{proposition}

\begin{proof}
Wick ordering is the subtraction in Lemma~\ref{lem:v6v65-Wick-spectrum}.
Theorem~\ref{thm:v6v65-Newton-log} computes the remaining connected Newton
stock and proves it divergent.
\end{proof}

\begin{firewall}[The scalar Wick square is not the SM stress tensor]
\label{fw:v6v65-test-scope}
The model source is \(:\!\phi^2\!:\), not the conserved improved stress
tensor.  Derivative terms, improvement terms, fermion and gauge sources,
ghosts, and complete diffeomorphism Ward identities can change tensor
numerators and local counterterms.  The theorem proves a sharp obstruction
to a naive source-stock claim and identifies the required renormalization
level; it does not compute the full SM gravitational beta functions.
\end{firewall}

\begin{firewall}[Subtracting a positive expectation is not automatically a
positive transfer]
\label{fw:v6v65-positive-subtraction}
A counterterm which renders a connected composite covariance finite can
alter the metric kinetic and higher-curvature blocks.  RNCC requires the
complete renormalized transfer to remain positive on the constraint quotient.
Finite covariance by itself does not prove this sign or reflection
positivity.
\end{firewall}

\begin{theorem}[Sixth-series V65 result]
\label{thm:v6v65-status}
The compulsory audit places complete composite renormalization before local
Newton flux refinement.  For a free massive scalar in three spatial
dimensions, the Wick-square spectrum is exactly
\(2B_m(k)\), with
\(B_m(k)=(4\pi|k|)^{-1}\arctan(|k|/(2m))\).  Its expected Newton stock is
\((8\pi^2)^{-1}\log(\Lambda/m)+O(1)\), with one positive constant payment
per ultraviolet shell.  Wick mean subtraction is therefore insufficient;
RNCC is the next upstream gravitational matter gate.
\end{theorem}

\begin{proof}
Use Theorem~\ref{thm:v6v65-audit-decision},
Lemma~\ref{lem:v6v65-Wick-spectrum},
Theorem~\ref{thm:v6v65-bubble},
Theorem~\ref{thm:v6v65-Newton-log}, and
Corollary~\ref{cor:v6v65-critical-shell}, subject to the three firewalls.
\end{proof}

V66 will test whether the conserved, improved free-scalar stress tensor
changes this logarithmic conclusion.  It will keep the complete tensor
numerator, prove its Ward transversality, contract it with the scalar
constraint projector, and separate removable local contact terms from the
physical nonlocal Newton kernel.

\clearpage
\section{V66: the conserved scalar stress tensor and the nonlocal Newton witness}
\label{sec:v6v66}

V65 proved a logarithmic Newton-stock divergence for the lower-dimensional
test composite \(:\!\phi^2\!:\).  The physical scalar source is instead the
energy density of a conserved stress tensor.  This version keeps its complete
two-particle numerator, including an arbitrary improvement coefficient.  The
Ward identity removes the spatial zero mode, but a positive two-particle
submeasure grows as \(|k|^5\) at high external momentum.  Consequently the
instantaneous Newton pairing of the unsmeared energy density has a
sixth-power ultraviolet obstruction.  Purely local gravitational contact
terms cannot cancel that off-vacuum spectral submeasure.

\subsection{Improved tensor and its exact on-shell Ward identity}

Work in four-dimensional Minkowski space with
\(\eta=\operatorname{diag}(1,-1,-1,-1)\).  Let \(m\ge0\), and normal order
with respect to the free scalar vacuum.  For \(\xi\in\mathbb R\), put
\begin{equation}
 T^{(\xi)}_{\mu\nu}
 =:\!\partial_\mu\phi\,\partial_\nu\phi
 -{1\over2}\eta_{\mu\nu}
   \bigl(\partial_\alpha\phi\,\partial^\alpha\phi-m^2\phi^2\bigr)
 +\xi(\eta_{\mu\nu}\Box-\partial_\mu\partial_\nu)\phi^2\!:\,.
 \label{eq:v6v66-improved-stress}
\end{equation}
The last term is identically conserved.  The equation of motion gives
\(\partial^\mu T^{(\xi)}_{\mu\nu}=0\) on the finite-particle core.

For on-shell future momenta
\begin{equation}
 p=(\omega_p,\mathbf p),\qquad
 q=(\omega_q,\mathbf q),\qquad
 \omega_p=(|\mathbf p|^2+m^2)^{1/2},
 \label{eq:v6v66-onshell-momenta}
\end{equation}
write \(K=p+q\).  After removal of the common nonzero normalization and an
irrelevant overall sign, the vacuum-to-two-particle matrix element of
\eqref{eq:v6v66-improved-stress} is
\begin{equation}
 V^{(\xi)}_{\mu\nu}(p,q)
 =p_\mu q_\nu+p_\nu q_\mu
  -\eta_{\mu\nu}(p\mathbin\cdot q+m^2)
  +2\xi(\eta_{\mu\nu}K^2-K_\mu K_\nu).
 \label{eq:v6v66-stress-vertex}
\end{equation}

\begin{theorem}[Exact two-particle stress Ward identity]
\label{thm:v6v66-stress-Ward}
For every \(m\ge0\), \(\xi\in\mathbb R\), and on-shell pair \(p,q\),
\begin{equation}
 K^\mu V^{(\xi)}_{\mu\nu}(p,q)=0.
 \label{eq:v6v66-vertex-Ward}
\end{equation}
In particular, the improvement term is transverse by itself.
\end{theorem}

\begin{proof}
Since \(p^2=q^2=m^2\), contraction of the first three terms in
\eqref{eq:v6v66-stress-vertex} gives
\[
 (m^2+p\mathbin\cdot q)q_\nu
 +(m^2+p\mathbin\cdot q)p_\nu
 -(p+q)_\nu(p\mathbin\cdot q+m^2)=0.
\]
The contraction of the improvement term is
\(2\xi(K_\nu K^2-K^2K_\nu)=0\).
\end{proof}

Take \(\mathbf q=\mathbf k-\mathbf p\).  The energy-density component of
the reduced vertex is
\begin{equation}
 F_{\xi,m}(\mathbf p,\mathbf k)
 :=V^{(\xi)}_{00}(p,q)
 =\omega_p\omega_q+\mathbf p\mathbin\cdot\mathbf q-m^2
  -2\xi|\mathbf k|^2.
 \label{eq:v6v66-energy-numerator}
\end{equation}

\begin{corollary}[The Ward identity removes only the total-energy mode]
\label{cor:v6v66-zero-mode}
For every \(\mathbf p\),
\begin{equation}
 F_{\xi,m}(\mathbf p,0)=0.
 \label{eq:v6v66-zero-energy-numerator}
\end{equation}
For every fixed \(\xi\), however,
\(F_{\xi,0}(\mathbf p,\mathbf k)\) is nonzero on a set of positive
\(\mathbf p\)-measure whenever \(\mathbf k\ne0\).
\end{corollary}

\begin{proof}
At \(\mathbf k=0\), one has \(\mathbf q=-\mathbf p\) and
\(\omega_q=\omega_p\), so the first three terms in
\eqref{eq:v6v66-energy-numerator} equal
\(\omega_p^2-|\mathbf p|^2-m^2=0\), while the improvement term vanishes.
For the second claim take \(\mathbf p=t\mathbf k\), \(0<t<1\).  Then
\[
 |\mathbf k|^{-2}F_{\xi,0}(t\mathbf k,\mathbf k)
 =2t(1-t)-2\xi,
\]
which is not identically zero as a function of \(t\).  Continuity supplies
an open set on which it is nonzero.
\end{proof}

Thus conservation correctly expresses the fact that the normal-ordered
total Hamiltonian does not create particles from the vacuum.  It does not
annihilate a local energy insertion at nonzero spatial momentum.

\subsection{A positive scheme-independent spectral witness}

With the standard one-particle normalization, the two-particle part of the
cutoff equal-time energy spectrum is
\begin{equation}
 \mathcal S^{(\Lambda)}_{\xi,m}(\mathbf k)
 ={1\over8}\int_{\substack{|\mathbf p|\le\Lambda\\
                    |\mathbf k-\mathbf p|\le\Lambda}}
 {d^3p\over(2\pi)^3}\,
 {F_{\xi,m}(\mathbf p,\mathbf k)^2\over
  \omega_p\omega_{\mathbf k-\mathbf p}}.
 \label{eq:v6v66-cutoff-spectrum}
\end{equation}
The factor \(1/2!\) for identical two-particle states and the two field
normalizations give the displayed factor \(1/8\).  Only positivity and the
momentum dependence will matter below.

Fix a unit vector \(e\) and a sufficiently small \(\epsilon>0\), and define
the bounded open relative-momentum cylinder
\begin{equation}
 \mathcal R_e
 =\{u=te+v:\ 1/4<t<3/4,\ v\mathbin\cdot e=0,
                         \ |v|<\epsilon\}.
 \label{eq:v6v66-relative-region}
\end{equation}
It stays a positive distance from both \(u=0\) and \(u=e\).  For
\(\mathbf k\ne0\), take \(e=\mathbf k/|\mathbf k|\) and restrict
\eqref{eq:v6v66-cutoff-spectrum} to
\(\mathbf p=|\mathbf k|u\), \(u\in\mathcal R_e\).  Denote the resulting
positive integral by \(\mathcal W_{\xi,m}(\mathbf k)\).

\begin{lemma}[Exact massless scaling and strict coefficient]
\label{lem:v6v66-witness-scaling}
For every \(\xi\in\mathbb R\),
\begin{equation}
 \mathcal W_{\xi,0}(\mathbf k)=c_\xi|\mathbf k|^5,
 \qquad c_\xi>0,
 \label{eq:v6v66-massless-witness}
\end{equation}
where
\begin{equation}
 c_\xi={1\over8(2\pi)^3}
 \int_{\mathcal R_e}
 {\bigl(|u||e-u|+u\mathbin\cdot(e-u)-2\xi\bigr)^2
  \over |u||e-u|}\,d^3u.
 \label{eq:v6v66-cxi}
\end{equation}
The value is independent of \(e\).
\end{lemma}

\begin{proof}
Substitute \(\mathbf p=|\mathbf k|u\) into
\eqref{eq:v6v66-cutoff-spectrum}.  The measure contributes
\(|\mathbf k|^3\), the squared numerator contributes
\(|\mathbf k|^4\), and the denominator contributes
\(|\mathbf k|^2\), proving the scaling and
\eqref{eq:v6v66-cxi}.  The integrand is nonnegative and locally bounded on
\(\mathcal R_e\).  If its integral vanished, continuity would imply
\[
 |u||e-u|+u\mathbin\cdot(e-u)=2\xi
\]
throughout that region.  On the line \(u=te\), the left side is
\(2t(1-t)\), which is not constant on \((1/4,3/4)\).  Hence
\(c_\xi>0\).  Rotational invariance proves independence of \(e\).
\end{proof}

\begin{theorem}[Massive high-momentum lower bound]
\label{thm:v6v66-massive-witness}
For every fixed \(m\ge0\) and \(\xi\in\mathbb R\), there are constants
\(k_{\xi,m}\ge0\) and \(b_\xi>0\), with \(b_\xi\) independent of \(m\),
such that
\begin{equation}
 \mathcal W_{\xi,m}(\mathbf k)\ge b_\xi|\mathbf k|^5
 \qquad (|\mathbf k|\ge k_{\xi,m}).
 \label{eq:v6v66-massive-lower}
\end{equation}
The witness is included in \(\mathcal S^{(\Lambda)}_{\xi,m}\) whenever
\(\Lambda\ge C_{\mathcal R}|\mathbf k|\), for a fixed geometric constant
\(C_{\mathcal R}\).
\end{theorem}

\begin{proof}
Put \(\mu=m/|\mathbf k|\).  After the same scaling as in
Lemma~\ref{lem:v6v66-witness-scaling}, the coefficient of
\(|\mathbf k|^5\) is
\begin{equation}
 c_\xi(\mu)={1\over8(2\pi)^3}
 \int_{\mathcal R_e}
 {G_{\xi,\mu}(u)^2\over
  (|u|^2+\mu^2)^{1/2}(|e-u|^2+\mu^2)^{1/2}}\,d^3u,
 \label{eq:v6v66-massive-coefficient}
\end{equation}
where
\[
 G_{\xi,\mu}(u)
 =( |u|^2+\mu^2)^{1/2}(|e-u|^2+\mu^2)^{1/2}
   +u\mathbin\cdot(e-u)-\mu^2-2\xi.
\]
The region is bounded away from both denominator singularities.  Dominated
convergence therefore gives \(c_\xi(\mu)\to c_\xi(0)=c_\xi>0\).  Choose
\(\mu_\xi>0\) so that \(c_\xi(\mu)\ge c_\xi/2\) for
\(0\le\mu\le\mu_\xi\), and set
\(b_\xi=c_\xi/2\) and
\(k_{\xi,m}=m/\mu_\xi\).  Boundedness of \(u\) and \(e-u\) on
\(\mathcal R_e\) proves the cutoff statement.
\end{proof}

\subsection{The scalar constraint retains the witness}

The normalization of the linearized scalar gravitational variable is not
needed.  Write its quadratic constrained energy in the convention-independent
form
\begin{equation}
 \mathcal E(\psi;\rho)
 ={a\over2}\|\nabla\psi\|_2^2-bG\langle\psi,\rho\rangle,
 \qquad a,b>0.
 \label{eq:v6v66-scalar-constraint-energy}
\end{equation}

\begin{lemma}[Exact Newton completion]
\label{lem:v6v66-Newton-completion}
On the mean-zero spatial sector, eliminating \(\psi\) gives
\begin{equation}
 \inf_\psi\mathcal E(\psi;\rho)
 =-{b^2G^2\over2a}
   \langle\rho,(-\Delta)^{-1}\rho\rangle.
 \label{eq:v6v66-Newton-completion}
\end{equation}
Thus the positive magnitude of the scalar-constraint source stock has
multiplier \(|\mathbf k|^{-2}\).
\end{lemma}

\begin{proof}
The Euler equation is
\(a(-\Delta)\psi=bG\rho\).  Substitute
\(\psi=(bG/a)(-\Delta)^{-1}\rho\) into
\eqref{eq:v6v66-scalar-constraint-energy}.
\end{proof}

\begin{theorem}[Sixth-power nonlocal Newton obstruction]
\label{thm:v6v66-Newton-obstruction}
Let \(L>k_{\xi,m}\), and take the internal cutoff large enough to contain
the relative-momentum region for every \(|\mathbf k|\le L\).  The positive
two-particle witness in the expected Newton source stock obeys
\begin{equation}
 \int_{k_{\xi,m}\le|\mathbf k|\le L}
 {d^3k\over(2\pi)^3}\,
 {\mathcal W_{\xi,m}(\mathbf k)\over|\mathbf k|^2}
 \ge {b_\xi\over12\pi^2}
       \bigl(L^6-k_{\xi,m}^6\bigr).
 \label{eq:v6v66-Newton-sixth}
\end{equation}
In particular, neither exact stress-tensor conservation nor any fixed
improvement coefficient makes the pointlike positive Newton stock uniformly
bounded.
\end{theorem}

\begin{proof}
Use Theorem~\ref{thm:v6v66-massive-witness} and polar coordinates:
\[
 {4\pi\over(2\pi)^3}
 \int_{k_{\xi,m}}^L b_\xi r^5\,dr
 ={b_\xi\over12\pi^2}
   (L^6-k_{\xi,m}^6).
\]
The factor \(r^2\) from the spatial measure is exactly canceled by the
Newton multiplier.  The remaining power is the integral of \(r^5\).
\end{proof}

The divergence in Theorem~\ref{thm:v6v66-Newton-obstruction} is stronger
than the V65 logarithm because a four-dimensional stress tensor contains
two derivatives and has a much larger high-frequency spectral weight.

\subsection{Contact terms cannot remove the positive cut}

The raw covariant stress two-point distribution also contains coincident
terms.  A complete metric differentiation includes the seagull obtained by
differentiating the action twice.  Renormalization may add the flat-space
quadratic variations of
\begin{equation}
 \int\sqrt{|g|}\,
 (c_0+c_1R+c_2R^2+c_3R_{\mu\nu}R^{\mu\nu}+c_4\Box R),
 \label{eq:v6v66-local-gravity-counterterms}
\end{equation}
as well as the chosen nonminimal scalar term which fixes \(\xi\).

\begin{proposition}[Separation of Ward contacts from the physical cut]
\label{prop:v6v66-contact-separation}
The seagull and the pure-gravity terms in
\eqref{eq:v6v66-local-gravity-counterterms} can change only local
vacuum/contact contributions to the stress two-point function.  They do not
change the vacuum-to-two-scalar matrix element
\eqref{eq:v6v66-stress-vertex}.  The nonminimal scalar improvement changes
that matrix element exactly through the displayed \(\xi\)-term, but for
every \(\xi\) its positive coefficient \(c_\xi\) remains strictly positive.
Consequently no Ward-completing local contact subtraction cancels the
two-particle witness of
Theorem~\ref{thm:v6v66-Newton-obstruction}.
\end{proposition}

\begin{proof}
The second metric variation and the pure-gravity counterterms are supported
at coincident insertions and act as multiples of the identity in the free
matter Hilbert space.  Their matrix elements between the vacuum and a
two-particle state vanish.  Equivalently, local polynomial terms have no
two-particle spectral discontinuity.  The only local curvature term in the
list that changes the one-stress/two-scalar vertex is the nonminimal
\(R\phi^2\) term; its flat first variation is precisely the improvement term
already retained in \eqref{eq:v6v66-stress-vertex}.  Strict positivity for
every value of its coefficient is Lemma~\ref{lem:v6v66-witness-scaling}.
\end{proof}

\begin{corollary}[Ward transversality is not ultraviolet Newton decay]
\label{cor:v6v66-Ward-not-decay}
The exact Ward identity \eqref{eq:v6v66-vertex-Ward} and the exact
zero-mode cancellation \eqref{eq:v6v66-zero-energy-numerator} coexist with
the lower bound \eqref{eq:v6v66-Newton-sixth}.  Therefore an RNCC proof
cannot infer ultraviolet Newton summability from diffeomorphism Ward
transversality.
\end{corollary}

\subsection{The necessary distributional replacement for RNCC}

For a spatial test function \(f\), the witness part of the Newton quadratic
form is
\begin{equation}
 \mathcal Q^{\rm wit}_{\xi,m}(f)
 =\int {d^3k\over(2\pi)^3}\,
  {|\widehat f(\mathbf k)|^2\over|\mathbf k|^2}
  \mathcal W_{\xi,m}(\mathbf k).
 \label{eq:v6v66-smeared-Newton-form}
\end{equation}

\begin{proposition}[Necessary Sobolev order]
\label{prop:v6v66-Sobolev-threshold}
If the renormalized Newton source form is to extend continuously to a
spatial Sobolev test space \(H^s(\mathbb R^3)\) while retaining the positive
two-particle spectrum, then necessarily
\begin{equation}
 s\ge {3\over2}.
 \label{eq:v6v66-Sobolev-order}
\end{equation}
This is a necessary free-scalar threshold; the endpoint upper bound is not
asserted here.
\end{proposition}

\begin{proof}
Choose \(\widehat f_R\) supported in \(R\le|\mathbf k|\le2R\), with
\(R\ge k_{\xi,m}\).  Theorem~\ref{thm:v6v66-massive-witness} gives
\[
 \mathcal Q^{\rm wit}_{\xi,m}(f_R)
 \ge c R^3\|\widehat f_R\|_2^2,
 \qquad
 \|f_R\|_{H^s}^2\asymp R^{2s}\|\widehat f_R\|_2^2.
\]
If \(s<3/2\), their ratio grows as \(R^{3-2s}\), contradicting continuity.
\end{proof}

\begin{definition}[Distributional renormalized Newton composite card]
\label{def:v6v66-dRNCC}
The distributional RNCC, abbreviated dRNCC, replaces the point-source
reading of RNCC\textup{(RN3)} by the following requirements:
\begin{enumerate}[label=\textup{(dRN\arabic*)},leftmargin=5.2em]
\item assemble the complete stress tensor, seagulls, constraints, and Ward
      contacts in one covariant regulator before subtraction;
\item absorb only local contact distributions into the cosmological,
      Einstein, nonminimal, and higher-curvature couplings;
\item retain the positive nonlocal spectral measure and prove continuity of
      the Newton source form on an explicitly declared test topology, whose
      spatial order is at least \(H^{3/2}\) in the free scalar sector; and
\item prove compatibility of that topology and subtraction with the
      reflected transfer, boundary source atlas, and metric constraint
      quotient.
\end{enumerate}
\end{definition}

\begin{assessment}[Upstream correction of the Newton-source target]
\label{ass:v6v66-upstream-correction}
RNCC remains correct if ``uniformly bounded in the required hybrid
topology'' is read distributionally.  It is false if read as a uniform
positive stock for an unsmeared pointlike \(T_{00}\).  V64's exact Newton
shell identity remains valid for each regulated source, but the continuum
source lives in a negative-regularity space and must be paired with test
data before its stock is finite.
\end{assessment}

\begin{firewall}[No cancellation claim for the interacting Standard Model]
\label{fw:v6v66-SM-scope}
The proof concerns one free real scalar and its full improved stress
numerator.  It does not compute the gauge, fermion, ghost, anomaly, or
interaction contributions to the Standard Model stress tensor.  It proves
that Ward conservation and scalar improvement alone cannot justify the
previous point-source target.  The complete interacting dRNCC and its
reflection-positive gravitational counterterms remain open.
\end{firewall}

\begin{firewall}[Positive spectral weight is not a positive conformal mode]
\label{fw:v6v66-spectral-conformal}
The witness is positive because it is a sum of squared matter matrix
elements.  The scalar gravitational completion
\eqref{eq:v6v66-Newton-completion} is attractive and has the opposite sign
in the reduced energy.  Neither fact repairs the Euclidean conformal-factor
problem isolated in V58.  Transfer positivity must be proved after the
complete constraint reduction and counterterm choice.
\end{firewall}

\begin{theorem}[Sixth-series V66 result]
\label{thm:v6v66-status}
The complete improved free-scalar two-particle stress vertex is exactly
transverse.  Its energy component vanishes at zero spatial momentum, but for
every improvement coefficient it has a strictly positive scheme-independent
spectral submeasure \(c_\xi|k|^5\) in the massless theory and a comparable
high-momentum submeasure in the massive theory.  Scalar-constraint
projection therefore gives a Newton witness growing at least as \(L^6\).
Local Ward contacts cannot remove it.  A continuum source card must be
distributional, with spatial test order at least \(H^{3/2}\) in this free
sector; dRNCC records the corrected target.
\end{theorem}

\begin{proof}
Use Theorem~\ref{thm:v6v66-stress-Ward},
Corollary~\ref{cor:v6v66-zero-mode},
Lemma~\ref{lem:v6v66-witness-scaling},
Theorem~\ref{thm:v6v66-massive-witness},
Lemma~\ref{lem:v6v66-Newton-completion},
Theorem~\ref{thm:v6v66-Newton-obstruction},
Proposition~\ref{prop:v6v66-contact-separation}, and
Proposition~\ref{prop:v6v66-Sobolev-threshold}, subject to the two
firewalls.
\end{proof}

V67 will move from the equal-time necessary obstruction to a covariantly
smeared free-field construction.  It will seek an explicit spacetime test
space on which the Ward-complete stress two-point distribution, its local
counterterms, and the scalar Newton constraint all converge, and will test
whether that free dRNCC is compatible with Osterwalder--Schrader reflection.

\clearpage
\section{V67: spacetime-smoothed Newton sources and the free reflected dRNCC}
\label{sec:v6v67}

V66 shows that the sharp-time, pointlike energy density has too much
positive spectral weight for an absolute Newton-stock bound.  The corrected
object is a spacetime-smeared stress distribution.  This version proves an
exact quadratic cancellation of the Newton pole, constructs the resulting
free scalar source as a finite Hilbert-space vector, and proves
Osterwalder--Schrader positivity when the Euclidean insertion stays a
positive distance from the reflection plane.  The bound deteriorates with
the sixth power of that distance, matching the V66 obstruction.

\subsection{Exact factorization at the Newton pole}

Retain the notation of V66: \(\mathbf q=\mathbf k-\mathbf p\),
\(\omega_p=(|\mathbf p|^2+m^2)^{1/2}\), and
\(F_{\xi,m}=V^{(\xi)}_{00}\).

\begin{lemma}[Exact energy-numerator factorization]
\label{lem:v6v67-energy-factorization}
For all \(m\ge0\), \(\xi\in\mathbb R\), and
\(\mathbf p,\mathbf k\in\mathbb R^3\),
\begin{equation}
 F_{\xi,m}(\mathbf p,\mathbf k)
 =\left({1\over2}-2\xi\right)|\mathbf k|^2
  -{1\over2}(\omega_p-\omega_q)^2.
 \label{eq:v6v67-energy-factorization}
\end{equation}
Consequently
\begin{equation}
 |F_{\xi,m}(\mathbf p,\mathbf k)|
 \le C_\xi|\mathbf k|^2,
 \qquad
 C_\xi=\left|{1\over2}-2\xi\right|+{1\over2}.
 \label{eq:v6v67-quadratic-zero}
\end{equation}
\end{lemma}

\begin{proof}
Using \(\omega_p^2=|\mathbf p|^2+m^2\) and the corresponding identity for
\(q\),
\begin{align*}
 |\mathbf k|^2-(\omega_p-\omega_q)^2
 &=|\mathbf p+\mathbf q|^2
   -\omega_p^2-\omega_q^2+2\omega_p\omega_q\\
 &=2(\omega_p\omega_q+\mathbf p\mathbin\cdot\mathbf q-m^2).
\end{align*}
Insert this in \eqref{eq:v6v66-energy-numerator}.  The dispersion relation
is one-Lipschitz, and the reverse triangle inequality gives
\[
 |\omega_p-\omega_q|
 \le\bigl||\mathbf p|-|\mathbf q|\bigr|
 \le|\mathbf p+\mathbf q|=|\mathbf k|.
\]
The displayed bound follows.
\end{proof}

\begin{corollary}[The complete vertex cancels the Newton singularity]
\label{cor:v6v67-Newton-pole-cancel}
With the value at \(\mathbf k=0\) defined by continuity,
\begin{equation}
 {F_{\xi,m}(\mathbf p,\mathbf k)^2\over|\mathbf k|^2}
 \le C_\xi^2|\mathbf k|^2.
 \label{eq:v6v67-Newton-pole-bound}
\end{equation}
The cancellation holds before integration and before an absolute value is
taken.  It uses the complete kinetic, gradient, mass, and improvement
numerator.
\end{corollary}

This infrared cancellation does not contradict V66.  In the scaling region
\(|\mathbf p|\asymp|\mathbf k|\), the right side of
\eqref{eq:v6v67-Newton-pole-bound} still has order \(|\mathbf k|^2\), and
the three-dimensional relative-momentum volume supplies the remaining
\(|\mathbf k|^3\).

\subsection{The Lorentzian spacetime-smeared source}

For \(f\in\mathcal S(\mathbb R)\) and
\(g\in\mathcal S(\mathbb R^3)\), define on the finite-particle core
\begin{equation}
 \mathsf X_{\xi,m}(f,g)
 =\int_{\mathbb R}dt\int_{\mathbb R^3}d^3x\,
 f(t)\,[(-\Delta)^{-1/2}g](\mathbf x)
 :\!T^{(\xi)}_{00}(t,\mathbf x)\!:.
 \label{eq:v6v67-smeared-source}
\end{equation}
The multiplier of \((-\Delta)^{-1/2}\) is \(|\mathbf k|^{-1}\); its
value at the single point \(\mathbf k=0\) is immaterial.

\begin{theorem}[Exact positive two-particle norm]
\label{thm:v6v67-positive-norm}
The vector \(\mathsf X_{\xi,m}(f,g)\Omega\) exists in the free Fock space,
and
\begin{align}
 \|\mathsf X_{\xi,m}(f,g)\Omega\|^2
 ={1\over8}\int {d^3p\,d^3q\over(2\pi)^6}\,
 &{|\widehat f(\omega_p+\omega_q)|^2
   |\widehat g(\mathbf p+\mathbf q)|^2
  \over |\mathbf p+\mathbf q|^2\omega_p\omega_q}
 \nonumber\\
 &\times
 |F_{\xi,m}(\mathbf p,\mathbf p+\mathbf q)|^2.
 \label{eq:v6v67-exact-positive-norm}
\end{align}
The apparent spatial zero-mode singularity is removable by
Corollary~\ref{cor:v6v67-Newton-pole-cancel}.
\end{theorem}

\begin{proof}
Expand the normal-ordered quadratic field in creation and annihilation
operators.  Acting on the vacuum leaves only its two-creation part.  Its
reduced tensor coefficient is \eqref{eq:v6v66-stress-vertex}; the time and
space integrations give respectively
\(\widehat f(\omega_p+\omega_q)\) and
\(|\mathbf p+\mathbf q|^{-1}\widehat g(\mathbf p+\mathbf q)\).
The identical-particle factor and field normalizations give \(1/8\), as in
\eqref{eq:v6v66-cutoff-spectrum}.

It remains to prove integrability.  Put
\(\Omega=\omega_p+\omega_q\) and \(\mathbf k=\mathbf p+\mathbf q\).
Then \(|\mathbf k|\le\Omega\), and
Corollary~\ref{cor:v6v67-Newton-pole-cancel} bounds the integrand by
\begin{equation}
 C_\xi^2{\Omega^2\over\omega_p\omega_q}
 |\widehat f(\Omega)|^2|\widehat g(\mathbf k)|^2.
 \label{eq:v6v67-integrable-majorant}
\end{equation}
For every \(N\), the two Fourier transforms are bounded by Schwartz
seminorms times \((1+\Omega)^{-N}\) and
\((1+|\mathbf k|)^{-N}\).  In radial variables,
\(d^3p/\omega_p\) is bounded by a constant times
\((1+|\mathbf p|)d|\mathbf p|\,d\sigma_p\), and similarly for \(q\).
Taking \(N\) sufficiently large makes the resulting six-dimensional
integral finite.  Near \(p=0\) or \(q=0\) in the massless case, the factors
\(d^3p/|\mathbf p|\) and \(d^3q/|\mathbf q|\) are locally integrable.
Truncated two-particle vectors are therefore Cauchy and converge to the
claimed vector and norm.
\end{proof}

\begin{corollary}[Continuity on the Schwartz source space]
\label{cor:v6v67-Schwartz-continuity}
For fixed \(m\ge0\) and \(\xi\), the map
\((f,g)\mapsto\mathsf X_{\xi,m}(f,g)\Omega\) is jointly continuous from
\(\mathcal S(\mathbb R)\times\mathcal S(\mathbb R^3)\) to the free Hilbert
space.
\end{corollary}

\begin{proof}
The majorant in \eqref{eq:v6v67-integrable-majorant} is controlled by a
fixed finite collection of Schwartz seminorms once one chooses any one
sufficiently large \(N\).  Apply dominated convergence to a convergent
sequence in those seminorms.
\end{proof}

\subsection{The nonlocal stress seminorm descends through the Ward quotient}

For a Schwartz tensor \(h^{\mu\nu}\), define the positive two-particle
stress seminorm
\begin{equation}
 \|h\|_{T,2}^2
 ={1\over8}\int {d^3p\,d^3q\over(2\pi)^6}
 {\left|\widehat h^{\mu\nu}(K)
 V^{(\xi)}_{\mu\nu}(p,q)\right|^2\over\omega_p\omega_q},
 \qquad K=(\omega_p+\omega_q,\mathbf p+\mathbf q).
 \label{eq:v6v67-stress-seminorm}
\end{equation}

\begin{theorem}[Exact diffeomorphism quotient of the physical cut]
\label{thm:v6v67-Ward-quotient}
If
\begin{equation}
 h_{\mu\nu}=\partial_\mu a_\nu+\partial_\nu a_\mu
 \label{eq:v6v67-pure-gauge-test}
\end{equation}
for a Schwartz one-form \(a\), then \(\|h\|_{T,2}=0\).  Hence
\eqref{eq:v6v67-stress-seminorm} descends to a positive seminorm on tensor
tests modulo linearized diffeomorphisms.
\end{theorem}

\begin{proof}
In Fourier space, \(\widehat h_{\mu\nu}(K)\) is a constant multiple of
\(K_\mu\widehat a_\nu(K)+K_\nu\widehat a_\mu(K)\).  Symmetry of the vertex
and Theorem~\ref{thm:v6v66-stress-Ward} make its contraction with
\(V^{(\xi)}_{\mu\nu}\) vanish pointwise in \(p,q\).
\end{proof}

\begin{firewall}[The quotient statement concerns the nonlocal cut]
\label{fw:v6v67-Ward-contact-scope}
Theorem~\ref{thm:v6v67-Ward-quotient} is an exact statement about
off-vacuum matrix elements.  A time-ordered or Euclidean differentiated
generating functional also contains coincident seagulls and metric contact
terms.  Those must be assembled before its distributional Ward identity is
stated; they cannot be inferred from the cut seminorm alone.
\end{firewall}

\subsection{Separated-half-space Osterwalder--Schrader form}

Let \(\phi_E\) be the centered Euclidean free scalar with covariance
\(( -\Delta_4+m^2)^{-1}\).  Reflect time by
\(\vartheta(\tau,\mathbf x)=(-\tau,\mathbf x)\).  Tensor reflection includes
one minus sign for each time index:
\begin{equation}
 (\Theta h)_{\mu\nu}(\tau,\mathbf x)
 =(-1)^{\mathbf1_{\mu=0}+\mathbf1_{\nu=0}}
  \overline{h_{\mu\nu}(-\tau,\mathbf x)}.
 \label{eq:v6v67-tensor-reflection}
\end{equation}

Take real \(f\in C_c^\infty((0,\infty))\) with
\(\operatorname{supp}f\subset[\tau_0,\tau_1]\), \(\tau_0>0\), and real
\(g\in\mathcal S(\mathbb R^3)\).  Define
\begin{equation}
 \widetilde f(E)=\int_0^\infty f(\tau)e^{-E\tau}\,d\tau.
 \label{eq:v6v67-Laplace-test}
\end{equation}

\begin{theorem}[Free reflected Newton-source certificate]
\label{thm:v6v67-OS-certificate}
The Euclidean Wick polynomial obtained from
\eqref{eq:v6v67-smeared-source} by analytic continuation is an element of
the positive-time free-field algebra.  Its Osterwalder--Schrader quadratic
form is
\begin{align}
 \mathcal Q^{\rm OS}_{\xi,m}(f,g)
 ={1\over8}\int {d^3p\,d^3q\over(2\pi)^6}\,
 &{|\widetilde f(\omega_p+\omega_q)|^2
   |\widehat g(\mathbf p+\mathbf q)|^2
  \over |\mathbf p+\mathbf q|^2\omega_p\omega_q}
 \nonumber\\
 &\times
 |F_{\xi,m}(\mathbf p,\mathbf p+\mathbf q)|^2
 \ge0.
 \label{eq:v6v67-OS-positive-form}
\end{align}
It is finite for every \(m\ge0\), \(\xi\in\mathbb R\), and such
\(f,g\).
\end{theorem}

\begin{proof}
The free Euclidean Gaussian measure is reflection positive.  The improved
stress tensor is a local Wick polynomial in the field and its derivatives;
\eqref{eq:v6v67-tensor-reflection} is exactly its tensor parity under time
reflection.  Therefore its smearing with data supported at positive time
belongs to the positive-time polynomial algebra, and the OS quadratic form
is nonnegative.

For an explicit calculation, use the standard reflected covariance
decomposition.  Each field crossing the reflection plane contributes
\((2\omega_p)^{-1}e^{-\omega_p(\tau+\sigma)}\).  Wick's theorem leaves the
two cross-pairings, and the complete derivative numerator is
\(F_{\xi,m}\).  The two positive-time integrations give
\(|\widetilde f(\omega_p+\omega_q)|^2\), proving
\eqref{eq:v6v67-OS-positive-form}.  Finally,
\[
 |\widetilde f(E)|\le\|f\|_{L^1}e^{-\tau_0E}.
\]
Combine this with \eqref{eq:v6v67-Newton-pole-bound}; exponential decay in
\(\omega_p+\omega_q\) makes the integral finite, including the massless
infrared region as in Theorem~\ref{thm:v6v67-positive-norm}.
\end{proof}

\begin{proposition}[Quantitative reflection-plane loss]
\label{prop:v6v67-reflection-loss}
For fixed \(m\ge0\) and \(\xi\),
\begin{equation}
 \mathcal Q^{\rm OS}_{\xi,m}(f,g)
 \le C_{\xi,m}\|f\|_{L^1}^2
       \|\widehat g\|_\infty^2(1+\tau_0^{-6}).
 \label{eq:v6v67-tau-loss}
\end{equation}
Thus the separated-half-space certificate is uniform at fixed
\(\tau_0>0\), but its elementary bound has precisely sixth-order
deterioration as the insertion approaches the reflection plane.
\end{proposition}

\begin{proof}
By \eqref{eq:v6v67-Newton-pole-bound} and
\(|\mathbf p+\mathbf q|\le\omega_p+\omega_q\), the integral is bounded by
a constant times
\[
 \|f\|_1^2\|\widehat g\|_\infty^2
 \int {d^3p\,d^3q\over\omega_p\omega_q}
 (\omega_p+\omega_q)^2
 e^{-2\tau_0(\omega_p+\omega_q)}.
\]
Use \((a+b)^2\le2a^2+2b^2\) to factor this into products of radial
one-particle integrals.  Since
\[
 \int_{\mathbb R^3}{d^3p\over\omega_p}
  \omega_p^r e^{-2\tau_0\omega_p}
 \le C_{r,m}(1+\tau_0^{-2-r}),
 \qquad r=0,2,
\]
their products are bounded by \(C_m(1+\tau_0^{-6})\).
\end{proof}

\begin{proposition}[Local contacts vanish between separated halves]
\label{prop:v6v67-separated-contacts}
Let \(C(x,y)\) be any finite sum of derivatives of
\(\delta^{(4)}(x-y)\) with smooth local coefficients.  If one test is
supported in \(\tau\ge\tau_0>0\) and its reflection is supported in
\(\tau\le-\tau_0\), then
\begin{equation}
 C(\Theta h,h)=0.
 \label{eq:v6v67-contact-vanish}
\end{equation}
Hence Ward-completing seagulls and the pure-gravity counterterms of
\eqref{eq:v6v66-local-gravity-counterterms} do not change the separated OS
quadratic form \eqref{eq:v6v67-OS-positive-form}.
\end{proposition}

\begin{proof}
The support of every derivative of \(\delta^{(4)}(x-y)\) is the diagonal
\(x=y\).  The Cartesian product of the two test supports stays a distance
at least \(2\tau_0\) from that diagonal.  The distributional pairing is
therefore zero.
\end{proof}

\subsection{Free dRNCC and the remaining boundary problem}

\begin{definition}[Free separated dRNCC certificate]
\label{def:v6v67-free-dRNCC}
The free separated dRNCC consists of:
\begin{enumerate}[label=\textup{(FdR\arabic*)},leftmargin=5.4em]
\item the complete improved vertex
      \eqref{eq:v6v66-stress-vertex} and its Ward quotient
      \eqref{eq:v6v67-stress-seminorm};
\item the exact pre-norm Newton-pole cancellation
      \eqref{eq:v6v67-Newton-pole-bound};
\item the finite Lorentzian source vector
      \eqref{eq:v6v67-exact-positive-norm};
\item the nonnegative separated OS form
      \eqref{eq:v6v67-OS-positive-form}; and
\item the explicit reflection-plane loss
      \eqref{eq:v6v67-tau-loss}, with all local contacts retained as a
      separate diagonal distribution bank.
\end{enumerate}
\end{definition}

\begin{theorem}[Free separated dRNCC closes]
\label{thm:v6v67-free-dRNCC}
For the free real scalar with arbitrary fixed improvement coefficient, all
five rows of Definition~\ref{def:v6v67-free-dRNCC} hold for Schwartz
spatial data and smooth Euclidean time data supported a positive distance
from the reflection plane.
\end{theorem}

\begin{proof}
Apply Theorem~\ref{thm:v6v66-stress-Ward},
Corollary~\ref{cor:v6v67-Newton-pole-cancel},
Theorem~\ref{thm:v6v67-positive-norm},
Theorem~\ref{thm:v6v67-Ward-quotient},
Theorem~\ref{thm:v6v67-OS-certificate},
Proposition~\ref{prop:v6v67-reflection-loss}, and
Proposition~\ref{prop:v6v67-separated-contacts}.
\end{proof}

\begin{firewall}[No sharp reflection-plane limit]
\label{fw:v6v67-no-sharp-limit}
The bound \eqref{eq:v6v67-tau-loss} is not uniform as
\(\tau_0\downarrow0\).  Local contact terms vanish for separated halves but
return at the reflection plane.  This version proves a positive interior
source domain, not the boundary trace of \(T_{00}\), an instantaneous
Hamiltonian-density operator, or a uniform time-zero Newton form.
\end{firewall}

\begin{firewall}[No dynamical-gravity OS theorem]
\label{fw:v6v67-no-gravity-OS}
Reflection positivity here is inherited from the free scalar Gaussian
measure with the metric treated as an external source.  Integrating the
metric, fixing diffeomorphisms, treating ghosts, and resolving the conformal
direction can alter the transfer form.  The theorem does not establish OS
positivity of interacting Einstein--Standard-Model theory.
\end{firewall}

\begin{theorem}[Sixth-series V67 result]
\label{thm:v6v67-status}
The complete improved scalar energy numerator factors exactly as
\eqref{eq:v6v67-energy-factorization}, cancels the Newton pole before norms,
and defines a finite positive spacetime-smeared source vector.  Its
nonlocal stress seminorm descends through the linearized diffeomorphism
quotient.  On Euclidean tests supported at distance \(\tau_0>0\) from the
reflection plane, the Newton-smeared stress has an explicit nonnegative OS
form; all local contacts vanish between the two halves, and the norm is
bounded by \(C(1+\tau_0^{-6})\).  The free separated dRNCC is therefore
closed.  The sharp boundary trace and the dynamical-gravity transfer remain
open.
\end{theorem}

\begin{proof}
Use Theorem~\ref{thm:v6v67-free-dRNCC} and the two firewalls.
\end{proof}

V68 will attack the reflection-plane boundary rather than assume it.  It
will classify which divergent terms in the \(\tau_0\downarrow0\) expansion
are local boundary contacts, test whether a finite nonlocal remainder
exists on a boundary Sobolev domain, and separate that question from the
negative conformal gravitational block.

\clearpage
\section{V68: reflection-plane asymptotics and the sharp-time source no-go}
\label{sec:v6v68}

V67 constructs a positive Newton-smeared stress source whenever its
Euclidean time support stays a distance \(\tau_0>0\) from the reflection
plane.  This version analyzes the collision with that plane.  The leading
divergence is an explicit local spatial-gradient contact with a strictly
positive coefficient for every improvement parameter.  After subtracting
it, the massless finite part is a homogeneous \(H^{3/2}\) boundary
distribution.  The subtraction does not produce a limit of the original
positive Hilbert-space vectors, so the interior dRNCC cannot be replaced by
an unsmeared time-zero source.

\subsection{The regularized slice kernel}

Fix \(\varepsilon>0\).  Formally take the V67 Euclidean time test to be a
unit point mass at \(\tau=\varepsilon\); equivalently, first use a smooth
approximation and then take its width to zero at fixed \(\varepsilon\).
Its Laplace transform is \(e^{-\varepsilon E}\).  For
\(\mathbf k\ne0\), define
\begin{equation}
 \mathcal K^{(m)}_{\xi,\varepsilon}(\mathbf k)
 ={1\over8}\int_{\mathbb R^3}{d^3p\over(2\pi)^3}\,
 {e^{-2\varepsilon(\omega_p+\omega_q)}
  F_{\xi,m}(\mathbf p,\mathbf k)^2
  \over |\mathbf k|^2\omega_p\omega_q},
 \qquad \mathbf q=\mathbf k-\mathbf p,
 \label{eq:v6v68-slice-kernel}
\end{equation}
and set its value at \(\mathbf k=0\) to zero.  For
\(g\in\mathcal S(\mathbb R^3)\), the reflected norm is
\begin{equation}
 \mathcal Q^{(m)}_{\xi,\varepsilon}(g)
 =\int_{\mathbb R^3}{d^3k\over(2\pi)^3}\,
  |\widehat g(\mathbf k)|^2
  \mathcal K^{(m)}_{\xi,\varepsilon}(\mathbf k)\ge0.
 \label{eq:v6v68-slice-form}
\end{equation}
Finiteness at every \(\varepsilon>0\) follows from
Theorem~\ref{thm:v6v67-OS-certificate}.

Put
\begin{equation}
 a_\xi={1\over2}-2\xi,
 \qquad
 A_\xi=a_\xi^2-{a_\xi\over3}+{1\over20}
       =\left(a_\xi-{1\over6}\right)^2+{1\over45}.
 \label{eq:v6v68-angular-coefficient}
\end{equation}
In particular, \(A_\xi>0\) for every real \(\xi\), and its minimum occurs
at the conformal value \(\xi=1/6\).

\subsection{Universal local divergence}

The exact identity
\begin{equation}
 \omega_p-\omega_q
 ={2\mathbf p\mathbin\cdot\mathbf k-|\mathbf k|^2
   \over\omega_p+\omega_q}
 \label{eq:v6v68-energy-difference}
\end{equation}
and Lemma~\ref{lem:v6v67-energy-factorization} give
\begin{equation}
 F_{\xi,m}
 =a_\xi|\mathbf k|^2
 -{(2\mathbf p\mathbin\cdot\mathbf k-|\mathbf k|^2)^2
   \over2(\omega_p+\omega_q)^2}.
 \label{eq:v6v68-exact-relative-form}
\end{equation}

\begin{theorem}[Exact leading reflection-plane asymptotic]
\label{thm:v6v68-leading-asymptotic}
For every fixed \(m\ge0\), \(\xi\in\mathbb R\), and
\(\mathbf k\ne0\),
\begin{equation}
 \lim_{\varepsilon\downarrow0}
 {\varepsilon\over|\mathbf k|^2}
 \mathcal K^{(m)}_{\xi,\varepsilon}(\mathbf k)
 ={A_\xi\over64\pi^2}.
 \label{eq:v6v68-leading-limit}
\end{equation}
The coefficient is independent of the mass and is strictly positive for
every improvement parameter.
\end{theorem}

\begin{proof}
Set \(\mathbf p=\mathbf u/\varepsilon\), keep \(\mathbf k\) fixed, and
write \(n=\mathbf u/|\mathbf u|\).  Away from \(\mathbf u=0\),
\begin{align*}
 \varepsilon\omega_p&\longrightarrow|\mathbf u|,&
 \varepsilon\omega_q&\longrightarrow|\mathbf u|,\\
 \varepsilon(\omega_p+\omega_q)&\longrightarrow2|\mathbf u|,&
 \omega_p-\omega_q&\longrightarrow n\mathbin\cdot\mathbf k.
\end{align*}
Thus the rescaled radial integrand converges to
\[
 {1\over8(2\pi)^3}e^{-4|\mathbf u|}
 |\mathbf k|^2
 \left(a_\xi-{1\over2}(n\mathbin\cdot e)^2\right)^2
 {d^3u\over|\mathbf u|^2},
 \qquad e={\mathbf k\over|\mathbf k|}.
\]
Split the integral into \(|\mathbf u|\le\delta\) and its complement.
Corollary~\ref{cor:v6v67-Newton-pole-cancel}, followed by the local
integrability of \(d^3p/(\omega_p\omega_q)\), makes the first part uniformly
small after \(\delta\downarrow0\).  On the complement, dominated convergence
applies because of the exponential.  Finally,
\begin{align*}
 {1\over4\pi}\int_{S^2}
 \left(a_\xi-{1\over2}(n\mathbin\cdot e)^2\right)^2d\sigma(n)
 &=a_\xi^2-{a_\xi\over3}+{1\over20}=A_\xi,\\
 \int_0^\infty e^{-4r}\,dr&={1\over4}.
\end{align*}
Together with
\(4\pi/[8(2\pi)^3]=1/(16\pi^2)\), this proves
\eqref{eq:v6v68-leading-limit}.
\end{proof}

\begin{corollary}[Local boundary contact]
\label{cor:v6v68-local-contact}
For every \(g\in\mathcal S(\mathbb R^3)\),
\begin{equation}
 \lim_{\varepsilon\downarrow0}
 \varepsilon\mathcal Q^{(m)}_{\xi,\varepsilon}(g)
 ={A_\xi\over64\pi^2}\|\nabla g\|_2^2.
 \label{eq:v6v68-form-leading}
\end{equation}
Hence the leading divergence is the local boundary distribution
\begin{equation}
 {A_\xi\over64\pi^2\varepsilon}(-\Delta)\delta^{(3)}.
 \label{eq:v6v68-boundary-contact}
\end{equation}
At \(\xi=1/6\), its coefficient is still
\(1/(2880\pi^2\varepsilon)\).
\end{corollary}

\begin{proof}
The pointwise limit is Theorem~\ref{thm:v6v68-leading-asymptotic}.  The
bound \(|F_{\xi,m}|\le C_\xi|\mathbf k|^2\) reduces domination to the
massive or massless convolution
\[
 \varepsilon|\mathbf k|^2
 \int {d^3p\over\omega_p\omega_q}
 e^{-2\varepsilon(\omega_p+\omega_q)}.
\]
Splitting into
\(|\mathbf p|\le2|\mathbf k|+2m\) and its complement gives a bound by
\(C_{m,\xi}|\mathbf k|^2(1+|\mathbf k|)^N\) for a fixed finite \(N\),
uniformly for \(0<\varepsilon\le1\).  The Schwartz decay of
\(\widehat g\) permits dominated convergence.  Plancherel identifies the
last integral with \(\|\nabla g\|_2^2\).
\end{proof}

\begin{corollary}[No nontrivial sharp-time Hilbert vector]
\label{cor:v6v68-no-boundary-vector}
Let \(Y_\varepsilon(g)\) be the positive OS Hilbert-space vector whose norm
is \eqref{eq:v6v68-slice-form}.  If \(\nabla g\ne0\), then
\begin{equation}
 \|Y_\varepsilon(g)\|^2
 ={A_\xi\over64\pi^2\varepsilon}\|\nabla g\|_2^2
 +o(\varepsilon^{-1}).
 \label{eq:v6v68-vector-divergence}
\end{equation}
In particular, \(Y_\varepsilon(g)\) has no norm-convergent sharp-time
limit.
\end{corollary}

\begin{proof}
This is \eqref{eq:v6v68-form-leading} and the strict positivity of
\(A_\xi\).  Every norm-convergent family is norm bounded.
\end{proof}

\subsection{Finite nonlocal boundary distribution}

The leading term is local.  To identify what remains, center the relative
momentum by
\begin{equation}
 \boldsymbol\ell=\mathbf p-{\mathbf k\over2},
 \qquad
 \omega_\pm=left(|\boldsymbol\ell\pm\mathbf k/2|^2+m^2\right)^{1/2}.
 \label{eq:v6v68-centered-momentum}
\end{equation}
For \(r=|\boldsymbol\ell|\gg |\mathbf k|+m\), Taylor expansion in the
centered variable gives, uniformly in angle,
\begin{align}
 \omega_++\omega_-&=2r+O((m^2+|\mathbf k|^2)/r),
 \label{eq:v6v68-sum-expansion}\\
 \omega_+-\omega_-&=widehat\ell\mathbin\cdot\mathbf k
 +O(|\mathbf k|(m^2+|\mathbf k|^2)/r^2),
 \label{eq:v6v68-difference-expansion}\\
 {r^2\over\omega_+\omega_-}&=1+O((m^2+|\mathbf k|^2)/r^2).
 \label{eq:v6v68-product-expansion}
\end{align}
There is no order-(r^{-1}) term in the difference or product.  This is the
parity gain produced by centering the two identical particles.

\begin{theorem}[Existence of the contact-subtracted finite part]
\label{thm:v6v68-finite-part}
For fixed \(m\ge0\), \(\xi\), and \(\mathbf k\), the limit
\begin{equation}
 \mathcal K^{(m),\mathrm{fin}}_\xi(\mathbf k)
 =\lim_{\varepsilon\downarrow0}
 \left[
 \mathcal K^{(m)}_{\xi,\varepsilon}(\mathbf k)
 -{A_\xi\over64\pi^2\varepsilon}|\mathbf k|^2
 \right]
 \label{eq:v6v68-finite-kernel}
\end{equation}
exists as a tempered, rotationally invariant spatial distribution.  In the
massless case there is a finite real constant \(B_\xi\) such that
\begin{equation}
 \mathcal K^{(0),\mathrm{fin}}_\xi(\mathbf k)
 =B_\xi|\mathbf k|^3.
 \label{eq:v6v68-massless-finite-kernel}
\end{equation}
No sign for \(B_\xi\) is asserted.
\end{theorem}

\begin{proof}
Use \(\boldsymbol\ell\) from
\eqref{eq:v6v68-centered-momentum}.  After angular integration, the radial
density in \eqref{eq:v6v68-slice-kernel} equals the constant responsible
for \eqref{eq:v6v68-leading-limit} plus
\begin{equation}
 O_{m,\xi,\mathbf k}(r^{-2}).
 \label{eq:v6v68-integrable-tail}
\end{equation}
Equations \eqref{eq:v6v68-sum-expansion}--
\eqref{eq:v6v68-product-expansion} prove this estimate; the potentially
nonintegrable order-(r^{-1}) term is absent.  Subtract the constant radial
tail using \(e^{-4\varepsilon r}\).  Its integral is exactly the local term
in \eqref{eq:v6v68-finite-kernel}.  The remainder
\eqref{eq:v6v68-integrable-tail} is in \(L^1(dr)\), so dominated convergence
gives a finite limit.  Changing
\(e^{-2\varepsilon(\omega_++\omega_-)}\) to
\(e^{-4\varepsilon r}\) contributes a term tending to zero, by
\eqref{eq:v6v68-sum-expansion} and another split at
\(r=\varepsilon^{-1/2}\).  The locally integrable finite-(r) part is
unchanged.  Polynomial growth in \(\mathbf k\) follows from
\eqref{eq:v6v67-quadratic-zero}, proving temperedness.

When \(m=0\), scaling
\(\boldsymbol\ell=|\mathbf k|u\) shows that the finite part is homogeneous
of degree three.  Rotational invariance then forces the form
\(B_\xi|\mathbf k|^3\).
\end{proof}

\begin{corollary}[The boundary Sobolev order is exactly the V66 order]
\label{cor:v6v68-boundary-Sobolev}
In the massless free theory, the contact-subtracted boundary form
\begin{equation}
 \mathcal Q^{\mathrm{fin}}_\xi(g)
 =\int {d^3k\over(2\pi)^3}
  |\widehat g(\mathbf k)|^2
  \mathcal K^{(0),\mathrm{fin}}_\xi(\mathbf k)
 \label{eq:v6v68-finite-form}
\end{equation}
is continuous on \(H^{3/2}(\mathbb R^3)\), with
\begin{equation}
 |\mathcal Q^{\mathrm{fin}}_\xi(g)|
 \le |B_\xi|\|g\|_{\dot H^{3/2}}^2.
 \label{eq:v6v68-H32-bound}
\end{equation}
This realizes the necessary order found in
Proposition~\ref{prop:v6v66-Sobolev-threshold} at the level of a finite
distribution.  It does not turn that distribution into a positive boundary
state norm.
\end{corollary}

\begin{proof}
Insert \eqref{eq:v6v68-massless-finite-kernel} in
\eqref{eq:v6v68-finite-form}; the resulting integral is exactly
\(B_\xi\|g\|_{\dot H^{3/2}}^2\) in the stated Fourier convention.
\end{proof}

\subsection{What may and may not be renormalized at the plane}

\begin{proposition}[A scalar finite-part subtraction is not Hilbert
renormalization]
\label{prop:v6v68-not-Hilbert-renormalization}
Subtracting the local number
\begin{equation}
 {A_\xi\over64\pi^2\varepsilon}\|\nabla g\|_2^2
 \label{eq:v6v68-number-subtraction}
\end{equation}
from \(\|Y_\varepsilon(g)\|^2\) can define the distributional finite part
of Theorem~\ref{thm:v6v68-finite-part}, but it does not produce a
norm-convergent family of the vectors \(Y_\varepsilon(g)\) in the original
OS Hilbert space.  Positivity of the finite part requires a separate proof.
\end{proposition}

\begin{proof}
The vectors themselves retain the divergent norms
\eqref{eq:v6v68-vector-divergence}; changing a reported scalar quadratic
form does not change those vectors or their mutual Hilbert distances.  In
addition, \eqref{eq:v6v68-number-subtraction} has the negative sign needed
to cancel a positive divergence.  A difference of positive forms need not
be positive, and Theorem~\ref{thm:v6v68-finite-part} makes no sign claim for
\(B_\xi\).
\end{proof}

\begin{assessment}[The reflection plane is not a physical counterterm
surface]
\label{ass:v6v68-plane-assessment}
The contact in \eqref{eq:v6v68-boundary-contact} appears when two reflected
insertions collide with the auxiliary OS plane.  At every positive
separation, Proposition~\ref{prop:v6v67-separated-contacts} says that bulk
local contacts vanish between the halves.  Absorbing
\eqref{eq:v6v68-boundary-contact} as though it were a bulk Einstein
counterterm would introduce reflection-plane-dependent data.  The
constructive target should retain positive-time smeared sources and recover
the Hamiltonian as a semigroup generator or closed form, rather than demand
a local energy-density vector at exact time zero.
\end{assessment}

\begin{firewall}[Boundary distribution versus Hamiltonian generator]
\label{fw:v6v68-generator}
The failure of a sharp-time local source vector does not imply failure of
the total Hamiltonian.  Indeed
\(F_{\xi,m}(\mathbf p,0)=0\) expresses that the normal-ordered total free
Hamiltonian annihilates the vacuum, and the OS semigroup has its usual
positive generator.  What fails is the stronger demand that the
Newton-smeared local density itself converge as a boundary Hilbert vector.
\end{firewall}

\begin{firewall}[The conformal metric problem is separate]
\label{fw:v6v68-conformal-separate}
The coefficient \(A_\xi\) concerns positive matter spectral weight and is
strictly positive even at conformal scalar improvement.  It neither creates
nor removes the negative Euclidean conformal metric block of V58.  The
matter boundary collision and the gravitational conformal direction require
different constructions and must not be canceled against one another.
\end{firewall}

\begin{theorem}[Sixth-series V68 result]
\label{thm:v6v68-status}
As a Newton-smeared improved scalar stress insertion approaches the OS
reflection plane, its positive norm diverges as
\(A_\xi\|\nabla g\|_2^2/(64\pi^2\varepsilon)\), where
\(A_\xi\ge1/45\) for every improvement coefficient.  The divergence is a
local boundary gradient contact and forbids a nontrivial sharp-time Hilbert
vector.  Its contact-subtracted finite part nevertheless exists as a
tempered distribution; in the massless theory it is
\(B_\xi|k|^3\) and is continuous on \(H^{3/2}\).  This finite part is not
automatically positive.  The correct constructive object remains the
interior, spacetime-smeared dRNCC of V67.
\end{theorem}

\begin{proof}
Use Theorem~\ref{thm:v6v68-leading-asymptotic},
Corollaries~\ref{cor:v6v68-local-contact} and
\ref{cor:v6v68-no-boundary-vector},
Theorem~\ref{thm:v6v68-finite-part},
Corollary~\ref{cor:v6v68-boundary-Sobolev}, and
Proposition~\ref{prop:v6v68-not-Hilbert-renormalization}, subject to the two
firewalls.
\end{proof}

V69 will keep the positive-time source domain and contract the complete
Ward-quotiented stress tensor with the positive transverse--traceless
linearized graviton carrier.  It will distinguish the TT exchange form from
the attractive scalar constraint and test which part admits a common
reflected one-graviton source certificate.

\clearpage
\section{V69: the reflected TT carrier and a joint free matter--graviton source}
\label{sec:v6v69}

V67--V68 determine the correct positive-time domain for the scalar stress
source.  This version couples that source to the propagating
transverse--traceless part of the linearized metric.  The reflected TT
kernel factorizes as a square and gives a positive one-graviton certificate.
The scalar Newton constraint is different: it is instantaneous, contributes
no separated cross-plane covariance, and reappears as the attractive Schur
completion of V62 and V66.  The two carriers must therefore remain distinct.

\subsection{The positive TT reflection kernel}

For \(\mathbf l\ne0\), put
\begin{equation}
 P_{ij}(\mathbf l)=\delta_{ij}-{l_il_j\over|\mathbf l|^2}
 \label{eq:v6v69-vector-projector}
\end{equation}
and
\begin{equation}
 P^{\mathrm{TT}}_{ij,ab}(\mathbf l)
 ={1\over2}\bigl(P_{ia}P_{jb}+P_{ib}P_{ja}-P_{ij}P_{ab}\bigr).
 \label{eq:v6v69-TT-projector}
\end{equation}
This is the orthogonal rank-two projector on symmetric spatial tensors that
are transverse and traceless.  Its value at \(\mathbf l=0\) may be chosen
arbitrarily on a null set.

The free Euclidean TT covariance from V59 is
\begin{equation}
 C^{\mathrm{TT}}_{ij,ab}(\tau-\sigma,\mathbf l)
 ={e^{-|\mathbf l||\tau-\sigma|}\over2|\mathbf l|}
 P^{\mathrm{TT}}_{ij,ab}(\mathbf l).
 \label{eq:v6v69-TT-covariance}
\end{equation}

\begin{theorem}[Factorized reflected TT source form]
\label{thm:v6v69-TT-factorization}
For every Schwartz symmetric source \(J_{ij}(\tau,\mathbf x)\) supported in
\(\tau>0\), its reflected TT quadratic form is
\begin{equation}
 \mathcal R_{\mathrm{TT}}(J)
 =\int_{\mathbb R^3}{d^3l\over(2\pi)^3}{1\over2|\mathbf l|}
 \left\|
 P^{\mathrm{TT}}(\mathbf l)
 \int_0^\infty e^{-|\mathbf l|\tau}
       \widehat J(\tau,\mathbf l)\,d\tau
 \right\|_{\mathrm{HS}}^2
 \ge0.
 \label{eq:v6v69-reflected-TT-form}
\end{equation}
The form is finite and vanishes on every source of the form
\begin{equation}
 J_{ij}=\partial_i a_j+\partial_j a_i+\delta_{ij}b.
 \label{eq:v6v69-spatial-gauge-trace}
\end{equation}
\end{theorem}

\begin{proof}
Reflecting the first time argument of
\eqref{eq:v6v69-TT-covariance} replaces
\(|\tau-\sigma|\) by \(\tau+\sigma\).  Hence its scalar factor is
\[
 {e^{-|\mathbf l|(\tau+\sigma)}\over2|\mathbf l|}
 ={1\over2|\mathbf l|}
  e^{-|\mathbf l|\tau}e^{-|\mathbf l|\sigma}.
\]
The projector is self-adjoint and idempotent, so the double source pairing
is exactly the squared norm in
\eqref{eq:v6v69-reflected-TT-form}.  Near \(\mathbf l=0\), the measure
\(d^3l/|\mathbf l|\) is locally integrable; Schwartz decay gives the large
momentum bound.  Finally,
\(P^{\mathrm{TT}}_{ij,ab}l_a=0\) and
\(P^{\mathrm{TT}}_{ij,ab}\delta_{ab}=0\), proving the quotient statement.
\end{proof}

\begin{corollary}[No gauge-projector norm loss]
\label{cor:v6v69-projector-no-loss}
The TT source norm is bounded by the same expression with
\(P^{\mathrm{TT}}\) deleted.  Thus the apparent denominators in
\eqref{eq:v6v69-vector-projector} produce no operator-norm loss after the
orthogonal projector is assembled.
\end{corollary}

\begin{proof}
An orthogonal projector has operator norm one.
\end{proof}

\subsection{The scalar constraint has no separated reflected carrier}

At linear order, the scalar constraint kernel is instantaneous.  Up to the
positive normalization used in Lemma~\ref{lem:v6v66-Newton-completion}, its
Euclidean distribution is
\begin{equation}
 C_{\mathrm N}(\tau-\sigma,\mathbf k)
 ={\delta(\tau-\sigma)\over|\mathbf k|^2}.
 \label{eq:v6v69-instantaneous-Newton}
\end{equation}

\begin{proposition}[Separated-plane vanishing of the Newton constraint]
\label{prop:v6v69-Newton-cross-zero}
If both source tests are supported in \(\tau\ge\tau_0>0\), then reflection
of one argument gives
\begin{equation}
 {\delta(-\tau-\sigma)\over|\mathbf k|^2}=0
 \label{eq:v6v69-reflected-Newton-zero}
\end{equation}
on their product support.  Therefore the scalar constraint supplies no
separated cross-plane oscillator norm.  Eliminating it within one time
slice instead gives the attractive energy
\eqref{eq:v6v66-Newton-completion}.
\end{proposition}

\begin{proof}
For \(\tau,\sigma\ge\tau_0\), one has
\(\tau+\sigma\ge2\tau_0\), which is disjoint from the support of the delta
distribution.  The Schur-complement statement is
Lemma~\ref{lem:v6v66-Newton-completion}.
\end{proof}

\begin{assessment}[Propagating and constrained gravitational rows]
\label{ass:v6v69-carrier-split}
The TT field contributes a positive reflected covariance because it is a
propagating oscillator.  The Newton scalar is fixed by a constraint and
contributes to the within-slice Hamiltonian with an attractive sign.  A
common ``positive graviton covariance'' that merges these rows would erase
the ADM reduction of V61.  Positivity must instead mean TT reflection
positivity together with lower-boundedness of the matter Hamiltonian after
the scalar constraint is eliminated.
\end{assessment}

\subsection{A joint scalar--TT interaction vector}

Let \(h^{\mathrm{TT}}_{ij}\) be an independent free TT Gaussian field with
covariance \eqref{eq:v6v69-TT-covariance}, and let
\(T^{(\xi)}_{ij}\) be the free scalar stress tensor from
\eqref{eq:v6v66-improved-stress}.  For real
\(f\in C_c^\infty((0,\infty))\) and real
\(g\in\mathcal S(\mathbb R^3)\), define the positive-time Wick insertion
\begin{equation}
 \mathsf Z_{\xi,m}(f,g)
 =\int_0^\infty d\tau\int d^3x\,
 f(\tau)g(\mathbf x)
 :\!h^{\mathrm{TT}}_{ij}(\tau,\mathbf x)
 T^{(\xi)}_{ij}(\tau,\mathbf x)\!:.
 \label{eq:v6v69-joint-insertion}
\end{equation}
Normal ordering is with respect to the product free vacuum.

For a graviton momentum \(\mathbf l\), scalar momenta
\(\mathbf p,\mathbf q\), and
\begin{equation}
 E=|\mathbf l|+\omega_p+\omega_q,
 \qquad
 \mathbf r=\mathbf l+\mathbf p+\mathbf q,
 \label{eq:v6v69-total-energy-momentum}
\end{equation}
let \(\widetilde f(E)=\int_0^\infty f(\tau)e^{-E\tau}d\tau\).

\begin{theorem}[Exact joint one-graviton/two-scalar OS norm]
\label{thm:v6v69-joint-norm}
The insertion \(\mathsf Z_{\xi,m}(f,g)\) defines a vector in the reflected
product Hilbert space, and
\begin{align}
 \|\mathsf Z_{\xi,m}(f,g)\|_{\mathrm{OS}}^2
 ={1\over16}\int {d^3l\,d^3p\,d^3q\over(2\pi)^9}
 {&|\widetilde f(E)|^2|\widehat g(\mathbf r)|^2
  \over |\mathbf l|\omega_p\omega_q}
 \nonumber\\
 &\times
 V^{(\xi)}_{ij}(p,q)
 P^{\mathrm{TT}}_{ij,ab}(\mathbf l)
 \overline{V^{(\xi)}_{ab}(p,q)}
 \ge0.
 \label{eq:v6v69-joint-positive-norm}
\end{align}
The integral is finite for every \(m\ge0\) and \(\xi\in\mathbb R\).
\end{theorem}

\begin{proof}
Acting on the product vacuum, the Wick monomial creates exactly one TT
graviton and two identical scalars.  The graviton polarization sum is
\(P^{\mathrm{TT}}/(2|\mathbf l|)\); the scalar pair normalization and
identical-particle factor are those in
Theorem~\ref{thm:v6v67-positive-norm}.  Their product gives the coefficient
\(1/16\).  Time reflection yields the Laplace transform at the sum \(E\),
while spatial smearing yields \(\widehat g(\mathbf r)\).  This proves the
formula.  Positivity follows because \(P^{\mathrm{TT}}\) is an orthogonal
projector.

If \(\operatorname{supp}f\subset[\tau_0,\tau_1]\), then
\(|\widetilde f(E)|\le\|f\|_1e^{-\tau_0E}\).  Moreover
\begin{equation}
 \left|P^{\mathrm{TT}}_{ij,ab}(\mathbf l)
 V^{(\xi)}_{ab}(p,q)\right|
 \le C_\xi(\omega_p+\omega_q)^2.
 \label{eq:v6v69-vertex-growth}
\end{equation}
Indeed every momentum in the vertex is bounded by the corresponding
on-shell energy and the projector has norm one.  Exponential decay then
makes the nine-dimensional ultraviolet integral finite.  At zero momentum,
the measures
\(d^3l/|\mathbf l|\),
\(d^3p/\omega_p\), and \(d^3q/\omega_q\) are locally integrable, including
when \(m=0\).
\end{proof}

\begin{corollary}[A nontrivial positive Einstein--matter source domain]
\label{cor:v6v69-positive-domain}
Finite linear combinations of insertions
\eqref{eq:v6v69-joint-insertion}, with time support bounded away from the
reflection plane, form a nontrivial positive source domain for the free
scalar plus free TT graviton product theory.
\end{corollary}

\begin{proof}
The free product Gaussian is reflection positive, and
Theorem~\ref{thm:v6v69-joint-norm} proves that each insertion has finite
norm.  Polarization gives the corresponding positive Gram matrices for
finite linear combinations.
\end{proof}

\subsection{Exact improvement cancellation and localization defect}

The stress Ward identity is tied to the scalar-pair momentum
\(K=p+q\), whereas the TT projector is tied to the graviton momentum
\(\mathbf l\).  The smearing momentum
\(\mathbf r=\mathbf l+\mathbf K\) measures their mismatch.

\begin{proposition}[TT improvement is a quadratic localization defect]
\label{prop:v6v69-improvement-defect}
For the spatial vertex in
\eqref{eq:v6v69-joint-positive-norm},
\begin{align}
 P^{\mathrm{TT}}_{ij,ab}(\mathbf l)V^{(\xi)}_{ab}(p,q)
 =P^{\mathrm{TT}}_{ij,ab}(\mathbf l)
 \bigl(&p_aq_b+p_bq_a
 \nonumber\\
 &-2\xi r_ar_b\bigr).
 \label{eq:v6v69-improvement-defect}
\end{align}
In particular, at exact momentum conservation \(\mathbf r=0\), both the
trace part and the scalar improvement term vanish from the TT amplitude.
For localized sources, the improvement defect obeys
\begin{equation}
 \left|2\xi P^{\mathrm{TT}}(\mathbf l)(\mathbf r\otimes\mathbf r)
ight|
 \le2|\xi||\mathbf r|^2.
 \label{eq:v6v69-improvement-defect-bound}
\end{equation}
\end{proposition}

\begin{proof}
The terms in \eqref{eq:v6v66-stress-vertex} proportional to the spatial
metric are traces and are killed by \(P^{\mathrm{TT}}\).  Its remaining
improvement term is \(-2\xi K_aK_b\).  Since
\(\mathbf K=\mathbf r-\mathbf l\) and the projector kills every tensor
factor containing \(\mathbf l\),
\[
 P^{\mathrm{TT}}(\mathbf l)(\mathbf K\otimes\mathbf K)
 =P^{\mathrm{TT}}(\mathbf l)(\mathbf r\otimes\mathbf r).
\]
The norm bound follows from
\(\|P^{\mathrm{TT}}\|=1\).
\end{proof}

\begin{corollary}[Localization pays ordinary test derivatives]
\label{cor:v6v69-localization-derivatives}
The improvement part of the joint norm is controlled by Schwartz seminorms
of two spatial derivatives of \(g\).  It introduces no inverse power of the
graviton momentum.
\end{corollary}

\begin{proof}
In \eqref{eq:v6v69-joint-positive-norm}, the defect is multiplied by
\(\widehat g(\mathbf r)\).  Equation
\eqref{eq:v6v69-improvement-defect-bound} replaces it by
\(|\mathbf r|^2|\widehat g(\mathbf r)|\), the Fourier transform magnitude
associated with \(-\Delta g\).  The sole graviton denominator remains
\(|\mathbf l|^{-1}\), whose infrared integrability was proved in
Theorem~\ref{thm:v6v69-joint-norm}.
\end{proof}

\begin{definition}[Free reduced Einstein--matter source certificate]
\label{def:v6v69-FREMSC}
The free reduced Einstein--matter source certificate, FREMSC, consists of:
\begin{enumerate}[label=\textup{(FR\arabic*)},leftmargin=4.8em]
\item the positive factorized TT reflected form
      \eqref{eq:v6v69-reflected-TT-form};
\item the separate instantaneous scalar-constraint identity
      \eqref{eq:v6v69-reflected-Newton-zero}, together with the lower-bound
      problem for its attractive Schur completion;
\item the finite joint matter--TT source norm
      \eqref{eq:v6v69-joint-positive-norm}; and
\item the exact momentum-conserving TT Ward reduction plus the controlled
      localization defect
      \eqref{eq:v6v69-improvement-defect}.
\end{enumerate}
\end{definition}

\begin{theorem}[The free reduced certificate closes]
\label{thm:v6v69-FREMSC}
For a free scalar of arbitrary mass and improvement coefficient coupled at
source level to the free linearized TT graviton, FREMSC holds on smooth
positive-time data supported away from the reflection plane.
\end{theorem}

\begin{proof}
Rows \textup{(FR1)}--\textup{(FR4)} are respectively
Theorem~\ref{thm:v6v69-TT-factorization},
Proposition~\ref{prop:v6v69-Newton-cross-zero},
Theorem~\ref{thm:v6v69-joint-norm}, and
Proposition~\ref{prop:v6v69-improvement-defect}.
\end{proof}

\begin{firewall}[A finite interaction insertion is not an interacting
measure]
\label{fw:v6v69-not-interacting-measure}
Theorem~\ref{thm:v6v69-joint-norm} constructs one Wick interaction
insertion and all of its finite Gram combinations.  It does not exponentiate
\(h^{\mathrm{TT}}T\), remove a regulator from the nonlinear Einstein
vertices, or prove stability of the resulting measure.  Those require
multilinear estimates and counterterms beyond FREMSC.
\end{firewall}

\begin{firewall}[The scalar constraint remains an energy problem]
\label{fw:v6v69-scalar-energy}
The vanishing separated cross-plane kernel
\eqref{eq:v6v69-reflected-Newton-zero} does not remove Newton attraction.
The constraint contributes inside the Hamiltonian, where V62 proves only
conditional lower bounds.  TT reflection positivity and scalar-constraint
stability are independent required rows of the full continuum proof.
\end{firewall}

\begin{theorem}[Sixth-series V69 result]
\label{thm:v6v69-status}
The linearized TT graviton has a factorized positive reflected source form
with no projector norm loss.  The Newton scalar is instantaneous and has no
separated reflected carrier; it remains an attractive within-slice Schur
completion.  On the V67 positive-time domain, a joint one-TT-graviton and
two-scalar insertion has a finite explicit nonnegative OS norm.  At exact
momentum conservation its improvement and trace terms vanish under the TT
projector; localization restores only a quadratic test-momentum defect.
Thus the free reduced Einstein--matter source certificate closes, while
interaction exponentiation and scalar-constraint stability remain open.
\end{theorem}

\begin{proof}
Use Theorem~\ref{thm:v6v69-FREMSC} subject to the two firewalls.
\end{proof}

V70 will perform the compulsory YM/NS companion audit.  It will compare
their newest complete-source and localization mechanisms with FREMSC, then
choose between a multilinear TT insertion estimate and an upstream
Hamiltonian stability attack for the instantaneous scalar constraint.

\clearpage
\section{V70: companion audit and the fixed-arity/exponential stability gap}
\label{sec:v6v70}

V69 closes a nontrivial free one-graviton/two-scalar source certificate,
but it explicitly does not exponentiate the interaction.  This compulsory
audit shows that exponentiation, rather than another isolated insertion, is
the next upstream gate.  The Yang--Mills companion now separates
fixed-arity closure from a uniform connected-grade estimate, while the
Navier--Stokes companion separates positive density at every singular point
from a common-scale counting mechanism.  Both distinctions apply directly
to the Einstein--matter source problem.

\subsection{Compulsory V70 audit}

The active companion files inspected were
\begin{center}
\begin{tabular}{p{0.23\linewidth}p{0.19\linewidth}p{0.47\linewidth}}
programme & source & audited mathematical endpoint\\ \hline
Yang--Mills & \texttt{V50} & the log-affine geometric heat--Wilson bridge
has density grade \(r\) of degree at most \(2r+2\); connected Wick graphs
obey \(\rho+2q\ge m-2\), closing the required thermal homogeneity at every
fixed arity, but the constants have not been bounded by one exponential
base \(K^m\), so the uniform geometric connected-grade card UGCGC remains
open\\
Navier--Stokes & \texttt{V52} & the upper critical density of the terminal
datum has a solution-dependent uniform positive lower bound on the type-I
singular set; nevertheless density bounds do not count singular points,
because activity may occur on different or lacunary scales, leaving a
common-scale activity and Liouville gate
\end{tabular}
\end{center}
At the settled read used for this audit, YM V50 had \(24\,092\) logical
lines, \(758\,240\) bytes, and SHA--256
\begin{equation}
 \texttt{226f4204f409a5aeae615738fda88bb644f9d465c6418a5d8d2fc0e7370a328f},
 \label{eq:v6v70-YM-hash}
\end{equation}
while NS V52 had \(30\,700\) logical lines, \(1\,254\,056\) bytes, and
SHA--256
\begin{equation}
 \texttt{99f6290540d0219c9cb23a7225d320664ba4269a4556d90e1bb28a354407c3db}.
 \label{eq:v6v70-NS-hash}
\end{equation}
Both were read only and neither was modified.

\begin{theorem}[Companion-audit transfer]
\label{thm:v6v70-audit-transfer}
The companion results imply the following proof discipline for the current
SM--Einstein programme:
\begin{enumerate}[label=\textup{(A\arabic*)},leftmargin=4.5em]
\item V69 finite Gram matrices are fixed-arity information and do not prove
      convergence or normalization of the exponential
      \(\exp(-\kappa\int h^{\mathrm{TT}}T)\).
\item A bound for each source insertion or each momentum shell does not
      provide a common all-arity/all-scale budget; the actual interaction
      ancestry must be summed before the norm.
\item The next attack should use exact Gaussian integration of the metric
      carrier and test the resulting quartic matter functional for
      stability, rather than estimate arbitrarily high interaction moments
      one at a time.
\end{enumerate}
No YM mass-gap or NS regularity estimate is imported as an
Einstein--matter estimate.
\end{theorem}

\begin{proof}
YM V50's fixed-arity theorem leaves precisely the exponential-base bound
open.  NS V52's density firewall shows that individually positive local
information need not share one scale at which it can be summed.  FREMSC in
Definition~\ref{def:v6v69-FREMSC} asserts only a single joint interaction
insertion and its finite polarizations.  These three statements give
\textup{(A1)}--\textup{(A3)}.
\end{proof}

\subsection{Every fixed moment can exist while the interaction exponential
diverges}

The obstruction already occurs in two Gaussian variables and therefore
cannot be repaired by a formal fixed-order expansion.

\begin{theorem}[Exact metric--quadratic-matter exponential no-go]
\label{thm:v6v70-Gaussian-no-go}
Let \(H,X\) be independent standard real Gaussians and put
\begin{equation}
 Z=H(X^2-1).
 \label{eq:v6v70-toy-interaction}
\end{equation}
Then
\begin{equation}
 \mathbb E|Z|^n<\infty
 \qquad\hbox{for every fixed }n\ge0,
 \label{eq:v6v70-all-moments}
\end{equation}
but
\begin{equation}
 \mathbb E e^{\kappa Z}=\infty
 \qquad\hbox{for every }\kappa\ne0.
 \label{eq:v6v70-mgf-diverges}
\end{equation}
Thus existence of all fixed interaction moments, even with exact centering,
does not imply exponentiation at any nonzero coupling.
\end{theorem}

\begin{proof}
Independence gives
\[
 \mathbb E|Z|^n
 =\mathbb E|H|^n\,\mathbb E|X^2-1|^n<\infty
\]
because Gaussian variables have moments of every order.  Conditional
Gaussian integration in \(H\) gives the exact identity
\begin{equation}
 \mathbb E_H e^{\kappa H(X^2-1)}
 =\exp\left({\kappa^2\over2}(X^2-1)^2\right).
 \label{eq:v6v70-conditional-Gaussian}
\end{equation}
The remaining integral has exponent
\(\kappa^2x^4/2+O(x^2)\) against the Gaussian density
\(e^{-x^2/2}\), and hence diverges at both tails whenever
\(\kappa\ne0\).
\end{proof}

\begin{corollary}[No fixed-arity route to the V69 exponential]
\label{cor:v6v70-no-fixed-arity-route}
Bounds for every fixed Gram arity of
\(\mathsf Z_{\xi,m}(f,g)\) cannot by themselves prove integrability of its
exponential.  A proof must control the complete metric-induced quadratic
form in the stress tensor or an equivalent resummed object.
\end{corollary}

\begin{proof}
Theorem~\ref{thm:v6v70-Gaussian-no-go} is a centered
linear-Gaussian times quadratic-Gaussian model with precisely this logical
structure.  It satisfies every fixed-moment premise while its exponential
fails.
\end{proof}

\subsection{The exact stabilizing threshold in the model}

The toy model also identifies the correct kind of additional estimate.
For \(\lambda\ge0\), tilt the \(X\)-Gaussian by the stabilizing weight
\(e^{-\lambda(X^2-1)^2}\).

\begin{theorem}[Exact quartic stability threshold]
\label{thm:v6v70-toy-threshold}
The unnormalized integral
\begin{equation}
 \mathcal Z(\kappa,\lambda)
 =\mathbb E\exp\bigl(
   \kappa H(X^2-1)-\lambda(X^2-1)^2\bigr)
 \label{eq:v6v70-stabilized-partition}
\end{equation}
is finite if and only if
\begin{equation}
 \lambda\ge{\kappa^2\over2}.
 \label{eq:v6v70-threshold}
\end{equation}
At equality, the induced quartic cancels exactly and
\(\mathcal Z=1\).
\end{theorem}

\begin{proof}
Integrate \(H\) using
\eqref{eq:v6v70-conditional-Gaussian}:
\begin{equation}
 \mathcal Z(\kappa,\lambda)
 =\mathbb E_X
   \exp\left[-\left(\lambda-{\kappa^2\over2}\right)(X^2-1)^2\right].
 \label{eq:v6v70-stabilized-reduction}
\end{equation}
If the coefficient is nonnegative, the integrand is bounded by one, and
at equality it is exactly one.  If it is negative, the exponent has a
positive quartic leading term and the Gaussian integral diverges.
\end{proof}

\begin{assessment}[Meaning for the reduced metric source]
\label{ass:v6v70-threshold-meaning}
Conditioning a positive Gaussian metric carrier on a quadratic matter
stress produces a positive quartic exponential
\(\exp(\kappa^2\langle T,C T\rangle/2)\).  Normalization therefore
requires the matter action, nonlinear gravitational action, or an allowed
counterterm to dominate that complete induced quartic form.  Centering the
stress changes lower-degree terms but not its amplitude-four instability.
\end{assessment}

\subsection{Finite-regulator exponential stability card}

Let \(H\) be a finite-dimensional centered Gaussian metric variable with
strictly positive covariance \(C\), and let \(T(\Phi)\) be a quadratic
matter source on a finite-dimensional regulator space.  Write the matter
action as
\begin{equation}
 S_{\rm mat}(\Phi)=S_2(\Phi)+S_4(\Phi)+S_{\ge6}(\Phi),
 \label{eq:v6v70-matter-action-split}
\end{equation}
where the subscript denotes amplitude degree.

\begin{lemma}[Exact finite-regulator metric integration]
\label{lem:v6v70-metric-integration}
For real \(\kappa\),
\begin{equation}
 \int \exp\bigl(-\tfrac12\langle H,C^{-1}H\rangle
                 -\kappa\langle H,T(\Phi)\rangle\bigr)\,dH
 =Z_C\exp\left({\kappa^2\over2}
     \langle T(\Phi),CT(\Phi)\rangle\right).
 \label{eq:v6v70-exact-metric-integration}
\end{equation}
The sign of the linear coupling does not change the induced term.
\end{lemma}

\begin{proof}
Complete the square:
\[
 {1\over2}\langle H,C^{-1}H\rangle+\kappa\langle H,T\rangle
 ={1\over2}\langle H+\kappa CT,C^{-1}(H+\kappa CT)\rangle
 -{\kappa^2\over2}\langle T,CT\rangle.
\]
Translate the finite-dimensional Gaussian integral.
\end{proof}

\begin{proposition}[Necessary quartic ray condition]
\label{prop:v6v70-quartic-ray}
Assume \(S_{\ge6}=0\), \(S_2(R\Phi_0)=O(R^2)\), and
\(T(R\Phi_0)=R^2T(\Phi_0)\).  If for some \(\Phi_0\)
\begin{equation}
 S_4(\Phi_0)
 -{\kappa^2\over2}
  \langle T(\Phi_0),CT(\Phi_0)\rangle<0,
 \label{eq:v6v70-negative-ray}
\end{equation}
then the metric-integrated matter partition function is divergent.
\end{proposition}

\begin{proof}
On the ray \(R\Phi_0\), the effective action after
Lemma~\ref{lem:v6v70-metric-integration} is
\[
 R^4\left[S_4(\Phi_0)-{\kappa^2\over2}
 \langle T(\Phi_0),CT(\Phi_0)\rangle\right]+O(R^2).
\]
If the bracket is negative, the integrand grows as
\(e^{cR^4+O(R^2)}\).  By continuity the same lower bound holds on a
fixed cone around the ray for all large \(R\), whose Lebesgue volume is
nonzero; the integral diverges.
\end{proof}

\begin{definition}[Exponential reduced metric-source card]
\label{def:v6v70-ERMSC}
ERMSC requires, uniformly in the regulator:
\begin{enumerate}[label=\textup{(ER\arabic*)},leftmargin=5.2em]
\item exact integration or an exact conditional description of every
      positive metric Gaussian carrier before expanding in arity;
\item nonnegativity, or a quantified lower bound, for the complete quartic
      form
      \begin{equation}
       S_4(\Phi)-{\kappa^2\over2}
       \langle T^{\mathrm{TT}}(\Phi),
                    C_{\mathrm{TT}}T^{\mathrm{TT}}(\Phi)\rangle;
       \label{eq:v6v70-ERMSC-quartic}
      \end{equation}
\item simultaneous control of the attractive instantaneous Newton Schur
      term, rather than treating it as a reflected oscillator; and
\item one common all-scale bound compatible with the dRNCC test topology,
      Ward quotient, local counterterms, and positive transfer.
\end{enumerate}
\end{definition}

\begin{theorem}[V70 audit decision]
\label{thm:v6v70-decision}
The next upstream SM--Einstein gate is ERMSC.  A further fixed-arity TT
insertion estimate cannot close the continuum programme unless the
metric-induced quartic and the Newton constraint are first shown stable in
one common regulator-uniform form.
\end{theorem}

\begin{proof}
Theorem~\ref{thm:v6v70-audit-transfer} rules out fixed-arity inference.
Theorems~\ref{thm:v6v70-Gaussian-no-go} and
\ref{thm:v6v70-toy-threshold} show that the missing datum is quartic
stability.  Lemma~\ref{lem:v6v70-metric-integration} and
Proposition~\ref{prop:v6v70-quartic-ray} give its exact finite-regulator
form.
\end{proof}

\begin{firewall}[No instability theorem for the full Standard Model]
\label{fw:v6v70-no-full-instability}
The toy theorem and finite-regulator ray criterion do not prove that the
complete Higgs, gauge, fermion, ghost, and nonlinear metric action has a
negative quartic ray.  The Standard Model Higgs potential and nonlinear
Einstein counterterms may contribute stabilizing terms.  ERMSC is the
required comparison; its sign has not yet been determined.
\end{firewall}

\begin{firewall}[A formal Gaussian completion is regulator dependent]
\label{fw:v6v70-regulator}
The covariance \(C_{\mathrm{TT}}\), stress composite, and local
counterterms depend on the common cutoff and boundary conditions.  Taking
the continuum limit in each factor separately can lose the Ward and
locality cancellations.  Equation \eqref{eq:v6v70-ERMSC-quartic} must be
tested on the complete regulated source and only then transported between
scales.
\end{firewall}

\begin{theorem}[Sixth-series V70 result]
\label{thm:v6v70-status}
The compulsory audit identifies the same structural gap in all three
programmes: fixed local information need not possess a common all-arity or
all-scale budget.  An exact two-Gaussian model proves that every fixed
metric--matter interaction moment may exist while the exponential diverges
for every nonzero coupling; a stabilizing quartic has the exact threshold
\(\lambda\ge\kappa^2/2\).  At a finite regulator, exact metric Gaussian
integration produces the quartic source form
\(\kappa^2\langle T,CT\rangle/2\), and any negative effective quartic ray
forces divergence.  ERMSC records the correct next gate without asserting
its Standard Model sign.
\end{theorem}

\begin{proof}
Use Theorem~\ref{thm:v6v70-decision}, subject to the two firewalls.
\end{proof}

V71 will compute the ERMSC quartic on explicit regulated scalar test
families.  It will first separate the TT-induced term from the Newton
constraint, then test amplitude and spatial-frequency rescalings against a
local Higgs quartic to determine whether a regulator-uniform coercive
inequality is even dimensionally possible without higher-derivative input.

\clearpage
\section{V71: high-frequency Higgs packets and failure of the linearized ERMSC}
\label{sec:v6v71}

V70 reduces exponentiation of the linearized metric to a regulator-uniform
quartic stability test.  This version performs that test on an explicit
two-mode scalar family.  Both the TT-induced quartic and the instantaneous
Newton quartic grow like two powers of the packet frequency relative to the
local Higgs quartic.  In addition, placing the Higgs potential itself in a
linearized Newton source produces an attractive amplitude-eight term.  Thus
the Gaussian linearized metric split is an interior perturbative source
construction, not a globally normalizable measure for the unbounded Higgs
field.  The next attack must retain a nonlinear gravitational constraint or
declare a genuine effective-field-theory horizon.

\subsection{A normalized torus packet}

Let \(\mathbb T^3=(\mathbb R/2\pi\mathbb Z)^3\), use normalized Haar
measure \(d\mu=(2\pi)^{-3}d^3x\), and set
\begin{equation}
 \phi_N(x)=A\bigl(\cos(Nx_1)+\cos(Nx_2)\bigr),
 \qquad A\in\mathbb R,\quad N\in\mathbb N.
 \label{eq:v6v71-packet}
\end{equation}
For a periodic function \(f\), write
\(\widehat f(k)=\int_{\mathbb T^3}f(x)e^{-ik\cdot x}d\mu(x)\).

\begin{lemma}[Exact local norms of the packet]
\label{lem:v6v71-local-packet}
One has
\begin{align}
 \int_{\mathbb T^3}\phi_N^2\,d\mu&=A^2,
 \label{eq:v6v71-L2}\\
 \int_{\mathbb T^3}|\nabla\phi_N|^2\,d\mu&=A^2N^2,
 \label{eq:v6v71-H1}\\
 \int_{\mathbb T^3}\phi_N^4\,d\mu&={9\over4}A^4.
 \label{eq:v6v71-L4}
\end{align}
In particular, the local quartic is independent of \(N\).
\end{lemma}

\begin{proof}
The two coordinate cosines are orthogonal and each has mean square
\(1/2\), proving the first two identities.  If
\(X=\cos(Nx_1)\) and \(Y=\cos(Nx_2)\), odd mixed moments vanish and
\[
 \mathbb E(X+Y)^4
 =2\mathbb EX^4+6\mathbb EX^2\mathbb EY^2
 =2{3\over8}+6{1\over2}{1\over2}={9\over4}.
\]
\end{proof}

\subsection{Exact TT-induced quartic}

For a static scalar, let
\begin{equation}
 \tau^{(\xi)}_{ij}(\phi)
 =\partial_i\phi\,\partial_j\phi
 -{1\over2}\delta_{ij}(|\nabla\phi|^2+m^2\phi^2)
 +\xi(\delta_{ij}\Delta-\partial_i\partial_j)\phi^2.
 \label{eq:v6v71-static-stress}
\end{equation}
Define the zero-frequency TT source form
\begin{equation}
 \mathfrak Q_{\mathrm{TT}}(\phi)
 =\sum_{k\in\mathbb Z^3\setminus\{0\}}
 {\widehat\tau_{ij}(k)^*
  P^{\mathrm{TT}}_{ij,ab}(k)\widehat\tau_{ab}(k)
  \over|k|^2}.
 \label{eq:v6v71-TT-form}
\end{equation}

\begin{lemma}[Static TT projection removes mass, trace, and improvement]
\label{lem:v6v71-TT-removal}
For every periodic scalar \(\phi\),
\begin{equation}
 P^{\mathrm{TT}}(k)\widehat\tau^{(\xi)}(k)
 =P^{\mathrm{TT}}(k)
  \widehat{\nabla\phi\otimes\nabla\phi}(k),
 \qquad k\ne0.
 \label{eq:v6v71-TT-removal}
\end{equation}
Hence \(\mathfrak Q_{\mathrm{TT}}\) is independent of \(m\) and \(\xi\).
\end{lemma}

\begin{proof}
The mass and scalar kinetic terms multiplying \(\delta_{ij}\) are traces.
The first improvement term is also a trace, and the Fourier transform of
\(-\partial_i\partial_j\phi^2\) is a multiple of \(k_ik_j\).  The TT
projector kills both \(\delta_{ij}\) and every tensor with a factor \(k_i\).
\end{proof}

\begin{theorem}[Exact TT packet value]
\label{thm:v6v71-TT-packet}
For \(\phi_N\) in \eqref{eq:v6v71-packet},
\begin{equation}
 \boxed{\quad
 \mathfrak Q_{\mathrm{TT}}(\phi_N)={A^4N^2\over16}.
 \quad}
 \label{eq:v6v71-TT-exact}
\end{equation}
\end{theorem}

\begin{proof}
The only nonzero TT Fourier modes of
\(\nabla\phi_N\otimes\nabla\phi_N\) are
\[
 k=(\pm N,\pm N,0).
\]
At \(k=(N,N,0)\), for example, the tensor coefficient is
\begin{equation}
 -{A^2N^2\over4}
 (e_1\otimes e_2+e_2\otimes e_1).
 \label{eq:v6v71-cross-coefficient}
\end{equation}
Indeed
\(\sin(Nx_1)\sin(Nx_2)
 =\tfrac12[\cos N(x_1-x_2)-\cos N(x_1+x_2)]\).
For the unit vector
\(n=(e_1+e_2)/\sqrt2\), choose the transverse orthonormal basis
\(a=(e_1-e_2)/\sqrt2\), \(b=e_3\).  The TT projection of
\(e_1\otimes e_2+e_2\otimes e_1\) has entries
\(-1/2,1/2\) in the \(a,b\) basis and squared Hilbert--Schmidt norm
\(1/2\).  The same calculation holds at the other three modes.  Each mode
therefore contributes
\[
 {A^4N^4/32\over2N^2}={A^4N^2\over64}.
\]
Summing the four modes proves \eqref{eq:v6v71-TT-exact}.
\end{proof}

\begin{theorem}[Finite-cutoff TT coercivity horizon]
\label{thm:v6v71-TT-horizon}
Let the local scalar quartic be
\(\lambda\int\phi^4d\mu\), \(\lambda>0\), and use the V70 metric coupling
normalization.  On the packet ray, the quartic after exact TT Gaussian
integration is
\begin{equation}
 \mathcal E^{\mathrm{TT}}_4(\phi_N)
 =A^4\left({9\lambda\over4}-{\kappa^2N^2\over32}\right).
 \label{eq:v6v71-TT-effective-quartic}
\end{equation}
It is negative whenever
\begin{equation}
 N^2>{72\lambda\over\kappa^2}.
 \label{eq:v6v71-TT-horizon}
\end{equation}
Consequently no regulator-uniform inequality of the form
\begin{equation}
 {\kappa^2\over2}\mathfrak Q_{\mathrm{TT}}(\phi)
 \le\lambda\int\phi^4d\mu
 \label{eq:v6v71-impossible-TT-bound}
\end{equation}
can hold on all scalar modes at fixed nonzero \(\kappa\) and finite
\(\lambda\).
\end{theorem}

\begin{proof}
Insert \eqref{eq:v6v71-L4} and
\eqref{eq:v6v71-TT-exact} into the quartic form
\eqref{eq:v6v70-ERMSC-quartic}.  If
\eqref{eq:v6v71-TT-horizon} holds, the negative quartic dominates the
positive quadratic scalar action along the amplitude ray
\(R\phi_N\); Proposition~\ref{prop:v6v70-quartic-ray} then gives
divergence of the linearized metric-integrated partition function.
\end{proof}

This is a frequency obstruction, not a large-volume effect: the torus and
the normalized local \(L^4\) norm are fixed while \(N\) increases.

\subsection{The Newton quadratic source is uniformly active}

The quadratic part of the static improved energy density is
\begin{equation}
 \rho^{(\xi)}_2(\phi)
 ={1\over2}(|\nabla\phi|^2+m^2\phi^2)-\xi\Delta(\phi^2).
 \label{eq:v6v71-quadratic-density}
\end{equation}
Let \(P_0f=f-\int f,d\mu\) and define
\begin{equation}
 \mathfrak Q_{\mathrm N,2}(\phi)
 =\left\langle P_0\rho^{(\xi)}_2(\phi),
 (-\Delta)^{-1}P_0\rho^{(\xi)}_2(\phi)\right\rangle.
 \label{eq:v6v71-Newton-quadratic}
\end{equation}

\begin{theorem}[Exact Newton packet value and uniform improvement bound]
\label{thm:v6v71-Newton-packet}
With \(r=m^2/N^2\),
\begin{align}
 \mathfrak Q_{\mathrm N,2}(\phi_N)
 =A^4N^2\left[
 { (8\xi-1+r)^2\over64}
 +2\left(\xi+{r\over4}\right)^2
 \right].
 \label{eq:v6v71-Newton-exact}
\end{align}
For every \(m\ge0\), \(N\ge1\), and \(\xi\in\mathbb R\),
\begin{equation}
 \boxed{\quad
 \mathfrak Q_{\mathrm N,2}(\phi_N)
 \ge {A^4N^2(1+r)^2\over96}
 \ge {A^4N^2\over96}.
 \quad}
 \label{eq:v6v71-Newton-uniform}
\end{equation}
Thus no scalar improvement coefficient removes the high-frequency Newton
source.
\end{theorem}

\begin{proof}
At each of the four modes \(k=\pm2Ne_1,\pm2Ne_2\), the Fourier coefficient
of \(P_0\rho_2^{(\xi)}\) is
\begin{equation}
 c_{\rm ax}={A^2\over8}\bigl[m^2+(8\xi-1)N^2\bigr].
 \label{eq:v6v71-axis-density}
\end{equation}
At each of the four modes
\(k=(\pm N,\pm N,0)\), it is
\begin{equation}
 c_{\rm cr}=A^2\left(\xi N^2+{m^2\over4}\right).
 \label{eq:v6v71-cross-density}
\end{equation}
The axis modes contribute \(c_{\rm ax}^2/N^2\) in total and the cross
modes contribute \(2c_{\rm cr}^2/N^2\), proving
\eqref{eq:v6v71-Newton-exact}.  The bracket is a quadratic polynomial in
\(\xi\).  Completing its square gives
\begin{equation}
 { (8\xi-1+r)^2\over64}
 +2\left(\xi+{r\over4}\right)^2
 =3\left(\xi-{1-5r\over24}\right)^2
  +{(1+r)^2\over96},
 \label{eq:v6v71-Newton-square}
\end{equation}
which proves the bound.
\end{proof}

\begin{corollary}[Finite-cutoff Newton coercivity horizon]
\label{cor:v6v71-Newton-horizon}
Let \(\gamma_{\mathrm N}>0\) be the coefficient of the attractive Newton
completion after the scalar constraint is eliminated.  The quadratic-source
part of the packet quartic satisfies
\begin{equation}
 \lambda\int\phi_N^4d\mu
 -\gamma_{\mathrm N}\mathfrak Q_{\mathrm N,2}(\phi_N)
 \le A^4\left({9\lambda\over4}
              -{\gamma_{\mathrm N}N^2\over96}\right).
 \label{eq:v6v71-Newton-effective}
\end{equation}
It is negative whenever
\begin{equation}
 N^2>{216\lambda\over\gamma_{\mathrm N}}.
 \label{eq:v6v71-Newton-horizon}
\end{equation}
\end{corollary}

\begin{proof}
Use \eqref{eq:v6v71-L4} and the lower bound
\eqref{eq:v6v71-Newton-uniform}.
\end{proof}

\subsection{The Higgs potential makes the linearized Newton problem octic}

The full static energy source also contains
\begin{equation}
 \rho_4(\phi)=\lambda\phi^4.
 \label{eq:v6v71-quartic-density}
\end{equation}
This term is a spatial trace and does not enter the TT form, but it does
source the scalar constraint.

\begin{theorem}[Large-amplitude no-go for a linearized Newton constraint]
\label{thm:v6v71-Newton-octic}
Let \(\phi\) be any nonconstant smooth periodic scalar for which
\(P_0(\phi^4)\ne0\), and let \(\lambda,\gamma_{\mathrm N}>0\).  If the
linearized Newton constraint is eliminated with the full density
\begin{equation}
 \rho(R\phi)=R^2\rho_2(\phi)+\lambda R^4\phi^4,
 \label{eq:v6v71-amplitude-density}
\end{equation}
then the reduced action along the amplitude ray has the asymptotic
\begin{equation}
 S_{\rm red}(R\phi)
 =-\gamma_{\mathrm N}\lambda^2R^8
 \left\|(-\Delta)^{-1/2}P_0(\phi^4)\right\|_2^2
 +O(R^6).
 \label{eq:v6v71-octic-asymptotic}
\end{equation}
It tends to \(-\infty\), so the corresponding finite-regulator integral is
not normalizable.  The packet \(\phi_N\) satisfies the nonconstancy premise.
\end{theorem}

\begin{proof}
The scalar matter action contains terms of amplitude order two and four.
The eliminated Newton term is
\[
 -\gamma_{\mathrm N}
 \left\|(-\Delta)^{-1/2}P_0
  \bigl(R^2\rho_2+\lambda R^4\phi^4\bigr)\right\|_2^2.
\]
Expanding the square gives the displayed negative order-eight coefficient,
an order-six cross term, and an order-four quadratic-source term.  The
order-eight coefficient is strictly negative under the premise.  Continuity
on a finite-dimensional regulator gives a cone around the ray on which the
same leading sign holds, so the integral diverges.  Finally,
\((\cos Nx_1+\cos Nx_2)^4\) has nonzero Fourier modes and hence is not
constant.
\end{proof}

\begin{assessment}[The linearized metric split has two distinct horizons]
\label{ass:v6v71-two-horizons}
Theorems~\ref{thm:v6v71-TT-horizon} and
\ref{thm:v6v71-Newton-packet} give a frequency horizon even when only the
quadratic stress is retained.  Theorem~\ref{thm:v6v71-Newton-octic} gives an
amplitude horizon once the Higgs potential is included in the gravitational
source.  A single global Gaussian metric completion therefore cannot be
used over the full unbounded Higgs integration domain.
\end{assessment}

\subsection{What kind of ultraviolet input has the correct scaling}

The packet makes the missing scale visible.  For example,
\begin{equation}
 \int_{\mathbb T^3}\phi_N^2|\nabla\phi_N|^2d\mu
 ={3\over4}A^4N^2.
 \label{eq:v6v71-derivative-quartic}
\end{equation}
Thus a positive dimension-six matter operator has the same packet scaling
as the induced TT and quadratic Newton quartics.  Alternatively, replace the
static TT covariance by
\begin{equation}
 C_{\alpha}(k)={1\over|k|^2+\alpha|k|^4},
 \qquad \alpha>0.
 \label{eq:v6v71-higher-curvature-covariance}
\end{equation}
On the four TT-active packet modes, the exact source form becomes
\begin{equation}
 \mathfrak Q_{\mathrm{TT},\alpha}(\phi_N)
 ={A^4N^2\over16(1+2\alpha N^2)}
 \le {A^4\over32\alpha}.
 \label{eq:v6v71-higher-curvature-packet}
\end{equation}

\begin{proof}[Proof of \eqref{eq:v6v71-derivative-quartic}--
\eqref{eq:v6v71-higher-curvature-packet}]
For the first identity, independence of the two torus coordinates gives
\[
 \mathbb E[(X+Y)^2(\sin^2Nx_1+\sin^2Nx_2)]
 =2\left({1\over8}+{1\over4}\right)={3\over4}.
\]
For the second, repeat Theorem~\ref{thm:v6v71-TT-packet} and replace the
denominator \(2N^2\) of each active mode by
\(2N^2+4\alpha N^4\).
\end{proof}

\begin{firewall}[Correct scaling is not a completed repair]
\label{fw:v6v71-no-repair}
Equations \eqref{eq:v6v71-derivative-quartic} and
\eqref{eq:v6v71-higher-curvature-packet} show only that derivative matter
operators or higher-curvature metric terms have the dimensional capacity to
remove the displayed packet obstruction.  They do not prove a global
coercive inequality.  Higher-derivative metric propagators can also violate
the positive spectral decomposition, as recorded in V58.
\end{firewall}

\begin{firewall}[Linearized large-field failure is not full Einstein
instability]
\label{fw:v6v71-linearized-scope}
The amplitude limit in Theorem~\ref{thm:v6v71-Newton-octic} drives the
metric shift generated by the source outside the linearized neighbourhood.
The theorem proves that linearized Gaussian gravity cannot be integrated
globally over the unbounded Higgs field.  It does not prove that the full
nonlinear Einstein constraint has the same octic asymptotic or that the
Lorentzian Standard Model coupled to gravity is unstable.
\end{firewall}

\begin{definition}[Nonlinear-or-EFT metric stability fork]
\label{def:v6v71-NEMSF}
The next metric stability card must choose and prove one of two branches:
\begin{enumerate}[label=\textup{(NE\arabic*)},leftmargin=5.0em]
\item \emph{nonlinear branch:} solve or control the full regulated lapse and
      scalar metric constraint on the unbounded matter domain, with a lower
      bound for the resulting reduced action; or
\item \emph{EFT branch:} declare simultaneous frequency and amplitude
      horizons below which the linearized completion is used, include every
      operator required at that scale, and prove uniform transport only
      inside that domain.
\end{enumerate}
Neither branch may use the global linearized ERMSC disproved above.
\end{definition}

\begin{theorem}[Sixth-series V71 result]
\label{thm:v6v71-status}
For the explicit two-mode packet \(\phi_N\), the TT source form is exactly
\(A^4N^2/16\), independent of mass and improvement, whereas the local
Higgs quartic is \(9\lambda A^4/4\).  The quadratic Newton source is at
least \(A^4N^2/96\) uniformly in mass and improvement.  Hence both
linearized reduced quartics cross finite frequency horizons at fixed
couplings.  Including the Higgs potential in the linearized Newton source
produces a strictly negative amplitude-eight term and destroys global
normalizability.  The global linearized ERMSC is false; the programme must
follow the nonlinear-or-EFT fork of Definition~\ref{def:v6v71-NEMSF}.
\end{theorem}

\begin{proof}
Use Lemma~\ref{lem:v6v71-local-packet},
Theorems~\ref{thm:v6v71-TT-packet},
\ref{thm:v6v71-TT-horizon},
\ref{thm:v6v71-Newton-packet}, and
\ref{thm:v6v71-Newton-octic}, subject to the two firewalls.
\end{proof}

V72 will go upstream to the nonlinear branch.  It will study the conformal
scalar constraint on a compact spatial slice, derive its exact amplitude
scaling on positive conformal factors, and determine whether solving the
constraint changes the linearized negative-octic conclusion or instead
produces a different finite obstruction.

\clearpage
\section{V72: nonlinear conformal constraint sectors and the Higgs-amplitude threshold}
\label{sec:v6v72}

V71 proves that global Gaussian integration of a linearized metric is
incompatible with an unbounded Higgs amplitude.  The nonlinear branch must
therefore begin with the actual Hamiltonian constraint.  On a compact flat
slice, the conformal method gives an exact alternative: at fixed mean
curvature, sufficiently large positive matter density is not assigned an
arbitrarily negative Newton energy; it leaves the constraint sector
altogether.  Allowing the mean curvature to grow restores solutions but
moves the amplitude into the negative DeWitt trace direction.  This version
derives those statements without claiming a nonlinear Euclidean measure.

\subsection{The CMC Lichnerowicz equation on a flat three-torus}

Let \((\mathbb T^3,\bar g)\) be a flat three-torus of normalized reference
volume.  Write the physical metric and extrinsic curvature as
\begin{equation}
 g=\psi^4\bar g,
 \qquad
 K_{ij}=\psi^{-2}\bar\sigma_{ij}+{\tau\over3}g_{ij},
 \qquad \psi>0,
 \label{eq:v6v72-conformal-data}
\end{equation}
where \(\tau\) is constant and \(\bar\sigma\) is transverse and traceless
with respect to \(\bar g\).  Assume the momentum constraint has already
been solved and regard the physical energy density \(\rho\) as prescribed
for the moment.

Use the Hamiltonian-constraint convention
\begin{equation}
 R(g)+K^2-K_{ij}K^{ij}=16\pi G\rho+2\Lambda.
 \label{eq:v6v72-Hamiltonian-constraint}
\end{equation}

\begin{lemma}[Exact flat CMC conformal equation]
\label{lem:v6v72-Lichnerowicz}
The data \eqref{eq:v6v72-conformal-data} satisfy
\eqref{eq:v6v72-Hamiltonian-constraint} exactly when
\begin{equation}
 -8\bar\Delta\psi
 -|\bar\sigma|_{\bar g}^2\psi^{-7}
 +\left({2\tau^2\over3}-2\Lambda-16\pi G\rho\right)\psi^5
 =0.
 \label{eq:v6v72-Lichnerowicz}
\end{equation}
\end{lemma}

\begin{proof}
In three dimensions, conformal covariance gives
\begin{equation}
 R(\psi^4\bar g)=\psi^{-5}(-8\bar\Delta\psi)
 \label{eq:v6v72-scalar-curvature}
\end{equation}
because \(R(\bar g)=0\).  Orthogonality of trace and TT parts gives
\begin{equation}
 K_{ij}K^{ij}=\psi^{-12}|\bar\sigma|_{\bar g}^2+{\tau^2\over3},
 \qquad K=\tau.
 \label{eq:v6v72-K-split}
\end{equation}
Insert these identities in
\eqref{eq:v6v72-Hamiltonian-constraint} and multiply by \(\psi^5\).
\end{proof}

Put
\begin{equation}
 b_\tau={2\tau^2\over3}-2\Lambda.
 \label{eq:v6v72-btau}
\end{equation}

\begin{theorem}[Exact integrated sector identity]
\label{thm:v6v72-integrated-identity}
Every positive solution of \eqref{eq:v6v72-Lichnerowicz} obeys
\begin{equation}
 \int_{\mathbb T^3}(b_\tau-16\pi G\rho)\psi^5,d\mu_{\bar g}
 =\int_{\mathbb T^3}|\bar\sigma|^2\psi^{-7},d\mu_{\bar g}
 \ge0.
 \label{eq:v6v72-integrated-identity}
\end{equation}
Consequently, if
\begin{equation}
 16\pi G\rho(x)\ge b_\tau\quad\hbox{for all }x
 \label{eq:v6v72-supercritical-density}
\end{equation}
and either the inequality is strict on a set of positive measure or
\(\bar\sigma\not\equiv0\), then no positive conformal solution exists.
\end{theorem}

\begin{proof}
Integrate \eqref{eq:v6v72-Lichnerowicz}; the Laplacian has zero integral on
the torus.  This gives \eqref{eq:v6v72-integrated-identity}.  Under
\eqref{eq:v6v72-supercritical-density}, its left side is nonpositive.  It
is strictly negative if the density inequality is strict on a set of
positive measure, because \(\psi>0\).  Its right side is nonnegative and is
strictly positive when \(\bar\sigma\not\equiv0\).  Either strict condition
is a contradiction.
\end{proof}

\begin{corollary}[Time-symmetric positive-density obstruction]
\label{cor:v6v72-time-symmetric}
If \(\tau=0\), \(\Lambda\ge0\), \(\rho\ge0\), and at least one of
\(\rho\) or \(\bar\sigma\) is nonzero, then the flat conformal class has no
positive solution of the Hamiltonian constraint.
\end{corollary}

\begin{proof}
Here \(b_\tau=-2\Lambda\le0\), so
\(16\pi G\rho\ge b_\tau\).  Apply
Theorem~\ref{thm:v6v72-integrated-identity}.
\end{proof}

Thus a time-symmetric flat torus is not an admissible background for a
nontrivial nonnegative matter source.  A perturbative Newton completion
around that datum was already expanding around a point outside the exact
constraint surface.

\subsection{Complete classification of the homogeneous zero-TT branch}

Set \(\bar\sigma=0\) and let \(\rho=\rho_0\) be constant.

\begin{theorem}[Homogeneous fixed-sector classification]
\label{thm:v6v72-homogeneous-classification}
There is a positive solution of
\eqref{eq:v6v72-Lichnerowicz} if and only if
\begin{equation}
 b_\tau=16\pi G\rho_0.
 \label{eq:v6v72-homogeneous-balance}
\end{equation}
When this equality holds, every solution is a positive constant.  Hence a
constant conformal rescaling does not absorb a mismatch between matter
density and the CMC sector.
\end{theorem}

\begin{proof}
The equation is
\[
 -8\bar\Delta\psi+(b_\tau-16\pi G\rho_0)\psi^5=0.
\]
Integration and positivity of \(\psi^5\) force the coefficient to vanish.
The remaining equation \(\bar\Delta\psi=0\) has only constants on a compact
connected torus.  Conversely every positive constant solves the equation
when the coefficient vanishes.
\end{proof}

Take a homogeneous real Higgs component of amplitude \(R\), zero canonical
momentum, and potential
\begin{equation}
 V(R)=V_0+{m^2\over2}R^2+\lambda R^4,
 \qquad \lambda>0.
 \label{eq:v6v72-Higgs-potential}
\end{equation}

\begin{corollary}[Exact CMC ancestry of a homogeneous Higgs amplitude]
\label{cor:v6v72-Higgs-CMC}
In the zero-TT homogeneous branch, amplitude \(R\) lies on the constraint
surface exactly when
\begin{equation}
 \tau^2=3\Lambda+24\pi G V(R).
 \label{eq:v6v72-tau-ancestry}
\end{equation}
For large \(|R|\), any such branch satisfies
\begin{equation}
 |\tau(R)|\sim\sqrt{24\pi G\lambda}\,R^2.
 \label{eq:v6v72-tau-asymptotic}
\end{equation}
At fixed \(\tau\), only amplitudes solving the algebraic equation
\eqref{eq:v6v72-tau-ancestry} are admitted; the full unbounded amplitude
line is not contained in that sector.
\end{corollary}

\begin{proof}
Insert \(\rho_0=V(R)\) and
\eqref{eq:v6v72-btau} in
\eqref{eq:v6v72-homogeneous-balance}.  The asymptotic follows from the
positive quartic leading term.
\end{proof}

\begin{assessment}[Nonlinear replacement of the negative-octic ray]
\label{ass:v6v72-octic-replacement}
The linearized Newton elimination in V71 assigned every amplitude a metric
shift and produced a negative order-eight action.  The exact homogeneous
constraint does something different: a fixed CMC sector simply ceases to
contain generic large amplitudes.  Restoring those amplitudes requires an
unbounded family of trace-momentum sectors satisfying
\eqref{eq:v6v72-tau-ancestry}.  This is an exact constraint-domain result,
not yet a positive measure on the union of sectors.
\end{assessment}

\subsection{A constant TT momentum delays but does not remove the threshold}

Assume \(|\bar\sigma|^2=s^2>0\) is constant and retain constant
\(\rho_0\).  A constant conformal factor solves the equation precisely when
\begin{equation}
 -s^2\psi^{-7}+(b_\tau-16\pi G\rho_0)\psi^5=0.
 \label{eq:v6v72-constant-sigma-equation}
\end{equation}

\begin{proposition}[Constant-TT branch and volume blow-up]
\label{prop:v6v72-TT-branch}
The constant ansatz has a positive solution if and only if
\begin{equation}
 c:=b_\tau-16\pi G\rho_0>0.
 \label{eq:v6v72-subcritical-c}
\end{equation}
It is then unique within the constant ansatz and satisfies
\begin{equation}
 \psi^{12}={s^2\over c},
 \qquad
 \operatorname{Vol}(g)=\psi^6\operatorname{Vol}(\bar g)
 ={s\over\sqrt c}\operatorname{Vol}(\bar g).
 \label{eq:v6v72-volume-blowup}
\end{equation}
As \(\rho_0\uparrow b_\tau/(16\pi G)\), the physical volume diverges like
\(c^{-1/2}\); the branch cannot cross the threshold.
\end{proposition}

\begin{proof}
Equation \eqref{eq:v6v72-constant-sigma-equation} is equivalent to
\(c\psi^{12}=s^2\).  It has the displayed unique positive solution exactly
when \(c>0\).  Since the volume form of \(g=\psi^4\bar g\) is
\(\psi^6d\mu_{\bar g}\), the volume identity follows.  Nonexistence for
\(c\le0\), even beyond the constant ansatz, is also a consequence of
Theorem~\ref{thm:v6v72-integrated-identity}.
\end{proof}

\begin{corollary}[Finite density ceiling in every fixed CMC sector]
\label{cor:v6v72-density-ceiling}
For fixed \(\tau,\Lambda\), every constant-density solution with nonzero
constant TT momentum satisfies
\begin{equation}
 \rho_0<{1\over16\pi G}
 \left({2\tau^2\over3}-2\Lambda\right).
 \label{eq:v6v72-density-ceiling}
\end{equation}
For the positive quartic potential \eqref{eq:v6v72-Higgs-potential}, this
is a finite amplitude ceiling.
\end{corollary}

\begin{proof}
This is \(c>0\) in Proposition~\ref{prop:v6v72-TT-branch}.  Since
\(V(R)\to\infty\) as \(|R|\to\infty\), only a bounded amplitude interval
can satisfy it.
\end{proof}

\subsection{The trace momentum carries the large-amplitude scaling}

The DeWitt kinetic combination of a pure-trace extrinsic curvature is
\begin{equation}
 K_{ij}K^{ij}-K^2={\tau^2\over3}-\tau^2
 =-{2\tau^2\over3}.
 \label{eq:v6v72-DeWitt-trace}
\end{equation}

\begin{proposition}[The Higgs amplitude is transferred to the DeWitt trace
direction]
\label{prop:v6v72-trace-transfer}
Along the homogeneous constraint branch
\eqref{eq:v6v72-tau-ancestry},
\begin{equation}
 K_{ij}K^{ij}-K^2
 =-2\Lambda-16\pi G V(R)
 \sim-16\pi G\lambda R^4.
 \label{eq:v6v72-trace-transfer}
\end{equation}
Thus allowing \(\tau\) to vary with \(R\) restores the constraint but does
not create a positive coercive trace-metric quadratic form.
\end{proposition}

\begin{proof}
Insert \eqref{eq:v6v72-tau-ancestry} in
\eqref{eq:v6v72-DeWitt-trace}.  The asymptotic follows from
\eqref{eq:v6v72-Higgs-potential}.
\end{proof}

This is the nonlinear canonical form of the conformal-sign issue isolated
in V58.  The Hamiltonian constraint uses the negative trace direction to
balance positive energy; it does not turn that direction into a positive
Euclidean Gaussian.

\subsection{A constraint-sector card for the nonlinear branch}

\begin{definition}[Nonlinear conformal constraint sector card]
\label{def:v6v72-NCCSC}
NCCSC requires, at every regulator:
\begin{enumerate}[label=\textup{(NC\arabic*)},leftmargin=5.2em]
\item a declared conformal class, CMC parameter, TT momentum, cosmological
      term, and matter conformal weights;
\item existence and quantitative bounds for a positive solution of the
      coupled Lichnerowicz and momentum constraints on the admitted matter
      domain;
\item an explicit account of sector boundaries such as
      \eqref{eq:v6v72-density-ceiling}, with no integration over matter
      configurations outside the constraint domain;
\item a lower bound or positive transfer construction after the trace
      momentum and conformal factor are treated, rather than before; and
\item compatibility of the sector union with the dRNCC source topology,
      gluing maps, and the ADM two-polarization quotient of V61.
\end{enumerate}
\end{definition}

\begin{theorem}[The fixed-sector part of NCCSC is nonvacuous but bounded]
\label{thm:v6v72-NCCSC-sector}
The homogeneous zero-TT branch satisfies NCCSC\textup{(NC1)}--
\textup{(NC3)} exactly through
Theorem~\ref{thm:v6v72-homogeneous-classification}.  The constant nonzero-TT
branch satisfies those rows below the strict density ceiling through
Proposition~\ref{prop:v6v72-TT-branch}.  Neither branch supplies
NCCSC\textup{(NC4)} or a measure on unbounded Higgs amplitudes.
\end{theorem}

\begin{proof}
The cited results explicitly declare the sector, solve its reduced
constraint, and identify its boundary.  Proposition~\ref{prop:v6v72-trace-transfer}
shows why positivity does not follow.
\end{proof}

\begin{firewall}[Prescribed density is not yet the conformally coupled
Standard Model]
\label{fw:v6v72-prescribed-density}
Equation \eqref{eq:v6v72-Lichnerowicz} treats \(\rho\) as a prescribed
physical density.  For an actual scalar field, gradient, canonical-momentum,
gauge, fermion, and improvement terms carry different powers of \(\psi\),
and the momentum constraint is coupled.  The homogeneous potential branch
is exact because all gradients and matter momenta vanish.  The general
Standard Model conformal constraint remains open.
\end{firewall}

\begin{firewall}[Lorentzian constraint solvability is not Euclidean OS
positivity]
\label{fw:v6v72-no-OS}
The Hamiltonian constraint characterizes canonical initial data.  It does
not define a Euclidean probability density, a reflection-positive lapse
integral, or a contour for the conformal factor.  The sector identities
cannot be used as an OS theorem without a separate transfer construction.
\end{firewall}

\begin{theorem}[Sixth-series V72 result]
\label{thm:v6v72-status}
On a flat compact slice, the nonlinear CMC Hamiltonian constraint is the
Lichnerowicz equation \eqref{eq:v6v72-Lichnerowicz} and obeys the exact
integrated identity \eqref{eq:v6v72-integrated-identity}.  A time-symmetric
nonnegative source is excluded.  With zero TT momentum, a homogeneous Higgs
amplitude is admitted exactly when
\(\tau^2=3\Lambda+24\pi GV(R)\), forcing \(|\tau|\sim R^2\); with constant
nonzero TT momentum, the solution exists only below a finite density ceiling
and its volume diverges at that ceiling.  Thus nonlinear constraint solving
replaces the V71 negative-octic ray by a sector-existence boundary and an
unbounded negative DeWitt trace direction.  NCCSC records the remaining
constraint and transfer requirements.
\end{theorem}

\begin{proof}
Use Lemma~\ref{lem:v6v72-Lichnerowicz},
Theorem~\ref{thm:v6v72-integrated-identity},
Theorem~\ref{thm:v6v72-homogeneous-classification},
Corollary~\ref{cor:v6v72-Higgs-CMC},
Proposition~\ref{prop:v6v72-TT-branch}, and
Proposition~\ref{prop:v6v72-trace-transfer}, subject to the two firewalls.
\end{proof}

V73 will insert the actual time-symmetric scalar density into the conformal
equation.  It will track separately the \(\psi\)-weights of the gradient and
potential energies, derive the resulting semilinear equation, and test
existence and a priori bounds by sub/supersolutions and the integrated
constraint before adding gauge and fermion sources.

\clearpage
\section{V73: the minimally coupled scalar constraint and conformal-volume escape}
\label{sec:v6v73}

V72 treated the physical matter density as prescribed and isolated the
fixed-sector potential ceiling.  This version inserts the actual
time-symmetric minimally coupled scalar density.  Gradient and potential
energies acquire different conformal weights.  Below the potential ceiling,
a nowhere-vanishing TT momentum supplies a positive subsolution and the
quintic CMC term supplies a supersolution, giving a positive conformal
factor with explicit bounds.  Applied to the V71 high-frequency packet, the
solutions exist at every frequency but their physical volumes diverge.  The
linearized frequency instability has therefore become a nonlinear loss of
geometric compactness.

\subsection{Exact scalar conformal weights}

Retain the flat reference torus and CMC data of V72.  Let \(\phi\) be a
smooth real scalar of conformal weight zero, set its canonical momentum to
zero, and take a smooth potential \(V(\phi)\).  Its physical energy density
is
\begin{equation}
 \rho_\phi(g)
 ={1\over2}|\nabla\phi|_g^2+V(\phi)
 ={1\over2}\psi^{-4}|\bar\nabla\phi|_{\bar g}^2+V(\phi).
 \label{eq:v6v73-scalar-density}
\end{equation}

\begin{lemma}[Exact minimally coupled scalar Lichnerowicz equation]
\label{lem:v6v73-scalar-Lichnerowicz}
The Hamiltonian constraint is equivalent to
\begin{equation}
 -8\bar\Delta\psi
 -|\bar\sigma|^2\psi^{-7}
 -8\pi G|\bar\nabla\phi|^2\psi
 +\bigl(b_\tau-16\pi GV(\phi)\bigr)\psi^5=0.
 \label{eq:v6v73-scalar-Lichnerowicz}
\end{equation}
\end{lemma}

\begin{proof}
Insert \eqref{eq:v6v73-scalar-density} into
\eqref{eq:v6v72-Lichnerowicz}.  Multiplication of the physical gradient
term by \(16\pi G\psi^5\) gives
\(8\pi G|\bar\nabla\phi|^2\psi\), whereas the potential contributes
\(16\pi GV(\phi)\psi^5\).
\end{proof}

For brevity put
\begin{equation}
 s(x)=|\bar\sigma(x)|^2,
 \qquad
 w(x)=8\pi G|\bar\nabla\phi(x)|^2,
 \qquad
 a(x)=b_\tau-16\pi GV(\phi(x)).
 \label{eq:v6v73-coefficients}
\end{equation}
Then the equation is
\begin{equation}
 \mathcal L(\psi):=-8\bar\Delta\psi-s\psi^{-7}-w\psi+a\psi^5=0.
 \label{eq:v6v73-abstract-equation}
\end{equation}

\begin{theorem}[Exact integrated scalar constraint]
\label{thm:v6v73-integrated}
Every positive solution satisfies
\begin{equation}
 \int a\psi^5,d\mu_{\bar g}
 =\int s\psi^{-7},d\mu_{\bar g}
  +\int w\psi,d\mu_{\bar g}.
 \label{eq:v6v73-integrated}
\end{equation}
In particular, if \(a\le0\) everywhere and either \(s\) or \(w\) is
nonzero, no positive solution exists.
\end{theorem}

\begin{proof}
Integrate \eqref{eq:v6v73-abstract-equation} on the torus.  The Laplacian
has zero integral, while both terms on the right of
\eqref{eq:v6v73-integrated} are nonnegative.
\end{proof}

The potential ceiling from V72 remains exact.  Scalar gradients do not
alter its sign because they occur on the positive right side of
\eqref{eq:v6v73-integrated}.

\subsection{Positive solutions below the potential ceiling}

Assume henceforth that the smooth coefficients obey
\begin{equation}
 0<s_0\le s(x)\le s_1,
 \qquad
 0<a_0\le a(x)\le a_1,
 \qquad
 0\le w(x)\le w_1.
 \label{eq:v6v73-uniform-coefficients}
\end{equation}
The condition \(a_0>0\) is the strict pointwise potential ceiling
\begin{equation}
 b_\tau>16\pi G\sup_xV(\phi(x)).
 \label{eq:v6v73-potential-ceiling}
\end{equation}

\begin{lemma}[Explicit constant barriers]
\label{lem:v6v73-barriers}
Define
\begin{align}
 c_-&=\left({s_0\over2a_1}\right)^{1/12},
 \label{eq:v6v73-lower-barrier}\\
 c_+&=\max\left\{
 c_-,
 \left({2s_1\over a_0}\right)^{1/12},
 \left({2w_1\over a_0}\right)^{1/4}
 \right\}.
 \label{eq:v6v73-upper-barrier}
\end{align}
Then
\begin{equation}
 \mathcal L(c_-)\le0,
 \qquad
 \mathcal L(c_+)\ge0.
 \label{eq:v6v73-barrier-signs}
\end{equation}
\end{lemma}

\begin{proof}
For the lower constant, the term \(-wc_-\) is nonpositive and
\[
 -sc_-^{-7}+ac_-^5
 \le c_-^{-7}(-s_0+a_1c_-^{12})
 =-{s_0\over2}c_-^{-7}<0.
\]
For the upper constant, its definition gives
\(s_1c_+^{-7}\le a_0c_+^5/2\) and
\(w_1c_+\le a_0c_+^5/2\).  Hence
\[
 -sc_+^{-7}-wc_++ac_+^5
 \ge-s_1c_+^{-7}-w_1c_++a_0c_+^5\ge0.
\]
\end{proof}

\begin{theorem}[Subcritical scalar conformal-factor existence]
\label{thm:v6v73-existence}
Under \eqref{eq:v6v73-uniform-coefficients}, there is a smooth positive
solution of \eqref{eq:v6v73-scalar-Lichnerowicz} satisfying
\begin{equation}
 c_-\le\psi\le c_+.
 \label{eq:v6v73-existence-bounds}
\end{equation}
\end{theorem}

\begin{proof}
On the compact interval \([c_-,c_+]\), the nonlinearity
\(-st^{-7}-wt+at^5\) is smooth with a uniformly bounded derivative.  Add a
constant multiple \(Mt\), with \(M\) larger than the negative part of that
derivative, so that the resulting right side is monotone in \(t\).  Starting
from the subsolution \(c_-\), solve successively the linear uniformly
elliptic equations with operator \(-8\bar\Delta+M\).  The maximum principle,
Lemma~\ref{lem:v6v73-barriers}, and induction keep the iterates increasing
and below \(c_+\).  Schauder estimates give compactness in
\(C^{2,\alpha}\); the monotone limit solves
\eqref{eq:v6v73-abstract-equation}.  Elliptic bootstrapping makes it smooth.
The lower barrier makes the singular power harmless throughout.
\end{proof}

\begin{proposition}[A priori bounds for every positive solution]
\label{prop:v6v73-apriori}
Every smooth positive solution under
\eqref{eq:v6v73-uniform-coefficients} obeys
\begin{align}
 \min\psi&\ge\left({s_0\over a_1}\right)^{1/12},
 \label{eq:v6v73-apriori-min}\\
 \max\psi&\le
 \max\left\{
 \left({2s_1\over a_0}\right)^{1/12},
 \left({2w_1\over a_0}\right)^{1/4}
 \right\}.
 \label{eq:v6v73-apriori-max}
\end{align}
\end{proposition}

\begin{proof}
At a minimum, \(\bar\Delta\psi\ge0\).  Rearranging the equation gives
\(a\psi^5=s\psi^{-7}+w\psi+8\bar\Delta\psi\), hence
\(a_1(\min\psi)^{12}\ge s_0\).  At a maximum,
\(\bar\Delta\psi\le0\), so
\[
 a_0(\max\psi)^{12}
 \le s_1+w_1(\max\psi)^8.
\]
If \((\max\psi)^4<2w_1/a_0\), the second bound follows directly.
Otherwise the last term is at most
\(a_0(\max\psi)^{12}/2\), giving
\((\max\psi)^{12}\le2s_1/a_0\).
\end{proof}

\begin{assessment}[What the existence theorem uses]
\label{ass:v6v73-existence-input}
The TT momentum floor controls \(\psi\downarrow0\), the strict CMC/potential
gap controls \(\psi\uparrow\infty\), and the scalar gradient changes only
the upper barrier.  Thus the nonlinear equation has a genuine subcritical
sector, but its bounds are not uniform when the potential gap closes or the
gradient frequency diverges.
\end{assessment}

\subsection{High-frequency packets force physical-volume escape}

Return to \(\phi_N\) from \eqref{eq:v6v71-packet} at fixed amplitude
\(A\).  Assume the CMC sector is chosen so that
\begin{equation}
 0<a_0\le b_\tau-16\pi GV(\phi_N(x))\le a_1
 \label{eq:v6v73-packet-gap}
\end{equation}
with constants independent of \(N\), and retain the TT bounds
\(0<s_0\le s\le s_1\).  This is possible whenever the fixed amplitude lies
strictly below the potential ceiling, because the range of \(\phi_N\) is
independent of \(N\).

\begin{theorem}[Nonlinear high-frequency conformal-volume lower bound]
\label{thm:v6v73-volume-escape}
Let \(\psi_N\) be any positive solution of the scalar constraint with
source \(\phi_N\).  There is a constant \(C>0\), independent of \(N\),
such that
\begin{equation}
 \operatorname{Vol}(g_N)
 =\int_{\mathbb T^3}\psi_N^6d\mu_{\bar g}
 \ge C N^{12/5}.
 \label{eq:v6v73-volume-lower}
\end{equation}
The a priori upper bound also gives
\begin{equation}
 \operatorname{Vol}(g_N)\le C'(1+N^3).
 \label{eq:v6v73-volume-upper}
\end{equation}
Thus the family has no fixed-volume compact limit inside this CMC sector.
\end{theorem}

\begin{proof}
Proposition~\ref{prop:v6v73-apriori} gives the uniform lower bound
\(\psi_N\ge m_0=(s_0/a_1)^{1/12}\).  By
\eqref{eq:v6v71-H1},
\begin{equation}
 \int w_Nd\mu=8\pi GA^2N^2.
 \label{eq:v6v73-gradient-total}
\end{equation}
The integrated constraint \eqref{eq:v6v73-integrated} therefore implies
\[
 a_1\int\psi_N^5d\mu
 \ge\int w_N\psi_Nd\mu
 \ge8\pi Gm_0A^2N^2.
\]
On the normalized torus, monotonicity of \(L^p\) norms yields
\[
 \int\psi_N^6d\mu
 \ge\left(\int\psi_N^5d\mu\right)^{6/5},
\]
which proves \eqref{eq:v6v73-volume-lower}.  Finally,
\(w_1\le16\pi GA^2N^2\).  The maximum bound
\eqref{eq:v6v73-apriori-max} gives
\(\max\psi_N\le C(1+N^{1/2})\), and its sixth power proves
\eqref{eq:v6v73-volume-upper}.
\end{proof}

\begin{corollary}[Frequency energy is transferred into geometry]
\label{cor:v6v73-frequency-transfer}
The nonlinear constraint does not yield a regulator-uniform conformal-factor
family for the V71 packet.  At fixed reference volume and fixed CMC/potential
gap, the growing scalar gradient forces the physical volume to escape at
least as \(N^{12/5}\).
\end{corollary}

\begin{proof}
Theorem~\ref{thm:v6v73-volume-escape}.
\end{proof}

\subsection{A variational identity for the constraint equation}

Equation \eqref{eq:v6v73-abstract-equation} is the Euler equation of
\begin{equation}
 \mathcal I_{\phi,\tau,\sigma}(\psi)
 =\int_{\mathbb T^3}\left[
 4|\bar\nabla\psi|^2+{s\over6}\psi^{-6}
 -{w\over2}\psi^2+{a\over6}\psi^6
 \right]d\mu_{\bar g}.
 \label{eq:v6v73-constraint-functional}
\end{equation}

\begin{proposition}[Coercivity inside the strict subcritical sector]
\label{prop:v6v73-functional-coercive}
Under \eqref{eq:v6v73-uniform-coefficients}, the functional
\eqref{eq:v6v73-constraint-functional} is bounded below and coercive in
\(H^1\cap L^6\) on positive functions with \(\psi^{-1}\in L^6\).
Every smooth critical point satisfies
\eqref{eq:v6v73-scalar-Lichnerowicz}.
\end{proposition}

\begin{proof}
Young's inequality gives, for a constant depending on \(a_0,w_1\),
\[
 {a\over6}\psi^6-{w\over2}\psi^2
 \ge {a_0\over12}\psi^6-C.
\]
The gradient and singular terms are nonnegative, and
\(s\ge s_0>0\), proving the lower bound and coercivity.  Differentiation in
a smooth compactly supported direction gives respectively
\(-8\bar\Delta\psi\), \(-s\psi^{-7}\), \(-w\psi\), and
\(a\psi^5\), which is \eqref{eq:v6v73-abstract-equation}.
\end{proof}

\begin{firewall}[Constraint-functional coercivity is not the physical
Euclidean action]
\label{fw:v6v73-functional-scope}
The auxiliary functional \eqref{eq:v6v73-constraint-functional} proves
elliptic solvability.  It is not the Einstein--Hilbert action after gauge
fixing, and its positive \(a\psi^6\) term does not resolve the DeWitt
conformal sign in Proposition~\ref{prop:v6v72-trace-transfer}.  No
reflection-positive metric measure follows from its coercivity.
\end{firewall}

\begin{firewall}[Minimal time-symmetric scalar scope]
\label{fw:v6v73-minimal-scope}
This version sets scalar canonical momentum and matter momentum density to
zero and uses minimal coupling.  A nonminimal \(R\phi^2\) term, Higgs gauge
covariant gradients, fermions, gauge electric and magnetic fields, and the
coupled momentum constraint change the coefficient system.  They must be
inserted as complete Standard Model sources before NCCSC is claimed.
\end{firewall}

\begin{definition}[Subcritical scalar conformal card]
\label{def:v6v73-SSCC}
SSCC consists of:
\begin{enumerate}[label=\textup{(SC\arabic*)},leftmargin=5.0em]
\item the exact weighted equation
      \eqref{eq:v6v73-scalar-Lichnerowicz};
\item a positive TT momentum floor and strict CMC/potential gap
      \eqref{eq:v6v73-uniform-coefficients};
\item positive conformal-factor existence with the a priori bounds
      \eqref{eq:v6v73-apriori-min}--\eqref{eq:v6v73-apriori-max}; and
\item a regulator-uniform physical-volume or normalized-volume replacement
      compatible with the source topology.
\end{enumerate}
\end{definition}

\begin{theorem}[The analytic SSCC rows close and the compactness row fails]
\label{thm:v6v73-SSCC}
Rows \textup{(SC1)}--\textup{(SC3)} hold for the minimally coupled
time-symmetric scalar in the strict subcritical sector.  Row
\textup{(SC4)} fails for the fixed-CMC V71 packet family by
\eqref{eq:v6v73-volume-lower}.
\end{theorem}

\begin{proof}
Use Lemma~\ref{lem:v6v73-scalar-Lichnerowicz},
Theorem~\ref{thm:v6v73-existence},
Proposition~\ref{prop:v6v73-apriori}, and
Theorem~\ref{thm:v6v73-volume-escape}.
\end{proof}

\begin{theorem}[Sixth-series V73 result]
\label{thm:v6v73-status}
The minimally coupled scalar constraint is
\eqref{eq:v6v73-scalar-Lichnerowicz}: gradient energy is linear in the
conformal factor, whereas potential energy shares the quintic CMC weight.
With a positive TT momentum floor and a strict potential ceiling, explicit
constant barriers give a smooth positive solution and uniform pointwise
bounds in terms of the coefficient extrema.  For the high-frequency packet
\(\phi_N\), solutions exist at every \(N\), but their physical volume grows
at least as \(N^{12/5}\).  Nonlinear constraint solving therefore converts
the linearized frequency instability into conformal-volume escape; the
compactness row of SSCC remains open.
\end{theorem}

\begin{proof}
Use Theorem~\ref{thm:v6v73-SSCC}, subject to the two firewalls.
\end{proof}

V74 will treat York time as a sector variable rather than hold it fixed.
It will impose a fixed physical-volume normalization, use the integrated
constraint to derive the necessary growth of \(\tau_N\) for the packet,
and test whether the resulting trace-momentum cost gives a compact
normalized family or merely moves the escape into the CMC label.

\clearpage
\section{V74: fixed-volume York-time escape and the conformal-collapse dichotomy}
\label{sec:v6v74}

V73 shows that fixed-CMC solutions for high-frequency scalar packets escape
through physical volume.  This version fixes that volume and instead solves
the integrated constraint for the York-time label.  A nonvanishing TT
momentum floor forces York time to diverge at a quantitative rate; if the
conformal metrics remain uniformly nondegenerate, the rate is linear in the
packet frequency.  Conversely, keeping York time bounded forces scalar
gradient energy into regions where the conformal factor collapses.  Thus
volume normalization moves the lack of compactness into either the CMC
sector label or metric degeneration.

\subsection{York time from the integrated constraint}

For a positive solution of \eqref{eq:v6v73-scalar-Lichnerowicz}, write
\begin{equation}
 I_5(\psi)=\int_{\mathbb T^3}\psi^5d\mu_{\bar g},
 \qquad
 \mathcal V(\psi)=\int_{\mathbb T^3}\psi^6d\mu_{\bar g}.
 \label{eq:v6v74-I5-volume}
\end{equation}

\begin{theorem}[Exact York-time identity]
\label{thm:v6v74-York-identity}
Every positive solution obeys
\begin{align}
 \tau^2
 =3\Lambda+{3\over2I_5(\psi)}
 \biggl[&\int |\bar\sigma|^2\psi^{-7}d\mu_{\bar g}
 +8\pi G\int|\bar\nabla\phi|^2\psi d\mu_{\bar g}
 \nonumber\\
 &+16\pi G\int V(\phi)\psi^5d\mu_{\bar g}\biggr].
 \label{eq:v6v74-York-identity}
\end{align}
\end{theorem}

\begin{proof}
Expand \(a=b_\tau-16\pi GV(\phi)\) in the integrated identity
\eqref{eq:v6v73-integrated} and solve for
\(b_\tau=2\tau^2/3-2\Lambda\).
\end{proof}

\begin{corollary}[A bounded York interval gives a bounded weighted gradient
budget]
\label{cor:v6v74-bounded-York-budget}
Assume \(V(\phi)\ge-V_-\).  Then every positive solution satisfies
\begin{equation}
 8\pi G\int|\bar\nabla\phi|^2\psi d\mu_{\bar g}
 \le\left({2\tau^2\over3}+2|\Lambda|+16\pi GV_-\right)I_5(\psi).
 \label{eq:v6v74-weighted-gradient-budget}
\end{equation}
If \(\mathcal V(\psi)=\mathcal V_*\), then
\begin{equation}
 I_5(\psi)\le\mathcal V_*^{5/6}.
 \label{eq:v6v74-I5-bound}
\end{equation}
\end{corollary}

\begin{proof}
Drop the nonnegative TT term from
\eqref{eq:v6v74-York-identity}, move the negative part of the potential to
the other side, and use \(b_\tau\le2\tau^2/3+2|\Lambda|\).  The second
statement is monotonicity of \(L^p\) norms on the normalized reference
torus.
\end{proof}

\subsection{Quantitative York-time escape for the packet}

Let \(\phi_N\) be the V71 packet with fixed \(A\ne0\).  Since its range is
the fixed interval \([-2|A|,2|A|]\), define finite constants
\begin{equation}
 V_- =\max_{|u|\le2|A|}(-V(u))_+,
 \qquad
 V_+ =\max_{|u|\le2|A|}|V(u)|.
 \label{eq:v6v74-potential-bounds}
\end{equation}

\begin{lemma}[York-dependent lower conformal bound]
\label{lem:v6v74-York-lower-psi}
Suppose \(|\bar\sigma|^2\ge s_0>0\), and let \(\psi_N\) be any positive
solution with York time \(\tau_N\).  There is a constant \(c_0>0\),
independent of \(N\) and \(\tau_N\), such that
\begin{equation}
 \min\psi_N\ge c_0(1+\tau_N^2)^{-1/12}.
 \label{eq:v6v74-York-lower-psi}
\end{equation}
\end{lemma}

\begin{proof}
At a minimum of \(\psi_N\), the argument of
Proposition~\ref{prop:v6v73-apriori} gives
\[
 a(x_{\min})(\min\psi_N)^{12}\ge s_0.
\]
The coefficient is bounded above by
\[
 a(x)={2\tau_N^2\over3}-2\Lambda-16\pi GV(\phi_N(x))
 \le C_{\Lambda,G,V_+}(1+\tau_N^2).
\]
Combining the two inequalities proves the claim.
\end{proof}

\begin{theorem}[Fixed-volume York-time lower rate]
\label{thm:v6v74-York-escape}
Assume \(|\bar\sigma|^2\ge s_0>0\), and suppose positive solutions
\(\psi_N\) exist with one fixed physical volume
\begin{equation}
 \mathcal V(\psi_N)=\mathcal V_*>0.
 \label{eq:v6v74-fixed-volume}
\end{equation}
Then there is \(c>0\), independent of \(N\), such that
\begin{equation}
 1+|\tau_N|\ge cN^{12/13}.
 \label{eq:v6v74-York-rate}
\end{equation}
In particular, no bounded CMC interval contains the complete fixed-volume
packet family.
\end{theorem}

\begin{proof}
By \eqref{eq:v6v71-H1} and
Lemma~\ref{lem:v6v74-York-lower-psi},
\begin{equation}
 8\pi G\int|\bar\nabla\phi_N|^2\psi_Nd\mu
 \ge c_1N^2(1+\tau_N^2)^{-1/12}.
 \label{eq:v6v74-gradient-lower}
\end{equation}
Corollary~\ref{cor:v6v74-bounded-York-budget} and
\eqref{eq:v6v74-I5-bound} bound the same expression above by
\(c_2(1+\tau_N^2)\mathcal V_*^{5/6}\).  Hence
\begin{equation}
 (1+\tau_N^2)^{13/12}\ge c_3N^2.
 \label{eq:v6v74-rate-intermediate}
\end{equation}
Taking the power \(6/13\) gives
\(1+|\tau_N|\ge cN^{12/13}\), after changing the constant.
\end{proof}

\begin{corollary}[Linear escape under metric nondegeneracy]
\label{cor:v6v74-linear-York}
If, in addition, \(\psi_N\ge c_*>0\) uniformly, then
\begin{equation}
 1+|\tau_N|\ge c'N.
 \label{eq:v6v74-linear-York}
\end{equation}
\end{corollary}

\begin{proof}
Replace \eqref{eq:v6v74-gradient-lower} by
\(8\pi Gc_*A^2N^2\).  The upper bound is still
\(C(1+\tau_N^2)\mathcal V_*^{5/6}\).
\end{proof}

\begin{assessment}[The rate is a necessary sector cost]
\label{ass:v6v74-rate-cost}
The exponent \(12/13\) comes from the weakest conformal lower bound forced
by the TT momentum floor.  It is not asserted to be sharp.  What is sharp
for the programme is the qualitative conclusion: fixed volume and
nonvanishing TT momentum require an unbounded York-time ancestry as the
matter frequency is removed.
\end{assessment}

\subsection{Bounded York time forces conformal collapse on the gradient
carrier}

The opposite branch can be stated without a TT momentum floor.

\begin{proposition}[Gradient-weighted conformal collapse]
\label{prop:v6v74-gradient-collapse}
Let \(\psi_N\) be positive solutions with fixed physical volume, and assume
\(|\tau_N|\le T\).  Then
\begin{equation}
 {\int |\bar\nabla\phi_N|^2\psi_Nd\mu
  \over\int |\bar\nabla\phi_N|^2d\mu}
 \le {C_T\over N^2}.
 \label{eq:v6v74-weighted-mean-collapse}
\end{equation}
For every \(\delta>0\),
\begin{equation}
 {\int_{\{\psi_N>\delta\}}|\bar\nabla\phi_N|^2d\mu
  \over\int |\bar\nabla\phi_N|^2d\mu}
 \le {C_T\over\delta N^2}\longrightarrow0.
 \label{eq:v6v74-gradient-fraction-collapse}
\end{equation}
Thus asymptotically all scalar gradient energy is carried where
\(\psi_N\) is arbitrarily small.
\end{proposition}

\begin{proof}
With \(|\tau_N|\le T\), the right side of
\eqref{eq:v6v74-weighted-gradient-budget} is uniformly bounded by
\eqref{eq:v6v74-I5-bound}.  Divide by
\(8\pi G\int|\bar\nabla\phi_N|^2d\mu=8\pi GA^2N^2\) to get the first
claim.  On \(\{\psi_N>\delta\}\),
\[
 \delta\int_{\{\psi_N>\delta\}}|\bar\nabla\phi_N|^2d\mu
 \le\int|\bar\nabla\phi_N|^2\psi_Nd\mu.
\]
Divide again by the total gradient energy.
\end{proof}

\begin{corollary}[No bounded-York nondegenerate fixed-volume family]
\label{cor:v6v74-no-compact-family}
For the packet family, fixed physical volume, bounded York time, and a
uniform positive lower bound on the conformal factor are mutually
incompatible.
\end{corollary}

\begin{proof}
If \(\psi_N\ge c_*>0\), the left side of
\eqref{eq:v6v74-weighted-mean-collapse} is at least \(c_*\), contradicting
its convergence to zero.
\end{proof}

\subsection{The trace-momentum ledger at fixed volume}

The integrated pure-trace DeWitt kinetic combination is
\begin{equation}
 \mathcal D_{\rm tr}(\tau,\psi)
 =\int_{\mathbb T^3}(K_{ij}K^{ij}-K^2)d\mu_g
 =-{2\tau^2\over3}\mathcal V(\psi).
 \label{eq:v6v74-trace-ledger}
\end{equation}

\begin{corollary}[Trace-direction escape]
\label{cor:v6v74-trace-escape}
Under the assumptions of Theorem~\ref{thm:v6v74-York-escape},
\begin{equation}
 |\mathcal D_{\rm tr}(\tau_N,\psi_N)|
 \ge c\mathcal V_*N^{24/13}-C\mathcal V_*.
 \label{eq:v6v74-trace-escape}
\end{equation}
Under the nondegeneracy assumption of
Corollary~\ref{cor:v6v74-linear-York}, the lower power improves to \(N^2\).
\end{corollary}

\begin{proof}
Insert \eqref{eq:v6v74-York-rate}, or
\eqref{eq:v6v74-linear-York}, into
\eqref{eq:v6v74-trace-ledger}.
\end{proof}

This ledger is signed.  It quantifies the sector label that escapes; it does
not define a positive Euclidean cost.

\subsection{Fixed-volume York cofinality card}

\begin{definition}[Fixed-volume York cofinality card]
\label{def:v6v74-FVYCC}
FVYCC requires a regulator family of constrained initial-data sectors with:
\begin{enumerate}[label=\textup{(Y\arabic*)},leftmargin=4.5em]
\item one fixed or tightly controlled physical-volume normalization;
\item quantitative control of the York-time labels and conformal factors in
      a topology strong enough to exclude the two alternatives
      \eqref{eq:v6v74-York-rate} and
      \eqref{eq:v6v74-gradient-fraction-collapse};
\item a positive transfer or lower-bound mechanism for the unbounded trace
      ledger \eqref{eq:v6v74-trace-ledger}; and
\item compatibility with refinement, gluing, the momentum constraint, and
      the complete SM source atlas.
\end{enumerate}
\end{definition}

\begin{theorem}[The minimal scalar packet violates naive FVYCC compactness]
\label{thm:v6v74-FVYCC-no-go}
For the V71 packet at fixed nonzero amplitude and fixed physical volume,
every constrained family either has unbounded York labels or has conformal
collapse on a set carrying asymptotically all scalar gradient energy.  With
a TT momentum floor, the first alternative is mandatory and satisfies
\eqref{eq:v6v74-York-rate}.  Hence FVYCC cannot be proved by volume
normalization alone.
\end{theorem}

\begin{proof}
If York time is bounded, apply
Proposition~\ref{prop:v6v74-gradient-collapse}.  If a TT momentum floor is
present, Theorem~\ref{thm:v6v74-York-escape} rules out bounded York time.
\end{proof}

\begin{firewall}[A deterministic packet does not determine path-integral
probability]
\label{fw:v6v74-packet-probability}
The theorem gives a necessary response to one explicit source family.  It
does not prove that a putative SM measure assigns that family positive
probability or that typical fields have the same escape rate.  Any claimed
uniform continuum construction must nevertheless control every regulated
configuration in its stated coercive domain, so the packet is a valid
counterexample to a global deterministic bound.
\end{firewall}

\begin{firewall}[No trace contour or sector measure]
\label{fw:v6v74-no-sector-measure}
FVYCC does not specify an integration contour for York time, a Faddeev--Popov
factor, or a probability law on CMC sectors.  The negative DeWitt ledger
cannot be converted into a decaying weight by taking its absolute value.
Those data remain part of the nonlinear Euclidean transfer problem.
\end{firewall}

\begin{theorem}[Sixth-series V74 result]
\label{thm:v6v74-status}
The integrated minimally coupled constraint determines York time by
\eqref{eq:v6v74-York-identity}.  At fixed physical volume, the high-frequency
packet with a TT momentum floor forces
\(1+|\tau_N|\gtrsim N^{12/13}\); uniform metric nondegeneracy strengthens
this to linear growth.  Bounded York time instead forces almost all gradient
energy into the region where the conformal factor collapses.  The associated
negative trace ledger diverges at least as \(N^{24/13}\).  Volume
normalization therefore moves, but does not remove, the nonlinear
compactness obstruction; FVYCC remains open.
\end{theorem}

\begin{proof}
Use Theorem~\ref{thm:v6v74-York-identity},
Theorem~\ref{thm:v6v74-York-escape},
Corollary~\ref{cor:v6v74-linear-York},
Proposition~\ref{prop:v6v74-gradient-collapse}, and
Corollary~\ref{cor:v6v74-trace-escape}, subject to the two firewalls.
\end{proof}

V75 will perform the compulsory YM/NS companion audit.  It will compare
their newest all-arity and common-scale mechanisms with the York-sector
escape, then decide whether the next attack should seek a cofinal sector
weight, a volume-normalized conformal compactness theorem, or an EFT cutoff
that explicitly excludes the packet obstruction.

\clearpage
\section{V75: companion audit and the conformal trace contour fork}
\label{sec:v6v75}

V74 proves that fixed-volume high-frequency sources escape through York time
or conformal collapse.  This compulsory audit shows why an arbitrary
moment estimate on the sector label is not the next step.  The Yang--Mills
companion has found an exact exponentially small wrapped branch whose
high derivatives nevertheless destroy a uniform moment card; the
Navier--Stokes companion has consolidated several cases in which a
normalization or rate along a sequence moves the compactness window without
reaching the limiting boundary.  For gravity, the complete York/conformal
sector must likewise be retained.  Before it can be weighted, the negative
DeWitt trace direction requires an explicit contour or a canonical
constraint quotient.

\subsection{Compulsory V75 audit}

The companion snapshots read for this audit were
\begin{center}
\begin{tabular}{p{0.22\linewidth}p{0.18\linewidth}p{0.49\linewidth}}
programme & source & audited endpoint\\ \hline
Yang--Mills & \texttt{V52} & the exact Hartman--Watson transform splits
into an entire parity sector and a square-root wrapped branch of weight
\(\delta_\alpha\asymp e^{-1/s}\); its high moments have the
\(r!r^{-3/2}\) law and cross the local thermal expansion at
\(r\asymp(s\log(1/s))^{-1}\), ruling out blockwise absolute
subordination while leaving a branch-resolved complete-source card open\\
Navier--Stokes & \texttt{V70} & rate-free results are organized into
type-I Liouville, intermediate, and strong type-II gates; local regularity
at the singular time needs a rate, and normalization along shorter windows
does not produce compactness at the endpoint
\end{tabular}
\end{center}
At the consistent snapshots used here, YM V52 had \(25\,280\) logical
lines, \(798\,567\) bytes, and SHA--256
\begin{equation}
 \texttt{674d2758caf95893ee5534a455afba81fe3e1f7ad1cf1a9d5742ce1c1aa21ff4},
 \label{eq:v6v75-YM-hash}
\end{equation}
while NS V70 had \(34\,157\) logical lines, \(1\,420\,002\) bytes, and
SHA--256
\begin{equation}
 \texttt{6aac28fe849c169a80e38f1a93ce497aabd8780545790a84cec83ec5aba4c2e6}.
 \label{eq:v6v75-NS-hash}
\end{equation}
The files were read only.  They were changing concurrently before these
snapshots, so the recorded hashes, rather than a later filename inventory,
define the audit evidence.

\begin{theorem}[V75 audit transfer]
\label{thm:v6v75-audit-transfer}
The applicable conclusions are structural:
\begin{enumerate}[label=\textup{(A\arabic*)},leftmargin=4.5em]
\item the escaping York/conformal branch must be retained as one complete
      sector until its exact weight and constraint cancellations are formed;
\item bounds on every fixed York moment cannot replace a common cofinal
      sector estimate;
\item fixed-volume normalization does not prove endpoint compactness, since
      V74 moves the escape into York time or conformal degeneration; and
\item no sector-weight attack is meaningful until the negative DeWitt trace
      direction is assigned either a convergent contour or a canonical
      constraint interpretation.
\end{enumerate}
No Bessel constant, Navier--Stokes rate, or companion theorem is imported
as a gravitational estimate.
\end{theorem}

\begin{proof}
YM V52's exact branch has exponentially small weight but a fixed analytic
radius, so differentiation before complete-source assembly loses the
nonperturbative factor.  NS V70 states explicitly that its sequential rate
information does not reach the singular endpoint.  V74 proves the analogous
gravity endpoint alternatives.  Finally,
\eqref{eq:v6v72-DeWitt-trace} is negative, so it cannot be the quadratic
action of a real positive Gaussian measure.
\end{proof}

\subsection{No positive real Gaussian for the DeWitt trace direction}

The obstruction is already one-dimensional.  Let \(a>0\) represent the
magnitude of one negative trace eigenvalue and let \(J\in\mathbb R\) be a
matter source.

\begin{theorem}[Real conformal Gaussian divergence]
\label{thm:v6v75-real-divergence}
The real-contour integral
\begin{equation}
 Z_{\mathbb R}(J)
 =\int_{\mathbb R}
   \exp\left({a\over2}y^2-Jy\right)dy
 \label{eq:v6v75-real-integral}
\end{equation}
diverges for every real \(J\).  In particular, a linear source cannot make
the negative quadratic trace direction normalizable.
\end{theorem}

\begin{proof}
Completing the square gives
\[
 {a\over2}y^2-Jy
 ={a\over2}\left(y-{J\over a}\right)^2-{J^2\over2a}.
\]
The integrand grows exponentially on both real tails.
\end{proof}

\begin{proposition}[Negative covariance cannot be a positive probability
covariance]
\label{prop:v6v75-no-positive-covariance}
There is no real positive probability measure and real square-integrable
random variable \(Y\) whose centered covariance is \(-a^{-1}\).
Consequently the inverse of a negative DeWitt block cannot be represented as
the covariance of a real positive Gaussian field.
\end{proposition}

\begin{proof}
Every real random variable satisfies
\(\operatorname{Var}(Y)=\mathbb E[(Y-\mathbb EY)^2]\ge0\), whereas
\(-a^{-1}<0\).
\end{proof}

The proposition is independent of gauge fixing.  Gauge quotienting can
remove null directions; it cannot turn a retained negative physical
quadratic coordinate into a positive covariance.

\subsection{Exact imaginary-contour completion}

Rotate the integration line to \(y=iz\), \(z\in\mathbb R\).  After
normalizing by the zero-source integral, define
\begin{equation}
 \widehat Z_{i\mathbb R}(J)
 ={\int_{\mathbb R}e^{-az^2/2-iJz}dz
   \over\int_{\mathbb R}e^{-az^2/2}dz}.
 \label{eq:v6v75-rotated-ratio}
\end{equation}

\begin{theorem}[Rotated trace characteristic functional]
\label{thm:v6v75-rotated-Gaussian}
For every real \(J\),
\begin{equation}
 \widehat Z_{i\mathbb R}(J)
 =\exp\left(-{J^2\over2a}\right).
 \label{eq:v6v75-rotated-result}
\end{equation}
Thus the rotated trace direction is convergent and induces a damping source
quadratic, with the opposite sign from the positive TT Gaussian completion
in Lemma~\ref{lem:v6v70-metric-integration}.
\end{theorem}

\begin{proof}
The numerator is the characteristic function of a centered real Gaussian
of variance \(a^{-1}\).  Equivalently, shift the integration contour in the
entire Gaussian integrand by \(iJ/a\); the vertical sides vanish at infinity,
and completion of the square yields the displayed factor.
\end{proof}

\begin{corollary}[Finite-dimensional negative-block formula]
\label{cor:v6v75-negative-block}
Let \(A\) be a positive definite \(d\)-by-\(d\) matrix.  Rotating every
coordinate of the negative quadratic block gives
\begin{equation}
 {\int_{\mathbb R^d}
  e^{-\langle z,Az\rangle/2-i\langle J,z\rangle}dz
  \over
  \int_{\mathbb R^d}e^{-\langle z,Az\rangle/2}dz}
 =e^{-\langle J,A^{-1}J\rangle/2}.
 \label{eq:v6v75-negative-block-formula}
\end{equation}
\end{corollary}

\begin{proof}
Diagonalize \(A\) orthogonally and apply
Theorem~\ref{thm:v6v75-rotated-Gaussian} to each coordinate.
\end{proof}

\begin{assessment}[The contour rotation changes the source problem]
\label{ass:v6v75-contour-sign}
For a positive Gaussian carrier, real source coupling produces
\(+\langle J,CJ\rangle/2\) in the exponential and the V70--V71 quartic
stability problem.  Rotating a negative carrier produces
\(-\langle J,A^{-1}J\rangle/2\), which is stabilizing as a scalar
functional.  This algebraic improvement is not free: the metric trace is
imaginary on the chosen contour, and the source insertion is oscillatory
before the trace variable is integrated out.
\end{assessment}

\subsection{Why convergence is not yet reflection positivity}

\begin{proposition}[A convergent complex contour is not a positive field
law]
\label{prop:v6v75-contour-not-positive-law}
The rotated integrand in \eqref{eq:v6v75-rotated-ratio} is not a positive
measure with a real trace field when \(J\ne0\).  The positivity of the final
number \eqref{eq:v6v75-rotated-result} does not by itself prove positivity
of Gram matrices for nonlinear observables containing the trace coordinate.
\end{proposition}

\begin{proof}
The factor \(e^{-iJz}\) is nonreal for generic \(z,J\), and the original
trace coordinate is \(y=iz\), not real.  A positive scalar characteristic
factor after one variable is integrated out determines only source
observables that depend on \(J\); it does not define a positive expectation
of arbitrary polynomials in \(y\).  Reflection positivity is a family of
matrix inequalities for all declared positive-time observables, which is
strictly stronger than convergence of one normalized integral.
\end{proof}

\begin{proposition}[The physical Newton sign is not licensed by contour
convergence]
\label{prop:v6v75-Newton-sign}
Equation \eqref{eq:v6v75-rotated-result} cannot be identified with the
attractive Lorentzian Newton completion
\eqref{eq:v6v66-Newton-completion} without an analytic-continuation theorem
that transports the constraint, source convention, and physical-state
inner product.  The two formulae have different derivations and source
signs.
\end{proposition}

\begin{proof}
The Newton formula is a real canonical minimization on the constraint
surface.  The rotated formula is a complex-contour evaluation of an
off-shell negative Euclidean quadratic coordinate.  Algebraic equality of
some absolute coefficients does not identify their domains, contours, or
physical inner products.
\end{proof}

\subsection{The contour-or-constraint card}

\begin{definition}[Conformal trace contour card]
\label{def:v6v75-CTCC}
CTCC requires, at every regulator:
\begin{enumerate}[label=\textup{(CT\arabic*)},leftmargin=5.2em]
\item an explicit spectral split of gauge, null, positive TT, and negative
      trace metric directions;
\item for every retained negative direction, either a specified convergent
      complex integration cycle with its orientation and source domain, or
      a canonical constraint quotient that removes it from the probabilistic
      carrier;
\item exact assembly of the induced matter functional before arity or York
      moments are taken;
\item a reflected involution and positive Gram theorem on the declared
      physical observable algebra; and
\item a cofinal estimate for the full York/conformal sector, including the
      V74 escape branch, without replacing it by absolute moments.
\end{enumerate}
\end{definition}

\begin{theorem}[Neither elementary branch closes CTCC]
\label{thm:v6v75-CTCC-status}
Real integration fails CTCC\textup{(CT2)} by
Theorem~\ref{thm:v6v75-real-divergence}.  Imaginary Gaussian rotation closes
finite-dimensional convergence and gives the exact induced functional
\eqref{eq:v6v75-negative-block-formula}, but it does not close
CTCC\textup{(CT4)} or identify the Lorentzian Newton constraint.  The
canonical constraint route has the sector-existence and cofinality problems
of V72--V74.  CTCC therefore remains open, with three sharply separated
subgates rather than one formal conformal-factor prescription.
\end{theorem}

\begin{proof}
Use Theorem~\ref{thm:v6v75-real-divergence},
Corollary~\ref{cor:v6v75-negative-block},
Propositions~\ref{prop:v6v75-contour-not-positive-law} and
\ref{prop:v6v75-Newton-sign}, and
Theorem~\ref{thm:v6v74-FVYCC-no-go}.
\end{proof}

\begin{firewall}[No unproved Picard--Lefschetz equivalence]
\label{fw:v6v75-no-thimble-equivalence}
The imaginary line is a valid contour for the isolated quadratic model.
No claim is made that it is homologous to a convergent integration cycle of
the nonlinear Einstein--Hilbert action, that nonlinear singularities do not
cross it, or that the same cycle works uniformly under refinement.  Those
are separate finite- and infinite-dimensional problems.
\end{firewall}

\begin{firewall}[The companion branch analogy is structural only]
\label{fw:v6v75-companion-scope}
The YM wrapped Bessel branch is an exact positive compact-group sector; the
gravity conformal trace is a negative quadratic direction.  They are not
the same random variable or measure.  The transferred rule is only to keep
the complete exceptional sector assembled before differentiation and
absolute values.
\end{firewall}

\begin{theorem}[Sixth-series V75 result]
\label{thm:v6v75-status}
The compulsory audit places an exact sector construction before York-moment
estimates.  A negative DeWitt trace coordinate has no positive real Gaussian
law and its real Euclidean integral diverges for every source.  Rotation to
an imaginary contour gives the exact convergent factor
\(e^{-\langle J,A^{-1}J\rangle/2}\), but does not provide a positive law for
trace observables or identify the attractive Lorentzian Newton constraint.
CTCC records the resulting contour-or-constraint fork and the still-open
reflection and cofinality rows.
\end{theorem}

\begin{proof}
Use Theorem~\ref{thm:v6v75-audit-transfer} and
Theorem~\ref{thm:v6v75-CTCC-status}, subject to the two firewalls.
\end{proof}

V76 will test the canonical alternative in a finite-dimensional constrained
model.  It will compare real-lapse group averaging, which imposes a delta
constraint, with imaginary conformal rotation, and determine the exact
conditions under which the induced physical inner product is positive and
normalizable before any field-theoretic limit is attempted.

\clearpage
\section{V76: real-lapse group averaging and the homogeneous Higgs coarea measure}
\label{sec:v6v76}

V75 leaves a contour-or-constraint fork.  This version develops the
canonical branch in finite dimensions.  Averaging the unitary flow of a
self-adjoint constraint over a real lapse produces a delta distribution,
not a Gaussian covariance.  For a regular multiplication constraint, the
resulting physical inner product is the positive coarea measure on the
constraint surface.  Applied to the homogeneous CMC--Higgs relation of V72,
the Faddeev--Popov/coarea Jacobian decays like the inverse York time.  That
decay controls one real amplitude but not a multi-component Higgs tail.

\subsection{Discrete-spectrum group averaging}

Let \(\mathcal H_{\rm kin}\) be a Hilbert space and let \(\widehat C\) be
self-adjoint.  First suppose zero is an isolated eigenvalue and write
\(P_0\) for the orthogonal projection onto \(\ker\widehat C\).

\begin{theorem}[Mean group averaging at an isolated constraint eigenvalue]
\label{thm:v6v76-mean-ergodic}
For every \(f\in\mathcal H_{\rm kin}\),
\begin{equation}
 {1\over2T}\int_{-T}^{T}e^{it\widehat C}f\,dt
 \longrightarrow P_0f
 \quad\hbox{strongly as }T\to\infty.
 \label{eq:v6v76-mean-ergodic}
\end{equation}
Consequently
\begin{equation}
 \langle f,g\rangle_{\rm phys}
 :=\langle P_0f,P_0g\rangle_{\rm kin}
 \label{eq:v6v76-discrete-physical-product}
\end{equation}
is positive and its null quotient is canonically isometric to
\(\ker\widehat C\).
\end{theorem}

\begin{proof}
By the spectral theorem, the averaged multiplier is
\[
 m_T(\lambda)={\sin(T\lambda)\over T\lambda},
 \qquad m_T(0)=1.
\]
It converges pointwise to \(\mathbf1_{\{0\}}(\lambda)\) and is bounded by
one.  Dominated convergence against the spectral measure of \(f\) proves
strong convergence.  The inner-product statement follows from orthogonality
of \(P_0\).
\end{proof}

\begin{warning}[Continuous zero spectrum needs a rigging normalization]
\label{warn:v6v76-continuous-zero}
If zero lies only in continuous spectrum, the mean in
\eqref{eq:v6v76-mean-ergodic} converges strongly to zero.  A nontrivial
constraint-surface inner product then requires a distributional delta
normalization.  Treating the zero result as absence of physical states would
discard the generalized constraint shell.
\end{warning}

\subsection{Gaussian-regulated real-lapse averaging}

For \(\epsilon>0\), define
\begin{equation}
 \delta_\epsilon(s)
 ={1\over\sqrt{2\pi\epsilon}}e^{-s^2/(2\epsilon)}
 ={1\over2\pi}\int_{\mathbb R}
   e^{its}e^{-\epsilon t^2/2}\,dt.
 \label{eq:v6v76-delta-regularizer}
\end{equation}
Let \(c\in C^\infty(\mathbb R^n;\mathbb R)\), let
\(\widehat C\) act by multiplication by \(c(x)\) on
\(L^2(\mathbb R^n,dx)\), and set
\begin{equation}
 \eta_\epsilon(f,g)
 =\langle f,\delta_\epsilon(\widehat C)g\rangle
 =\int_{\mathbb R^n}\overline{f(x)}g(x)
   \delta_\epsilon(c(x))\,dx.
 \label{eq:v6v76-regularized-rigging}
\end{equation}

\begin{theorem}[Positive coarea rigging limit]
\label{thm:v6v76-coarea-rigging}
Suppose
\begin{equation}
 \Sigma=c^{-1}(0),qquad \nabla c\ne0\quad\hbox{on }\Sigma,
 \label{eq:v6v76-regular-shell}
\end{equation}
and take \(f,g\in C_c^\infty(\mathbb R^n)\).  Then
\begin{equation}
 \lim_{\epsilon\downarrow0}\eta_\epsilon(f,g)
 =\int_\Sigma {\overline{f}g\over|\nabla c|}\,d\sigma
 =:\eta(f,g).
 \label{eq:v6v76-coarea-product}
\end{equation}
The limiting form is positive semidefinite, and its completion after the
null quotient is
\begin{equation}
 \mathcal H_{\rm phys}
 =L^2\left(\Sigma,{d\sigma\over|\nabla c|}\right).
 \label{eq:v6v76-coarea-Hilbert}
\end{equation}
\end{theorem}

\begin{proof}
The coarea formula gives, in a neighbourhood of the regular level,
\begin{equation}
 \eta_\epsilon(f,g)
 =\int_{\mathbb R}\delta_\epsilon(s)
  \left[\int_{c^{-1}(s)}{\overline f g\over|\nabla c|},d\sigma_s\right]ds.
 \label{eq:v6v76-coarea-epsilon}
\end{equation}
The bracket is continuous near zero for compactly supported smooth data and
regular nearby levels.  Since \(\delta_\epsilon\) is an approximate
identity, the limit is its value at zero.  Positivity holds already at each
\(\epsilon\) when \(f=g\), and the limit is the squared
\(L^2(d\sigma/|\nabla c|)\) norm of the trace of \(f\).  Quotienting traces
of zero norm and completing gives \eqref{eq:v6v76-coarea-Hilbert}.
\end{proof}

\begin{corollary}[The real lapse imposes rather than fluctuates the
constraint]
\label{cor:v6v76-lapse-delta}
The real-lapse formula
\begin{equation}
 {1\over2\pi}\int_{\mathbb R}
 \langle f,e^{it\widehat C}g\rangle,dt
 =\eta(f,g)
 \label{eq:v6v76-real-lapse}
\end{equation}
holds distributionally on the test domain of
Theorem~\ref{thm:v6v76-coarea-rigging}.  It creates the positive shell
measure \eqref{eq:v6v76-coarea-product}; it does not assign a covariance to
the negative conformal coordinate.
\end{corollary}

\begin{proof}
Equation \eqref{eq:v6v76-delta-regularizer} identifies the regulated lapse
integral with \(\delta_\epsilon(\widehat C)\).  Take the coarea limit.
\end{proof}

\subsection{Several commuting regular constraints}

Let \(F=(c_1,\ldots,c_r):\mathbb R^n\to\mathbb R^r\) with \(r\le n\), and
put
\begin{equation}
 J_F(x)=\sqrt{\det(DF(x)DF(x)^{\mathsf T})}.
 \label{eq:v6v76-normal-Jacobian}
\end{equation}

\begin{theorem}[Abelian multi-constraint coarea product]
\label{thm:v6v76-multiconstraint}
If \(DF\) has rank \(r\) on \(\Sigma=F^{-1}(0)\), then simultaneous
Gaussian-regulated averaging of the commuting multiplication constraints
has the limit
\begin{equation}
 \eta_F(f,g)
 =\int_\Sigma{\overline f g\over J_F}\,d\sigma_{n-r}\ge0
 \quad(f=g).
 \label{eq:v6v76-multicoarea}
\end{equation}
\end{theorem}

\begin{proof}
Use the codimension-(r) coarea formula with the product approximate delta
\((2\pi\epsilon)^{-r/2}e^{-|F|^2/(2\epsilon)}\).
\end{proof}

The normal Jacobian is the exact finite-dimensional source of the
Faddeev--Popov/coarea weight.  It is not optional: dropping it changes the
physical inner product.

\subsection{Application to the homogeneous CMC--Higgs constraint}

Let \(x\in\mathbb R^d\) be a homogeneous scalar multiplet and let
\(\tau\in\mathbb R\) be York time.  The zero-TT homogeneous constraint of
V72 is
\begin{equation}
 c(x,\tau)
 ={2\tau^2\over3}-2\Lambda-16\pi GV(x)=0.
 \label{eq:v6v76-homogeneous-constraint}
\end{equation}
Put
\begin{equation}
 B(x)=2\Lambda+16\pi GV(x).
 \label{eq:v6v76-Bx}
\end{equation}
Where \(B(x)>0\), the shell has two branches
\begin{equation}
 \tau_\pm(x)=\pm\sqrt{{3\over2}B(x)}.
 \label{eq:v6v76-tau-branches}
\end{equation}

\begin{theorem}[Exact homogeneous physical coarea measure]
\label{thm:v6v76-Higgs-coarea}
For compactly supported smooth kinematical data away from singular zeros of
the constraint, real-lapse group averaging gives
\begin{align}
 \eta(f,g)
 =\int_{\{B>0\}}d^dx\sum_{\varepsilon=\pm}
 {3\over4|\tau_\varepsilon(x)|}
 \overline{f(x,\tau_\varepsilon(x))}
 g(x,\tau_\varepsilon(x)).
 \label{eq:v6v76-Higgs-shell-product}
\end{align}
The factor \(3/(4|\tau|)\) is exact.
\end{theorem}

\begin{proof}
At fixed \(x\), the zeros in \(\tau\) are simple when \(B(x)>0\), and
\begin{equation}
 \left|{\partial c\over\partial\tau}\right|_{\tau_\pm}
 ={4|\tau_\pm|\over3}.
 \label{eq:v6v76-tau-Jacobian}
\end{equation}
The one-variable identity
\(\delta(c(x,\tau))=sum_\pm
\delta(\tau-\tau_\pm)/|\partial_\tau c(\tau_\pm)|\), followed by
integration in \(\tau\), proves the formula.  It is equivalent to the full
coarea measure in Theorem~\ref{thm:v6v76-coarea-rigging} after the shell is
parameterized by \(x\).
\end{proof}

\begin{corollary}[York escape produces an inverse-square amplitude weight]
\label{cor:v6v76-tail-weight}
Suppose
\begin{equation}
 V(x)=\lambda|x|^4+o(|x|^4),
 \qquad \lambda>0.
 \label{eq:v6v76-quartic-tail}
\end{equation}
Then
\begin{equation}
 |\tau_\pm(x)|
 \sim\sqrt{24\pi G\lambda}\,|x|^2,
 \qquad
 {3\over4|\tau_\pm(x)|}\asymp |x|^{-2}.
 \label{eq:v6v76-coarea-tail}
\end{equation}
Thus the total shell-measure tail in \(d\) scalar dimensions has radial
order
\begin{equation}
 r^{d-3}\,dr.
 \label{eq:v6v76-radial-tail}
\end{equation}
It is finite at infinity exactly when \(d=1\); it is logarithmically
divergent for \(d=2\) and power divergent for \(d\ge3\).
\end{corollary}

\begin{proof}
Insert \eqref{eq:v6v76-quartic-tail} into
\eqref{eq:v6v76-tau-branches}.  Multiplying the resulting \(r^{-2}\)
weight by the Euclidean radial Jacobian \(r^{d-1}dr\) gives
\eqref{eq:v6v76-radial-tail}.  The convergence criterion is
\(d-3<-1\).
\end{proof}

\begin{assessment}[What the constraint Jacobian does and does not cure]
\label{ass:v6v76-Jacobian-reading}
For one homogeneous real amplitude, the coarea factor makes the raw
large-amplitude shell measure finite.  For a four-real-component Higgs
doublet before gauge quotient, the tail is of order \(r\,dr\) and diverges.
Passing to a radial gauge coordinate does not automatically remove the
orbit Jacobian; the complete gauge quotient, ghosts, and stabilizer strata
must determine the actual measure.  Constraint group averaging supplies a
real positive sector weight, but that weight alone does not normalize the
full Standard Model homogeneous sector.
\end{assessment}

\subsection{Regular and singular threshold strata}

At \(B(x)=0\), the two York branches meet at \(\tau=0\).  The parameterized
factor in \eqref{eq:v6v76-Higgs-shell-product} diverges, but the full shell
can remain regular.

\begin{proposition}[Simple threshold integrability]
\label{prop:v6v76-threshold-integrability}
Suppose \(B(x_0)=0\) and \(\nabla B(x_0)\ne0\).  Then
\(\nabla_{x,\tau}c(x_0,0)\ne0\), so the constraint shell is regular at the
branch meeting.  In a coordinate \(u=B(x)\ge0\), the parameterized measure
has order
\begin{equation}
 u^{-1/2}du,
 \label{eq:v6v76-simple-threshold}
\end{equation}
which is locally integrable.
\end{proposition}

\begin{proof}
At \((x_0,0)\),
\(\nabla c=(-\nabla B,0)\ne0\).  The implicit function theorem supplies
\(u=B(x)\) as one local coordinate.  Since
\(|\tau_\pm|=\sqrt{3u/2}\), the branch weight is a constant multiple of
\(u^{-1/2}\), whose integral at zero is finite.
\end{proof}

\begin{proposition}[Critical threshold is outside the regular coarea
theorem]
\label{prop:v6v76-critical-threshold}
If \(B(x_0)=0\), \(\nabla B(x_0)=0\), and \(\tau=0\), then
\(dc=0\).  The regular-shell hypothesis
\eqref{eq:v6v76-regular-shell} fails, and neither
\eqref{eq:v6v76-Higgs-shell-product} nor its apparent singular factor is by
itself a proof of a finite physical inner product at that stratum.
\end{proposition}

\begin{proof}
Both the \(x\)-gradient \(-\nabla B\) and the derivative
\(4\tau/3\) vanish.  The coarea denominator has a critical zero, so the
regular-level proof does not apply.
\end{proof}

\subsection{Observables and the canonical positivity domain}

\begin{proposition}[Only constraint-preserving observables descend]
\label{prop:v6v76-Dirac-observables}
Let \(O\) be an operator on the test domain which commutes with the
multiplication constraint and preserves traces on \(\Sigma\).  Then it
descends to the physical coarea Hilbert space.  A general kinematical
operator which changes the value of \(c\) does not define an operator on the
null quotient without an additional relational or gauge-invariant
construction.
\end{proposition}

\begin{proof}
If \([O,\widehat C]=0\), functional calculus gives
\(O\delta_\epsilon(\widehat C)=\delta_\epsilon(\widehat C)O\), so the
rigging form intertwines \(O\) and its adjoint.  In the limit, the action
depends only on the trace of a state on \(\Sigma\) and preserves zero-trace
null vectors.  Without constraint preservation, two kinematical
representatives with the same shell trace can be sent to different shell
traces, so no quotient operator is defined.
\end{proof}

\begin{definition}[Canonical group-averaging card]
\label{def:v6v76-CGAC}
CGAC requires:
\begin{enumerate}[label=\textup{(GA\arabic*)},leftmargin=5.0em]
\item self-adjoint regulated constraints and a declared common invariant
      test domain;
\item a normalized group-averaging or rigging limit, including the complete
      coarea/Faddeev--Popov Jacobian and all singular strata;
\item a positive, nonzero physical inner product after the null quotient;
\item a physical observable algebra that preserves the constraint shell;
      and
\item cofinal bounds on the physical sector measure, including Higgs,
      gauge-orbit, York, boundary, and refinement multiplicities.
\end{enumerate}
\end{definition}

\begin{theorem}[Finite regular multiplication constraints close the local
CGAC rows]
\label{thm:v6v76-CGAC-local}
For one or several commuting regular multiplication constraints,
Theorems~\ref{thm:v6v76-coarea-rigging} and
\ref{thm:v6v76-multiconstraint} close CGAC\textup{(GA1)}--
\textup{(GA4)} on compactly supported smooth data.  The homogeneous
CMC--Higgs constraint has the exact positive shell product
\eqref{eq:v6v76-Higgs-shell-product}; its raw cofinal measure already fails
CGAC\textup{(GA5)} for \(d\ge2\) scalar amplitude dimensions.
\end{theorem}

\begin{proof}
The coarea theorems give the rigging limit, positivity, and null quotient;
Proposition~\ref{prop:v6v76-Dirac-observables} gives the observable row.
Corollary~\ref{cor:v6v76-tail-weight} gives the tail criterion.
\end{proof}

\begin{firewall}[The Einstein constraints are not an Abelian multiplication
family]
\label{fw:v6v76-Dirac-algebra}
The field-theoretic Hamiltonian and momentum constraints have
structure functions, domain questions, and regulator-dependent commutators.
They are not simultaneously multiplication operators on a fixed Euclidean
space.  The coarea theorem is a rigorous local model and identifies the
required Jacobian; it does not prove anomaly-free closure or group averaging
for the Dirac algebra.
\end{firewall}

\begin{firewall}[Canonical positivity is not yet OS reconstruction]
\label{fw:v6v76-no-OS}
The physical coarea product is positive, but no Euclidean-time field measure
or reflection operation has been constructed from it.  Equivalence with the
V67 TT Osterwalder--Schrader Hilbert space, and recovery of the Lorentzian
Hamiltonian and local observables, remain separate requirements.
\end{firewall}

\begin{theorem}[Sixth-series V76 result]
\label{thm:v6v76-status}
Real-lapse group averaging gives a positive delta-constraint inner product.
For regular finite-dimensional constraints it is exactly the coarea measure
\(d\sigma/|\nabla c|\), and its multi-constraint form contains the normal
Jacobian \(J_F^{-1}\).  The homogeneous CMC--Higgs shell has exact branch
weight \(3/(4|\tau|)\).  A quartic Higgs tail therefore gives radial measure
\(r^{d-3}dr\), finite only for one real amplitude dimension.  Simple York
branch meetings are integrable; critical meetings require a separate
singular-stratum analysis.  CGAC closes locally for regular Abelian models
but its Standard Model cofinal and Dirac-algebra rows remain open.
\end{theorem}

\begin{proof}
Use Theorems~\ref{thm:v6v76-mean-ergodic},
\ref{thm:v6v76-coarea-rigging},
\ref{thm:v6v76-multiconstraint}, and
\ref{thm:v6v76-Higgs-coarea}, together with
Corollary~\ref{cor:v6v76-tail-weight} and
Propositions~\ref{prop:v6v76-threshold-integrability}--
\ref{prop:v6v76-Dirac-observables}, subject to the two firewalls.
\end{proof}

V77 will apply the same rigging construction to a finite Fourier-mode
linearized ADM system.  It will average the scalar and longitudinal
constraints, compare the resulting Jacobian and null quotient with the
explicit TT reduction of V61, and determine whether the matter source
preserves the physical domain at that regulated level.

\clearpage
\section{V77: finite-mode ADM rigging and recovery of the TT physical space}
\label{sec:v6v77}

V76 proves positivity of real-lapse group averaging for regular finite
constraints.  This version applies that construction to a finite set of
nonzero Fourier modes of linearized ADM gravity.  The Hamiltonian and
momentum constraints form an Abelian linear system at this order.  Their
rigging map removes four gauge pairs per mode and leaves exactly the two TT
canonical pairs found by direct reduction in V61.  A conserved matter source
shifts the constraint surface by the Newton and momentum solutions without
changing its local Jacobian or positivity.  Lower-boundedness of the reduced
Hamiltonian remains a separate question.

\subsection{Linear canonical group averaging}

Let
\begin{equation}
 \Gamma=\mathbb R^{2(d+r)}
 \label{eq:v6v77-linear-phase-space}
\end{equation}
be a finite-dimensional symplectic vector space with canonical coordinates
\((z,p_z;u,p_u)\), where \(z\in\mathbb R^d\) are intended physical
coordinates and \(u\in\mathbb R^r\) are gauge coordinates.

\begin{lemma}[Linear coisotropic normal form]
\label{lem:v6v77-coisotropic-normal-form}
Let \(C_1,\ldots,C_r\) be independent real linear functions on \(\Gamma\)
with
\begin{equation}
 \{C_a,C_b\}=0.
 \label{eq:v6v77-Abelian-linear}
\end{equation}
There is a real linear symplectic transformation after which
\begin{equation}
 C_a=p_{u_a},qquad 1\le a\le r.
 \label{eq:v6v77-momentum-normal-form}
\end{equation}
The reduced symplectic space is the \((z,p_z)\) space of dimension \(2d\).
\end{lemma}

\begin{proof}
The Hamiltonian vectors \(X_{C_a}\) are independent and span an isotropic
subspace because their pairwise symplectic products are the brackets
\eqref{eq:v6v77-Abelian-linear}.  Extend this isotropic family to a
symplectic basis by the linear symplectic Gram--Schmidt construction.  Choose
dual vectors \(U_a\) with \(\omega(U_a,X_{C_b})=\delta_{ab}\), and complete
their symplectic orthogonal complement by \((z,p_z)\) pairs.  In the dual
coordinates, \(C_a\) are the momenta \(p_{u_a}\).  Setting them to zero and
quotienting the flows generated by them removes \((u,p_u)\).
\end{proof}

Quantize in the Schrödinger representation
\(\mathcal H_{\rm kin}=L^2(\mathbb R^r_u\times\mathbb R^d_z)\), so
\(\widehat C_a=-i\partial_{u_a}\).  Let \(\widetilde f(p_u,z)\) denote the
unitary Fourier transform in \(u\).

\begin{theorem}[Exact Abelian linear rigging product]
\label{thm:v6v77-linear-rigging}
For Schwartz kinematical states,
\begin{align}
 \eta(f,g)
 &:={1\over(2\pi)^r}
 \int_{\mathbb R^r}d^rt\,
 \langle f,e^{it^a\widehat C_a}g\rangle_{\rm kin}
 \label{eq:v6v77-linear-average}\\
 &=\int_{\mathbb R^d}
 \overline{\widetilde f(0,z)}\widetilde g(0,z)\,d^dz.
 \label{eq:v6v77-linear-physical-product}
\end{align}
The physical completion is canonically \(L^2(\mathbb R^d,dz)\).  The null
space consists exactly of states whose gauge-momentum Fourier transform
vanishes at \(p_u=0\).
\end{theorem}

\begin{proof}
In the Fourier representation,
\(e^{it^a\widehat C_a}\) multiplies by \(e^{it\cdot p_u}\).  The real-lapse
integral gives \(\delta^{(r)}(p_u)\), and integration in \(p_u\) yields
\eqref{eq:v6v77-linear-physical-product}.  Its positivity, null space, and
completion are immediate.
\end{proof}

\begin{corollary}[Polarization independence]
\label{cor:v6v77-polarization}
If some linear constraints are configuration coordinates rather than
momenta, a partial Fourier transform exchanges each such coordinate with
its conjugate momentum and reduces the system to
Theorem~\ref{thm:v6v77-linear-rigging}.  The resulting physical Hilbert
space is unitarily equivalent.
\end{corollary}

\begin{proof}
The unitary Fourier transform implements the symplectic exchange
\((u,p_u)\mapsto(p_u,-u)\).
\end{proof}

\subsection{One nonzero linearized ADM mode}

Fix a nonzero spatial wave vector \(k\), put \(n=k/|k|\), and let
\(P_{ij}=\delta_{ij}-n_in_j\).  A real sine or cosine mode has symmetric
metric and momentum tensors \((h_{ij},\pi^{ij})\), a twelve-dimensional
real phase space.  Up to one fixed positive convention factor, the
linearized vacuum constraints are
\begin{align}
 \mathcal H(k)
 &=k_ik_jh_{ij}-|k|^2h
 =-|k|^2P_{ij}h_{ij},
 \label{eq:v6v77-Hamiltonian-mode}\\
 \mathcal H_i(k)&=k_j\pi^{ij}.
 \label{eq:v6v77-momentum-mode}
\end{align}

\begin{lemma}[Rank and Abelian closure at nonzero momentum]
\label{lem:v6v77-ADM-rank}
For \(k\ne0\), the four functions
\(\mathcal H,\mathcal H_i\) are independent and satisfy
\begin{equation}
 \{\mathcal H,\mathcal H_i\}=0,
 \qquad
 \{\mathcal H_i,\mathcal H_j\}=0
 \label{eq:v6v77-linear-ADM-Abelian}
\end{equation}
on the linearized phase space.
\end{lemma}

\begin{proof}
The Hamiltonian constraint depends only on the transverse trace
\(P_{ij}h_{ij}\).  The three momentum constraints depend on the independent
components \(n_j\pi^{ij}\).  Their differentials therefore have rank four.
The momentum constraints mutually commute because they are linear in
momenta.  Their bracket with \(\mathcal H\) is the linearization about a
background satisfying the constraint of the ADM relation in which a
spatial diffeomorphism differentiates the Hamiltonian constraint; at a
fixed flat-background mode the right side is quadratic in perturbations and
vanishes at linear order.  Directly, the contraction of the canonical
bracket is proportional to
\(k_j(k_ik_j-|k|^2\delta_{ij})=0\).
\end{proof}

Decompose a symmetric tensor into:
\begin{enumerate}[label=\textup{(\roman*)},leftmargin=3.5em]
\item two TT components;
\item two transverse-vector components \(h_{na}\);
\item the longitudinal scalar \(h_{nn}\); and
\item the transverse trace \(P_{ij}h_{ij}\).
\end{enumerate}
The momentum tensor has the dual decomposition.

\begin{theorem}[Finite-mode ADM physical space]
\label{thm:v6v77-ADM-physical-space}
Real-lapse and real-shift group averaging of
\eqref{eq:v6v77-Hamiltonian-mode}--
\eqref{eq:v6v77-momentum-mode}, followed by the null quotient, leaves a
four-dimensional reduced phase space consisting of two TT canonical pairs.
In a configuration polarization, the physical Hilbert space of this mode is
\begin{equation}
 \mathcal H_{\rm phys}(k)\simeq L^2(\mathbb R^2_{h^{\rm TT}}).
 \label{eq:v6v77-mode-Hilbert}
\end{equation}
This is unitarily equivalent to the direct canonical reduction of V61.
\end{theorem}

\begin{proof}
Lemma~\ref{lem:v6v77-ADM-rank} and
Lemma~\ref{lem:v6v77-coisotropic-normal-form} provide four gauge pairs.
Theorem~\ref{thm:v6v77-linear-rigging} removes them and leaves phase
dimension \(12-8=4\).  The tensor decomposition shows that the surviving
coordinates and momenta are exactly the two TT components.  Quantizing the
two configuration coordinates gives \eqref{eq:v6v77-mode-Hilbert}.
\end{proof}

\begin{corollary}[Finite mode sets]
\label{cor:v6v77-finite-mode-set}
For any finite set \(K\) of nonzero real Fourier modes, the group-averaged
physical Hilbert space is the finite tensor product
\begin{equation}
 \mathcal H_{\rm phys}(K)
 \simeq\bigotimes_{k\in K}L^2(\mathbb R^2_{h^{\rm TT}(k)}).
 \label{eq:v6v77-finite-product}
\end{equation}
The free reduced Hamiltonian is the sum of the two positive oscillators per
mode constructed in V59.
\end{corollary}

\begin{proof}
Distinct linear Fourier modes have vanishing symplectic brackets, so their
constraint averages and null quotients factor.  Apply
Theorem~\ref{thm:v6v77-ADM-physical-space} mode by mode and use the V59
oscillator Hamiltonian.
\end{proof}

\subsection{The inherited constraint Jacobian}

Use the declared component coordinates
\begin{equation}
 u_0=P_{ij}h_{ij},
 \qquad
 u_i=n_j\pi^{ij}.
 \label{eq:v6v77-constraint-coordinates}
\end{equation}
Then
\begin{equation}
 (\mathcal H,\mathcal H_i)
 =(-|k|^2u_0,\,|k|u_i).
 \label{eq:v6v77-diagonal-constraint-map}
\end{equation}

\begin{proposition}[Modewise ADM coarea scale]
\label{prop:v6v77-ADM-Jacobian}
With the component Lebesgue measure declared in
\eqref{eq:v6v77-constraint-coordinates},
\begin{equation}
 \delta(\mathcal H)\prod_{i=1}^3\delta(\mathcal H_i)
 =|k|^{-5}\delta(u_0)\prod_{i=1}^3\delta(u_i).
 \label{eq:v6v77-mode-Jacobian}
\end{equation}
Thus the raw constraint coarea factor has frequency degree \(-5\) per real
mode in the standard ADM normalization.
\end{proposition}

\begin{proof}
The Hamiltonian row contributes \(|k|^{-2}\), and the three momentum rows
contribute \(|k|^{-1}\) each.
\end{proof}

\begin{warning}[Constraint rescaling changes the rigging normalization]
\label{warn:v6v77-constraint-rescaling}
Replacing a constraint \(C_a\) by a nonzero multiple \(r_a(k)C_a\) leaves
its classical zero set unchanged but multiplies the delta rigging measure by
\(|r_a(k)|^{-1}\).  Therefore
\eqref{eq:v6v77-mode-Jacobian} is meaningful only with the lapse, shift,
symplectic, and Fourier normalizations inherited from one regulated action.
An arbitrary rescaling cannot be used to manufacture cofinal summability.
\end{warning}

\subsection{Affine matter sources}

At linearized level, let \(\rho(k)\) and \(j_i(k)\) be real matter source
coordinates which commute with the four gauge coordinates in this finite
model.  The sourced constraints have the form
\begin{align}
 \mathcal H^{\rm tot}(k)&=-|k|^2u_0-\kappa\rho(k),
 \label{eq:v6v77-sourced-H}\\
 \mathcal H_i^{\rm tot}(k)&=|k|u_i-\kappa j_i(k).
 \label{eq:v6v77-sourced-M}
\end{align}

\begin{theorem}[Source-shifted rigging surface]
\label{thm:v6v77-source-shift}
If the sourced constraints remain commuting self-adjoint multiplication
coordinates on a common test domain, their group-averaged shell is
\begin{equation}
 u_0=-{\kappa\rho(k)\over|k|^2},
 \qquad
 u_i={\kappa j_i(k)\over|k|},
 \label{eq:v6v77-source-shell}
\end{equation}
with the same positive coarea factor \(|k|^{-5}\) as in
\eqref{eq:v6v77-mode-Jacobian}.  The source shifts the shell but does not
change its rank or local Jacobian.
\end{theorem}

\begin{proof}
The constraint map is affine with the same derivative in the \(u\)
variables as \eqref{eq:v6v77-diagonal-constraint-map}.  Solving its zero
set gives \eqref{eq:v6v77-source-shell}, and the coarea determinant is
unchanged.  Positivity follows from
Theorem~\ref{thm:v6v76-coarea-rigging}.
\end{proof}

\begin{corollary}[Canonical origin of the Newton multiplier]
\label{cor:v6v77-Newton-origin}
The \(|k|^{-2}\) displacement in
\eqref{eq:v6v77-source-shell} is the finite-mode canonical origin of the
Newton solution and agrees with the scalar completion in
Lemma~\ref{lem:v6v66-Newton-completion}.  It arises before any probabilistic
norm or absolute value.
\end{corollary}

\begin{proof}
The Hamiltonian constraint is elliptic of degree two in spatial momentum;
solving it divides its source by \(|k|^2\).  Substitution into the quadratic
constraint energy gives the same inverse Laplacian as V66.
\end{proof}

\begin{proposition}[Rigging positivity does not imply energy stability]
\label{prop:v6v77-positivity-not-stability}
The positive physical product on the affine shell is compatible with an
unbounded-below reduced Hamiltonian after the constrained coordinate is
substituted.  In particular, Theorem~\ref{thm:v6v77-source-shift} does not
repair the V62, V71, or V74 Newton and trace-energy obstructions.
\end{proposition}

\begin{proof}
The rigging product depends on the delta shell and its coarea Jacobian.  The
reduced Hamiltonian is a separate function obtained by restricting and
solving the constraint.  Its scalar Schur term has the attractive sign in
\eqref{eq:v6v66-Newton-completion}; positivity of a measure does not change
the sign of that function.
\end{proof}

\subsection{The zero mode and the constraint algebra}

\begin{firewall}[Zero momentum is a singular stratum]
\label{fw:v6v77-zero-mode}
At \(k=0\), all four linear coefficients in
\eqref{eq:v6v77-diagonal-constraint-map} vanish and the rank-four coarea
hypothesis fails.  Homogeneous volume, York time, total energy, and global
gauge data must remain as the separate sectors studied in V72--V76.  They
cannot be obtained by inserting \(k=0\) into
\eqref{eq:v6v77-mode-Jacobian}.
\end{firewall}

\begin{firewall}[Quantum Standard Model sources are not commuting numbers]
\label{fw:v6v77-source-algebra}
The full stress and momentum densities are composite operators.  Their
commutators, Ward contacts, regulator anomalies, and boundary terms must
combine with the gravitational constraints to realize the Dirac algebra.
The affine theorem assumes a joint self-adjoint commuting source model and
therefore does not prove anomaly-free closure for the interacting
SM--Einstein constraints.
\end{firewall}

\begin{definition}[Finite-mode ADM group-averaging card]
\label{def:v6v77-FAGAC}
FAGAC consists of:
\begin{enumerate}[label=\textup{(FA\arabic*)},leftmargin=5.0em]
\item rank-four, anomaly-free linear constraints at each retained nonzero
      mode;
\item real-lapse/shift rigging with the action-normalized coarea Jacobian;
\item a null quotient unitarily equal to the two-polarization TT space;
\item source shifts formed from the complete regulated stress tensor while
      preserving the common domain and constraint algebra; and
\item separate zero-mode sectors and a regulator-uniform reduced-Hamiltonian
      lower bound.
\end{enumerate}
\end{definition}

\begin{theorem}[Vacuum finite-mode FAGAC closes]
\label{thm:v6v77-FAGAC}
For every finite nonzero Fourier-mode regulator of vacuum linearized ADM
gravity, FAGAC\textup{(FA1)}--\textup{(FA3)} close and reproduce V61's TT
phase space and V59's positive oscillator Hilbert space.  Commuting affine
matter coordinates close the local shell and Jacobian part of
FAGAC\textup{(FA4)}.  The composite-source algebra, zero modes, and
Hamiltonian lower bound in FAGAC\textup{(FA4)}--\textup{(FA5)} remain open.
\end{theorem}

\begin{proof}
Use Lemma~\ref{lem:v6v77-ADM-rank},
Theorem~\ref{thm:v6v77-ADM-physical-space},
Corollary~\ref{cor:v6v77-finite-mode-set},
Proposition~\ref{prop:v6v77-ADM-Jacobian}, and
Theorem~\ref{thm:v6v77-source-shift}, subject to the two firewalls.
\end{proof}

\begin{theorem}[Sixth-series V77 result]
\label{thm:v6v77-status}
Real-lapse and shift group averaging of each nonzero linearized ADM mode is
positive and removes four gauge pairs, leaving exactly two TT canonical
pairs.  In the inherited ADM normalization the raw constraint coarea factor
has degree \(|k|^{-5}\).  A commuting affine matter source shifts the scalar
constraint by \(-\kappa\rho/|k|^2\) and the momentum constraints by
\(\kappa j_i/|k|\), without changing the coarea determinant.  This recovers
the Newton multiplier and the V61 reduced vacuum space from one rigging
construction.  Zero modes, composite-source Dirac closure, and reduced
energy stability remain open.
\end{theorem}

\begin{proof}
Use Theorem~\ref{thm:v6v77-FAGAC} and
Corollary~\ref{cor:v6v77-Newton-origin}.
\end{proof}

V78 will compute the reduced Hamiltonian on the affine finite-mode shell.
It will separate the positive TT oscillators, the attractive Newton term,
and source-dependent Jacobian factors, and derive an exact finite-regulator
criterion under which the group-averaged physical inner product supports a
lower-bounded Hamiltonian.

\clearpage
\section{V78: the affine-shell Jacobian cancellation and reduced-Hamiltonian stability}
\label{sec:v6v78}

V77 proves positivity of the finite-mode ADM rigging product but leaves
energy stability open.  This version shows that an affine matter shift does
not generate a stabilizing coarea weight: the surface-element dependence on
the source cancels exactly against the normal Jacobian.  The reduced
Hamiltonian is lower bounded only if its attractive Newton and induced
metric terms are controlled as quadratic forms relative to the physical
kinetic energy.  The high-frequency packet of V71 violates this condition
beyond its explicit cutoff horizon, so positive group averaging alone does
not close the metric--matter theory.

\subsection{Exact cancellation for affine constraint shells}

Let \(u\in\mathbb R^r\), \(z\in\mathbb R^d\), let
\(L:\mathbb R^r\to\mathbb R^r\) be invertible, and let
\(J:\mathbb R^d\to\mathbb R^r\) be smooth.  Consider
\begin{equation}
 C(u,z)=Lu-J(z).
 \label{eq:v6v78-affine-constraint}
\end{equation}

\begin{theorem}[Affine-shell delta identity]
\label{thm:v6v78-affine-delta}
For every compactly supported continuous \(F\),
\begin{equation}
 \int_{\mathbb R^{r+d}}F(u,z)\delta^{(r)}(Lu-J(z))\,du\,dz
 ={1\over|\det L|}
 \int_{\mathbb R^d}F(L^{-1}J(z),z)\,dz.
 \label{eq:v6v78-affine-delta}
\end{equation}
In particular, the induced measure on the parameter \(z\) is independent of
\(DJ(z)\).
\end{theorem}

\begin{proof}
For each fixed \(z\), make the linear change of variables
\(v=Lu-J(z)\).  Then \(du=|\det L|^{-1}dv\), and the delta distribution
evaluates \(v\) at zero.  Fubini's theorem on compact support gives the
identity.
\end{proof}

\begin{corollary}[Geometric coarea cancellation]
\label{cor:v6v78-geometric-cancellation}
If the same shell is parameterized by
\(z\mapsto(L^{-1}J(z),z)\), the growth of its surface element caused by
\(DJ\) cancels exactly against the normal coarea Jacobian.  Hence replacing
the delta calculation by a surface integral cannot create an additional
matter-tail suppression.
\end{corollary}

\begin{proof}
Both calculations represent the unique distributional pairing with
\(\delta^{(r)}(C)\).  Directly, the Gram determinant of the tangent graph
and the normal determinant of \(DC\) contain the same
\(DJ\)-dependent Schur factor, which cancels, leaving
\(|\det L|^{-1}\).
\end{proof}

\begin{corollary}[No source Jacobian in the nonzero ADM modes]
\label{cor:v6v78-no-source-Jacobian}
For the sourced constraints
\eqref{eq:v6v77-sourced-H}--\eqref{eq:v6v77-sourced-M}, the physical
measure on the remaining TT and matter coordinates differs from the vacuum
measure only by the inherited factor \(|k|^{-5}\).  Nonlinear dependence of
\(\rho\) and \(j_i\) on matter coordinates does not supply an additional
amplitude weight as long as the constraints remain affine in the four
eliminated ADM coordinates.
\end{corollary}

\begin{proof}
Use Theorem~\ref{thm:v6v78-affine-delta} with the diagonal map in
\eqref{eq:v6v77-diagonal-constraint-map}; its determinant has magnitude
\(|k|^5\) in the declared component convention.
\end{proof}

This contrasts with V76's homogeneous CMC shell, which is quadratic in
\(\tau\) and therefore has the genuine source-dependent factor
\(3/(4|\tau|)\).

\subsection{The physical quadratic form after reduction}

At a finite regulator, write the group-averaged physical Hilbert space as
\begin{equation}
 \mathcal H_{\rm phys}
 =\mathcal H_{\rm TT}\otimes L^2(\mathbb R^d_\Phi,d\Phi),
 \label{eq:v6v78-physical-Hilbert}
\end{equation}
where constant action-normalization factors have been absorbed once.  Let
\(K_{\rm TT}\ge0\) be the sum of TT oscillators, let \(M>0\) be a matter
mass matrix, and put
\begin{equation}
 \mathfrak k[\Psi]
 =\langle\Psi,K_{\rm TT}\Psi\rangle
 +{1\over2}\int
  \langle\nabla_\Phi\Psi,M^{-1}\nabla_\Phi\Psi\rangle,d\Phi.
 \label{eq:v6v78-kinetic-form}
\end{equation}
The reduced source potential has the schematic but exact finite-regulator
form
\begin{equation}
 W_{\rm red}(\Phi)
 =U(\Phi)
 -{1\over2}\langle\rho(\Phi),G_{\mathrm N}\rho(\Phi)\rangle
 -{1\over2}\langle T^{\rm TT}(\Phi),G_{\rm TT}T^{\rm TT}(\Phi)\rangle
 +W_{\rm loc}(\Phi),
 \label{eq:v6v78-reduced-potential}
\end{equation}
where the positive matrices \(G_{\mathrm N},G_{\rm TT}\) include the
coupling constants and regulated inverse kinetic operators, and
\(W_{\rm loc}\) denotes the common local counterterms.  Define
\begin{equation}
 \mathfrak q[\Psi]
 =\mathfrak k[\Psi]
 +\int W_{\rm red}(\Phi)|\Psi(\Phi)|^2d\Phi.
 \label{eq:v6v78-reduced-form}
\end{equation}

\begin{theorem}[Uniform form-stability criterion]
\label{thm:v6v78-form-stability}
Suppose the negative part \(W_-\) of \(W_{\rm red}\) obeys
\begin{equation}
 \int W_-(\Phi)|\Psi|^2d\Phi
 \le a\,\mathfrak k[\Psi]+b\|\Psi\|^2,
 \qquad 0\le a<1,quad b<\infty,
 \label{eq:v6v78-relative-form-bound}
\end{equation}
on a core for \(\mathfrak k\).  Then \(\mathfrak q\) is closable and lower
bounded by \(-b\); its closure defines a self-adjoint physical Hamiltonian
by the Friedrichs representation theorem.  If \(a,b\) are uniform in the
regulator, so is the lower bound.
\end{theorem}

\begin{proof}
The positive part of the potential can be added to the closed nonnegative
kinetic form by monotone form convergence.  Inequality
\eqref{eq:v6v78-relative-form-bound} makes the negative part relatively form
bounded with relative bound strictly below one.  The KLMN argument gives a
closed lower-bounded sum satisfying
\[
 \mathfrak q[\Psi]
 \ge(1-a)\mathfrak k[\Psi]-b\|\Psi\|^2.
\]
The representation theorem supplies the self-adjoint operator.
\end{proof}

\begin{corollary}[Pointwise coercivity is sufficient but not necessary]
\label{cor:v6v78-pointwise-sufficient}
If \(W_{\rm red}(\Phi)\ge-B\) pointwise with a regulator-independent
\(B\), then \eqref{eq:v6v78-relative-form-bound} holds with \(a=0,b=B\).
More singular negative potentials may still be allowed if their relative
form bound is below one.
\end{corollary}

\begin{proof}
Immediate from integration; the last statement is
Theorem~\ref{thm:v6v78-form-stability}.
\end{proof}

\subsection{A negative polynomial ray destroys the physical lower bound}

\begin{theorem}[Trial-state ray no-go]
\label{thm:v6v78-ray-no-go}
Suppose for a unit vector \(v\in\mathbb R^d\), constants \(c>0,p>0\),
and all sufficiently large \(R\),
\begin{equation}
 \sup_{|y-Rv|\le1}W_{\rm red}(y)\le-cR^p.
 \label{eq:v6v78-negative-ball}
\end{equation}
Then the quadratic form \(\mathfrak q\) is unbounded below.  In particular,
no estimate \eqref{eq:v6v78-relative-form-bound} with \(a<1\) and finite
\(b\) can hold.
\end{theorem}

\begin{proof}
Choose a fixed normalized smooth bump \(\chi\) supported in the unit ball
and a normalized TT vacuum vector \(\Omega_{\rm TT}\).  Put
\(\Psi_R(\Phi)=\Omega_{\rm TT}\chi(\Phi-Rv)\).  Translation invariance of
the finite-dimensional matter kinetic term makes
\(\mathfrak k[\Psi_R]\) independent of \(R\), while
\eqref{eq:v6v78-negative-ball} gives
\[
 \mathfrak q[\Psi_R]\le C_\chi-cR^p\longrightarrow-\infty.
\]
\end{proof}

\begin{corollary}[A decaying shell measure cannot repair a multiplication
ray]
\label{cor:v6v78-weight-not-stability}
The conclusion remains true for a positive smooth physical density
\(w(\Phi)d\Phi\) whose ratio of maximum to minimum on the unit balls
\(B(Rv,1)\) is uniformly bounded and whose logarithmic derivative is
bounded there.  Normalizing a translated bump cancels the small total
weight of the ball; it does not change the diverging negative value of the
potential.
\end{corollary}

\begin{proof}
Normalize \(\chi(\Phi-Rv)\) in \(L^2(wd\Phi)\).  The local comparability of
\(w\) bounds its weighted kinetic energy uniformly, while the potential
expectation is still at most \(-cR^p\).
\end{proof}

This includes polynomial coarea tails such as the homogeneous
\(|\tau(R)|^{-1}\asymp R^{-2}\) factor of V76.  Such a weight can make the
total measure finite in one dimension, but it cannot make an unbounded-below
multiplication Hamiltonian lower bounded.

\subsection{Application to the high-frequency packet}

Fix a packet frequency \(N\) and use its real amplitude \(R\) as one matter
coordinate.  The local scalar action contributes a positive quadratic term
and the quartic coefficient computed in V71.  After TT integration,
\begin{equation}
 W_N^{\rm TT}(R)
 =c_{2,N}R^2
 +\left({9\lambda\over4}-{\kappa^2N^2\over32}\right)R^4
 +O(1),
 \label{eq:v6v78-packet-potential}
\end{equation}
with \(c_{2,N}>0\) at a finite regulator.

\begin{theorem}[Packet instability on the group-averaged Hilbert space]
\label{thm:v6v78-packet-instability}
If
\begin{equation}
 N^2>{72\lambda\over\kappa^2},
 \label{eq:v6v78-packet-threshold}
\end{equation}
then the reduced Hamiltonian on the positive group-averaged physical Hilbert
space is unbounded below along the packet-amplitude coordinate.  The same
conclusion holds when the quadratic Newton source makes the effective
quartic negative beyond \eqref{eq:v6v71-Newton-horizon}.  No affine-shell
coarea Jacobian changes either conclusion.
\end{theorem}

\begin{proof}
Under \eqref{eq:v6v78-packet-threshold}, the quartic coefficient in
\eqref{eq:v6v78-packet-potential} is negative.  On every unit interval
around sufficiently large \(R\), the potential is at most \(-cR^4\).
Apply Theorem~\ref{thm:v6v78-ray-no-go}.  The Newton statement uses
Corollary~\ref{cor:v6v71-Newton-horizon}.  Source independence of the
affine coarea measure is
Corollary~\ref{cor:v6v78-no-source-Jacobian}.
\end{proof}

\begin{corollary}[Positive rigging and Hamiltonian stability are logically
independent]
\label{cor:v6v78-rigging-versus-energy}
FAGAC\textup{(FA1)}--\textup{(FA4)} can hold at a finite regulator while
FAGAC\textup{(FA5)} fails.  The physical inner product is then positive and
nonzero, but the would-be time generator has no lower bound.
\end{corollary}

\begin{proof}
V77 proves the first four local rows for commuting affine sources;
Theorem~\ref{thm:v6v78-packet-instability} supplies a regulator for which
the lower-bound row fails.
\end{proof}

\subsection{Complete finite-regulator Hamiltonian card}

\begin{definition}[Group-averaged Hamiltonian stability card]
\label{def:v6v78-GAHSC}
GAHSC requires:
\begin{enumerate}[label=\textup{(HS\arabic*)},leftmargin=5.0em]
\item a positive CGAC/FAGAC physical inner product with its exact inherited
      coarea density;
\item a densely defined physical kinetic form and an observable source
      potential on the constraint quotient;
\item a regulator-uniform relative bound
      \eqref{eq:v6v78-relative-form-bound} for the complete negative source
      potential, including TT, Newton, counterterm, and zero-mode pieces;
\item a self-adjoint lower-bounded Hamiltonian preserving the physical
      domain; and
\item compatibility of that Hamiltonian with reflection reconstruction and
      cofinal refinement.
\end{enumerate}
\end{definition}

\begin{theorem}[Exact finite-regulator implication]
\label{thm:v6v78-GAHSC-implication}
At any fixed regular regulator, GAHSC\textup{(HS1)}--\textup{(HS3)} imply
GAHSC\textup{(HS4)} by
Theorem~\ref{thm:v6v78-form-stability}.  The V71 packet disproves a uniform
GAHSC\textup{(HS3)} for the globally integrated linearized metric with only
a local Higgs quartic.  A nonlinear constraint restriction or explicit EFT
horizon is therefore still mandatory.
\end{theorem}

\begin{proof}
The first assertion is the KLMN/Friedrichs conclusion.  The second is
Theorem~\ref{thm:v6v78-packet-instability} for arbitrarily large retained
frequencies.
\end{proof}

\begin{firewall}[The schematic reduced potential must come from one action]
\label{fw:v6v78-one-action}
Equation \eqref{eq:v6v78-reduced-potential} is a ledger for terms obtained
after a declared finite-regulator reduction.  In the full theory, seagulls,
nonlinear constraints, gauge fixing, determinant phases, and counterterms
can change it.  They must be derived together from one regulated
SM--Einstein action; inserting independent favorable signs would not prove
GAHSC.
\end{firewall}

\begin{firewall}[A lower-bounded Hamiltonian is not yet the continuum]
\label{fw:v6v78-no-continuum}
Even a uniform GAHSC would still require tightness or operator convergence,
local observable reconstruction, the constraint algebra, and the continuum
Ward identities.  This version isolates the energy gate and does not claim
those later rows.
\end{firewall}

\begin{theorem}[Sixth-series V78 result]
\label{thm:v6v78-status}
For affine sourced constraints, the source-dependent surface Jacobian
cancels exactly and the physical matter measure retains only the inherited
linear constraint determinant.  On that positive physical space, a uniform
relative form bound on the complete negative reduced potential is sufficient
for a self-adjoint lower-bounded Hamiltonian.  A polynomial negative ray is
an exact obstruction.  The V71 packet produces such a ray beyond its TT or
Newton frequency horizon, so positive ADM group averaging does not repair
the global linearized ERMSC.  GAHSC records the required nonlinear-or-EFT
energy estimate.
\end{theorem}

\begin{proof}
Use Theorem~\ref{thm:v6v78-affine-delta},
Theorem~\ref{thm:v6v78-form-stability},
Theorem~\ref{thm:v6v78-ray-no-go}, and
Theorem~\ref{thm:v6v78-packet-instability}, subject to the two firewalls.
\end{proof}

V79 will pursue the EFT branch quantitatively.  It will impose explicit
frequency and amplitude cutoffs below the V71 horizons, derive a finite-mode
relative form bound with its dependence on those cutoffs, and determine
which constants must remain uniform for a controlled low-energy
SM--Einstein effective theory even though the ultraviolet continuum remains
open.

\clearpage
\section{V79: finite-cutoff EFT bounds and the high--high to low stress obstruction}
\label{sec:v6v79}

V78 shows that the reduced physical Hamiltonian needs a regulator-uniform
relative form bound.  This version quantifies what is available at a fixed
spectral and amplitude cutoff.  A new packet with two almost opposite high
momenta produces a fixed low-momentum TT stress.  Its induced gravitational
quartic grows as the fourth power of the matter cutoff, while the local
Higgs quartic remains fixed.  This realizes the common-scale obstruction
that a one-shell estimate misses.  Finite-cutoff stability is elementary
after an amplitude bound is imposed, but none of its constants is uniform
under cutoff removal.

\subsection{Regulated quartic constants}

On the normalized flat three-torus, let
\begin{equation}
 \mathcal B_\Lambda
 =\{\phi:\widehat\phi(k)=0\ \hbox{for }|k|>\Lambda\}
 \label{eq:v6v79-bandlimited-space}
\end{equation}
and remove the zero mode from inverse Laplacians.  Define
\begin{align}
 C_{\mathrm{TT}}(\Lambda)
 &=\sup_{\substack{\phi\in\mathcal B_\Lambda\\\phi\ne0}}
 {\mathfrak Q_{\mathrm{TT}}(\phi)
  \over\int|\phi|^4d\mu},
 \label{eq:v6v79-CTT}\\
 C_{\mathrm N}(\Lambda;m,\xi)
 &=\sup_{\substack{\phi\in\mathcal B_\Lambda\\\phi\ne0}}
 {\mathfrak Q_{\mathrm N,2}(\phi)
  \over\int|\phi|^4d\mu}.
 \label{eq:v6v79-CN}
\end{align}

\begin{lemma}[Finite spectral-cutoff upper bounds]
\label{lem:v6v79-upper-bounds}
For fixed \(m,\xi\), there are constants depending only on the torus
Fourier convention such that, for \(\Lambda\ge1\),
\begin{align}
 C_{\mathrm{TT}}(\Lambda)&\le C\Lambda^4,
 \label{eq:v6v79-TT-upper}\\
 C_{\mathrm N}(\Lambda;m,\xi)
 &\le C_{m,\xi}(1+\Lambda^4).
 \label{eq:v6v79-N-upper}
\end{align}
\end{lemma}

\begin{proof}
The first nonzero eigenvalue of \(-\Delta\) is one, and the TT projector
has norm one.  Hence
\[
 \mathfrak Q_{\mathrm{TT}}(\phi)
 \le\|\nabla\phi\otimes\nabla\phi\|_2^2
 =\int|\nabla\phi|^4d\mu.
\]
The finite-band Bernstein multiplier estimate on \(L^4\) gives
\(\|\nabla\phi\|_4\le C\Lambda\|\phi\|_4\), proving
\eqref{eq:v6v79-TT-upper}.

For the Newton source,
\[
 \rho_2={1\over2}|\nabla\phi|^2+{m^2\over2}\phi^2
 -\xi\Delta(\phi^2).
\]
The inverse Laplacian is bounded by one on the mean-zero sector.  Bernstein
and the product cutoff \(2\Lambda\) give
\[
 \|P_0\rho_2\|_2
 \le C_{m,\xi}(1+\Lambda^2)\|\phi^2\|_2
 =C_{m,\xi}(1+\Lambda^2)\|\phi\|_4^2.
\]
Squaring proves \eqref{eq:v6v79-N-upper}.
\end{proof}

The upper power four is not merely a crude use of the lowest inverse-
Laplacian eigenvalue.  It is attained by a high--high to low interaction.

\subsection{An exact beat-frequency TT source}

For \(N\ge2\), set
\begin{equation}
 p_N=(N,0,0),
 \qquad q_N=(-N,1,0),
 \qquad k_0=p_N+q_N=(0,1,0),
 \label{eq:v6v79-beat-momenta}
\end{equation}
and define
\begin{equation}
 \varphi_N(x)
 =A\bigl(\cos(p_N\cdot x)+\cos(q_N\cdot x)\bigr).
 \label{eq:v6v79-beat-packet}
\end{equation}
It lies in \(\mathcal B_{\sqrt{N^2+1}}\).

\begin{lemma}[Local quartic remains fixed]
\label{lem:v6v79-beat-L4}
For every \(N\),
\begin{equation}
 \int_{\mathbb T^3}|\varphi_N|^4d\mu={9\over4}A^4.
 \label{eq:v6v79-beat-L4}
\end{equation}
\end{lemma}

\begin{proof}
The two integer momenta are linearly independent.  Haar integration of any
function of their two phases equals Haar integration on the two-torus.
The moment computation is therefore the same as in
Lemma~\ref{lem:v6v71-local-packet}.
\end{proof}

\begin{theorem}[Exact low-momentum TT witness]
\label{thm:v6v79-beat-TT}
At the two modes \(\pm k_0\),
\begin{equation}
 \left\|P^{\mathrm{TT}}(k_0)
 \widehat{\nabla\varphi_N\otimes\nabla\varphi_N}(k_0)
 \right\|_{\mathrm{HS}}^2
 ={A^4N^4\over8}.
 \label{eq:v6v79-beat-TT-mode}
\end{equation}
Consequently
\begin{equation}
 \mathfrak Q_{\mathrm{TT}}(\varphi_N)
 \ge {A^4N^4\over4},
 \qquad
 C_{\mathrm{TT}}(\sqrt{N^2+1})\ge {N^4\over9}.
 \label{eq:v6v79-beat-TT-lower}
\end{equation}
\end{theorem}

\begin{proof}
The cross term at \(k_0=p_N+q_N\) has tensor coefficient
\begin{equation}
 -{A^2\over4}
 (p_N\otimes q_N+q_N\otimes p_N).
 \label{eq:v6v79-beat-coefficient}
\end{equation}
Relative to \(k_0=e_2\), TT projection kills every tensor factor containing
\(e_2\).  Since
\[
 p_N\otimes q_N+q_N\otimes p_N
 =-2N^2e_1\otimes e_1
 +N(e_1\otimes e_2+e_2\otimes e_1),
\]
the surviving transverse tensor is the trace-free part of
\(-2N^2e_1\otimes e_1\) on the \(e_1,e_3\) plane.  Its squared norm is
\(2N^4\).  Multiplication by \(A^4/16\) proves
\eqref{eq:v6v79-beat-TT-mode}.  Both \(k_0\) and \(-k_0\) contribute and
\(|k_0|=1\), giving the first inequality in
\eqref{eq:v6v79-beat-TT-lower}.  Divide by
\eqref{eq:v6v79-beat-L4} for the second.
\end{proof}

\begin{corollary}[Sharp cutoff order for the TT quartic constant]
\label{cor:v6v79-TT-sharp-order}
There are constants \(0<c<C<\infty\) such that for all sufficiently large
\(\Lambda\),
\begin{equation}
 c\Lambda^4\le C_{\mathrm{TT}}(\Lambda)\le C\Lambda^4.
 \label{eq:v6v79-TT-sharp-order}
\end{equation}
\end{corollary}

\begin{proof}
The upper bound is Lemma~\ref{lem:v6v79-upper-bounds}.  Given \(\Lambda\),
choose an integer \(N\) comparable to \(\Lambda\) with
\(\sqrt{N^2+1}\le\Lambda\), and use
Theorem~\ref{thm:v6v79-beat-TT}.
\end{proof}

\subsection{The beat-frequency Newton source}

At \(k_0\), the Fourier coefficient of the quadratic energy density
\eqref{eq:v6v71-quadratic-density} is
\begin{equation}
 \widehat\rho_2^{(\xi)}(k_0)
 ={A^2\over4}\bigl(N^2+m^2+2\xi\bigr).
 \label{eq:v6v79-beat-density}
\end{equation}

\begin{lemma}[Beat-density coefficient]
\label{lem:v6v79-beat-density}
Equation \eqref{eq:v6v79-beat-density} holds exactly.
\end{lemma}

\begin{proof}
The cross part of \(\tfrac12|\nabla\varphi_N|^2\) is
\(A^2(p_N\cdot q_N)\sin(p_N\cdot x)\sin(q_N\cdot x)\).
Since \(p_N\cdot q_N=-N^2\), its complex Fourier coefficient at
\(p_N+q_N\) is \(A^2N^2/4\).  The coefficient of \(\varphi_N^2\) there is
\(A^2/2\); the mass term contributes \(A^2m^2/4\), and
\(-\xi\Delta(\varphi_N^2)\) contributes \(\xi A^2/2\) because
\(|k_0|=1\).  Summing gives the result.
\end{proof}

\begin{theorem}[Newton cutoff constant also has fourth-order growth]
\label{thm:v6v79-Newton-fourth}
For fixed \(m,\xi\), there are \(c_{m,\xi}>0\) and \(N_0\) such that
\begin{equation}
 \mathfrak Q_{\mathrm N,2}(\varphi_N)
 \ge c_{m,\xi}A^4N^4
 \qquad(N\ge N_0).
 \label{eq:v6v79-Newton-beat-lower}
\end{equation}
Consequently
\begin{equation}
 C_{\mathrm N}(\Lambda;m,\xi)\asymp_{m,\xi} \Lambda^4
 \label{eq:v6v79-Newton-sharp-order}
\end{equation}
for large \(\Lambda\).
\end{theorem}

\begin{proof}
The two modes \(\pm k_0\) alone give
\begin{equation}
 \mathfrak Q_{\mathrm N,2}(\varphi_N)
 \ge {A^4\over8}(N^2+m^2+2\xi)^2.
 \label{eq:v6v79-Newton-two-mode}
\end{equation}
For fixed \(m,\xi\), this is at least \(A^4N^4/32\) for sufficiently
large \(N\).  Divide by \eqref{eq:v6v79-beat-L4} and combine with
Lemma~\ref{lem:v6v79-upper-bounds}.
\end{proof}

\begin{assessment}[The obstruction is a scale merger]
\label{ass:v6v79-scale-merger}
V71 used an output momentum comparable to the input frequency and found
quadratic cutoff growth.  The new pair has input frequencies of order
\(N\) but output momentum one, so the inverse gravitational carrier supplies
no ultraviolet decay.  The complete stress convolution therefore has
fourth-order cutoff growth.  A shellwise proof that assigns the output only
to its low shell would miss the two high-frequency ancestors; a proof that
assigns it only to the high shells would miss the long gravitational
carrier.  The interaction must retain both labels until after the source is
assembled.
\end{assessment}

\subsection{A finite-cutoff coercivity condition}

Let \(\kappa_{\mathrm{TT}}^2/2\) and \(\gamma_{\mathrm N}\) be the positive
coefficients multiplying the two attractive source forms.  The quadratic-
source quartic obeys
\begin{equation}
 \lambda\int|\phi|^4d\mu
 -{\kappa_{\mathrm{TT}}^2\over2}\mathfrak Q_{\mathrm{TT}}(\phi)
 -\gamma_{\mathrm N}\mathfrak Q_{\mathrm N,2}(\phi)
 \ge\lambda_{\rm eff}(\Lambda)\int|\phi|^4d\mu,
 \label{eq:v6v79-cutoff-coercivity}
\end{equation}
where
\begin{equation}
 \lambda_{\rm eff}(\Lambda)
 =\lambda-{\kappa_{\mathrm{TT}}^2\over2}C_{\mathrm{TT}}(\Lambda)
         -\gamma_{\mathrm N}C_{\mathrm N}(\Lambda;m,\xi).
 \label{eq:v6v79-effective-lambda}
\end{equation}

\begin{theorem}[Finite-cutoff quartic stability]
\label{thm:v6v79-finite-cutoff-stability}
If
\begin{equation}
 \lambda_{\rm eff}(\Lambda)>0,
 \label{eq:v6v79-positive-effective-lambda}
\end{equation}
then the quadratic-source part of the reduced scalar potential is coercive
at the spectral cutoff \(\Lambda\).  For fixed nonzero gravitational
couplings and fixed \(\lambda\), condition
\eqref{eq:v6v79-positive-effective-lambda} fails for all sufficiently large
\(\Lambda\).
\end{theorem}

\begin{proof}
The first statement is the definition of the two regulator constants.  The
second follows from
Corollary~\ref{cor:v6v79-TT-sharp-order} and
Theorem~\ref{thm:v6v79-Newton-fourth}.
\end{proof}

\begin{corollary}[Required bare-coupling scaling in this normalization]
\label{cor:v6v79-coupling-scaling}
If \(\lambda\) remains bounded and positive while
\(\Lambda\to\infty\), the sufficient condition requires
\begin{equation}
 \kappa_{\mathrm{TT}}^2C_{\mathrm{TT}}(\Lambda)
 +\gamma_{\mathrm N}C_{\mathrm N}(\Lambda;m,\xi)=O(1).
 \label{eq:v6v79-coupling-scaling}
\end{equation}
In the declared torus normalization this forces the corresponding bare
gravitational coefficients to decay at least at fourth inverse cutoff order
as quadratic coefficients.  Keeping the physical gravitational coupling
fixed requires additional operators or a nonlinear source mechanism.
\end{corollary}

\begin{proof}
Use the two-sided fourth-order growth of the regulator constants in
\eqref{eq:v6v79-effective-lambda}.
\end{proof}

\subsection{Amplitude cutoffs give only cutoff-dependent lower bounds}

Let the finite-dimensional regulated scalar coordinates satisfy
\(|\Phi|\le R\).  Every polynomial reduced potential is then bounded below.

\begin{proposition}[Finite spectral-and-amplitude box]
\label{prop:v6v79-finite-box}
At fixed \((\Lambda,R)\), suppose all coefficients and local counterterms
are finite.  Then there is a finite number \(B(\Lambda,R)\) such that
\begin{equation}
 W_{\rm red}(\Phi)\ge-B(\Lambda,R)
 \qquad(|\Phi|\le R).
 \label{eq:v6v79-box-lower}
\end{equation}
The finite-box physical Hamiltonian has a Friedrichs lower-bounded
realization.  If the full Higgs potential is included in the linearized
Newton source, one necessarily permits a contribution of order
\begin{equation}
 B(\Lambda,R)\gtrsim
 \gamma_{\mathrm N}\lambda^2R^8
 \label{eq:v6v79-octic-box-loss}
\end{equation}
along nonconstant amplitude directions, so this argument is not uniform as
\(R\to\infty\).
\end{proposition}

\begin{proof}
Continuity on the compact finite-dimensional ball gives a finite minimum.
Corollary~\ref{cor:v6v78-pointwise-sufficient} gives the Hamiltonian.  The
order-eight loss is the leading term in
\eqref{eq:v6v71-octic-asymptotic} evaluated at the amplitude boundary.
\end{proof}

\begin{definition}[Controlled finite-cutoff EFT card]
\label{def:v6v79-CEFTC}
CEFTC requires:
\begin{enumerate}[label=\textup{(EF\arabic*)},leftmargin=5.0em]
\item declared spectral, amplitude, volume, and source-domain cutoffs;
\item the exact regulator constants
      \eqref{eq:v6v79-CTT}--\eqref{eq:v6v79-CN}, including high--high to
      low source mergers;
\item positivity of \(\lambda_{\rm eff}(\Lambda)\), or a stronger complete
      relative form bound including all EFT operators;
\item a lower-bounded physical Hamiltonian with constants explicitly
      depending on \((\Lambda,R)\); and
\item matching and error estimates only for observables below the declared
      EFT horizon, with no assertion of cutoff removal.
\end{enumerate}
\end{definition}

\begin{theorem}[The finite box closes but the continuum row does not]
\label{thm:v6v79-CEFTC}
At every fixed finite \((\Lambda,R)\), CEFTC\textup{(EF1)},
\textup{(EF2)}, and \textup{(EF4)} hold for the polynomial linearized
reduced model; CEFTC\textup{(EF3)} holds when
\eqref{eq:v6v79-positive-effective-lambda} and the remaining finite
counterterm inequalities hold.  The beat packet proves that this route
cannot make the constants uniform as \(\Lambda\to\infty\) at fixed
gravitational and Higgs couplings.
\end{theorem}

\begin{proof}
Use Lemma~\ref{lem:v6v79-upper-bounds},
Theorem~\ref{thm:v6v79-finite-cutoff-stability}, and
Proposition~\ref{prop:v6v79-finite-box}; nonuniformity is
Corollary~\ref{cor:v6v79-TT-sharp-order}.
\end{proof}

\begin{firewall}[Cutoff stability is not ultraviolet completion]
\label{fw:v6v79-no-UV-completion}
CEFTC is a mathematically controlled finite effective theory.  It does not
prove that the physical Newton coupling vanishes, that a particular
higher-derivative counterterm is allowed, or that an infinite tower of EFT
operators converges.  The full SM--Einstein continuum still requires a
nonlinear or ultraviolet-complete mechanism beyond this card.
\end{firewall}

\begin{firewall}[The beat packet is a scalar subsector test]
\label{fw:v6v79-beat-scope}
Gauge, fermion, ghost, and nonlinear metric contributions are absent.  The
packet disproves any proposed deterministic bound whose hypotheses include
this scalar subsector with only the displayed local quartic.  It does not
prove that no complete Standard Model source assembly can add a compensating
term; such a term must be exhibited with its sign and common ancestry.
\end{firewall}

\begin{theorem}[Sixth-series V79 result]
\label{thm:v6v79-status}
At spectral cutoff \(\Lambda\), the worst TT and quadratic Newton source
constants grow exactly at fourth order.  The lower bound is realized by two
nearly opposite high-frequency scalar modes whose product has fixed output
momentum, while their local Higgs quartic stays constant.  Hence a fixed
local quartic and fixed gravitational couplings cannot satisfy the global
linearized relative form bound as the cutoff is removed.  Finite spectral
and amplitude boxes have lower-bounded Friedrichs Hamiltonians with explicit
cutoff losses, defining CEFTC but not an ultraviolet continuum.
\end{theorem}

\begin{proof}
Use Corollary~\ref{cor:v6v79-TT-sharp-order},
Theorem~\ref{thm:v6v79-Newton-fourth},
Theorem~\ref{thm:v6v79-finite-cutoff-stability}, and
Theorem~\ref{thm:v6v79-CEFTC}, subject to the two firewalls.
\end{proof}

V80 will perform the compulsory YM/NS audit.  It will compare their newest
complete-source handling of rare or scale-merged sectors with the beat
packet, then decide whether the next attack should retain a nonlinear metric
source through the high--high to low merger or formulate a precise
ultraviolet-completion obstruction card.

\clearpage
\chapter{Sixth-series V80: companion audit, the merger-symbol theorem,
and the nonlinear beat-response gate}
\label{ch:v6v80-merger}

\section{Purpose of this continuation}
\label{sec:v6v80}

V79 proved that the attractive TT and Newton source squares have sharp
fourth-order spectral-cutoff growth.  The witness was not a large output
momentum: two almost antipodal high-frequency scalar modes merged into the
fixed mode \(k_0=e_2\).  This continuation first performs the compulsory
read-only audit of the newest Yang--Mills and Navier--Stokes notes.  It then
abstracts the V79 computation into a general necessary symbol condition,
strengthens the finite-cutoff conclusion to a fixed-amplitude spectral
instability of the displayed linearized reduced model, and finally matches
the amplitude to the frequency to identify the point at which nonlinear
metric response, rather than another linear estimate, becomes the next
mathematical gate.

No assertion about a full Standard Model--Einstein continuum is made here.
The purpose is to replace a vague ultraviolet difficulty by an exact
two-frequency test that every proposed completion must pass.

\section{Compulsory V80 companion audit}
\label{sec:v6v80-audit}

The companion folder is shared with concurrently running programmes.  Each
file below was therefore selected by its largest active version number and
read into memory once.  The byte count, logical source-line count, and
SHA--256 digest refer to those same bytes; the digest, rather than a mutable
directory listing, defines the evidence used in this audit.

\begin{definition}[Hash-defined V80 audit snapshots]
\label{def:v6v80-audit-snapshots}
The Yang--Mills snapshot is
\begin{align}
 &\texttt{Renewal\_Yang\_Mills\_Mass\_Gap\_Research\_Notes\_V56.tex},
 \notag\\
 &860\,041\ \text{bytes},\qquad 27\,071\ \text{logical source lines},
 \notag\\
 &\texttt{2145bd724f845006021db608ebb476f8f78048d3c9a764389eb218055ed71507}.
 \label{eq:v6v80-YM-snapshot}
\end{align}
The Navier--Stokes snapshot is
\begin{align}
 &\texttt{Renewal\_NS\_Stability\_Research\_Notes\_V87.tex},
 \notag\\
 &1\,557\,431\ \text{bytes},\qquad 36\,999\ \text{logical source lines},
 \notag\\
 &\texttt{27615589d9e2ccf997823c1043b1eaec27bfa15c49df4d45cc45adf3909604d5}.
 \label{eq:v6v80-NS-snapshot}
\end{align}
\end{definition}

\subsection{Yang--Mills endpoint}

YM V56 starts from the exact positive Hartman--Watson local/tail split of
V55.  It proves local conditional source invertibility on the V44 weighted
simplex, computes the exact singleton and conjugate-singlet zeros, and proves
an exponentially small noncommutative rare-log correction on that baseline
domain.  Its generic multivariable Cauchy extraction gives
\[
 e^{-q/s}(Km)^m\prod_i b_i.
\]
For equal minimal marks \(b_i=\sqrt{s}\), comparison with the desired
Gaussian-renormalized scale
\(K^m(\prod_i b_i)(\sum_i b_i)^{m-2}\) leaves a factor proportional to
\(e^{-q/s}C^m m^2s^{1-m/2}\), unbounded in the arity \(m\) for fixed
\(s<1\).  Thus positivity and zero-freeness have been established on the
baseline domain, but estimating the rare sector after separating it still
loses the strong connected homogeneity.  Its GRLTC gate requires the local
forest and rare logarithm to remain assembled through squarefree coefficient
extraction.  No cutoff removal or mass gap is claimed there.

\subsection{Navier--Stokes endpoint}

NS V87 transfers the Gaussian-weighted energy identity from an ancient
limit to the original torus solution by a periodized backward heat weight.
At a homogeneous candidate singular point, integration over scale proves
that the weighted radial energy-pressure flux tends to \(-\infty\) at a
linear rate in the logarithmic scale ratio.  The sign appears only after the
positive weighted dissipation, the scale derivative, and the nonlinear flux
are kept in one identity and integrated.  A pointwise-in-scale lower bound
for the flux is not proved, and at lacunary scales the local dissipation may
vanish while distant satellites are Gaussian-suppressed.  No global
regularity conclusion is claimed there.

\begin{theorem}[Companion-audit transfer decision]
\label{thm:v6v80-audit-decision}
No numerical heat-kernel, source-radius, dissipation, or flux constant from
\eqref{eq:v6v80-YM-snapshot}--\eqref{eq:v6v80-NS-snapshot} transfers to the
SM--Einstein estimates.  The transferable proof rule is structural:
exceptional conditional sectors and scale-merging convolution pairs must be
kept inside the exact nonlinear operation until the cancellation, sign, or
coercive quantity has been formed.
\end{theorem}

\begin{proof}
The YM constants concern a positive heat-time mixture and completely
bounded tensor sources; the NS constants concern local energy, pressure,
and viscous dissipation.  Neither set satisfies the hypotheses of the
Hamiltonian constraint or TT source estimates, so a numerical transfer
would be invalid.  The logical obstruction is nevertheless the same.  In
YM, separate absolute Cauchy extraction destroys the needed arity
homogeneity even though the rare probability is exponentially small.  In
NS, neither the nonlinear flux nor the remaining terms have the required
sign separately, while the integrated assembled identity has a signed
consequence.  V79 likewise shows that assigning a source only to its low
output shell erases its two high-frequency ancestors.  These three facts
prove the stated proof rule and nothing stronger.
\end{proof}

\section{A general high--high to low merger theorem}
\label{sec:v6v80-general-merger}

Let \(E\) be a finite-dimensional real or complex Hilbert space.  For a
real trigonometric polynomial \(\phi\), consider a translation-invariant
quadratic source \(\mathcal B(\phi,\phi)\) with Fourier representation
\begin{equation}
 \widehat{\mathcal B(\phi,\phi)}(k)
 =\sum_{r+s=k}\mathsf b(k;r,s)\,widehat\phi(r)\widehat\phi(s)
 \in E.
 \label{eq:v6v80-bilinear-source}
\end{equation}
Assume the reality compatibility
\begin{equation}
 \mathsf b(-k;-r,-s)
 =\overline{\mathsf b(k;r,s)}
 \label{eq:v6v80-reality}
\end{equation}
and let \(w(k)=w(-k)>0\) for \(k\ne0\).  Define
\begin{equation}
 \mathfrak Q_{\mathcal B}(\phi)
 :=\sum_{k\ne0}w(k)
 \left\|\widehat{\mathcal B(\phi,
 \phi)}(k)\right\|_E^2.
 \label{eq:v6v80-QB}
\end{equation}

Use the V79 momenta
\begin{equation}
 p_N=(N,0,0),\qquad q_N=(-N,1,0),\qquad
 k_0=p_N+q_N=e_2,
 \label{eq:v6v80-momenta}
\end{equation}
and set
\begin{equation}
 \beta_N
 :=\mathsf b(k_0;p_N,q_N)+\mathsf b(k_0;q_N,p_N).
 \label{eq:v6v80-sym-symbol}
\end{equation}

\begin{lemma}[Two-character Haar law]
\label{lem:v6v80-Haar-law}
For every Borel function \(F\) integrable on \([-2|A|,2|A|]\),
\begin{equation}
 \int_{\mathbb T^3}
 F\!\left(A\cos(p_N\cdot x)+A\cos(q_N\cdot x)\right)d\mu(x)
 ={1\over(2\pi)^2}\int_0^{2\pi}\!\int_0^{2\pi}
 F(A\cos\theta+A\cos\eta)\,d\theta d\eta.
 \label{eq:v6v80-Haar-pushforward}
\end{equation}
In particular the left side is independent of \(N\).
\end{lemma}

\begin{proof}
The map
\[
 x\longmapsto(p_N\cdot x,q_N\cdot x)\pmod{2\pi}
\]
is a continuous homomorphism from \(\mathbb T^3\) onto \(\mathbb T^2\):
its integer character matrix has rank two, and the equations
\(Nx_1=\theta\), \(-Nx_1+x_2=\eta\) can be solved modulo \(2\pi\).
The pushforward of normalized Haar measure under a surjective compact-group
homomorphism is normalized Haar measure.  Applying this to the displayed
integrable function proves the identity.
\end{proof}

\begin{theorem}[Merger-symbol lower bound]
\label{thm:v6v80-merger-symbol}
Let
\[
 \phi_N(x)=A\bigl(\cos(p_N\cdot x)+\cos(q_N\cdot x)\bigr).
\]
Then
\begin{equation}
 \mathfrak Q_{\mathcal B}(\phi_N)
 \ge {A^4\over8}\,w(k_0)\|\beta_N\|_E^2.
 \label{eq:v6v80-merger-lower}
\end{equation}
Consequently, if
\begin{equation}
 \mathfrak Q_{\mathcal B}(\phi)
 \le C_\Lambda\int_{\mathbb T^3}|\phi|^4d\mu
 \label{eq:v6v80-local-quartic-control}
\end{equation}
for all real \(\phi\in\mathcal B_\Lambda\), then, whenever
\(\sqrt{N^2+1}\le\Lambda\),
\begin{equation}
 C_\Lambda\ge {w(k_0)\|\beta_N\|_E^2\over18}.
 \label{eq:v6v80-necessary-symbol-bound}
\end{equation}
Thus uniform local-quartic control requires the symmetrized merger symbol
\(w(k_0)^{1/2}\beta_N\) to remain uniformly bounded along every pair of
high inputs with fixed nonzero output.
\end{theorem}

\begin{proof}
The Fourier support of \(\phi_N\) is
\(\{\pm p_N,\pm q_N\}\), with coefficient \(A/2\) at each point.
The only ordered pairs in this support whose sum is \(k_0\) are
\((p_N,q_N)\) and \((q_N,p_N)\).  Hence
\[
 \widehat{\mathcal B(\phi_N,\phi_N)}(k_0)
 ={A^2\over4}\beta_N.
\]
Reality gives the conjugate coefficient at \(-k_0\).  Keeping just those
two terms in the nonnegative sum \eqref{eq:v6v80-QB} proves
\eqref{eq:v6v80-merger-lower}.  Lemma~\ref{lem:v6v80-Haar-law}, or the
direct fourth-moment computation, gives
\[
 \int|\phi_N|^4d\mu={9A^4\over4}.
\]
Insert this and \eqref{eq:v6v80-merger-lower} into
\eqref{eq:v6v80-local-quartic-control} and cancel \(A^4\).
\end{proof}

\begin{corollary}[Derivative merger obstruction]
\label{cor:v6v80-derivative-merger}
If \(w(k_0)\) is fixed and positive and
\(\|\beta_N\|_E\ge cN^r\) along an infinite sequence, then
\begin{equation}
 C_{\sqrt{N^2+1}}\ge {c^2w(k_0)\over18}N^{2r}.
 \label{eq:v6v80-symbol-growth}
\end{equation}
In particular a second-order quadratic source, for which \(r=2\),
forces fourth-order cutoff growth unless its complete symmetrized symbol
cancels on all such pairs.
\end{corollary}

\begin{proof}
This is \eqref{eq:v6v80-necessary-symbol-bound}.  The last statement is
the specialization \(r=2\).
\end{proof}

\section{Exact specialization to the Einstein source squares}
\label{sec:v6v80-Einstein-specialization}

\begin{theorem}[The V79 constants are symbol obstructions]
\label{thm:v6v80-Einstein-symbols}
For the TT stress source of V79,
\begin{equation}
 w(k_0)=1,qquad \|\beta_N^{\mathrm{TT}}\|_{\mathrm{HS}}^2=2N^4,
 \label{eq:v6v80-TT-beta}
\end{equation}
and Theorem~\ref{thm:v6v80-merger-symbol} gives exactly
\begin{equation}
 C_{\mathrm{TT}}(\sqrt{N^2+1})\ge {N^4\over9}.
 \label{eq:v6v80-TT-recovered}
\end{equation}
For the quadratic Newton density,
\begin{equation}
 \beta_N^{\mathrm N}=N^2+m^2+2\xi,
 \label{eq:v6v80-Newton-beta}
\end{equation}
and therefore
\begin{equation}
 C_{\mathrm N}(\sqrt{N^2+1};m,\xi)
 \ge{(N^2+m^2+2\xi)^2\over18}.
 \label{eq:v6v80-Newton-symbol-lower}
\end{equation}
\end{theorem}

\begin{proof}
For TT, the symmetrized derivative symbol before its scalar Fourier
coefficient is
\(P^{\mathrm{TT}}(k_0)(p_N\otimes q_N+q_N\otimes p_N)\).
The computation in Theorem~\ref{thm:v6v79-beat-TT} gives its squared
Hilbert--Schmidt norm \(2N^4\), while the inverse carrier has
\(|k_0|^{-2}=1\).  Substitution in
\eqref{eq:v6v80-necessary-symbol-bound} proves
\eqref{eq:v6v80-TT-recovered}.

For Newton, Lemma~\ref{lem:v6v79-beat-density} says that the Fourier
coefficient at \(k_0\) is \(A^2\beta_N^{\mathrm N}/4\).  The output weight
is again one.  Applying the same theorem gives
\eqref{eq:v6v80-Newton-symbol-lower}.
\end{proof}

\begin{assessment}[What a shell estimate must remember]
\label{ass:v6v80-shell-memory}
The obstruction is visible before any norm inequality: it is the value of
the complete symmetrized source symbol on
\((p_N,q_N,k_0)\).  A dyadic argument is admissible only if its bookkeeping
retains both input labels and the output label until after projection and
carrier inversion.  Replacing the pair by the low output scale before
applying derivatives discards \(N^2\); replacing it by the high input scale
before inversion incorrectly inserts ultraviolet decay into the fixed
carrier \(|k_0|^{-2}=1\).
\end{assessment}

\section{Fixed-amplitude spectral instability of the displayed reduced
model}
\label{sec:v6v80-fixed-amplitude}

The fourth-order cutoff loss is stronger than a failure of one sufficient
inequality.  Within the scalar quadratic-source reduced model itself, it
produces a sequence of uniformly bounded scalar fields whose energy tends
to \(-\infty\) as the spectral cutoff is removed.

\begin{theorem}[No derivative-free local potential repairs the linearized
beat sequence]
\label{thm:v6v80-local-potential-no-go}
Let \(A\ne0\), let \(W:\mathbb R\to\mathbb R\) be continuous, and let
\(a_2\ge0\), \(a_{\mathrm{TT}},a_{\mathrm N}\ge0\), with
\(a_{\mathrm{TT}}+a_{\mathrm N}>0\).  On the V79 beat packet define
\begin{equation}
 \mathcal E_N
 :=a_2\int|\nabla\phi_N|^2d\mu
 +\int W(\phi_N)d\mu
 -a_{\mathrm{TT}}\mathfrak Q_{\mathrm{TT}}(\phi_N)
 -a_{\mathrm N}\mathfrak Q_{\mathrm N,2}(\phi_N).
 \label{eq:v6v80-reduced-energy}
\end{equation}
Then
\begin{equation}
 \mathcal E_N\longrightarrow-\infty
 \qquad(N\to\infty),
 \label{eq:v6v80-energy-minus-infinity}
\end{equation}
although \(\|\phi_N\|_{L^\infty}\le2|A|\) for every \(N\).
\end{theorem}

\begin{proof}
Orthogonality of the two characters gives
\begin{equation}
 \int|\nabla\phi_N|^2d\mu
 ={A^2\over2}(|p_N|^2+|q_N|^2)
 =A^2\left(N^2+{1\over2}\right).
 \label{eq:v6v80-gradient-exact}
\end{equation}
By Lemma~\ref{lem:v6v80-Haar-law},
\begin{equation}
 \int W(\phi_N)d\mu
 ={1\over(2\pi)^2}\int_0^{2\pi}\!\int_0^{2\pi}
 W(A\cos\theta+A\cos\eta)\,d\theta d\eta=:M_{A,W},
 \label{eq:v6v80-potential-fixed}
\end{equation}
a finite number independent of \(N\).  The V79 two-output-mode bounds give
\begin{align}
 \mathfrak Q_{\mathrm{TT}}(\phi_N)&\ge {A^4N^4\over4},
 \label{eq:v6v80-QTT-fixed}\\
 \mathfrak Q_{\mathrm N,2}(\phi_N)&\ge
 {A^4\over8}(N^2+m^2+2\xi)^2.
 \label{eq:v6v80-QN-fixed}
\end{align}
The positive part of \eqref{eq:v6v80-reduced-energy} is at most quadratic
in \(N\), up to the constant \(M_{A,W}\), whereas at least one negative
term has a strictly negative fourth-order coefficient.  This proves
\eqref{eq:v6v80-energy-minus-infinity}.  The pointwise bound follows from
\(|\cos|\le1\).
\end{proof}

\begin{corollary}[An amplitude box alone does not give spectral uniformity]
\label{cor:v6v80-amplitude-box}
Fix \(R>0\).  For every \(0<|A|\le R/2\), the packets in
Theorem~\ref{thm:v6v80-local-potential-no-go} lie in the amplitude box
\(\|\phi\|_\infty\le R\), yet their displayed reduced energies have no
lower bound uniform in the spectral cutoff.
\end{corollary}

\begin{proof}
The amplitude inclusion is immediate and the sequence belongs to
\(\mathcal B_{\sqrt{N^2+1}}\).  Apply
Theorem~\ref{thm:v6v80-local-potential-no-go}.
\end{proof}

\begin{firewall}[Scope of the spectral no-go]
\label{fw:v6v80-linearized-scope}
Theorem~\ref{thm:v6v80-local-potential-no-go} is an exact theorem for the
displayed quadratic-source reduced scalar model.  Along its fixed-amplitude
sequence, the induced low-mode metric coefficient predicted by the
linearized constraint grows like the source, namely \(A^2N^2\).  Thus the
sequence eventually leaves every fixed weak-metric neighbourhood.  The
theorem rules out a uniform continuum lower bound obtained by extending
that linearized reduced polynomial unchanged; it does not rule out a
nonlinear Einstein response, a constraint branch that terminates before
those data, or a compensating complete SM source term.
\end{firewall}

\section{Matched weak-response scaling}
\label{sec:v6v80-matched-scaling}

The preceding firewall dictates an upstream step.  Scale the scalar
amplitude down while the two input frequencies increase:
\begin{equation}
 A_N={a\over N},\qquad
 \phi_N^{\mathrm{match}}(x)
 ={a\over N}\bigl(\cos(p_N\cdot x)+\cos(q_N\cdot x)\bigr).
 \label{eq:v6v80-matched-packet}
\end{equation}

\begin{proposition}[Exact matched-scaling ledger]
\label{prop:v6v80-matched-ledger}
For fixed \(a\),
\begin{align}
 \|\phi_N^{\mathrm{match}}\|_\infty&\le {2|a|\over N},
 \label{eq:v6v80-match-sup}\\
 \int|\nabla\phi_N^{\mathrm{match}}|^2d\mu
 &=a^2\left(1+{1\over2N^2}\right),
 \label{eq:v6v80-match-gradient}\\
 \mathfrak Q_{\mathrm{TT}}(\phi_N^{\mathrm{match}})
 &\ge {a^4\over4},
 \label{eq:v6v80-match-TT}\\
 \mathfrak Q_{\mathrm N,2}(\phi_N^{\mathrm{match}})
 &\ge {a^4\over8}
 \left(1+{m^2+2\xi\over N^2}\right)^2.
 \label{eq:v6v80-match-Newton}
\end{align}
The two high-frequency inputs therefore converge uniformly to zero, while
their fixed-output quadratic gravitational sources remain of order
\(a^2\) and their attractive source squares remain of order \(a^4\).
\end{proposition}

\begin{proof}
Substitute \(A=a/N\) in
\eqref{eq:v6v80-gradient-exact},
\eqref{eq:v6v80-QTT-fixed}, and
\eqref{eq:v6v80-QN-fixed}.  The source amplitudes are the square roots of
the corresponding two-mode source-square contributions and hence remain
of order \(a^2\).
\end{proof}

\begin{assessment}[The nonlinear variable is the merged source amplitude]
\label{ass:v6v80-nonlinear-variable}
Pointwise smallness of the scalar field is not the weak-field parameter for
this merger.  The low-output Einstein constraint sees
\(z\sim\kappa a^2\), because two derivatives cancel the factor \(N^{-2}\)
from the matched amplitudes while the output inverse Laplacian has unit
size.  For small \(|z|\), perturbative elimination must reproduce the
negative quadratic source square already computed.  Whether the branch
continues, folds, changes sign, or acquires a compensating positive metric
term for finite \(z\) is a nonlinear constraint question; no estimate in
V71--V80 answers it.
\end{assessment}

\section{The nonlinear beat-response card}
\label{sec:v6v80-NBRC}

\begin{definition}[Nonlinear beat-response card, NBRC]
\label{def:v6v80-NBRC}
For the matched family \eqref{eq:v6v80-matched-packet}, NBRC requires:
\begin{enumerate}[label=\textup{(NB\arabic*)},leftmargin=5.2em]
\item a declared gauge and a finite-cutoff solution of the full Hamiltonian
      and momentum constraints, retaining the fixed output \(k_0\), both
      high-frequency ancestors, and every harmonic generated by the
      nonlinear metric coefficients;
\item a branch parameterized by the invariant merged amplitude
      \(z=\kappa a^2\), with a proof of its domain, uniqueness or controlled
      multiplicity, and regularity;
\item agreement of the branch's small-\(z\) expansion with the V79 negative
      TT and Newton source squares;
\item an exact physical energy or quadratic form on that branch, including
      lapse, shift, conformal, TT, gauge, Higgs, fermion, and ghost terms
      that occur at the same order, with no term estimated before common
      ancestry is assembled;
\item either a lower bound uniform in the input frequency \(N\), or an
      explicit branch obstruction identifying the finite value of \(z\) at
      which the constraint solution ceases to exist; and
\item compatibility with the positive real-lapse/shift rigging map and the
      zero-mode treatment before any continuum claim.
\end{enumerate}
\end{definition}

\begin{theorem}[Why NBRC is now necessary]
\label{thm:v6v80-NBRC-necessary}
Any completion that keeps the physical gravitational couplings fixed,
contains the V79 scalar subsector, and claims a cutoff-uniform lower bound
must supply NBRC or a stronger result that implies its frequency-uniform
energy row.  A derivative-free local scalar potential and the linearized
constraint elimination cannot supply that row.
\end{theorem}

\begin{proof}
The local potential has an \(N\)-independent value on the fixed-amplitude
beat family by Lemma~\ref{lem:v6v80-Haar-law}; the scalar kinetic energy is
only quadratic in \(N\), while the attractive linearized source squares are
quartic by Theorem~\ref{thm:v6v80-local-potential-no-go}.  Thus those terms
cannot give a uniform lower bound.  Firewall~\ref{fw:v6v80-linearized-scope}
shows why extrapolating that conclusion directly to the full Einstein
system would be invalid: the fixed-amplitude family exits the linearized
metric regime.  Proposition~\ref{prop:v6v80-matched-ledger} keeps the merged
metric source finite and isolates its invariant amplitude.  A claimed
completion must therefore control the nonlinear constraint and energy on
that family, which is exactly the frequency-uniform row of NBRC.
\end{proof}

\begin{programme}[V81: nonlinear low-output metric response]
\label{prog:v6v80-V81}
\begin{enumerate}[leftmargin=3em]
\item Return upstream to the finite-cutoff CMC conformal constraint system
      rather than extending the linearized source square beyond its domain.
\item Insert the matched beat family and identify the closed set of metric
      Fourier modes generated at the first nonlinear orders, beginning with
      \(0\), \(\pm k_0\), and \(\pm2k_0\).
\item Derive the exact finite-dimensional constraint map and its derivative
      at \(z=0\); use the implicit-function theorem only on a quantitatively
      stated neighbourhood.
\item Compute the reduced physical energy through the first order at which
      nonlinear metric terms can compete with the negative source square.
\item If harmonic proliferation prevents a closed finite system, formulate
      the corresponding weighted convolution algebra and prove the mapping
      estimates before continuing the branch.
\item If the branch or its coercivity becomes circular, go further upstream
      to the conformal-factor sign and zero-mode solvability conditions.
\end{enumerate}
\end{programme}

\begin{firewall}[Continuum claim boundary after V80]
\label{fw:v6v80-continuum-boundary}
V80 proves a general merger-symbol theorem, its exact TT and Newton
specializations, and a fixed-amplitude spectral no-go for the displayed
linearized reduced scalar energy.  It also proves the matched-scaling ledger
that isolates the finite nonlinear source parameter \(z\).  It does not
solve NBRC, the nonlinear Einstein constraints on the beat family, the
zero-mode constraint, the full SM source assembly, renormalization, cutoff
removal, reflection positivity, Lorentzian reconstruction, or a complete
SM--Einstein continuum.
\end{firewall}

\begin{theorem}[Sixth-series V80 result]
\label{thm:v6v80-status}
The V80 audit establishes the common exact-assembly rule without importing
foreign constants.  For every translation-invariant quadratic source, the
two-frequency test gives the necessary bound
\(C_\Lambda\ge w(k_0)\|\beta_N\|^2/18\).  The Einstein TT and Newton
symbols grow quadratically in the input frequency, recovering their sharp
fourth-order regulator loss.  Consequently the displayed linearized reduced
energy tends to \(-\infty\) on a uniformly amplitude-bounded spectral
sequence for every fixed attractive gravitational coefficient and every
derivative-free local scalar potential.  Matching amplitude as \(N^{-1}\)
keeps the nonlinear merged source finite and identifies NBRC, rather than a
further linearized estimate, as the next gate.
\end{theorem}

\begin{proof}
Use Theorem~\ref{thm:v6v80-audit-decision},
Theorem~\ref{thm:v6v80-merger-symbol},
Theorem~\ref{thm:v6v80-Einstein-symbols},
Theorem~\ref{thm:v6v80-local-potential-no-go},
Proposition~\ref{prop:v6v80-matched-ledger}, and
Theorem~\ref{thm:v6v80-NBRC-necessary}, subject to
Firewalls~\ref{fw:v6v80-linearized-scope} and
\ref{fw:v6v80-continuum-boundary}.
\end{proof}

The next compulsory YM/NS audit is V85.

\clearpage
\chapter{Sixth-series V81: matched-beat conformal compactness and the
homogenized nonlinear response}
\label{ch:v6v81-response}

\section{V81 entry point}
\label{sec:v6v81}

V80 isolated the merger amplitude by taking the beat-packet amplitude to
be \(A_N=a/N\).  This version inserts that family into the exact minimally
coupled CMC Lichnerowicz equation of V73.  Its gradient coefficient is now
uniformly bounded, so the existing barrier theorem gives conformal factors
with frequency-independent bounds.  The coefficient nevertheless retains a
fixed low Fourier mode.  Compactness therefore produces a nonconstant
homogenized constraint rather than the homogeneous vacuum equation.

\section{Exact matched gradient density}
\label{sec:v6v81-gradient}

Retain the flat normalized torus, assume
\begin{equation}
 |\bar\sigma|^2=s_*>0,
 \qquad a_*:=b_\tau-16\pi G V(0)>0,
 \label{eq:v6v81-sector}
\end{equation}
and let \(V\) be smooth near zero.  Set
\begin{equation}
 \phi_N(x)={a\over N}
 \left(\cos(Nx_1)+\cos(-Nx_1+x_2)\right),
 \quad w_N=8\pi G|\bar\nabla\phi_N|^2,
 \quad a_N=b_\tau-16\pi G V(\phi_N).
 \label{eq:v6v81-packet}
\end{equation}

\begin{lemma}[Exact slow/fast decomposition]
\label{lem:v6v81-wN}
One has
\begin{align}
 {w_N\over8\pi G a^2}
 ={}&1+{1\over2N^2}+\cos x_2-{1\over2}\cos(2Nx_1)
 \nonumber\\
 &-{1\over2}\left(1+{1\over N^2}\right)
       \cos(2Nx_1-2x_2)-\cos(2Nx_1-x_2).
 \label{eq:v6v81-wN}
\end{align}
Consequently
\begin{equation}
 w_N\stackrel{*}{\rightharpoonup}
 w_\infty(x_2):=8\pi G a^2(1+\cos x_2)
 \quad\hbox{in }L^\infty(\mathbb T^3),
 \qquad a_N\longrightarrow a_*\quad\hbox{uniformly}.
 \label{eq:v6v81-coefficient-limits}
\end{equation}
\end{lemma}

\begin{proof}
With \(p_N=(N,0,0)\), \(q_N=(-N,1,0)\),
\begin{equation}
 |\bar\nabla\phi_N|^2={a^2\over N^2}
 \left[N^2\sin^2(p_N\cdot x)+(N^2+1)\sin^2(q_N\cdot x)
 -2N^2\sin(p_N\cdot x)\sin(q_N\cdot x)\right].
 \label{eq:v6v81-before-trig}
\end{equation}
Use \(\sin^2u=(1-\cos2u)/2\),
\(2\sin u\sin v=\cos(u-v)-\cos(u+v)\), and
\(p_N+q_N=e_2\).  This gives \eqref{eq:v6v81-wN}.  The three fast
characters converge weak-star to zero by the Riemann--Lebesgue lemma.  The
uniform bound \(\|\phi_N\|_\infty\le2|a|/N\) gives the last convergence.
\end{proof}

\begin{assessment}[The nonlinear beat survives]
\label{ass:v6v81-beat}
Although \(\phi_N\to0\) uniformly, its quadratic gradient density does not
converge to its spatial mean: the common-ancestry pair leaves the exact low
mode \(8\pi Ga^2\cos x_2\).  Averaging the density before solving the
constraint would erase the merger found in V79--V80.
\end{assessment}

\section{Uniform solvability and homogenized compactness}
\label{sec:v6v81-compactness}

The exact scalar constraint is
\begin{equation}
 -8\bar\Delta\psi_N-s_*\psi_N^{-7}-w_N\psi_N+a_N\psi_N^5=0,
 \qquad\psi_N>0.
 \label{eq:v6v81-Lichnerowicz-N}
\end{equation}

\begin{theorem}[Frequency-uniform conformal solutions]
\label{thm:v6v81-uniform}
There are \(N_0\) and \(0<c<C<\infty\), independent of \(N\), such that
for every \(N\ge N_0\), equation \eqref{eq:v6v81-Lichnerowicz-N} has a
smooth solution with
\begin{equation}
 c\le\psi_N\le C.
 \label{eq:v6v81-bounds}
\end{equation}
Every positive solution obeys bounds of this form.  In particular
\begin{equation}
 c^6\le\operatorname{Vol}(\psi_N^4\bar g)
 =\int\psi_N^6d\mu_{\bar g}\le C^6.
 \label{eq:v6v81-volume}
\end{equation}
\end{theorem}

\begin{proof}
For large \(N\), uniform convergence in
\eqref{eq:v6v81-coefficient-limits} gives
\(a_*/2\le a_N\le3a_*/2\).  Directly from the gradient formula,
\begin{equation}
 0\le w_N\le8\pi Ga^2(2+N^{-1})^2.
 \label{eq:v6v81-w-bound}
\end{equation}
Apply Lemma~\ref{lem:v6v73-barriers} and
Theorem~\ref{thm:v6v73-existence} with \(s_0=s_1=s_*\) and these common
coefficient bounds.  Proposition~\ref{prop:v6v73-apriori} gives the same
type of bounds for every solution.  Integrating their sixth powers proves
\eqref{eq:v6v81-volume}.
\end{proof}

\begin{theorem}[Exact nonlinear homogenized limit]
\label{thm:v6v81-limit}
From every sequence of positive solutions of
\eqref{eq:v6v81-Lichnerowicz-N}, one may extract a subsequence converging
strongly in \(C^{1,\alpha}(\mathbb T^3)\), for every \(\alpha<1\), and
weakly in every finite \(W^{2,p}\), to a smooth positive solution of
\begin{equation}
 -8\bar\Delta\psi_\infty-s_*\psi_\infty^{-7}
 -8\pi Ga^2(1+\cos x_2)\psi_\infty+a_*\psi_\infty^5=0.
 \label{eq:v6v81-limit-equation}
\end{equation}
The pointwise and volume bounds of Theorem~\ref{thm:v6v81-uniform} pass to
the limit.
\end{theorem}

\begin{proof}
The right side of
\begin{equation}
 -8\bar\Delta\psi_N
 =s_*\psi_N^{-7}+w_N\psi_N-a_N\psi_N^5
 \label{eq:v6v81-elliptic}
\end{equation}
is uniformly bounded in \(L^\infty\).  Torus elliptic estimates, including
the already bounded zero mode, give uniform \(W^{2,p}\) bounds.  Rellich
compactness and Sobolev embedding give the asserted subsequence.  All
nonlinear powers and \(a_N\psi_N^5\) then converge uniformly.  For a smooth
test function \(\zeta\),
\begin{align}
 \int(w_N\psi_N-w_\infty\psi_\infty)\zeta
 ={}&\int w_N(\psi_N-\psi_\infty)\zeta
 \nonumber\\
 &+\int(w_N-w_\infty)\psi_\infty\zeta\longrightarrow0
 \label{eq:v6v81-product-limit}
\end{align}
by strong convergence, the uniform bound on \(w_N\), and
\eqref{eq:v6v81-coefficient-limits}.  Passage to the weak equation proves
\eqref{eq:v6v81-limit-equation}; smoothness follows by bootstrapping.
\end{proof}

\begin{corollary}[A sufficient uniqueness condition]
\label{cor:v6v81-unique}
If all positive solutions of \eqref{eq:v6v81-limit-equation} lie in
\([c,C]\) and
\begin{equation}
 7s_*C^{-8}+5a_*c^4-16\pi Ga^2>0,
 \label{eq:v6v81-unique-condition}
\end{equation}
then that equation has at most one positive solution.
\end{corollary}

\begin{proof}
The derivative in \(t\) of
\(-s_*t^{-7}-8\pi Ga^2(1+\cos x_2)t+a_*t^5\) is bounded below by the
left side of \eqref{eq:v6v81-unique-condition}.  Subtract two solutions,
multiply the difference equation by their difference, and integrate.  The
gradient term and the positive divided-difference term must both vanish.
\end{proof}

\section{Ancestry indices replace a finite-mode closure}
\label{sec:v6v81-ancestry}

Put \(\theta=p_N\cdot x\), \(\eta=q_N\cdot x\).  The ancestry character
\(e^{i(r\theta+s\eta)}\) has physical momentum and Laplacian cost
\begin{equation}
 K_N(r,s)=((r-s)N,s,0),
 \qquad |K_N(r,s)|^2=N^2(r-s)^2+s^2.
 \label{eq:v6v81-ancestry-symbol}
\end{equation}

\begin{lemma}[Diagonal ancestry and harmonic proliferation]
\label{lem:v6v81-ancestry}
The diagonal indices \((r,r)\) are precisely the slow modes
\(rk_0=(0,r,0)\), while every fixed off-diagonal index has cost at least
\(N^2\).  No finite set of ancestry indices containing \((1,1)\) is closed
under the products occurring in \(\psi^5\) and \(\psi^{-7}\).
\end{lemma}

\begin{proof}
The first statement is immediate from
\eqref{eq:v6v81-ancestry-symbol}.  If a finite set \(S\subset\mathbb Z^2\)
contains a nonzero vector \(v\) and is closed under addition, then it must
contain \(nv\) for every positive integer \(n\), a contradiction.  Fourier
supports multiply by Minkowski addition, so a nontrivial finite support
cannot be an algebra for the conformal powers.
\end{proof}

\begin{proposition}[An exact ancestry-lattice solution]
\label{prop:v6v81-ancestry-solution}
For each large \(N\), at least one solution in
Theorem~\ref{thm:v6v81-uniform} has the form
\begin{equation}
 \psi_N(x)=\Psi_N(p_N\cdot x,q_N\cdot x)
 \label{eq:v6v81-phase-form}
\end{equation}
and is independent of \(x_3\).  Multiplication is convolution in
\((r,s)\), and elliptic inversion uses the exact denominator
\eqref{eq:v6v81-ancestry-symbol}.
\end{proposition}

\begin{proof}
The coefficients have this phase form.  Start the V73 monotone iteration
from constant barriers.  Pointwise nonlinearities preserve the phase
subspace, and \((-8\bar\Delta+M)^{-1}\) preserves it because it is diagonal
on the momenta \(K_N(r,s)\).  Hence every iterate and its limit have the
stated form.
\end{proof}

\begin{assessment}[Exact replacement of the finite-mode proposal]
\label{ass:v6v81-ancestry}
The modes \(0,\pm k_0,\pm2k_0\) locate the first responses but cannot form
an exact nonlinear system.  The two-index lattice is the exact replacement:
its diagonal contains the low-output tower, while its off-diagonal records
the high-frequency ancestors.  The compactness theorem sums this infinite
algebra without an exponentially weighted norm that would grow with \(N\).
\end{assessment}

\section{Analytic homogenized response}
\label{sec:v6v81-response}

Write
\begin{equation}
 \varepsilon=8\pi Ga^2,
 \qquad \psi_0=(s_*/a_*)^{1/12},
 \qquad \mu_0=12a_*\psi_0^4.
 \label{eq:v6v81-parameters}
\end{equation}
On functions of \(y=x_2\), the limit equation is
\begin{equation}
 -8\psi''-s_*\psi^{-7}-\varepsilon(1+\cos y)\psi+a_*\psi^5=0.
 \label{eq:v6v81-one-dimensional}
\end{equation}

\begin{theorem}[Analytic small-merger branch]
\label{thm:v6v81-branch}
For each \(0<\alpha<1\), equation
\eqref{eq:v6v81-one-dimensional} has a unique real-analytic positive branch
\(\varepsilon\mapsto\psi_\varepsilon\in C^{2,\alpha}(\mathbb T)\) near
\((0,\psi_0)\), and
\begin{equation}
 \boxed{\quad
 \psi_\varepsilon(y)=\psi_0+\varepsilon\psi_0
 \left({1\over\mu_0}+{\cos y\over\mu_0+8}\right)
 +O_{C^{2,\alpha}}(\varepsilon^2).
 \quad}
 \label{eq:v6v81-branch-expansion}
\end{equation}
In particular
\begin{equation}
 \widehat\psi_\varepsilon(\pm k_0)
 ={\varepsilon\psi_0\over2(\mu_0+8)}+O(\varepsilon^2).
 \label{eq:v6v81-k0}
\end{equation}
\end{theorem}

\begin{proof}
The derivative of the left side of
\eqref{eq:v6v81-one-dimensional} at \((\psi_0,0)\) is
\begin{equation}
 L_0=-8{d^2\over dy^2}+7s_*\psi_0^{-8}+5a_*\psi_0^4
 =-8{d^2\over dy^2}+\mu_0.
 \label{eq:v6v81-L0}
\end{equation}
It is an isomorphism from \(C^{2,\alpha}\) to \(C^{0,\alpha}\), so the
analytic implicit-function theorem applies.  Differentiation gives
\(L_0\dot\psi_0=(1+\cos y)\psi_0\).  The constant and cosine eigenvalues
are \(\mu_0\) and \(\mu_0+8\), which proves
\eqref{eq:v6v81-branch-expansion}--\eqref{eq:v6v81-k0}.
\end{proof}

\begin{proposition}[Auxiliary critical-value expansion]
\label{prop:v6v81-critical-value}
The one-dimensional V73 constraint functional is
\begin{equation}
 \mathcal I_\varepsilon(\psi)=\int\left[
 4|\psi'|^2+{s_*\over6}\psi^{-6}
 -{\varepsilon\over2}(1+\cos y)\psi^2+{a_*\over6}\psi^6\right]d\mu.
 \label{eq:v6v81-functional}
\end{equation}
For \(J(\varepsilon)=\mathcal I_\varepsilon(\psi_\varepsilon)\),
\begin{align}
 J(\varepsilon)={}&J(0)-{\varepsilon\psi_0^2\over2}
 \nonumber\\
 &-{\varepsilon^2\psi_0^2\over2}
 \left({1\over\mu_0}+{1\over2(\mu_0+8)}\right)+O(\varepsilon^3).
 \label{eq:v6v81-critical-value}
\end{align}
\end{proposition}

\begin{proof}
Stationarity gives
\(J'(\varepsilon)=-\frac12\int(1+\cos y)\psi_\varepsilon^2d\mu\).
Use \eqref{eq:v6v81-branch-expansion},
\(\int\cos y=0\), and \(\int\cos^2y=1/2\), then apply Taylor's theorem.
\end{proof}

\begin{firewall}[Auxiliary versus physical energy]
\label{fw:v6v81-auxiliary}
The negative coefficient in \eqref{eq:v6v81-critical-value} belongs to the
Lichnerowicz solvability functional.  That functional is not the gauge-fixed
Einstein--Hilbert action or the physical ADM Hamiltonian.  Its sign cannot be
identified with the V79 source-square energy without deriving the York,
canonical, and measure terms on the same branch.
\end{firewall}

\begin{theorem}[NBRC progress in V81]
\label{thm:v6v81-NBRC}
In the minimally coupled time-symmetric scalar sector with constant positive
TT floor and strict CMC gap, the matched constraint has frequency-uniform
solutions and volumes; every solution sequence is precompact and has the
exact homogenized equation \eqref{eq:v6v81-limit-equation}; the ancestry
lattice gives the exact source algebra; and the homogenized small-source
branch has the explicit response \eqref{eq:v6v81-k0}.  These results close
the existence, frequency-compactness, ancestry, and conformal-response rows
of NBRC in this sector.  They do not close its physical-energy row.
\end{theorem}

\begin{proof}
Use Theorems~\ref{thm:v6v81-uniform},
\ref{thm:v6v81-limit}, Proposition~\ref{prop:v6v81-ancestry-solution}, and
Theorem~\ref{thm:v6v81-branch}, subject to
Firewall~\ref{fw:v6v81-auxiliary}.
\end{proof}

\begin{programme}[V82: canonical energy on the homogenized branch]
\label{prog:v6v81-V82}
\begin{enumerate}[leftmargin=3em]
\item Keep reference volume, physical volume, York time, and TT floor as
      distinct variables; derive the ADM canonical energy restricted to
      \(\psi_\varepsilon\).
\item Compute through order \(\varepsilon^2\), including the trace-momentum
      term required by the integrated Hamiltonian constraint.
\item Determine whether fixed physical volume forces a York-time response
      already at order \(\varepsilon\).
\item Reinsert scalar canonical momentum and solve the momentum constraint,
      extending the ancestry lattice if its vector inverse creates another
      merger.
\item Keep Lorentzian constraint energy, positive real-lapse rigging, and
      any Euclidean transfer operator distinct.
\end{enumerate}
\end{programme}

\begin{firewall}[Claim boundary after V81]
\label{fw:v6v81-boundary}
V81 treats one minimally coupled, time-symmetric scalar sector and assumes a
constant positive TT floor and strict CMC gap.  It does not fix physical
volume, solve the coupled momentum constraint, derive the physical reduced
Hamiltonian, include the remaining SM fields, construct a positive transfer
operator, remove a regulator, or prove a full SM--Einstein continuum.
\end{firewall}

\begin{theorem}[Sixth-series V81 result]
\label{thm:v6v81-status}
The matched beat has exact gradient-density limit
\(8\pi Ga^2(1+\cos x_2)\).  Its conformal solutions and volumes are uniform
in frequency and converge subsequentially to a nonlinear homogenized
Lichnerowicz equation.  The two-index ancestry lattice is the exact source
algebra, and the limiting small-merger branch has the explicit \(k_0\)
response \eqref{eq:v6v81-k0}.  The remaining bottleneck is the canonical
physical energy and York response on this same branch.
\end{theorem}

\begin{proof}
Theorem~\ref{thm:v6v81-NBRC}, subject to
Firewalls~\ref{fw:v6v81-auxiliary} and \ref{fw:v6v81-boundary}.
\end{proof}

The next compulsory YM/NS audit remains V85.
\clearpage
\chapter{Sixth-series V82: compact-slice Hamiltonian reduction and the
fixed-volume York response}
\label{ch:v6v82-York}

\section{V82 entry point}
\label{sec:v6v82}

V81 closed the conformal-solvability and frequency-compactness part of the
matched beat problem.  Its next programme asked for the ADM canonical
energy on the homogenized branch.  On the closed spatial torus this request
must first be corrected: the ADM Hamiltonian has no boundary charge and is
a sum of constraints, hence it vanishes on the exact constraint surface.
A nonzero physical generator appears only after a clock or a boundary is
chosen.  This version declares CMC York time as the clock, derives its
canonical conjugate from the ADM one-form, and computes the response of the
deparametrized volume Hamiltonian.  It then proves that fixing physical
volume absorbs precisely the constant part of the beat source into York
time while leaving the \(k_0\) conformal response intact.

\section{Why there is no on-shell ADM energy on the closed torus}
\label{sec:v6v82-closed-Hamiltonian}

Use the canonical convention
\begin{equation}
 \pi^{ij}={\sqrt g\over16\pi G}(K^{ij}-Kg^{ij}),
 \qquad K=\tau.
 \label{eq:v6v82-pi-convention}
\end{equation}

\begin{theorem}[Closed-slice ADM Hamiltonian]
\label{thm:v6v82-ADM-zero}
On \(\mathbb T^3\), before gauge fixing,
\begin{equation}
 H_{\mathrm{ADM}}[N,X]
 =\int_{\mathbb T^3}
 \left(N\mathcal H+X^i\mathcal H_i\right)d^3x,
 \label{eq:v6v82-ADM-Hamiltonian}
\end{equation}
with no boundary term.  Hence
\begin{equation}
 H_{\mathrm{ADM}}[N,X]=0
 \label{eq:v6v82-ADM-onshell}
\end{equation}
on every solution of the Hamiltonian and momentum constraints, for every
lapse \(N\) and shift \(X\).  Substituting a Lichnerowicz solution into the
constraint functional therefore does not produce a nonzero physical ADM
energy.
\end{theorem}

\begin{proof}
The Legendre transform of the Einstein--matter action gives the bulk
Hamiltonian \eqref{eq:v6v82-ADM-Hamiltonian}; integrations by parts add only
the standard spatial boundary charge.  The torus is compact and has empty
boundary, so that charge is absent.  Variation of \(N\) and \(X\) imposes
\(\mathcal H=0\) and \(\mathcal H_i=0\).  Inserting them in the bulk
Hamiltonian proves \eqref{eq:v6v82-ADM-onshell}.  The V73 functional has the
Lichnerowicz equation as its Euler equation, but is not a boundary charge or
a deparametrized Hamiltonian; its critical value is therefore not changed
into an ADM energy by imposing its Euler equation.
\end{proof}

\begin{correction}[The NBRC physical-energy row needs a clock]
\label{corr:v6v82-clock}
On a closed universe, a request for a lower bound on ``the on-shell ADM
energy'' is vacuous because that quantity is zero.  The nontrivial choices
are a deparametrized Hamiltonian after an internal clock is fixed, a
quasilocal energy after a spatial boundary is introduced, or a transfer
generator derived from a reflection-positive Euclidean construction.  V82
uses the first choice and does not identify it with the other two.
\end{correction}

\section{York time and physical volume are a canonical pair}
\label{sec:v6v82-canonical-pair}

Let
\begin{equation}
 \mathcal V(g)=\int_{\mathbb T^3}\sqrt g\,d^3x.
 \label{eq:v6v82-volume}
\end{equation}

\begin{lemma}[Trace part of the ADM one-form]
\label{lem:v6v82-one-form}
On the CMC slice, the trace part of
\(\Theta_{\mathrm{ADM}}=\int\pi^{ij}\delta g_{ij}\,d^3x\) is
\begin{equation}
 \Theta_{\mathrm{tr}}
 =-{\tau\over12\pi G}\,\delta\mathcal V.
 \label{eq:v6v82-trace-one-form}
\end{equation}
Equivalently, if
\begin{equation}
 T=-{\tau\over12\pi G},
 \label{eq:v6v82-York-clock}
\end{equation}
then, modulo the exact variation \(\delta(T\mathcal V)\),
\begin{equation}
 \Theta_{\mathrm{tr}}=-\mathcal V\,\delta T.
 \label{eq:v6v82-York-pair}
\end{equation}
Thus the momentum conjugate to the York clock \(T\) is
\(P_T=-\mathcal V\).
\end{lemma}

\begin{proof}
Taking the trace of \eqref{eq:v6v82-pi-convention} gives
\begin{equation}
 \pi:=g_{ij}\pi^{ij}=-{\sqrt g\over8\pi G}\tau.
 \label{eq:v6v82-pi-trace}
\end{equation}
The trace projection of a metric variation is
\((g^{ij}\delta g_{ij})g_{kl}/3\), so its contribution to the canonical
one-form is
\[
 \int {1\over3}\pi g^{ij}\delta g_{ij},d^3x
 =-{\tau\over24\pi G}\int\sqrt g\,g^{ij}\delta g_{ij},d^3x
 =-{\tau\over12\pi G}\delta\mathcal V.
\]
This is \(T\delta\mathcal V\).  The identity
\(T\delta\mathcal V=\delta(T\mathcal V)-\mathcal V\delta T\) proves the
canonical-pair statement.
\end{proof}

\begin{definition}[York deparametrized Hamiltonian]
\label{def:v6v82-York-Hamiltonian}
After the Hamiltonian constraint is solved for
\(P_T=-\mathcal V\) and \(T\) is used as time, define
\begin{equation}
 H_Y:=\mathcal V.
 \label{eq:v6v82-HY}
\end{equation}
The opposite time orientation reverses the generator's sign.  All
statements below use the orientation in \eqref{eq:v6v82-HY}.
\end{definition}

\section{Two-parameter homogenized constraint}
\label{sec:v6v82-two-parameter}

Write \(h(y)=1+\cos y\) and promote the positive CMC gap to a parameter
\(\mathfrak a\).  The homogenized equation becomes
\begin{equation}
 \mathcal F(\psi,\epsilon,\mathfrak a)
 :=-8\psi''-s_*\psi^{-7}-\epsilon h\psi
 +\mathfrak a\psi^5=0,
 \label{eq:v6v82-F}
\end{equation}
where \(\epsilon=8\pi Ga^2\).  Fix
\begin{equation}
 \mathfrak a_0>0,
 \qquad \psi_0=(s_*/\mathfrak a_0)^{1/12},
 \qquad \mu_0=12\mathfrak a_0\psi_0^4.
 \label{eq:v6v82-base}
\end{equation}

\begin{theorem}[Analytic two-parameter conformal branch]
\label{thm:v6v82-two-parameter}
In a neighbourhood of \((0,\mathfrak a_0)\), there is a unique analytic
branch \(\psi(\epsilon,\mathfrak a)\) near \(\psi_0\) solving
\eqref{eq:v6v82-F}.  At the base point,
\begin{align}
 \partial_\epsilon\psi
 &=\psi_0\left({1\over\mu_0}+{\cos y\over\mu_0+8}\right),
 \label{eq:v6v82-psi-epsilon}\\
 \partial_{\mathfrak a}\psi
 &=-{\psi_0^5\over\mu_0}
 =-{\psi_0\over12\mathfrak a_0}.
 \label{eq:v6v82-psi-a}
\end{align}
\end{theorem}

\begin{proof}
The derivative in \(\psi\) at the base point is the positive isomorphism
\[
 L_0=-8{d^2\over dy^2}+\mu_0
\]
from V81, so the analytic implicit-function theorem applies.  Parameter
differentiation gives
\[
 L_0\partial_\epsilon\psi=h\psi_0,
 \qquad
 L_0\partial_{\mathfrak a}\psi=-\psi_0^5.
\]
The constant and cosine eigenvalues of \(L_0\) are \(\mu_0\) and
\(\mu_0+8\), proving both formulas.
\end{proof}

Define the volume of this branch by
\begin{equation}
 \mathcal V(\epsilon,\mathfrak a)
 :=\int_{\mathbb T}\psi(\epsilon,\mathfrak a)^6d\mu.
 \label{eq:v6v82-branch-volume}
\end{equation}

\begin{corollary}[Exact first derivatives of the York Hamiltonian]
\label{cor:v6v82-volume-derivatives}
At \((0,\mathfrak a_0)\),
\begin{align}
 \mathcal V_0&=\psi_0^6,
 \label{eq:v6v82-V0}\\
 \partial_\epsilon\mathcal V
 &={6\psi_0^6\over\mu_0}>0,
 \label{eq:v6v82-Vepsilon}\\
 \partial_{\mathfrak a}\mathcal V
 &=-{6\psi_0^{10}\over\mu_0}
 =-{\psi_0^6\over2\mathfrak a_0}<0.
 \label{eq:v6v82-Va}
\end{align}
Thus at fixed York time and fixed remaining sector data,
\begin{equation}
 H_Y(\epsilon)
 =\psi_0^6+{6\epsilon\psi_0^6\over\mu_0}+O(\epsilon^2).
 \label{eq:v6v82-HY-response}
\end{equation}
\end{corollary}

\begin{proof}
Differentiate \eqref{eq:v6v82-branch-volume} and insert
\eqref{eq:v6v82-psi-epsilon}--\eqref{eq:v6v82-psi-a}.  The cosine has zero
mean.  Definition~\ref{def:v6v82-York-Hamiltonian} gives the last formula.
\end{proof}

\section{Fixed physical volume forces a York-time response}
\label{sec:v6v82-fixed-volume}

\begin{theorem}[Fixed-volume gap and conformal response]
\label{thm:v6v82-fixed-volume}
There is a unique analytic function \(\mathfrak a(\epsilon)\) near zero
such that
\begin{equation}
 \mathfrak a(0)=\mathfrak a_0,
 \qquad
 \mathcal V(\epsilon,\mathfrak a(\epsilon))=\psi_0^6.
 \label{eq:v6v82-fixed-volume-condition}
\end{equation}
Its expansion and the corresponding conformal factor are
\begin{align}
 \mathfrak a(\epsilon)
 &=\mathfrak a_0+\epsilon\psi_0^{-4}+O(\epsilon^2),
 \label{eq:v6v82-gap-response}\\
 \psi_{\mathrm{vol},\epsilon}(y)
 &=\psi_0+{\epsilon\psi_0\over\mu_0+8}\cos y
 +O_{C^{2,\alpha}}(\epsilon^2).
 \label{eq:v6v82-fixed-volume-psi}
\end{align}
The gap shift cancels exactly the constant part of the fixed-York conformal
response, but it does not cancel the \(k_0\) mode.
\end{theorem}

\begin{proof}
Equation \eqref{eq:v6v82-Va} is nonzero, so the analytic implicit-function
theorem applied to \eqref{eq:v6v82-fixed-volume-condition} gives the branch.
Differentiating that condition and using
\eqref{eq:v6v82-Vepsilon}--\eqref{eq:v6v82-Va} yields
\begin{equation}
 \mathfrak a'(0)
 =-{\partial_\epsilon\mathcal V\over
        \partial_{\mathfrak a}\mathcal V}
 =\psi_0^{-4}.
 \label{eq:v6v82-a-prime}
\end{equation}
The total conformal derivative is
\[
 \partial_\epsilon\psi
 +\partial_{\mathfrak a}\psi\,\mathfrak a'(0)
 =\psi_0\left({1\over\mu_0}+{\cos y\over\mu_0+8}\right)
 -{\psi_0^5\over\mu_0}\psi_0^{-4},
\]
whose two constant terms cancel.  This proves
\eqref{eq:v6v82-fixed-volume-psi}.
\end{proof}

Return to
\begin{equation}
 \mathfrak a={2\tau^2\over3}-2\Lambda-16\pi GV(0).
 \label{eq:v6v82-a-tau}
\end{equation}

\begin{corollary}[Exact first York-square response]
\label{cor:v6v82-tau-response}
At fixed physical volume,
\begin{align}
 \tau(\epsilon)^2
 &=\tau_0^2+{3\epsilon\over2\psi_0^4}+O(\epsilon^2),
 \label{eq:v6v82-tau-square}\\
 -{2\tau(\epsilon)^2\over3}
 &=-{2\tau_0^2\over3}-\epsilon\psi_0^{-4}+O(\epsilon^2).
 \label{eq:v6v82-trace-response}
\end{align}
If \(\tau_0\ne0\), a chosen local sign branch also satisfies
\begin{equation}
 \tau'(0)={3\over4\tau_0\psi_0^4}.
 \label{eq:v6v82-tau-prime}
\end{equation}
\end{corollary}

\begin{proof}
Differentiate \eqref{eq:v6v82-a-tau} and use
\eqref{eq:v6v82-a-prime}.  The trace identity is multiplication by
\(-2/3\); differentiating a local square-root branch gives the last line.
\end{proof}

\begin{assessment}[Zero mode and beat mode separate]
\label{ass:v6v82-mode-separation}
The matched scalar source has a constant component and a \(k_0\) component.
On a fixed-York slice both increase the conformal factor.  On a fixed-volume
slice the integrated constraint transfers the constant component into the
negative DeWitt trace direction through
\eqref{eq:v6v82-trace-response}; the nonzero beat component remains as
\eqref{eq:v6v82-fixed-volume-psi}.  Thus the zero mode cannot be discarded
or treated by the nonzero-mode inverse Laplacian.
\end{assessment}

\section{The compact-slice generator card}
\label{sec:v6v82-CSGC}

\begin{definition}[Compact-slice generator card, CSGC]
\label{def:v6v82-CSGC}
CSGC requires:
\begin{enumerate}[label=\textup{(CG\arabic*)},leftmargin=5.2em]
\item a declared physical clock, boundary condition, or Euclidean transfer
      construction, since the unreduced closed-slice ADM Hamiltonian
      vanishes on shell;
\item derivation of the corresponding canonical pair and generator, such as
      \((T,P_T)=(-\tau/(12\pi G),-\mathcal V)\);
\item separate treatment of the integrated zero-mode constraint and every
      nonzero source merger;
\item a self-adjoint lower-bounded quantization of the chosen reduced
      generator, not of the auxiliary Lichnerowicz functional; and
\item compatibility of that generator with the real-lapse/shift rigging map
      and with transitions between CMC charts.
\end{enumerate}
\end{definition}

\begin{theorem}[CSGC progress]
\label{thm:v6v82-CSGC}
For the homogenized scalar CMC branch, V82 closes
CSGC\textup{(CG1)}--\textup{(CG3)} locally: York time is declared, its
canonical volume momentum is derived, and the fixed-volume calculation
separates the zero and \(k_0\) responses.  Rows \textup{(CG4)}--
\textup{(CG5)} remain open.
\end{theorem}

\begin{proof}
Use Theorem~\ref{thm:v6v82-ADM-zero},
Lemma~\ref{lem:v6v82-one-form},
Theorem~\ref{thm:v6v82-fixed-volume}, and
Corollary~\ref{cor:v6v82-tau-response}.
\end{proof}

\begin{firewall}[Local York reduction is not a continuum Hamiltonian]
\label{fw:v6v82-local-only}
The York branch is local in \((\epsilon,\mathfrak a)\), assumes CMC
admissibility and a positive TT floor, and uses a classical constraint
solution.  No global CMC atlas, quantum self-adjoint generator, unitary
evolution, reflection-positive transfer operator, or regulator removal has
been constructed.  The sign of the auxiliary critical value in V81 is not
used as a physical energy sign.
\end{firewall}

\begin{programme}[V83: matter momentum and the vector constraint]
\label{prog:v6v82-V83}
\begin{enumerate}[leftmargin=3em]
\item Add a matched scalar canonical momentum and derive its exact momentum
      density with the same two ancestry indices.
\item Solve the flat CMC vector constraint on the orthogonal complement of
      its Killing zero modes and compute the high--high to low longitudinal
      response.
\item State the exact mean-current compatibility condition before applying
      the vector inverse.
\item Insert the longitudinal momentum norm into the Lichnerowicz equation
      and determine whether the V81 compactness bounds remain uniform.
\item If the current creates a nonzero mean or destroys the TT floor, return
      upstream to the global momentum and conformal-weight assignments.
\end{enumerate}
\end{programme}

\begin{firewall}[Claim boundary after V82]
\label{fw:v6v82-boundary}
V82 treats a local, classical, time-symmetric scalar CMC branch on a closed
torus.  It does not solve the matter momentum constraint, quantize York
evolution, assemble the remaining SM sources, prove positivity of a transfer
operator, remove a cutoff, or construct the full SM--Einstein continuum.
\end{firewall}

\begin{theorem}[Sixth-series V82 result]
\label{thm:v6v82-status}
The unreduced ADM Hamiltonian on the closed torus vanishes on the constraint
surface, so a physical generator requires an additional choice.  With York
time as clock, physical volume is the reduced Hamiltonian in the declared
orientation.  At fixed York time the matched beat increases that volume at
first order.  At fixed physical volume the required response is
\(\mathfrak a'(0)=\psi_0^{-4}\), equivalently
\((\tau^2)'(0)=3/(2\psi_0^4)\); this absorbs the constant source into the
DeWitt trace while leaving the explicit \(k_0\) conformal mode.  CSGC records
the remaining quantum-generator requirements.
\end{theorem}

\begin{proof}
Use Theorem~\ref{thm:v6v82-ADM-zero},
Lemma~\ref{lem:v6v82-one-form},
Corollary~\ref{cor:v6v82-volume-derivatives},
Theorem~\ref{thm:v6v82-fixed-volume}, and
Corollary~\ref{cor:v6v82-tau-response}, subject to
Firewalls~\ref{fw:v6v82-local-only} and \ref{fw:v6v82-boundary}.
\end{proof}

The next compulsory YM/NS audit remains V85.
\clearpage
\chapter{Sixth-series V83: matched scalar momentum, the vector constraint,
and the momentum-coupled homogenized branch}
\label{ch:v6v83-momentum}

\section{V83 entry point}
\label{sec:v6v83}

V81--V82 used time-symmetric scalar data.  This continuation introduces a
canonical scalar momentum whose two phase ancestors create a nonzero
low-output matter current, verifies the torus zero-mode compatibility
condition, solves the CMC vector constraint exactly in Fourier variables,
and then inserts both the longitudinal gravitational momentum and scalar
kinetic momentum into the Lichnerowicz equation.  The matched scaling again
has uniform barriers and a compact homogenized limit, now with an explicit
momentum-dependent singular coefficient.

\section{Conformal momentum conventions}
\label{sec:v6v83-conventions}

Write
\begin{equation}
 K^{ij}=\psi^{-10}(\bar\sigma+\bar L W)^{ij}
 +{\tau\over3}g^{ij},
 \label{eq:v6v83-K}
\end{equation}
where
\begin{equation}
 (\bar L W)_{ij}=\partial_iW_j+\partial_jW_i
 -{2\over3}\delta_{ij}\partial_kW_k,
 \qquad
 (\bar\Delta_LW)_i=\Delta W_i+{1\over3}\partial_i\partial_jW_j.
 \label{eq:v6v83-vector-operator}
\end{equation}
Let \(P\) be the scalar canonical momentum density relative to the flat
reference volume.  Then the physical normal derivative and contravariant
momentum density are
\begin{equation}
 \Pi_\phi=\psi^{-6}P,
 \qquad
 j^i=-\Pi_\phi\nabla^i\phi
 =\psi^{-10}\bar j^i,
 \qquad
 \bar j^i=-P\,\bar\nabla^i\phi.
 \label{eq:v6v83-current-weight}
\end{equation}

\begin{lemma}[Flat CMC vector constraint]
\label{lem:v6v83-vector-constraint}
With the convention above and constant \(\tau\), the physical momentum
constraint is equivalent to
\begin{equation}
 \bar\Delta_LW=8\pi G\bar j.
 \label{eq:v6v83-vector-constraint}
\end{equation}
It has a mean-zero solution on \(\mathbb T^3\) exactly when
\begin{equation}
 \int_{\mathbb T^3}\bar j\,d\mu_{\bar g}=0.
 \label{eq:v6v83-mean-current}
\end{equation}
The solution is then unique among mean-zero vector fields.
\end{lemma}

\begin{proof}
The conformal divergence identity for a trace-free tensor gives
\(\nabla_j[\psi^{-10}(\bar\sigma+\bar LW)^{ij}]
=\psi^{-10}\bar\nabla_j(\bar\sigma+\bar LW)^{ij}\).
The CMC trace term has zero divergence, and \(\bar\sigma\) is transverse.
Using \eqref{eq:v6v83-current-weight} gives
\eqref{eq:v6v83-vector-constraint}.  The kernel of
\(\bar\Delta_L\) on the flat torus consists of constant vector fields.
Fredholm orthogonality is therefore exactly
\eqref{eq:v6v83-mean-current}; restriction to their orthogonal complement
gives uniqueness.
\end{proof}

\section{A matched current with fixed low output}
\label{sec:v6v83-matched-current}

Retain \(p_N=(N,0,0)\), \(q_N=(-N,1,0)\) and the matched scalar
\begin{equation}
 \phi_N={a\over N}\bigl(\cos(p_N\cdot x)+\cos(q_N\cdot x)\bigr).
 \label{eq:v6v83-phi}
\end{equation}
Choose the canonical momentum
\begin{equation}
 P_N=b\bigl(\cos(p_N\cdot x)-\cos(q_N\cdot x)\bigr).
 \label{eq:v6v83-P}
\end{equation}

\begin{lemma}[Exact current decomposition and compatibility]
\label{lem:v6v83-current}
The reference current \(\bar j_N=-P_N\bar\nabla\phi_N\) is
\begin{align}
 \bar j_N={ab\over2N}\bigl[&p_N\sin(2p_N\cdot x)
 -q_N\sin(2q_N\cdot x)
 +(q_N-p_N)\sin x_2
 \nonumber\\
 &-(p_N+q_N)\sin((p_N-q_N)\cdot x)\bigr].
 \label{eq:v6v83-current-decomposition}
\end{align}
It has zero mean.  Its low-output part is
\begin{equation}
 \bar j_N^{\mathrm{low}}
 =\left(-ab\,e_1+{ab\over2N}e_2\right)\sin x_2,
 \label{eq:v6v83-current-low}
\end{equation}
which converges uniformly to \(-ab e_1\sin x_2\).
\end{lemma}

\begin{proof}
Expand
\(-P_N\bar\nabla\phi_N\) and use
\(2\cos u\sin v=\sin(u+v)+\sin(v-u)\).
The self terms give the first two modes; the cross terms give the last two.
Every displayed term is a nonzero sine character, proving the mean-current
condition.  Since \(p_N+q_N=e_2\) and
\(q_N-p_N=(-2N,1,0)\), its \(e_2\)-frequency term is
\eqref{eq:v6v83-current-low}.
\end{proof}

\begin{theorem}[Exact low vector response]
\label{thm:v6v83-vector-response}
Let \(W_N\) be the unique mean-zero solution of
\eqref{eq:v6v83-vector-constraint}.  Its low-output component is
\begin{equation}
 W_N^{\mathrm{low}}
 =\left(8\pi Gab\,e_1-{3\pi Gab\over N}e_2\right)\sin x_2.
 \label{eq:v6v83-W-low}
\end{equation}
Moreover
\begin{equation}
 \|\bar L(W_N-W_N^{\mathrm{low}})\|_{L^\infty}
 \le {C_G|ab|\over N},
 \label{eq:v6v83-high-response}
\end{equation}
and therefore
\begin{equation}
 \bar LW_N\longrightarrow\Sigma_{ab}(x_2),
 \qquad
 (\Sigma_{ab})_{12}=(\Sigma_{ab})_{21}
 =8\pi Gab\cos x_2,
 \label{eq:v6v83-Sigma-limit}
\end{equation}
uniformly, with all other limiting components zero.
\end{theorem}

\begin{proof}
On a nonzero Fourier mode \(k\), the negative symbol of
\(\bar\Delta_L\) has eigenvalue \(|k|^2\) on \(k^\perp\) and
\(4|k|^2/3\) on the span of \(k\).  Relative to \(k_0=e_2\), the two
components in \eqref{eq:v6v83-current-low} are transverse and longitudinal.
Inverting their eigenvalues after multiplication by \(8\pi G\) gives
\eqref{eq:v6v83-W-low}.  The self-current modes at \(2p_N,2q_N\) have
bounded amplitude and frequency comparable to \(N\); applying
\(\bar L\bar\Delta_L^{-1}\) costs \(O(N^{-1})\).  The remaining
\(p_N-q_N\) mode has amplitude \(O(N^{-1})\) and hence contributes
\(O(N^{-2})\) after the same operator.  There are only three high modes,
which proves \eqref{eq:v6v83-high-response}.  Applying \(\bar L\) to
\eqref{eq:v6v83-W-low} and taking \(N\to\infty\) gives
\eqref{eq:v6v83-Sigma-limit}.
\end{proof}

\begin{assessment}[The vector carrier retains the beat]
\label{ass:v6v83-vector-beat}
The high self-currents are suppressed by the order-minus-one operator
\(\bar L\bar\Delta_L^{-1}\), but the current at fixed output \(k_0\)
is not.  Its limiting longitudinal gravitational momentum is finite and
nonzero.  The mean-current test is logically prior to this inversion; here
it holds because every current component is a sine character.
\end{assessment}

\section{The momentum-coupled Lichnerowicz equation}
\label{sec:v6v83-coupled-equation}

Choose a constant diagonal transverse-traceless tensor \(\bar\sigma_*\)
with
\begin{equation}
 |\bar\sigma_*|^2=s_*>0.
 \label{eq:v6v83-sigma}
\end{equation}
It is pointwise orthogonal to the off-diagonal tensor
\(\Sigma_{ab}\).  Define
\begin{equation}
 S_N(x)=|\bar\sigma_*+\bar LW_N|^2+8\pi G P_N^2.
 \label{eq:v6v83-SN}
\end{equation}

\begin{lemma}[Exact scalar-momentum conformal equation]
\label{lem:v6v83-full-Lichnerowicz}
The Hamiltonian constraint with the scalar canonical momentum
\eqref{eq:v6v83-P} is
\begin{equation}
 -8\bar\Delta\psi_N-S_N\psi_N^{-7}
 -w_N\psi_N+a_N\psi_N^5=0,
 \qquad\psi_N>0,
 \label{eq:v6v83-full-Lichnerowicz}
\end{equation}
where \(w_N,a_N\) are those of V81.
\end{lemma}

\begin{proof}
The scalar energy density is
\[
 \rho_\phi={1\over2}\psi^{-12}P_N^2
 +{1\over2}\psi^{-4}|\bar\nabla\phi_N|^2+V(\phi_N).
\]
Insert this and \(\bar\sigma_*+\bar LW_N\) in the CMC Hamiltonian
constraint.  After multiplication by \(\psi^5\), the momentum energy
contributes \(-8\pi GP_N^2\psi^{-7}\), the gradient energy contributes
\(-w_N\psi\), and the potential contributes to \(a_N\psi^5\).
\end{proof}

\begin{lemma}[Limit of the singular coefficient]
\label{lem:v6v83-S-limit}
For all sufficiently large \(N\), there are constants
\(0<S_-<S_+<\infty\), independent of \(N\), such that
\begin{equation}
 S_-\le S_N(x)\le S_+.
 \label{eq:v6v83-S-bounds}
\end{equation}
Furthermore
\begin{equation}
 S_N\stackrel{*}{\rightharpoonup}S_\infty(x_2)
 \quad\hbox{in }L^\infty,
 \label{eq:v6v83-S-weak}
\end{equation}
where
\begin{equation}
 S_\infty(x_2)=s_*+2(8\pi Gab)^2\cos^2x_2
 +8\pi Gb^2(1-\cos x_2).
 \label{eq:v6v83-S-infinity}
\end{equation}
\end{lemma}

\begin{proof}
Theorem~\ref{thm:v6v83-vector-response} gives uniform convergence of
\(\bar LW_N\) to the off-diagonal \(\Sigma_{ab}\).  Orthogonality to the
diagonal \(\bar\sigma_*\) implies
\(|\bar\sigma_*+\bar LW_N|^2\ge s_*/2\) for large \(N\), and also gives a
uniform upper bound.  The scalar term is nonnegative and
\(|P_N|\le2|b|\), proving \eqref{eq:v6v83-S-bounds}.  Finally
\begin{equation}
 P_N^2=b^2\left[1+{1\over2}\cos(2p_N\cdot x)
 +{1\over2}\cos(2q_N\cdot x)-\cos x_2
 -\cos((p_N-q_N)\cdot x)\right].
 \label{eq:v6v83-P-square}
\end{equation}
The three fast cosines converge weak-star to zero.  Uniform convergence of
the tensor term and
\(|\Sigma_{ab}|^2=2(8\pi Gab)^2\cos^2x_2\) prove
\eqref{eq:v6v83-S-infinity}.
\end{proof}

\section{Uniform solutions and the coupled homogenized limit}
\label{sec:v6v83-homogenized}

\begin{theorem}[Momentum-coupled frequency compactness]
\label{thm:v6v83-compactness}
For all large \(N\), equation
\eqref{eq:v6v83-full-Lichnerowicz} has a smooth positive solution, and all
positive solutions obey frequency-independent bounds
\begin{equation}
 0<c\le\psi_N\le C<\infty.
 \label{eq:v6v83-psi-bounds}
\end{equation}
Every solution sequence has a subsequence converging strongly in every
\(C^{1,\alpha}\), \(\alpha<1\), to a smooth positive solution of
\begin{equation}
 -8\bar\Delta\psi_\infty-S_\infty(x_2)\psi_\infty^{-7}
 -8\pi Ga^2(1+\cos x_2)\psi_\infty+a_*\psi_\infty^5=0.
 \label{eq:v6v83-homogenized}
\end{equation}
Their physical volumes remain between positive constants independent of
\(N\).
\end{theorem}

\begin{proof}
Use \eqref{eq:v6v83-S-bounds}, the uniform positive bounds on \(a_N\), and
the uniform upper bound on \(w_N\).  The V73 constant barriers and monotone
iteration apply with \(s(x)=S_N(x)\), giving existence and
\eqref{eq:v6v83-psi-bounds}; its maximum-principle argument applies to every
solution.  The right side of the rearranged equation is uniformly bounded
in \(L^\infty\), so the V81 elliptic compactness proof gives strong
\(C^{1,\alpha}\) convergence after a subsequence.  Products with \(S_N\)
and \(w_N\) pass to their weak-star limits because the powers of \(\psi_N\)
converge uniformly.  This proves \eqref{eq:v6v83-homogenized}.  Integrating
the sixth powers of \eqref{eq:v6v83-psi-bounds} gives the volume bounds.
\end{proof}

\begin{corollary}[The vector constraint does not reopen volume escape]
\label{cor:v6v83-no-escape}
For fixed \(a,b\) in the strict CMC sector, the matched scalar momentum
creates a finite low gravitational-momentum response but does not destroy
the regulator-uniform conformal and volume bounds.
\end{corollary}

\begin{proof}
Combine Theorems~\ref{thm:v6v83-vector-response} and
\ref{thm:v6v83-compactness}.
\end{proof}

\begin{definition}[Momentum-coupled beat card, MCBC]
\label{def:v6v83-MCBC}
MCBC requires:
\begin{enumerate}[label=\textup{(MB\arabic*)},leftmargin=5.2em]
\item the exact conformal weights of canonical matter momenta and currents;
\item orthogonality of the total reference current to conformal Killing
      zero modes before vector inversion;
\item an ancestry-resolved solution of the vector constraint, including
      every fixed low output;
\item insertion of both longitudinal gravitational momentum and matter
      momentum energy in the Hamiltonian constraint before taking limits;
\item frequency-uniform conformal and physical-volume bounds; and
\item a physical York generator and positive rigging construction for the
      coupled, non-time-symmetric branch.
\end{enumerate}
\end{definition}

\begin{theorem}[MCBC progress]
\label{thm:v6v83-MCBC}
The matched scalar sector closes MCBC\textup{(MB1)}--\textup{(MB5)} under
the constant diagonal TT-floor and strict-gap hypotheses.  Row
\textup{(MB6)} remains open.
\end{theorem}

\begin{proof}
Use Lemmas~\ref{lem:v6v83-vector-constraint},
\ref{lem:v6v83-current}, Theorems~\ref{thm:v6v83-vector-response},
\ref{thm:v6v83-compactness}, and Lemma~\ref{lem:v6v83-full-Lichnerowicz}.
\end{proof}

\begin{programme}[V84: analytic momentum response and York reduction]
\label{prog:v6v83-V84}
\begin{enumerate}[leftmargin=3em]
\item Treat \(\epsilon=8\pi Ga^2\) and the scalar-momentum parameter
      \(\beta=8\pi Gb^2\) as independent nonnegative variables in
      \eqref{eq:v6v83-homogenized}.
\item Separate the direct scalar-momentum term, of order \(\beta\), from
      the vector-response term, of order \(\epsilon\beta\).
\item Derive the analytic conformal and physical-volume responses at the
      homogeneous base point and compute the fixed-volume York-gap shift.
\item Determine the first order at which the vector constraint changes the
      York generator rather than only the scalar kinetic momentum.
\item Formulate the remaining self-adjoint reduced-generator row before the
      compulsory V85 companion audit.
\end{enumerate}
\end{programme}

\begin{firewall}[Claim boundary after V83]
\label{fw:v6v83-boundary}
V83 treats one scalar canonical momentum with a specially chosen matched
phase and a constant diagonal TT floor.  It does not treat arbitrary total
SM current, nonzero Killing-mode charge, gauge or fermion momentum,
quantization of the York generator, reflection positivity, regulator
removal, or the full SM--Einstein continuum.
\end{firewall}

\begin{theorem}[Sixth-series V83 result]
\label{thm:v6v83-status}
The matched scalar canonical momentum has zero mean current and creates the
finite low vector response \eqref{eq:v6v83-W-low}.  After the longitudinal
gravitational momentum and scalar kinetic momentum are inserted into the
Hamiltonian constraint, uniform barriers and elliptic compactness persist.
The exact homogenized equation is \eqref{eq:v6v83-homogenized}, with
singular coefficient \eqref{eq:v6v83-S-infinity}.  Thus the vector
constraint closes classically in this ancestry-resolved sector without
reopening conformal-volume escape; the coupled York generator remains the
next bottleneck.
\end{theorem}

\begin{proof}
Use Theorems~\ref{thm:v6v83-vector-response} and
\ref{thm:v6v83-compactness}, subject to
Firewall~\ref{fw:v6v83-boundary}.
\end{proof}

The next compulsory YM/NS audit remains V85.
\clearpage
\chapter{Sixth-series V84: two-source York response, the first-order beat
cancellation line, and the mixed vector forcing}
\label{ch:v6v84-two-source}

\section{V84 entry point}
\label{sec:v6v84}

V83 produced the momentum-coupled homogenized Lichnerowicz equation.  Its
two matter amplitudes enter differently.  The scalar-gradient source is
linear in
\(\epsilon=8\pi Ga^2\); the direct scalar kinetic-momentum source is linear
in \(\beta=8\pi Gb^2\); and the longitudinal gravitational momentum created
by the vector constraint first enters at the mixed order
\(\epsilon\beta\).  This version constructs the analytic two-source branch,
computes its conformal and York-volume derivatives, and proves that fixed
physical volume has a unique first-order line on which the two \(k_0\)
conformal responses cancel.  The vector source survives on that line at
second order and generates a \(2k_0\) harmonic.

\section{The two-source homogenized equation}
\label{sec:v6v84-equation}

Set
\begin{equation}
 h_+(y)=1+\cos y,
 \qquad h_-(y)=1-\cos y.
 \label{eq:v6v84-hpm}
\end{equation}
The coefficient \eqref{eq:v6v83-S-infinity} can be written exactly as
\begin{equation}
 S_\infty(y)=s_*+\beta h_-(y)+2\epsilon\beta\cos^2y.
 \label{eq:v6v84-S}
\end{equation}
Indeed
\(\epsilon\beta=(8\pi Gab)^2\).  Promote the CMC gap to
\(\mathfrak a>0\) and define
\begin{equation}
 \mathcal F(\psi,\epsilon,\beta,\mathfrak a)
 =-8\psi''-
 \bigl(s_*+\beta h_-+2\epsilon\beta\cos^2y\bigr)\psi^{-7}
 -\epsilon h_+\psi+\mathfrak a\psi^5.
 \label{eq:v6v84-F}
\end{equation}
At a base gap \(\mathfrak a_0>0\), put
\begin{equation}
 \psi_0=(s_*/\mathfrak a_0)^{1/12},
 \qquad \mu_0=12\mathfrak a_0\psi_0^4,
 \qquad L_0=-8{d^2\over dy^2}+\mu_0.
 \label{eq:v6v84-base}
\end{equation}

\begin{theorem}[Analytic momentum-coupled branch]
\label{thm:v6v84-analytic}
There is a unique real-analytic positive branch
\begin{equation}
 (\epsilon,\beta,\mathfrak a)
 \longmapsto\psi(\epsilon,\beta,\mathfrak a)
 \label{eq:v6v84-branch}
\end{equation}
near \((0,0,\mathfrak a_0,\psi_0)\) satisfying
\(\mathcal F=0\).  Its derivatives at the base point are
\begin{align}
 u:=\partial_\epsilon\psi
 &=\psi_0\left({1\over\mu_0}
 +{\cos y\over\mu_0+8}\right),
 \label{eq:v6v84-u}\\
 v:=\partial_\beta\psi
 &=\psi_0^{-7}\left({1\over\mu_0}
 -{\cos y\over\mu_0+8}\right),
 \label{eq:v6v84-v}\\
 r:=\partial_{\mathfrak a}\psi
 &=-{\psi_0^5\over\mu_0}.
 \label{eq:v6v84-r}
\end{align}
\end{theorem}

\begin{proof}
The derivative of \(\mathcal F\) in \(\psi\) at the base point is the
positive isomorphism \(L_0:C^{2,\alpha}\to C^{0,\alpha}\).  The analytic
implicit-function theorem gives the branch.  Parameter differentiation
gives
\begin{equation}
 L_0u=h_+\psi_0,
 \qquad L_0v=h_-\psi_0^{-7},
 \qquad L_0r=-\psi_0^5.
 \label{eq:v6v84-linear-responses}
\end{equation}
The constant and first-cosine eigenvalues of \(L_0\) are \(\mu_0\) and
\(\mu_0+8\), which yields the displayed formulas.
\end{proof}

\begin{assessment}[Opposite nonzero-mode signs]
\label{ass:v6v84-opposite-signs}
Both matter sources increase the conformal zero mode at fixed CMC gap.  At
\(k_0\), however, the gradient density is proportional to \(1+\cos y\),
whereas the squared canonical momentum is proportional to
\(1-\cos y\).  Their first conformal responses therefore have opposite
signs.  This is an exact common-ancestry effect, not an inequality.
\end{assessment}

\section{York-volume derivatives}
\label{sec:v6v84-volume}

Let
\begin{equation}
 \mathcal V(\epsilon,\beta,\mathfrak a)
 =\int_{\mathbb T}\psi(\epsilon,\beta,\mathfrak a)^6d\mu.
 \label{eq:v6v84-volume}
\end{equation}

\begin{theorem}[First response of the York Hamiltonian]
\label{thm:v6v84-York-response}
At the base point,
\begin{align}
 \partial_\epsilon\mathcal V
 &={6\psi_0^6\over\mu_0},
 \label{eq:v6v84-Vepsilon}\\
 \partial_\beta\mathcal V
 &={6\psi_0^{-2}\over\mu_0},
 \label{eq:v6v84-Vbeta}\\
 \partial_{\mathfrak a}\mathcal V
 &=-{6\psi_0^{10}\over\mu_0}.
 \label{eq:v6v84-Va}
\end{align}
In the York orientation of V82, at fixed gap
\begin{equation}
 H_Y=\psi_0^6+{6\psi_0^6\over\mu_0}\epsilon
 +{6\psi_0^{-2}\over\mu_0}\beta
 +O\bigl((|\epsilon|+|\beta|)^2\bigr).
 \label{eq:v6v84-HY}
\end{equation}
Thus both direct matter sources increase the local classical York
Hamiltonian to first order.
\end{theorem}

\begin{proof}
Differentiate \eqref{eq:v6v84-volume}, insert
\eqref{eq:v6v84-u}--\eqref{eq:v6v84-r}, and use the zero mean of
\(\cos y\).  The York Hamiltonian is physical volume by
Definition~\ref{def:v6v82-York-Hamiltonian}.
\end{proof}

\section{Fixed-volume response and cancellation}
\label{sec:v6v84-fixed-volume}

\begin{theorem}[Two-source fixed-volume York gap]
\label{thm:v6v84-fixed-volume}
There is a unique analytic gap
\(\mathfrak a=\mathfrak a(\epsilon,\beta)\) near the base point for which
\begin{equation}
 \mathcal V(\epsilon,\beta,\mathfrak a(\epsilon,\beta))
 =\psi_0^6.
 \label{eq:v6v84-fixed-volume}
\end{equation}
It satisfies
\begin{equation}
 \boxed{\quad
 \mathfrak a(\epsilon,\beta)
 =\mathfrak a_0+\psi_0^{-4}\epsilon+\psi_0^{-12}\beta
 +O\bigl((|\epsilon|+|\beta|)^2\bigr).
 \quad}
 \label{eq:v6v84-gap}
\end{equation}
The corresponding conformal factor is
\begin{equation}
 \boxed{\quad
 \psi_{\mathrm{vol}}(y)
 =\psi_0+{\psi_0\epsilon-\psi_0^{-7}\beta\over\mu_0+8}
 \cos y
 +O_{C^{2,\alpha}}\bigl((|\epsilon|+|\beta|)^2\bigr).
 \quad}
 \label{eq:v6v84-fixed-psi}
\end{equation}
\end{theorem}

\begin{proof}
The derivative \eqref{eq:v6v84-Va} is nonzero, so the analytic
implicit-function theorem solves \eqref{eq:v6v84-fixed-volume} for the gap.
Its two derivatives are
\begin{align}
 \partial_\epsilon\mathfrak a
 &=-{\partial_\epsilon\mathcal V\over
          \partial_{\mathfrak a}\mathcal V}=\psi_0^{-4},
 \label{eq:v6v84-aepsilon}\\
 \partial_\beta\mathfrak a
 &=-{\partial_\beta\mathcal V\over
          \partial_{\mathfrak a}\mathcal V}=\psi_0^{-12}.
 \label{eq:v6v84-abeta}
\end{align}
This proves \eqref{eq:v6v84-gap}.  Add to \(u\) and \(v\) the gap
responses \(r\partial_\epsilon\mathfrak a\) and
\(r\partial_\beta\mathfrak a\).  In each case the constant mode cancels,
leaving the two cosine terms in \eqref{eq:v6v84-fixed-psi}.
\end{proof}

\begin{corollary}[First-order beat-cancellation line]
\label{cor:v6v84-cancellation}
The nonzero \(k_0\) conformal response at fixed volume vanishes to first
order exactly on
\begin{equation}
 \boxed{\qquad \beta=\psi_0^8\epsilon.\qquad}
 \label{eq:v6v84-cancellation-line}
\end{equation}
The York-gap response does not vanish there; rather,
\begin{equation}
 \mathfrak a(\epsilon,\psi_0^8\epsilon)
 =\mathfrak a_0+2\psi_0^{-4}\epsilon+O(\epsilon^2).
 \label{eq:v6v84-gap-on-line}
\end{equation}
\end{corollary}

\begin{proof}
The cosine coefficient in \eqref{eq:v6v84-fixed-psi} is zero exactly when
\(\psi_0\epsilon=\psi_0^{-7}\beta\).  Substitute that relation in
\eqref{eq:v6v84-gap}.
\end{proof}

Since
\(\mathfrak a=2\tau^2/3-2\Lambda-16\pi GV(0)\), the same theorem gives
\begin{equation}
 \tau(\epsilon,\beta)^2
 =\tau_0^2+{3\over2}
 \left(\psi_0^{-4}\epsilon+\psi_0^{-12}\beta\right)
 +O\bigl((|\epsilon|+|\beta|)^2\bigr).
 \label{eq:v6v84-tau}
\end{equation}

\section{The vector constraint first appears at mixed order}
\label{sec:v6v84-mixed}

\begin{theorem}[Exact mixed-response equation]
\label{thm:v6v84-mixed}
At fixed gap, let
\begin{equation}
 z:=\partial_\epsilon\partial_\beta\psi(0,0,\mathfrak a_0).
 \label{eq:v6v84-z}
\end{equation}
Then
\begin{align}
 L_0z={}&36\mathfrak a_0\psi_0^3uv
 -7h_-\psi_0^{-8}u+h_+v
 +2\cos^2y\,\psi_0^{-7}.
 \label{eq:v6v84-mixed-equation}
\end{align}
The last term is the direct longitudinal gravitational-momentum response.
It contains the harmonics \(0\) and \(2k_0\), because
\begin{equation}
 2\cos^2y=1+\cos2y.
 \label{eq:v6v84-second-harmonic}
\end{equation}
It is absent from every first derivative and survives on the cancellation
line \eqref{eq:v6v84-cancellation-line} at order \(\epsilon^2\).
\end{theorem}

\begin{proof}
Differentiate \(\mathcal F=0\) once in each source variable.  The second
derivative of the base nonlinearity
\(-s_*\psi^{-7}+\mathfrak a_0\psi^5\) at \(\psi_0\) is
\begin{equation}
 -56s_*\psi_0^{-9}+20\mathfrak a_0\psi_0^3
 =-36\mathfrak a_0\psi_0^3,
 \label{eq:v6v84-base-second}
\end{equation}
using \(s_*=\mathfrak a_0\psi_0^{12}\).  The mixed derivatives of
\(-\beta h_-\psi^{-7}\), \(-\epsilon h_+\psi\), and
\(-2\epsilon\beta\cos^2y\psi^{-7}\) are respectively
\(7h_-\psi_0^{-8}u\), \(-h_+v\), and
\(-2\cos^2y\psi_0^{-7}\) on the left side.  Move them and the base
second-derivative term to the right to obtain
\eqref{eq:v6v84-mixed-equation}.  The harmonic statement is elementary,
and \(\epsilon\beta=\psi_0^8\epsilon^2\) on the cancellation line.
\end{proof}

\begin{assessment}[Cancellation is order-specific]
\label{ass:v6v84-order-specific}
The line \eqref{eq:v6v84-cancellation-line} cancels only the first
nonzero-mode conformal response after the volume zero mode has been removed.
It does not cancel the York-time shift, the direct matter energies, or the
mixed longitudinal gravitational momentum.  At the next order the latter
forces \(0\) and \(2k_0\) components, and the remaining terms in
\eqref{eq:v6v84-mixed-equation} also contribute.  No all-order cancellation
has been proved.
\end{assessment}

\begin{definition}[Local York response card, LYRC]
\label{def:v6v84-LYRC}
LYRC requires:
\begin{enumerate}[label=\textup{(YR\arabic*)},leftmargin=5.2em]
\item analytic dependence of the constrained conformal and vector solutions
      on every independent matter source amplitude;
\item separate fixed-York and fixed-volume response matrices;
\item exact treatment of zero modes through the integrated constraint;
\item ancestry-resolved higher responses, beginning with
      \eqref{eq:v6v84-mixed-equation};
\item a self-adjoint lower-bounded quantization of the volume generator on a
      declared reduced measure; and
\item transition control beyond the local CMC chart.
\end{enumerate}
\end{definition}

\begin{theorem}[LYRC progress]
\label{thm:v6v84-LYRC}
The matched scalar sector closes LYRC\textup{(YR1)}--\textup{(YR4)} through
mixed order locally at the homogeneous base point.  Rows
\textup{(YR5)}--\textup{(YR6)} remain open.
\end{theorem}

\begin{proof}
Use Theorems~\ref{thm:v6v84-analytic},
\ref{thm:v6v84-York-response},
\ref{thm:v6v84-fixed-volume}, and
\ref{thm:v6v84-mixed}.
\end{proof}

\begin{programme}[V85: compulsory audit and reduced-measure gate]
\label{prog:v6v84-V85}
\begin{enumerate}[leftmargin=3em]
\item Audit the newest active YM and NS notes by hash-defined snapshots and
      import no numerical constants without matching hypotheses.
\item Compare their newest complete-source and scale-assembly mechanisms
      with the first-order cancellation and mixed forcing above.
\item Specify the finite-cutoff reduced configuration space, symplectic
      density, and volume-generator quadratic form required for
      LYRC\textup{(YR5)}.
\item Prove closability and lower boundedness if the density is controlled;
      otherwise isolate the exact Jacobian or chart-transition obstruction.
\item Set the next direction only after this measure gate is explicit.
\end{enumerate}
\end{programme}

\begin{firewall}[Claim boundary after V84]
\label{fw:v6v84-boundary}
V84 is a local classical response calculation in one matched scalar CMC
sector.  It does not prove an all-order cancellation, construct the reduced
quantum measure, establish self-adjoint York evolution, cover non-CMC data,
assemble the remaining Standard Model fields, remove a regulator, or prove
the full SM--Einstein continuum.
\end{firewall}

\begin{theorem}[Sixth-series V84 result]
\label{thm:v6v84-status}
The momentum-coupled homogenized constraint has an analytic two-source
branch.  At fixed York time both scalar-gradient and kinetic-momentum
sources increase the volume Hamiltonian.  At fixed volume their zero modes
shift York time according to \eqref{eq:v6v84-tau}, while their \(k_0\)
conformal responses cancel to first order exactly on
\(\beta=\psi_0^8\epsilon\).  The longitudinal vector response begins at
mixed order, obeys \eqref{eq:v6v84-mixed-equation}, and creates a
\(2k_0\) harmonic.  The remaining bottleneck is the reduced quantum measure
and self-adjoint volume generator.
\end{theorem}

\begin{proof}
Use Theorems~\ref{thm:v6v84-analytic},
\ref{thm:v6v84-York-response},
\ref{thm:v6v84-fixed-volume}, and
\ref{thm:v6v84-mixed}, subject to
Firewall~\ref{fw:v6v84-boundary}.
\end{proof}

The compulsory YM/NS audit is due in V85.
\clearpage
\chapter{Sixth-series V85: compulsory companion audit and the regular
finite-cutoff reduced-measure theorem}
\label{ch:v6v85-measure}

\section{Purpose and hash-defined companion audit}
\label{sec:v6v85}

V84 isolated a first-order cancellation line but also proved that the
longitudinal vector response survives at mixed order.  The next task is not
to extrapolate that cancellation.  It is to identify the measure on which a
York-volume operator could be defined.  This version first performs the
compulsory audit and then proves a local finite-dimensional reduction theorem
for a regular CMC chart.  The result supplies a positive coarea density and a
positive self-adjoint volume multiplication operator on compact regular
source boxes.  It is deliberately limited to a local sector-label
representation; the full canonical York evolution remains open.

The companion folder is concurrently modified by the other programmes.  The
following metadata refer to single in-memory byte snapshots, and their
SHA--256 digests define the evidence used here.

\begin{definition}[V85 companion snapshots]
\label{def:v6v85-snapshots}
The Yang--Mills snapshot is
\begin{align}
 &\texttt{Renewal\_Yang\_Mills\_Mass\_Gap\_Research\_Notes\_V60.tex},
 \notag\\
 &931\,699\ \text{bytes},\qquad29\,241\ \text{logical source lines},
 \notag\\
 &\texttt{c4d0867f8234c879921dccdabc19f62dbbc5a4a324884f889bbcf8acb8104efa}.
 \label{eq:v6v85-YM-snapshot}
\end{align}
The active Navier--Stokes programme has opened its second cycle.  Its
snapshot is
\begin{align}
 &\texttt{Renewal\_NS\_Stability\_Research\_Notes\_V1.tex},
 \notag\\
 &27\,351\ \text{bytes},\qquad567\ \text{logical source lines},
 \notag\\
 &\texttt{27a17aaf6b17e18ccfa7c4dccaaf54679ac85f7eebacb81aebd37129e571a8df}.
 \label{eq:v6v85-NS-snapshot}
\end{align}
\end{definition}

\subsection{Yang--Mills endpoint}

YM V60 begins from an exact positive L\'evy split of the Wilson clock.  It
proves all positive small-jump moments at the required scale, an all-arity
same-jump hyperedge bound in weighted Pr\"ufer form, and an exact
endpoint-factorizing interpolation whose BKAR forest equals the connected
small-clock source.  It then proves a precise obstruction: applying
Fr\'echet--Fa\`a di Bruno differentiation to the outer exponential and
taking norms term by term produces a Bell-number majorant that cannot be
absorbed into one exponential arity base.  Its OCRC gate requires all jump
roots sharing a label to be resummed before norms.  This is a no-go for that
proof order, not a no-go for the Yang--Mills source or a mass-gap theorem.

\subsection{Navier--Stokes endpoint}

NS second-cycle V1 records the first-cycle results and rewrites the
Gaussian-weighted identity in logarithmic scale.  Its exact scalar ledger is
\(\varphi'=2\varphi-\xi\), and the average of \(\xi\) on every log interval
equals twice the average of \(\varphi\), up to the endpoint difference
divided by the interval length.  At homogeneous points this gives a positive
long-window average of the excess inward flux.  The pointwise inequality
needed for a monotone integrating factor remains open.  Thus an averaged
assembled identity does not imply the pointwise sign required by the final
regularity gate.

\begin{theorem}[V85 audit decision]
\label{thm:v6v85-audit-decision}
No numerical YM jump estimate or NS flux constant transfers to the
SM--Einstein constraint measure.  Two structural conclusions do transfer:
\begin{enumerate}[label=\textup{(\roman*)},leftmargin=3.7em]
\item the V84 first-order cancellation cannot be iterated by differentiating
      each nonlinear source contribution separately and taking absolute
      values; and
\item an averaged or integrated positive statement cannot replace the
      pointwise Jacobian nondegeneracy and lower bound required for a
      reduced Hilbert-space measure.
\end{enumerate}
\end{theorem}

\begin{proof}
The companion estimates concern compact-group heat jumps and viscous flux,
not the CMC constraint derivative, so their constants have incompatible
hypotheses.  YM V60 explicitly exhibits superexponential Bell overlap after
premature termwise differentiation; V84 has the same potential proliferation
through mixed ancestry derivatives.  NS V1 distinguishes its exact
log-average identity from the unproved pointwise sign; a coarea density
likewise needs pointwise rank.  These facts prove only the two stated proof
rules.
\end{proof}

\section{A finite-dimensional constraint disintegration theorem}
\label{sec:v6v85-disintegration}

At a fixed spectral cutoff all fields have finitely many real Fourier
coordinates.  Separate them as \(z\in U\subset\mathbb R^d\), the proposed
reduced coordinates, and \(y\in V\subset\mathbb R^m\), the conformal,
longitudinal, York-zero-mode, and gauge coordinates to be solved.  Let
\begin{equation}
 F:U\times V\longrightarrow\mathbb R^m
 \label{eq:v6v85-F}
\end{equation}
collect the constraint and gauge equations after Killing zero modes have
been removed or fixed.

\begin{theorem}[Local delta/coarea reduction]
\label{thm:v6v85-coarea}
Suppose \(F\in C^1\),
\begin{equation}
 F(z_0,y_0)=0,
 \qquad \det D_yF(z_0,y_0)\ne0.
 \label{eq:v6v85-regular-point}
\end{equation}
After shrinking neighbourhoods, there is a unique \(C^1\) branch
\(y=Y(z)\) with \(F(z,Y(z))=0\), and for every compactly supported
continuous \(A\),
\begin{equation}
 \int A(z,y)\,\delta_0(F(z,y))\,dy\,dz
 =\int A(z,Y(z)){dz\over|\det D_yF(z,Y(z))|}.
 \label{eq:v6v85-delta-reduction}
\end{equation}
If a positive smooth canonical density \(J_{\rm can}(z,y)\) and a regular
Faddeev--Popov matrix \(M_{\rm FP}(z,y)\) are included, the local reduced
density is
\begin{equation}
 \rho_{\rm red}(z)
 ={J_{\rm can}(z,Y(z))\,
   |\det M_{\rm FP}(z,Y(z))|
  \over|\det D_yF(z,Y(z))|}.
 \label{eq:v6v85-reduced-density}
\end{equation}
It is positive and continuous wherever the three displayed factors are
finite and nonzero.
\end{theorem}

\begin{proof}
The implicit-function theorem gives \(Y\).  For each fixed \(z\), the map
\(y\mapsto u=F(z,y)\) is a local diffeomorphism.  Change variables from
\(y\) to \(u\) and use
\(\int B(u)\delta_0(u)du=B(0)\).  The Jacobian is
\(|\det D_yF|^{-1}\), giving
\eqref{eq:v6v85-delta-reduction}.  Multiplication by the canonical density
and the absolute Faddeev--Popov determinant gives
\eqref{eq:v6v85-reduced-density}.  Positivity and continuity follow from
the hypotheses and continuity of determinants.
\end{proof}

\begin{corollary}[Uniform density on a compact regular box]
\label{cor:v6v85-density-bounds}
Let \(K\Subset U\) be compact and suppose the branch stays in one regular
chart over \(K\).  Then there are \(0<c_K<C_K<\infty\) such that
\begin{equation}
 c_K\le\rho_{\rm red}(z)\le C_K
 \qquad(z\in K).
 \label{eq:v6v85-density-bounds}
\end{equation}
\end{corollary}

\begin{proof}
The positive continuous numerator and denominator in
\eqref{eq:v6v85-reduced-density} attain positive minima and finite maxima
on the compact graph of \(Y|_K\).
\end{proof}

\section{Regularity of the matched CMC block}
\label{sec:v6v85-CMC-block}

At the homogeneous base point, solve the regulated equations in the order
\begin{equation}
 \text{vector constraint}\ \longrightarrow\
 \text{Lichnerowicz equation}\ \longrightarrow\
 \text{fixed-volume York equation}.
 \label{eq:v6v85-solve-order}
\end{equation}

\begin{theorem}[Invertible matched constraint Jacobian]
\label{thm:v6v85-Jacobian}
At every finite spectral cutoff containing the matched ancestors, the
derivative of the three-block map in
\eqref{eq:v6v85-solve-order}, restricted to mean-zero vector fields and a
fixed local York sign branch, is block triangular with invertible diagonal
blocks
\begin{equation}
 \bar\Delta_L,
 \qquad L_0=-8{d^2\over dy^2}+\mu_0,
 \qquad
 \partial_{\mathfrak a}\mathcal V
 =-{6\psi_0^{10}\over\mu_0}.
 \label{eq:v6v85-diagonal-blocks}
\end{equation}
It is therefore invertible.  If the chosen finite-mode spatial gauge has an
invertible Faddeev--Popov derivative at the base point after translational
zero modes are fixed, the complete constraint-plus-gauge map remains
regular on a sufficiently small matched-source neighbourhood.
\end{theorem}

\begin{proof}
Lemma~\ref{lem:v6v83-vector-constraint} proves that
\(\bar\Delta_L\) is invertible after its constant kernel is removed.
Equation \eqref{eq:v6v84-base} gives \(\mu_0>0\), so every Fourier
eigenvalue of \(L_0\) is positive.  Equation \eqref{eq:v6v84-Va} gives the
nonzero scalar derivative in the third block.  The selected solve order
makes dependencies on earlier solved variables lower triangular; a block
triangular matrix with invertible diagonal blocks is invertible.  Appending
an invertible gauge block preserves invertibility.  Determinants vary
continuously, so regularity persists on a smaller neighbourhood.
\end{proof}

\begin{assessment}[What the determinant contains]
\label{ass:v6v85-determinant}
The reduced density is not merely the conformal determinant.  It contains
the mean-zero vector determinant, the conformal determinant, the scalar
fixed-volume/York derivative, the spatial gauge determinant, and the
canonical density in one product.  Removing the zero modes before forming
that product is essential.  V82's nonzero derivative
\(\partial_{\mathfrak a}\mathcal V\) is exactly the missing zero-mode
Jacobian.
\end{assessment}

\section{A positive local volume operator}
\label{sec:v6v85-volume-operator}

Let \(K\) be a compact regular reduced source box and define
\begin{equation}
 \mathcal H_K=L^2(K,\rho_{\rm red}(z)dz).
 \label{eq:v6v85-HK}
\end{equation}
Let \(\mathcal V(z)\) be the physical volume of the constrained branch.

\begin{theorem}[Sector-label York-volume operator]
\label{thm:v6v85-volume-operator}
Assume the matched conformal bounds hold throughout \(K\):
\begin{equation}
 0<c\le\psi(z,x)\le C<\infty.
 \label{eq:v6v85-uniform-psi}
\end{equation}
Then
\begin{equation}
 c^6\le\mathcal V(z)=\int\psi(z,x)^6d\mu_{\bar g}\le C^6.
 \label{eq:v6v85-volume-bounds}
\end{equation}
The multiplication operator
\begin{equation}
 (\widehat H_Yf)(z)=\mathcal V(z)f(z)
 \label{eq:v6v85-HY-operator}
\end{equation}
is bounded, self-adjoint, and satisfies
\begin{equation}
 \widehat H_Y\ge c^6I
 \label{eq:v6v85-positive-HY}
\end{equation}
on \(\mathcal H_K\).
\end{theorem}

\begin{proof}
The volume bounds follow by integration on the normalized reference torus.
A bounded real measurable multiplication function defines a bounded
self-adjoint operator on the full \(L^2\) space.  The pointwise lower bound
gives
\(\langle f,\widehat H_Yf\rangle\ge c^6\|f\|^2\).
\end{proof}

\begin{proposition}[Local chart covariance]
\label{prop:v6v85-chart-covariance}
Suppose \(z'=\chi(z)\) is a diffeomorphic change between two regular
reduced charts obtained from the same constrained density.  Then
\begin{equation}
 \rho'_{\rm red}(z')
 =\rho_{\rm red}(\chi^{-1}(z'))
 |\det D\chi^{-1}(z')|,
 \label{eq:v6v85-density-transform}
\end{equation}
and pullback is a unitary map between the two local Hilbert spaces that
intertwines their volume multiplication operators.
\end{proposition}

\begin{proof}
Equation \eqref{eq:v6v85-density-transform} is the ordinary change-of-
variables law for the same reduced measure.  For
\((Uf)(z)=f(\chi(z))\), that law gives equality of the two \(L^2\) norms.
Physical volume is a scalar on the constraint surface, so its two coordinate
representatives compose with \(\chi\), proving the intertwining identity.
\end{proof}

\begin{firewall}[A sector-label observable is not full York dynamics]
\label{fw:v6v85-sector-label}
The operator \eqref{eq:v6v85-HY-operator} quantizes volume as a function of
local reduced sector labels.  The full York Hamiltonian generally depends
on canonical coordinates and momenta, and its quantization may contain
derivatives, ordering terms, and nonlocal constraint dependence.  The
multiplication theorem does not construct that operator, prove unitary York
evolution, or identify a Euclidean transfer generator.  It proves that the
regular local coarea measure itself creates no sign obstruction to the
volume observable.
\end{firewall}

\section{Regular reduced-measure card}
\label{sec:v6v85-RRMC}

\begin{definition}[Regular reduced-measure card, RRMC]
\label{def:v6v85-RRMC}
RRMC requires at each cutoff:
\begin{enumerate}[label=\textup{(RM\arabic*)},leftmargin=5.2em]
\item removal or separate integration of all Killing, gauge, York, and
      global charge zero modes;
\item a full-rank constraint-plus-gauge map on each admitted chart;
\item the positive coarea/Faddeev--Popov density
      \eqref{eq:v6v85-reduced-density};
\item compatible transition densities and operators on chart overlaps;
\item control as a chart approaches a determinant zero or CMC boundary;
\item a canonical self-adjoint York generator, including momentum and
      ordering dependence; and
\item regulator-uniform estimates sufficient to pass to a continuum
      Hilbert space and dynamics.
\end{enumerate}
\end{definition}

\begin{theorem}[RRMC progress and exact remaining gate]
\label{thm:v6v85-RRMC}
On every compact box inside the regular matched finite-cutoff CMC chart,
rows RRMC\textup{(RM1)}--\textup{(RM4)} hold locally, and the positive
volume multiplication observable is self-adjoint.  Rows
RRMC\textup{(RM5)}--\textup{(RM7)} remain open.  In particular, the first
unresolved measure obstruction is quantitative control of
\begin{equation}
 |\det D_yF|^{-1}|\det M_{\rm FP}|
 \label{eq:v6v85-boundary-factor}
\end{equation}
as the regular chart approaches a constraint fold, gauge horizon, or CMC
sector boundary.
\end{theorem}

\begin{proof}
Theorem~\ref{thm:v6v85-Jacobian} removes the declared zero modes and gives
local rank.  Theorem~\ref{thm:v6v85-coarea} and
Proposition~\ref{prop:v6v85-chart-covariance} give the positive density and
overlap covariance.  Theorem~\ref{thm:v6v85-volume-operator} gives the local
self-adjoint observable.  None supplies a lower determinant bound at the
boundary, a derivative-valued canonical Hamiltonian, or cutoff-uniform
control, which are exactly the remaining rows.
\end{proof}

\begin{programme}[V86: determinant approach to the first constraint fold]
\label{prog:v6v85-V86}
\begin{enumerate}[leftmargin=3em]
\item Follow the analytic two-source branch until the lowest eigenvalue of
      the conformal derivative reaches zero or the positive conformal bounds
      fail.
\item Derive a lower bound for that eigenvalue from the pointwise derivative
      of the full momentum-coupled nonlinearity.
\item Relate the finite-cutoff determinant to the product of vector,
      conformal, and York blocks; isolate which block can vanish first.
\item Test the V84 cancellation line at second order using the exact mixed
      equation rather than separate absolute estimates.
\item If determinant control becomes circular, return upstream to the
      maximum-principle uniqueness condition and restrict the admitted
      source domain explicitly.
\end{enumerate}
\end{programme}

\begin{firewall}[Claim boundary after V85]
\label{fw:v6v85-boundary}
V85 proves a local finite-dimensional coarea measure and a positive
sector-label volume observable on compact regular boxes.  It does not prove
global gauge fixing, a CMC atlas, determinant control at chart boundaries,
the full canonical York Hamiltonian, unitary or reflection-positive
dynamics, complete Standard Model source assembly, regulator removal, or a
full SM--Einstein continuum.
\end{firewall}

\begin{theorem}[Sixth-series V85 result]
\label{thm:v6v85-status}
The companion audit confirms that nonlinear ancestry must be resummed before
termwise norms and that averaged positivity cannot replace pointwise rank.
At every finite cutoff, the matched vector, conformal, and fixed-volume York
blocks have an invertible triangular derivative at the homogeneous base
point.  The resulting local coarea/Faddeev--Popov density is positive and
bounded above and below on compact regular source boxes, is covariant under
regular chart changes, and supports a positive self-adjoint volume
multiplication observable.  The next gate is determinant control at the
first constraint fold or gauge horizon, followed by construction of the
full canonical York generator and regulator-uniform bounds.
\end{theorem}

\begin{proof}
Use Theorem~\ref{thm:v6v85-audit-decision},
Theorems~\ref{thm:v6v85-coarea},
\ref{thm:v6v85-Jacobian},
\ref{thm:v6v85-volume-operator}, and
\ref{thm:v6v85-RRMC}, subject to
Firewalls~\ref{fw:v6v85-sector-label} and \ref{fw:v6v85-boundary}.
\end{proof}

The next compulsory YM/NS audit is V90.
\clearpage
\chapter{Sixth-series V86: conformal spectral gap, raw-determinant
divergence, and the modified determinant}
\label{ch:v6v86-determinant}

\section{V86 entry point}
\label{sec:v6v86}

V85 reduced the local finite-cutoff measure to an inverse constraint
Jacobian.  Two distinct failures can occur as the cutoff or source domain is
enlarged.  A genuine constraint fold occurs when the lowest eigenvalue of
the conformal derivative reaches zero.  Even when that eigenvalue stays
strictly positive, the ordinary determinant has an ultraviolet product over
infinitely many conformal modes.  This version proves a pointwise spectral-
gap criterion excluding the first failure, computes the field-dependent
linear divergence of the second, and replaces the raw relative determinant
by the Hilbert--Schmidt modified determinant \(\det_2\).

\section{The conformal derivative and a no-fold domain}
\label{sec:v6v86-gap}

For the momentum-coupled homogenized equation, write
\begin{equation}
 \mathcal F(\psi)
 =-8\bar\Delta\psi-S(x)\psi^{-7}
 -\epsilon h_+(x)\psi+\mathfrak a\psi^5,
 \qquad h_+=1+\cos x_2.
 \label{eq:v6v86-F}
\end{equation}
Its derivative at a positive solution is
\begin{equation}
 \mathcal L_\psi
 =-8\bar\Delta+7S(x)\psi^{-8}
 -\epsilon h_+(x)+5\mathfrak a\psi^4.
 \label{eq:v6v86-linearization}
\end{equation}

\begin{theorem}[Pointwise spectral-gap criterion]
\label{thm:v6v86-gap}
Assume
\begin{equation}
 S(x)\ge S_->0,
 \qquad c\le\psi(x)\le C,
 \qquad \mathfrak a\ge\mathfrak a_->0,
 \label{eq:v6v86-coefficient-bounds}
\end{equation}
and define
\begin{equation}
 \delta:=7S_-C^{-8}+5\mathfrak a_-c^4-2|\epsilon|.
 \label{eq:v6v86-delta}
\end{equation}
If \(\delta>0\), then
\begin{equation}
 \langle u,\mathcal L_\psi u\rangle
 \ge8\|\bar\nabla u\|_2^2+\delta\|u\|_2^2
 \label{eq:v6v86-coercive}
\end{equation}
for every real \(u\in H^1(\mathbb T^3)\).  Hence
\(\mathcal L_\psi:H^2\to L^2\) is invertible,
\begin{equation}
 \|\mathcal L_\psi^{-1}\|_{L^2\to L^2}\le\delta^{-1},
 \label{eq:v6v86-inverse-bound}
\end{equation}
and no conformal constraint fold occurs inside a parameter region on which
the same positive \(\delta\) holds.
\end{theorem}

\begin{proof}
Since \(0\le h_+\le2\), the multiplication part of
\eqref{eq:v6v86-linearization} is bounded below by \(\delta\).  Integration
by parts gives \eqref{eq:v6v86-coercive}.  Lax--Milgram gives a weak inverse,
elliptic regularity maps its range into \(H^2\), and the quadratic-form
bound gives \eqref{eq:v6v86-inverse-bound}.  A fold requires a nontrivial
kernel of the constraint derivative, which the same estimate excludes.
\end{proof}

\begin{corollary}[A quantitative regular source box]
\label{cor:v6v86-regular-box}
Every compact matched-source box on which the V83 barriers give common
constants \(S_-,c,C,\mathfrak a_-\) and
\begin{equation}
 2|\epsilon|<7S_-C^{-8}+5\mathfrak a_-c^4
 \label{eq:v6v86-box-condition}
\end{equation}
lies strictly before the first conformal fold.  Its finite-cutoff conformal
determinants are positive.
\end{corollary}

\begin{proof}
Theorem~\ref{thm:v6v86-gap} gives a uniform positive lower eigenvalue.  A
finite-dimensional compression of a positive operator has positive
determinant.
\end{proof}

\section{Relative conformal determinants}
\label{sec:v6v86-relative}

At the homogeneous base point use
\begin{equation}
 \mathcal L_0=-8\bar\Delta+\mu_0,
 \qquad \mu_0>0,
 \label{eq:v6v86-L0}
\end{equation}
and write a nearby positive conformal derivative as
\begin{equation}
 \mathcal L=\mathcal L_0+U(x),
 \label{eq:v6v86-L}
\end{equation}
where \(U\) is real and smooth.  Let \(P_\Lambda\) project onto torus
Fourier modes \(|k|\le\Lambda\), and put
\begin{equation}
 K=\mathcal L_0^{-1/2}U\mathcal L_0^{-1/2},
 \qquad K_\Lambda=P_\Lambda KP_\Lambda.
 \label{eq:v6v86-K}
\end{equation}

\begin{lemma}[The relative perturbation is Hilbert--Schmidt]
\label{lem:v6v86-HS}
On \(L^2(\mathbb T^3)\), \(K\) is self-adjoint and Hilbert--Schmidt.  For a
nonzero constant \(U\), it is not trace class.  Moreover
\begin{equation}
 \|K_\Lambda-K\|_{\mathfrak S_2}\longrightarrow0.
 \label{eq:v6v86-HS-convergence}
\end{equation}
\end{lemma}

\begin{proof}
In the Fourier basis, with
\(\lambda_k=8|k|^2+\mu_0\),
\begin{equation}
 K_{k\ell}={\widehat U(k-\ell)\over\sqrt{\lambda_k\lambda_\ell}}.
 \label{eq:v6v86-K-matrix}
\end{equation}
Thus
\begin{equation}
 \|K\|_{\mathfrak S_2}^2
 =\sum_q|\widehat U(q)|^2
 \sum_\ell{1\over\lambda_{\ell+q}\lambda_\ell}<\infty.
 \label{eq:v6v86-HS-sum}
\end{equation}
The inner sum converges because its summand is \(O(|\ell|^{-4})\) in
dimension three; smoothness of \(U\) controls the outer sum.  If \(U=u\ne0\)
is constant, the singular values are
\(|u|/\lambda_k\), whose sum diverges like
\(\sum_{|k|\le\Lambda}|k|^{-2}\).  Hence trace class fails in general.
Finite-rank Fourier projections converge to a Hilbert--Schmidt operator in
Hilbert--Schmidt norm, proving \eqref{eq:v6v86-HS-convergence}.
\end{proof}

\begin{theorem}[Exact finite-cutoff determinant split]
\label{thm:v6v86-det-split}
Let \(\mathcal L_{0,\Lambda}\) and \(\mathcal L_\Lambda\) be the Fourier
compressions.  Then
\begin{equation}
 {\det\mathcal L_\Lambda\over\det\mathcal L_{0,\Lambda}}
 =\det(I+K_\Lambda)
 =\exp\{\operatorname{Tr}K_\Lambda\}
   \det{}_2(I+K_\Lambda),
 \label{eq:v6v86-det-split}
\end{equation}
where
\begin{equation}
 \operatorname{Tr}K_\Lambda
 =\overline U\,\mathscr S_\Lambda,
 \qquad
 \overline U:=\int_{\mathbb T^3}U,d\mu,
 \qquad
 \mathscr S_\Lambda:=\sum_{|k|\le\Lambda}{1\over8|k|^2+\mu_0}.
 \label{eq:v6v86-trace}
\end{equation}
There are constants \(0<c<C<\infty\) such that
\begin{equation}
 c\Lambda\le\mathscr S_\Lambda\le C\Lambda
 \label{eq:v6v86-linear-divergence}
\end{equation}
for all large \(\Lambda\).  In contrast,
\begin{equation}
 \det{}_2(I+K_\Lambda)\longrightarrow\det{}_2(I+K)
 \label{eq:v6v86-det2-limit}
\end{equation}
whenever \(I+K\) is invertible; if \(\mathcal L>0\), the limit is strictly
positive.
\end{theorem}

\begin{proof}
Because \(P_\Lambda\) commutes with \(\mathcal L_0\), conjugation by
\(\mathcal L_{0,\Lambda}^{-1/2}\) gives \(I+K_\Lambda\), proving the first
equality.  The finite-dimensional definition
\(\det_2(I+A)=\det(I+A)e^{-\operatorname{Tr}A}\) gives the second.  The
diagonal matrix entry in \eqref{eq:v6v86-K-matrix} is
\(\widehat U(0)/\lambda_k=\overline U/\lambda_k\), proving
\eqref{eq:v6v86-trace}.  Comparing lattice shells with
\(r^2\) points and summand of order \(r^{-2}\) proves
\eqref{eq:v6v86-linear-divergence}.  Modified Fredholm determinants are
continuous in Hilbert--Schmidt norm, so
Lemma~\ref{lem:v6v86-HS} proves \eqref{eq:v6v86-det2-limit}.  When
\(I+K=\mathcal L_0^{-1/2}\mathcal L\mathcal L_0^{-1/2}>0\), every factor in
the defining modified eigenvalue product is positive.
\end{proof}

\begin{corollary}[The raw coarea determinant has a local ultraviolet term]
\label{cor:v6v86-raw-density}
After the field-independent base determinant is divided out, the conformal
factor in the inverse coarea density is
\begin{equation}
 {\det\mathcal L_{0,\Lambda}\over\det\mathcal L_\Lambda}
 =\exp\{-\overline U\mathscr S_\Lambda\}
  \det{}_2(I+K_\Lambda)^{-1}.
 \label{eq:v6v86-inverse-det}
\end{equation}
Thus positivity and a uniform lowest eigenvalue do not make the raw
relative density cutoff-independent.  When \(\overline U\ne0\), it contains
a field-dependent exponential of order \(\Lambda\).
\end{corollary}

\begin{proof}
Invert \eqref{eq:v6v86-det-split} and use
\eqref{eq:v6v86-linear-divergence}.
\end{proof}

\section{The divergence on the fixed-volume beat branch}
\label{sec:v6v86-beat-divergence}

Let \(U_{\mathrm{vol}}(\epsilon,\beta;y)\) be the potential in
\eqref{eq:v6v86-L} evaluated on the fixed-volume branch of V84.  Recall
\begin{equation}
 \delta\psi={\psi_0\epsilon-\psi_0^{-7}\beta\over\mu_0+8}\cos y,
 \qquad
 \delta\mathfrak a=\psi_0^{-4}\epsilon+\psi_0^{-12}\beta
 \label{eq:v6v86-fixed-first}
\end{equation}
to first order.

\begin{theorem}[Exact first divergent coefficient]
\label{thm:v6v86-mean-U}
The spatial mean of the conformal derivative perturbation is
\begin{equation}
 \boxed{\quad
 \overline U_{\mathrm{vol}}(\epsilon,\beta)
 =4\epsilon+12\psi_0^{-8}\beta
 +O\bigl((|\epsilon|+|\beta|)^2\bigr).
 \quad}
 \label{eq:v6v86-mean-U}
\end{equation}
On the V84 first-order conformal-cancellation line
\(\beta=\psi_0^8\epsilon\),
\begin{equation}
 \overline U_{\mathrm{vol}}
 =16\epsilon+O(\epsilon^2).
 \label{eq:v6v86-mean-U-cancellation}
\end{equation}
Hence cancellation of the first \(k_0\) conformal coefficient does not
cancel the linear ultraviolet divergence of the raw coarea determinant.
\end{theorem}

\begin{proof}
The multiplication part of the derivative is
\begin{equation}
 Q=7S\psi^{-8}-\epsilon h_++5\mathfrak a\psi^4.
 \label{eq:v6v86-Q}
\end{equation}
At the base point \(Q=\mu_0\).  The variation of
\(7s_*\psi^{-8}+5\mathfrak a_0\psi^4\) due to \(\delta\psi\) is
\begin{equation}
 (-56s_*\psi_0^{-9}+20\mathfrak a_0\psi_0^3)\delta\psi
 =-36\mathfrak a_0\psi_0^3\delta\psi,
 \label{eq:v6v86-Q-psi}
\end{equation}
which has zero mean by \eqref{eq:v6v86-fixed-first}.  For the
\(\epsilon\)-variation, the explicit term \(-h_+\) has mean \(-1\), while
the gap term \(5\psi_0^4\partial_\epsilon\mathfrak a\) equals \(5\).
Their mean is \(4\).  For the \(\beta\)-variation, the direct singular
coefficient gives \(7h_-\psi_0^{-8}\), of mean
\(7\psi_0^{-8}\), and the gap term gives
\(5\psi_0^4\partial_\beta\mathfrak a=5\psi_0^{-8}\).  Their mean is
\(12\psi_0^{-8}\).  The longitudinal vector term begins at
\(\epsilon\beta\) and does not enter this first-order calculation.  This
proves \eqref{eq:v6v86-mean-U}; substitution of the cancellation line gives
\eqref{eq:v6v86-mean-U-cancellation}.
\end{proof}

\begin{assessment}[A mode cancellation is not a determinant cancellation]
\label{ass:v6v86-not-determinant-cancellation}
The V84 line removes one nonzero Fourier coefficient of the constrained
conformal factor at first order.  The determinant samples every conformal
fluctuation mode, and its leading ultraviolet term depends on the spatial
mean of the derivative potential.  That mean is strictly positive to first
order for nonnegative \(\epsilon,\beta\), including on the cancellation
line.  The two cancellations are mathematically different.
\end{assessment}

\section{Renormalized local conformal Jacobian}
\label{sec:v6v86-renormalized}

\begin{definition}[Modified conformal coarea factor]
\label{def:v6v86-renormalized-factor}
Define the finite-cutoff counterterm and renormalized inverse determinant by
\begin{align}
 C_\Lambda(U)&=\overline U\mathscr S_\Lambda,
 \label{eq:v6v86-counterterm}\\
 \mathcal D_{\mathrm{conf},\Lambda}^{\mathrm{ren}}(U)
 &:=e^{C_\Lambda(U)}
 {\det\mathcal L_{0,\Lambda}\over\det\mathcal L_\Lambda}
 =\det{}_2(I+K_\Lambda)^{-1}.
 \label{eq:v6v86-renormalized-factor}
\end{align}
\end{definition}

\begin{theorem}[Positive continuum limit of the local conformal factor]
\label{thm:v6v86-renormalized-limit}
If \(U\) is smooth and \(\mathcal L_0+U\ge\delta I>0\), then
\begin{equation}
 \mathcal D_{\mathrm{conf},\Lambda}^{\mathrm{ren}}(U)
 \longrightarrow\det{}_2(I+K)^{-1}\in(0,\infty).
 \label{eq:v6v86-renormalized-limit}
\end{equation}
The convergence is locally uniform on sets for which \(K\) is bounded in
Hilbert--Schmidt norm and \(I+K\) stays uniformly invertible.
\end{theorem}

\begin{proof}
The identity is Definition~\ref{def:v6v86-renormalized-factor}; convergence
and positivity are Theorem~\ref{thm:v6v86-det-split}.  Continuity of
\(\det_2\) on bounded Hilbert--Schmidt sets, together with a uniform spectral
gap preventing zeros, gives local uniformity.
\end{proof}

\begin{firewall}[A declared subtraction, not a derived SM counterterm]
\label{fw:v6v86-counterterm}
The factor \(e^{C_\Lambda(U)}\) is the exact subtraction needed to convert
the ordinary relative conformal determinant into \(\det_2\).  V86 does not
derive this subtraction from the bare Einstein--Standard Model action,
prove compatibility with BRST symmetry, or combine it with gauge, ghost,
fermion, and TT determinants.  Those sectors may cancel, reinforce, or alter
the local trace term only after their common regulator and measure are
specified.
\end{firewall}

\begin{definition}[Conformal determinant card, CDC]
\label{def:v6v86-CDC}
CDC requires:
\begin{enumerate}[label=\textup{(CD\arabic*)},leftmargin=5.2em]
\item a uniform positive spectral gap for the conformal derivative on the
      admitted source domain;
\item normalization by a fixed reference determinant;
\item explicit subtraction or cancellation of
      \(\overline U\mathscr S_\Lambda\);
\item convergence of the remaining modified determinant;
\item a common regulator computation with vector, gauge, ghost, fermion,
      and TT Jacobians; and
\item compatibility of the resulting density with chart overlaps and the
      York generator.
\end{enumerate}
\end{definition}

\begin{theorem}[CDC progress]
\label{thm:v6v86-CDC}
V86 closes CDC\textup{(CD1)} on every box satisfying
\eqref{eq:v6v86-box-condition}, closes CDC\textup{(CD2)}--
\textup{(CD4)} for the isolated conformal block through
\eqref{eq:v6v86-renormalized-factor}, and leaves
CDC\textup{(CD5)}--\textup{(CD6)} open.
\end{theorem}

\begin{proof}
Use Theorem~\ref{thm:v6v86-gap},
Theorem~\ref{thm:v6v86-det-split}, and
Theorem~\ref{thm:v6v86-renormalized-limit}, subject to
Firewall~\ref{fw:v6v86-counterterm}.
\end{proof}

\begin{programme}[V87: common-regulator determinant ledger]
\label{prog:v6v86-V87}
\begin{enumerate}[leftmargin=3em]
\item Put the conformal, mean-zero vector, spatial gauge, and ghost
      derivatives on the same Fourier cutoff and normalize all determinants
      at the same homogeneous base point.
\item Separate source-independent determinants, linear local trace terms,
      and Hilbert--Schmidt modified determinants block by block.
\item Test whether the vector/gauge/ghost supertrace cancels the conformal
      coefficient \eqref{eq:v6v86-mean-U}; do not assume BRST cancellation
      before matching operators and multiplicities.
\item Add matter Hessians only after their conformal weights and statistics
      are fixed.
\item If the linear supertrace does not cancel, record the exact local
      counterterm required before studying chart-boundary limits.
\end{enumerate}
\end{programme}

\begin{firewall}[Claim boundary after V86]
\label{fw:v6v86-boundary}
V86 treats the isolated scalar conformal Jacobian near one regular CMC
branch.  It does not prove a common BRST determinant cancellation, global
gauge fixing, control at a constraint fold, a canonical quantum York
Hamiltonian, regulator-independent SM source assembly, or the full
SM--Einstein continuum.
\end{firewall}

\begin{theorem}[Sixth-series V86 result]
\label{thm:v6v86-status}
A positive pointwise bound \eqref{eq:v6v86-delta} excludes conformal folds
and controls the inverse operator.  It does not control the raw determinant:
the relative inverse coarea factor contains the exact local divergence
\(\exp\{-\overline U\mathscr S_\Lambda\}\), with
\(\mathscr S_\Lambda\asymp\Lambda\).  On the fixed-volume beat branch,
\(\overline U=4\epsilon+12\psi_0^{-8}\beta+O(2)\), and it remains nonzero
on the first-order beat-cancellation line.  Subtracting that trace converts
the factor exactly to a positive inverse \(\det_2\) with a continuum limit.
The remaining gate is the common-regulator determinant ledger of all
constraint, gauge, ghost, and matter blocks.
\end{theorem}

\begin{proof}
Use Theorems~\ref{thm:v6v86-gap},
\ref{thm:v6v86-det-split},
\ref{thm:v6v86-mean-U}, and
\ref{thm:v6v86-renormalized-limit}, subject to
Firewalls~\ref{fw:v6v86-counterterm} and \ref{fw:v6v86-boundary}.
\end{proof}

The next compulsory YM/NS audit is V90.
\clearpage
\chapter{Sixth-series V87: the common-regulator supertrace ledger and the
principal-symbol gate}
\label{ch:v6v87-supertrace}

\section{V87 entry point}
\label{sec:v6v87}

V86 renormalized the isolated conformal Jacobian by subtracting its linear
trace.  A full constrained Standard Model--Einstein measure cannot choose
that subtraction block by block without first testing gauge, ghost, vector,
TT, and matter contributions on one regulator.  This continuation proves a
common Fourier-cutoff ledger for every Laplace-type block with fixed
principal part and order-zero source perturbation.  Its linear divergence is
a local weighted supertrace, while all Hilbert--Schmidt remainders converge
as modified determinants.  Applied to the matched CMC sector, the vector and
fixed background-gauge blocks are source independent and do not cancel the
conformal trace.  The calculation also identifies a prior obstruction for a
nonlinear physical-metric gauge: a changing principal symbol is not a
Hilbert--Schmidt perturbation and cannot be inserted into the V86
\(\det_2\) formula.

\section{A common class of determinant blocks}
\label{sec:v6v87-blocks}

Let \(E_a\) be a finite-dimensional Hermitian fibre of rank \(r_a\), and
on \(L^2(\mathbb T^3;E_a)\) put
\begin{equation}
 L_{a,0}=-c_a\Delta+\mu_a,
 \qquad c_a,\mu_a>0,
 \qquad L_a=L_{a,0}+U_a(x),
 \label{eq:v6v87-block}
\end{equation}
where \(U_a\) is a smooth self-adjoint endomorphism.  Assign a real exponent
\(\sigma_a\): for example an inverse constraint Jacobian has exponent
\(-1\), a square root from a real bosonic Gaussian has exponent
\(-1/2\), and a determinant in a gauge numerator has positive exponent.
The exponent and determinant orientation must be declared for each block.

Use the same Fourier ball \(|k|\le\Lambda\) in every block, and define
\begin{equation}
 K_a=L_{a,0}^{-1/2}U_aL_{a,0}^{-1/2},
 \qquad K_{a,\Lambda}=P_\Lambda K_aP_\Lambda.
 \label{eq:v6v87-Ka}
\end{equation}

\begin{theorem}[Common-regulator determinant ledger]
\label{thm:v6v87-ledger}
Assume each \(L_a\) is invertible.  The normalized determinant product
\begin{equation}
 \mathcal R_\Lambda
 :=\prod_a\left({\det L_{a,\Lambda}
                         \over\det L_{a,0,\Lambda}}
          \right)^{\sigma_a}
 \label{eq:v6v87-R}
\end{equation}
obeys the exact identity
\begin{align}
 \log\mathcal R_\Lambda
 ={}&\mathcal T_\Lambda
 +\sum_a\sigma_a\log\det{}_2(I+K_{a,\Lambda}),
 \label{eq:v6v87-ledger}\\
 \mathcal T_\Lambda
 :={}&\sum_a\sigma_a\,
 \operatorname{tr}_{E_a}(\overline U_a)\,
 \mathscr S_{a,\Lambda},
 \qquad
 \mathscr S_{a,\Lambda}
 :=\sum_{|k|\le\Lambda}{1\over c_a|k|^2+\mu_a},
 \label{eq:v6v87-supertrace}
\end{align}
where \(\overline U_a=\int U_a(x)d\mu\).  Every modified-determinant term
has a finite limit.  Moreover there is a torus constant \(C_{\rm lat}>0\),
independent of the block, such that
\begin{equation}
 {\mathscr S_{a,\Lambda}\over\Lambda}
 \longrightarrow {C_{\rm lat}\over c_a}.
 \label{eq:v6v87-S-asymptotic}
\end{equation}
Consequently a necessary condition for the unrenormalized product to have a
finite nonzero limit for a fixed source is
\begin{equation}
 \boxed{\quad
 \sum_a{\sigma_a\over c_a}
 \operatorname{tr}_{E_a}(\overline U_a)=0.
 \quad}
 \label{eq:v6v87-supertrace-cancellation}
\end{equation}
\end{theorem}

\begin{proof}
For every block, the V86 finite-dimensional identity gives
\[
 \log{\det L_{a,\Lambda}\over\det L_{a,0,\Lambda}}
 =\operatorname{Tr}K_{a,\Lambda}
 +\log\det{}_2(I+K_{a,\Lambda}).
\]
The Fourier diagonal of \(K_a\) is
\(\overline U_a/(c_a|k|^2+\mu_a)\), and taking the fibre trace gives
\eqref{eq:v6v87-supertrace}.  The V86 Hilbert--Schmidt proof applies to
matrix-valued smooth \(U_a\), so each \(K_a\) is Hilbert--Schmidt and its
modified determinant converges.  Comparing the lattice sum with
\(\int_{|\xi|\le\Lambda}(c_a|\xi|^2+\mu_a)^{-1}d\xi\) gives
\eqref{eq:v6v87-S-asymptotic}; the leading radial integral is proportional
to \(\Lambda/c_a\), and the lattice error is lower order.  Divide
\eqref{eq:v6v87-supertrace} by \(\Lambda\).  If
\eqref{eq:v6v87-supertrace-cancellation} fails, the logarithm of
\(\mathcal R_\Lambda\) has a nonzero linear divergence while its modified
part converges, excluding a finite nonzero limit.
\end{proof}

\begin{corollary}[Exact cancellation with identical principal blocks]
\label{cor:v6v87-exact-cancellation}
If all blocks have the same \((c_a,\mu_a)=(c,\mu)\) and
\begin{equation}
 \sum_a\sigma_a\operatorname{tr}_{E_a}(\overline U_a)=0,
 \label{eq:v6v87-identical-supertrace}
\end{equation}
then \(\mathcal T_\Lambda=0\) at every cutoff, and
\begin{equation}
 \mathcal R_\Lambda\longrightarrow
 \prod_a\det{}_2(I+K_a)^{\sigma_a}.
 \label{eq:v6v87-product-limit}
\end{equation}
If all \(I+K_a\) are positive, the limit is positive.
\end{corollary}

\begin{proof}
The common lattice sum factors out of
\eqref{eq:v6v87-supertrace}; use
\eqref{eq:v6v87-identical-supertrace} and then pass to the modified-
determinant limits in \eqref{eq:v6v87-ledger}.
\end{proof}

\begin{definition}[Linear local supertrace counterterm]
\label{def:v6v87-counterterm}
When \eqref{eq:v6v87-supertrace-cancellation} is not supplied by the field
content, define at finite cutoff
\begin{equation}
 C_\Lambda^{\rm loc}:=\mathcal T_\Lambda,
 \qquad
 \mathcal R_\Lambda^{\rm ren}
 :=e^{-C_\Lambda^{\rm loc}}\mathcal R_\Lambda.
 \label{eq:v6v87-renormalized-product}
\end{equation}
Then
\begin{equation}
 \mathcal R_\Lambda^{\rm ren}
 \longrightarrow\prod_a\det{}_2(I+K_a)^{\sigma_a}.
 \label{eq:v6v87-renormalized-limit}
\end{equation}
\end{definition}

\begin{assessment}[Why the regulator must be common]
\label{ass:v6v87-common-regulator}
The coefficient in \eqref{eq:v6v87-supertrace-cancellation} contains the
principal constant \(c_a^{-1}\), the fibre multiplicity, and the determinant
exponent.  Independent cutoffs or independently normalized traces can
manufacture or destroy a cancellation.  A claimed BRST or supersymmetric
cancellation is meaningful only after all of these data are matched on the
same spectral counting rule.
\end{assessment}

\section{Application to the matched CMC constraint blocks}
\label{sec:v6v87-application}

Normalize every constraint block at the same homogeneous CMC base point.
In the restricted matched model:
\begin{enumerate}[label=\textup{(\roman*)},leftmargin=3.7em]
\item the conformal coarea factor has exponent \(\sigma_{\rm conf}=-1\),
      principal constant \(c_{\rm conf}=8\), and perturbation mean from
      V86;
\item the mean-zero vector constraint derivative is the fixed operator
      \(\bar\Delta_L\), independent of \(\epsilon,\beta\); and
\item a fixed linear Fourier spatial gauge and its ghost derivative at this
      restricted background are source independent.
\end{enumerate}

\begin{theorem}[No cancellation inside the fixed constraint/gauge blocks]
\label{thm:v6v87-no-cancellation}
For the three blocks just listed, the source-dependent linear trace in the
logarithm of the reduced density is entirely conformal.  On the fixed-volume
branch,
\begin{equation}
 \mathcal T_\Lambda^{\rm matched}
 =-\bigl(4\epsilon+12\psi_0^{-8}\beta+O(2)\bigr)
 \mathscr S_\Lambda,
 \qquad
 \mathscr S_\Lambda=\sum_{|k|\le\Lambda}{1\over8|k|^2+\mu_0}.
 \label{eq:v6v87-matched-trace}
\end{equation}
On \(\beta=\psi_0^8\epsilon\), this becomes
\begin{equation}
 \mathcal T_\Lambda^{\rm matched}
 =-\bigl(16\epsilon+O(\epsilon^2)\bigr)\mathscr S_\Lambda.
 \label{eq:v6v87-cancellation-line-trace}
\end{equation}
Thus neither the vector constraint determinant nor the fixed linear
gauge/ghost pair cancels the V86 conformal divergence.
\end{theorem}

\begin{proof}
Source-independent normalized determinants have \(U_a=0\) and contribute
nothing to \eqref{eq:v6v87-supertrace}.  The inverse conformal coarea
determinant contributes with exponent \(-1\).  Insert
Theorem~\ref{thm:v6v86-mean-U} and its cancellation-line specialization.
\end{proof}

\begin{firewall}[Restricted gauge result]
\label{fw:v6v87-restricted-gauge}
Theorem~\ref{thm:v6v87-no-cancellation} concerns the fixed linear
background gauge used in the local Fourier chart.  A nonlinear gauge tied to
the physical metric has a source-dependent ghost operator.  Its determinant
must be recomputed; it cannot be declared cancelling merely because a linear
background ghost is source independent.
\end{firewall}

\section{The principal-symbol gate}
\label{sec:v6v87-principal-symbol}

The preceding theorem assumes that every source perturbation is order zero
relative to a fixed Laplace principal part.  Nonlinear gauge and matter
operators often violate that assumption.

\begin{theorem}[Changing a second-order principal symbol is not a
Hilbert--Schmidt perturbation]
\label{thm:v6v87-principal-no-go}
Let \(L_0\) and \(L\) be positive elliptic second-order operators on a
closed three-manifold.  If their principal symbols differ on a nonempty open
conic set, then
\begin{equation}
 L_0^{-1/2}(L-L_0)L_0^{-1/2}
 \label{eq:v6v87-order-zero-relative}
\end{equation}
is a pseudodifferential operator of order zero with nonzero principal
symbol.  It is not compact on \(L^2\), hence is not Hilbert--Schmidt, and the
\(\det_2\) construction of V86--V87 does not apply.
\end{theorem}

\begin{proof}
Each factor \(L_0^{-1/2}\) has order \(-1\), while the nonzero principal
part of \(L-L_0\) has order two.  Their composition has order zero and
principal symbol
\(\sigma_2(L_0)^{-1/2}[\sigma_2(L)-\sigma_2(L_0)]
\sigma_2(L_0)^{-1/2}\), nonzero on the stated set.  A classical order-zero
pseudodifferential operator with nonvanishing principal symbol has nonzero
essential norm: applying it to normalized wave packets with frequencies
tending to infinity in that conic set prevents any subsequence from
converging after the operator is applied.  It is therefore not compact.
Every Hilbert--Schmidt operator is compact.
\end{proof}

\begin{corollary}[Physical-metric ghost operators require heat-kernel
renormalization]
\label{cor:v6v87-ghost-gate}
If a nonlinear gauge ghost or matter Hessian is defined using the physical
metric \(g=\psi^4\bar g\), comparison with a fixed flat reference Laplacian
generally changes the principal symbol.  Its determinant cannot be combined
with \eqref{eq:v6v87-ledger} until the principal symbols are aligned by a
field-dependent unitary/conformal transformation or the determinant is
renormalized covariantly by heat-kernel coefficients.
\end{corollary}

\begin{proof}
The inverse physical metric occurs in the leading symbol.  Unless
\(\psi\) is the reference constant, it differs from the flat leading
symbol.  Apply Theorem~\ref{thm:v6v87-principal-no-go}.
\end{proof}

\begin{definition}[Common determinant ledger card, CDLC]
\label{def:v6v87-CDLC}
CDLC requires:
\begin{enumerate}[label=\textup{(DL\arabic*)},leftmargin=5.2em]
\item one regulator, reference density, zero-mode convention, and
      determinant orientation for every block;
\item alignment of principal symbols, or a covariant heat-kernel
      subtraction through every divergent coefficient;
\item computation of the linear local supertrace
      \eqref{eq:v6v87-supertrace-cancellation};
\item proof of its cancellation by field content or declaration of the
      corresponding local counterterm;
\item convergence and nonvanishing of the modified determinant product; and
\item BRST/gauge independence and chart compatibility of the resulting
      reduced measure.
\end{enumerate}
\end{definition}

\begin{theorem}[CDLC progress]
\label{thm:v6v87-CDLC}
V87 closes CDLC\textup{(DL1)}--\textup{(DL4)} for the fixed-principal-symbol
model as a conditional algebraic ledger: the exact cancellation condition
and counterterm are \eqref{eq:v6v87-supertrace-cancellation} and
\eqref{eq:v6v87-renormalized-product}.  It closes
CDLC\textup{(DL5)} when all blocks are positive fixed-principal-symbol
perturbations.  The physical-metric gauge and full SM Hessians do not yet
satisfy CDLC\textup{(DL2)}, and CDLC\textup{(DL6)} remains open.
\end{theorem}

\begin{proof}
Use Theorem~\ref{thm:v6v87-ledger},
Definition~\ref{def:v6v87-counterterm}, and
Theorem~\ref{thm:v6v87-principal-no-go}, subject to
Firewall~\ref{fw:v6v87-restricted-gauge}.
\end{proof}

\begin{programme}[V88: covariant heat-kernel subtraction]
\label{prog:v6v87-V88}
\begin{enumerate}[leftmargin=3em]
\item Put a scalar Laplace-type operator for the physical conformal metric
      into a fixed reference Hilbert space by the exact volume-density
      unitary map.
\item Compute which principal and first-order coefficients remain after
      conjugation.
\item Derive the first heat-kernel divergences of its relative determinant
      on the three-torus and identify their local geometric densities.
\item Compare those densities with the conformal coarea trace term before
      adding vector ghosts.
\item If the leading-symbol mismatch cannot be controlled uniformly on the
      matched source box, restrict the chart by the V83 conformal bounds and
      formulate the precise ellipticity constants required.
\end{enumerate}
\end{programme}

\begin{firewall}[Claim boundary after V87]
\label{fw:v6v87-boundary}
V87 proves a determinant ledger only for finitely many fixed-principal-
symbol Laplace blocks and a no-go for applying it directly to changing
principal symbols.  It does not compute the nonlinear gravitational ghost,
full SM Hessians, BRST supertrace, heat-kernel counterterms, global reduced
measure, quantum York dynamics, or the full SM--Einstein continuum.
\end{firewall}

\begin{theorem}[Sixth-series V87 result]
\label{thm:v6v87-status}
For fixed-principal-symbol Laplace blocks on one Fourier regulator, every
ordinary determinant splits into a local linear trace and a convergent
modified determinant.  The necessary no-counterterm condition is the
weighted supertrace \eqref{eq:v6v87-supertrace-cancellation}.  In the
matched CMC constraint plus fixed background gauge, source-independent
vector and ghost blocks do not cancel the conformal divergence, including
on the V84 beat-cancellation line.  A nonlinear physical-metric gauge changes
the leading symbol, making the relative perturbation noncompact and forcing
a covariant heat-kernel treatment before a full determinant cancellation can
be assessed.
\end{theorem}

\begin{proof}
Use Theorems~\ref{thm:v6v87-ledger},
\ref{thm:v6v87-no-cancellation}, and
\ref{thm:v6v87-principal-no-go}, subject to
Firewalls~\ref{fw:v6v87-restricted-gauge} and \ref{fw:v6v87-boundary}.
\end{proof}

The next compulsory YM/NS audit is V90.
\clearpage
\chapter{Sixth-series V88: physical-metric scalar determinants, covariant
heat subtraction, and the fixed-volume curvature term}
\label{ch:v6v88-heat}

\section{V88 entry point}
\label{sec:v6v88}

V87 proved that a changing principal symbol is not a Hilbert--Schmidt
perturbation of a fixed flat Laplacian.  The appropriate determinant is then
defined covariantly from its heat trace.  This version carries out that
construction for a massive scalar Laplacian on the conformal metric
\(g=\psi^4\bar g\).  A density-unitary map gives its exact flat-space
differential expression and uniform ellipticity.  The heat expansion shows
that in three dimensions only the volume and integrated scalar-curvature
coefficients diverge in proper time.  Fixing physical volume removes the
leading divergence.  The remaining linear divergence is proportional to
\(\int R_gd\mu_g=8\int|\bar\nabla\psi|^2d\mu_{\bar g}\), whose first
nonzero response is computed on the V84 branch.

\section{Exact conjugation of the scalar Laplacian}
\label{sec:v6v88-conjugation}

Let \(\bar g\) be flat and normalized, let \(g=\psi^4\bar g\), and define
\begin{equation}
 D_g=-\Delta_g+m^2,
 \qquad m^2>0.
 \label{eq:v6v88-Dg}
\end{equation}
The positive mass is included to remove the scalar zero mode; a massless
determinant requires a separately declared primed determinant.

\begin{lemma}[Density-unitary conjugation]
\label{lem:v6v88-conjugation}
The map
\begin{equation}
 \mathcal U_\psi:L^2(\mathbb T^3,d\mu_g)
 \longrightarrow L^2(\mathbb T^3,d\mu_{\bar g}),
 \qquad \mathcal U_\psi f=\psi^3f,
 \label{eq:v6v88-unitary}
\end{equation}
is unitary.  On smooth functions,
\begin{align}
 \mathcal U_\psi D_g\mathcal U_\psi^{-1}
 ={}&-\bar\nabla_i(\psi^{-4}\bar\nabla_i)
 +3\psi^{-5}\bar\Delta\psi
 -6\psi^{-6}|\bar\nabla\psi|^2+m^2
 \label{eq:v6v88-conjugated}\\
 ={}&-\psi^{-4}\bar\Delta
 +4\psi^{-5}\bar\nabla\psi\cdot\bar\nabla
 +3\psi^{-5}\bar\Delta\psi
 -6\psi^{-6}|\bar\nabla\psi|^2+m^2.
 \nonumber
\end{align}
Its principal symbol is \(\psi^{-4}|\xi|_{\bar g}^2\).
\end{lemma}

\begin{proof}
Since \(d\mu_g=\psi^6d\mu_{\bar g}\), multiplication by \(\psi^3\)
preserves the norm.  For a scalar,
\begin{equation}
 \Delta_gf=\psi^{-6}\bar\nabla_i
 (\psi^2\bar\nabla_i f).
 \label{eq:v6v88-conformal-Laplacian}
\end{equation}
Insert \(f=\psi^{-3}u\), expand both derivatives, and multiply the result by
\(\psi^3\).  The second-order and first-order terms combine into
\(-\bar\nabla_i(\psi^{-4}\bar\nabla_i)\), yielding
\eqref{eq:v6v88-conjugated}.  Its leading coefficient gives the principal
symbol.
\end{proof}

\begin{corollary}[Uniform ellipticity on the matched source box]
\label{cor:v6v88-ellipticity}
If \(0<c\le\psi\le C\), then
\begin{equation}
 C^{-4}|\xi|^2\le
 \sigma_2(D_g)(x,\xi)\le c^{-4}|\xi|^2.
 \label{eq:v6v88-ellipticity}
\end{equation}
Thus the V83 uniform conformal bounds give cutoff-independent ellipticity
constants for this block.
\end{corollary}

\begin{proof}
Apply the pointwise bounds to the principal symbol in
Lemma~\ref{lem:v6v88-conjugation}.
\end{proof}

\section{Proper-time determinant and its divergent coefficients}
\label{sec:v6v88-proper-time}

For \(s>0\), define the proper-time regulated logarithm
\begin{equation}
 \log\det_sD_g:=-\int_s^\infty {dt\over t}
 \operatorname{Tr}(e^{-tD_g}).
 \label{eq:v6v88-proper-time-det}
\end{equation}
The mass makes the large-time integral finite.

\begin{theorem}[Three-dimensional scalar heat divergences]
\label{thm:v6v88-heat-divergences}
As \(t\downarrow0\),
\begin{equation}
 \operatorname{Tr}(e^{-tD_g})
 =(4\pi t)^{-3/2}
 \left[A_0(g)+tA_1(g,m)+O(t^2)\right],
 \label{eq:v6v88-heat-expansion}
\end{equation}
where
\begin{align}
 A_0(g)&=\operatorname{Vol}(g),
 \label{eq:v6v88-A0}\\
 A_1(g,m)&={1\over6}\int_{\mathbb T^3}R_gd\mu_g
 -m^2\operatorname{Vol}(g).
 \label{eq:v6v88-A1}
\end{align}
For a reference metric \(g_0\), the relative regulated determinant has
divergent part
\begin{align}
 \log{\det_sD_g\over\det_sD_{g_0}}
 =-{1\over(4\pi)^{3/2}}\biggl[&{2\over3}s^{-3/2}
     \bigl(A_0(g)-A_0(g_0)\bigr)
 \nonumber\\
 &+2s^{-1/2}
     \bigl(A_1(g,m)-A_1(g_0,m)\bigr)\biggr]+O(1).
 \label{eq:v6v88-relative-divergence}
\end{align}
After subtracting the two displayed local terms, the relative proper-time
integral has a finite limit as \(s\downarrow0\).
\end{theorem}

\begin{proof}
The scalar Laplace heat parametrix on a closed manifold gives
\(\operatorname{Tr}(e^{t\Delta_g})
=(4\pi t)^{-3/2}[\operatorname{Vol}(g)+
t\int R_g/6+O(t^2)]\).  Since \(m^2\) is constant,
\(e^{-tD_g}=e^{-tm^2}e^{t\Delta_g}\); multiplying by
\(1-tm^2+O(t^2)\) gives
\eqref{eq:v6v88-heat-expansion}--\eqref{eq:v6v88-A1}.  Insert the relative
expansion in \eqref{eq:v6v88-proper-time-det}.  The first two integrals are
\(\int_s t^{-5/2}dt=(2/3)s^{-3/2}+O(1)\) and
\(\int_s t^{-3/2}dt=2s^{-1/2}+O(1)\), proving
\eqref{eq:v6v88-relative-divergence}.  The next heat term contributes
\(O(t^{-1/2})\) after division by \(t\), which is integrable at zero; the
positive mass controls infinity.
\end{proof}

\begin{definition}[Covariantly renormalized scalar determinant]
\label{def:v6v88-renormalized-det}
Define
\begin{align}
 \log\det_{\rm ren}(D_g,D_{g_0})
 :=\lim_{s\downarrow0}\biggl\{
 \log{\det_sD_g\over\det_sD_{g_0}}
 +{1\over(4\pi)^{3/2}}\biggl[&{2\over3}s^{-3/2}\Delta A_0
 \nonumber\\
 &+2s^{-1/2}\Delta A_1\biggr]\biggr\}.
 \label{eq:v6v88-renormalized-det}
\end{align}
Theorem~\ref{thm:v6v88-heat-divergences} proves that this limit exists.
\end{definition}

\section{Exact conformal curvature identity}
\label{sec:v6v88-curvature}

\begin{lemma}[Integrated curvature in the flat conformal class]
\label{lem:v6v88-curvature}
For \(g=\psi^4\bar g\) with flat \(\bar g\),
\begin{align}
 R_g&=-8\psi^{-5}\bar\Delta\psi,
 \label{eq:v6v88-R-pointwise}\\
 \int_{\mathbb T^3}R_gd\mu_g
 &=8\int_{\mathbb T^3}|\bar\nabla\psi|^2d\mu_{\bar g}\ge0.
 \label{eq:v6v88-R-integrated}
\end{align}
Equality holds exactly when \(\psi\) is constant.
\end{lemma}

\begin{proof}
The three-dimensional conformal scalar-curvature formula gives the first
line.  Multiply by \(d\mu_g=\psi^6d\mu_{\bar g}\) and integrate:
\(-8\int\psi\bar\Delta\psi=8\int|\bar\nabla\psi|^2\).  The last integral
vanishes precisely for a constant function on the connected torus.
\end{proof}

\begin{corollary}[Fixed volume removes the leading heat divergence]
\label{cor:v6v88-fixed-volume}
On the V84 fixed-physical-volume branch, with constant reference
\(g_0=\psi_0^4\bar g\),
\begin{equation}
 \Delta A_0=0,
 \qquad
 \Delta A_1={4\over3}
 \int|\bar\nabla\psi|^2d\mu_{\bar g}.
 \label{eq:v6v88-fixed-heat-coefficients}
\end{equation}
Hence its relative scalar determinant has no \(s^{-3/2}\) divergence; its
remaining \(s^{-1/2}\) divergence is local, nonnegative, and vanishes only
at a constant conformal factor.
\end{corollary}

\begin{proof}
Fixed volume gives \(A_0(g)=A_0(g_0)\), so the mass contribution to
\(\Delta A_1\) also cancels.  Use
Lemma~\ref{lem:v6v88-curvature} in \(\Delta A_1=\int R_g/6\).
\end{proof}

\section{Response on the two-source branch}
\label{sec:v6v88-response}

V84 gives, at fixed volume,
\begin{equation}
 \psi(y)=\psi_0+A_{\epsilon,\beta}\cos y+O(2),
 \qquad
 A_{\epsilon,\beta}
 ={\psi_0\epsilon-\psi_0^{-7}\beta\over\mu_0+8}.
 \label{eq:v6v88-A-response}
\end{equation}

\begin{theorem}[First nonzero curvature divergence]
\label{thm:v6v88-response}
On the fixed-volume branch,
\begin{align}
 \int R_gd\mu_g
 &=4A_{\epsilon,\beta}^2
 +O\bigl((|\epsilon|+|\beta|)^3\bigr),
 \label{eq:v6v88-R-response}\\
 \Delta A_1
 &={2\over3}A_{\epsilon,\beta}^2
 +O\bigl((|\epsilon|+|\beta|)^3\bigr).
 \label{eq:v6v88-A1-response}
\end{align}
The curvature part of the scalar determinant divergence therefore vanishes
through quadratic order on the V84 cancellation line
\(\beta=\psi_0^8\epsilon\).  The conformal inverse-coarea divergence of V86
does not vanish there and remains linear in \(\epsilon\).
\end{theorem}

\begin{proof}
Differentiate \eqref{eq:v6v88-A-response}:
\(\bar\nabla\psi=-A_{\epsilon,\beta}\sin y\,e_2+O(2)\).
Normalized Haar integration gives
\(\int\sin^2y=1/2\), so
\(8\int|\bar\nabla\psi|^2=4A_{\epsilon,\beta}^2+O(3)\).
This proves \eqref{eq:v6v88-R-response}; divide by six for
\eqref{eq:v6v88-A1-response}.  The response amplitude is zero to first
order on the cancellation line.  In contrast,
Theorem~\ref{thm:v6v86-mean-U} gives
\(\overline U=16\epsilon+O(\epsilon^2)\) there.
\end{proof}

\begin{assessment}[Two local counterterms have different ancestry]
\label{ass:v6v88-two-counterterms}
The conformal coarea trace divergence is linear in the mean perturbation of
the constraint derivative.  The physical scalar Laplacian's remaining
fixed-volume divergence is quadratic in the gradient of the conformal
factor.  The first-order beat-cancellation line removes the latter through
second order but not the former.  They cannot be cancelled against one
another by matching only their cutoff power; their local field densities
and source orders differ.
\end{assessment}

\section{Uniformity and the covariant determinant card}
\label{sec:v6v88-CHDC}

\begin{theorem}[Uniform heat construction on a compact matched box]
\label{thm:v6v88-uniform}
Let \(K\) be a compact regular matched-source box on which
\(0<c\le\psi\le C\) and finitely many derivatives of \(\psi\) are uniformly
bounded.  Then the operators \(D_g\) form a uniformly elliptic family, their
small-time heat remainders after \(A_0,A_1\) subtraction are locally uniform
on \(K\), and \(\det_{\rm ren}(D_g,D_{g_0})\) is a finite continuous
function of the source labels.
\end{theorem}

\begin{proof}
Corollary~\ref{cor:v6v88-ellipticity} gives common ellipticity constants.
Uniform coefficient bounds in the conjugated expression
\eqref{eq:v6v88-conjugated} give the standard heat-parametrix remainder
uniformly on the compact parameter set.  The mass gives a uniform large-time
spectral gap after decreasing the box if necessary.  Dominated convergence
in the subtracted proper-time integral then proves finiteness and
continuity.
\end{proof}

\begin{definition}[Covariant heat-determinant card, CHDC]
\label{def:v6v88-CHDC}
CHDC requires:
\begin{enumerate}[label=\textup{(HD\arabic*)},leftmargin=5.2em]
\item exact unitary placement of every physical-metric operator in a
      declared reference Hilbert space;
\item uniform ellipticity and zero-mode control on the admitted source box;
\item computation of every divergent heat coefficient with statistics and
      multiplicities;
\item cancellation by the full field supertrace or declaration of the
      corresponding local covariant counterterms;
\item convergence and nonvanishing of the subtracted determinant product;
      and
\item BRST, chart, and regulator compatibility of the resulting measure.
\end{enumerate}
\end{definition}

\begin{theorem}[CHDC progress]
\label{thm:v6v88-CHDC}
V88 closes CHDC\textup{(HD1)}--\textup{(HD5)} for one positive massive
scalar Laplacian on a compact regular matched-source box, with the explicit
subtractions \eqref{eq:v6v88-A0}--\eqref{eq:v6v88-A1}.  It does not close the
common field supertrace or CHDC\textup{(HD6)}.
\end{theorem}

\begin{proof}
Use Lemma~\ref{lem:v6v88-conjugation},
Corollary~\ref{cor:v6v88-ellipticity},
Theorem~\ref{thm:v6v88-heat-divergences}, and
Theorem~\ref{thm:v6v88-uniform}.
\end{proof}

\begin{firewall}[Scalar test block only]
\label{fw:v6v88-scope}
The scalar Laplacian is a test determinant, not the full Higgs, gauge,
fermion, ghost, or graviton Hessian.  Its heat coefficients do not establish
a Standard Model supertrace cancellation.  The proper-time subtraction is a
declared renormalization prescription and has not been matched to a bare
Einstein--Standard Model action or an Osterwalder--Schrader transfer
construction.
\end{firewall}

\begin{programme}[V89: heat supertrace of minimal field blocks]
\label{prog:v6v88-V89}
\begin{enumerate}[leftmargin=3em]
\item Place scalar, one-form gauge, scalar ghost, and Dirac-square test
      operators on the same physical conformal metric and proper-time
      cutoff.
\item Record their ranks, statistics exponents, zero-mode conventions, and
      \(A_0,A_1\) coefficients without assuming cancellation.
\item Test the volume and curvature supertraces for one gauge generator and
      one Dirac fermion representation.
\item Separate physical Standard Model multiplicities from the test-block
      algebra before drawing any conclusion.
\item Add the conformal coarea determinant only after converting its
      reference-Fourier subtraction to the same covariant scheme.
\end{enumerate}
\end{programme}

\begin{firewall}[Claim boundary after V88]
\label{fw:v6v88-boundary}
V88 treats one positive massive scalar Laplacian on a compact regular CMC
box.  It does not compute the full gauge-fixed gravitational or Standard
Model heat supertrace, establish BRST cancellation, construct the global
reduced measure or quantum York generator, remove all regulators, or prove
the full SM--Einstein continuum.
\end{firewall}

\begin{theorem}[Sixth-series V88 result]
\label{thm:v6v88-status}
Density-unitary conjugation of the physical scalar Laplacian yields the
exact uniformly elliptic operator \eqref{eq:v6v88-conjugated}.  In three
dimensions its proper-time determinant requires only volume and integrated-
curvature subtractions.  Fixed physical volume removes the leading
\(s^{-3/2}\) term, and the remaining \(s^{-1/2}\) coefficient is
\((1/6)\int R_g=(4/3)\int|\nabla\psi|^2\).  On the two-source branch it is
quadratic in the uncancelled \(k_0\) response and vanishes through quadratic
order on the V84 cancellation line, while the conformal coarea divergence
remains.  The next gate is the common heat supertrace of minimal gauge,
ghost, fermion, and scalar blocks.
\end{theorem}

\begin{proof}
Use Lemma~\ref{lem:v6v88-conjugation},
Theorem~\ref{thm:v6v88-heat-divergences},
Lemma~\ref{lem:v6v88-curvature}, and
Theorem~\ref{thm:v6v88-response}, subject to
Firewalls~\ref{fw:v6v88-scope} and \ref{fw:v6v88-boundary}.
\end{proof}

The next compulsory YM/NS audit is V90.
\clearpage
\chapter{Sixth-series V89: minimal heat-supertrace test blocks and the
uncancelled curvature coefficient}
\label{ch:v6v89-supertrace}

\section{V89 entry point}
\label{sec:v6v89}

V88 treated one scalar determinant with changing principal symbol.  This
version places four standard test blocks on the same closed conformal
three-manifold and the same proper-time regulator: a real scalar, a
gauge-fixed one-form, its complex scalar Faddeev--Popov ghost, and one
complex two-component Dirac fermion through the modulus of its determinant.
Their first heat coefficients and determinant exponents are computed
explicitly.  One scalar, one gauge generator, and one Dirac field cancel the
leading volume coefficient, but their curvature coefficient is
\((1/4)\int R_g\), not zero.  This test also proves that the covariant field
supertrace cannot cancel the matched conformal coarea divergence near the
homogeneous base point: the former begins quadratically in the beat
response, whereas the latter begins linearly.

\section{Operators, exponents, and zero modes}
\label{sec:v6v89-operators}

Let \(g\) be a smooth metric on a closed oriented three-manifold.  Use
primed determinants, with every harmonic or Dirac zero mode removed and its
finite-dimensional integration recorded separately.  The priming changes
large-time behaviour and finite factors but not the local small-time heat
coefficients.

The test operators and their determinant exponents in the partition
function are:
\begin{center}
\begin{tabular}{c|c|c|c}
block & operator & complex fibre rank & exponent \(\sigma\)\\
\hline
real scalar & \(\Delta_0=-\nabla^2\) & \(1\) & \(-1/2\)\\
gauge one-form & \(\Delta_1=d\delta+\delta d\) & \(3\) & \(-1/2\)\\
complex scalar ghost & \(\Delta_0\) & \(1\) & \(+1\)\\
complex Dirac modulus & \(\slashed D^{,2}\) & \(2\) & \(+1/2\)
\end{tabular}
\end{center}
The last exponent follows from
\begin{equation}
 |\det{}'\slashed D|=det{}'(\slashed D^{,2})^{1/2};
 \label{eq:v6v89-Dirac-modulus}
\end{equation}
its phase and eta invariant are not included.

\begin{definition}[Heat supertrace coefficients]
\label{def:v6v89-supertrace}
If
\begin{equation}
 \operatorname{Tr}(e^{-tL_a})
 =(4\pi t)^{-3/2}
 \left[A_0(L_a)+tA_1(L_a)+O(t^2)\right],
 \label{eq:v6v89-heat}
\end{equation}
define
\begin{equation}
 A_j^{\mathrm{sup}}:=\sum_a\sigma_aA_j(L_a),
 \qquad j=0,1.
 \label{eq:v6v89-supertrace-def}
\end{equation}
\end{definition}

\section{The first coefficients}
\label{sec:v6v89-coefficients}

\begin{lemma}[Scalar, one-form, and Dirac-square coefficients]
\label{lem:v6v89-coefficients}
For the massless test operators above,
\begin{align}
 A_0(\Delta_0)&=\operatorname{Vol}(g),
 &A_1(\Delta_0)&={1\over6}\int R_gd\mu_g,
 \label{eq:v6v89-scalar-coefficients}\\
 A_0(\Delta_1)&=3\operatorname{Vol}(g),
 &A_1(\Delta_1)&=-{1\over2}\int R_gd\mu_g,
 \label{eq:v6v89-oneform-coefficients}\\
 A_0(\slashed D^{,2})&=2\operatorname{Vol}(g),
 &A_1(\slashed D^{,2})&=-{1\over6}\int R_gd\mu_g.
 \label{eq:v6v89-Dirac-coefficients}
\end{align}
\end{lemma}

\begin{proof}
For a Laplace-type operator written as
\(L=\nabla^*\nabla+E\), the first integrated coefficient is
\begin{equation}
 A_1(L)=\int\operatorname{tr}
 \left({R_g\over6}I-E\right)d\mu_g.
 \label{eq:v6v89-heat-A1-formula}
\end{equation}
The scalar has \(E=0\) and rank one.  The Weitzenb\"ock formula for one-
forms is \(\Delta_1=\nabla^*\nabla+\operatorname{Ric}\); its fibre trace is
\(R_g\), so \(3R_g/6-R_g=-R_g/2\).  The Lichnerowicz formula is
\(\slashed D^{,2}=\nabla^*\nabla+R_g/4\) on a rank-two complex spinor
bundle; hence \(2(R_g/6-R_g/4)=-R_g/6\).  The zeroth coefficient is the
fibre rank times volume.
\end{proof}

\begin{remark}[Mass terms]
\label{rem:v6v89-masses}
Adding a constant \(m_a^2\) to a block leaves \(A_0\) unchanged and replaces
\(A_1\) by \(A_1-m_a^2A_0\).  Unequal masses therefore add a volume-type
term to the first supertrace and generally destroy any massless
cancellation.  V89 keeps the massless primed test algebra so that field-
content and curvature effects are separated.
\end{remark}

\section{One-scalar, one-gauge, one-Dirac test}
\label{sec:v6v89-test-content}

\begin{theorem}[Leading cancellation and curvature remainder]
\label{thm:v6v89-test-supertrace}
For one real scalar, one gauge one-form with one complex scalar ghost, and
one complex two-component Dirac modulus, with the exponents listed above,
\begin{align}
 A_0^{\mathrm{sup}}&=0,
 \label{eq:v6v89-A0-cancel}\\
 A_1^{\mathrm{sup}}&={1\over4}\int R_gd\mu_g.
 \label{eq:v6v89-A1-remain}
\end{align}
Thus the proper-time determinant product has no volume divergence but has a
nonzero curvature divergence unless \(\int R_g=0\).
\end{theorem}

\begin{proof}
For the volume coefficient, the four contributions are
\begin{equation}
 -{1\over2},
 \qquad -{3\over2},
 \qquad +1,
 \qquad +1,
 \label{eq:v6v89-A0-ledger}
\end{equation}
whose sum is zero.  For the curvature coefficient, in units of
\(\int R_gd\mu_g\), they are
\begin{equation}
 -{1\over12},
 \qquad +{1\over4},
 \qquad +{1\over6},
 \qquad -{1\over12},
 \label{eq:v6v89-A1-ledger}
\end{equation}
whose sum is \(1/4\).  Insert these supertraces in the relative proper-time
formula of Theorem~\ref{thm:v6v88-heat-divergences}.
\end{proof}

\begin{corollary}[Gauge-plus-ghost block alone]
\label{cor:v6v89-gauge-ghost}
For one gauge generator,
\begin{align}
 A_0^{\mathrm{g+gh}}&=-{1\over2}\operatorname{Vol}(g),
 \label{eq:v6v89-gauge-A0}\\
 A_1^{\mathrm{g+gh}}&={5\over12}\int R_gd\mu_g.
 \label{eq:v6v89-gauge-A1}
\end{align}
The incomplete cancellation of \(A_0\) is the single physical bosonic
degree of freedom of a massless gauge field in three Euclidean dimensions,
expressed at determinant level.
\end{corollary}

\begin{proof}
Combine the one-form contribution with the scalar ghost contribution in
\eqref{eq:v6v89-A0-ledger}--\eqref{eq:v6v89-A1-ledger}.
\end{proof}

\section{Evaluation on the fixed-volume conformal branch}
\label{sec:v6v89-fixed-volume}

For \(g=\psi^4\bar g\) in the flat conformal class, V88 proved
\begin{equation}
 \int R_gd\mu_g=8\int|\bar\nabla\psi|^2d\mu_{\bar g}.
 \label{eq:v6v89-conformal-R}
\end{equation}

\begin{theorem}[Test-supertrace response]
\label{thm:v6v89-response}
On the V84 fixed-volume two-source branch, let
\begin{equation}
 A_{\epsilon,\beta}
 ={\psi_0\epsilon-\psi_0^{-7}\beta\over\mu_0+8}.
 \label{eq:v6v89-A}
\end{equation}
Then the test content of Theorem~\ref{thm:v6v89-test-supertrace} obeys
\begin{equation}
 A_1^{\mathrm{sup}}
 =A_{\epsilon,\beta}^2
 +O\bigl((|\epsilon|+|\beta|)^3\bigr).
 \label{eq:v6v89-supertrace-response}
\end{equation}
It vanishes through quadratic order on
\(\beta=\psi_0^8\epsilon\).
\end{theorem}

\begin{proof}
Theorem~\ref{thm:v6v88-response} gives
\(\int R_g=4A_{\epsilon,\beta}^2+O(3)\).  Multiply by the coefficient
\(1/4\) in \eqref{eq:v6v89-A1-remain}.
\end{proof}

\begin{theorem}[Order mismatch with the conformal coarea determinant]
\label{thm:v6v89-order-mismatch}
Near the homogeneous base point, the heat-supertrace curvature divergence
\eqref{eq:v6v89-supertrace-response} cannot cancel the conformal inverse-
coarea trace divergence of V86.  The former is quadratic in
\((\epsilon,\beta)\), whereas the latter has mean coefficient
\begin{equation}
 -4\epsilon-12\psi_0^{-8}\beta+O(2)
 \label{eq:v6v89-coarea-linear}
\end{equation}
in the logarithm of the reduced density.  On the first-order cancellation
line the heat curvature term vanishes through order two while the coarea
term equals \(-16\epsilon+O(\epsilon^2)\) times its linearly divergent
lattice sum.
\end{theorem}

\begin{proof}
Use Theorem~\ref{thm:v6v89-response} and
Theorem~\ref{thm:v6v86-mean-U}.  Analytic functions whose lowest nonzero
homogeneous orders differ cannot cancel in a neighbourhood of the origin.
The cancellation-line values are their stated specializations.
\end{proof}

\begin{assessment}[Rank balance is insufficient]
\label{ass:v6v89-rank-insufficient}
The zeroth heat coefficient counts determinant-weighted fibre ranks, and the
test field content cancels it exactly.  The next coefficient sees curvature
endomorphisms and does not cancel.  Even its cancellation would not remove
the conformal coarea trace because that Jacobian is not one of the covariant
field Hessians and begins at a different source order.  A continuum measure
requires a ledger of operator coefficients, statistics, and constraint
Jacobians, not a count of nominal bosonic and fermionic components.
\end{assessment}

\section{General multiplicity equations}
\label{sec:v6v89-multiplicities}

Let \(n_s,n_g,n_D\) be respectively the numbers of real scalars, gauge
generators with their ghosts, and complex rank-two Dirac moduli in this
three-dimensional test algebra.

\begin{proposition}[Two heat-supertrace equations]
\label{prop:v6v89-multiplicity-equations}
The massless coefficients are
\begin{align}
 {A_0^{\mathrm{sup}}\over\operatorname{Vol}(g)}
 &=-{n_s\over2}-{n_g\over2}+n_D,
 \label{eq:v6v89-general-A0}\\
 {A_1^{\mathrm{sup}}\over\int R_gd\mu_g}
 &=-{n_s\over12}+{5n_g\over12}-{n_D\over12}.
 \label{eq:v6v89-general-A1}
\end{align}
Both vanish exactly when
\begin{equation}
 n_D={n_s+n_g\over2},
 \qquad
 3n_g=n_s.
 \label{eq:v6v89-two-cancellations}
\end{equation}
Equivalently \((n_s,n_g,n_D)=(3r,r,2r)\) for some nonnegative integer
\(r\).
\end{proposition}

\begin{proof}
Sum the scalar, gauge-plus-ghost, and Dirac entries above.  Setting the first
line to zero gives \(2n_D=n_s+n_g\).  Substitution in twelve times the
second line gives \(-n_s+5n_g-n_D=0\), hence
\(-3n_s+9n_g=0\).  This is \(n_s=3n_g\), and then
\(n_D=2n_g\).
\end{proof}

\begin{firewall}[Not a Standard Model multiplicity test]
\label{fw:v6v89-not-SM}
Equations \eqref{eq:v6v89-two-cancellations} belong to massless spatial
three-dimensional test operators with Dirac phases omitted.  The physical
Standard Model is chiral in four spacetime dimensions, has symmetry
breaking, masses, nonminimal curvature terms under renormalization, and
gauge representations with different multiplicities.  No claim about
Standard Model cancellation follows by inserting informal field counts into
these equations.  The actual regulated Hessians must first be derived.
\end{firewall}

\begin{definition}[Minimal heat-supertrace card, MHSC]
\label{def:v6v89-MHSC}
MHSC requires:
\begin{enumerate}[label=\textup{(HS\arabic*)},leftmargin=5.2em]
\item one covariant regulator and primed-zero-mode convention;
\item the exact gauge-fixed Hessian and determinant exponent of each field;
\item computation of \(A_0,A_1\) and all higher coefficients divergent in
      the chosen spacetime or spatial dimension;
\item separation of determinant modulus, phase, and anomaly;
\item addition of constraint coarea and York Jacobians by their own source
      expansion; and
\item cancellation or declared local counterterms compatible with BRST and
      chart transitions.
\end{enumerate}
\end{definition}

\begin{theorem}[MHSC progress]
\label{thm:v6v89-MHSC}
V89 closes MHSC\textup{(HS1)}--\textup{(HS3)} for the stated massless test
content and exhibits the two multiplicity equations
\eqref{eq:v6v89-two-cancellations}.  The omitted Dirac phase leaves
MHSC\textup{(HS4)} open, and no full Standard Model Hessian or BRST
counterterm has closed MHSC\textup{(HS5)}--\textup{(HS6)}.
\end{theorem}

\begin{proof}
Use Lemma~\ref{lem:v6v89-coefficients},
Theorem~\ref{thm:v6v89-test-supertrace}, and
Proposition~\ref{prop:v6v89-multiplicity-equations}, subject to
Firewall~\ref{fw:v6v89-not-SM}.
\end{proof}

\begin{programme}[V90: compulsory audit and coarea/heat scheme matching]
\label{prog:v6v89-V90}
\begin{enumerate}[leftmargin=3em]
\item Perform the compulsory newest-YM/newest-NS audit by immutable digest.
\item Convert the conformal inverse-coarea trace subtraction from the sharp
      Fourier scheme to proper time on the same reference operator.
\item Prove scheme matching of its linear divergence and identify the finite
      local ambiguity.
\item Place that Jacobian beside the covariant heat supertrace without
      claiming cancellation across different source orders.
\item Choose the next attack between the Dirac phase/anomaly and the first
      nonlinear gravitational ghost coefficient after the audit.
\end{enumerate}
\end{programme}

\begin{firewall}[Claim boundary after V89]
\label{fw:v6v89-boundary}
V89 computes only a massless three-dimensional test supertrace.  It does not
derive the physical Standard Model or gravitational Hessians, include the
fermion determinant phase, establish anomaly cancellation or BRST
invariance, construct the global reduced measure or York evolution, remove
all regulators, or prove the full SM--Einstein continuum.
\end{firewall}

\begin{theorem}[Sixth-series V89 result]
\label{thm:v6v89-status}
For one real scalar, one gauge generator with scalar ghost, and one complex
rank-two Dirac modulus, the common proper-time supertrace cancels the volume
coefficient but leaves \((1/4)\int R_g\).  General test multiplicities cancel
both first coefficients exactly in the ratio \((n_s,n_g,n_D)=(3r,r,2r)\).
On the matched fixed-volume branch, the test curvature divergence is
quadratic in the uncancelled beat amplitude and cannot cancel the conformal
coarea trace, which begins linearly and survives the V84 cancellation line.
The next step is to put both subtractions in one proper-time scheme before
adding phases or nonlinear ghosts.
\end{theorem}

\begin{proof}
Use Theorems~\ref{thm:v6v89-test-supertrace},
\ref{thm:v6v89-response},
\ref{thm:v6v89-order-mismatch}, and
Proposition~\ref{prop:v6v89-multiplicity-equations}, subject to
Firewalls~\ref{fw:v6v89-not-SM} and \ref{fw:v6v89-boundary}.
\end{proof}

The compulsory YM/NS audit is due in V90.

% =============================================================================
% BEGIN APPENDED SECOND-SERIES V90 MATERIAL
% =============================================================================

\clearpage
\chapter{Second-series V90: companion audit and regulator matching for the
conformal coarea determinant}
\label{sec:v2v90}

\section{Checkpoint and scope}
\label{sec:newV90-entry}

V90 is a compulsory companion checkpoint.  It first records immutable
snapshots of the active Yang--Mills and Navier--Stokes notes, and then closes
one item left open by V86--V89: the comparison between the sharp Fourier
subtraction used for the conformal coarea determinant and a covariant
proper-time subtraction.  The comparison is made for the positive
fixed-principal-symbol operator already isolated in V86.  It does not identify
that scalar determinant with the complete gravitational Faddeev--Popov
factor.

\section{Companion audit at V90}
\label{sec:newV90-audit}

The companion files were read as byte snapshots.  The settled Yang--Mills
snapshot was
\begin{center}
\begin{tabular}{ll}
file &
\texttt{Renewal\_Yang\_Mills\_Mass\_Gap\_Research\_Notes\_V63.tex},\\
bytes & \(977\,631\),\\
logical source lines & \(30\,718\),\\
SHA--256 &
\texttt{3DE7432BFB9A747C38A7624519F6E83788B71476103FA8C4F8279B5593893F3C}.
\end{tabular}
\end{center}
The first read occurred while that file was being replaced and returned the
V62 body under a V63 filename.  No conclusion below uses that transitional
read.  A second read returned the V63 end marker, parent record, and the hash
displayed above.

The Navier--Stokes snapshot was
\begin{center}
\begin{tabular}{ll}
file &
\texttt{Renewal\_NS\_Stability\_Research\_Notes\_V15.tex},\\
bytes & \(209\,670\),\\
logical source lines & \(4\,134\),\\
SHA--256 &
\texttt{691BCE8817E34F0508C084336C6C7D4CB412801577489878DC9ADBF842EBADC3}.
\end{tabular}
\end{center}

\begin{audit}[Mathematical information imported at V90]
\label{audit:newV90-transfer}
The Yang--Mills V63 note proves an exact Poisson add-one covariance identity
for the small clock and uses it before estimates to recover the enhanced
\(s^3\) variance scale.  It closes the first marked-clock row but leaves the
iterated-difference card open.  The Navier--Stokes V15 note organizes rescaled
limits as invariant probability measures and reduces its homogeneous point
gate to a strict mean inequality; equality at the endpoint is exact, while the
strict subendpoint sign remains open.

The transferable rule is structural: form the connected or relative object
before differentiation and before removal of the regulator.  Here that means
forming the relative heat trace
\(\operatorname{Tr}(e^{-tL}-e^{-tL_0})\) before taking
\(t\downarrow0\).  No estimate, exponent, or numerical constant is imported
from either companion programme.
\end{audit}

\section{The translation-invariant reference operator}
\label{sec:newV90-reference}

Let
\[
 \mathbb T^3=(\mathbb R/2\pi\mathbb Z)^3
\]
carry normalized Haar measure.  The functions \(e_k(x)=e^{ik\cdot x}\),
\(k\in\mathbb Z^3\), are then an orthonormal Fourier basis.  Fix
\(\mu_0>0\) and put
\begin{equation}
 L_0=-8\Delta+\mu_0,
 \qquad
 L=L_0+U,                                                    \label{eq:newV90-L}
\end{equation}
where \(U\in C^\infty(\mathbb T^3;\mathbb R)\) acts by
multiplication.  Write
\begin{equation}
 \overline U=\int_{\mathbb T^3}U(x)\,dx,
 \qquad
 K=L_0^{-1/2}UL_0^{-1/2}.                                  \label{eq:newV90-K}
\end{equation}
The integral uses normalized Haar measure.

\begin{lemma}[Schatten class and positivity]
\label{lem:newV90-Schatten}
The operator \(K\) is self-adjoint and Hilbert--Schmidt.  If
\(\|K\|<1\), then
\[
 L_q=L_0+qU=L_0^{1/2}(I+qK)L_0^{1/2}>0,
 \qquad 0\le q\le1,
\]
and \(\det_2(I+K)>0\).
\end{lemma}

\begin{proof}
The operator \(L_0^{-1/2}\) is a pseudodifferential operator of order
\(-1\).  Hence \(K\) has order \(-2\).  On a three-dimensional compact
manifold an operator of order strictly less than \(-3/2\) is
Hilbert--Schmidt.  Self-adjointness follows from the real-valuedness of
\(U\).  The norm hypothesis gives
\(I+qK\ge(1-\|K\|)I\).  Its eigenvalues are positive, and the convergent
product defining the second Fredholm determinant is therefore positive.
\end{proof}

\section{The proper-time trace and its exact leading coefficient}
\label{sec:newV90-proper-time}

Define
\begin{equation}
 S_s^{\rm pt}
 :=\operatorname{Tr}\left(L_0^{-1}e^{-sL_0}\right)
 =\sum_{k\in\mathbb Z^3}
 {e^{-s(8|k|^2+\mu_0)}\over8|k|^2+\mu_0},
 \qquad s>0.                                                \label{eq:newV90-Spt}
\end{equation}

\begin{lemma}[Proper-time asymptotics]
\label{lem:newV90-Spt-asymptotic}
There is a finite constant \(C_{\rm pt}(\mu_0)\) such that
\begin{equation}
 S_s^{\rm pt}
 =c_{\rm pt}s^{-1/2}+C_{\rm pt}(\mu_0)+O(s^{1/2}),
 \qquad
 c_{\rm pt}={\pi^{3/2}\over8\sqrt2}.                       \label{eq:newV90-Spt-asymptotic}
\end{equation}
\end{lemma}

\begin{proof}
Functional calculus gives
\begin{equation}
 S_s^{\rm pt}=\int_s^\infty H_0(t)\,dt,
 \qquad
 H_0(t)=\operatorname{Tr}(e^{-tL_0}).                       \label{eq:newV90-Spt-integral}
\end{equation}
Poisson summation yields
\[
 H_0(t)
 =e^{-\mu_0t}\left({\pi\over8t}\right)^{3/2}
 \sum_{n\in\mathbb Z^3}e^{-\pi^2|n|^2/(8t)}.
\]
Consequently, as \(t\downarrow0\),
\[
 H_0(t)
 =\left({\pi\over8}\right)^{3/2}t^{-3/2}
 -\mu_0\left({\pi\over8}\right)^{3/2}t^{-1/2}
 +O(t^{1/2})+O(t^{-3/2}e^{-c/t}).
\]
The difference between \(H_0(t)\) and its first term is integrable
at zero.  Integrating the first term in
\eqref{eq:newV90-Spt-integral} gives
\[
 2\left({\pi\over8}\right)^{3/2}s^{-1/2}
 ={\pi^{3/2}\over8\sqrt2}s^{-1/2}.
\]
The exponentially decaying large-time tail is integrable because
\(\mu_0>0\).  These observations prove both the existence of the finite
constant and the stated remainder.
\end{proof}

\begin{lemma}[A multiplication insertion sees only the mean]
\label{lem:newV90-mean-insertion}
For every \(s>0\),
\begin{equation}
 \operatorname{Tr}\left(UL_0^{-1}e^{-sL_0}\right)
 =\overline U\,S_s^{\rm pt}.                                \label{eq:newV90-mean-insertion}
\end{equation}
\end{lemma}

\begin{proof}
The integral kernel of \(L_0^{-1}e^{-sL_0}\) is translation invariant.
Its diagonal is therefore constant and equals its trace with respect to
normalized Haar measure.  Multiplication by \(U\) and integration of the
diagonal give \eqref{eq:newV90-mean-insertion}.  Equivalently, in the
Fourier basis every diagonal matrix element of \(U\) is \(\overline U\).
\end{proof}

\section{Relative proper-time determinant}
\label{sec:newV90-relative}

For \(s>0\), define the regulated inverse relative determinant by
\begin{equation}
 \log {\mathcal R}_s(U)
 :=\int_s^\infty {1\over t}
 \operatorname{Tr}\left(e^{-tL}-e^{-tL_0}\right)\,dt.       \label{eq:newV90-relative}
\end{equation}
The integral is finite for fixed \(s\): its lower endpoint is cut off,
and both positive operators have an exponential large-time bound.

\begin{theorem}[Exact trace subtraction]
\label{thm:newV90-exact-subtraction}
Assume \(\|K\|<1\).  Then
\begin{equation}
 \log {\mathcal R}_s(U)
 =-\overline U S_s^{\rm pt}
   -\log\det{}_2(I+K)+o(1)\qquad(s\downarrow0).              \label{eq:newV90-relative-expansion}
\end{equation}
In particular,
\begin{equation}
 \lim_{s\downarrow0}
 e^{\overline U S_s^{\rm pt}}{\mathcal R}_s(U)
 ={1\over\det{}_2(I+K)}>0.                                 \label{eq:newV90-exact-limit}
\end{equation}
\end{theorem}

\begin{proof}
Introduce \(L_q=L_0+qU\) and let \(r_s(q)\) be the right side of
\eqref{eq:newV90-relative} with \(L\) replaced by \(L_q\).  Duhamel's
formula and cyclicity of the trace give
\begin{equation}
 r_s'(q)
 =-\int_s^\infty\operatorname{Tr}(Ue^{-tL_q})\,dt
 =-\operatorname{Tr}(UL_q^{-1}e^{-sL_q}).                  \label{eq:newV90-rprime}
\end{equation}
At \(q=0\), \Cref{lem:newV90-mean-insertion} gives
\begin{equation}
 r_s'(0)=-\overline U S_s^{\rm pt}.                         \label{eq:newV90-rprime-zero}
\end{equation}

For \(0\le q\le1\), the resolvent identity gives
\[
 L_q^{-1}-L_0^{-1}=-qL_q^{-1}UL_0^{-1}.
\]
After multiplication by \(U\), the right side has order \(-4\) and is
trace class in dimension three.  Positivity from
\Cref{lem:newV90-Schatten} makes the trace-class bounds uniform in
\(q\).  Therefore the heat cutoffs in the difference
\(r_s'(q)-r_s'(0)\) may be removed in trace norm, and
\begin{align}
 \lim_{s\downarrow0}\{r_s'(q)-r_s'(0)\}
 &=-\operatorname{Tr}\{UL_q^{-1}-UL_0^{-1}\}               \notag\\
 &=-\operatorname{Tr}\{(I+qK)^{-1}K-K\}.                  \label{eq:newV90-rprime-limit}
\end{align}
The last equality follows by cyclicity applied only to the trace-class
difference.  The right side is the negative derivative of
\(\log\det_2(I+qK)\).  Uniform domination permits integration over
\(q\in[0,1]\).  Since \(r_s(0)=0\) and \(\det_2(I)=1\),
\[
 \lim_{s\downarrow0}
 \{r_s(1)-r_s'(0)\}=-\log\det_2(I+K).
\]
Substitution of \eqref{eq:newV90-rprime-zero} proves
\eqref{eq:newV90-relative-expansion} and
\eqref{eq:newV90-exact-limit}.
\end{proof}

\begin{corollary}[Leading heat subtraction and finite local normalization]
\label{cor:newV90-leading-subtraction}
Under the hypotheses of \Cref{thm:newV90-exact-subtraction},
\begin{equation}
 \lim_{s\downarrow0}
 \exp\left(c_{\rm pt}\overline U s^{-1/2}\right)
 {\mathcal R}_s(U)
 =
 {e^{-C_{\rm pt}(\mu_0)\overline U}\over\det_2(I+K)}.       \label{eq:newV90-leading-limit}
\end{equation}
Thus exact trace subtraction and leading heat subtraction differ by the
finite local factor
\(e^{-C_{\rm pt}(\mu_0)\overline U}\).
\end{corollary}

\begin{proof}
Insert \eqref{eq:newV90-Spt-asymptotic} into
\eqref{eq:newV90-relative-expansion} and exponentiate.
\end{proof}

\section{Comparison with the sharp Fourier scheme}
\label{sec:newV90-sharp}

Let \(P_\Lambda\) project onto the Fourier modes
\(|k|\le\Lambda\) and set
\begin{equation}
 S_\Lambda^{\rm sh}
 :=\sum_{|k|\le\Lambda}{1\over8|k|^2+\mu_0}.                \label{eq:newV90-Ssharp}
\end{equation}

\begin{lemma}[Sharp leading coefficient]
\label{lem:newV90-sharp-leading}
One has
\begin{equation}
 {S_\Lambda^{\rm sh}\over\Lambda}\longrightarrow
 c_{\rm sh}:={\pi\over2}.                                  \label{eq:newV90-sharp-leading}
\end{equation}
\end{lemma}

\begin{proof}
After the change of variables \(k=\Lambda\xi\), the nonzero-mode
part is a Riemann sum for
\[
 \int_{|\xi|\le1}{d^3\xi\over8|\xi|^2}={\pi\over2}.
\]
To justify the integrable singularity, remove a ball of radius
\(\delta\), apply ordinary Riemann-sum convergence on its complement,
and bound both the integral and the rescaled lattice sum inside the ball
by \(C\delta\), uniformly in \(\Lambda\).  The zero mode contributes
\((\mu_0\Lambda)^{-1}\) after division by \(\Lambda\), and hence
vanishes.
\end{proof}

\begin{proposition}[Exact sharp and proper-time subtractions agree]
\label{prop:newV90-scheme-agreement}
Let
\(K_\Lambda=P_\Lambda K P_\Lambda\).  Then
\begin{equation}
 \lim_{\Lambda\to\infty}
 {\exp(\overline U S_\Lambda^{\rm sh})
  \det(I_{\Lambda})
  \over
  \det(I_{\Lambda}+K_\Lambda)}
 ={1\over\det_2(I+K)},                                    \label{eq:newV90-sharp-det2}
\end{equation}
where the harmless \(\det(I_\Lambda)=1\) records the reference
normalization.  Consequently the exact sharp trace subtraction in V86 and
the exact proper-time trace subtraction
\eqref{eq:newV90-exact-limit} define the same positive
\(\det_2\)-relative factor.
\end{proposition}

\begin{proof}
In the Fourier basis,
\[
 \operatorname{Tr}K_\Lambda
 =\overline U S_\Lambda^{\rm sh}.
\]
The finite-dimensional identity
\[
 \det{}_2(I_\Lambda+K_\Lambda)
 =\det(I_\Lambda+K_\Lambda)e^{-\operatorname{Tr}K_\Lambda}
\]
turns the left side of \eqref{eq:newV90-sharp-det2} into
\(\det_2(I_\Lambda+K_\Lambda)^{-1}\).  Since
\(P_\Lambda K P_\Lambda\to K\) in Hilbert--Schmidt norm and
\(\det_2\) is continuous in that norm, the asserted limit follows.
The proper-time equality is \eqref{eq:newV90-exact-limit}.
\end{proof}

\begin{corollary}[Leading-cutoff matching]
\label{cor:newV90-leading-match}
If
\begin{equation}
 s={\pi\over32}\Lambda^{-2},                                \label{eq:newV90-cutoff-map}
\end{equation}
then
\[
 c_{\rm pt}s^{-1/2}=c_{\rm sh}\Lambda.
\]
This matches the leading linear divergence.  It does not assert that
\(S_\Lambda^{\rm sh}-c_{\rm sh}\Lambda\) has a limit.
\end{corollary}

\begin{proof}
By \Cref{lem:newV90-Spt-asymptotic,lem:newV90-sharp-leading},
\[
 \left({c_{\rm pt}\over c_{\rm sh}}\right)^2
 =\left({\pi^{3/2}/(8\sqrt2)\over\pi/2}\right)^2
 ={\pi\over32}.
\]
The displayed identity follows directly from
\eqref{eq:newV90-cutoff-map}.
\end{proof}

\begin{assessment}[What is universal]
\label{ass:newV90-universal}
The divergent coefficient multiplying the source is always the spatial
mean \(\overline U\).  The numerical relation between a length cutoff and a
proper-time cutoff depends on the regulator.  Once the complete trace
\(\overline U S\) is subtracted, both schemes leave the same nonlocal,
positive relative factor \(\det_2(I+K)^{-1}\).  Retaining only the leading
heat coefficient leaves a finite local normalization, explicitly shown in
\eqref{eq:newV90-leading-limit}.
\end{assessment}

\section{Insertion of the fixed-volume beat branch}
\label{sec:newV90-beat}

V86 computed the mean perturbing potential on the fixed-volume branch as
\begin{equation}
 \overline U
 =4\varepsilon+12\psi_0^{-8}\beta
 +O((|\varepsilon|+|\beta|)^2).                             \label{eq:newV90-Umean}
\end{equation}
The first-order beat-cancellation line from V84 is
\(\beta=\psi_0^8\varepsilon\).

\begin{theorem}[Regulator-independent source direction]
\label{thm:newV90-source-direction}
For the conformal coarea block, both exact subtraction schemes remove a
linear divergence in the source direction
\[
 4\varepsilon+12\psi_0^{-8}\beta.
\]
On the first-order beat-cancellation line this direction is
\begin{equation}
 \overline U=16\varepsilon+O(\varepsilon^2).                \label{eq:newV90-line-mean}
\end{equation}
Hence cancellation of the first metric beat does not cancel the conformal
coarea trace.
\end{theorem}

\begin{proof}
Equation \eqref{eq:newV90-Umean} is the V86 expansion.  The sharp and
proper-time first trace insertions are respectively
\(\overline U S_\Lambda^{\rm sh}\) and
\(\overline U S_s^{\rm pt}\), by the Fourier diagonal calculation and
\Cref{lem:newV90-mean-insertion}.  Substitution of
\(\beta=\psi_0^8\varepsilon\) into
\eqref{eq:newV90-Umean} gives \eqref{eq:newV90-line-mean}.
\end{proof}

\begin{corollary}[Order mismatch with the tested matter heat blocks]
\label{cor:newV90-order-mismatch}
On the beat-cancellation line, the conformal coarea counterterm begins at
order \(\varepsilon\).  The fixed-volume scalar curvature term computed in
V88 and the massless test supertrace computed in V89 begin at order
\(\varepsilon^2\).  Those tested heat blocks therefore cannot cancel the
linear conformal coarea counterterm.
\end{corollary}

\begin{proof}
The linear statement is \eqref{eq:newV90-line-mean}.  V88 gives
\(\int R_g\,d\operatorname{vol}_g=O(\varepsilon^2)\) on this line, and
V89 gives the corresponding test supertrace coefficient at the same order.
Terms of different homogeneous degree cannot cancel in a small-source
expansion.
\end{proof}

\begin{definition}[Scheme-matched conformal coarea card]
\label{def:newV90-SMCCC}
The scheme-matched conformal coarea card, SMCCC, consists of the following
data on a small positive branch:
\begin{enumerate}[label=\textup{(\roman*)},leftmargin=3.5em]
\item the reference operator \(L_0\), perturbation \(U\), and
      \(K=L_0^{-1/2}UL_0^{-1/2}\in\mathfrak S_2\);
\item positivity of \(I+K\);
\item the sharp trace \(\overline U S_\Lambda^{\rm sh}\) and the
      proper-time trace \(\overline U S_s^{\rm pt}\); and
\item the common renormalized value \(\det_2(I+K)^{-1}\) after exact
      trace subtraction.
\end{enumerate}
\end{definition}

\begin{theorem}[SMCCC is closed on the perturbative beat branch]
\label{thm:newV90-SMCCC}
If the beat-branch parameters are sufficiently small that \(\|K\|<1\),
then SMCCC holds.  Its value is finite, strictly positive, continuous in
the branch parameters in Hilbert--Schmidt topology, and independent of the
choice between exact sharp and exact proper-time trace subtraction.
\end{theorem}

\begin{proof}
Use \Cref{lem:newV90-Schatten,thm:newV90-exact-subtraction,
prop:newV90-scheme-agreement}.  Continuity follows from continuity of
\(\det_2\) on \(\mathfrak S_2\).  Strict positivity follows because
\(I+K>0\).
\end{proof}

\section{Next obstruction: the fermionic phase}
\label{sec:newV90-next}

V89 retained only the modulus of a three-dimensional Dirac determinant.
That is insufficient for a chiral Standard Model sector.  The next note will
separate the modulus from the phase at finite spectral cutoff, identify the
spectral-flow obstruction to a globally continuous phase, and state the
precise eta-invariant or four-dimensional anomaly-inflow datum needed for a
cutoff-independent choice.  This is upstream of any claimed Standard Model
measure: without a coherent fermionic phase, determinant-modulus cancellation
does not define the chiral functional integral.

\begin{firewall}[Claim boundary after V90]
\label{fire:newV90-boundary}
V90 performs the compulsory YM/NS audit, proves the exact proper-time
asymptotics of the reference conformal trace, proves that exact proper-time
trace subtraction converges to the inverse second Fredholm determinant, and
matches that value to the exact sharp Fourier subtraction.  It identifies the
finite local normalization left by leading-only heat subtraction and proves
that the fixed-volume beat-cancellation line does not cancel the linear
conformal coarea counterterm.

V90 does not derive the complete gravitational ghost operator, a BRST
complex, a Standard Model gauge-fixed measure, the fermionic determinant
phase, anomaly cancellation, reflection positivity, Lorentzian reconstruction,
Einstein evolution, continuum Yang--Mills existence or mass gap, Navier--Stokes
regularity, or the full SM--Einstein continuum.  The finite local factor in
\eqref{eq:newV90-leading-limit} is a renormalization choice, not a value fixed
by the present model.
\end{firewall}

\begin{theorem}[V90 result ledger]
\label{thm:newV90-results}
V90:
\begin{enumerate}[label=\textup{(V90.\arabic*)},leftmargin=6.5em]
\item records settled, hashed YM V63 and NS V15 companion snapshots;
\item transfers only their assemble-before-limiting structural rule;
\item computes the exact leading proper-time trace coefficient;
\item proves that a multiplication insertion depends only on
      \(\overline U\);
\item proves the exact proper-time subtraction theorem;
\item identifies its common \(\det_2\) value with sharp Fourier
      subtraction;
\item computes the leading cutoff conversion \(s=(\pi/32)\Lambda^{-2}\);
\item isolates the finite local ambiguity of leading-only subtraction;
\item proves the noncancellation of the linear coarea trace on the beat line;
      and
\item closes SMCCC on the small positive branch.
\end{enumerate}
\end{theorem}

\begin{proof}
Use \Cref{audit:newV90-transfer,
lem:newV90-Spt-asymptotic,
lem:newV90-mean-insertion,
thm:newV90-exact-subtraction,
cor:newV90-leading-subtraction,
prop:newV90-scheme-agreement,
cor:newV90-leading-match,
thm:newV90-source-direction,
cor:newV90-order-mismatch,
thm:newV90-SMCCC}.
\end{proof}

\paragraph{Parent-version record.}
V90 was formed from
\texttt{Renewal\_SM\_Einstein\_Continuum\_Research\_Notes\_V89.tex}.
The V89 parent had \(30\,330\) logical source lines,
\(1\,165\,109\) bytes, and SHA--256
\[
 \texttt{B6571A30FD005D1DDDE8D5663452F8B5B3F165A2FC0A8928713AD5942CD82C23}.
\]
The next compulsory companion audit is V95.

% =============================================================================
% END APPENDED SECOND-SERIES V90 MATERIAL
% =============================================================================


% =============================================================================
% BEGIN APPENDED SECOND-SERIES V91 MATERIAL
% =============================================================================

\clearpage
\chapter{Second-series V91: eta phases, spectral-flow transport, and the
mapping-torus anomaly gate}
\label{sec:v2v91}

\section{Nonduplication boundary and analytic input}
\label{sec:newV91-entry}

V18 already proves local complex trace-log estimates on simply connected
high-gap charts.  V19 already proves the abstract determinant-line patching
criterion, and V38 already verifies the usual perturbative Standard Model
anomaly sums and the ordinary mod-two \(SU(2)\) count.  The Grand Tensor
phase panel also contains the finite-dimensional determinant-sign formula
in terms of spectral-flow parity.  None of those results identifies the
regularized phase of the three-dimensional Dirac modulus used in V89, or
turns the abstract V19 holonomy into the spectral invariant of a mapping
torus.

V91 supplies that missing bridge.  Two established analytic results are
used with their hypotheses stated explicitly: regularity of the eta function
of a Dirac-type operator at the origin, and the Atiyah--Patodi--Singer and
Dai--Freed index/holonomy theorems.  The algebraic determinant formulas and
all consequences drawn from those theorems are proved below.  No new proof
of the full APS or Dai--Freed theorem is claimed.

\section{Finite spectral cutoff and the missing rank phase}
\label{sec:newV91-finite}

Let \(D\) be an invertible Hermitian matrix.  Write
\[
 n_+(D)=\#\{\lambda>0\},
 \qquad
 n_-(D)=\#\{\lambda<0\},
\]
with algebraic multiplicity, and set
\begin{equation}
 N(D)=n_+(D)+n_-(D),
 \qquad
 \eta_{\rm fin}(D)=n_+(D)-n_-(D).                          \label{eq:newV91-finite-eta}
\end{equation}

\begin{lemma}[Finite determinant phase]
\label{lem:newV91-finite-phase}
For every invertible Hermitian matrix,
\begin{align}
 {\det D\over|\det D|}
 &=\exp\left\{{i\pi\over2}
       \bigl(N(D)-\eta_{\rm fin}(D)\bigr)\right\},          \label{eq:newV91-sign}\
 e^{-i\pi N(D)/2}{\det D\over|\det D|}
 &=e^{-i\pi\eta_{\rm fin}(D)/2}.                           \label{eq:newV91-rank-renormalized}
\end{align}
\end{lemma}

\begin{proof}
The determinant has sign \((-1)^{n_-(D)}\).  Solving
\eqref{eq:newV91-finite-eta} gives
\(n_-=(N-\eta_{\rm fin})/2\), which proves
\eqref{eq:newV91-sign}.  Multiplication by the displayed rank phase gives
\eqref{eq:newV91-rank-renormalized}.
\end{proof}

\begin{assessment}[What V89 omitted]
\label{ass:newV91-rank-phase}
The V89 factor \(\det(D^2)^{1/2}\) is the modulus.  Even at finite cutoff it
does not contain either the sign in \eqref{eq:newV91-sign} or the rank phase
needed to compare different cutoffs.  The rank-renormalized expression
\eqref{eq:newV91-rank-renormalized} is the finite spectral precursor of the
eta phase.
\end{assessment}

\section{Zeta determinant and eta invariant on a closed three-manifold}
\label{sec:newV91-zeta}

Let \(Y\) be a closed oriented Riemannian spin three-manifold and let
\(D\) be an invertible self-adjoint Dirac-type operator on a finite-rank
Hermitian Clifford module.  For \(\operatorname{Re}z\) large, define
\begin{align}
 \zeta_+(z)&=\sum_{\lambda>0}\lambda^{-z},
 &
 \zeta_-(z)&=\sum_{\lambda<0}|\lambda|^{-z},                 \label{eq:newV91-half-zetas}\
 \eta_D(z)&=\zeta_+(z)-\zeta_-(z),
 &
 \zeta_{|D|}(z)&=\zeta_+(z)+\zeta_-(z).                    \label{eq:newV91-eta-zeta}
\end{align}
Choose the upper spectral prescription
\begin{equation}
 \zeta_D^\uparrow(z)
 :=\zeta_+(z)+e^{-i\pi z}\zeta_-(z),                        \label{eq:newV91-upper-cut}
\end{equation}
so a negative eigenvalue has logarithm
\(\log|\lambda|+i\pi\).  Put
\begin{equation}
 \det{}^\uparrow_\zeta D
 :=\exp\{-({\zeta_D^\uparrow})'(0)\}.                       \label{eq:newV91-zeta-det}
\end{equation}

\begin{theorem}[Eta completion of the Dirac modulus]
\label{thm:newV91-eta-completion}
The functions in \eqref{eq:newV91-half-zetas}--
\eqref{eq:newV91-upper-cut} admit the standard meromorphic continuations
which are regular at zero, and
\begin{equation}
 \det{}^\uparrow_\zeta D
 =\det{}_{\zeta}(D^2)^{1/2}
   \exp\left\{{i\pi\over2}
   \bigl(\zeta_{D^2}(0)-\eta_D(0)\bigr)\right\}.            \label{eq:newV91-general-phase}
\end{equation}
On a closed odd-dimensional manifold with \(D\) invertible,
\(\zeta_{D^2}(0)=0\).  Hence in the present three-dimensional setting
\begin{equation}
 \boxed{
 \det{}^\uparrow_\zeta D
 =\det{}_{\zeta}(D^2)^{1/2}e^{-i\pi\eta_D(0)/2}.}            \label{eq:newV91-eta-phase}
\end{equation}
\end{theorem}

\begin{proof}
From \eqref{eq:newV91-upper-cut},
\[
 -({\zeta_D^\uparrow})'(0)
 =-\zeta_{|D|}'(0)+i\pi\zeta_-(0).
\]
Moreover
\[
 \zeta_-(0)
 ={1\over2}\{\zeta_{|D|}(0)-\eta_D(0)\},
 \qquad
 \zeta_{|D|}(z)=\zeta_{D^2}(z/2).
\]
It follows that
\(-\zeta_{|D|}'(0)=-\tfrac12\zeta_{D^2}'(0)\), proving
\eqref{eq:newV91-general-phase}.

For a Laplace-type operator on a closed manifold the heat expansion has
terms \(t^{j-\dim Y/2}\) with integral \(j\ge0\).  When \(\dim Y\) is odd,
there is no \(t^0\) term.  The Mellin formula therefore gives
\(\zeta_{D^2}(0)=-\dim\ker D^2\), which vanishes under invertibility.
Substitution proves \eqref{eq:newV91-eta-phase}.
\end{proof}

\begin{corollary}[The two spectral cuts differ by a nonlocal phase]
\label{cor:newV91-cut-ambiguity}
If the lower prescription uses
\(\zeta_D^\downarrow(z)=\zeta_+(z)+e^{i\pi z}\zeta_-(z)\), then
\begin{equation}
 {\det{}^\uparrow_\zeta D\over
  \det{}^\downarrow_\zeta D}
 =e^{-i\pi\eta_D(0)}.                                      \label{eq:newV91-cut-ratio}
\end{equation}
Thus the odd-dimensional cut ambiguity is governed by the eta invariant,
not by the positive heat determinant \(\det_\zeta D^2\).
\end{corollary}

\begin{proof}
Repeat the proof of \Cref{thm:newV91-eta-completion} with the opposite sign
of \(i\pi\zeta_-(0)\), and use \(\zeta_{D^2}(0)=0\).
\end{proof}

\begin{warning}[A sharp eta count is not automatically convergent]
\label{warn:newV91-sharp-eta}
For the spectral projector
\(P_\Lambda=\mathbf1_{[-\Lambda,\Lambda]}(D)\), the finite expression in
\eqref{eq:newV91-rank-renormalized} contains
\(\eta_\Lambda=n_+(P_\Lambda D)-n_-(P_\Lambda D)\).
The analytic value \(\eta_D(0)\) is not asserted to be the ordinary limit of
\(\eta_\Lambda\).  The eta function is the regulator prescription, including
the local subtractions implicit in meromorphic continuation.
\end{warning}

\section{Zero crossings and the square-root datum}
\label{sec:newV91-crossing}

For a possibly singular self-adjoint Dirac operator define the reduced eta
invariant
\begin{equation}
 \xi(D):={\eta_D(0)+h(D)\over2},
 \qquad
 h(D)=\dim\ker D,                                           \label{eq:newV91-xi}
\end{equation}
and its exponentiation
\begin{equation}
 \tau(D):=e^{2\pi i\xi(D)}.                                \label{eq:newV91-tau}
\end{equation}

\begin{lemma}[Eta phase as a square root of the inverse xi phase]
\label{lem:newV91-square-root}
If \(D\) is invertible and
\(\Phi(D)=e^{-i\pi\eta_D(0)/2}\), then
\begin{equation}
 \Phi(D)^2=\tau(D)^{-1}.                                   \label{eq:newV91-square-root}
\end{equation}
The choice of \(\Phi\) therefore contains a square-root sign which is not
specified by \(\tau\) alone.
\end{lemma}

\begin{proof}
Invertibility gives \(h(D)=0\) and
\(2\pi i\xi(D)=i\pi\eta_D(0)\).  Squaring \(\Phi\) proves the identity.
\end{proof}

Let \(D_t\), \(0\le t\le1\), be a \(C^1\) path of self-adjoint
Dirac-type operators with invertible endpoints and only regular crossings.
Use the convention that a negative-to-positive crossing contributes
\(+1\) to \(\operatorname{sf}(D_t)\).

\begin{proposition}[Crossing transport of the eta square root]
\label{prop:newV91-crossing-phase}
Across the crossings of \(D_t\), the phase square root changes by
\begin{equation}
 \Phi(D_1)=(-1)^{\operatorname{sf}(D_t)}
 \Phi(D_0)                                                   \label{eq:newV91-crossing-phase}
\end{equation}
after the smooth local variation between crossings is removed.  At an
isolated simple crossing, the determinant vanishes and there is no
continuous nowhere-zero scalar phase through that point.
\end{proposition}

\begin{proof}
A regular negative-to-positive crossing changes the eta count by \(+2\);
the reverse crossing changes it by \(-2\).  In either case
\(e^{-i\pi\eta/2}\) acquires one factor \(-1\).  Multiplication over all
crossings proves \eqref{eq:newV91-crossing-phase}.  At a simple crossing one
eigenvalue is zero, so the determinant itself vanishes.  A nowhere-zero
scalar phase is undefined there.  This refines the earlier determinant-sign
spectral-flow formula by identifying the eta square root which carries that
sign.
\end{proof}

\begin{assessment}[Why exact low-mode transfer remains compulsory]
\label{ass:newV91-low-transfer}
The jump in \Cref{prop:newV91-crossing-phase} cannot be repaired by a
continuous branch of the high determinant logarithm.  The crossing eigenspace
must enter the retained Grassmann block, together with its determinant-line
orientation, before the high block loses its gap.  This is exactly the low-mode
rule in V18--V19 and the exact Feshbach transfer in V38--V39.
\end{assessment}

\section{APS cylinder transport}
\label{sec:newV91-APS}

Suppose the family \(D_t\) is the boundary family of a Dirac operator on
the cylinder \(X=[0,1]\times Y\), with product form near the boundary.
Let
\begin{equation}
 {\mathscr D}=\partial_t+D_t                              \label{eq:newV91-cylinder-operator}
\end{equation}
carry APS boundary conditions.  Fix signs by defining
\(\operatorname{sf}(D_t)=\operatorname{ind}{\mathscr D}_{\rm APS}\).
The APS index theorem then reads
\begin{equation}
 \operatorname{sf}(D_t)
 =\int_X\left[\widehat A(TX)\operatorname{ch}(E)\right]_4
  -\xi(D_1)+\xi(D_0).                                      \label{eq:newV91-APS}
\end{equation}
Write
\begin{equation}
 I_4(X,E):=int_X
 \left[\widehat A(TX)\operatorname{ch}(E)\right]_4.         \label{eq:newV91-I4}
\end{equation}

\begin{theorem}[Exact cylinder phase transport]
\label{thm:newV91-cylinder-phase}
If the endpoint operators are invertible, then
\begin{align}
 {\Phi(D_1)\over\Phi(D_0)}
 &=(-1)^{\operatorname{sf}(D_t)}e^{-i\pi I_4(X,E)},          \label{eq:newV91-cylinder-phase}\
 {\Phi(D_1)^2\over\Phi(D_0)^2}
 &=e^{-2\pi i I_4(X,E)}.                                   \label{eq:newV91-cylinder-square}
\end{align}
Thus the APS density controls the smooth square of the phase transport,
while spectral-flow parity fixes its missing sign.
\end{theorem}

\begin{proof}
With invertible endpoints, \(\xi(D_j)=\eta_{D_j}(0)/2\).  Rearranging
\eqref{eq:newV91-APS} gives
\[
 \xi(D_1)-\xi(D_0)=I_4(X,E)-\operatorname{sf}(D_t).
\]
Substitute this identity into
\(\Phi(D_1)/\Phi(D_0)=
e^{-i\pi(\xi(D_1)-\xi(D_0))}\).  Since spectral flow is integral,
the first formula follows.  Squaring removes its parity sign and proves the
second.
\end{proof}

\begin{corollary}[Multiplet addition law]
\label{cor:newV91-multiplet-cylinder}
For Dirac families in representations \(R_s\) with multiplicities and
chirality signs absorbed into integers \(m_s\), the product phase obeys
\begin{equation}
 {\Phi_{\rm tot}(1)\over\Phi_{\rm tot}(0)}
 =(-1)^{\sum_s m_s\operatorname{sf}(D_{s,t})}
  \exp\left\{-i\pi\sum_s m_s I_4(X,E_s)\right\}.            \label{eq:newV91-multiplet-phase}
\end{equation}
\end{corollary}

\begin{proof}
Raise \eqref{eq:newV91-cylinder-phase} to the powers \(m_s\) and multiply.
Both spectral flow and the local index density are additive under direct sum
and virtual difference.
\end{proof}

\section{The five-dimensional mapping-torus gate}
\label{sec:newV91-mapping-torus}

The preceding cylinder concerns the phase of a self-adjoint
three-dimensional time-slice operator.  The global anomaly of a
four-dimensional chiral determinant line is one dimension higher.
Let \(Z\) be a closed spin four-manifold and let
\(D_z^+:H_z^+\to H_z^-\) be a smooth family of chiral Dirac operators over a
parameter space \(\mathcal P\).  For a loop
\(\gamma:S^1\to\mathcal P\), let \(M_\gamma\) be the associated
five-dimensional mapping torus and let \(D_{M_\gamma,\epsilon}\) be its
self-adjoint Dirac operator with the base-circle metric scaled by
\(\epsilon^{-2}\).

\begin{theorem}[Dai--Freed holonomy formula]
\label{thm:newV91-DF-holonomy}
For the inverse determinant line with its Bismut--Freed connection and the
nonbounding spin structure on the base circle,
\begin{equation}
 \operatorname{Hol}_{\det^{-1}D^+}(\gamma)
 =(-1)^{\operatorname{ind}D_{\gamma(0)}^+}
  \lim_{\epsilon\downarrow0}
  \exp\{2\pi i\xi(D_{M_\gamma,\epsilon})\}.                 \label{eq:newV91-DF-holonomy}
\end{equation}
The adiabatic limit exists.  With the bounding spin structure the displayed
index sign is absent.
\end{theorem}

\begin{proof}
This is the Dai--Freed holonomy theorem, including the graded determinant-line
sign convention.  Its proof constructs parallel transport from the
adiabatic limit of the exponentiated xi invariant, proves the gluing law for
cylinders, and identifies the resulting connection with the Bismut--Freed
connection by their common variational formula.  We use that established
theorem as analytic input; the formula fixes exactly which five-dimensional
spectral datum the V19 abstract holonomy requires.
\end{proof}

For a chiral multiplet indexed by \(s\), define
\begin{equation}
 \mathfrak H_{\rm tot}(\gamma)
 :=\prod_s
 \left[
 (-1)^{\operatorname{ind}D_{s,\gamma(0)}^+}
 \lim_{\epsilon\downarrow0}
 e^{2\pi i\xi(D_{s,M_\gamma,\epsilon})}
 \right]^{m_s}.                                             \label{eq:newV91-total-holonomy}
\end{equation}

\begin{theorem}[Computable form of the V19 global gate]
\label{thm:newV91-global-gate}
On a connected good component of configuration space, the tensor-product
determinant line admits a parallel nowhere-zero section if and only if
\begin{enumerate}[label=\textup{(\roman*)},leftmargin=3.5em]
\item its Bismut--Freed curvature vanishes; and
\item \(\mathfrak H_{\rm tot}(\gamma)=1\) for every loop \(\gamma\) in
      that component.
\end{enumerate}
When these conditions hold, a fibre value at one base point determines the
parallel section uniquely.  This topological section is locality compatible
only if it also satisfies the rooted bounds required in V19.
\end{theorem}

\begin{proof}
The V19 flatness-and-holonomy lemma gives the equivalence between a parallel
nowhere-zero section and vanishing curvature plus trivial loop holonomy.
\Cref{thm:newV91-DF-holonomy} identifies the holonomy of each inverse
determinant line; tensor products multiply holonomies, giving
\eqref{eq:newV91-total-holonomy}.  The final sentence is the V19 distinction
between topological triviality and a cutoff-uniform local trivialization.
\end{proof}

\begin{corollary}[Exact residual Standard Model anomaly datum]
\label{cor:newV91-SM-residual}
For the one-generation Standard Model representation treated in V38, the
usual local gauge and mixed gauge--gravitational curvature coefficients
vanish, and the ordinary \(SU(2)\) mod-two loop has even doublet multiplicity.
After those established cancellations, the remaining global phase question
on a chosen Einstein background and chosen global form of the gauge group is
exactly
\begin{equation}
 \mathfrak H_{\rm SM}(\gamma)=1
 \quad\hbox{for every admissible mapping-torus loop }\gamma. \label{eq:newV91-SM-global-gate}
\end{equation}
V38 does not by itself prove \eqref{eq:newV91-SM-global-gate} for all
topologies or all central quotients of
\(SU(3)\times SU(2)\times U(1)\).
\end{corollary}

\begin{proof}
The local statement and ordinary mod-two count are
Theorem~\ref{thm:v6v38-SM-anomaly-cancellation} and
Corollary~\ref{cor:v6v38-local-curvature}.  Apply
\Cref{thm:newV91-global-gate}.  Its second condition ranges over the complete
fundamental group of the good orbit component and therefore contains more
data than one perturbative curvature calculation or one selected
\(SU(2)\) loop.
\end{proof}

\begin{definition}[Eta--mapping-torus phase card, EMTPC]
\label{def:newV91-EMTPC}
For a cutoff chiral multiplet, EMTPC consists of:
\begin{enumerate}[label=\textup{(\roman*)},leftmargin=3.5em]
\item the positive modulus \(\det_\zeta(D^2)^{1/2}\) and the selected eta
      square root \(\Phi=e^{-i\pi\eta/2}\) on every high-gap chart;
\item exact transfer of every zero crossing to the retained Grassmann
      determinant line, with spectral-flow parity recorded;
\item the APS cylinder identity \eqref{eq:newV91-cylinder-phase};
\item vanishing of the total local determinant-line curvature;
\item the mapping-torus identity
      \(\mathfrak H_{\rm tot}(\gamma)=1\) on every loop; and
\item a cutoff-compatible rooted local trivialization satisfying the V19
      derivative and complex polymer bounds.
\end{enumerate}
\end{definition}

\begin{proposition}[What is closed and what remains in EMTPC]
\label{prop:newV91-EMTPC-status}
For a smooth high-gap three-dimensional Dirac family, items
\textup{(i)} and \textup{(iii)} have the exact formulas proved in V91, and
item \textup{(ii)} has the exact crossing rule
\eqref{eq:newV91-crossing-phase}.  V38 supplies the standard local
coefficient part of item \textup{(iv)} for one Standard Model generation.
Items \textup{(v)} and \textup{(vi)}, and a single regulator realizing all
items simultaneously, remain unproved.
\end{proposition}

\begin{proof}
Use \Cref{thm:newV91-eta-completion,prop:newV91-crossing-phase,
thm:newV91-cylinder-phase,cor:newV91-SM-residual} and the V19 locality
criterion.
\end{proof}

\section{Next acquisition step}
\label{sec:newV91-next}

V92 will put EMTPC on an explicit finite three-torus regulator.  It will
define a gauge-covariant Hermitian Wilson--Dirac kernel, prove its
self-adjointness and gauge covariance, locate a quantitative high-mass gap,
and test exact conjugate-representation phase pairing.  That calculation will
separate phases which cancel algebraically at finite cutoff from the genuinely
chiral mapping-torus residue.  If the Wilson kernel introduces a circular use
of the desired continuum anomaly cancellation, the programme will move
upstream to an overlap/Ginsparg--Wilson spectral projector before estimating
any determinant.

\begin{firewall}[Time-slice eta is not the four-dimensional Weyl measure]
\label{fire:newV91-time-slice}
Equation \eqref{eq:newV91-eta-phase} completes the modulus of one
self-adjoint three-dimensional Dirac operator.  A four-dimensional Weyl
operator maps between different chiral bundles and defines a determinant
line rather than such a scalar.  Its global anomaly is the five-dimensional
mapping-torus quantity in \eqref{eq:newV91-DF-holonomy}.  The time-slice eta
phase cannot be substituted for that datum.
\end{firewall}

\begin{firewall}[Claim boundary after V91]
\label{fire:newV91-boundary}
V91 proves the finite rank-phase identity, derives the exact zeta
determinant modulus--eta decomposition, proves the spectral-flow square-root
rule, and derives the cylinder phase formula from the APS theorem.  It uses
the Dai--Freed theorem to turn the V19 abstract holonomy into the explicit
mapping-torus invariant \eqref{eq:newV91-total-holonomy}.

V91 does not construct a lattice chiral determinant, prove mapping-torus
triviality for every global form of the Standard Model gauge group, control
all Einstein backgrounds, construct a locality-compatible phase section,
establish a fermionic reflection-positive transfer kernel, or construct the
full SM--Einstein continuum.  The APS and Dai--Freed theorems are invoked as
established analytic inputs, not reproved here.
\end{firewall}

\begin{theorem}[V91 result ledger]
\label{thm:newV91-results}
V91:
\begin{enumerate}[label=\textup{(V91.\arabic*)},leftmargin=6.5em]
\item identifies the cutoff-rank phase missing from a finite Dirac modulus;
\item proves the zeta determinant decomposition into
      \(\det_\zeta(D^2)^{1/2}\) and an eta phase;
\item proves the upper/lower spectral-cut phase ratio;
\item identifies the eta phase as a square root of the inverse exponentiated
      xi invariant;
\item proves its spectral-flow parity rule at zero crossings;
\item derives the exact APS cylinder transport formula;
\item records the Dai--Freed mapping-torus holonomy with its graded sign;
\item converts V19's abstract global obstruction into
      \(\mathfrak H_{\rm tot}(\gamma)=1\) for every loop;
\item isolates the residual Standard Model mapping-torus datum after the V38
      local cancellations; and
\item formulates EMTPC and records its unclosed regulator and locality rows.
\end{enumerate}
\end{theorem}

\begin{proof}
Use \Cref{lem:newV91-finite-phase,
thm:newV91-eta-completion,
cor:newV91-cut-ambiguity,
lem:newV91-square-root,
prop:newV91-crossing-phase,
thm:newV91-cylinder-phase,
thm:newV91-DF-holonomy,
thm:newV91-global-gate,
cor:newV91-SM-residual,
prop:newV91-EMTPC-status}.
\end{proof}

\paragraph{Parent-version record.}
V91 was formed from
\texttt{Renewal\_SM\_Einstein\_Continuum\_Research\_Notes\_V90.tex}.
The V90 parent had \(30\,876\) logical source lines,
\(1\,185\,228\) bytes, and SHA--256
\[
 \texttt{D822F8D46F0619A01E352D8F98A9C085F9CAEBF3CB25F4DC440F9141ABB3C784}.
\]
The next compulsory companion audit remains V95.

% =============================================================================
% END APPENDED SECOND-SERIES V91 MATERIAL
% =============================================================================


% =============================================================================
% BEGIN APPENDED SECOND-SERIES V92 MATERIAL
% =============================================================================

\clearpage
\chapter{Second-series V92: an explicit Hermitian Wilson test kernel,
phase pairing, and the chiral-window obstruction}
\label{sec:v2v92}

\section{Purpose and relation to the inherited programme}
\label{sec:newV92-entry}

V91 formulated the eta--mapping-torus phase card.  The Grand Tensor field
shell assumes a gauge-covariant Hermitian Wilson kernel with a spectral
window, but does not prove that window.  This note constructs a literal
finite three-torus kernel and proves its elementary spectral and locality
properties uniformly in every link configuration.  The resulting window is
a large-mass chamber.  A free-symbol calculation then proves that this
chamber has constant sign and hence cannot supply the nontrivial chiral
projector required by an overlap regulator.  The failure identifies the next
upstream problem without borrowing the desired chiral conclusion.

\section{Gauge-covariant shifts on the finite three-torus}
\label{sec:newV92-shifts}

Let
\[
 \Lambda_{a,N}=(a\mathbb Z/Na\mathbb Z)^3
\]
with nearest-neighbour graph distance \(d_{\rm lat}\), measured in lattice
edges.  Let \(G\) be compact and let
\(R:G\to U(V_R)\) be a finite-dimensional unitary representation.  A link
field consists of matrices
\(U_j(x)\in R(G)\), \(j=1,2,3\).  On
\begin{equation}
 \mathcal H_{a,R}
 =\ell^2(\Lambda_{a,N};\mathbb C^2\otimes V_R)               \label{eq:newV92-Hilbert}
\end{equation}
use the inner product \(a^3\sum_x\langle\psi(x),\phi(x)\rangle\).
Define the covariant forward shift
\begin{equation}
 (T_j^U\psi)(x)=U_j(x)\psi(x+ae_j).                          \label{eq:newV92-shift}
\end{equation}

\begin{lemma}[Unitary shift and its adjoint]
\label{lem:newV92-shift-unitary}
Each \(T_j^U\) is unitary and
\begin{equation}
 ((T_j^U)^*\psi)(x)
 =U_j(x-ae_j)^*\psi(x-ae_j).                                \label{eq:newV92-shift-adjoint}
\end{equation}
\end{lemma}

\begin{proof}
Translation by \(ae_j\) permutes the finite periodic lattice and every
\(U_j(x)\) is unitary.  Changing variables \(y=x+ae_j\) in the inner
product gives \eqref{eq:newV92-shift-adjoint}; its product with
\eqref{eq:newV92-shift} in either order is the identity.
\end{proof}

For a lattice gauge transformation \(g:\Lambda_{a,N}\to G\), put
\begin{align}
 U_j^g(x)&=R(g(x))U_j(x)R(g(x+ae_j))^{-1},                   \label{eq:newV92-link-gauge}\
 (G_g\psi)(x)&=R(g(x))\psi(x).                              \label{eq:newV92-field-gauge}
\end{align}

\begin{lemma}[Exact covariance of the shift]
\label{lem:newV92-shift-covariance}
One has
\begin{equation}
 T_j^{U^g}=G_gT_j^UG_g^{-1}.                                \label{eq:newV92-shift-covariance}
\end{equation}
\end{lemma}

\begin{proof}
Apply the right side to \(\psi\) at \(x\).  The two factors at
\(x+ae_j\) cancel and the remaining expression is
\(U_j^g(x)\psi(x+ae_j)\).
\end{proof}

\section{The Hermitian Wilson test kernel}
\label{sec:newV92-kernel}

Let \(\sigma_1,\sigma_2,\sigma_3\) be the Pauli matrices and let
\(r\ge0\).  Define
\begin{align}
 Q_j(U)&={T_j^U-(T_j^U)^*\over2ia},                         \label{eq:newV92-Q}\
 D_{\rm n}(U)&=\sum_{j=1}^3\sigma_jQ_j(U),                  \label{eq:newV92-naive}\
 W_r(U)&={r\over2a}\sum_{j=1}^3
       (I-T_j^U)^*(I-T_j^U),                                \label{eq:newV92-Wilson}\
 H_{m,r}(U)&=mI+D_{\rm n}(U)+W_r(U).                        \label{eq:newV92-H}
\end{align}

\begin{theorem}[Self-adjointness, positivity of the Wilson term, and
uniform norms]
\label{thm:newV92-basic-bounds}
For every link field,
\begin{align}
 Q_j(U)^*&=Q_j(U),
 &D_{\rm n}(U)^*&=D_{\rm n}(U),                             \label{eq:newV92-selfadjoint-D}\
 W_r(U)&\succeq0,
 &H_{m,r}(U)^*&=H_{m,r}(U),                                 \label{eq:newV92-selfadjoint-H}\
 \|D_{\rm n}(U)\|&\le {3\over a},
 &\|W_r(U)\|&\le {6r\over a}.                             \label{eq:newV92-norms}
\end{align}
All four operators are gauge covariant; in particular
\begin{equation}
 H_{m,r}(U^g)=G_gH_{m,r}(U)G_g^{-1}.                        \label{eq:newV92-H-covariance}
\end{equation}
\end{theorem}

\begin{proof}
The numerator in \eqref{eq:newV92-Q} is skew-adjoint, so division by
\(i\) makes \(Q_j\) self-adjoint.  Pauli matrices act only on the spin
factor and commute with the shifts; hence every \(\sigma_jQ_j\) is
self-adjoint.  Formula \eqref{eq:newV92-Wilson} is a sum of squares.
This proves \eqref{eq:newV92-selfadjoint-D}--
\eqref{eq:newV92-selfadjoint-H}.

By \Cref{lem:newV92-shift-unitary},
\(\|T_j-(T_j)^*\|\le2\), whence \(\|Q_j\|\le a^{-1}\).
Also \(\|I-T_j\|\le2\), so each square in
\eqref{eq:newV92-Wilson} has norm at most four.  Summation gives
\eqref{eq:newV92-norms}.  Finally
\Cref{lem:newV92-shift-covariance} and its adjoint version pass through
every polynomial in the shifts, proving gauge covariance.
\end{proof}

\begin{corollary}[Gauge invariance of spectrum and determinant]
\label{cor:newV92-gauge-spectrum}
For every \(g\), the matrices \(H_{m,r}(U^g)\) and \(H_{m,r}(U)\) have the
same spectrum, eta count, determinant, and spectral-flow data along
gauge-related paths.
\end{corollary}

\begin{proof}
Equation \eqref{eq:newV92-H-covariance} is unitary equivalence.  Each listed
quantity is invariant under unitary conjugation.
\end{proof}

\section{Two uniform large-mass chambers}
\label{sec:newV92-gap}

Set
\begin{equation}
 B_{a,r}:={3+6r\over a}.                                    \label{eq:newV92-B}
\end{equation}

\begin{theorem}[Configuration-uniform spectral windows]
\label{thm:newV92-uniform-window}
For every link field:
\begin{enumerate}[label=\textup{(\roman*)},leftmargin=3.5em]
\item if \(m>3/a\), then
\begin{equation}
 H_{m,r}(U)\succeq(m-3/a)I;                                \label{eq:newV92-positive-window}
\end{equation}
\item if \(m<-B_{a,r}\), then
\begin{equation}
 H_{m,r}(U)\preceq-(|m|-B_{a,r})I.                         \label{eq:newV92-negative-window}
\end{equation}
\end{enumerate}
Consequently both chambers are free of zero crossings, uniformly in the
gauge field.
\end{theorem}

\begin{proof}
Since \(W_r\succeq0\) and
\(D_{\rm n}\succeq-\|D_{\rm n}\|I\),
\[
 H_{m,r}\succeq(m-3/a)I,
\]
which proves the first claim.  For the second, write
\(H=mI+A\) with \(A=D_{\rm n}+W_r\).  By
\eqref{eq:newV92-norms}, \(\|A\|\le B_{a,r}\), and therefore
\(H\preceq(-|m|+B_{a,r})I\).  The strict inequalities give the stated
uniform gaps.
\end{proof}

\begin{corollary}[Exact eta and determinant signs in the uniform chambers]
\label{cor:newV92-chamber-eta}
Let \(N_{a,R}=\dim\mathcal H_{a,R}\).  In the positive chamber,
\[
 \eta_{\rm fin}(H_{m,r})=N_{a,R},
 \qquad \det H_{m,r}>0.
\]
In the negative chamber,
\[
 \eta_{\rm fin}(H_{m,r})=-N_{a,R},
 \qquad
 {\det H_{m,r}\over|\det H_{m,r}|}=(-1)^{N_{a,R}}.
\]
\end{corollary}

\begin{proof}
All eigenvalues have the sign specified by
\Cref{thm:newV92-uniform-window}.  Count them and multiply their signs.
\end{proof}

\section{A direct inverse-locality bound}
\label{sec:newV92-locality}

The stronger condition \(|m|>B_{a,r}\) gives locality without invoking a
general Combes--Thomas theorem.

\begin{theorem}[Nearest-neighbour Neumann locality]
\label{thm:newV92-Neumann-locality}
Assume \(|m|>B_{a,r}\) and set \(q=B_{a,r}/|m|<1\).  Then
\begin{equation}
 H_{m,r}(U)^{-1}
 ={1\over m}\sum_{n=0}^\infty
 \left[-{D_{\rm n}(U)+W_r(U)\over m}\right]^n             \label{eq:newV92-Neumann}
\end{equation}
in operator norm.  Its site blocks obey
\begin{equation}
 \|H_{m,r}(U)^{-1}(x,y)\|
 \le {q^{d_{\rm lat}(x,y)}\over |m|(1-q)}.                 \label{eq:newV92-kernel-decay}
\end{equation}
The constants are uniform in the lattice volume and link field.
\end{theorem}

\begin{proof}
With \(A=D_{\rm n}+W_r\), \Cref{thm:newV92-basic-bounds} gives
\(\|A/m\|\le q<1\), so the Neumann series converges and proves
\eqref{eq:newV92-Neumann}.  The operator \(A\) has only diagonal and
nearest-neighbour site blocks.  Hence \((A^n)(x,y)=0\) when
\(n<d_{\rm lat}(x,y)\).  Every remaining block norm is at most
\(\|A\|^n\).  Summing the geometric tail proves
\eqref{eq:newV92-kernel-decay}.
\end{proof}

\begin{corollary}[Input for the V18 local trace-log chart]
\label{cor:newV92-V18-input}
Let \(U\) and \(U'\) lie in the same chamber
\(|m|>B_{a,r}\), set \(\Delta H=H(U')-H(U)\), and assume
\begin{equation}
 \|H(U)^{-1}\Delta H\|_\mu<1                              \label{eq:newV92-relative-small}
\end{equation}
in a weighted kernel norm whose weight is smaller than the decay rate in
\eqref{eq:newV92-kernel-decay}.  Then the relative determinant has the V18
absolutely convergent rooted complex trace-log expansion.
\end{corollary}

\begin{proof}
The perturbation \(\Delta H\) is nearest-neighbour and is supported on the
changed links and their endpoints.  The decay estimate makes
\(H(U)^{-1}\Delta H\) an exponentially weighted kernel.  Under
\eqref{eq:newV92-relative-small}, apply
Theorem~\ref{thm:v6v18-rooted-factorization} with
\(R=H(U)^{-1}\Delta H\).
\end{proof}

\section{Conjugate representations and phase orientation}
\label{sec:newV92-conjugation}

Let \(\overline R\) denote the conjugate representation and let
\(\overline U\) be the entrywise conjugate link field in a fixed unitary
basis.  Define the antiunitary map
\begin{equation}
 \Theta=(i\sigma_2)\,K:
 \mathcal H_{a,R}\longrightarrow\mathcal H_{a,\overline R}, \label{eq:newV92-Theta}
\end{equation}
where \(K\) is entrywise complex conjugation on spin and representation
coordinates.

\begin{lemma}[Antiunitary conjugate-representation equivalence]
\label{lem:newV92-antiunitary}
For real \(m,r\),
\begin{equation}
 \Theta H_{m,r}^{R}(U)\Theta^{-1}
 =H_{m,r}^{\overline R}(\overline U).                       \label{eq:newV92-antiunitary}
\end{equation}
Consequently the two operators have identical spectra and eta counts.
\end{lemma}

\begin{proof}
The Pauli identity
\((i\sigma_2)\overline{\sigma_j}(i\sigma_2)^{-1}=-\sigma_j\)
holds for \(j=1,2,3\).  Complex conjugation sends
\(T_j^U\) to \(T_j^{\overline U}\) and sends \(i^{-1}\) to
\(-i^{-1}\).  Thus it contributes one minus sign to \(Q_j\), while the
spin conjugation contributes a second minus sign to \(\sigma_j\); their
product is unchanged.  The Wilson term has real coefficients and contains
no Pauli matrix, so it is carried directly to the conjugate-representation
Wilson term.  The mass is real.  This proves
\eqref{eq:newV92-antiunitary}.  Antiunitary equivalence preserves real
eigenvalues with multiplicity.
\end{proof}

For an invertible finite kernel define the two rank-normalized phase
orientations
\begin{align}
 \widehat\Phi_R^\uparrow(U)
 &:=e^{-i\pi N_{a,R}/2}{\det H_R(U)\over|\det H_R(U)|}
   =e^{-i\pi\eta_{\rm fin}(H_R(U))/2},                      \label{eq:newV92-upper-phase}\
 \widehat\Phi_R^\downarrow(U)
 &:=e^{ i\pi N_{a,R}/2}{\det H_R(U)\over|\det H_R(U)|}
   =e^{ i\pi\eta_{\rm fin}(H_R(U))/2}.                    \label{eq:newV92-lower-phase}
\end{align}

\begin{theorem}[Exact opposite-orientation phase pairing]
\label{thm:newV92-phase-pairing}
For every invertible link configuration,
\begin{equation}
 \widehat\Phi_R^\uparrow(U)
 \widehat\Phi_{\overline R}^\downarrow(\overline U)=1.     \label{eq:newV92-phase-pairing}
\end{equation}
The equality is pointwise in the gauge field and does not use an anomaly
expansion.
\end{theorem}

\begin{proof}
By \Cref{lem:newV92-antiunitary}, the two eta counts are equal.  The
exponents in \eqref{eq:newV92-upper-phase} and
\eqref{eq:newV92-lower-phase} therefore cancel exactly.
\end{proof}

\begin{warning}[Conjugation alone does not choose the cancelling orientation]
\label{warn:newV92-orientation}
Equation \eqref{eq:newV92-phase-pairing} uses opposite upper/lower spectral
orientations.  Two same-orientation factors do not obey that identity.  In
the zeta limit their product contains \(e^{-i\pi\eta_D(0)}\), which is
generally nontrivial.  A vectorlike Dirac pair supplies the opposite
orientations through adjoint chiral blocks; an unpaired chiral multiplet does
not acquire cancellation merely because the conjugate representation exists
as an abstract representation.
\end{warning}

\section{The free symbol and failure of the uniform chamber to be chiral}
\label{sec:newV92-free-symbol}

For the trivial link field, Fourier momentum satisfies
\(k_j=2\pi n_j/(Na)\).  Put
\begin{align}
 M_{m,r}(k)&=m+{r\over a}\sum_{j=1}^3(1-\cos ak_j),          \label{eq:newV92-M-symbol}\
 s(k)&={1\over a}(\sin ak_1,\sin ak_2,\sin ak_3).           \label{eq:newV92-s-symbol}
\end{align}

\begin{lemma}[Exact free spectrum]
\label{lem:newV92-free-spectrum}
The free symbol and its eigenvalues are
\begin{align}
 h_{m,r}(k)&=M_{m,r}(k)I+\sigma\cdot s(k),                  \label{eq:newV92-free-symbol}\
 \lambda_\pm(k)&=M_{m,r}(k)\pm|s(k)|.                      \label{eq:newV92-free-eigenvalues}
\end{align}
At a Brillouin corner with exactly \(\ell\) components equal to
\(\pi/a\), one has \(s(k)=0\) and
\begin{equation}
 \lambda_\pm(k)=m+{2r\ell\over a}.                        \label{eq:newV92-corner-mass}
\end{equation}
\end{lemma}

\begin{proof}
At trivial links, the shift has multiplier \(e^{iak_j}\).  Hence
\(Q_j\) has multiplier \(\sin(ak_j)/a\), and the square in
\eqref{eq:newV92-Wilson} has multiplier \(2(1-\cos ak_j)\).
This proves \eqref{eq:newV92-free-symbol}.  Since
\((\sigma\cdot s)^2=|s|^2I\), its eigenvalues are \(\pm|s|\), proving
\eqref{eq:newV92-free-eigenvalues}.  The corner formula is immediate.
\end{proof}

\begin{theorem}[Uniform large-mass windows are topologically trivial]
\label{thm:newV92-trivial-window}
Let
\(\varepsilon(H)=\operatorname{sgn}H\) and
\(P_H=(I+\varepsilon(H))/2\).  Throughout the positive uniform chamber of
\Cref{thm:newV92-uniform-window},
\begin{equation}
 \varepsilon(H)=I,
 \qquad P_H=I.                                              \label{eq:newV92-positive-sign}
\end{equation}
Throughout the negative uniform chamber,
\begin{equation}
 \varepsilon(H)=-I,
 \qquad P_H=0.                                              \label{eq:newV92-negative-sign}
\end{equation}
Thus neither configuration-uniform large-mass chamber can be the nonconstant
moving chiral projector assumed by the overlap field shell.
\end{theorem}

\begin{proof}
The functional calculus sign is \(+1\) on a strictly positive spectrum and
\(-1\) on a strictly negative spectrum.  Apply
\Cref{thm:newV92-uniform-window}.  A constant projector has zero index
variation and no gauge-field-dependent determinant line, so it cannot supply
the required moving chiral subspace.
\end{proof}

\begin{corollary}[The actual Wilson-window acquisition problem]
\label{cor:newV92-actual-window}
A nontrivial overlap projector must use an intermediate Wilson mass lying
between free corner crossings in \eqref{eq:newV92-corner-mass}.  On that
branch the configuration-independent estimates
\eqref{eq:newV92-positive-window}--
\eqref{eq:newV92-negative-window} do not prove a gap.  A physical construction
must instead specify a gauge-field admissibility class and prove a uniform
lower bound for \(|H_W(U)|\) on that class.
\end{corollary}

\begin{proof}
Outside the intermediate region the sign projector is constant by
\Cref{thm:newV92-trivial-window}.  The free corner masses list the parameter
values at which sign sectors can change.  A continuous sign operator is
defined only while zero is excluded, so a nontrivial chamber needs a separate
gap estimate restricted to its admissible gauge fields.
\end{proof}

\begin{definition}[Hermitian Wilson test card, HWTC]
\label{def:newV92-HWTC}
HWTC consists of:
\begin{enumerate}[label=\textup{(\roman*)},leftmargin=3.5em]
\item exact gauge covariance and self-adjointness;
\item a configuration-uniform spectral window;
\item exponentially decaying inverse blocks;
\item exact conjugate-representation phase pairing with opposite spectral
      orientation; and
\item a nonconstant sign projector on a declared admissible gauge-field
      class.
\end{enumerate}
\end{definition}

\begin{proposition}[HWTC status]
\label{prop:newV92-HWTC-status}
Items \textup{(i)}--\textup{(iv)} are proved in the stronger chamber
\(|m|>B_{a,r}\).  Item \textup{(v)} fails there identically.  The overlap
programme must therefore move to an intermediate mass and prove a new
admissibility-dependent gap rather than reuse the large-mass estimate.
\end{proposition}

\begin{proof}
Use \Cref{thm:newV92-basic-bounds,thm:newV92-uniform-window,
thm:newV92-Neumann-locality,thm:newV92-phase-pairing,
thm:newV92-trivial-window}.
\end{proof}

\section{Next upstream step}
\label{sec:newV92-next}

V93 will pass to a four-dimensional \(\gamma_5\)-Hermitian Wilson kernel and
derive the overlap/Ginsparg--Wilson algebra under an explicit spectral-window
hypothesis.  It will prove smoothness and quantitative locality of the sign
projector from that window, then isolate the separate plaquette-admissibility
estimate needed to populate the hypothesis in the nontrivial mass interval.
This order avoids using anomaly cancellation to assume the chiral regulator
whose anomaly must later be computed.

\begin{firewall}[Claim boundary after V92]
\label{fire:newV92-boundary}
V92 constructs an explicit finite three-torus Hermitian Wilson test kernel,
proves gauge covariance, self-adjointness, two configuration-uniform
large-mass gaps, a direct exponential inverse bound, and exact
opposite-orientation phase pairing for conjugate representations.  It also
proves that these uniform chambers have constant sign projectors and are
therefore topologically trivial.

V92 does not prove an admissibility-dependent gap in the physical overlap
mass interval, define a four-dimensional chiral lattice determinant, prove a
Ginsparg--Wilson relation, orient the Standard Model determinant line, prove
the mapping-torus gate, establish fermionic reflection positivity, or
construct the full SM--Einstein continuum.  The mass bounds scale as
\(a^{-1}\) and do not describe a fixed physical fermion mass.
\end{firewall}

\begin{theorem}[V92 result ledger]
\label{thm:newV92-results}
V92:
\begin{enumerate}[label=\textup{(V92.\arabic*)},leftmargin=6.5em]
\item constructs the covariant finite-torus shifts and their adjoints;
\item constructs a literal Hermitian Wilson test kernel;
\item proves its self-adjointness, Wilson positivity, and uniform norms;
\item proves exact gauge covariance and spectral gauge invariance;
\item proves positive and negative configuration-uniform mass chambers;
\item proves a volume-uniform exponential inverse-kernel bound;
\item verifies the input of the earlier local trace-log theorem;
\item proves antiunitary equivalence of conjugate representations;
\item proves exact opposite-orientation phase cancellation;
\item computes the complete free symbol and corner masses;
\item proves that the large-mass sign projectors are constant; and
\item isolates the nontrivial admissibility-window gap as the next regulator
      acquisition.
\end{enumerate}
\end{theorem}

\begin{proof}
Use \Cref{lem:newV92-shift-unitary,
thm:newV92-basic-bounds,
cor:newV92-gauge-spectrum,
thm:newV92-uniform-window,
thm:newV92-Neumann-locality,
cor:newV92-V18-input,
lem:newV92-antiunitary,
thm:newV92-phase-pairing,
lem:newV92-free-spectrum,
thm:newV92-trivial-window,
cor:newV92-actual-window,
prop:newV92-HWTC-status}.
\end{proof}

\paragraph{Parent-version record.}
V92 was formed from
\texttt{Renewal\_SM\_Einstein\_Continuum\_Research\_Notes\_V91.tex}.
The V91 parent had \(31\,404\) logical source lines,
\(1\,206\,823\) bytes, and SHA--256
\[
 \texttt{ADA5B90CF154B93A3A74C22C321784FF71CDA4A927ED4C91B8F2496D30A32B0D}.
\]
The next compulsory companion audit remains V95.

% =============================================================================
% END APPENDED SECOND-SERIES V92 MATERIAL
% =============================================================================


% =============================================================================
% BEGIN APPENDED SECOND-SERIES V93 MATERIAL
% =============================================================================

\clearpage
\chapter{Second-series V93: the overlap algebra, sign-projector locality,
and the admissibility gap}
\label{sec:v2v93}

\section{Upstream correction}
\label{sec:newV93-entry}

V92 proves that an unrestricted large-mass Wilson chamber has a uniform gap
but a constant sign projector.  The chiral construction must instead use a
four-dimensional \(\gamma_5\)-Hermitian Wilson kernel at an intermediate
mass.  V93 separates two questions which must not be conflated:
\begin{enumerate}[label=\textup{(\roman*)},leftmargin=3.5em]
\item what follows rigorously from a two-sided spectral window for the
      Hermitian Wilson kernel; and
\item how that window is proved on a gauge-invariant admissible field class.
\end{enumerate}
The first question is closed below, including a direct exponential locality
bound for the sign operator.  For the second, a quantitative orbit-small
chart is constructed.  The extension from that chart to a full
plaquette-admissible topological sector remains a separate acquisition.

\section{Four-dimensional Wilson kernel}
\label{sec:newV93-Wilson}

Let \(\Lambda_{a,N}^{(4)}=(a\mathbb Z/Na\mathbb Z)^4\), let
\(T_\mu^U\) be the covariant unitary shifts defined as in V92, and choose
Hermitian Euclidean gamma matrices satisfying
\begin{equation}
 \{\gamma_\mu,\gamma_\nu\}=2\delta_{\mu\nu}I,
 \qquad
 \gamma_5=\gamma_1\gamma_2\gamma_3\gamma_4,
 \qquad
 \gamma_5^* =\gamma_5,
 \quad \gamma_5^2=I.                                       \label{eq:newV93-gamma}
\end{equation}
Use Wilson parameter \(r=1\) and define
\begin{align}
 D_W(U)
 &:={1\over2a}\sum_{\mu=1}^4
 \left\{
 \gamma_\mu(T_\mu^U-(T_\mu^U)^*)
 +(2I-T_\mu^U-(T_\mu^U)^*)
 \right\},                                                 \label{eq:newV93-DW}\
 A_\rho(U)&:=D_W(U)-{\rho\over a}I,                         \label{eq:newV93-A}\
 H_W(U)&:=\gamma_5A_\rho(U),
 \qquad 0<\rho<2.                                          \label{eq:newV93-HW}
\end{align}

\begin{lemma}[Gamma-five Hermiticity]
\label{lem:newV93-gamma5-Hermitian}
For every link field,
\begin{equation}
 \gamma_5D_W(U)\gamma_5=D_W(U)^*,
 \qquad
 H_W(U)^*=H_W(U).                                           \label{eq:newV93-gamma5-Hermitian}
\end{equation}
Moreover, for every lattice gauge transformation \(g\),
\begin{equation}
 H_W(U^g)=G_gH_W(U)G_g^{-1}.                                \label{eq:newV93-HW-gauge}
\end{equation}
\end{lemma}

\begin{proof}
The first summand in braces in \eqref{eq:newV93-DW} is
anti-self-adjoint.  Conjugation by \(\gamma_5\) changes its sign because
\(\gamma_5\gamma_\mu\gamma_5=-\gamma_\mu\).  The Wilson summand is
self-adjoint, has no spin matrix, and commutes with \(\gamma_5\).  This
proves the first identity, also after subtracting the real scalar
\(\rho/a\).  Therefore
\[
 H_W^*=A_\rho^*\gamma_5
 =\gamma_5A_\rho\gamma_5\gamma_5=H_W.
\]
Gauge covariance follows from the exact shift covariance proved in V92;
gamma matrices and scalar terms commute with \(G_g\).
\end{proof}

\section{The exact free window}
\label{sec:newV93-free}

At the trivial link field, momentum \(k\) has components in the Brillouin
zone.  Put \(\theta_\mu=ak_\mu\) and
\begin{equation}
 b(\theta)=\sum_{\mu=1}^4(1-\cos\theta_\mu).                \label{eq:newV93-b}
\end{equation}

\begin{lemma}[Free Wilson square]
\label{lem:newV93-free-square}
The free Hermitian kernel satisfies
\begin{equation}
 H_W(1;k)^2
 ={1\over a^2}
 \left\{
 \sum_{\mu=1}^4\sin^2\theta_\mu
 +(b(\theta)-\rho)^2
 \right\}I.                                                \label{eq:newV93-free-square}
\end{equation}
For \(0<\rho<2\),
\begin{equation}
 |H_W(1)|\succeq {g_0(\rho)\over a}I,
 \qquad
 g_0(\rho):=\min\{\rho,2-\rho\}>0.                       \label{eq:newV93-free-gap}
\end{equation}
\end{lemma}

\begin{proof}
The Fourier symbol of \(A_\rho\) is
\[
 {1\over a}
 \left\{i\sum_\mu\gamma_\mu\sin\theta_\mu
 +b(\theta)-\rho\right\}.
\]
Since \(H_W^2=A_\rho^*A_\rho\), the Clifford relations eliminate the
cross terms and give \eqref{eq:newV93-free-square}.

Set \(x_\mu=1-\cos\theta_\mu\in[0,2]\).  Then
\(\sin^2\theta_\mu=2x_\mu-x_\mu^2\), and the expression in braces is
\begin{equation}
 F(x)=\sum_\mu(2x_\mu-x_\mu^2)
       +\left(\sum_\mu x_\mu-\rho\right)^2.                \label{eq:newV93-F}
\end{equation}
With all other variables fixed, the \(x_\mu^2\) terms cancel and
\(F\) is affine in \(x_\mu\).  Its minimum on the cube therefore occurs
at a vertex.  If exactly \(\ell\) coordinates equal two, the value is
\((2\ell-\rho)^2\).  For \(0<\rho<2\), the least value over
\(\ell=0,\ldots,4\) is
\(\min\{\rho^2,(2-\rho)^2\}\), proving
\eqref{eq:newV93-free-gap}.
\end{proof}

\section{A populated orbit-small spectral chart}
\label{sec:newV93-orbit-small}

Define the gauge-orbit link distance
\begin{equation}
 \delta_{\rm orb}(U)
 :=\inf_g\max_{x,\mu}\|I-U_\mu^g(x)\|.                    \label{eq:newV93-orbit-distance}
\end{equation}

\begin{lemma}[Wilson perturbation bound]
\label{lem:newV93-Wilson-perturbation}
If \(\max_{x,\mu}\|I-U_\mu(x)\|\le\delta\), then
\begin{equation}
 \|H_W(U)-H_W(1)\|\le {8\delta\over a}.                   \label{eq:newV93-perturbation}
\end{equation}
\end{lemma}

\begin{proof}
For each direction,
\(\|T_\mu^U-T_\mu^1\|\le\delta\), and the same holds for the adjoints.
The naive term therefore changes by at most \(\delta/a\) per direction;
the Wilson term has the same bound.  Four directions give
\(8\delta/a\) for \(D_W\).  Multiplication by the unitary
\(\gamma_5\) does not change the norm.
\end{proof}

\begin{theorem}[Explicit orbit-small gap]
\label{thm:newV93-orbit-gap}
If
\begin{equation}
 8\delta_{\rm orb}(U)<g_0(\rho),                            \label{eq:newV93-orbit-condition}
\end{equation}
then
\begin{equation}
 |H_W(U)|\succeq
 {g_0(\rho)-8\delta_{\rm orb}(U)\over a}I.                 \label{eq:newV93-orbit-gap}
\end{equation}
\end{theorem}

\begin{proof}
For every \(\epsilon>0\), choose \(g\) for which all links of \(U^g\)
are within \(\delta_{\rm orb}(U)+\epsilon\) of the identity.  The least
singular value is one-Lipschitz in operator norm, so
\Cref{lem:newV93-free-square,lem:newV93-Wilson-perturbation} give
\[
 s_{\min}(H_W(U^g))
 \ge {g_0(\rho)-8\delta_{\rm orb}(U)-8\epsilon\over a}.
\]
Gauge covariance preserves singular values.  Let
\(\epsilon\downarrow0\).
\end{proof}

\begin{assessment}[Why this is not yet plaquette admissibility]
\label{ass:newV93-orbit-vs-plaquette}
The condition \eqref{eq:newV93-orbit-condition} is gauge invariant, but it
requires one global gauge in which every link is small.  A bound on
plaquette curvature alone does not supply that gauge on a torus without
additional control of Polyakov holonomies, bundle sector, and patching.
Therefore \Cref{thm:newV93-orbit-gap} populates a genuine chart but does not
prove the full admissibility window assumed by the Grand Tensor field shell.
\end{assessment}

\section{Sign operator and exponential locality}
\label{sec:newV93-sign}

Assume on a gauge-invariant chart \(\mathscr K\) that
\begin{equation}
 0<g_*I\preceq |H_W(U)|\preceq M_*I,
 \qquad U\in\mathscr K.                                    \label{eq:newV93-window}
\end{equation}
Define
\begin{equation}
 \varepsilon(U)=\operatorname{sgn}H_W(U)
 =H_W(U)|H_W(U)|^{-1}.                                     \label{eq:newV93-epsilon}
\end{equation}

\begin{theorem}[Exact sign algebra and gauge covariance]
\label{thm:newV93-sign-algebra}
On \(\mathscr K\),
\begin{equation}
 \varepsilon^*=\varepsilon,
 \qquad \varepsilon^2=I,
 \qquad
 \varepsilon(U^g)=G_g\varepsilon(U)G_g^{-1}.               \label{eq:newV93-sign-algebra}
\end{equation}
\end{theorem}

\begin{proof}
The first two statements follow from the spectral theorem and the exclusion
of zero.  Functional calculus commutes with unitary conjugation, so
\eqref{eq:newV93-HW-gauge} proves the third.
\end{proof}

\begin{theorem}[Direct volume-uniform sign-kernel decay]
\label{thm:newV93-sign-locality}
Let
\begin{equation}
 q:=1-{g_*^2\over M_*^2}\in[0,1).                           \label{eq:newV93-q}
\end{equation}
If \(0<q<1\), then with
\[
 n_0(\ell)=\max\left\{0,
 \left\lceil{\ell-1\over2}\right\rceil\right\}
\]
the site blocks satisfy
\begin{equation}
 \|\varepsilon(U)(x,y)\|
 \le {q^{n_0(d_{\rm lat}(x,y))}\over1-q}.                 \label{eq:newV93-sign-decay}
\end{equation}
The bound is uniform in lattice volume and \(U\in\mathscr K\).  If
\(q=0\), the sign operator has the finite propagation of \(H_W\).
\end{theorem}

\begin{proof}
Set \(X=I-H_W^2/M_*^2\).  Its spectrum lies in \([0,q]\).  The binomial
series gives
\begin{equation}
 \varepsilon
 ={H_W\over M_*}
 \sum_{n=0}^\infty c_nX^n,
 \qquad
 c_n={\binom{2n}{n}\over4^n}\le1.                          \label{eq:newV93-binomial-sign}
\end{equation}
The kernel \(H_W\) has propagation one and \(X\) has propagation two.
Thus the \(n\)-th summand vanishes between sites at graph distance greater
than \(2n+1\).  Every remaining block norm is at most \(c_nq^n\).
Summation from \(n_0(d_{\rm lat}(x,y))\) proves
\eqref{eq:newV93-sign-decay}.  If \(q=0\), then \(X=0\) by positivity, so
only the first term remains.
\end{proof}

\begin{theorem}[Analytic stability of the sign projector]
\label{thm:newV93-sign-derivatives}
On a norm-open set satisfying \eqref{eq:newV93-window}, the map
\(H\mapsto\operatorname{sgn}H\) is real analytic.  For each \(n\ge1\)
there is \(C_n=C_n(g_*,M_*)<\infty\) such that
\begin{equation}
 \|D^n\varepsilon_H[K_1,\ldots,K_n]\|
 \le C_n\prod_{j=1}^n\|K_j\|,                             \label{eq:newV93-sign-derivatives}
\end{equation}
and one may take
\(C_n\le n!C(M_*/g_*)g_*^{-n}\) with a numerical \(C\) after a fixed
choice of spectral contours.
\end{theorem}

\begin{proof}
Let \(\Gamma_+\) be a contour enclosing \([g_*,M_*]\) and staying at
distance at least \(g_*/2\) from the spectrum.  Then
\[
 P_+(H)={1\over2\pi i}\int_{\Gamma_+}(z-H)^{-1}\,dz,
 \qquad \varepsilon(H)=2P_+(H)-I.
\]
Differentiate the resolvent identity.  The \(n\)-th derivative is the sum
over \(n!\) orders of
\[
 (z-H)^{-1}K_{\sigma(1)}(z-H)^{-1}\cdots
 K_{\sigma(n)}(z-H)^{-1}.
\]
Every resolvent has norm at most \(2/g_*\), and a rectangular contour may
be chosen with length at most \(C(M_*+g_*)\).  Integration gives the stated
bound and the norm-convergent resolvent formula proves analyticity.
\end{proof}

\section{Overlap and Ginsparg--Wilson algebra}
\label{sec:newV93-overlap}

Define
\begin{equation}
 D_{\rm ov}={1\over a}(I+\gamma_5\varepsilon),
 \qquad
 \widehat\gamma_5:=\gamma_5(I-aD_{\rm ov})=-\varepsilon.    \label{eq:newV93-overlap}
\end{equation}

\begin{theorem}[Complete conditional overlap algebra]
\label{thm:newV93-GW}
Under \eqref{eq:newV93-window},
\begin{align}
 D_{\rm ov}^*&=\gamma_5D_{\rm ov}\gamma_5,                 \label{eq:newV93-overlap-Hermitian}\
 \gamma_5D_{\rm ov}+D_{\rm ov}\gamma_5
 &=aD_{\rm ov}\gamma_5D_{\rm ov},                         \label{eq:newV93-GW}\
 \widehat\gamma_5^*&=\widehat\gamma_5,
 &\widehat\gamma_5^2&=I,                                  \label{eq:newV93-hatgamma}\
 D_{\rm ov}\widehat\gamma_5&=-\gamma_5D_{\rm ov}.        \label{eq:newV93-intertwine}
\end{align}
The operator is gauge covariant and exponentially local with the same
off-diagonal rate as \(\varepsilon\).
\end{theorem}

\begin{proof}
Since \(\gamma_5\) and \(\varepsilon\) are self-adjoint unitaries,
\[
 D_{\rm ov}^*=a^{-1}(I+\varepsilon\gamma_5)
 =\gamma_5D_{\rm ov}\gamma_5.
\]
Expanding both sides of \eqref{eq:newV93-GW}, the right side is
\[
 {1\over a}(I+\gamma_5\varepsilon)
 \gamma_5(I+\gamma_5\varepsilon)
 ={1\over a}(2\gamma_5+\varepsilon+\gamma_5\varepsilon\gamma_5),
\]
which is also the left side.  The identities for
\(\widehat\gamma_5=-\varepsilon\) follow from
\Cref{thm:newV93-sign-algebra}; rearranging the Ginsparg--Wilson identity
gives \eqref{eq:newV93-intertwine}.  Gauge covariance follows from
\eqref{eq:newV93-sign-algebra}.  Off the diagonal, the identity term is
absent, so \Cref{thm:newV93-sign-locality} gives locality.
\end{proof}

Let
\begin{equation}
 P_\pm={I\pm\gamma_5\over2},
 \qquad
 \widehat P_-={I-\widehat\gamma_5\over2}
 ={I+\varepsilon\over2}.                                   \label{eq:newV93-projectors}
\end{equation}

\begin{corollary}[Exact moving-source chiral map]
\label{cor:newV93-chiral-map}
One has
\begin{equation}
 D_{\rm ov}\widehat P_-=P_+D_{\rm ov}.                    \label{eq:newV93-chiral-map}
\end{equation}
Hence
\begin{equation}
 D_\chi:=P_+D_{\rm ov}\widehat P_-:
 \operatorname{Ran}\widehat P_-longrightarrow
 \operatorname{Ran}P_+                                    \label{eq:newV93-Dchi}
\end{equation}
is the coordinate-free chiral block used by the Grand Tensor field shell.
\end{corollary}

\begin{proof}
Multiply \eqref{eq:newV93-intertwine} into the definitions of the two
projectors:
\[
 D_{\rm ov}(I-\widehat\gamma_5)
 =(I+\gamma_5)D_{\rm ov}.
\]
Divide by two and insert the projectors.
\end{proof}

\begin{proposition}[Finite index and rank stability]
\label{prop:newV93-index}
On the periodic lattice \(\operatorname{Tr}\gamma_5=0\), and the finite
Fredholm index of the map \eqref{eq:newV93-Dchi} is
\begin{equation}
 \operatorname{ind}D_\chi
 =\dim\operatorname{Ran}\widehat P_-
  -\dim\operatorname{Ran}P_+
 ={1\over2}\operatorname{Tr}\varepsilon.                  \label{eq:newV93-index}
\end{equation}
It is constant on every connected chart satisfying
\eqref{eq:newV93-window}.  It can change only when an eigenvalue of
\(H_W\) crosses zero.
\end{proposition}

\begin{proof}
For a linear map between finite-dimensional spaces, kernel dimension minus
cokernel dimension equals domain dimension minus codomain dimension.
Taking traces in \eqref{eq:newV93-projectors} gives
\[
 \dim\operatorname{Ran}\widehat P_-
 -\dim\operatorname{Ran}P_+
 ={1\over2}(\operatorname{Tr}\varepsilon-
             \operatorname{Tr}\gamma_5).
\]
The spin trace of \(\gamma_5\) vanishes at every lattice site.
Under a spectral gap, \(\varepsilon\) varies continuously while its trace is
integral, hence its trace is constant.  A change requires loss of the gap;
the sign trace then changes by twice the signed crossing multiplicity.
\end{proof}

\begin{corollary}[Determinant-line rather than scalar output]
\label{cor:newV93-det-line}
The top exterior power of \(D_\chi\) is naturally a section of
\begin{equation}
 \det(P_+\mathcal H)\otimes
 \det(\widehat P_-\mathcal H)^*.                            \label{eq:newV93-det-line}
\end{equation}
It is a nonzero scalar determinant only after an index-zero chart, an
invertibility bound for \(D_\chi\), and a compatible trivialization of the
moving determinant line have all been supplied.
\end{corollary}

\begin{proof}
The top exterior power of a map lies in the determinant of its target
tensor the dual determinant of its source.  Equal ranks are necessary for a
square determinant, invertibility is necessary for nonvanishing, and V19
proves that scalarization requires a determinant-line trivialization.
\end{proof}

\section{The exact remaining admissibility statement}
\label{sec:newV93-admissibility}

For a plaquette \(p\), let \(U_p\) be its oriented holonomy and define
\begin{equation}
 \mathscr A_\delta
 :=\left\{U:\max_p\|I-U_p\|<\delta\right\}.                \label{eq:newV93-admissible}
\end{equation}
This set is gauge invariant.

\begin{definition}[Admissible spectral-window gate, ASWG]
\label{def:newV93-ASWG}
ASWG at \((\rho,\delta)\) is a volume-independent constant
\(c(\rho,\delta)>0\) for which
\begin{equation}
 H_W(U)^2\succeq {c(\rho,\delta)^2\over a^2}I
 \qquad\hbox{for every }U\in\mathscr A_\delta              \label{eq:newV93-ASWG}
\end{equation}
in every admitted bundle and Polyakov sector.
\end{definition}

\begin{theorem}[ASWG closes the algebraic overlap regulator]
\label{thm:newV93-ASWG-closes}
If ASWG holds and \(\|H_W\|\le M_*/a\) uniformly, then the sign operator,
overlap operator, moving chiral projector, and finite index are all
gauge-covariant, analytic on each admissible component, and exponentially
local with volume-independent lattice-unit constants.  The remaining phase
problem is then the determinant-line EMTPC of V91, not existence of the sign
operator.
\end{theorem}

\begin{proof}
Set \(g_*=c(\rho,\delta)/a\) in
\Cref{thm:newV93-sign-algebra,thm:newV93-sign-locality,
thm:newV93-sign-derivatives,thm:newV93-GW,prop:newV93-index}.  The ratio
\(g_*/(M_*/a)=c/M_*\) is independent of \(a\) and volume, so the lattice-unit
decay rate in \eqref{eq:newV93-sign-decay} is uniform.  V91 describes the
additional orientation and mapping-torus rows after this operator exists.
\end{proof}

\begin{assessment}[Current acquisition boundary]
\label{ass:newV93-ASWG-status}
The orbit-small theorem \ref{thm:newV93-orbit-gap} proves ASWG on the
strict subchart where one global gauge makes every link small.  It does not
prove \eqref{eq:newV93-ASWG} on all of
\(\mathscr A_\delta\).  Conversely, the overlap identities do not imply
ASWG: each uses \(\varepsilon=H_W|H_W|^{-1}\), which is undefined when the
gap fails.  Using the Ginsparg--Wilson relation to justify its own sign
operator would therefore be circular.
\end{assessment}

\begin{definition}[Overlap algebra and locality card, OALC]
\label{def:newV93-OALC}
OALC comprises gamma-five Hermiticity, gauge covariance, a two-sided
spectral window, sign-projector locality and analytic stability, the
Ginsparg--Wilson identity, the moving-source intertwining identity, and the
finite determinant-line index formula.
\end{definition}

\begin{theorem}[OALC status]
\label{thm:newV93-OALC-status}
OALC is proved on every chart carrying the window
\eqref{eq:newV93-window}; in particular it is populated on the orbit-small
chart \eqref{eq:newV93-orbit-condition}.  Population on the full
plaquette-admissible class is equivalent here to the remaining ASWG estimate.
\end{theorem}

\begin{proof}
Use \Cref{lem:newV93-gamma5-Hermitian,
thm:newV93-orbit-gap,
thm:newV93-sign-algebra,
thm:newV93-sign-locality,
thm:newV93-sign-derivatives,
thm:newV93-GW,
cor:newV93-chiral-map,
prop:newV93-index,
cor:newV93-det-line,
thm:newV93-ASWG-closes}.
\end{proof}

\section{Next acquisition step}
\label{sec:newV93-next}

V94 will attack ASWG directly.  It will expand \(H_W^2\) into its free
positive part plus commutators of covariant shifts, express every commutator
through a plaquette defect, and determine whether the resulting norm bound
is small enough to preserve the free gap for an explicit
\(\delta<\delta_*(\rho)\).  Global flat holonomies will be retained exactly
rather than gauged away.  If the direct square estimate loses the free gap,
the programme will isolate the offending mixed Wilson--naive term before
trying a stronger topological argument.

\begin{firewall}[Claim boundary after V93]
\label{fire:newV93-boundary}
V93 constructs the four-dimensional gamma-five-Hermitian Wilson kernel,
computes its exact free gap, proves a nonempty gauge-invariant orbit-small
window, proves exponential locality and analytic stability of the sign
operator under any uniform window, and derives the full overlap,
Ginsparg--Wilson, moving-projector, and finite-index algebra.

V93 does not prove ASWG on the complete plaquette-admissible class, a lower
singular-value bound for the Yukawa-perturbed chiral block, determinant-line
trivialization, the mapping-torus phase gate, reflection positivity, cutoff
transport, or the full SM--Einstein continuum.  The orbit-small chart is a
proper local chart and is not substituted for all topological sectors.
\end{firewall}

\begin{theorem}[V93 result ledger]
\label{thm:newV93-results}
V93:
\begin{enumerate}[label=\textup{(V93.\arabic*)},leftmargin=6.5em]
\item proves gamma-five Hermiticity and gauge covariance of a literal
      four-dimensional Wilson kernel;
\item computes its exact free square and free gap for \(0<\rho<2\);
\item proves an explicit gauge-orbit-small interacting gap;
\item separates orbit-small control from full plaquette admissibility;
\item proves exact sign algebra and gauge covariance;
\item proves a volume-uniform exponential sign-kernel bound;
\item proves analytic sign-projector derivative bounds;
\item derives gamma-five Hermiticity and the Ginsparg--Wilson identity for
      the overlap operator;
\item proves the moving-source chiral intertwining identity;
\item proves the finite chiral index and rank-stability formula;
\item identifies the determinant-line output before scalarization;
\item formulates ASWG and proves that it closes the conditional overlap
      construction; and
\item closes OALC on every populated spectral chart while leaving full ASWG
      as the next acquisition.
\end{enumerate}
\end{theorem}

\begin{proof}
Use \Cref{lem:newV93-gamma5-Hermitian,
lem:newV93-free-square,
lem:newV93-Wilson-perturbation,
thm:newV93-orbit-gap,
thm:newV93-sign-algebra,
thm:newV93-sign-locality,
thm:newV93-sign-derivatives,
thm:newV93-GW,
cor:newV93-chiral-map,
prop:newV93-index,
cor:newV93-det-line,
thm:newV93-ASWG-closes,
thm:newV93-OALC-status}.
\end{proof}

\paragraph{Parent-version record.}
V93 was formed from
\texttt{Renewal\_SM\_Einstein\_Continuum\_Research\_Notes\_V92.tex}.
The V92 parent had \(31\,916\) logical source lines,
\(1\,226\,868\) bytes, and SHA--256
\[
 \texttt{7ECCA8294E64B8DCD4B206737FD0E0718DFFBC0260F299D89D0C986837C96E18}.
\]
The next compulsory companion audit remains V95.

% =============================================================================
% END APPENDED SECOND-SERIES V93 MATERIAL
% =============================================================================


% =============================================================================
% BEGIN APPENDED SECOND-SERIES V94 MATERIAL
% =============================================================================

\clearpage
\chapter{Second-series V94: closure of the plaquette-admissible Wilson gap}
\label{sec:v2v94}

\section{Acquisition target and normalization}
\label{sec:newV94-entry}

V93 left ASWG open on the complete plaquette-admissible class.  At the
central overlap mass \(\rho=1\), the required estimate follows from an exact
square decomposition due to the standard locality analysis of the Neuberger
operator.  This note reproduces the operator ledger and its norm proof in the
normalization of V93.  The result is
\[
 a^2H_W(U)^2\succeq(1-30\delta)I
\]
whenever every represented plaquette obeys
\(\|I-U_p\|\le\delta\).  No global gauge choice and no small-link
hypothesis is used.

Let
\begin{equation}
 \nabla_\mu={T_\mu-I\over a},
 \qquad
 \nabla_\mu^*={I-T_\mu^*\over a}.                           \label{eq:newV94-differences}
\end{equation}
Thus the Hilbert adjoint of \(\nabla_\mu\) is
\(-\nabla_\mu^*\).  The V93 Wilson operator may be written
\begin{equation}
 D_W={1\over2}\sum_\mu
 \left\{\gamma_\mu(\nabla_\mu^*+\nabla_\mu)
       -a\nabla_\mu^*\nabla_\mu\right\}.                   \label{eq:newV94-DW-difference}
\end{equation}
At \(\rho=1\), put
\begin{equation}
 A:=I-aD_W=-aA_{\rho=1},
 \qquad
 a^2H_W^2=A^*A.                                             \label{eq:newV94-A}
\end{equation}

\section{Plaquette control of covariant commutators}
\label{sec:newV94-commutators}

Assume
\begin{equation}
 \max_p\|I-U_p\|\le\delta.                                 \label{eq:newV94-admissibility}
\end{equation}

\begin{lemma}[All forward--backward commutators]
\label{lem:newV94-commutators}
For \(\mu\ne\nu\), and for either choice of a forward or backward
covariant difference in each slot,
\begin{equation}
 a^2\|[\nabla_\mu^{\#},\nabla_\nu^{\flat}]\|
 \le\delta,                                                \label{eq:newV94-all-commutators}
\end{equation}
where \(\# ,\flat\in\{\phantom{*},*\}\).
\end{lemma}

\begin{proof}
For two forward differences, the identity and one-shift terms cancel and
\begin{align*}
 a^2[\nabla_\mu,\nabla_\nu]\psi(x)
 ={}&\{U_\mu(x)U_\nu(x+a e_\mu)
       -U_\nu(x)U_\mu(x+a e_\nu)\}\\
 &\hspace{8em}\cdot\psi(x+ae_\mu+ae_\nu).
\end{align*}
The matrix in braces is a product of unitaries times \(I-U_p\) for the
corresponding oriented plaquette, so its norm is at most \(\delta\).
Replacing a forward shift by its adjoint changes the plaquette orientation
and translates its base point; the same norm bound applies.  Division by
\(a^2\) proves every case.
\end{proof}

\begin{corollary}[Two commutator combinations]
\label{cor:newV94-combination-bounds}
For \(\mu\ne\nu\),
\begin{align}
 a^2\|[\nabla_\mu^*+\nabla_\mu,
          \nabla_\nu^*+\nabla_\nu]\|&\le4\delta,            \label{eq:newV94-plus-plus}\
 a^2\|[\nabla_\mu^*+\nabla_\mu,
          \nabla_\nu^*-\nabla_\nu]\|&\le4\delta.          \label{eq:newV94-plus-minus}
\end{align}
\end{corollary}

\begin{proof}
Expand each bracket into four commutators and use
\Cref{lem:newV94-commutators} and the triangle inequality.
\end{proof}

\section{Exact Wilson-square decomposition}
\label{sec:newV94-square}

For \(\mu\ne\nu\), set
\begin{align}
 B_{\mu\nu}
 &:=a^4\nabla_\mu^*\nabla_\mu
          \nabla_\nu^*\nabla_\nu,                          \label{eq:newV94-B}\
 C_{\mu\nu}
 &:={i\over2}\sigma_{\mu\nu}a^2
 [\nabla_\mu^*+\nabla_\mu,
  \nabla_\nu^*+\nabla_\nu],                                \label{eq:newV94-C}\
 D_{\mu\nu}
 &:=-\gamma_\mu a^2
 [\nabla_\mu^*+\nabla_\mu,
  \nabla_\nu^*-\nabla_\nu],                               \label{eq:newV94-D}\
 \sigma_{\mu\nu}&={i\over2}[\gamma_\mu,\gamma_\nu].      \label{eq:newV94-sigma}
\end{align}

\begin{lemma}[Exact square identity]
\label{lem:newV94-square-identity}
At \(\rho=1\),
\begin{equation}
 A^*A
 =I+{1\over4}\sum_{\mu\ne\nu}
 \{B_{\mu\nu}+C_{\mu\nu}+D_{\mu\nu}\}.                 \label{eq:newV94-square-identity}
\end{equation}
\end{lemma}

\begin{proof}
Insert \eqref{eq:newV94-DW-difference} into
\(A=I-aD_W\) and multiply \(A^*A\).  The diagonal
\(\mu=\nu\) terms cancel at \(\rho=1\) after using
\begin{equation}
 \nabla_\mu^*+\nabla_\mu
 =-a\nabla_\mu^*\nabla_\mu.                                \label{eq:newV94-difference-identity}
\end{equation}
For \(\mu\ne\nu\), the Clifford anticommutator separates the scalar
Wilson--Wilson product \(B_{\mu\nu}\), the antisymmetric
gamma--gamma product \(C_{\mu\nu}\), and the two mixed
naive--Wilson orders \(D_{\mu\nu}\).  Pairing the ordered terms
\((\mu,\nu)\) and \((\nu,\mu)\) gives the factor \(1/4\) and precisely
\eqref{eq:newV94-B}--\eqref{eq:newV94-D}.  This is a finite algebraic
identity in the covariant shifts; no commutation of distinct directions is
assumed.
\end{proof}

\section{The three lower bounds}
\label{sec:newV94-three-bounds}

\begin{lemma}[The spin commutator rows]
\label{lem:newV94-CD-bounds}
For every ordered pair \(\mu\ne\nu\),
\begin{equation}
 \|C_{\mu\nu}\|\le2\delta,
 \qquad
 \|D_{\mu\nu}\|\le4\delta.                               \label{eq:newV94-CD-bounds}
\end{equation}
\end{lemma}

\begin{proof}
For distinct indices, \(\|\gamma_\mu\|=1\) and
\(\|\sigma_{\mu\nu}\|=1\).  Apply
\Cref{cor:newV94-combination-bounds} to
\eqref{eq:newV94-C}--\eqref{eq:newV94-D}, including the factor one half
in \(C_{\mu\nu}\).
\end{proof}

\begin{lemma}[The Wilson--Wilson row]
\label{lem:newV94-B-bound}
For every ordered pair \(\mu\ne\nu\),
\begin{equation}
 B_{\mu\nu}\succeq-4\delta I.                              \label{eq:newV94-B-bound}
\end{equation}
\end{lemma}

\begin{proof}
Move \(\nabla_\mu\) through the \(\nu\)-factor:
\begin{align}
 B_{\mu\nu}
 ={}&a^4\nabla_\mu^*\nabla_\nu^*\nabla_\nu\nabla_\mu
 \notag\\
 &-a^3\nabla_\mu^*
 [\nabla_\mu,\nabla_\nu^*-\nabla_\nu].                    \label{eq:newV94-B-rewrite}
\end{align}
The first term is
\(a^4(\nabla_\nu\nabla_\mu)^*
        (\nabla_\nu\nabla_\mu)\) and is nonnegative, because the two
minus signs in the adjoints cancel.  Moreover
\(a\|\nabla_\mu^*\|\le2\), while
\Cref{lem:newV94-commutators} gives
\[
 a^2\|[\nabla_\mu,\nabla_\nu^*-\nabla_\nu]\|
 \le2\delta.
\]
The remainder in \eqref{eq:newV94-B-rewrite} therefore has norm at most
\(4\delta\), proving the operator lower bound.
\end{proof}

\section{Closure of ASWG}
\label{sec:newV94-ASWG}

\begin{theorem}[Plaquette-admissible Wilson gap]
\label{thm:newV94-admissible-gap}
For \(\rho=1\), every link field satisfying
\eqref{eq:newV94-admissibility} obeys
\begin{equation}
 A^*A\succeq(1-30\delta)I,
 \qquad
 H_W^2\succeq{1-30\delta\over a^2}I.                       \label{eq:newV94-gap}
\end{equation}
In particular, if \(0\le\delta<1/30\), then
\begin{equation}
 |H_W|\succeq{\sqrt{1-30\delta}\over a}I.                  \label{eq:newV94-gap-root}
\end{equation}
The constants are independent of lattice volume, gauge orbit, and flat
holonomy.
\end{theorem}

\begin{proof}
There are \(4\cdot3=12\) ordered pairs \(\mu\ne\nu\).  By
\Cref{lem:newV94-CD-bounds,lem:newV94-B-bound}, each brace in
\eqref{eq:newV94-square-identity} is bounded below by
\(-(4+2+4)\delta I=-10\delta I\).  The prefactor \(1/4\) therefore gives
\[
 A^*A\succeq
 \left(1-{12\over4}10\delta\right)I
 =(1-30\delta)I.
\]
Use \eqref{eq:newV94-A} for the second inequality and functional calculus
for its square root.  The proof used only plaquette defects, so no global
gauge or holonomy enters.
\end{proof}

\begin{lemma}[Uniform upper window]
\label{lem:newV94-upper-window}
At \(\rho=1\),
\begin{equation}
 \|A\|\le7,
 \qquad
 \|H_W\|\le {7\over a}.                                    \label{eq:newV94-upper-window}
\end{equation}
\end{lemma}

\begin{proof}
Let \(P_\mu^\pm=(I\pm\gamma_\mu)/2\).  Directly from the shift form of
\(D_W\),
\begin{equation}
 A=-3I+\sum_{\mu=1}^4
 \{P_\mu^-T_\mu+P_\mu^+T_\mu^*\}.                         \label{eq:newV94-A-shift}
\end{equation}
For fixed \(\mu\), the operator in braces is unitary: the two orthogonal
spin projectors commute with the shifts and its adjoint product is
\(P_\mu^-+P_\mu^+=I\).  The triangle inequality gives
\(\|A\|\le3+4=7\).  Equation \eqref{eq:newV94-A} gives the bound for
\(H_W\).
\end{proof}

\begin{corollary}[Explicit sign-locality constants]
\label{cor:newV94-locality}
For \(\delta<1/30\), V93's sign-kernel theorem applies with
\begin{equation}
 g_*={\sqrt{1-30\delta}\over a},
 \qquad
 M_*={7\over a},
 \qquad
 q_\delta=1-{1-30\delta\over49}.                            \label{eq:newV94-locality-constants}
\end{equation}
Hence
\begin{equation}
 \|\varepsilon(U)(x,y)\|
 \le {q_\delta^{,n_0(d_{\rm lat}(x,y))}\over1-q_\delta}  \label{eq:newV94-sign-locality}
\end{equation}
uniformly in volume and throughout the admissible class.
\end{corollary}

\begin{proof}
Combine \Cref{thm:newV94-admissible-gap,
lem:newV94-upper-window} with
Theorem~\ref{thm:newV93-sign-locality}.
\end{proof}

\begin{theorem}[ASWG and OALC are populated on the full admissible class]
\label{thm:newV94-ASWG-closed}
At \(\rho=1\) and \(\delta<1/30\), ASWG holds on
\(\mathscr A_\delta\) with
\(c(1,\delta)=\sqrt{1-30\delta}\).  Consequently OALC holds on every
connected component of \(\mathscr A_\delta\): the overlap sign,
Ginsparg--Wilson operator, moving chiral projector, and finite index are
gauge covariant, analytic, and exponentially local with volume-independent
lattice-unit constants.
\end{theorem}

\begin{proof}
The first statement is \Cref{thm:newV94-admissible-gap}; the second is
Theorem~\ref{thm:newV93-ASWG-closes}.  Analyticity is componentwise because
the uniform gap excludes crossings inside \(\mathscr A_\delta\).
\end{proof}

\begin{corollary}[Admissible topological sectors are separated]
\label{cor:newV94-sector-separation}
The integer
\begin{equation}
 \nu(U):={1\over2}\operatorname{Tr}\varepsilon(U)           \label{eq:newV94-index}
\end{equation}
is constant on every connected component of
\(\mathscr A_\delta\), \(\delta<1/30\).  Configurations with different
values of \(\nu\) cannot be joined by a path remaining in the admissible
class.
\end{corollary}

\begin{proof}
V93 proves that \(\nu\) is integral and locally constant wherever the
Wilson gap is open.  V94 opens that gap on all of
\(\mathscr A_\delta\).  A continuous path has connected image, so its
integer value cannot change.
\end{proof}

\section{Representation and Standard Model scope}
\label{sec:newV94-SM-scope}

\begin{assessment}[Representationwise admissibility]
\label{ass:newV94-representation}
The norm in \eqref{eq:newV94-admissibility} is taken in the unitary
representation carried by the fermion kernel.  For a Standard Model
multiplet, the estimate applies simultaneously if the represented plaquette
defect is below \(1/30\) in every fermion representation.  Passing from one
preferred fundamental norm on
\((SU(3)\times SU(2)\times U(1))/\mathbb Z_6\) to all finitely many
multiplet norms requires only finite representation-dependent continuity
constants, but their common numerical admissibility radius is additional
bookkeeping and is not set here.
\end{assessment}

\begin{assessment}[What the gap now licenses]
\label{ass:newV94-license}
The Grand Tensor spectral-window hypothesis for the massless overlap kernel
is no longer empty on the declared admissible class: V94 supplies explicit
lower and upper bounds and V93 supplies the resulting algebra and locality.
This licenses construction of the finite chiral determinant line.  It does
not choose a global phase section, prove the chiral-block singular-value
floor after Yukawa coupling, or show that a chosen gauge action concentrates
on \(\mathscr A_\delta\) as the cutoff is removed.
\end{assessment}

\begin{definition}[Admissible overlap packet, AOP]
\label{def:newV94-AOP}
AOP consists of the plaquette condition
\(\delta<1/30\), the bounds
\eqref{eq:newV94-gap-root} and
\eqref{eq:newV94-upper-window}, the sign-locality estimate
\eqref{eq:newV94-sign-locality}, and all OALC identities of V93.
\end{definition}

\begin{theorem}[AOP is closed]
\label{thm:newV94-AOP}
For every finite periodic four-lattice, every unitary representation, and
every represented link field obeying
\(\max_p\|I-U_p\|<1/30\), AOP holds with constants independent of lattice
volume.
\end{theorem}

\begin{proof}
Use \Cref{thm:newV94-admissible-gap,
lem:newV94-upper-window,
cor:newV94-locality,
thm:newV94-ASWG-closed}.
\end{proof}

\section{V95 audit and next gate}
\label{sec:newV94-next}

V95 is a compulsory companion checkpoint.  After reading the newest settled
YM and NS notes, it will combine AOP with the V91 eta/mapping-torus ledger.
The next local operator question is the singular-value floor for the
Yukawa-perturbed chiral block on an index-zero admissible component.  The
next global question is the locality-compatible phase section.  The audit
will decide which has the more usable upstream input; neither will be
declared solved by the Wilson gap alone.

\begin{firewall}[Claim boundary after V94]
\label{fire:newV94-boundary}
V94 proves the full plaquette-admissible Wilson gap at \(\rho=1\), with
explicit threshold \(\delta<1/30\), and hence closes ASWG and OALC on every
component of that admissible class.  It proves the estimate by an exact
Wilson-square decomposition and plaquette commutator bounds, so flat
holonomies and disconnected admissible sectors are retained.

V94 does not prove that the interacting gauge measure is supported or
concentrated on this strict admissible set, construct the Standard Model
determinant phase, prove the Yukawa chiral floor, trivialize the determinant
line, establish reflection positivity, remove the cutoff, or construct the
full SM--Einstein continuum.  The constant \(30\) is sufficient and is not
claimed optimal.
\end{firewall}

\begin{theorem}[V94 result ledger]
\label{thm:newV94-results}
V94:
\begin{enumerate}[label=\textup{(V94.\arabic*)},leftmargin=6.5em]
\item rewrites the V93 Wilson kernel in covariant-difference form;
\item proves that every forward/backward commutator is controlled by one
      plaquette defect;
\item states and verifies the exact Wilson-square decomposition at
      \(\rho=1\);
\item proves the separate \(B\), \(C\), and \(D\) lower bounds;
\item proves \(A^*A\succeq(1-30\delta)I\);
\item proves the uniform upper bound \(\|A\|\le7\);
\item obtains explicit volume-independent sign-locality constants;
\item closes ASWG and OALC on the complete admissible class;
\item proves separation of admissible index sectors; and
\item closes AOP while isolating measure concentration, the Yukawa floor,
      and the determinant phase as distinct remaining gates.
\end{enumerate}
\end{theorem}

\begin{proof}
Use \Cref{lem:newV94-commutators,
lem:newV94-square-identity,
lem:newV94-CD-bounds,
lem:newV94-B-bound,
thm:newV94-admissible-gap,
lem:newV94-upper-window,
cor:newV94-locality,
thm:newV94-ASWG-closed,
cor:newV94-sector-separation,
thm:newV94-AOP}.
\end{proof}

\paragraph{Parent-version record.}
V94 was formed from
\texttt{Renewal\_SM\_Einstein\_Continuum\_Research\_Notes\_V93.tex}.
The V93 parent had \(32\,504\) logical source lines,
\(1\,248\,654\) bytes, and SHA--256
\[
 \texttt{B9BD64B92EB685031F87E36BA23C0A14F1B961946CB73084A4574A52CA926463}.
\]
The next compulsory companion audit is V95.

% =============================================================================
% END APPENDED SECOND-SERIES V94 MATERIAL
% =============================================================================


% =============================================================================
% BEGIN APPENDED SECOND-SERIES V95 MATERIAL
% =============================================================================

\clearpage
\chapter{Second-series V95: companion audit, the zero-Higgs obstruction,
and an exact massive-overlap floor}
\label{sec:v2v95}

\section{Compulsory companion audit}
\label{sec:newV95-audit}

The companion notes were read as immutable byte snapshots.  The settled
Yang--Mills snapshot was
\begin{center}
\begin{tabular}{ll}
file &
\texttt{Renewal\_Yang\_Mills\_Mass\_Gap\_Research\_Notes\_V69.tex},\\
bytes & \(1\,048\,914\),\\
logical source lines & \(32\,996\),\\
SHA--256 &
\texttt{1B97BE72E2AB16377E0B1ADC4539A231A177A194436CA7EAB27F1DB13371B622}.
\end{tabular}
\end{center}
The first concurrent read saw the completed V68 body while the V69 filename
was being replaced.  It was discarded.  The displayed snapshot contains the
V69 end marker and parent record.

The settled Navier--Stokes snapshot was
\begin{center}
\begin{tabular}{ll}
file &
\texttt{Renewal\_NS\_Stability\_Research\_Notes\_V29.tex},\\
bytes & \(372\,742\),\\
logical source lines & \(7\,327\),\\
SHA--256 &
\texttt{3BE7398455BA117FC89A243639ECDAD037FEDDF5C91FAB969547C7C65598C2A6}.
\end{tabular}
\end{center}

\begin{audit}[Transfer at V95]
\label{audit:newV95-transfer}
Yang--Mills V69 finds an unabsorbed superexponential factor in the earlier
V66 closure, retracts that closure, and repairs it by a global heat-semigroup
representation and one joint partition--derivative generating function.  It
then closes the one-complete-centre tangent row without separating the
factors whose cancellation is needed.  Navier--Stokes V29 derives the exact
equation for a profile divided by its Gaussian norm and proves that this
normalization can destroy the inherited \(H^1\) and \(L^6\) bounds when the
denominator vanishes.

The common lesson is used literally below: a determinant norm or a normalized
phase cannot replace a lower bound on the smallest singular value.  The
kinetic and Yukawa terms are assembled before that singular value is
estimated.  No numerical estimate is imported from either companion note.
\end{audit}

\section{The zero-Higgs obstruction}
\label{sec:newV95-zero-Higgs}

Retain the V93--V94 overlap normalization
\begin{equation}
 D_{\rm ov}(U)={1\over a}(I+\gamma_5\varepsilon(U)),
 \qquad
 \varepsilon(U)=\operatorname{sgn}H_W(U),                  \label{eq:newV95-Dov}
\end{equation}
at \(\rho=1\).  Assume periodic spin structure for the moment.  Let
\(U=1\) and denote the zero-momentum spinor subspace by \(E_0\).

\begin{lemma}[The free overlap zero mode]
\label{lem:newV95-free-zero}
On \(E_0\),
\begin{equation}
 D_W(1)=0,
 \qquad
 H_W(1)=-{1\over a}\gamma_5,
 \qquad
 \varepsilon(1)=-\gamma_5,
 \qquad
 D_{\rm ov}(1)=0.                                          \label{eq:newV95-free-zero}
\end{equation}
\end{lemma}

\begin{proof}
At zero momentum every shift is the identity.  Both the antisymmetric
difference and the Wilson term in \eqref{eq:newV93-DW} vanish.  The remaining
three identities follow successively from
\(H_W=\gamma_5(D_W-a^{-1}I)\), functional calculus, and
\eqref{eq:newV95-Dov}.
\end{proof}

Let \(\mathscr Y(H)\) be a finite-cutoff Yukawa map which is linear in the
Higgs field and hence obeys \(\mathscr Y(0)=0\).  On an index-zero component,
write the chiral block as
\begin{equation}
 D_\chi(U,H)
 =P_+D_{\rm ov}(U)\widehat P_-(U)+\mathscr Y(H).             \label{eq:newV95-Dchi}
\end{equation}

\begin{theorem}[No uniform chiral floor on a Higgs ball containing zero]
\label{thm:newV95-zero-Higgs-no-go}
Suppose a protected field set contains the configuration \((U,H)=(1,0)\)
and contains a nonzero vector in the zero-momentum source space on which the
Yukawa map vanishes.  Then
\begin{equation}
 \inf_{(U,H)}s_{\min}(D_\chi(U,H))=0.                       \label{eq:newV95-zero-floor}
\end{equation}
In particular, compactness of a Higgs ball containing the origin does not
imply the positive chiral-floor hypothesis used by the Grand Tensor field
shell.
\end{theorem}

\begin{proof}
At \((1,0)\), \Cref{lem:newV95-free-zero} makes the kinetic term vanish on
the zero-momentum subspace, while \(\mathscr Y(0)=0\).  By hypothesis there
is a nonzero source vector annihilated by the full chiral block.  Its least
singular value is zero at that configuration, so the infimum over any set
containing it is zero.
\end{proof}

\begin{corollary}[Changing the spin structure does not give a volume floor]
\label{cor:newV95-spin-structure}
An antiperiodic direction removes the exact finite-volume zero momentum, but
does not produce a positive lower bound uniform in physical side length
\(L=Na\).  Along free lattices with \(L\to\infty\),
\begin{equation}
 s_{\min}(D_{\rm ov}(1))\le {C\over L}.                     \label{eq:newV95-antiperiodic}
\end{equation}
\end{corollary}

\begin{proof}
The least antiperiodic momentum has magnitude \(\pi/L\).  The free overlap
symbol is smooth near zero and vanishes at zero by
\Cref{lem:newV95-free-zero}.  Its derivative is bounded in a fixed
neighbourhood after physical units are restored, so the mean-value theorem
bounds the corresponding singular value by \(C|k|=C\pi/L\).
\end{proof}

\begin{assessment}[Exact choices forced by the obstruction]
\label{ass:newV95-zero-choices}
A uniform finite-cutoff floor requires at least one of the following typed
operations: retain the massless modes as an exact low Grassmann sector;
restrict the Higgs field to a broken-phase set separated from zero and prove
that every intended Yukawa mass is nonzero; or add an independently justified
infrared regulator and later prove a controlled removal theorem.  Simply
normalizing the determinant or excluding its zero set pointwise does not
supply a uniform floor.
\end{assessment}

\section{The massive overlap operator}
\label{sec:newV95-massive}

For \(0<m\le2/a\), define the standard massive overlap deformation
\begin{equation}
 D_m
 :=\left(1-{am\over2}\right)D_{\rm ov}+mI.                \label{eq:newV95-Dm}
\end{equation}
Put \(V=\gamma_5\varepsilon\).  Since \(\gamma_5\) and \(\varepsilon\)
are self-adjoint unitaries, \(V\) is unitary and
\begin{equation}
 D_{\rm ov}={1\over a}(I+V).                                \label{eq:newV95-Dov-V}
\end{equation}

\begin{lemma}[Overlap spectral circle]
\label{lem:newV95-circle}
The operator \(D_{\rm ov}\) is normal and
\begin{equation}
 \operatorname{spec}(D_{\rm ov})
 \subset\left\{\lambda:\left|\lambda-{1\over a}\right|
 ={1\over a}\right\}.                                     \label{eq:newV95-circle}
\end{equation}
The massive operator is normal and its spectrum lies on the circle with
centre and radius
\begin{equation}
 c_m={1\over a}+{m\over2},
 \qquad
 r_m={1\over a}-{m\over2}.                                 \label{eq:newV95-mass-circle}
\end{equation}
\end{lemma}

\begin{proof}
Equation \eqref{eq:newV95-Dov-V} is an affine polynomial in the unitary
operator \(V\), so it is normal and its spectral image is the stated circle.
Substitution in \eqref{eq:newV95-Dm} gives
\[
 D_m=c_mI+r_mV.
\]
Because \(0<m\le2/a\), one has \(r_m\ge0\), proving the second circle
statement.
\end{proof}

\begin{theorem}[Exact massive-overlap singular-value floor]
\label{thm:newV95-mass-floor}
For every link field on which the overlap sign is defined and every
\(0<m\le2/a\),
\begin{equation}
 s_{\min}(D_m)\ge m,
 \qquad
 \|D_m^{-1}\|\le{1\over m}.                                \label{eq:newV95-mass-floor}
\end{equation}
The constants are uniform in lattice volume and gauge configuration and are
sharp on every class containing the free periodic field.
\end{theorem}

\begin{proof}
For a normal operator, singular values are the moduli of its spectral
values.  The distance from the origin to the circle in
\eqref{eq:newV95-mass-circle} is
\[
 c_m-r_m=m.
\]
Hence every spectral value has modulus at least \(m\), which proves both
inequalities in \eqref{eq:newV95-mass-floor}.  The spectral circle of a finite
unitary \(V\) need not contain its leftmost point, so equality is not asserted
for each individual configuration.  On the free periodic branch from
\Cref{lem:newV95-free-zero}, \(D_{\rm ov}\) has a zero mode and \(D_m\) acts
there as multiplication by \(m\).  Thus both constants are sharp uniformly.
\end{proof}

\begin{proposition}[Positivity of the vectorlike determinant]
\label{prop:newV95-vectorlike-positive}
For \(m>0\),
\begin{equation}
 \det D_m>0.                                                \label{eq:newV95-vectorlike-positive}
\end{equation}
Thus the massive vectorlike overlap determinant has no residual complex
phase.
\end{proposition}

\begin{proof}
V93 gives \(D_{\rm ov}^*=\gamma_5D_{\rm ov}\gamma_5\), and the same
identity holds for \(D_m\).  Since \(D_m\) is normal, every nonreal
eigenvalue is paired with its complex conjugate with the same multiplicity;
each pair has positive product.  The only possible real points of the circle
\eqref{eq:newV95-mass-circle} are \(m\) and \(2/a\), both positive.  Their
product with the nonreal pairs is strictly positive.
\end{proof}

\begin{assessment}[Phase scope]
\label{ass:newV95-vectorlike-scope}
The positivity in \Cref{prop:newV95-vectorlike-positive} is the operator form
of V92's opposite-orientation pairing.  It applies to a vectorlike Dirac
pair.  A chiral Standard Model determinant remains a section of the moving
line in V93 and still requires EMTPC; the vectorlike determinant cannot be
substituted for it.
\end{assessment}

\section{Perturbative tubes around an assembled massive block}
\label{sec:newV95-tube}

\begin{lemma}[Sharp singular-value perturbation]
\label{lem:newV95-Weyl}
For finite matrices \(B\) and \(E\),
\begin{equation}
 s_{\min}(B+E)\ge s_{\min}(B)-\|E\|.                       \label{eq:newV95-Weyl}
\end{equation}
\end{lemma}

\begin{proof}
For every unit vector \(v\),
\[
 \|(B+E)v\|\ge\|Bv\|-\|Ev\|
 \ge s_{\min}(B)-\|E\|.
\]
Take the infimum over \(v\).
\end{proof}

Let \(D_{\mathbf m}=\bigoplus_{f=1}^F D_{m_f}\) be an assembled
vectorlike flavour block with
\begin{equation}
 m_*:=\min_fm_f>0.                                          \label{eq:newV95-mstar}
\end{equation}
Let \(E(U,H)\) contain the deviation of a gauge--Higgs--Yukawa block from
this reference, after source and target spaces have been identified by the
same admissible spectral chart.

\begin{theorem}[Protected massive-overlap tube]
\label{thm:newV95-protected-tube}
If, on a field set \(\mathscr T_\theta\),
\begin{equation}
 \|E(U,H)\|\le\theta m_*,
 \qquad 0\le\theta<1,                                      \label{eq:newV95-tube-condition}
\end{equation}
then
\begin{equation}
 s_{\min}(D_{\mathbf m}+E(U,H))
 \ge(1-\theta)m_*                                          \label{eq:newV95-tube-floor}
\end{equation}
uniformly on \(\mathscr T_\theta\).
\end{theorem}

\begin{proof}
The direct sum and \Cref{thm:newV95-mass-floor} give
\(s_{\min}(D_{\mathbf m})\ge m_*\).  Apply
\Cref{lem:newV95-Weyl} and \eqref{eq:newV95-tube-condition} to the complete
assembled block.
\end{proof}

\begin{corollary}[Resolvent and determinant response on the tube]
\label{cor:newV95-tube-resolvent}
On \(\mathscr T_\theta\),
\begin{equation}
 \|(D_{\mathbf m}+E)^{-1}\|
 \le{1\over(1-\theta)m_*}.                                 \label{eq:newV95-tube-resolvent}
\end{equation}
Every finite derivative of the determinant logarithm is therefore bounded
by the corresponding product of operator derivatives and powers of
\(((1-\theta)m_*)^{-1}\), with the trace taken only after the complete
cyclic product has been formed.
\end{corollary}

\begin{proof}
The inverse norm is the reciprocal of the least singular value.  Jacobi's
formula and repeated differentiation express each logarithmic derivative as
a finite sum of traces of cyclic products of inverse and differentiated
blocks.  Insert \eqref{eq:newV95-tube-resolvent} only after each product is
assembled, in accordance with the V95 audit.
\end{proof}

\section{Standard Model mass rank}
\label{sec:newV95-SM-rank}

On a constant broken Higgs background, singular-value decomposition of each
Yukawa matrix identifies its nonzero singular values with vectorlike masses.
The previous theorem then applies to the direct sum of those massive pairs,
provided the actual overlap--Yukawa discretization is unitarily equivalent to
the standard massive form plus the declared perturbation \(E\).

\begin{proposition}[Massless rows cannot enter the positive tube]
\label{prop:newV95-massless-row}
If the assembled Yukawa mass matrix has a zero singular value, then
\(m_*=0\) and \Cref{thm:newV95-protected-tube} supplies no positive floor.
For the minimal one-generation Standard Model without a right-handed
neutrino, the neutrino row is of this type.  It must be retained as a
massless Grassmann sector or completed by an independently specified
right-handed-neutrino/Yukawa block before a full massive tube is asserted.
\end{proposition}

\begin{proof}
A zero singular value leaves a flavour vector on which the constant Yukawa
mass vanishes.  At the free periodic field the overlap kinetic operator also
vanishes at zero momentum by \Cref{lem:newV95-free-zero}.  Thus that flavour
has no positive uniform mass floor.  The representation list in V38 contains
no right-handed neutrino in its minimal form; the optional sterile conjugate
is recorded separately there.
\end{proof}

\begin{definition}[Higgs--Yukawa floor card, HYFC]
\label{def:newV95-HYFC}
HYFC consists of:
\begin{enumerate}[label=\textup{(\roman*)},leftmargin=3.5em]
\item an AOP admissible gauge component;
\item a declared decomposition into massive vectorlike pairs and retained
      massless chiral rows;
\item a positive minimum reference mass \(m_*\) on the massive block;
\item an exact common source/target identification on the moving overlap
      projectors;
\item the assembled perturbation bound
      \(\|E(U,H)\|\le\theta m_*\), \(\theta<1\); and
\item exact Grassmann retention of every excluded zero or crossing row.
\end{enumerate}
\end{definition}

\begin{theorem}[HYFC closes the massive local singular-value gate]
\label{thm:newV95-HYFC}
If HYFC holds, the massive block has the uniform floor
\((1-\theta)m_*\), bounded inverse
\eqref{eq:newV95-tube-resolvent}, and a nonvanishing vectorlike determinant.
The retained massless block and the chiral determinant line remain separate
outputs and are not integrated into that positive determinant.
\end{theorem}

\begin{proof}
Use \Cref{thm:newV95-protected-tube,
cor:newV95-tube-resolvent,
prop:newV95-vectorlike-positive}.  Item \textup{(vi)} prevents a zero row
from being hidden in the determinant of the massive complement.
\end{proof}

\section{Next acquisition step}
\label{sec:newV95-next}

V96 will determine how a nonempty HYFC tube can enter the field measure.
It will distinguish hard restriction to the admissible/broken chart from a
probabilistic bad-set estimate under the gauge--Higgs action.  The immediate
question is whether the Wilson plaquette action and Higgs potential can give
a cutoff-uniform probability bound for the complements of
\(\mathscr A_\delta\) and \(\mathscr T_\theta\) without assuming the desired
continuum measure.  If no such bound is available, those charts will remain
conditional sectors and the programme will move upstream to a modified
admissibility-enforcing lattice action.

\begin{firewall}[Claim boundary after V95]
\label{fire:newV95-boundary}
V95 performs the compulsory companion audit, proves that any field set
containing the trivial gauge field and zero Higgs configuration fails the
uniform chiral floor, constructs the massive overlap operator, proves its
sharp uniform singular floor and vectorlike determinant positivity, and
closes the local massive-block gate on an explicit perturbative tube.

V95 does not prove that the Standard Model gauge--Higgs measure concentrates
on that tube, give a mass to the minimal-model neutrino, trivialize the chiral
determinant line, prove the mapping-torus gate, establish reflection
positivity, remove the lattice cutoff, or construct the full SM--Einstein
continuum.  Equation \eqref{eq:newV95-mass-floor} gives a uniform
lower bound; equality is claimed only as a sharp class-wide constant when the
class contains the free periodic field.
\end{firewall}

\begin{theorem}[V95 result ledger]
\label{thm:newV95-results}
V95:
\begin{enumerate}[label=\textup{(V95.\arabic*)},leftmargin=6.5em]
\item records settled, hashed YM V69 and NS V29 snapshots;
\item imports their assemble-before-estimating and vanishing-normalization
      warnings without importing numerical bounds;
\item proves the free periodic overlap zero-mode identity;
\item proves the zero-Higgs no-floor theorem;
\item proves that an antiperiodic spin choice has no volume-uniform floor;
\item derives the exact massive-overlap spectral circle;
\item proves the sharp configuration-uniform lower singular-value bound;
\item proves positivity of the vectorlike massive determinant;
\item proves a uniform floor on an assembled massive perturbation tube;
\item isolates the minimal-model massless neutrino row; and
\item formulates HYFC and closes its massive local singular-value gate.
\end{enumerate}
\end{theorem}

\begin{proof}
Use \Cref{audit:newV95-transfer,
lem:newV95-free-zero,
thm:newV95-zero-Higgs-no-go,
cor:newV95-spin-structure,
lem:newV95-circle,
thm:newV95-mass-floor,
prop:newV95-vectorlike-positive,
thm:newV95-protected-tube,
prop:newV95-massless-row,
thm:newV95-HYFC}.
\end{proof}

\paragraph{Parent-version record.}
V95 was formed from
\texttt{Renewal\_SM\_Einstein\_Continuum\_Research\_Notes\_V94.tex}.
The V94 parent had \(32\,941\) logical source lines,
\(1\,264\,229\) bytes, and SHA--256
\[
 \texttt{87CE0AE7BDD5645BC5CB1579AAB1511C4A2E52C59D984C1C9BFAA8D6AC771003}.
\]
The next compulsory companion audit is V100.

% =============================================================================
% END APPENDED SECOND-SERIES V95 MATERIAL
% =============================================================================


% =============================================================================
% BEGIN APPENDED SECOND-SERIES V96 MATERIAL
% =============================================================================

\clearpage
\chapter{Second-series V96: finite-action support, a positive-type
admissibility weight, and the local-norm bridge}
\label{sec:v2v96}

\section{What an ordinary finite action can and cannot enforce}
\label{sec:newV96-soft-support}

Let \(G\) be the compact Standard Model gauge group, with normalized Haar
measure, and let \(\Lambda\) be a finite lattice.  The bosonic reference
measure is the product of Haar measure on \(G^{E(\Lambda)}\) and Lebesgue
measure on the finite-dimensional Higgs field space.  The next statement
applies to the Wilson plaquette action and the usual polynomial
gauge--Higgs action whenever their finite-volume partition function exists.

\begin{theorem}[Finite continuous actions have full bosonic support]
\label{thm:newV96-full-support}
Let \(S(U,H)\) be finite and continuous at every finite bosonic
configuration, and suppose
\[
 Z=\int\exp\{-S(U,H)\}\,dU\,dH
\]
is finite and nonzero.  The probability measure
\[
 d\mu_S=Z^{-1}\exp\{-S(U,H)\}\,dU\,dH
\]
assigns positive probability to every nonempty open subset of bosonic
configuration space.
\end{theorem}

\begin{proof}
The product Haar--Lebesgue reference measure has full support.  Hence every
nonempty open set \(O\) has positive reference measure.  The integrand
\(\exp(-S)\) is strictly positive at every point, so its integral over \(O\)
is positive.  Division by \(Z\) proves the claim.
\end{proof}

For the finitely many unitary representations carried by Standard Model
fermions, put
\begin{equation}
 d_{\rm SM}(g):=\max_{R\in\mathcal R_{\rm SM}}
                  \|I-R(g)\|.                              \label{eq:newV96-dSM}
\end{equation}
This is continuous, conjugation invariant, and symmetric under inversion.
Fix \(0<\delta_0<1/30\), and define
\begin{equation}
 \mathscr A_{\delta_0}
 :=\{U:\max_p d_{\rm SM}(U_p)<\delta_0\}.                  \label{eq:newV96-Adelta}
\end{equation}

\begin{corollary}[Soft Wilson--Higgs weights do not give exact AOP or HYFC
support]
\label{cor:newV96-soft-no-exact}
Assume the represented group contains an element with
\(d_{\rm SM}(g)>\delta_0\), and the finite lattice contains a plaquette.
Then the ordinary finite Wilson--Higgs measure gives positive probability
to \(\mathscr A_{\delta_0}^{\,c}\).

Suppose also that the intended chiral block is continuous in the
finite-dimensional bosonic variables and vanishes on a nonzero source vector
at \((U,H)=(1,0)\), as in \Cref{thm:newV95-zero-Higgs-no-go}.  For every
\(\kappa>0\), the set
\[
 \{(U,H):s_{\min}(D_\chi(U,H))<\kappa\}
\]
has positive probability.  Thus neither AOP nor a positive HYFC floor is an
almost-sure consequence of a finite continuous Wilson--Higgs action at a
fixed cutoff.
\end{corollary}

\begin{proof}
Assign the chosen group element to one link and the identity to the
remaining links.  At least one adjacent plaquette then has defect exceeding
\(\delta_0\).  By continuity, an open neighborhood of that configuration is
contained in \(\mathscr A_{\delta_0}^{\,c}\).  Apply
\Cref{thm:newV96-full-support}.

The least singular value of a finite matrix is continuous in its entries.
It is zero at \((1,0)\) on the declared source vector, so a neighborhood of
that point has least singular value below \(\kappa\).  Apply the same
full-support theorem.
\end{proof}

\begin{assessment}[The support result is not a continuum no-go]
\label{ass:newV96-not-continuum-no-go}
Positive bad-set probability at each finite cutoff does not prevent that
probability from tending to zero along a tuned cutoff sequence.  In
particular, \Cref{cor:newV96-soft-no-exact} does not disprove concentration
of the Wilson action as the bare coupling changes.  It proves only that an
exact use of AOP or HYFC needs either a hard-support measure or a quantitative
bad-set estimate.
\end{assessment}

\section{The quantitative conditioning bridge}
\label{sec:newV96-conditioning}

\begin{lemma}[Cost of conditioning on a protected set]
\label{lem:newV96-conditioning}
Let \(\mu\) be a probability measure, let \(G\) be measurable with
\(p:=\mu(G^c)<1\), and let
\(\mu_G(\,\cdot\,)=\mu(\,\cdot\,\mid G)\).  For every bounded observable
\(F\),
\begin{equation}
 \left|\mathbb E_\mu F-\mathbb E_{\mu_G}F\right|
 \le2\|F\|_\infty p.                                      \label{eq:newV96-conditioning}
\end{equation}
\end{lemma}

\begin{proof}
Write \(B=G^c\).  The total expectation formula gives
\[
 \mathbb E_\mu F
 =(1-p)\mathbb E_{\mu_G}F+p\mathbb E_{\mu(\,\cdot\,|B)}F.
\]
Subtract \(\mathbb E_{\mu_G}F\) and use the triangle inequality.
\end{proof}

\begin{corollary}[A precise soft-to-hard universality obligation]
\label{cor:newV96-soft-hard}
Let \(G_a\) be the intersection of AOP and a declared HYFC tube at cutoff
\(a\).  If
\begin{equation}
 \mu_a(G_a^c)\longrightarrow0,                             \label{eq:newV96-bad-prob}
\end{equation}
then the original and conditioned measures have the same limit for every
uniformly bounded observable whose original expectation has a limit.
For unbounded observables, the additional required input is uniform
integrability; it is not supplied by \eqref{eq:newV96-bad-prob}.
\end{corollary}

\begin{proof}
Apply \Cref{lem:newV96-conditioning} at each cutoff.  The right side tends to
zero for uniformly bounded observables.  Truncation extends the argument to
a uniformly integrable family.
\end{proof}

\begin{assessment}[Two distinct routes]
\label{ass:newV96-two-routes}
Equation \eqref{eq:newV96-bad-prob} would justify conditioning the original
Wilson--Higgs measure.  A new hard-support action can instead make AOP exact
at every cutoff, but then agreement with the target continuum is a separate
universality theorem.  These routes must not be conflated.
\end{assessment}

\section{A positive-type plaquette weight with admissible support}
\label{sec:newV96-positive-type}

The naive indicator of \(\mathscr A_{\delta_0}\) gives hard support but does
not by itself display the character positivity used in the standard
lattice reflection argument.  The following autocorrelation construction
supplies both properties.

\begin{theorem}[Positive-type admissibility bump]
\label{thm:newV96-positive-bump}
There exists a nonzero smooth central function
\(w_{\delta_0}:G\to[0,\infty)\) such that
\begin{enumerate}[label=\textup{(\roman*)},leftmargin=3.5em]
\item
\(\operatorname{supp}w_{\delta_0}\subset
 \{g:d_{\rm SM}(g)<\delta_0\}\);
\item \(w_{\delta_0}(g^{-1})=w_{\delta_0}(g)\);
\item \(w_{\delta_0}\) is of positive type; and
\item its Peter--Weyl character expansion is
\begin{equation}
 w_{\delta_0}(g)=\sum_{\pi\in\widehat G}
                   c_\pi\chi_\pi(g),
 \qquad c_\pi\ge0,                                         \label{eq:newV96-character}
\end{equation}
with smooth convergence.
\end{enumerate}
\end{theorem}

\begin{proof}
Continuity of multiplication at the identity gives a conjugation-invariant,
symmetric open neighborhood \(O\) of the identity such that
\[
 \overline O\,\overline O^{-1}
 \subset\{g:d_{\rm SM}(g)<\delta_0\}.
\]
Choose a nonzero, nonnegative, smooth central bump \(f\) supported in \(O\).
Such a bump is obtained from a bump near the identity by averaging over
conjugation.  With normalized Haar measure, set
\begin{equation}
 w_{\delta_0}(g)
 :=\int_G f(h)f(g^{-1}h)\,dh
 =(f*\widetilde f)(g),
 \qquad \widetilde f(g):=f(g^{-1}).                        \label{eq:newV96-autocorrelation}
\end{equation}
The integrand is nonnegative and
\(w_{\delta_0}(e)=\int f^2>0\).  If the integrand is nonzero, then
\(h\in\overline O\) and \(g^{-1}h\in\overline O\), hence
\(g\in\overline O\,\overline O^{-1}\); this proves the support statement.
Centrality and inversion symmetry follow from Haar invariance and centrality
of \(f\).

Convolution turns each Fourier coefficient into
\[
 \widehat w_{\delta_0}(\pi)
 =\widehat f(\pi)\widehat f(\pi)^*.
\]
Because \(f\) is central, Schur's lemma makes
\(\widehat f(\pi)\) a scalar matrix.  Consequently every coefficient in
\eqref{eq:newV96-character} is a nonnegative scalar multiple of its squared
modulus.  This proves positive type and \(c_\pi\ge0\).  Smoothness of \(f\)
gives rapid Fourier decay and therefore smooth convergence.
\end{proof}

Define the pure-gauge finite-lattice measure
\begin{equation}
 d\nu_{\delta_0}(U)
 ={1\over Z_{\delta_0}}
  \prod_p w_{\delta_0}(U_p)\prod_{\ell}dU_\ell.            \label{eq:newV96-hard-measure}
\end{equation}

\begin{theorem}[Exact AOP support with pure-gauge reflection positivity]
\label{thm:newV96-hard-AOP-RP}
The partition function in \eqref{eq:newV96-hard-measure} is finite and
strictly positive.  Every configuration with nonzero density lies in
\(\mathscr A_{\delta_0}\), so all AOP conclusions of
\Cref{thm:newV94-AOP} hold \(\nu_{\delta_0}\)-almost surely with
volume-independent constants.  Moreover, on a reflection-symmetric
hypercubic lattice, the pure-gauge measure
\eqref{eq:newV96-hard-measure} obeys the standard link-reflection positivity
inequality
\begin{equation}
 \int \overline{F(\vartheta U)}F(U)\,d\nu_{\delta_0}(U)
 \ge0                                                       \label{eq:newV96-RP}
\end{equation}
for every polynomial cylinder function \(F\) supported in the positive
half-lattice.
\end{theorem}

\begin{proof}
The density is bounded on the compact link space, so the partition function
is finite.  Since \(w_{\delta_0}(e)>0\), continuity supplies a link
neighborhood of the identity on which every plaquette factor is positive;
this set has positive Haar measure, proving \(Z_{\delta_0}>0\).
The support inclusion in \Cref{thm:newV96-positive-bump} makes every
represented plaquette defect smaller than \(\delta_0<1/30\), proving exact
AOP support.

For reflection positivity, split every plaquette crossing the reflection
plane into negative and positive half-transporters.  Insert
\eqref{eq:newV96-character}.  The representation identity
\[
 \chi_\pi(g_-^{-1}g_+)
 =\sum_{i,j}\overline{\pi(g_-)_{ij}}\pi(g_+)_{ij}
\]
writes a crossing factor as a sum of reflected squares with coefficients
\(c_\pi\ge0\).  Products of crossing factors retain this form.  Plaquettes
strictly inside the two halves are exchanged by reflection, and Haar
integration over reflected link pairs turns the quadratic form into a sum
of squared moduli.  This proves \eqref{eq:newV96-RP}; smooth convergence
justifies termwise insertion, and density extends the result to the stated
cylinder functions.
\end{proof}

\begin{assessment}[What the positive bump does not prove]
\label{ass:newV96-bump-scope}
The construction gives a local, gauge-invariant, reflection-positive
pure-gauge regulator with exact AOP support.  It does not prove that this
regulator has the Yang--Mills continuum limit, lies in the Wilson
universality class, admits the full chiral Standard Model measure, or remains
reflection positive after the determinant phase, Higgs coupling, and
gravitational sector are included.
\end{assessment}

\section{From kernel decay to the global HYFC norm}
\label{sec:newV96-Schur}

V95's HYFC uses the global operator norm of the assembled perturbation.
The next lemma turns that norm into a volume-uniform family of kernel
estimates.

\begin{lemma}[Block Schur test]
\label{lem:newV96-block-Schur}
Let \(E=(E(x,y))_{x,y\in\Lambda}\) act between finite direct sums of
finite-dimensional fibres.  If
\begin{equation}
 R:=\sup_x\sum_y\|E(x,y)\|<\infty,
 \qquad
 C:=\sup_y\sum_x\|E(x,y)\|<\infty,                          \label{eq:newV96-row-column}
\end{equation}
then
\begin{equation}
 \|E\|\le\sqrt{RC}.                                        \label{eq:newV96-Schur}
\end{equation}
\end{lemma}

\begin{proof}
For fibre norms \(u_y=\|v_y\|\),
\[
 \|(Ev)_x\|
 \le\sum_y\|E(x,y)\|u_y.
\]
The scalar Schur test applied to the nonnegative matrix
\((\|E(x,y)\|)\) gives its \(\ell^2\) norm at most \(\sqrt{RC}\), and
therefore gives \eqref{eq:newV96-Schur}.
\end{proof}

On a four-dimensional periodic lattice with graph distance, the number of
sites at distance \(r\ge1\) from a fixed site is at most
\(16(r+1)^3\).  Define
\begin{equation}
 K_4(\alpha)
 :=1+16\sum_{r=1}^{\infty}(r+1)^3e^{-\alpha r}<\infty.
                                                               \label{eq:newV96-K4}
\end{equation}

\begin{corollary}[Exponential local-block criterion for HYFC]
\label{cor:newV96-local-HYFC}
Suppose the fully assembled perturbation in V95 satisfies
\begin{equation}
 \|E(U,H;x,y)\|\le C_Ee^{-\alpha d_{\rm lat}(x,y)}
                                                               \label{eq:newV96-E-decay}
\end{equation}
throughout a field chart, with \(C_E,\alpha>0\) independent of volume.
Then
\begin{equation}
 \|E(U,H)\|\le C_EK_4(\alpha).                              \label{eq:newV96-E-global}
\end{equation}
If
\begin{equation}
 C_EK_4(\alpha)\le\theta m_*,
 \qquad \theta<1,                                          \label{eq:newV96-local-card}
\end{equation}
the HYFC singular floor and resolvent bounds of V95 hold throughout that
chart.
\end{corollary}

\begin{proof}
Sum \eqref{eq:newV96-E-decay} over distance shells.  The same bound applies
to rows and columns, so
\(R,C\le C_EK_4(\alpha)\).  Apply
\Cref{lem:newV96-block-Schur} and then
\Cref{thm:newV95-protected-tube}.
\end{proof}

\begin{definition}[Exponential local-block acquisition card, ELBAC]
\label{def:newV96-ELBAC}
ELBAC asks for a common admissible spectral chart, an exact kernel formula
for the fully assembled overlap--Yukawa perturbation, and constants
\(C_E,\alpha\) satisfying
\eqref{eq:newV96-E-decay}--\eqref{eq:newV96-local-card} in every intended
massive Standard Model flavour block.  The massless rows remain the separate
Grassmann sector required by HYFC.
\end{definition}

\begin{assessment}[Why sign locality alone is insufficient]
\label{ass:newV96-sign-not-enough}
V94 bounds the kernel of the overlap sign itself.  ELBAC instead bounds the
difference between the complete moving-projector Yukawa block and a fixed
massive reference after source and target fibres have been identified.
The former does not automatically imply the latter: the chart
identification, Higgs fluctuation, flavour matrix, and all projector
differences must enter one assembled kernel before
\Cref{lem:newV96-block-Schur} is applied.
\end{assessment}

\section{Next acquisition step}
\label{sec:newV96-next}

V97 will attack ELBAC.  It will first derive a quantitative Lipschitz bound
for the sign function from the V94 Wilson gap, then transport that bound to
the moving overlap projectors and assemble it with an ultralocal Higgs
fluctuation.  If a gauge-invariant local field chart cannot control the
needed kernel constant without choosing a global gauge, the programme will
move upstream to a covariant projector transport rather than hide the
dependence in \(\|E\|\).

\begin{firewall}[Claim boundary after V96]
\label{fire:newV96-boundary}
V96 proves full support for ordinary finite continuous Wilson--Higgs
densities and therefore excludes exact finite-cutoff AOP or positive-floor
support from such densities.  It gives the precise bad-probability condition
under which conditioning is harmless for bounded observables.  It constructs
a local central positive-type plaquette weight with support inside the V94
admissibility ball and proves pure-gauge link-reflection positivity for the
resulting hard-support regulator.  It also proves a volume-uniform block
Schur bridge from exponential kernel control to HYFC.

V96 does not prove Wilson-action concentration, universality of the new
positive-type regulator, ELBAC for the actual Standard Model
overlap--Yukawa block, positivity of the chiral determinant measure,
reflection positivity after matter or gravity is added, cutoff removal, or
the full SM--Einstein continuum.
\end{firewall}

\begin{theorem}[V96 result ledger]
\label{thm:newV96-results}
V96:
\begin{enumerate}[label=\textup{(V96.\arabic*)},leftmargin=6.5em]
\item proves the finite-action full-support theorem;
\item proves that ordinary soft Wilson--Higgs weights do not exactly enforce
      AOP or HYFC at fixed cutoff;
\item isolates the bad-set probability and uniform-integrability
      obligations for a soft-to-hard comparison;
\item constructs a smooth positive-type plaquette bump supported strictly
      inside the common Standard Model admissibility ball;
\item proves exact AOP support and pure-gauge link-reflection positivity for
      its product plaquette measure;
\item proves a volume-uniform block Schur estimate; and
\item reduces the global HYFC perturbation bound to ELBAC.
\end{enumerate}
\end{theorem}

\begin{proof}
Use \Cref{thm:newV96-full-support,
cor:newV96-soft-no-exact,
lem:newV96-conditioning,
cor:newV96-soft-hard,
thm:newV96-positive-bump,
thm:newV96-hard-AOP-RP,
lem:newV96-block-Schur,
cor:newV96-local-HYFC}.
\end{proof}

\paragraph{Parent-version record.}
V96 was formed from
\texttt{Renewal\_SM\_Einstein\_Continuum\_Research\_Notes\_V95.tex}.
The V95 parent had \(33\,397\) logical source lines,
\(1\,282\,339\) bytes, and SHA--256
\[
 \texttt{B67EF00B3A7B43348C3DA97752B123C910ECEFDD2E726EE79079C0B4A310C808}.
\]
The next compulsory companion audit is V100.

% =============================================================================
% END APPENDED SECOND-SERIES V96 MATERIAL
% =============================================================================

% =============================================================================
% BEGIN APPENDED SECOND-SERIES V97 MATERIAL
% =============================================================================

\clearpage
\chapter{Second-series V97: sign Lipschitz transport and the
fibrewise vectorlike split}
\label{sec:v2v97}

\section{A sharp gap-to-sign perturbation estimate}
\label{sec:newV97-sign-Lipschitz}

\begin{lemma}[Resolvent formula for a sign difference]
\label{lem:newV97-sign-resolvent}
Let \(H\) and \(K\) be finite-dimensional self-adjoint operators with
\[
 H^2\succeq g^2I,\qquad K^2\succeq g^2I
\]
for some \(g>0\).  Then
\begin{equation}
 \|\operatorname{sgn}H-\operatorname{sgn}K\|
 \le {1\over g}\|H-K\|.                                   \label{eq:newV97-sign-Lipschitz}
\end{equation}
\end{lemma}

\begin{proof}
Spectral calculus gives the norm-convergent difference formula
\begin{align}
 \operatorname{sgn}H-\operatorname{sgn}K
 ={1\over\pi}\int_{-\infty}^{\infty}
 \big\{(H-it)^{-1}-(K-it)^{-1}\big\}\,dt.                 \label{eq:newV97-sign-integral}
\end{align}
Indeed, the corresponding scalar integral is
\(\pi^{-1}\int_{\mathbb R}(\lambda-it)^{-1}dt
=\operatorname{sgn}\lambda\).  The resolvent identity and self-adjointness
give
\[
 \|(H-it)^{-1}-(K-it)^{-1}\|
 \le{\|H-K\|\over g^2+t^2}.
\]
Integration uses
\(\pi^{-1}\int_{\mathbb R}(g^2+t^2)^{-1}dt=g^{-1}\), proving
\eqref{eq:newV97-sign-Lipschitz}.
\end{proof}

\begin{assessment}[No hidden spectral derivative]
\label{ass:newV97-no-hidden-derivative}
The proof estimates the difference of the two resolvents before integrating.
The individual resolvent integrals are not absolutely norm integrable, while
their difference is.  Thus \eqref{eq:newV97-sign-Lipschitz} does not rely on
differentiating eigenvectors or choosing eigenvalue labels.
\end{assessment}

\section{The explicit Wilson-link constant}
\label{sec:newV97-Wilson-link}

For two link fields \(U,V\) in the same unitary representation, put
\begin{equation}
 \epsilon_\ell(U,V)
 :=\max_{x,\mu}\|U_\mu(x)-V_\mu(x)\|.                      \label{eq:newV97-link-distance}
\end{equation}
This number is not asserted to be gauge invariant; its role is to quantify a
single chosen local chart.

\begin{lemma}[Wilson-kernel link Lipschitz bound]
\label{lem:newV97-Wilson-link}
At \(\rho=1\),
\begin{equation}
 \|H_W(U)-H_W(V)\|
 \le {4\over a}\epsilon_\ell(U,V).                         \label{eq:newV97-HW-link}
\end{equation}
\end{lemma}

\begin{proof}
Use the shift identity \eqref{eq:newV94-A-shift},
\[
 A(U)=-3I+\sum_{\mu=1}^4
 \{P_\mu^-T_\mu(U)+P_\mu^+T_\mu(U)^*\}.
\]
For fixed \(\mu\), the two spin projectors are orthogonal and commute with
the colour shifts.  Therefore
\[
 \|P_\mu^-(T_\mu(U)-T_\mu(V))
   +P_\mu^+(T_\mu(U)^*-T_\mu(V)^*)\|
 \le\epsilon_\ell(U,V).
\]
Sum over four directions.  Since
\(H_W=\gamma_5(D_W-a^{-1}I)\) and \(D_W=a^{-1}(I-A)\), multiplication by
\(\gamma_5\) does not change the norm and
\eqref{eq:newV97-HW-link} follows.
\end{proof}

Put
\begin{equation}
 c_\delta:=\sqrt{1-30\delta}.                              \label{eq:newV97-cdelta}
\end{equation}

\begin{theorem}[Overlap and moving-projector chart bounds]
\label{thm:newV97-overlap-chart}
If \(U,V\in\mathscr A_\delta\) with \(0\le\delta<1/30\), then
\begin{align}
 \|\varepsilon(U)-\varepsilon(V)\|
 &\le {4\over c_\delta}\epsilon_\ell(U,V),                 \label{eq:newV97-sign-link}\\
 \|D_{\rm ov}(U)-D_{\rm ov}(V)\|
 &\le {4\over ac_\delta}\epsilon_\ell(U,V),                \label{eq:newV97-Dov-link}\\
 \|\widehat P_\pm(U)-\widehat P_\pm(V)\|
 &\le {2\over c_\delta}\epsilon_\ell(U,V).                 \label{eq:newV97-projector-link}
\end{align}
All constants are independent of lattice volume.
\end{theorem}

\begin{proof}
V94 gives the common gap \(g=c_\delta/a\).  Combine
\Cref{lem:newV97-sign-resolvent,lem:newV97-Wilson-link} for the first
bound.  The definitions
\[
 D_{\rm ov}=a^{-1}(I+\gamma_5\varepsilon),
 \qquad
 \widehat P_\pm={1\mp\varepsilon\over2}
\]
give the remaining two.
\end{proof}

\section{Canonical transport of the moving source fibre}
\label{sec:newV97-Kato}

\begin{lemma}[Kato--Nagy projector transport with an explicit norm bound]
\label{lem:newV97-Kato}
Let \(P,Q\) be orthogonal projections and let
\(d=\|P-Q\|<1\).  Define
\begin{equation}
 \mathcal U(Q,P)
 :=\{QP+(I-Q)(I-P)\}\{I-(P-Q)^2\}^{-1/2}.                 \label{eq:newV97-Kato}
\end{equation}
Then \(\mathcal U(Q,P)\) is unitary,
\[
 \mathcal U(Q,P)P\mathcal U(Q,P)^*=Q,
\]
and
\begin{equation}
 \|\mathcal U(Q,P)-I\|
 \le\omega(d):={1+d\over\sqrt{1-d^2}}-1.                  \label{eq:newV97-Kato-bound}
\end{equation}
\end{lemma}

\begin{proof}
Set \(B=QP+(I-Q)(I-P)\).  Direct multiplication gives
\[
 B^*B=I-(P-Q)^2.
\]
The right side is positive and invertible because \(d<1\), so
\eqref{eq:newV97-Kato} is the polar unitary of \(B\).  On
\(\operatorname{Ran}P\), \(B\) takes values in
\(\operatorname{Ran}Q\), and on the orthogonal complement it takes values in
the orthogonal complement of \(Q\).  The positive factor
\(\{I-(P-Q)^2\}^{-1/2}\) commutes with \(P\), which proves the intertwining
identity.

Moreover,
\[
 B-I=(2Q-I)(P-Q),
 \qquad
 \|B-I\|\le d,
\]
and
\[
 \|\{I-(P-Q)^2\}^{-1/2}\|\le(1-d^2)^{-1/2}.
\]
Writing \(\mathcal U-I=(B-I)C+(C-I)\), where
\(C=\{I-(P-Q)^2\}^{-1/2}\), gives
\[
 \|\mathcal U-I\|
 \le {d\over\sqrt{1-d^2}}
      +{1\over\sqrt{1-d^2}}-1
 =\omega(d).
\]
\end{proof}

\begin{corollary}[Overlap source transport]
\label{cor:newV97-overlap-transport}
Under the hypotheses of \Cref{thm:newV97-overlap-chart}, if
\begin{equation}
 d_{UV}:={2\epsilon_\ell(U,V)\over c_\delta}<1,            \label{eq:newV97-dUV}
\end{equation}
then \eqref{eq:newV97-Kato}, with
\(P=\widehat P_-(V)\) and \(Q=\widehat P_-(U)\), canonically identifies
their source spaces and differs from the identity by at most
\(\omega(d_{UV})\).
\end{corollary}

\begin{proof}
Apply \Cref{thm:newV97-overlap-chart,lem:newV97-Kato}.
\end{proof}

\section{Why a fixed free reference has the wrong cutoff scaling}
\label{sec:newV97-fixed-reference}

\begin{proposition}[The direct Lipschitz route produces an \(am\)-scale
link chart]
\label{prop:newV97-am-chart}
Let \(m_*>0\) be fixed in physical units.  If one uses
\(D_{\rm ov}(V)\) as a fixed reference and controls the kinetic perturbation
only by \eqref{eq:newV97-Dov-link}, then the sufficient condition
\[
 \|D_{\rm ov}(U)-D_{\rm ov}(V)\|\le\theta m_*
\]
requires
\begin{equation}
 \epsilon_\ell(U,V)
 \le{\theta ac_\delta m_*\over4}.                          \label{eq:newV97-am-radius}
\end{equation}
Thus this proof route shrinks the link-coordinate chart linearly with
\(a\).
\end{proposition}

\begin{proof}
Insert \eqref{eq:newV97-Dov-link} and solve the resulting sufficient
inequality for \(\epsilon_\ell(U,V)\).
\end{proof}

\begin{assessment}[Scope of the shrinking-chart statement]
\label{ass:newV97-am-scope}
Equation \eqref{eq:newV97-am-radius} is a limitation of the displayed norm
estimate, not a no-go theorem for all gauges or all field decompositions.
Smooth continuum links naturally differ from the identity by order \(a\),
but a global gauge with that property is obstructed by flat holonomy and
topological sector data.  The estimate therefore cannot be promoted to a
global admissible-sector argument.
\end{assessment}

\section{The fibrewise vectorlike reference}
\label{sec:newV97-fibrewise}

The massive-overlap theorem of V95 is uniform in the gauge field.  It is
therefore wasteful to perturb an unbroken vectorlike gauge connection back
to the free field.  Let \(K\) denote a proposed unbroken compact subgroup,
and let \(V\) be a \(K\)-valued admissible link field.  For every intended
massive flavour, retain
\begin{equation}
 D_{m_f}(V)
 =\left(1-{am_f\over2}\right)D_{\rm ov}(V)+m_fI           \label{eq:newV97-fibre-base}
\end{equation}
inside the reference block itself.

\begin{definition}[Vectorlike reduction and matching card, VRMC]
\label{def:newV97-VRMC}
VRMC consists of:
\begin{enumerate}[label=\textup{(\roman*)},leftmargin=3.5em]
\item a nonvanishing Higgs chart that reduces the Standard Model gauge
      bundle to an unbroken subgroup \(K\);
\item explicit unitary intertwiners between the left and right restrictions
      to \(K\) for every declared massive flavour;
\item a \(K\)-connection \(V\) and a covariant remainder containing the
      broken-direction gauge variables;
\item a Ginsparg--Wilson-compatible constant Yukawa block which, under the
      intertwiners, is exactly the direct sum of
      \eqref{eq:newV97-fibre-base};
\item retention of the unmatched massless chiral rows as Grassmann
      variables; and
\item a bound on the complete transported remainder relative to the
      fibrewise base.
\end{enumerate}
\end{definition}

\begin{lemma}[Ultralocal Higgs multiplication bound]
\label{lem:newV97-Higgs-bound}
Let \(Y\) be a finite flavour matrix and let \(\eta(x)\) be a Higgs
fluctuation acting by sitewise multiplication.  Then
\begin{equation}
 \|Y\eta\|
 \le\|Y\|\sup_x\|\eta(x)\|.                               \label{eq:newV97-Higgs-bound}
\end{equation}
\end{lemma}

\begin{proof}
At each site,
\(\|Y\eta(x)v(x)\|\le\|Y\|\|\eta(x)\|\|v(x)\|\).
Square, sum over sites, and take the supremum.
\end{proof}

Let the fully transported massive Standard Model block be written, on a
VRMC chart, as
\begin{equation}
 \mathcal D(U,H)
 =\bigoplus_fD_{m_f}(V)+E_{\rm br}+E_H+E_{\rm tr}+E_{\rm rem}.       \label{eq:newV97-assembled}
\end{equation}
Here \(E_{\rm br}\) is the broken-direction kinetic difference, \(E_H\) the
Higgs fluctuation, \(E_{\rm tr}\) the source-fibre transport contribution,
and \(E_{\rm rem}\) a declared discretization remainder.  These are
definitions by exact subtraction, so no term may be omitted from their sum.

\begin{theorem}[Fibrewise VRMC singular-value criterion]
\label{thm:newV97-fibrewise-floor}
Let \(m_*=\min_fm_f>0\).  Suppose the four assembled terms in
\eqref{eq:newV97-assembled} obey
\begin{align}
 \|E_{\rm br}\|&\le {4\epsilon_{\rm br}\over ac_\delta},
                                                               \label{eq:newV97-Ebr}\\
 \|E_H\|&\le y_*h,                                             \label{eq:newV97-EH}\\
 \|E_{\rm tr}\|&\le M_{\rm tr}\omega(d_{\rm tr}),              \label{eq:newV97-Etr}\\
 \|E_{\rm rem}\|&\le r_*                                      \label{eq:newV97-Erem}
\end{align}
with \(d_{\rm tr}<1\).  If
\begin{equation}
 \Theta:={1\over m_*}
 \left\{{4\epsilon_{\rm br}\over ac_\delta}
       +y_*h+M_{\rm tr}\omega(d_{\rm tr})+r_*\right\}<1,       \label{eq:newV97-Theta}
\end{equation}
then
\begin{equation}
 s_{\min}(\mathcal D(U,H))\ge(1-\Theta)m_*,
 \qquad
 \|\mathcal D(U,H)^{-1}\|\le{1\over(1-\Theta)m_*}.             \label{eq:newV97-fibre-floor}
\end{equation}
The unbroken connection \(V\) is not required to be close to the trivial
connection.
\end{theorem}

\begin{proof}
For every admissible \(V\), \Cref{thm:newV95-mass-floor} gives
\[
 s_{\min}\left(\bigoplus_fD_{m_f}(V)\right)\ge m_*.
\]
Estimate the norm of the complete remainder in
\eqref{eq:newV97-assembled} only after assembling all four terms:
\[
 \|E_{\rm br}+E_H+E_{\rm tr}+E_{\rm rem}\|
 \le\Theta m_*.
\]
Now apply \Cref{lem:newV95-Weyl}; the inverse bound is the reciprocal of the
least singular value.
\end{proof}

\begin{corollary}[Necessary scaling in the broken-direction estimate]
\label{cor:newV97-broken-scaling}
Within the sufficient bound \eqref{eq:newV97-Theta}, a cutoff-independent
positive floor requires
\begin{equation}
 \epsilon_{\rm br}=O(am_*),                                \label{eq:newV97-broken-scaling}
\end{equation}
unless the broken kinetic term is absorbed into a stronger reference
operator rather than treated by \eqref{eq:newV97-Ebr}.
\end{corollary}

\begin{proof}
All terms in \eqref{eq:newV97-Theta} are nonnegative.  Boundedness of the
first term as \(a\downarrow0\) requires the displayed scaling.
\end{proof}

\begin{assessment}[The representation-matching bottleneck]
\label{ass:newV97-representation-bottleneck}
Before symmetry reduction, the left and right Standard Model multiplets
carry different \(SU(2)\times U(1)\) representations.  They cannot simply
be inserted into one vectorlike operator \(D_m(U)\).  VRMC makes the missing
operation explicit: the nonzero Higgs field must define a covariant
reduction, the restrictions to the unbroken subgroup must be intertwined,
and the Ginsparg--Wilson Yukawa term must reproduce the standard massive
overlap normalization.  V95's vectorlike positivity applies only after
these typed steps.
\end{assessment}

\section{Relation to ELBAC and the next acquisition}
\label{sec:newV97-next}

V97 closes a global operator-norm chart once VRMC and
\eqref{eq:newV97-Theta} are supplied.  It does not yet close V96's ELBAC,
because the Kato transport and sign difference were bounded globally rather
than by a pointwise exponentially decaying kernel.

V98 will construct the Higgs-defined reduction at the algebraic bundle
level for the one-doublet electroweak representation.  It will compute the
unbroken electromagnetic charges and the exact left--right intertwiners,
identify the unmatched neutrino row in the minimal model, and then test
whether the resulting lattice link split is local and covariant.  If a
global normalized Higgs direction is unavailable because the Higgs field
vanishes, the construction will remain confined to a declared broken-phase
chart, consistently with V95.

\begin{firewall}[Claim boundary after V97]
\label{fire:newV97-boundary}
V97 proves a gap-controlled sign Lipschitz estimate, an explicit Wilson-link
constant, overlap and projector chart bounds, and a canonical unitary
transport between nearby moving chiral source spaces.  It proves that the
direct fixed-reference estimate has an \(am\)-scale sufficient radius and
replaces that split by a fibrewise vectorlike criterion which keeps the
unbroken gauge connection in the massive base.

V97 does not construct the electroweak Higgs reduction, prove VRMC or ELBAC
for the Standard Model, prove concentration in its chart, trivialize the
chiral determinant line, establish the mapping-torus gate, prove full
matter--gravity reflection positivity, remove the cutoff, or construct the
full SM--Einstein continuum.
\end{firewall}

\begin{theorem}[V97 result ledger]
\label{thm:newV97-results}
V97:
\begin{enumerate}[label=\textup{(V97.\arabic*)},leftmargin=6.5em]
\item proves the resolvent sign-difference bound;
\item computes the four-dimensional Wilson-link Lipschitz constant;
\item obtains volume-independent overlap and moving-projector bounds;
\item gives an explicit Kato--Nagy source-fibre transport;
\item isolates the \(am\)-scale limitation of a fixed free reference;
\item formulates VRMC so that representation matching is not assumed;
\item proves the ultralocal Higgs multiplication bound; and
\item proves the assembled fibrewise singular-value criterion
      \eqref{eq:newV97-Theta}.
\end{enumerate}
\end{theorem}

\begin{proof}
Use \Cref{lem:newV97-sign-resolvent,
lem:newV97-Wilson-link,
thm:newV97-overlap-chart,
lem:newV97-Kato,
prop:newV97-am-chart,
lem:newV97-Higgs-bound,
thm:newV97-fibrewise-floor,
cor:newV97-broken-scaling}.
\end{proof}

\paragraph{Parent-version record.}
V97 was formed from
\texttt{Renewal\_SM\_Einstein\_Continuum\_Research\_Notes\_V96.tex}.
The V96 parent had \(33\,832\) logical source lines,
\(1\,299\,867\) bytes, and SHA--256
\[
 \texttt{753E708FB33B4C7FAD15376B5659A436B2ADAFFD9D0F5345BF2C192793F1305A}.
\]
The next compulsory companion audit is V100.

% =============================================================================
% END APPENDED SECOND-SERIES V97 MATERIAL
% =============================================================================

% =============================================================================
% BEGIN APPENDED SECOND-SERIES V98 MATERIAL
% =============================================================================

\clearpage
\chapter{Second-series V98: Higgs reduction, electromagnetic matching,
and a covariant tubular link split}
\label{sec:v2v98}

\section{The one-doublet stabilizer}
\label{sec:newV98-stabilizer}

Work first at the Lie-algebra level, where finite central quotients do not
alter the calculation.  Let
\[
 G_{\rm EW}=SU(2)_L\times U(1)_Y,
 \qquad
 T_j={\sigma_j\over2},
\]
and use the hypercharge convention of
\eqref{eq:v6v51-hypercharge-table}, so that
\(Q=T_3+Y\).  The Higgs doublet has \(Y_H=1/2\).  Choose the unit reference
vector
\begin{equation}
 n_0=\binom{0}{1}.                                         \label{eq:newV98-n0}
\end{equation}

\begin{lemma}[Connected stabilizer of the Higgs direction]
\label{lem:newV98-stabilizer}
The infinitesimal stabilizer of \(n_0\) in the Higgs representation is
one-dimensional and generated by
\begin{equation}
 Q=T_3+Y.                                                   \label{eq:newV98-Q}
\end{equation}
Consequently a nonzero one-doublet Higgs field reduces the connected
electroweak group to \(U(1)_{\rm em}\); with colour restored, the connected
unbroken Lie algebra is
\(\mathfrak{su}(3)_c\oplus\mathfrak u(1)_{\rm em}\).
\end{lemma}

\begin{proof}
Let
\(X=\sum_{j=1}^3\alpha_jT_j+\beta Y\).
The first component of \(Xn_0\) vanishes only if
\(\alpha_1=\alpha_2=0\).  Since
\(T_3n_0=-n_0/2\) and \(Yn_0=n_0/2\), the second component vanishes exactly
when \(\beta=\alpha_3\).  Thus \(X=\alpha_3(T_3+Y)\).  Exponentiation gives
the connected stabilizer.  Colour acts trivially on the Higgs and is
unchanged.
\end{proof}

\begin{assessment}[Global group convention]
\label{ass:newV98-global-quotient}
The physical Standard Model gauge group may be presented with a finite
central quotient.  This changes the global period and common centre of the
embedded electromagnetic circle, but it does not change
\eqref{eq:newV98-Q} or any charge calculation below.  A global principal
bundle construction must retain that quotient data.
\end{assessment}

\section{A nonzero Higgs field as a reduction}
\label{sec:newV98-reduction}

Let \(H_x\in\mathbb C^2\) be the lattice Higgs field.  On the broken chart
\begin{equation}
 \|H_x\|\ge r_->0\quad\hbox{for every site }x,              \label{eq:newV98-nonzero-chart}
\end{equation}
define
\begin{equation}
 r_x:=\|H_x\|,
 \qquad n_x:={H_x\over r_x}.                               \label{eq:newV98-normalized-Higgs}
\end{equation}
The orbit of \(n_0\) is \(G_{\rm EW}/U(1)_{\rm em}\), so \(n_x\) is
equivalently a pointwise reduction to the electromagnetic stabilizer.

\begin{lemma}[Lipschitz control of the normalized Higgs direction]
\label{lem:newV98-normalized-Lipschitz}
If \(\|H\|,\|H'\|\ge r_->0\), then
\begin{equation}
 \left\|{H\over\|H\|}-{H'\over\|H'\|}\right\|
 \le {2\over r_-}\|H-H'\|.                                \label{eq:newV98-normalized-Lipschitz}
\end{equation}
For the rank-one projectors \(P_H=nn^*\),
\begin{equation}
 \|P_H-P_{H'}\|
 \le {4\over r_-}\|H-H'\|.                                \label{eq:newV98-Higgs-projector}
\end{equation}
\end{lemma}

\begin{proof}
Insert and subtract \(H'/\|H\|\):
\[
 \left\|{H\over\|H\|}-{H'\over\|H'\|}\right\|
 \le{\|H-H'\|\over\|H\|}
 +\|H'\|\left|{1\over\|H\|}-{1\over\|H'\|}\right|.
\]
The reverse triangle inequality bounds the second term by
\(\|H-H'\|/\|H\|\), proving the first estimate.  Next,
\[
 nn^*-n'n'^*=n(n-n')^*+(n-n')n'^*,
\]
so the projector difference is at most \(2\|n-n'\|\).
\end{proof}

\begin{proposition}[The Higgs-zero obstruction is exactly the reduction
obstruction]
\label{prop:newV98-zero-reduction}
The normalized reduction \eqref{eq:newV98-normalized-Higgs} extends
continuously over a field configuration only if its Higgs direction has a
continuous limit at every zero.  No such extension exists on a full
neighborhood of the zero field.  Hence a VRMC construction based on the
Higgs reduction must remain on a chart such as
\eqref{eq:newV98-nonzero-chart}.
\end{proposition}

\begin{proof}
Along \(H=t n\), \(t\downarrow0\), the normalized direction is the arbitrary
unit vector \(n\).  Different rays have different limits in
\(G_{\rm EW}/U(1)_{\rm em}\), so no continuous extension exists on a full
neighborhood of zero.  This is the bundle form of the zero-Higgs
singular-floor obstruction in \Cref{thm:newV95-zero-Higgs-no-go}.
\end{proof}

\section{Electromagnetic charge matching}
\label{sec:newV98-charges}

Use the all-left-handed one-generation multiplets and hypercharges from
\eqref{eq:v6v51-hypercharge-table}.  On an \(SU(2)\) doublet,
\(Q=T_3+Y\); on a singlet, \(Q=Y\).

\begin{lemma}[One-generation electromagnetic ledger]
\label{lem:newV98-charge-ledger}
The electric charges are
\begin{equation}
 \begin{array}{c|ccccccc}
 \text{field}
 &u_L&d_L&\nu_L&e_L&u^c&d^c&e^c\\ \hline
 Q
 &{2\over3}&-{1\over3}&0&-1&-{2\over3}&{1\over3}&1 .
 \end{array}                                                \label{eq:newV98-charge-table}
\end{equation}
Equivalently, the physical right-handed fields
\(u_R,d_R,e_R\), obtained by conjugating
\(u^c,d^c,e^c\), have charges
\(2/3,-1/3,-1\), equal to those of their left-handed partners.
There is no right-handed partner for \(\nu_L\) in the minimal
representation list.
\end{lemma}

\begin{proof}
For \(Q_L\), add \(Y=1/6\) to \(T_3=1/2,-1/2\).  For \(L\), add
\(Y=-1/2\) to the same two eigenvalues.  The three conjugate fields are
\(SU(2)\) singlets, so their charges equal their hypercharges
\(-2/3,1/3,1\).  Charge conjugation reverses these signs.  The absence of a
minimal \(\nu^c\) is the representation statement already isolated in
\Cref{prop:newV95-massless-row}.
\end{proof}

\begin{theorem}[Algebraic left--right matching after Higgs reduction]
\label{thm:newV98-vectorlike-matching}
After restriction to
\(SU(3)_c\times U(1)_{\rm em}\), the physical left and right components of
each charged massive species carry unitarily equivalent representations:
\[
 (u_L,u_R),\qquad(d_L,d_R),\qquad(e_L,e_R).
\]
For three generations, the statement tensors these representation
intertwiners with flavour space; singular-value decomposition of the Yukawa
matrices changes only the flavour bases.  The minimal neutrino row remains
unmatched.  If a sterile \(\nu_R\) of charge zero and a nonzero neutrino
Yukawa singular value are added, that row is matched as well.
\end{theorem}

\begin{proof}
The charge equalities are
\Cref{lem:newV98-charge-ledger}.  Both quark chiralities are colour
triplets in the physical right-handed convention, while the leptons are
colour singlets.  Thus identity maps on the common colour and charge spaces
give unitary intertwiners.  A flavour unitary commutes with the gauge
representation, so the two unitaries in a Yukawa singular-value
decomposition preserve the intertwining property.  The neutrino conclusions
follow from charge zero and the presence or absence of the right-handed
singlet.
\end{proof}

\begin{assessment}[All-left-handed notation]
\label{ass:newV98-left-handed-notation}
In the anomaly ledger, \(u^c,d^c,e^c\) carry the dual electromagnetic
representations.  Their charges therefore have the opposite signs from the
physical right-handed fields.  The gauge-invariant two-Weyl mass bilinear
pairs a representation with its dual; the corresponding four-component
Dirac operator pairs physical left and right fields with equal charge.
These are the same matching statement in two conventions.
\end{assessment}

\section{A covariant tubular link split}
\label{sec:newV98-tubular}

The reduction must act on links without choosing a global unitary gauge.
The following compact-group construction provides the required local split.

\begin{lemma}[Equivariant tubular projection onto a compact subgroup]
\label{lem:newV98-tubular}
Let \(G\) be a compact Lie group, \(K\subset G\) a closed subgroup, and give
\(G\) a bi-invariant Riemannian metric.  There is an open
\(K\times K\)-invariant neighborhood \(\mathcal N_K\) of \(K\) and a smooth
nearest-point projection
\begin{equation}
 \pi_K:\mathcal N_K\longrightarrow K                       \label{eq:newV98-piK}
\end{equation}
such that
\begin{equation}
 \pi_K(k_1gk_2)=k_1\pi_K(g)k_2
 \quad(k_1,k_2\in K).                                      \label{eq:newV98-piK-equivariant}
\end{equation}
\end{lemma}

\begin{proof}
The tubular-neighborhood theorem applied to the embedded compact submanifold
\(K\subset G\) gives a neighborhood on which the normal exponential map is a
diffeomorphism and nearest points are unique.  Left and right multiplication
by elements of \(K\) preserve both the bi-invariant metric and \(K\).
They therefore preserve normal geodesics and carry the unique nearest point
of \(g\) to the unique nearest point of \(k_1gk_2\), which proves
\eqref{eq:newV98-piK-equivariant}.  Intersecting the tubular neighborhood
over the compact \(K\times K\) action makes it invariant.
\end{proof}

At each site choose a local lift \(s_x\in G_{\rm EW}\) satisfying
\(s_xn_0=n_x\).  A different lift has the form \(s_xk_x\) with \(k_x\in K\).
For an electroweak link \(U_{xy}\), define the dressed link
\begin{equation}
 A_{xy}:=s_x^{-1}U_{xy}s_y.                                \label{eq:newV98-dressed-link}
\end{equation}
Whenever \(A_{xy}\in\mathcal N_K\), put
\begin{equation}
 V_{xy}:=\pi_K(A_{xy})\in K,
 \qquad
 b_{xy}:=\operatorname{dist}(A_{xy},K).                    \label{eq:newV98-reduced-link}
\end{equation}

\begin{theorem}[Lift independence and gauge covariance of the reduced link]
\label{thm:newV98-reduced-link}
The family \(V_{xy}\) in \eqref{eq:newV98-reduced-link} defines a
\(K\)-connection on the Higgs-reduced lattice bundle.  It is independent of
the chosen lifts in the precise transition-law sense:
if \(s_x'=s_xk_x\), then
\begin{equation}
 V_{xy}'=k_x^{-1}V_{xy}k_y.                                \label{eq:newV98-K-transition}
\end{equation}
The number \(b_{xy}\) is lift independent and gauge invariant.  Both
quantities depend only on the two endpoint Higgs fields and the intervening
link, and hence are lattice local on the nonzero tubular chart.
\end{theorem}

\begin{proof}
Changing lifts gives
\[
 A_{xy}'=k_x^{-1}A_{xy}k_y.
\]
Apply \eqref{eq:newV98-piK-equivariant} to obtain
\eqref{eq:newV98-K-transition}; bi-invariance of the distance gives
\(b_{xy}'=b_{xy}\).  Under a gauge transformation, transformed lifts can be
chosen by multiplying the old lifts by the gauge transformations; the same
calculation leaves the reduced geometric connection and \(b_{xy}\)
unchanged up to their required bundle transition laws.  The formula uses no
field outside \(x,y\), and the smoothness follows from
\Cref{lem:newV98-tubular} and the local smooth lifts.
\end{proof}

\begin{lemma}[Representationwise distance bound]
\label{lem:newV98-representation-distance}
For the finite set \(\mathcal R_{\rm SM}\), after shrinking
\(\mathcal N_K\) if necessary, there is a finite \(L_{\rm SM}\) such that
\begin{equation}
 \|R(A)-R(\pi_K(A))\|
 \le L_{\rm SM}\operatorname{dist}(A,K)                   \label{eq:newV98-R-distance}
\end{equation}
for every \(R\in\mathcal R_{\rm SM}\) and \(A\in\mathcal N_K\).
\end{lemma}

\begin{proof}
For fixed \(R\), the derivative of the smooth matrix map \(R:G\to U(d_R)\)
is bounded on the compact closure of a sufficiently small tubular
neighborhood.  Integrate that derivative along the unique normal geodesic
from \(\pi_K(A)\) to \(A\).  Take the maximum over the finite representation
set.
\end{proof}

\begin{lemma}[Four-link plaquette stability]
\label{lem:newV98-plaquette-stability}
Let \(A_1,\ldots,A_4\) and \(V_1,\ldots,V_4\) be unitary matrices, with
inverse orientations inserted in the same positions in both plaquette
products.  If \(\|A_j-V_j\|\le\epsilon\) for every \(j\), then
\begin{equation}
 \|A_1A_2A_3A_4-V_1V_2V_3V_4\|\le4\epsilon.               \label{eq:newV98-four-link}
\end{equation}
Consequently, if the full dressed connection has represented plaquette
defect at most \(\delta_A\), then its reduced connection has defect at most
\begin{equation}
 \delta_A+4L_{\rm SM}b_*.                                  \label{eq:newV98-reduced-defect}
\end{equation}
\end{lemma}

\begin{proof}
Telescope the difference of the two four-factor products, replacing one
factor at a time.  Every untouched factor is unitary, so each of the four
summands has norm at most \(\epsilon\).  An inverse-oriented factor obeys
\(\|A^{-1}-V^{-1}\|=\|A-V\|\).  Apply
\Cref{lem:newV98-representation-distance} and maximize over the plaquettes
and the finite representation set.
\end{proof}

\begin{corollary}[Explicit broken-link input for the V97 criterion]
\label{cor:newV98-broken-input}
Assume that the full dressed connection lies in
\(\mathscr A_{\delta_A}\) and that
\begin{equation}
 \delta_A+4L_{\rm SM}b_*\le\delta<1/30,
 \qquad
 b_*:=\max_{\langle xy\rangle}b_{xy}.                     \label{eq:newV98-bstar}
\end{equation}
Then the reduced \(K\)-connection and the full connection both lie in
\(\mathscr A_\delta\), and their link-coordinate difference is at most
\(L_{\rm SM}b_*\) in every Standard Model fermion representation.  Thus the
broken kinetic row in \eqref{eq:newV97-Theta} may be taken as
\begin{equation}
 {4L_{\rm SM}b_*\over ac_\delta}.                          \label{eq:newV98-broken-row}
\end{equation}
\end{corollary}

\begin{proof}
Use \Cref{lem:newV98-representation-distance,
lem:newV98-plaquette-stability} for admissibility and link distance, then
apply \Cref{lem:newV97-Wilson-link,thm:newV97-overlap-chart}.
\end{proof}

\section{Radial Yukawa variation and the partially closed VRMC}
\label{sec:newV98-VRMC}

In a Higgs-reduced frame, the angular Higgs variable is absorbed into the
choice of reduced bundle.  Write the remaining radial fluctuation as
\[
 r_x=v+\eta_x.
\]
Let \(y_*\) be the largest singular value among the declared Yukawa matrices.
The ultralocal row of V97 gives
\begin{equation}
 \|E_H\|\le y_*\|\eta\|_\infty.                            \label{eq:newV98-radial-row}
\end{equation}

\begin{theorem}[Algebraic and geometric rows of VRMC]
\label{thm:newV98-VRMC-rows}
On the chart
\[
 r_x\ge r_->0,\qquad A_{xy}\in\mathcal N_K,\qquad
 A\in\mathscr A_{\delta_A},\qquad
 \delta_A+4L_{\rm SM}b_*\le\delta<1/30,
\]
the following VRMC rows are closed:
\begin{enumerate}[label=\textup{(\roman*)},leftmargin=3.5em]
\item the Higgs field defines a local covariant reduction to
      \(K=U(1)_{\rm em}\), with colour carried as a spectator;
\item the charged left and right Standard Model species have explicit
      \(K\)-representation intertwiners;
\item the link field has the lift-independent local split
      \eqref{eq:newV98-reduced-link};
\item the broken kinetic contribution obeys
      \eqref{eq:newV98-broken-row}; and
\item the radial Yukawa fluctuation obeys
      \eqref{eq:newV98-radial-row}.
\end{enumerate}
If, in addition, the chosen chiral lattice Yukawa operator reproduces the
standard massive-overlap block at \(r_x=v\), and its transport and
discretization remainders obey the remaining V97 bounds, then
\begin{equation}
 {1\over m_*}\left\{
 {4L_{\rm SM}b_*\over ac_\delta}
 +y_*\|\eta\|_\infty
 +M_{\rm tr}\omega(d_{\rm tr})+r_*\right\}<1              \label{eq:newV98-final-chart}
\end{equation}
implies the floor \eqref{eq:newV97-fibre-floor}.
\end{theorem}

\begin{proof}
The five rows are respectively
\Cref{lem:newV98-stabilizer,
prop:newV98-zero-reduction,
thm:newV98-vectorlike-matching,
thm:newV98-reduced-link,
cor:newV98-broken-input,
lem:newV97-Higgs-bound}.
The final implication is \Cref{thm:newV97-fibrewise-floor} with
\eqref{eq:newV98-broken-row} and
\eqref{eq:newV98-radial-row} substituted into
\eqref{eq:newV97-Theta}.
\end{proof}

\begin{assessment}[Remaining exact-normalization row]
\label{ass:newV98-normalization-open}
The theorem does not assume that a continuum-looking Yukawa multiplication
automatically equals the massive-overlap deformation
\((1-am/2)D_{\rm ov}+m\).  That equality depends on the precise
Ginsparg--Wilson chiral Yukawa discretization and its source/target
projectors.  Establishing it, or computing its remainder, is the next
operator-level VRMC row.
\end{assessment}

\section{Next acquisition step}
\label{sec:newV98-next}

V99 will write a finite-dimensional Ginsparg--Wilson Yukawa block in the
Higgs-reduced frames and compare its Schur assembly with
\eqref{eq:newV97-fibre-base}.  It will either derive the missing factor
\(1-am/2\) and close the constant-Higgs normalization row, or record the
exact defect operator and test it against
\eqref{eq:newV98-final-chart}.  The chiral determinant-line phase will remain
separate.  V100 will then perform the compulsory YM/NS audit before choosing
the next gate.

\begin{firewall}[Claim boundary after V98]
\label{fire:newV98-boundary}
V98 derives the one-doublet electromagnetic stabilizer and complete
one-generation charge ledger, proves left--right representation matching for
the charged massive species, identifies the minimal neutrino mismatch,
constructs a lift-independent covariant tubular reduction of electroweak
links, and, when both the full and reduced connections obey the common V94
admissibility bound, inserts its explicit broken-link constant into the V97
floor criterion.

V98 does not extend the reduction through Higgs zeros, prove concentration
in the nonzero tubular chart, establish the exact Ginsparg--Wilson Yukawa
normalization, close ELBAC, construct the chiral determinant phase, prove
matter--gravity reflection positivity, remove the cutoff, or construct the
full SM--Einstein continuum.
\end{firewall}

\begin{theorem}[V98 result ledger]
\label{thm:newV98-results}
V98:
\begin{enumerate}[label=\textup{(V98.\arabic*)},leftmargin=6.5em]
\item computes the connected Higgs stabilizer \(U(1)_{\rm em}\);
\item proves quantitative control of the normalized Higgs reduction;
\item proves that the reduction cannot extend over a full Higgs-zero
      neighborhood;
\item computes all one-generation electromagnetic charges;
\item proves the charged left--right intertwiners after reduction;
\item constructs the equivariant tubular projection and reduced link;
\item proves the representationwise broken-link and plaquette-stability
      bounds; and
\item closes the algebraic and geometric VRMC rows subject to the exact
      overlap--Yukawa normalization row.
\end{enumerate}
\end{theorem}

\begin{proof}
Use \Cref{lem:newV98-stabilizer,
lem:newV98-normalized-Lipschitz,
prop:newV98-zero-reduction,
lem:newV98-charge-ledger,
thm:newV98-vectorlike-matching,
lem:newV98-tubular,
thm:newV98-reduced-link,
lem:newV98-representation-distance,
lem:newV98-plaquette-stability,
cor:newV98-broken-input,
thm:newV98-VRMC-rows}.
\end{proof}

\paragraph{Parent-version record.}
V98 was formed from
\texttt{Renewal\_SM\_Einstein\_Continuum\_Research\_Notes\_V97.tex}.
The V97 parent had \(34\,267\) logical source lines,
\(1\,315\,683\) bytes, and SHA--256
\[
 \texttt{534150582BA302223A088DE2ACB0F31CA11622FDDFF027B9F50FB82BC2A5FDF8}.
\]
The next compulsory companion audit is V100.

% =============================================================================
% END APPENDED SECOND-SERIES V98 MATERIAL
% =============================================================================

% =============================================================================
% BEGIN APPENDED SECOND-SERIES V99 MATERIAL
% =============================================================================

\clearpage
\chapter{Second-series V99: exact Ginsparg--Wilson Yukawa normalization}
\label{sec:v2v99}

\section{The improved-field intertwiner}
\label{sec:newV99-improved}

Let \(D=D_{\rm ov}(V)\) be the overlap operator for one vectorlike
\(K\)-representation on an AOP connection.  Recall
\begin{equation}
 \gamma_5D+D\gamma_5=aD\gamma_5D,
 \qquad
 \widehat\gamma_5=\gamma_5(I-aD),
 \qquad
 \widehat P_\pm={I\pm\widehat\gamma_5\over2}.              \label{eq:newV99-GW-data}
\end{equation}
Set
\begin{equation}
 A_D:=I-{a\over2}D.                                        \label{eq:newV99-AD}
\end{equation}

\begin{lemma}[Two exact Ginsparg--Wilson intertwiners]
\label{lem:newV99-intertwiners}
The operators \(D\) and \(A_D\) obey
\begin{align}
 D\widehat\gamma_5&=-\gamma_5D,
 &D\widehat P_\pm&=P_\mp D,                                \label{eq:newV99-D-intertwine}\\
 A_D\widehat\gamma_5&=\gamma_5A_D,
 &A_D\widehat P_\pm&=P_\pm A_D.                            \label{eq:newV99-A-intertwine}
\end{align}
Thus \(D\) maps a moving source chirality to the opposite fixed row
chirality, while \(A_D\) maps it to the same fixed chirality.
\end{lemma}

\begin{proof}
Using \eqref{eq:newV99-GW-data},
\[
 D\widehat\gamma_5
 =D\gamma_5-aD\gamma_5D
 =-\gamma_5D,
\]
which gives \eqref{eq:newV99-D-intertwine}.  For the second identity,
\begin{align*}
 A_D\widehat\gamma_5
 &=(I-aD/2)\gamma_5(I-aD)\\
 &=\gamma_5-a\gamma_5D-{a\over2}D\gamma_5
   +{a^2\over2}D\gamma_5D\\
 &=\gamma_5-{a\over2}\gamma_5D
 =\gamma_5A_D,
\end{align*}
where the Ginsparg--Wilson relation was used in the third line.
Multiplying by \((I\pm\widehat\gamma_5)/2\) proves
\eqref{eq:newV99-A-intertwine}.
\end{proof}

\begin{lemma}[The improved-field contraction]
\label{lem:newV99-AD-contraction}
With \(V_D=\gamma_5\varepsilon(V)\), one has
\begin{equation}
 D={1\over a}(I+V_D),
 \qquad
 A_D={1\over2}(I-V_D),
 \qquad
 \|A_D\|\le1.                                              \label{eq:newV99-AD-unitary}
\end{equation}
\end{lemma}

\begin{proof}
The first identity is \eqref{eq:newV95-Dov-V}.  Substitution in
\eqref{eq:newV99-AD} gives the second.  Since \(V_D\) is unitary,
\(\|I-V_D\|\le2\).
\end{proof}

\section{Exact assembly of the massive overlap block}
\label{sec:newV99-massive-assembly}

Decompose the source as
\(\operatorname{Ran}\widehat P_-\oplus
  \operatorname{Ran}\widehat P_+\)
and the row space as
\(\operatorname{Ran}P_+\oplus\operatorname{Ran}P_-\).
The ordering reflects the usual convention
\(\psi_L=\widehat P_-\psi\),
\(\psi_R=\widehat P_+\psi\),
\(\bar\psi_L=\bar\psi P_+\), and
\(\bar\psi_R=\bar\psi P_-\).

\begin{theorem}[Exact chiral block assembly]
\label{thm:newV99-exact-assembly}
For \(m\ge0\), the block operator
\begin{equation}
 \mathcal B_m:=
 \begin{pmatrix}
 P_+D\widehat P_-&
 mP_+A_D\widehat P_+\\
 mP_-A_D\widehat P_-&
 P_-D\widehat P_+
 \end{pmatrix}                                             \label{eq:newV99-block}
\end{equation}
is, under the displayed orthogonal source and row decompositions, exactly
the full-space operator
\begin{equation}
 D+mA_D
 =\left(1-{am\over2}\right)D+mI
 =D_m.                                                      \label{eq:newV99-exact-Dm}
\end{equation}
Consequently, for \(0<m\le2/a\),
\begin{equation}
 s_{\min}(\mathcal B_m)\ge m,
 \qquad
 \|\mathcal B_m^{-1}\|\le m^{-1}.                          \label{eq:newV99-block-floor}
\end{equation}
\end{theorem}

\begin{proof}
By \eqref{eq:newV99-D-intertwine}, the two omitted kinetic corners vanish
and the sum of the displayed kinetic blocks is
\[
 P_+D\widehat P_-+P_-D\widehat P_+=D.
\]
By \eqref{eq:newV99-A-intertwine}, the two omitted improved-field corners
vanish and the sum of the displayed mass blocks is
\[
 mP_+A_D\widehat P_++mP_-A_D\widehat P_-=mA_D.
\]
This proves \eqref{eq:newV99-exact-Dm}.  Orthogonal decompositions preserve
singular values, so \eqref{eq:newV99-block-floor} is
\Cref{thm:newV95-mass-floor}.
\end{proof}

\begin{assessment}[Why the labels in the mass corners agree]
\label{ass:newV99-mass-corners}
The mass term is off-diagonal in the physical left--right fields even though
\(A_D\) maps equal projector signs in
\eqref{eq:newV99-A-intertwine}.  The row convention reverses the sign:
\(\bar\psi_L\) occupies \(P_+\), while \(\psi_R\) occupies
\(\widehat P_+\).  Thus the upper-right entry of
\eqref{eq:newV99-block} is precisely the left--right mass coupling.
\end{assessment}

\section{Flavour matrices}
\label{sec:newV99-flavour}

Let \(M:\mathbb C^{n_R}\to\mathbb C^{n_L}\) be a constant Yukawa mass
matrix after the V98 reduction and representation intertwiners.  The kinetic
operator is the identity in flavour space.  Use the improved field \(A_D\)
in each Yukawa corner.

\begin{theorem}[Yukawa singular values are massive-overlap floors]
\label{thm:newV99-flavour-SVD}
Let
\[
 M=U_L\Sigma U_R^*,
 \qquad
 \Sigma=\operatorname{diag}(m_1,\ldots,m_r,0,\ldots)
\]
be a singular-value decomposition.  Constant unitary changes of left and
right flavour frames transform the assembled Ginsparg--Wilson
kinetic--Yukawa block into the orthogonal direct sum
\begin{equation}
 \bigoplus_{j=1}^rD_{m_j}
 \ \oplus\ \mathcal D_{\rm unmatched}.                    \label{eq:newV99-flavour-sum}
\end{equation}
If every retained massive singular value satisfies
\(0<m_j\le2/a\), then its direct summand has least singular value at least
\(m_j\).  Zero singular values and rectangular unmatched rows belong to
\(\mathcal D_{\rm unmatched}\) and cannot be included in a positive massive
floor.
\end{theorem}

\begin{proof}
The constant flavour unitaries commute with \(D\), \(A_D\), and all spin,
site, and common \(K\)-representation projectors.  They therefore diagonalize
only the two Yukawa corners, with \(M\) in one corner and \(M^*\) in the
other.  Each nonzero diagonal entry produces exactly
\eqref{eq:newV99-block}, hence \(D_{m_j}\), by
\Cref{thm:newV99-exact-assembly}.  The remaining source or row directions
have zero Yukawa singular value and form the unmatched block.  Apply
\eqref{eq:newV99-block-floor} to every nonzero summand.
\end{proof}

\begin{corollary}[Constant-Higgs normalization row of VRMC]
\label{cor:newV99-VRMC-normalization}
On the V98 Higgs-reduced chart, suppose the left--right
\(K\)-representation intertwiners are used and the chiral Yukawa action is
defined on the improved field \(A_D\psi\).  At constant radial Higgs value
\(r_x=v\), every nonzero Yukawa singular value gives exactly the fibrewise
massive-overlap reference \eqref{eq:newV97-fibre-base}.  Thus VRMC item
\textup{(iv)} is closed for this discretization.
\end{corollary}

\begin{proof}
V98 supplies the common \(K\)-representation and its unitary intertwiners.
In those frames the constant Yukawa matrices act only in flavour space.
Apply \Cref{thm:newV99-flavour-SVD}.
\end{proof}

\section{Variable radial Higgs field}
\label{sec:newV99-variable-Higgs}

Let \(M_x=M_0+\Delta M_x\) be a sitewise Yukawa mass matrix, with \(M_0\)
the constant reduced value.  Define the variable Yukawa term with the
operator order
\begin{equation}
 \mathcal M_H:=M_HA_D,                                     \label{eq:newV99-MHAD}
\end{equation}
together with the chiral corner projections of
\eqref{eq:newV99-block}.  No commutation of \(M_H\) with \(A_D\) is assumed.

\begin{lemma}[Sharp variable-mass perturbation bound]
\label{lem:newV99-variable-mass}
The deviation from the constant Yukawa block satisfies
\begin{equation}
 \|(M_H-M_0)A_D\|
 \le\sup_x\|\Delta M_x\|.                                 \label{eq:newV99-variable-bound}
\end{equation}
For \(M_x=Y r_x\), this gives
\begin{equation}
 \|(M_H-M_0)A_D\|
 \le\|Y\|\,\|r-v\|_\infty.                                \label{eq:newV99-radial-bound}
\end{equation}
\end{lemma}

\begin{proof}
Sitewise multiplication has norm
\(\sup_x\|\Delta M_x\|\).  Multiply by
\(\|A_D\|\le1\) from \Cref{lem:newV99-AD-contraction}.  The last statement
uses \(\Delta M_x=Y(r_x-v)\).
\end{proof}

\begin{assessment}[Operator ordering]
\label{ass:newV99-ordering}
For nonconstant \(M_H\), the operators \(M_HA_D\) and \(A_DM_H\) differ.
The proof closes the first declared ordering, in which the improved fermion
field is formed before Yukawa multiplication.  Any alternative ordering
must carry its commutator as an explicit remainder; cyclicity of a
determinant does not identify the two chiral maps.
\end{assessment}

\section{The naive mass insertion loses the certified margin}
\label{sec:newV99-naive}

\begin{proposition}[Exact defect of an unimproved mass]
\label{prop:newV99-naive-defect}
The naive full-space insertion \(D+mI\) differs from the standard massive
overlap operator by
\begin{equation}
 (D+mI)-D_m={am\over2}D.                                  \label{eq:newV99-naive-defect}
\end{equation}
Since \(\|D\|\le2/a\),
\begin{equation}
 \|(D+mI)-D_m\|\le m.                                     \label{eq:newV99-naive-bound}
\end{equation}
The V95 perturbation theorem applied only with this bound reaches
\(\theta=1\) and certifies no positive residual floor.
\end{proposition}

\begin{proof}
Subtract \eqref{eq:newV99-exact-Dm}.  The spectral circle
\eqref{eq:newV95-circle} gives \(\|D\|\le2/a\).  Hence
\eqref{eq:newV99-naive-bound}.  In
\eqref{eq:newV95-tube-condition}, the reference floor is \(m\), so an error
bound equal to \(m\) lies exactly at the excluded endpoint \(\theta=1\).
\end{proof}

\begin{assessment}[This is a proof limitation, not a singularity claim]
\label{ass:newV99-naive-scope}
The proposition does not assert that \(D+mI\) is singular.  It proves that
the already established massive-overlap floor cannot be transferred to that
operator by the sharp general perturbation inequality.  The improved-field
normalization avoids this loss by the exact identity
\eqref{eq:newV99-exact-Dm}.
\end{assessment}

\section{Updated VRMC ledger and next acquisition}
\label{sec:newV99-next}

\begin{theorem}[VRMC constant-background closure]
\label{thm:newV99-VRMC-closed}
For the improved-field Yukawa discretization
\eqref{eq:newV99-MHAD}, the constant-background algebraic rows of VRMC are
closed on the V98 nonzero tubular chart:
\begin{enumerate}[label=\textup{(\roman*)},leftmargin=3.5em]
\item the Higgs reduction and \(K\)-connection are
      \Cref{thm:newV98-reduced-link};
\item the charged left--right intertwiners are
      \Cref{thm:newV98-vectorlike-matching};
\item each nonzero constant Yukawa singular value is exactly a massive
      overlap block by \Cref{thm:newV99-flavour-SVD};
\item the radial fluctuation has the sharp bound
      \eqref{eq:newV99-radial-bound}; and
\item the minimal neutrino row remains in the unmatched Grassmann block.
\end{enumerate}
Subject to the V98 common-admissibility and tubular hypotheses and to the
remaining transport/discretization estimates, the explicit sufficient
floor condition is still \eqref{eq:newV98-final-chart}, now with no
constant-Higgs normalization remainder.
\end{theorem}

\begin{proof}
The first two items are V98.  The third and fourth are
\Cref{thm:newV99-flavour-SVD,lem:newV99-variable-mass}; the fifth is
\Cref{lem:newV98-charge-ledger}.  Insert them into
\Cref{thm:newV98-VRMC-rows}.
\end{proof}

V100 is the compulsory companion audit.  It will read the newest settled
Yang--Mills and Navier--Stokes notes, compare their current bottlenecks with
the remaining transport/discretization and measure-concentration rows, and
choose the next upstream gate.  Because V100 also reaches the active-series
rollover threshold, it will then be frozen with the next
\(\texttt{\_fN}\) suffix and a new active V1 will begin with a compact
result map and context, as required by the version protocol.

\begin{firewall}[Claim boundary after V99]
\label{fire:newV99-boundary}
V99 proves the exact improved-field chiral intertwiners, reconstructs the
standard massive overlap operator from its four chiral kinetic and mass
corners, diagonalizes constant flavour Yukawa matrices into massive-overlap
summands, controls a variable radial Higgs mass, and closes the
constant-background normalization row of VRMC for the declared
improved-field discretization.

V99 does not prove the remaining covariant transport or discretization
bounds, concentration in the V98 chart, ELBAC, the chiral determinant-line
phase, the mapping-torus gate, full reflection positivity, cutoff removal,
or the full SM--Einstein continuum.
\end{firewall}

\begin{theorem}[V99 result ledger]
\label{thm:newV99-results}
V99:
\begin{enumerate}[label=\textup{(V99.\arabic*)},leftmargin=6.5em]
\item proves the exact \(D\) and \(A_D\) chirality intertwiners;
\item proves \(\|A_D\|\le1\);
\item assembles the four chiral corners exactly into \(D_m\);
\item proves the massive floor for every nonzero Yukawa singular value;
\item closes the VRMC constant-Higgs normalization row;
\item proves the variable radial-Higgs perturbation bound; and
\item computes the full defect of an unimproved mass insertion.
\end{enumerate}
\end{theorem}

\begin{proof}
Use \Cref{lem:newV99-intertwiners,
lem:newV99-AD-contraction,
thm:newV99-exact-assembly,
thm:newV99-flavour-SVD,
cor:newV99-VRMC-normalization,
lem:newV99-variable-mass,
prop:newV99-naive-defect,
thm:newV99-VRMC-closed}.
\end{proof}

\paragraph{Parent-version record.}
V99 was formed from
\texttt{Renewal\_SM\_Einstein\_Continuum\_Research\_Notes\_V98.tex}.
The V98 parent had \(34\,750\) logical source lines,
\(1\,335\,098\) bytes, and SHA--256
\[
 \texttt{9554E4F95DB5DE5495FCF7BE3998381123136EA0DA0230E454D6708AB16E71C5}.
\]
The next compulsory companion audit is V100.

% =============================================================================
% END APPENDED SECOND-SERIES V99 MATERIAL
% =============================================================================

% =============================================================================
% BEGIN APPENDED SECOND-SERIES V100 MATERIAL
% =============================================================================

\clearpage
\chapter{Second-series V100: companion audit and assembled determinant
transport}
\label{sec:v2v100}

\section{Compulsory companion audit}
\label{sec:newV100-audit}

The companion files were read as immutable byte snapshots.  The settled
Yang--Mills body was
\begin{center}
\begin{tabular}{ll}
file &
\texttt{Renewal\_Yang\_Mills\_Mass\_Gap\_Research\_Notes\_V74.tex},\\
bytes & \(1\,110\,936\),\\
logical source lines & \(34\,855\),\\
SHA--256 &
\texttt{5424CC6C7C96795590AC0CE8AE546F22F9670BA417D57CEF422D2B2523FC8F15}.
\end{tabular}
\end{center}
A simultaneously visible file named
\texttt{Renewal\_Yang\_Mills\_Mass\_Gap\_Research\_Notes\_V75.tex}
had exactly the same byte count and digest and ended with the V74 marker and
V74 parent record.  It contained no V75 body and was therefore discarded as
a transitional duplicate.

The settled Navier--Stokes snapshot was
\begin{center}
\begin{tabular}{ll}
file &
\texttt{Renewal\_NS\_Stability\_Research\_Notes\_V43.tex},\\
bytes & \(534\,014\),\\
logical source lines & \(10\,498\),\\
SHA--256 &
\texttt{4C34D19C9557EC3801C897E3856A13708AC8C11F3A66CC7C75E720441B61EB10}.
\end{tabular}
\end{center}

\begin{audit}[Transfer at V100]
\label{audit:newV100-transfer}
Yang--Mills V74 proves that every fixed centre hypergraph has the correct
thermal degree and that incidence cycles carry a two-power reserve.  It also
exhibits an acyclic Bell family which defeats termwise summation.  The repair
for centre stars is to retain the unrelaxed clock reserve until the complete
partition has been summed; non-star hypertree pruning remains open.

Navier--Stokes V43 proves a window trichotomy.  Type-I window dissipation can
concentrate at finitely many active centres with nonzero strong limits, or it
can spread over a diverging number of cells while every fixed-centre window
limit vanishes.  Local compactness therefore does not prevent escape through
a growing spatial allocation.

For the present programme these are two forms of the same warning.  A
determinant phase or modulus cannot be controlled by bounding each
projector, flavour, derivative, or spatial block separately and then summing
an uncontrolled allocation.  The exact chiral block and its frame transport
must first be assembled.  A continuum claim must then control its trace-norm
length and exclude escape to a growing number of lattice cells.  No
numerical bound is imported from either companion note.
\end{audit}

\section{Exact determinant transport in moving chiral frames}
\label{sec:newV100-determinant-transport}

Let \(E_t,F_t\) be equal-rank Hermitian vector spaces over
\(t\in[0,1]\), let \(B_t:E_t\to F_t\) be a \(C^1\) family of invertible
maps, and choose \(C^1\) unitary transports
\[
 S_t:E_0\to E_t,\qquad R_t:F_0\to F_t.
\]
Put
\begin{equation}
 \widetilde B_t:=R_t^*B_tS_t:E_0\longrightarrow F_0.       \label{eq:newV100-Btilde}
\end{equation}

\begin{theorem}[Exponentiated Jacobi transport]
\label{thm:newV100-Jacobi-transport}
The assembled transported block obeys the exact identity
\begin{equation}
 {\det\widetilde B_1\over\det\widetilde B_0}
 =
 \exp\left\{\int_0^1
 \operatorname{Tr}\bigl(\widetilde B_t^{-1}
                         \dot{\widetilde B}_t\bigr)\,dt
      \right\}.                                             \label{eq:newV100-Jacobi}
\end{equation}
If \(s_{\min}(\widetilde B_t)\ge\kappa_t>0\), then
\begin{equation}
 \left|\operatorname{Tr}\bigl(\widetilde B_t^{-1}
                              \dot{\widetilde B}_t\bigr)\right|
 \le {1\over\kappa_t}\|\dot{\widetilde B}_t\|_1,            \label{eq:newV100-trace-bound}
\end{equation}
where \(\|\cdot\|_1\) is the trace norm.  Hence, with
\begin{equation}
 L_1(B,S,R):=\int_0^1{\|\dot{\widetilde B}_t\|_1\over\kappa_t}\,dt,
                                                               \label{eq:newV100-L1}
\end{equation}
both the logarithmic modulus change and every continuous lift of the phase
change have absolute value at most \(L_1(B,S,R)\).
\end{theorem}

\begin{proof}
Jacobi's formula gives
\[
 {d\over dt}\det\widetilde B_t
 =\det\widetilde B_t\,
  \operatorname{Tr}(\widetilde B_t^{-1}\dot{\widetilde B}_t).
\]
Divide by the nonzero determinant and integrate.  Exponentiation removes
the choice of logarithm and proves \eqref{eq:newV100-Jacobi}.  Trace duality
and the singular floor give
\[
 |\operatorname{Tr}(XY)|
 \le\|X\|\|Y\|_1,\qquad
 \|\widetilde B_t^{-1}\|\le\kappa_t^{-1},
\]
which proves \eqref{eq:newV100-trace-bound}.  The real and imaginary parts
of the integrand are respectively the derivatives of a continuous
logarithmic modulus and phase lift.
\end{proof}

\begin{assessment}[Assembly order]
\label{ass:newV100-assembly-order}
The derivative in \eqref{eq:newV100-L1} is taken after the exact operator
\(\widetilde B_t=R_t^*B_tS_t\) has been formed.  It includes the derivatives
of both Kato transports, the kinetic block, the improved-field Yukawa block,
and every declared remainder.  Estimating these pieces before their sum can
discard cancellations in the trace and recreates the allocation loss seen
in the V100 audit.
\end{assessment}

\section{Endpoint transitions and the finite-cutoff monodromy}
\label{sec:newV100-endpoint}

Suppose the physical data form a loop, so \(E_1=E_0\), \(F_1=F_0\), and
\(B_1=B_0\), while the chosen transports close with unitaries
\begin{equation}
 S_1=S_0g_E,\qquad R_1=R_0g_F.                             \label{eq:newV100-frame-return}
\end{equation}

\begin{corollary}[Exact frame-transition phase]
\label{cor:newV100-frame-phase}
Under \eqref{eq:newV100-frame-return},
\begin{equation}
 \exp\left\{\int_0^1
 \operatorname{Tr}(\widetilde B_t^{-1}
                    \dot{\widetilde B}_t)\,dt\right\}
 =
 \overline{\det g_F}\,\det g_E.                            \label{eq:newV100-frame-phase}
\end{equation}
\end{corollary}

\begin{proof}
Equation \eqref{eq:newV100-Btilde} gives
\[
 \widetilde B_1=g_F^*\widetilde B_0g_E.
\]
Take determinants and use \Cref{thm:newV100-Jacobi-transport}.
\end{proof}

\begin{assessment}[Relation to the mapping-torus gate]
\label{ass:newV100-mapping-torus}
Equation \eqref{eq:newV100-frame-phase} is the finite-dimensional
frame-transition factor of the determinant line.  It is not by itself the
complete chiral Standard Model phase: the chosen regulator and eta or
counterterm contribution must be combined with it.  V91's mapping-torus gate
asks that the complete combined holonomy be trivial for every gauge loop
after local anomaly cancellation.  V100 supplies the exact assembled
finite-cutoff transport identity, not that global trivialization.
\end{assessment}

\section{Why an operator floor does not control an extensive determinant}
\label{sec:newV100-extensive}

\begin{proposition}[The raw operator-norm estimate carries the matrix
dimension]
\label{prop:newV100-dimension}
If \(\widetilde B_t\) acts on an \(N_t\)-dimensional space, then
\begin{equation}
 \left|\operatorname{Tr}(\widetilde B_t^{-1}
                         \dot{\widetilde B}_t)\right|
 \le {N_t\over\kappa_t}\|\dot{\widetilde B}_t\|.            \label{eq:newV100-dimension}
\end{equation}
The factor \(N_t\) is sharp for scalar dilation
\(\widetilde B_t=b(t)I\).  Therefore the V95 singular floor and the V97
operator-norm chart do not alone give a cutoff-uniform determinant
transport as the lattice volume in sites diverges.
\end{proposition}

\begin{proof}
Use
\(\|X\|_1\le N_t\|X\|\) in
\eqref{eq:newV100-trace-bound}.  For
\(\widetilde B_t=b(t)I\),
\[
 \operatorname{Tr}(\widetilde B_t^{-1}\dot{\widetilde B}_t)
 =N_t{\dot b(t)\over b(t)},
\]
so the dimensional factor cannot be removed from a general operator-norm
bound.
\end{proof}

\begin{corollary}[Required continuum replacement]
\label{cor:newV100-continuum-replacement}
A cutoff-uniform determinant phase or relative modulus requires at least one
of the following additional structures:
\begin{enumerate}[label=\textup{(\roman*)},leftmargin=3.5em]
\item a uniform trace-norm bound on the assembled relative derivative;
\item an exact local counterterm subtraction after which the remainder has
      such a bound;
\item a localized connected expansion whose complete allocation is summable;
      or
\item a determinant-line holonomy theorem which evaluates the combined
      phase without taking an extensive absolute trace bound.
\end{enumerate}
A uniform least singular value without one of these rows is insufficient.
\end{corollary}

\begin{proof}
The sufficiency of the first row is
\Cref{thm:newV100-Jacobi-transport}; the other rows are mechanisms for
producing or replacing the same finite assembled quantity.  The
insufficiency statement is \Cref{prop:newV100-dimension}.
\end{proof}

\section{The no-escape card}
\label{sec:newV100-no-escape}

\begin{definition}[Assembled transport and no-escape card, ATNEC]
\label{def:newV100-ATNEC}
For a cutoff family, ATNEC consists of:
\begin{enumerate}[label=\textup{(\roman*)},leftmargin=3.5em]
\item the AOP and VRMC/HYFC protected chart with all unmatched massless rows
      retained;
\item one exact transported chiral block \(\widetilde B_a(t)\);
\item a renormalized or relative connection one-form whose integrated
      trace norm is bounded uniformly in \(a\) and physical volume on the
      intended local observable support;
\item a proof that complement probabilities and the corresponding
      determinant weights are uniformly integrable, so mass cannot escape
      through a growing number of bad cells; and
\item trivial combined determinant-line holonomy on every Standard Model
      gauge loop, including the eta or regulator contribution.
\end{enumerate}
\end{definition}

\begin{theorem}[Finite-cutoff consequence of ATNEC rows
\textup{(i)}--\textup{(iii)}]
\label{thm:newV100-ATNEC-local}
At a fixed cutoff, ATNEC rows \textup{(i)}--\textup{(iii)} give a
nonvanishing transported determinant on the massive block and a finite
path-independent modulus ratio whenever the renormalized connection is
exact.  They give a controlled phase along each chosen path with variation
bounded by \eqref{eq:newV100-L1}.  Path independence of the phase requires
the separate loop-holonomy row \textup{(v)}.
\end{theorem}

\begin{proof}
HYFC gives invertibility.  Apply
\Cref{thm:newV100-Jacobi-transport}.  Exactness of the real connection gives
path-independent modulus ratios.  Two phase paths differ by a loop, whose
factor is \eqref{eq:newV100-frame-phase} combined with the regulator
contribution; hence row \textup{(v)} is precisely the additional condition.
\end{proof}

\section{Rollover decision}
\label{sec:newV100-rollover}

The companion transfer changes the next acquisition order.  Chasing each
higher determinant derivative would create the open hierarchy seen in the
Navier--Stokes moment ledger, while bounding each local contribution before
allocation risks the Yang--Mills Bell loss.  The next active series will
therefore begin with a compact result map and attack the assembled relative
connection in \eqref{eq:newV100-L1}.  The first target is an exact
finite-cutoff subtraction formula for the vectorlike massive base versus
the chiral electroweak block, followed by a locality/trace-class estimate
for the remainder.  The global mapping-torus phase and bad-set
uniform-integrability rows remain explicit.

\begin{firewall}[Claim boundary after V100]
\label{fire:newV100-boundary}
V100 performs the compulsory settled companion audit, proves exact
determinant transport in moving chiral frames, identifies the endpoint
transition factor, proves the unavoidable dimension loss in a raw
operator-norm argument, and formulates ATNEC.

V100 does not prove a cutoff-uniform trace-norm bound, determinant
renormalization, bad-set uniform integrability, global determinant-line
trivialization, the mapping-torus gate, full reflection positivity,
cutoff removal, or the full SM--Einstein continuum.  The complete
second-series file is frozen after this note; the next active file restarts
at V1 with a summary and context.
\end{firewall}

\begin{theorem}[V100 result ledger]
\label{thm:newV100-results}
V100:
\begin{enumerate}[label=\textup{(V100.\arabic*)},leftmargin=7em]
\item records the settled YM V74 and NS V43 byte snapshots and rejects a
      transitional duplicate YM filename;
\item transfers their global-allocation and spatial-no-escape warnings;
\item proves the exponentiated Jacobi transport identity;
\item proves the trace-norm length bound;
\item computes the exact endpoint frame-transition phase;
\item proves the sharp matrix-dimension loss of raw operator-norm control;
      and
\item formulates ATNEC and its finite-cutoff consequence.
\end{enumerate}
\end{theorem}

\begin{proof}
Use \Cref{audit:newV100-transfer,
thm:newV100-Jacobi-transport,
cor:newV100-frame-phase,
prop:newV100-dimension,
cor:newV100-continuum-replacement,
def:newV100-ATNEC,
thm:newV100-ATNEC-local}.
\end{proof}

\paragraph{Parent-version record.}
V100 was formed from
\texttt{Renewal\_SM\_Einstein\_Continuum\_Research\_Notes\_V99.tex}.
The V99 parent had \(35\,122\) logical source lines,
\(1\,349\,056\) bytes, and SHA--256
\[
 \texttt{53CBAA0772DBA73AE19618614C0D47D98197436759668EF6743D26570C74B28E}.
\]

% =============================================================================
% END APPENDED SECOND-SERIES V100 MATERIAL
% =============================================================================
\end{document}
